% !TEX program = lualatex
% Continuous Chinese translation without facing-page English source plates.
\documentclass[UTF8,11pt,a4paper,openany,fontset=fandol]{ctexbook}
\usepackage[margin=2.35cm,headheight=25pt]{geometry}
\usepackage{amsmath,amssymb,bm,mathtools}
\usepackage{fancyhdr,microtype,xcolor,hyperref,xurl,needspace}
\setlength{\parindent}{2em}
\setlength{\parskip}{0.55em}
\setlength{\emergencystretch}{3em}
\linespread{1.22}\selectfont
\clubpenalty=10000
\widowpenalty=10000
\displaywidowpenalty=10000
\raggedbottom
\pagestyle{fancy}
\fancyhf{}
\fancyhead[C]{\small\nouppercase{\leftmark}}
\fancyfoot[C]{\thepage}
\renewcommand{\headrulewidth}{0.3pt}
\hypersetup{hypertexnames=false,colorlinks=true,linkcolor=black,urlcolor=blue!60!black,pdftitle={计算流体动力学前沿（2002）},pdfauthor={D. A. Caughey、M. M. Hafez 主编}}
\title{计算流体动力学前沿（2002）\\[0.45em]\Large 中文译本}
\author{D. A. Caughey、M. M. Hafez 主编\\\small 中文初译经 CFD 专业术语表统一}
\date{}
\begin{document}
\frontmatter
\maketitle
\chapter*{译本说明}
本译本按原书章节顺序翻译并连续排版，正文不附英文原文，供中文阅读和检索使用。作者姓名、文献题名、专业缩写以及数学符号按学术写作惯例保留必要的拉丁字符。

中文初稿由离线神经机器翻译批量生成，并通过 CFD 专业术语表统一 Navier-Stokes 方程、有限体积法、伴随方法、多重网格、预处理、湍流、磁流体动力学等词汇。由于全书篇幅大且原 PDF 文本层存在断行与公式混排，本文件属于可继续校订的学习译本，不替代正式授权译本。
\tableofcontents
\mainmatter
Ationa • s L nami editei L。 抓住M.M. 科学

2002年Q计算流体动力学前沿

由D.A.编辑的计算流体动力学的前沿。 考伊康奈尔大学理学硕士 加州大学哈菲兹分校，戴维斯Yfe世界科学wB新泽西•伦敦•新泽西•伦敦•新加坡•香港

由世界科学出版公司出版。 Pte。 有限公司 P O Box 128, Farrer Road, Singapore 912805 USA 办公室：Suite IB, 1060 Main Street, River Edge, NJ 07661 英国办公室：57 Shelton Street, Covent Garden, London WC2H 9HE 大英图书馆编目出版数据 本书的目录记录可从大英图书馆获得。 计算流体动力学 2002的前沿 版权所有 © 2002 世界科学出版公司。 Pte。 有限公司 保留所有权利。 未经出版商书面许可，不得以任何形式或任何方式（电子或机械）复制本书或其部分，包括复印、录音或任何已知或即将发明的信息存储和检索系统。 如需复印本卷材料，请通过版权清关中心支付复印费，地址为222 Rosewood Drive, Danvers, MA 01923, USA。 在这种情况下，不需要出版商的复印许可。 国际标准书号981-02-4849-0由大陆出版社在新加坡印刷

\chapter*{奉献}
奉献 本卷包括在纪念罗伯特·W的研讨会上发表的论文。 MacCormack，并表彰他三十多年来对计算流体动力学（CFD）领域的开创性贡献。 题为《计算未来III：计算流体动力学前沿》的研讨会于2000年6月26日至28日在加利福尼亚州半月湾举行。 作者是从国际知名的空气动力学和CFD研究人员中选出的，MacCormack的贡献影响非常重要。 作者和编辑很高兴将这本书献给Bob，以表彰他在我们的技术和生活中发挥的重要作用。 Bob MacCormack于1940年2月21日出生在纽约布鲁克林。 他在那里长大，在布鲁克林学院接受本科教育，主修物理和数学。 他于1961年加入美国宇航局艾姆斯研究中心，最初在高超音速自由飞行部门工作。 在艾姆斯时，他完成了M。 斯坦福大学数学理学学位，1971年，转为新成立的计算流体动力学分部的助理主任。 随后，他担任艾姆斯热和气体动力学部门的高级职员科学家，然后于1981年开始他的学术生涯，担任西雅图华盛顿大学航空航天系的教授。 1985年，他回到湾区，接受了斯坦福大学航空航天系的教授职位。 Bob曾在意大利、日本和（前）苏联以及美国的国际会议上发表过主题演讲。 他曾在布鲁塞尔的冯·卡曼研究所、印度特里凡得琅的维克拉姆·萨拉巴伊航天中心、台湾国立成空大学、北京的曲华大学、西安西北理工大学、四川的中国空气动力学研发中心以及中国空气动力学研发中心多次讲授CFD的短期课程。 他向十几个U提供建议和咨询。 S.航空航天公司和计算流体动力学的前沿-2002年编辑：David A. Caughey和Mohamed M. 哈菲兹 ©2002 世界科学

VI 奉献政府机构。 Bob对CFD的贡献得到了许多主要奖项的认可。 他获得了美国宇航局艾姆斯研究中心H。 1973年获得朱利安·艾伦奖，1981年获得美国宇航局杰出科学成就奖章，并于2001年被选中发表西奥多森讲座。 他在美国航空航天研究所（AIAA）有很长的服务历史，包括1977-79年担任流体动力学技术委员会的成员，并担任AIAA杂志的副编辑。 他于1988年当选为AIAA院士，并于1996年获得该协会的流体动力学奖。 1970-76年，他担任《计算物理学杂志》的副编辑，并在1982年和1985年以及1983-84年担任国家工程学院委员会的成员，评估CFD的现状和增长。 自1992年以来，他一直是美国国家工程学院的成员。 在本书的第一章中，将更详细地讨论Bob的技术贡献，特别是他们对高超音速空气动力学和CFD的影响。 几乎所有研讨会的与会者都讲述了他们第一次使用“MacCormack计划”的故事。 第二章重印了鲍勃介绍该计划的著名论文，第三章总结了他与美国宇航局艾姆斯CFD分部几位研究人员的互动。 其余章节介绍了该领域领先专家撰写的当前感兴趣的主题。 但Bob的贡献不仅限于他的技术理念、他的国家领导力、他教授的课程，或者他对华盛顿大学和斯坦福大学许多有才华的学生的监督。 Bob是最真实意义上的绅士，他的优雅和幽默丰富了所有在他非凡职业生涯中认识他的人。 鲍勃在研讨会宴会上拍摄的照片显示在正面页面上。 研讨会的第二天，一些与会者与鲍勃一起在加州海岸外的太平洋钓了一天的鲑鱼。 下一页显示了Bob站在新Captain Pete的船尾甲板上，等待下一条鲑鱼咬一口的照片。

5？ Jy Robert W. 麦科马克

4 \textbackslash\{\}£ Bob MacCormack在新皮特船长的甲板上。

\chapter*{内容}
内容奉献v 1 Robert W的贡献。 MacCormack到计算流体动力学 Caughey \& Hafez 1 1.1 导言 1 1.2 概述 2 1.3 技术贡献 3 1.4 结论性 17 参考文献 18 2 黏度对超高速冲击坑洞的影响 MacCormack 27 2.1 摘要 27 2.2 导言 27 2.3 数值方法 28 2.4 数值计算 37 2.5 结论性 40 参考文献 42 3 MacCormack方法 - 历史视角 Hung, Deiwert \& Inouye 45 3.1 导言 45 3.2 MacCormack方法的演变 46 3.3 应用 50 3.4 结束语 57 参考文献 58

X 目录 4 实现教科书多网效率的一般框架：一维欧拉示例 Thomas, Diskin, Brandt \& South 61 摘要 61 介绍 61 一般框架 63 准一维方程 65 放松方案 66 分布式放松 68 计算结果 71 超声速流动 74 结论性评论 75 参考文献 76 附录 77 I - 保守通量 77 II - 分布矩阵 78 III - 跨声速冲击 - ENO 微分 79 5 二维和三维流动的 Cauchy-Riemann 方程 数值解 Hafez \& Houseman 81 介绍 5.1 82 5.2 控制方程 和边界条件 83 5.3 数值方法 84 5.4 数值结果 85 5.5 结论性备注 86 5.6 附录：多重网格收敛 结果 86 参考资料 86 6 多维可压缩流的非均匀网格上的高效高阶方案 Lerat、Corre和Hanss 89 6.1 介绍 89 6.2 正笛卡尔网格上的欧拉求解器 90 6.3 不规则笛卡尔网格上的欧拉求解器 93 6.4 Navier-Stokes求解器 98 6.5 数值实验 100 6.6 结论 105 参考文献 105 7 计算可压缩流的未来方向：高阶居中与多维上绕 Napolitano等人 113

目录 xi 7.1 介绍 113 7.2 高阶中心数值方法 115 7.3 波动分裂 方法 116 7.4 结果和讨论 118 7.5 结论 124 7.6 致谢 125 参考资料 125 8 三维曲线移动网格的高效低耗散高阶方案的扩展 Vinokur \& Fee 129 8.1 导言 130 8.2 方程的制定 134 8.3 数值方法 143 8.4 结论性 155 致谢 156 附录 A：一类数值混合偏导数的交换性 156 附录 B：非平衡流动的 Riemann 求解器 160 参考资料 163 9 四阶方法 对于交错网格上的斯托克斯和Navier-Stokes 方程 Gustafsson \& Nilsson 165 9.1 介绍 165 9.2 稳定斯托克斯方程和交错网格 167 9.3 斯托克斯方程的四阶方法 171 9.4 Navier-Stokes 方程的四阶方法...... 175 参考 178 10 非稳定不可压缩流动的可扩展并行隐式多网格解决方案 Pankajakshan等人 181 10.1 摘要 181 10.2 介绍 182 10.3 基本非定常流动求解器 182 10.4 可扩展并行隐式算法 185 10.5 并行性能估计和可扩展性 188 10.6 演示：舵诱导机动模拟...... 193 10.7 致谢 195 参考 195

十二 内容 11 涡量约束在预测复杂体上流动的应用 Steinhoff 197 11.1 介绍 198 11.2 常规欧拉方法 199 11.3 涡量约束 200 11.4 当前结果 206 11.5 结论 213 参考 214 12 晶格玻尔兹曼不可压缩流动模拟 Satofuka \& Ishikura 227 12.1 介绍 227 12.2 格子 Boltzmann 方法 用于二维 228 12.3 二维均匀各向同性 湍流 231 12.4 突然膨胀的二维通道 233 12.5 格子 Boltzmann 方法 用于三维 235 12.6 三维 均质各向同性湍流 236 12.7 三维管道流 237 12.8 平行化 239 12.9 结论240 参考240 13 入口Agarwal和Deb 243中弱电离气体的MHD效应的数值模拟 243 13.1摘要243 13.2命名法244 13.3 介绍246 13.4 Electro-磁流体动力学的控制方程.... 247 13.5 控制方程在弱守恒律形式249 13.6 控制方程在广义坐标 252 13.7数值方法254 13.8重要引数259 13.9入口超声速流动的数值模拟260 13.10结论263 13.11致谢263参考263 14 计算磁流体空气动力学 Shang, Canupp和Gaitonde的进展273 14.1介绍273 14.2 控制方程 275

内容 xiii 14.3 数值程序 277 14.4 兰金-胡戈尼奥特跳跃条件 282 14.5 理想的MHD 激波管 模拟 285 14.6 高超音速MHD 钝体 模拟 289 14.7 结论 293 14.8 致谢 294 14.9 参考 294 15 复杂几何三维龙网格方法的开发 Liou \& Zheng 299 15.1 介绍 299 15.2 DRAGON网格 301 15.3 三维DRAGON 网格生成 303 15.4 流动求解器 308 15.5 测试案例 309 15.6 结论性备注 312 致谢 313 参考 314 16 多块、修补网格拓扑应用于Navier-Stokes 预测 陆军炮弹的空气动力学 Sturek \& Haroldsen 319 16.1 介绍 319 16.2 导弹配置 320 16.3 边界/初始条件 322 16.4 性能/收敛 标准 323 16.5 结果 323 16.6 结论性评论 324 16.7 致谢 324 参考 324 17 通过雷诺方程预测-平均Navier-Stokes 方程 大厅 333 17.1 介绍 333 17.2 RANS方案和门特湍流模型 336 17.3 门特湍流模型338 17.4 对门特湍流模型的RANS结果 341 17.5 结论 345 参考资料 346

Xiv 目录 18 计算空气动力学流动算法的进步 Zingg, De Rango \& Pueyo 347 18.1 导言 347 18.2 Newton-Krylov算法 349 18.3 高阶空间 离散化 356 18.4 结论性 366 致谢 367 参考 367 19 高超音速边界层 稳定性和接受性的数值模拟 Zhong, Whang \& Ma 381 19.1 导言 381 19.2 控制方程和数值方法 382 19.3 结果和讨论 383 19.4 结论性 395 参考资料 396 20 使用重叠网格方法在涡轮泵中对不可压缩流动的依赖模拟 399 20.1 介绍 399 20.2 数值方法 400 20.3 方法和计算模型 402 20.4 计算结果 406 20.5 摘要 413 20.6 致谢 414 参考 414 21 德尔塔翼涡流模拟方面 Rizzi, Gortz \& LeMoigne 415 21.1 介绍 415 21.2 计算方法 419 21.3 测试案例和网格 421 21.4 固定翼计算和结果 426 21.5 Pitching Delta 的初步结果 434 21.6 结论和展望 438 21.7 致谢 439 参考 439 22 DLR Kordulla 的选定CFD 能力 443

目录 xv 22.1 介绍 443 22.2 CFD 发展 444 22.3 最近应用 449 22.4 去哪里 454 22.5 致谢 455 参考 455 23 CFD 太空运输系统的应用 Fujii 459 23.1 介绍 459 23.2 数值方法 460 23.3 结果和讨论 460 23.4 结论 471 23.5 致谢 472 参考 472 24 使用进化算法对超音速机翼进行多点最佳设计 Obayashi, Takeguchi \& Sasaki 475 24.1 介绍 475 24.2 优化方法 476 24.3 当前优化问题的制定 477 24.4 超音速运输的优化 Wing 478 24.5 结论 480 参考资料 481 25 信息科学- CFD 大岛和大岛的新前沿 489 25.1 确定性系统进入复杂系统 489 25.2 计算机与人脑 490 25.3 信息科学 491 26 将CFD 整合到空气动力学教育 Murman \& Rizzi 493 26.1 介绍 493 26.2 从1981年到2000年的变化 494 26.3 教育考虑和问题 497 26.4 非正式调查结果 499 26.5 整合的例子 503 26.6 摘要 505 26.7 致谢 506 参考文献 506

\clearpage
\chapter{1 Robert W的贡献。 MacCormack到计算流体动力学}
\section{1.1 介绍}
Robert W的贡献。 MacCormack到计算流体动力学 David A. Caughey1和Mohamed M. Hafez2 1.1 介绍 Robert W. 自计算流体动力学（CFD）的起步阶段以来，MacCormack一直是计算流体动力学（CFD）开发的主要力量。 他为求解可压缩流体流方程的基本数值方法做出了重大和开创性的贡献，包括非平衡化学高速流，并将这些方法应用于重要的基本问题，包括激波边界层相互作用和压缩斜坡上的超音速流，以及更多应用问题，包括完整的航空航天飞行器的流。 大多数CFD研究人员熟悉Bob对明确的Lax-Wendroff方法的高效修改，但许多人不知道可以追溯到Bob论文的其他重要概念的数量。 这些包括有限体积法，使用第二和第四差分数值耗散，他的隐式方案，使用线松弛技术将可压缩方程迭代到稳态，以及他修改后的近似因式分解方案。 他还使用子迭代来消除（或最小化）分裂错误，他主张以类似于自然黏度的形式引入数值黏性，以保持1西比利机械和航空航天工程学院的框架独立性，康奈尔大学，伊萨卡，纽约14853-7501。 2 加州大学戴维斯分校机械与航空工程系，加利福尼亚州戴维斯，95616。 Frontiers of 计算流体动力学 - 2002 编辑：David A. Caughey和Mohamed M. 哈菲兹 ©2002 世界科学

\section{1.2 概述}
2 CAUGHEY \& HAFEZ Navier-Stokes 方程，即使在任意曲线坐标系中解决。 在他的论文中也可以找到许多其他想法。 Bob将他的方法应用于许多物理问题，包括高超音速层流和湍流流，以及热化学非平衡流，以及超音速和超音速流。 他研究了连续极限和各种形式的伯内特方程及其稳定。 最近，他研究了磁流体力学（1978年还研究了使用分裂和重新分区技术的非线性光学传播）。 本章的目的是更详细地描述这些技术贡献，并在此过程中提供CFD中重要发展的历史概述。 1.2 概述 Bob MacCormack对计算流体动力学领域的贡献可以通过计算他论文标题中各种单词的出现来看出。 本章末尾包含Bob的部分出版物列表。 表1总结了应用于此列表的此类单词搜索的结果。 最常见的单词是数字，出现在超过1/3的标题中，这也许不足为奇。 更有趣的是，计算一词只出现了4次，首字母缩写词CFD只出现了两次（一次是在对该领域25年进步的个人回顾中，自1993年在佛罗里达州奥兰多举行的第11届AIAA CFD会议上提出的“MacCormack方法”开发以来）。 鲍勃工作的实际重点表现在这样一个事实中：Navier-Stokes短语出现了22次，黏性一词出现了12次，湍流（或湍流）出现了11次，而层次只出现了4次。 此外，复合词三维出现7次，而二维只出现一次。 值得注意的是，inviscid一词只出现过一次，而Euler完全不存在；对于20世纪80年代积极开发可压缩流体流动问题的数值方法的研究人员来说，这是一个了不起的壮举。 最后，Bob的许多工作重点都说明了高超音速一词出现19次，而可压缩和休克（或激波）总共出现33次。 “超高速”一词只出现了三次，但值得注意的是，因为它出现在Bob最常被引用的论文《黏度对超高速撞击坑3的影响》中，该论文向世界介绍了MacCormack方案。 这篇经典论文作为本卷的第2章转载。

\section{1.3 技术贡献}
R的贡献。 W。 MACCORMACK 3字数值Navier-Stokes高超音速可压缩冲击（-波）黏性湍流（t/ce）边界层三维非平衡层计算超速度CFD非黏性二维尤拉标题数表1罗伯特·W.近100篇论文的标题中选定单词的出现。 麦科马克。 1.3 技术贡献 鲍勃在美国宇航局艾姆斯高超音速自由飞行分支的最早工作是在高速撞击坑洞领域。 他的报告[3]总结了旨在解决与超高速冲击相关的可见辐射闪光是否需要气体大气的争议的实验。 这些实验的动机是通过对弹丸对月球暗面表面的超高速撞击进行光谱分析来确定月球表面的化学成分。 Bob对可能导致闪光的各种机制所需的能量的分析表明，发光率的观测值对环境压力的依赖性与高速喷射物与大气层的相互作用是一致的，这是观测到的辐射的主要原因。 Bob在总结中指出，对于更高的撞击速度，或者对于他们测试中的铝和玄武岩以外的材料，火山口机制本身可能足以产生可观测的辐射，尽管编辑们不知道月球实验是否曾经进行过。

4 CAUGHEY和HAFEZ 对物理机制的不确定性，该物理机制负责观测到穿透深度与弹丸直径的比率以及最终目标动量与初始弹丸动量的比率的缩放，促使Bob分析了黏度对超速冲击现象的影响。 当然，即使对于轴对称的情况，也有必要用数值解决问题，鲍勃对这个问题的数值解需要一个高效、二阶准确的方案，因此他开发了Lax-Wendroff方案的交替后空间预测器正向空间校正器版本，自[5]以来，该方案一直是他的名字的代名词。 在论文中，Bob分析了方案的稳定性，并推测在向后/向前步骤的应用顺序之间交替将允许对时间步骤进行全（一维）CFL限制。 这一猜想当然在该方法的多次应用中得到了一次又一次的验证，但尚未得到证实（至少在编辑们的知识中是这样）。 Bob在效率方面描述了他的方法的权衡，认为只有能够达到完整的CFL限制，它才会有效。 在包含32 x 33网格单元的网格上，解决这个历史性问题，在IBM 7094上需要大约15分钟的CPU时间。 Bob使用Sakharov对铝黏度的值，计算了圆柱形铝弹丸（I = d）对铝靶子（包括半无限和薄板）的影响。 结果表明，总轴向和径向动量“......表现出与......测量结果一致的效果”，鲍勃继续建议“可以证实黏度在超速度影响中的重要性”的具体实验。 这项研究完成后，鲍勃显然认为流体力学包含足够多的未解决的问题，可以成为一名职业。 大约在同一时间，迪安·查普曼看到了流体力学和空气动力学中数值方法发展的潜力，并在美国宇航局艾姆斯成立了CFD分部，鲍勃被邀请加入。 他的下一篇论文将他的新数值方案应用于平板上斜激波与层边界层相互作用的问题。 为了解决涉及边界层的问题，他开发了一个空间分割版本的方案，具有两个优点[6]。 首先，由于每个一维运算符都保持稳定，直到完整的CFL极限，因此之前关于这方面的稳定性的问题变得有问题。 其次，拆分允许在垂直于边界的方向（即边界层的「正向」方向）上执行多个时间步，每个时间步都与边界平行的方向。 计算是在多块网格上进行的，以便在墙附近实现更精细的分辨率；网格分别包含34 x 32和34 x 22网格点，需要256A£X时间步骤的典型计算在IBM 360/67上需要约4小时的CPU时间。 与广泛使用的实验测量相比，非分离和分离流动情况的表面压力和皮肤摩擦结果都相当好

R的贡献。 W。 Hakkinen et a/.的MACCORMACK 5，特别是当网格在分离流案的再循环区中精炼时。 有趣的是，应用于CFD的“计算机图形”的第一批演示之一出现在本文中：由阴极射线显示管生成的流线图叠加在边界层速度曲线上的复制。 Bob在一篇带有A的论文中成为空间分割的支持者。 J。 Paullay [7]，在1972年1月的AIAA航空航天科学会议上发表演讲。 作者在这里提出了数值方法的准确性和效率之间的正相关性，指出以最大允许时间步数运行的显式方案具有推进解决方案所需的所有数据，并且有最少的无关数据。 本文的目的是展示早期用于激波/层边界层相互作用问题的运算符分割到更一般的流体动力学问题的优势。 分割显式MacCormack方案应用于可压缩流动的非黏性方程，以解决超声速流动过去对称菱形翼型和使用简单、非正交、剪切网的双压缩角。 它们取得了与这些问题的精确（非黏性）解决方案非常一致的结果，表明与未分割方法相比，计算时间减少了两倍以上。 拆分方法允许1）在每个空间维度的全一维CFL极限下推进解，2）在较粗的方向上推进每个时间步的较小网格间距的多个时间步进解，允许更好地匹配依赖的数值和物理域。 在随后的一篇论文[13]中，MacCormack和Paullay讨论了计算网格对解决方案准确性的影响，介绍了“网格拟合”的概念，用于准确处理激波。 他们向CFD社区引入了弱解的概念，并引入了时间分割、显式MacCormack方法的有限体积形式--这是第一次在文献中发现这种现在运动方程的近似标准类。 作者用三个问题来说明三个不同的要点。 作者认为，由于时空网格与该值Courant 数的解对齐，线性波（对流）方程用于表明MacCormack 显式方法再现了统一Courant 数的精确解。 其次，无黏性Burgers 方程用于表明，如果没有纠正措施，数值方案可能会捕获（物理不正确的）膨胀“休克”。 他们为这个问题提供了两种补救措施；1）对通量计算进行简单的修改，以确保膨胀区域速度的连续性，2）添加一个第四差分耗散项（Jameson、Schmidt和Turkel广泛使用的混合第二和第四差分耗散中的一个元素）。 作者指出，制造这些是危险的

6 CAUGHEY和HAFEZ在解决方案不连续的区域进行修改，因为额外的截断误差可能很大；这与Jameson、Schmidt和Turkel的策略一致，即使用非线性开关关闭激波附近的第四个差分耗散。 应该注意的是，MacCormack和Paullay的策略与Jameson、Schmidt和Turkel的策略之间的一个重大区别：前者建议仅在需要避免扩张冲击的地方添加第四个差异项，而后者建议在除近不连续性外的任何地方添加它们。 最后，作者考虑了Euler 方程对几个二维超声速流动的解决方案，包括通过楔形、钻石翼型和球体的流动。 对于这些流程，结果表明，当网格与激波位置对齐时，数值误差会减少。 当先验的激波位置未知时，这需要一个解决方案自适应程序--即对于球体的情况。 使用Rankine-Hugoniot条件采用网格位置校正方案，但应该强调的是，冲击跳跃关系仍然由数值方案捕获（而不是拟合）。 对于钻石翼型的例子，使用第四差分耗散来避免膨胀冲击，否则会从机身的膨胀角发出。 鲍勃接下来将注意力转向了与巴雷特·鲍德温的一系列论文中斜电击与湍流边界层相互作用的更困难的案例[10，11，14]。 计算了在M\textasciicircum{} = 8.47处通过平板的流量，在与雷诺数 Rex = 22.5 x 106相对应的点上，强度Sp/p\textasciicircum{} = 83的斜冲击冲击边界层。 使用了MacCormack的显式方案的空间分割版本，现在包含四个连续更精细网格区域的网格，最精细的网格紧挨着墙壁。 板附近的细网允许平面法线因子在流向因子的每个时间步骤中前进96个时间步。 这项工作引入了用一个与\$xxP、OxxW fJ-xxP成比例的项来增加数值黏性的想法，其中p是流体压力，w是守恒变量的向量，5XX和fixx分别是3点微分和平均运算符。 作者发现有必要添加这种额外的耗散，以稳定具有如此强激波的案例方案。 当然，这个术语后来将成为詹姆森、施密特、图尔克尔混合第2/4差分自适应耗散的一个重要因素。 作者还引入了特殊处理，以在黏性子层中实现指数精度，在黏性子层中，湍流动能和耗散率变化数级。 计算的皮肤摩擦和沿板的传热分布与实验进行了比较；一致性是公平的-不是

R的贡献。 W。 MACCORMACK 7几乎和之前在层压案例中取得的一样好（令人惊讶！）。 Bob还与Art Rizzi合作，开发了他在广义坐标中超音速流的空间行进方法[12]。 本文代表了广义坐标中非黏性运动方程的最早呈现之一，描述了反变速度的通量。 MacCormack的两步、维度分割的显式方案在空间上进行超声速流动，在身体贴合的网格上也“适合”到激波上，以减少那里的振荡（另见[13]）。 提出了Moo = 14.9流过氦中钝锥体和MQO = 21.7流过空气中光滑三维体的计算结果。 MacCormack在第19届AIAA航空航天科学会议上发表的论文描述了隐式方法的重大进步[37]。 这篇论文还因包含方程而引人名，该方程也许最能描述鲍勃对CFD的方法：\{NUMERICS\} 5U? +l = \{PHYSICS\} (1.1) 换句话说，驱动解决方案更新的方程的右侧应该是控制问题物理学的方程的准确局部近似值，而左侧的责任是以稳定的方式在全球范围内传播局部确定的解决方案变化，以允许快速收敛的解决方案。 对于Navier-Stokes 方程，以Eq的右侧的8G n的紧凑向量形式9U书写。 （1.2）变成AF AP\textbackslash\{\}" A\textasciicircum{}+AST），。 （L3）对于稳态问题，方程中的\{NUMERICS\}。 （1.2）可以解释为预处理算子，而Eq的准线性形式的微分。 （1.2）对于一般时间依赖问题，给出8\{d\textbackslash\{\}J/dt）dA\{d\textbackslash\{\}J/dt）d\&\{d\textbackslash\{\}J/dt）dt dx dy（1.4），其中A-dF/dXJ和B=9G/9U分别是通量向量F和G的雅各布。 该方程描述了解决方案中At(dU/dt)的变化应该如何在整个域中传播。 等式的隐式近似。 （1.4）产量，

8 CAUGHEY \& HAFEZ，该方程分子中的点表明x和y的偏导数也作用于校正JU"\textasciicircum{}1。 在Eq中。 （1.5）假设At独立于x和y。 这些考虑因素表明，MacCormack预测校正器方案的有效实现可以在8月71（1.6）编写，其中A+和A\_分别是适当坐标方向上的两点正向和向后差。 矩阵|A|和|B|是具有非负特征值的矩阵，从相应的雅各布矩阵中计算出来，只有当违反局部（显式）CFL条件时，它们才是非零。 因此，由于等式的形式。 （1.6）该方案的隐含性被纳入近似的LU因数分解，该方案可以在空间上进行；由于|A|和|B|的构建，当满足本地CFL条件时，该方案简化为MacCormack的原始显式预测校正器-校正器Lax-Wendroff方案，当不需要它们时，不会在局部块反转上浪费精力。4 3 x 107的雷诺数的湍流边界层/冲击交互的结果与早期的显式方案的结果几乎相同，所需的CPU时间减少了1000倍以上。 该方案仅比MacCormack早期的显式隐式特征方案[19、22、23]略有效率，但更容易实现。 Kordulla和MacCormack[39]将隐式-显式LU因子方案应用于预测通过机翼的跨声速流动。 应用了显式-隐式、预测器-校正器方案的有限体积形式V Ax Ay J uii + svff1 AUgT请注意，如果没有在每个步骤中按维度划分的维度，上述安排变得类似于稳态计算的点隐式对称高斯-塞德尔迭代。

R的贡献。 W。 MACCORMACK 9在车身安装网格上解决雷诺平均 Navier-Stokes 方程。 对基本数值方案进行了几项修改。 首先，在最终解决方案中尽可能多地保留显式贡献是有利的，并引入了基于CFL的隐式运算符加权，以便在显式运算符和隐式运算符之间提供尽可能平滑的混合。 其次，发现有必要为更复杂的翼型问题增加更多的耗散；这种耗散被描述为“......三阶小，扫荡方向的导数作为系数。” 第三，修改了实心壁的边界条件，以立即抵消那里的通量（而不是像原始方法中建议的那样将其延续到校正器步骤）。 与传统的完全隐式（ADI）方法相比，显式-隐式、预测器-校正器方案的优点是：1）只需要双对角因子的解；2）只需要使用（修改的）欧拉雅各布，以及3）当满足局部CFL条件时，恢复到显式预测器-校正器方案。 在0.30 < M\textasciicircum{} < 0.73范围内的马赫数和4 x 106 < Rec < 6.5 x 106范围内的雷诺数的三个不同机翼的计算与其他计算以及使用足够精细的网格（210 x 60个单元格的顺序）时，与实验结果有很好的一致性。 Gupta、Gnoffo和MacCormack [42，50]将新的explicit-隐式方法应用于钝锥体上的黏性冲击层。 船头减震器再次安装，并在由此产生的车身坐标系中开发了隐式运算符。 结果表明，这些结果对Courant 数相对不敏感，表明隐式方法作为收敛加速技术解决稳定流量问题的益处。 Kneile和MacCormack [45]将显式隐式方法应用于Navier-Stokes 方程，用于三维内部流动。 这项工作展示了开发隐式技术的好处，该技术可以作为现有显式代码的“附加”实现。 鲍勃的显式欧拉代码的一个版本（相对）很容易被修改，以包括Navier-Stokes 方程的黏性项和双对角线隐式算法。 提出了几个测试案例的结果，包括三维收敛发散喷管流动。 Bob研究了使用多网格来加速Navier-Stokes 方程解决方案的收敛，以实现[47]中的稳定流动。 该方法基于他早期的显式-隐式算法[37]，以有限体积形式应用。 多电网的实施基于Ni方案（由Johnson为Navier-Stokes 方程实施）。 该方法应用于层压冲击诱导的分离问题，导致收敛速率加速仅为3（与预期系数8 1/2相比）。 结果确实表明，多重网格法能够大大加快双曲信号传播的速度

10个CAUGHEY和HAFEZ问题，但要实现多电网的全部潜力，必须等待进一步的发展。 在[46]中，MacCormack将Chakravarthy用于Euler 方程的通量分割、类型依赖差分方案开发的高斯-塞德尔线松弛方法扩展到Navier-Stokes 方程。 注意到Dean R做出的估计。 查普曼（MacCormack [48]）回顾了与安装完整飞机配置的计算网格相关的问题，查普曼（MacCormack [48]）审查了与安装完整飞机配置的计算网格相关的问题。5 在这篇调查文章中，鲍勃认为，雷诺平均Navier-Stokes 方程将很快将以与当前（1984年）的高雷诺数计算相同的准确度解决，通过相对简单的配置，如二维机翼和旋转体。 本文提供了通过完整的飞机配置实现计算黏性流动目标所需的技术的完整配方，包括（多块、结构化）网格生成和求解雷诺-平均Navier-Stokes（RANS）方程的算法，包括通量-向量拆分、有限体积公式、隐式算法和多网格的发展。 在影响深远的调查论文[49]中，鲍勃超越了对过去的工作的总结，还强调了他的高斯-塞德尔隐含方案。 在巧妙地将重要思想置于历史背景下后，他重复了Eq中表达的哲学。 （1.2）：\{NUMERICS\} SUfj-1 = \{PHYSICS\} u\textasciicircum{}1 = uTj + supj-1作为该方案发展的动机。 他提出了几个示例计算，包括超声速流动经过球形瓽子锥体和收敛-发散喷嘴中的跨声速流动。 他表明，Navier-Stokes 方程的充分收敛结果可以在大约10次迭代中获得，但指出，结果只有一阶精度，重要的工作仍然是实现具有更高阶精度的可比迭代效率。 在[55]中，MacCormack、Chapman和Gogken为Navier-Stokes 方程引入了新的滑移边界条件，该滑移在小的Knudsen数下减少到Maxwell的滑移边界条件，并在自由分子流动的极限情况下产生正确的剪切应力。 为二维高超声速流动通过平板计算的皮肤摩擦和传热率的比较，与实验结果和整个过渡流动制度的直接模拟蒙特卡洛计算结果相比出奇的好，5 有趣的是，这些计算资源目前在高端笔记本电脑上可用。

R的贡献。 W。 当使用新的边界条件时，MACCORMACK 11从连续体到自由分子流动。 在[51]中，Viegas、MacCormack和Rubesin讨论了循环控制机翼后缘区域湍流流量的预测。 提出了使用两个代数涡黏性模型的计算结果。 在与他的学生Graham Candler [54]的第一篇论文中，Bob扩展了他的高斯-塞德尔方法[49]来处理超过三维配置的高超音速流。 该方法是完全隐含的，使用高斯-塞德尔线松弛，使用通量依赖微分，并使用激波拟合作为弓波。 然而，该方法是完全保守的，允许捕获嵌入式交叉流冲击。 雷诺平均 Navier-Stokes 方程被解决，使用Baldwin-Lomax 湍流模型关闭系统。 为Shang和Scherr之前计算的Bionic体和X-24C-10D升降体提出了解决方案。为了将他们的计算适应Cray X-MP 48的内存限制内，MacCormack和Candler使用了更粗的网格（在子午线方向上是2的系数），但与早期的解决方案保持了良好的一致性，与显式方法相比，计算时间减少了两个数。 在[56、57、58、63、65]中，MacCormack和Candler开发并提出了解决高超声速流动问题的方法，包括有限速率化学和热非平衡的影响。 这种nowfleds由物种守恒、质量、质量平均动量、每个双原子物种的振动能、电子能和总质量平均能量的耦合、时间依赖、偏微分方程来描述。 解程序是完全隐含的，将流体流动方程与气体物理和化学关系耦合。 欧拉通量使用通量分裂近似，而黏性项是中心差分的。 该方法保留了强化学源项中的元素，并使用高斯-塞德尔线松弛求解方程。 该方法只需要几百个时间步骤就能解决经过简单身体形状的轴对称流，并预计扩展到更复杂的二维身体几何是直接的。 在[72]中，该方法被扩展到包括弱电离流动的电子数密度。 为高超声速流动计算的经过球形钝锥体的电子密度与在一系列高度的飞行测量结果非常一致。 在[59]和[62] Viegas中，Rubesin和MacCormack描述了他们在风洞测试部分解决流过循环控制翼型的流动的计算机代码。 在引入代码验证的想法后，使用Baldwin-Lomax和Jones-Launder 湍流 B的变体计算结果。该计算在里诺第23届航空航天科学会议上在AIAA论文85-1509中提出，被广泛认为是Navier-Stokes方程的第一个解决方案，用于完整的航空航天飞行器配置。

将12个CAUGHEY和HAFEZ模型与低亚声速流动的大量实验数据进行了比较，该模型经过二维循环控制翼型。 添加到湍流模型的变体包括一种解释喷气式飞机发展历史和溪流曲率影响的方法。 在[60]Candler和MacCormack总结了斯坦福大学的超音速研究，强调了辐射、反应和热激发流动的数值模拟、高超音速激波物理学的调查和数值模拟、连续流体动力学方程扩展到连续体和自由分子流动之间的过渡体系，以及开发用于稀疏流场的高效颗粒模拟的新型数值算法。 Bob鼓励他的学生Candler成为这篇概述论文的第一作者，这是Bob指导学生（和其他年轻研究人员）的能力的衡量标准。 在[61] MacCormack中回顾了构建三维流的高效算法的困难。 对一些候选人进行了分析和测试，发现大多数候选人有缺点。 然而，Bob得出结论，有希望在严格时间步幅限制的显式或近似因子算法和需要大稀疏矩阵的计算密集型直接反演算法之间可以找到一类高效的算法。 他在本文中用大部分文字和方程来说明因数算法如何不一定遵循旧的观点，即“扩展到三维是直接的”。 尽管如此，他提供了一个高斯-塞德尔算法，该算法在大约50次迭代中收敛到三维超音速级联问题（诚然是涡轮喷嘴，而不是压气机叶片通道）的解。 在[64]中，MacCormack和Gogken描述了一种用于高超音速、过渡性空气流动的热化学非平衡配方。 假定空气有五种化学物种（JV2、O2、NO、N和O），以及对应于能量的平移、旋转和振动模式的三种温度。 引入了速度和温度的滑动边界条件，以扩展低密度流动连续配方的有效性。 稀疏的高超声速流动经过5度、稀疏的锥体的解与DSMC结果进行比较，以表明连续体结果仍然有效的过渡Knudsen数的范围。 在[67]中，Bob讨论了计算流体动力学对流体流动装置设计的影响。 他回顾了他解决三维Navier-Stokes 方程的高效数值程序，基于两个方向的块状三对角线反转和第三个方向的高斯-塞德尔松弛[61]，提出了高超音速（Mach 20）流过一个有翅膀的再入车辆的结果，在一个廉价的桌面工作站上计算。 在[68]和[78]中，MacCormack和Wilson提出了二维轴对称流的完全隐式有限体积算法的耦合

R的贡献。 W。 MACCORMACK 13详细介绍了氢气反应机制（由13个物种中的33个反应表示），以研究超音速燃烧现象。 他们将数值计算的结果与弹道范围阴影图进行了比较，这些阴影图在钝体通过氢气和空气的预混合化学计量混合物时表现出两种不连续性。 他们讨论了模拟这些双锋现象的数值程序的适用性，并检查了这些流场对关键反应速率的敏感性。 在[69][81]中，MacCormack和他的学生Zhong使用线性化稳定性分析开发了一组（稳定）增强伯内特方程。 新方程适用于一维激波结构和经过钝前缘的二维流动。 对常规方程组和增强方程组的稳定性进行了数值测试，证实了增强方程组始终保持稳定，并保持与常规集相同的准确性。 他们表明，在高海拔地区，伯内特方程的解和Navier-Stokes 方程的解之间的差异是显著的，特别是对于对流场细节敏感的参数，如辐射。 在[74]中，Zhong和MacCormack评估了增强伯内特方程的表面滑移边界条件的一些模型。 在[70]和[76]中，MacCormack和Conti将Bob的Navier-Stokes 方程 [49]的隐式数值方法与碳材料的材料响应技术相结合，为耦合流动材料问题提供了二维瞬态解决方案。 计算了车辆表面温度和隔热屏蔽消融率，并考虑了由此产生的车辆形状变化。 测试计算结果显示了典型的弹道再入飞行器，其高度范围从43公里到海平面。 此外，在[77] MacCormack和Conti中，将MacCormack的隐式方法应用于带有气体喷射的层压、轴对称近尾迹问题。 通过将冷却惰性气体瞬态注入两个不同位置的注入端口的平衡空气来计算在22马赫时通过球形通过7度锥体的流量。 在[66，75]中，MacCormack和Rostand将完全隐式技术应用于热力学非平衡中氮等离子体的模拟。 这需要整合最先进的物理模型，以及MacCormack和Candler的数值技术。 结果与电弧加热氮等离子体射流进行了比较，总体上有很好的一致性。 在[71] Moreau、Chapman和MacCormack中，提出了一种用于轴对称流动的完全隐式有限体积算法，包括完整的热和化学非平衡以及高阶简化的伯内特应力张量，加上改进的详细非平衡辐射代码。 低速弓形冲击紫外线飞行实验用于基准测试稀疏建模对高海拔辐射的影响。 他们表明，包括旋转非平衡和简化的伯内特项并不能改善低速测试的趋势，但确实使

14 CAUGHEY和HAFEZ在更高速度下的差异。 结果表明，结果对辐射建模非常敏感，对于获得准确的辐射结果，最大振动温度和氮氧化物浓度比最大平移温度更关键。 Menon和MacCormack [73]将隐式高斯-塞德尔线松弛求解器应用于向后台阶下游区域空气和氦的超音速混合问题。 与同一案例的实验结果一致是公平的，毫无疑问，被使用的相当简单的（代数）湍流模型显著退化。 Moreau、Laux、Chapman和MacCormack [79]根据两项实验的结果描述了NEQAIR计算机代码的改进：在斯坦福大学进行的等离子炬实验和SDIO/IST弓形冲击超紫外导弹飞行的测量结果。 计算机代码也被扩展为处理气体中任意数量的物种和辐射带，其热力学状态可以通过多达四个温度来描述。 它为计算非常精细的光谱提供了更高的效率，并包括沿视线的运输现象。 Moreau、Chapman和MacCormack [80]开发了一种准一维通量分割、有限体积的计算机代码，包括额外的旋转松弛和单独的振动模式。 在Sharma和Gillespie进行的辐射实验中，该代码用于计算激波。 结果表明，常用的帕克旋转模型不足以模拟峰值放疗下观察到的旋转温度，对帕克模型进行校正，以解释松弛过程的扩散性质，能够恢复较大的初始差异。 Comeau、Chapman和MacCormack [82]研究了事故激波撞击钝体时产生的冲击相互作用，例如高超音速汽车的发动机入口罩唇。 Steger和Warming的通量向量分割方案用于在从连续条件到过渡流动条件的高度求解Navier-Stokes 方程的完美气体。 作者表明，随着流体密度的降低，相互作用变得根本不同，其对过热问题的影响也相应减弱。 他们发现，只有当入射冲击完全错过罩唇时，才能达到最高高度的最大驻点加热，并且与确实发生的与罩弓冲击的任何相互作用都发生在下游（因此，对驻点的条件影响不大）。 在[83]中，Welder、Chapman和MacCormack研究了Burnett方程的替代形式，其中材料导数的非黏性、等𤆡近似，存在于方程的原始形式的黏性应力和热传导表达式中，被精确材料导数所取代，并使用基于的改进近似值

R的贡献。 W。 MACCORMACK 15 Navier-Stokes（而不是inviscid）方程。 研究了各种伯内特方程，以确定它们对小波长扰动的稳定性，以及一维冲击结构的数值精度。 人们发现，不对材料衍生物进行一些近似的配方会导致非物理热传导。 开发了两种修改后的配方，大大减少了这个问题，同时提高了冲击结构计算的准确性。 在[84]中，Bob以非常热情和个人的方式总结了四分之一世纪的CFD研究。 他讨论了跨声速势流计算的发展，包括小扰动理论的发展，以及尤拉和Navier-Stokes 方程的数值方法，并一路上提供了许多令人愉快的个人轶事。 他还对未来在计算网格、计算机架构、算法和湍流研究方面的工作进行了预测。 在[85]中，鲍勃回顾了CFD二十年研究的进展，并为该领域完全成熟必须解决的问题指了方向。 Bob预测，未来的决策将涉及结构化、多块网格与非结构化网格，湍流的建模与湍流现象的直接模拟，以及间接放松（或近似分解）方案与直接解程序。 回想起来，这些是高度准确的预测，因为这些战斗中的大多数至今仍在进行。 在[86] Comeaux、Chapman和MacCormack中，从两个角度看待伯内特方程的熵平衡关系：从使用吉布斯方程和质量、动量和能量的连续守恒关系的经典热力学理论；以及从使用玻尔兹曼H定理与查普曼-恩斯科格展开的动力学理论。 他们发现，在这两种情况下，不可逆熵的产生都不是正半确定的，违反了热力学第二定律。 他们还表明，这两种公式是完全等价的（与克努森数中的二阶），表明吉布斯方程与伯内特方程一致（与早期研究人员的结果相矛盾，他们没有将推导到其高潮）。 与第二定律的不一致被认为是过去五十年来试图解决伯内特方程的研究人员所经历的数值问题的根源。 在[87] Kao中，von Ellenrieder、MacCormack和Bershader用实验和数值研究了二维可压缩涡旋与激波的相互作用。 非稳定Navier-Stokes 方程使用二阶精确的激波捕捉总变异递减（TVD）方案求解，计算结果与物理实验在定性上具有良好的一致性。 在[88]莫罗、查普曼和麦克科马克提出了一个新的温度

16 ZEldovich反应速率的CAUGHEY和HAFEZ依赖表达式解释了观察到的因振动激发而延迟形成NO。 新的交换反应模型自然地扩展了Park的解离反应的多温度反应速率模型。 对能量交换机制的详细分析也强调了扩展的Schwartz-Slawsky-Herzfeld（SSH）模型（与更简单的模型相比）更好的物理行为。 这些简单的修改有望大大提高用于研究低密度气体流动的最先进的详细热化学非平衡代码的预测能力。 在[89]中，Melville和MacCormack提出了一种为常微分方程设计最佳积分方案的方法。 线性分析用于构建广义的两步、预测-校正方法，然后针对双曲和抛物线系统对其进行优化。 在这两种情况下，计算效率都比以前的（标准）方案提高了，准确性没有显著损失。 不需要额外的内存，但需要初始化。 Melville和MacCormack [90]使用三维可压缩Euler 方程来研究涡流分解的双螺旋模式的不稳定行为。 纵向涡流通过负压梯度的对流表明，非定常流动场由单个空间均匀的频率主导，与螺旋涡流结构的旋转有关。 在[91]中，Bob使用高效的矩阵分解来构建隐式算法。 他首先分析了求解隐式矩阵方程的三种策略。 通过强隐式程序\{SIP）的近似LU分解，其中LU矩阵通过向下的对角线向前消除反转，然后向后倒置[/-矩阵的对角线。 在正向消除过程中，计算并存储二维流大小为4 x 4的2N矩阵元素（三维流为5 x 5的3N矩阵元素），其中N是网格点的数量。 近似分解在计算中引入了不对称性，可以通过在交替时间步骤上重新排序矩阵方程（或通过平均原始不对称算子）来最小化。 高斯塞德尔线松弛（GSLR）是另一种引入首选方向的策略，通常通过边界层与块状三对角线反转。 因此，它不适合包含交叉边界层的角落的域，尽管它通常非常高效。 第三个策略是微分方程的近似因数分解（AF）。 GSLR和AF都具有在二维或三维中反转与“一次一条线”相关的矩阵的优势。 （另一方面，近似LU分解或SIP可用于完全非结构化的网格）。

\section{1.4 结束的讲话}
R的贡献。 W。 MACCORMACK 17 经过此分析后，Bob通过在L和U矩阵中添加额外的元素来修改SIP，以消除计算中的不对称性，这样可以返回原始矩阵的完全相同的元素，但矩阵现在被分解为三个因子，其中一个是块对角矩阵（在三维中，有五个因子有两个块对角矩阵）。 这种修改后的SIP与特殊的近似因数分解（AF）有关，其中只有与“一次一条线”相关的矩阵是反转的。7 Bob还建议迭代拆分误差（两次或三次）来改进近似值，见[92]。 在[92、93、97]中，鲍勃介绍了他最新的隐式方案，用于解决不稳定的欧拉或Navier-Stokes 方程。 该方法基于使用（有限）迭代，以大幅降低基于近似因数分解的隐式方案的因式分解误差，并允许解决方案的收敛在三个测试问题（包括钝体后的超声速流动，以及跨声速和亚音速流通过喷嘴）的“工程精度”内达到约50-100个时间步程内的“工程精度”。 新方法与MacCormack和Pulliam [94]的标准近似因数分解方案进行了比较，表明新程序的效率提高了约五倍。 Pulliam、MacCormack和Venkateswaran在[96]中检查了许多隐式近似方案的收敛特征，包括DDADI方案，他们还展示了分级的好处，并得出结论，ADI和D3ADI的性能同样好，D3ADI可能更稳健。 有两种方法可以提高近似分解方案的性能，要么循环参数，要么循环网格。 在[94]中，MacCormack和Pulliam使用新的修改近似因式分解，结合了两个分级（AF2）与多网格相结合，并取得了令人印象深刻的结果。 在[98, 99]中，MacCormack展示了如何修改磁流体力学方程，使通量向量均匀一度。 这允许以守恒形式求解，并允许使用改良的Steger-Warming通量向量分割。 1.4 总结 前面关于他的贡献的摘要清楚地表明了Bob MacCormack对计算流体动力学做出的许多原创贡献，以及他对CFD的开发及其在工程中实际问题的应用产生了巨大影响。 7 另一种方法是交替方向对称高斯·塞德尔线松弛或交替方向斑马松弛，每个方向都有交替的奇数和偶数线。

\section{参考资料}
18 CAUGHEY \& HAFEZ Bob是第一批使用定向分割（x和y的向后差，然后是两个方向的正向差）、维度分割或维度分解和物理分割（对流，或双曲和扩散，或抛物线）的人之一，正如他在本卷美国宇航局的同事（Hung、Deiwert和Inouye）在论文中讨论的那样。 Bob和很多人一起工作过，例如，包括H。 洛马克斯，R。 变暖，B。 鲍德温，A。 Paullay，M. 伊努耶，G。 Deiwert，C. 洪，A。 里兹，J。 维加斯，M。 Rubesin和T。 美国宇航局艾姆斯的普利亚姆；与J。 Shang在Wright Field的空军研究实验室；与W。 Kordulla在Ames担任NRC高级研究助理时。 在斯坦福大学，他和D一起工作。 查普曼和D。 Bershader，他是许多学生的顾问，包括T。 戈肯，G。 坎德勒，X。 钟，S。 莫罗，K。 科莫，R. 梅尔维尔，G。 Wilson，C。 劳克斯，W。 焊工，P。 Bourqin，C. Kao、K. von Ellenreider和其他许多人。 他经常访问美国宇航局兰利的科学与工程计算机应用研究所（ICASE），并于2001年在那里发表了Theodorsen讲座。 编辑们很高兴能够将为本卷做出贡献的研究人员聚集在一起，向Bob表示感谢，并提供他的技术贡献摘要。 如前所述，这些技术贡献令人印象深刻，但它们只是Bob影响的一个方面；他的个人存在和精力，以及他帮助他人的意愿，特别是年轻研究人员，尤其值得注意。 他受到美国和国外这个社区的每个人的尊重。 我们祝愿Bob在未来的许多年里继续取得成功。 参考资料1。 麦科马克，R. W.，低环境压力下冲击闪光的调查，Proc. 第六届超高速冲击研讨会，1963年4月30日至5月2日，俄亥俄州克利夫兰。 2. 摩尔，H。 J.，MacCormack，R. W.，和Gault，D.. E.，金属和岩石中的流体撞击坑和超高速撞击实验，Proc。 第六届超高速冲击研讨会，1963年4月30日至5月2日，俄亥俄州克利夫兰。 3. 麦科马克，R. W.，低环境压力下的冲击闪光，美国宇航局TND-2232，1964年3月。 4. 麦科马克，R. W.，超高速撞击问题的数值解决方案，美国宇航局OART流星体撞击穿透研讨会，载人航天器中心，1968年10月8日至9日，pp. 180-193。 5. 麦科马克，R. W.，黏度对超高速撞击坑的影响，AIAA论文69-354，AIAA超高速撞击会议，俄亥俄州辛辛那提，1969年4月30日至5月2日。 6. 麦科马克，R. W.，激波与层流边界层相互作用的数值解，第二届流体力学数值方法会议，加利福尼亚州伯克利，1970年9月15日至19日，物理学讲义，卷。 8，Springer-Verlag，1971年，pp. 151-163。

R的贡献。 W。 MACCORMACK 19 7. 麦科马克，R. W。 \& Paullay，A。 J.，有限差分算子的时间分割实现的计算效率，AIAA论文72-154，第10届航空航天科学会议，加利福尼亚州圣地亚哥，1972年1月17日至19日。 8. 麦科马克，R. W。 \& 温暖，R。 F.，带冲击的三维超音速无黏流动的计算方法调查，AGARD论文LS-64，系列讲座编号。 64关于数值流体力学的进步，冯·卡曼研究所，布鲁塞尔，1973年3月5日至9日。 9. 奥尔森，L。 E.，McGowan，P. R。 \& MacCormack，R. W.，入口区域中时间依赖性可压缩Navier-Stokes 方程的数值解，美国宇航局-TM-X-62SS8，1974年3月。 10. 鲍德温，B。 S。 \& MacCormack，R. W.，强激波与湍流边界层的相互作用，AIAA论文74-558，流体和等离子动力学会议，加利福尼亚州帕洛阿尔托，1974年6月17日至19日。 11. 鲍德温，B。 S。 \& MacCormack，R. W.，强激波与湍流边界层的相互作用，第四届流体力学数值方法国际会议，1974年6月24日至28日，科罗拉多州博尔德，物理学讲义，卷。 35，Springer-Verlag，1975年，pp. 51-56。 12. Rizzi，A。 W.，Klavins，A. \& MacCormack，R. W.，带冲击的三维超声速流动的广义双曲行进技术，第四届流体力学数值方法国际会议，科罗拉多州博尔德，1974年6月24日至28日。 在物理学讲义中，卷。 35，Springer-Verlag，纽约，1975年，p. 341-346。 13. 麦科马克，R. W。 \& Paullay，A。 J.，计算网格对不连续或非唯一解的初始值问题的准确性的影响，计算机与流体，卷。 1974年12月2日，pp. 339-361。 14. 鲍德温，B。 S。 \& MacCormack，R. W.，用于冲击-边界层相互作用的求解Navier-Stokes 方程的数值方法，Sandia实验室预印本SLA-74-5009，新墨西哥州阿尔伯克基，1974年。 15. 麦科马克，R. W。 和鲍德温，B。 S.，用于冲击-边界层相互作用的求解Navier-Stokes 方程的数值方法，AIAA论文75-1，第13届航空航天科学会议，加利福尼亚州帕萨迪纳，1975年1月。 16. 亨，C。 M。 \& MacCormack，R. W.，超音速和高超音速层压角流的数值解，AIAA论文75-2，第13届航空航天科学会议，加利福尼亚州帕萨迪纳，1975年1月20日至22日。 17. Baldwin，Barrett S.，MacCormack，Robert W.和Deiwert，George S.，可压缩Navier-Stokes 方程解决方案和湍流模型实现的数值技术，AGARD系列讲座编号。 73，比利时布鲁塞尔，1975年2月17日至22日，pp. 2-1 - 2-24。 18. 亨，C。 M。 \& MacCormack，R. W.，超音速和高超音速层压角流的数值解，AIAA J.，卷。 1976年4月14日，pp. 475-481。 19. 麦科马克，R. W.，双曲方程组的快速求解器，第五届流体力学数值方法国际会议，荷兰恩斯赫德，1976年6月28日至7月2日，物理学讲义，卷。 59，Springer-Verlag，1976年，pp. 307-317。 20. 鲍德温，B。 S。 \& MacCormack，R. W.，分离边界层的墙和代数湍流建模的修改，AIAA论文76-350，第9届流体和等离子动力学会议，加利福尼亚州圣地亚哥，1976年7月14日至16日。 21. 亨，C。 M。 \& MacCormack，R. W.，超音速的数值模拟和

20 CAUGHEY \& HAFEZ高超音速湍流使用松弛模型的压缩角流，AIAA论文76-410，第9届流体和等离子动力学会议，加利福尼亚州圣地亚哥，1976年7月14日至16日。 22. 麦科马克，R. W.，在高雷诺数下解决时间依赖性可压缩Navier-Stokes 方程的高效数值方法，美国宇航局TM X-73129，1976年7月。 23. 麦科马克，R. W.，在高雷诺数下求解时间依赖可压缩Navier-Stokes 方程的有效数值方法，应用力学计算，AMD卷。 18，ASME，纽约，1976年。 24. 麦科马克，R. W.，Rizzi，A. W.，和Inouye，M.，具有嵌入式超音速区域的稳定超音速流场，在航空流体力学中的计算方法和问题中，B。 L。 Hewitt等人，编辑。 学术出版社，纽约，1976年，pp. 424-447。 25. 麦科马克，R. W。 \& 史蒂文斯，K. G。 Jr.，ILLIAC IV计算机的流体动力学应用，航空流体动力学中的计算方法和问题，B. L。 Hewitt等人，编辑。 学术出版社，纽约，1976年，pp. 448-465。 26. 亨，C。 M。 \& MacCormack，R. W.，超音速和高超音速湍流压缩角流的数值模拟，AIAA J.，卷。 1977年3月15日，pp. 410-416。 27. 亨，C。 M。 \& MacCormack，R. W.，三维压缩角上的超音速层流的数值解，AIAA论文77-694，流体和等离子动力学会议，新墨西哥州阿尔伯克基，1977年6月。 28. 亨，C。 M。 \& MacCormack，R. W.，三维激波和湍流边界层相互作用的数值解，AIAA文件78-161，第16届航空航天科学会议，阿拉巴马州亨茨维尔，1978年1月16日至18日。 29. 麦科马克，R. W.，使用数值方法研究高雷诺数下的复杂流动的现状和未来前景，AGARD论文编号 LS-94，系列讲座编号 94关于高雷诺数的三维非稳分离，冯·卡曼研究所，布鲁塞尔，1978年2月20-24日。 30. 麦科马克，R. W.，高雷诺数黏度流动的数值解，Proc。 第26传热和流体力学研究所，华盛顿州普尔曼，1978年6月26日至28日，斯坦福大学出版社，pp. 218-221。 31. 亨，C。 M。 \& MacCormack，R. W.，三维激波和湍流边界层相互作用的数值解，AIAA J.，卷。 1978年10月16日，pp. 1090-1096。 32. 马塔尔，F。 P.，Teichmann，J.，Bissonnette，L. R。 \& MacCormack，R. W.，使用分割和重新分区技术进行非线性光学传播的流体方法的显式算法，在Proc中。 第二届国际气流和化学激光研讨会，比利时罗德-圣日塞，1978年，半球，华盛顿特区，pp. 437-448。 33. 麦科马克，R. W.，解决可压缩Navier-Stokes 方程的有效显式-隐式特征方法，在计算流体动力学中，SIAM-AMS会议记录，卷。 XI，美国数学学会，1978年，pp. 130-155。 34. 麦科马克，R. W。 \& Lomax，H.，可压缩黏性流动的数值解，Ann。 牧师。 流体力学，卷。 1979年11月，pp. 289-316。 35. 雷诺，W。 C. \& MacCormack，R. W.，编辑，第七届流体力学数值方法国际会议，斯坦福，加利福尼亚州，1980年6月，物理学讲义，卷。 141，Springer-Verlag，1981年。 36. 侯赛尼，M. Y.，鲍德温，B. S.和MacCormack，R. W.，渐近特征

R的贡献。 W。 MACCORMACK 21激波边界层相互作用，AIAA J.，1980年8月，pp. 1014-1016。 37. 麦科马克，R. W.，求解可压缩方程的数值方法，黏性流动，AIAA论文81-0110，第19届航空航天科学会议，St. 密苏里州路易斯，1981年1月。 38. 麦科马克，R. W.，高雷诺数下可压缩黏性流动的数值解NASA-TM-81279，1981年3月。 39. 科杜拉，W。 \& MacCormack，R. W.，跨声速流动使用显式显式方法计算，Proc。 第八届流体力学数值方法国际会议，德国亚琛，1982年6月至7月，Springer-Verlag，pp. 286-295。 40. 麦科马克，R. W.，求解可压缩方程的数值方法，黏性流动，AIAA J.，卷。 1982年9月20日，第1275-1281页。 41. 麦科马克，R. W.，可压缩黏性流动方程的数值解，在跨声速、冲击和多维流动中：科学计算的进步，学术出版社，纽约，1982年，pp. 161-179。 42. 古普塔，R。 N.，Gnoffo，P. A.和MacCormack，R. W.，显式隐式方法的黏性冲击层流场分析，AIAA论文83-1423，第18届热物理学会议，加拿大蒙特利尔，1983年6月1日至3日。 43. 尚，J。 S。 \& MacCormack，R. W.，带有后体压缩襟翼的双形配置的流动-比较数值研究，AIAA论文83-1668，第16届流体和等离子动力学会议，马萨诸塞州丹弗斯，1983年7月12日至14日。 44. 麦科马克，R. W.，麦克马斯特，D. L.，Kao，T. J。 \& Imlay，S. T.，Navier-Stokes 方程在二维推力反转喷嘴内流动的解决方案，AIAA论文84-0344>第22届航空航天科学会议，内华达州里诺，1984年1月9日至12日。 45. Kneile，K。 R。 \& MacCormack，R. W.，用于内部流动的3D可压缩Navier-Stokes 方程的隐式解决方案，Proc。 第九届流体力学数值方法国际会议，1984年6月25日至29日，法国萨克莱，物理学讲义，卷。 218，pp. 302-307。 46. 麦科马克，R. W.，Navier-Stokes 方程的数值方法，超级计算和计算流体动力学的进步，美国。 /以色列研讨会，耶路撒冷，以色列，1984年12月，Birkhaeuser，波士顿，pp. 143-153。 47. 麦科马克，R. W.，Navier-Stokes计算的收敛加速度，在大规模科学计算中，S. Parter，Ed.，学术出版社，1984年，pp. 161-193。 48. 麦科马克，R. W.，关于飞行中完整飞机的可压缩黏性流动场的数值解，在数值方法的最新进展中，卷。 Ill：黏性流动，W。 G。 Habashi，Ed.，Pineridge Press，Swansea，1984年，pp. 225-254。 49. 麦科马克，R. W.，Navier-Stokes方程数值解的现状，AIAA论文85-0032，第23届航空航天科学会议，内华达州里诺，1985年1月14日至17日。 50. 古普塔，R。 N.，Gnoffo，P. A.和MacCormack，R. W.，显式隐式方法的黏性冲击层流场分析，AIAA J.，卷。 1985年5月23日，pp. 723-732。 51. Viegas，John R.，Rubesin，Morris W.和MacCormack，R. W.，Navier-Stokes计算和湍流循环控制翼型的后缘区域建模，循环控制研讨会会议记录，美国宇航局艾姆斯研究中心，加利福尼亚州莫菲特菲尔德，1986年2月19日至21日。 52. 里贝，F。 L.，克里斯蒂安森，W. H.，MacCormack，R. W.，桑卡兰，L. \& Yaghmaee，

22 CAUGHEY \& HAFEZ S.，冲击-融合目标动力学的数值研究，ICENES会议，西班牙马德里，1986年7月7日。 53. 麦科马克，R. W.，可压缩黏性流动的有限体积法，可压缩流动的数值方法-有限差分、元素和体积技术，ASME冬季年会，加利福尼亚州阿纳海姆，1986年12月7日，AMD卷。 78，pp. 159ff。 54. 坎德勒，G。 V。 \& MacCormack，R. W.，高超声速流动过去的3D配置，AIAA论文87-0480，第25届航空航天科学会议，内华达州里诺，1987年1月12日至15日。 55. MacCormack，Robert W.，Chapman，Dean R.，k. Gocken，Tahir，计算流体动力学靠近连续极限，Proc。 AIAA第八届计算流体动力学会议，夏威夷檀香山，1987年6月9日至11日，pp. 153-158。 56. 坎德勒，G。 V。 \& MacCormack，R. W.，化学和热不平衡中高超声速流动的计算，论文编号 107，Proc。 1987年6月，美国宇航局艾姆斯研究中心，加利福尼亚州莫菲特菲尔德，第三届国家航空航天飞机技术研讨会。 57. 麦科马克，R. W。 \& 坎德勒，G。 V.，在传感、辨别、信号处理和超导材料和仪器中预测高超音速流场的数值方法，光光学仪器工程师协会，加利福尼亚州洛杉矶，1988年1月12-14日，pp. 123-129。 58. 坎德勒，G。 V。 \& MacCormack，R. W.，化学和热不平衡中高超音速电子流动的计算，AIAA论文88-0511，第26届航空航天科学会议，内华达州里诺，1988年1月。 59. 维加斯，J。 R.，Rubesin，M. W.，和MacCormack，R. W.，关于适合循环控制机翼的代码和湍流模型的验证，AGARD，计算流体动力学的验证。 第1卷：研讨会论文和圆桌讨论，葡萄牙里斯本，1988年5月2日至5日。 60. 坎德勒，G。 V。 \& MacCormack，R. W.，斯坦福大学高超音速研究，高级航空航天空气动力学；Proc。 航空航天技术会议和博览会，加利福尼亚州阿纳海姆，1988年10月3日至6日，pp. 257-265。 61. 麦科马克，R. W.，关于三维流体流动的高效算法的开发，计算流体动力学的最新发展，ASME冬季年会，伊利诺伊州芝加哥，1988年11月27日至12月2日，pp. 117-137。 62. 维加斯，J。 R.，Rubesin，M. W。 \& MacCormack，R. W.，关于适合循环控制机翼的代码和湍流模型的验证，美国宇航局TM-100090，1988年。 63. 麦科马克，R. W.，和坎德勒，G。 V。 预测高超音速流场的数值方法，第2届欧洲/美国联合 高超音速短期课程，科罗拉多州科罗拉多斯普林斯，1989年1月16日至20日。 64. 戈格肯，T。 \& MacCormack，R. W.，使用连续体方法的高超音速过渡流动的非平衡效应，AIAA论文89-0461，第27届航空航天科学会议，内华达州里诺，1989年1月9日至12日。 65. 麦科马克，R. W.，和坎德勒，G。 V。 使用高斯-塞德尔松弛的Navier-Stokes 方程解决方案，计算机和流体，卷。 1989年17月，pp. 135-150。 66. 罗斯斯坦，P。 \& MacCormack，R. W.，CFD电弧加热射流的建模，AIAA论文90-1475，第21届流体力学、等离子动力学和激光会议，华盛顿州西雅图，1990年6月18日至20日。 67. 麦科马克，R. W.，Navier-Stokes 方程的三解决方案

R的贡献。 W。 MACCORMACK 23尺寸，AIAA论文90-1520，第21届流体力学、等离子动力学和激光会议，华盛顿州西雅图，1990年6月18日至20日。 68. 威尔逊，格雷戈里·J。 \& MacCormack，Robert W.，使用全隐式数值方法对超音速燃烧进行建模，AIAA论文90-2307，第26届联合推进会议，佛罗里达州奥兰多，1990年7月16日至18日。 69. 钟，X.，MacCormack，R. W.，和查普曼，D. R.，伯内特方程的稳定和对高海拔高超声速流动的应用，AIAA论文91-0770，第29届航空航天科学会议，1991年1月7日至10日。 70. 康蒂，劳尔·J。 \& MacCormack，R. W.，高超声速流动的廉价Navier-Stokes计算，AIAA论文91-1391，第26届热物理学会议，夏威夷檀香山，1991年6月24-26日。 71. 莫罗，斯蒂芬，查普曼，D。 R。 \& MacCormack，R. W.，旋转松弛和近似伯内特项对高海拔高超音速流场辐射的影响，AIAA论文91-1702，第22届流体力学、等离子体动力学和激光会议，夏威夷檀香山，1991年6月24日至26日。 72. 坎德勒，G。 V。 \& MacCormack，R. W.，热化学非平衡中弱电子高超声速流动的计算，J. 热物理学和热传递，卷。 5，不。 3，1991年7月，pp. 266-273。 73. Menon, Suresh, \& MacCormack, Robert W., 使用隐式Navier-Stokes求解器对三维火焰保持器附近超音速混合的数值研究，Proc J\textasciicircum{}th 计算流体动力学国际研讨会，加利福尼亚州戴维斯，1991年9月9日至12日，pp. 801-806。 74. Zhong, Xiaolin, MacCormack, Robert W., \& Chapman, Dean R., Burnett方程的Slip 边界条件的评估，应用于高超音速领先边缘流，Proc J\textasciicircum{}th 计算流体动力学国际研讨会，加利福尼亚州戴维斯，1991年9月9日至12日，pp. 1360-1366。 75. 罗斯斯坦，P。 \& MacCormack，R. W.，弧射流中的非平衡流，再入问题的高超音速流，卷。 2，Springer-Verlag，柏林，1991年，pp. 1102-1115。 76. Conti, Raul J., MacCormack, Robert W., Groener, Liam S., \& Fryer, Jack M., Practical Navier-Stokes 计算耦合消融和形状变化的轴对称再入流场，AIAA论文92-0752，第30届航空航天科学会议，内华达州里诺，1992年1月6日至9日。 77. 康蒂，劳尔·J。 \& MacCormack，Robert W.，Navier-Stokes用外来气体注入高超音速近波的计算，AIAA论文92-0838，第30届航空航天科学会议，内华达州里诺，1992年1月6日至9日。 78. 威尔逊，格雷戈里·J。 \& MacCormack，Robert W.，使用完全隐式数值方法对超音速燃烧进行建模，AIAA J.，卷。 1992年4月30日，pp. 1008-1015。 79. 莫罗，斯蒂芬，劳克斯，克里斯托夫·O.，查普曼，迪恩·R. \& MacCormack，Robert W.，更准确的非平衡空气辐射代码-NEQAIR第二代，AIAA论文92-2968，第23届等离子动力学和激光会议，田纳西州纳什维尔，1992年7月6日至8日。 80. Moreau，S.，Bourquin，P. Y.，查普曼，Dean R. \& MacCormack，Robert W.，夏尔马冲击管实验的数值模拟，AIAA论文93-0273第31届航空航天科学会议，里诺。 内华达州，1993年1月11日至14日。 81. Zhong, Xiaolin, MacCormack, Robert W., \& Chapman, Dean R., 伯内特方程的稳定和对高超声速流动的应用，AIAA J., Vol. 1993年6月31日，pp. 1036-1043。 82. Comeaux，Keith A.，Chapman，Dean R. \& MacCormack，Robert W.，Viscous

24 高海拔钝体上的CAUGHEY和HAFEZ高超音速冲击-冲击相互作用，AIAA论文93-2722，第28届热物理学会议，佛罗里达州奥兰多，1993年7月6日至9日。 83. 焊工、华莱士T.、查普曼、迪恩·R. \& MacCormack，Robert W.，Burnett方程的各种形式的评估，AIAA论文93-3094，第24届流体力学会议，佛罗里达州奥兰多，1993年7月6日至8日。 84. MacCormack，Robert W.，CFD研究四分之一世纪的视角，Proc。 AIAA第11届计算流体动力学会议，佛罗里达州奥兰多，1993年7月6日至9日，pp. 1-15。 85. 麦科马克，R. W.，求解关于航空航天飞行器的可压缩黏性流动方程，在航空航天科学与工程普莱姆出版社的应用数学中，纽约，1994年，pp. 25-34。 86. Comeaux，Keith A.，Chapman，Dean R. \& MacCormack，Robert W.，基于热力学第二定律的伯内特方程分析，AIAA论文95-0415，第34届航空航天科学会议，内华达州里诺，1996年1月15日至18日。 87. 考，C。 T.，von Ellenrieder，K.，MacCormack，R. W.和Bershader，D.，二维可压缩涡旋冲击相互作用的物理分析，AIAA论文96-0044，第34届航空航天科学会议，内华达州里诺，1996年1月15日至18日。 88. Moreau，Stephane，Chapman，Dean R. \& MacCormack，Robert W.，I的数值模拟。 R。 冲击管实验中的辐射，AIAA论文96-0108，第34届航空航天科学会议，内华达州里诺，1996年1月15日至18日。 89. 梅尔维尔，R. \& MacCormack，R. W.，双曲和抛物线系统的优化、显式时间积分方法，AIAA论文96-0531，第34届航空航天科学会议，内华达州里诺，1996年1月15日至18日。 90. 梅尔维尔，R. \& MacCormack，R. W.，可压缩流动的自由涡流突发模拟，AIAA论文96-0805，第34届航空航天科学会议，内华达州里诺，1996年1月15日至18日。 91. MacCormack，Robert W.，隐式算法的高效矩阵分解，Proc。 第15届流体力学数值方法国际会议，1996年6月24日至28日，加利福尼亚州蒙特雷，pp. 237-242。 92. MacCormack，Robert W.，流体流动的新隐式算法，Proc。 AIAA第13届计算流体动力学会议，科罗拉多州斯诺马斯，1997年6月29日至7月2日，pp. 112-119。 93. MacCormack，Robert W.，快速导航-斯托克斯求解器的考虑因素，Proc。 流动模拟技术的进步，加利福尼亚州戴维斯，1997年5月2日至4日，pp. 107-117。 94. MacCormack，Robert W. \& Pulliam，Thomas，流体动力学新数值程序的评估，AIAA论文98-2821，第29届流体力学会议，新墨西哥州阿尔伯克基，1998年6月15日至18日。 95. MacCormack，Robert W.，在流动计算中增加的消散，在计算流体动力学的前沿-1998年，世界科学，新加坡，D。 A. 考伊和M。 M。 哈菲兹，编辑，pp. 171-185。 96. 普利姆，T。 H.，MacCormack，Robert W. \& Venkateswaren，S.，收敛近似因数分解方法的特征，第十六届流体力学数值方法国际会议，法国阿拉雄，1998年7月6日至10日，物理学讲义卷。 515，pp. 409-414。 97. MacCormack，Robert W.，解决Navier-Stokes 方程的快速准确方法，WAS论文98-2,7,2，第21届ICAS大会，澳大利亚墨尔本，1998年9月13日至18日。 98. MacCormack，Robert W.，一种上风保护形式方法

R的贡献。 W。 MACCORMACK 25磁流体动力学，AIAA论文99-3609，第30届等离子动力学和激光会议，弗吉尼亚州诺福克，1999年6月28日至7月1日。 99. MacCormack，Robert W.，磁流体力学中的数值计算，21世纪的计算流体动力学，日本京都，2000年7月5日至7日。

\clearpage
\chapter{2 黏度对超速度撞击坑的影响}
\section{2.1 摘要}
\section{2.2 介绍}
黏度对超高速撞击坑的影响 Robert W. MacCormack1 2.1 摘要 提出了一种时间和空间二阶的数值方法，用于求解与时间相关的可压缩Navier-Stokes 方程。 讨论了稳定的条件。 该方法已被应用于计算超高速冲击产生的轴对称流场。 给出了铝圆筒（直径为0.16、0.32和0.64厘米）对铝靶的冲击结果。 假定黏度为零和104的平衡。 板和半无限目标都以10公里/秒的撞击速度被考虑。 结论是，随着弹丸尺寸的缩小，黏性效应变得越来越重要，在直径小于0.5厘米的弹丸的陨石坑形成的初始阶段不能忽视。 2.2 导言Denardo [1, 2]在1964年报告了直径为0.16、0.32、0.64和1.27厘米的铝球体对铝靶的超高速冲击的简单线性缩放的偏差。 渗透和动量传递 1 美国宇航局艾姆斯研究中心，莫菲特菲尔德，加利福尼亚州94035。 本文最初在俄亥俄州辛辛那提举行的AIAA超高速冲击会议上作为AIAA论文69-354发表。 感谢AIAA允许在本卷中重新出版这篇经典论文。 Frontiers of 计算流体动力学 - 2002 编辑：David A. Caughey和Mohamed M. 哈菲兹 ©2002 世界科学

\section{2.3 数值方法}
28 MACCORMACK对目标的减少速度比简单的缩放规则所暗示的要快，因为弹丸尺寸被缩小了。 更明确地说，穿透与弹丸直径的比率以及撞击后的目标动量与弹丸动量的比率分别因弹丸直径的1/18和1/6的幂而变化。 这种现象的两个可能来源是：（a）如Kuhn和Figge [3]所示，材料静态强度的纯尺度效应；（b）与速率相关的应力效应，例如黏度。 前者将在陨石坑形成结束时起作用；后者在陨石坑形成的早期部分起作用，那里存在高剪切率。 本研究的目的是观察火山口初始阶段黏度的影响（定义为从撞击开始到静态强度效应变得显著）。 在此期间，超高速冲击坑可以用流体力学的Navier-Stokes 方程来描述，其解对黏度与零不同的非线性缩放。 描述了一种求解这些时间依赖方程的数值方法。 该方法在时间和空间上都是二阶的，因此比Los Alamos粒子细胞代码[4]得出的方法更准确，后者以前曾用于[5,6, 7]来计算超高速冲击现象。 检查了几种影响案例的计算结果，以确定黏度的影响。 2.3 数值方法 2.3.1 微分方程 流体动力学的Navier-Stokes 方程，忽略体力和热源，可以写成 [8] dp dpu dxj = 0（连续方程）dpUi dpuiUj dp dai\}j - 1 - + - - - = 0（动量方程）在oxj oxi axj de d(e + p)uj d(uia\textasciicircum{}j - qj) - H = 0（能量方程）dt dxj dxj e \textbackslash\{\}u\textbackslash\{\}2 p = f(e,p) = f I 2\textasciitilde{}'W（e<状态e<luation），其中使用了爱因斯坦符号（重复指数相加）和和t = 时间p = 时间p = 单位体积质量 Ui = Xi方向辅助速度u的分量 = 黏性应力 = Ac\textasciicircum{}f

超速冲击坑 29 Sij - 克罗内克三角洲p，X 分别是黏度的第一和第二系数，e - 每单位体积总能，其中k是导热系数，温度T和比内能e由e = g（p，T）形式的方程相关。 最后，压力p表示为e和p的函式。 这组方程可以用其他坐标系轻松表达[9]。 已经设计了一种两步差分法来数值求解上述集合。 为了简单起见，也为了与文献中对其他数值技术的描述保持一致，该方法将在二维笛卡尔坐标系中说明。 稍后将简要评论对轴对称系统的应用。 此外，对于数值稳定性的分析，只对非黏性非导热方程进行详细处理。 黏性和热传导项，无论是物理还是数值（如果它们在微分方程中的幅度不是太大），都倾向于抑制解的高频分量。 这些组件通常是导致数值不稳定的组件。 因此，就稳定性分析而言，这里选择p和k等于零的方法是保守的。 向量形式的方程是dU dF dG,„，其中U m \textbackslash\{\} / n \textbackslash\{\} F=\textbackslash\{\} m2lp + P, G=\textbackslash\{\} ™/p mn/p I I n /p + p.(e + p)m/p/ \textbackslash\{\}(e + p)n/p/ P = /(e./o) = / e m2 + n2,P 2p2 "7和m = pu和n = pv分别是x和y方向的单位体积的矩。 2.3.2 差分方程 求解方程的两步差分法。 (2.1)由U?,k = Ulk \textasciitilde{} -\textasciicircum{} \{Fj+i,k - F",k) \textasciitilde{} -£y (G",fe+1 - G",fe) (2-2) UT = \textbackslash\{\} fa + uJ\textasciicircum{}-\textasciicircum{}(FJf- Ff\textasciicircum{}>k) - \textasciicircum{} (G\&T - G\textasciicircum{}) \}(2.3)

30 MACCORMACK，其中F™k和G"fc分别等于F(U™k)和G(U"k)。 现在和从这里开始，下标指的是空间点网（xj，yk），间距为Ax和Ay，上标指的是时间t = nAt，其中At是方程的每个周期中解推进的时间增量。 （2.2，2.3）。 该方法首先通过Eq获得近似值，U"\textasciicircum{}1在每个点。 （2.2），它使用两个正向差来近似两个空间导数。 然后，近似解用于计算F"kx和G"\textasciicircum{}1，它们用于等式。（ 2.3）有两个向后差异，以获得新的接受值U\textasciicircum{}t1。 与细胞中的粒子方法一样，流场可以被视为分为细胞，力作用在细胞面上，质量、动量和能量在细胞之间传输。 微分方程组为守恒律形式[10]。 同样，检查一组差分方程。 （2.2，2.3），表明质量、动量和能量在计算过程中是守恒的；即等式中的差量。 （2.2,2.3）在网格的内部点，在扫过网格时正好出现两次，每次都是相反的符号。 因此，向量和，\textasciicircum{}2j k U? kl或Sj fc \textasciicircum{}Tfc"1代表边界处的净运输和应力效应。 2.3.3 准确性和稳定性 这种类型的方法的数值稳定性，即Lax和Wendroff [10]的方法，目前无法以一般非线性形式完全分析。 迄今为止最成功的分析尝试是首先将微分方程集线性化。 （2.1）然后通过应用于线性集合的差分法研究解的傅里叶分量的放大。 然后，新的方程集在ox ay（2.4）处为dU T dU T dU n，其中Jp和JQ是F和G相对于U的雅各布矩阵，被认为是常数。 JF JG m2, §£ P2 \textasciicircum{} dp mn m e+p, m dp P P P p mn P n2, dp p2 ""f" dp \_ n e+p, n dp P P P dp 1m, dp p dm n P e+p. m dp p p dm n P dp 2n dm p dp e+p p p dm P dp dn P dn P dn dp dn dn dp dn dp dp dn dp de de de 0 n n de n dp p p p /

超速度冲击坑 31 dp dp dp dp dp d\textasciitilde{}e dp dp dp dn dp de e de dp dp de 1 dp de de dp de m dm p2 de dp de n dp de dn p2 de dp "dp <+B+ m2 + n2 dp \textasciitilde{}d\textasciitilde{}e 方程集。 （2.4）近似等式。 （2.1）局部，在一般非线性情况下，发现不稳定的差分方法可以预期是不稳定的。 这种分析中固有的两个条件是：（a）边界条件对稳定性没有影响，以及（b）Eq的确切解。 （2.1）是光滑的。 后一个条件允许矩阵Jp和JQ被视为常数。 因此，分析结果主要适用于远离边界且没有不连续的流动区域，如激波。 在接下来的几段中，将使用这种分析方法来获得在满足这些条件的区域中解的任何傅里叶分量的最大放大的边界。 当前方法的放大矩阵应用于等式。 （2.4）对于解决方案的单个傅里叶分量，变成[11]，G = I+ W\{t)exp[i\{k1x + k2y)] i\textasciicircum{}(JFsin£, + JGsmr]) l(\textasciicircum{})2((l-\textasciicircum{})\textasciicircum{}+(l-e--)JG) ((l - e*) JF + (l - J") JG))) (2.5) 其中£ = kiAx，r\textbackslash\{\} - \textasciicircum{}Ay，并假定等距（Ax = Ay）。 对于£和n « 1，近似正弦；= £，1 - cos£ = £2/2，以及1 + cos£ = 2 - £2/2以及n的类似近似值可用于表明G r • Af / T /• T \textasciicircum{} 1 / At /\textasciicircum{}-(JK+J\textasciicircum{})--\textasciicircum{}- (JFC + JGV) + 0(£W) 因此G = exp \textbackslash\{\}\}\textasciicircum{}.\{ JF£，+ JGV）] f和-q中第三阶的模项。 微分方程的精确解。 （2.4）对于上述傅里叶分量是exp[it(k1 JF + k2JG)] exp[i(fcix + fc2y)]W(0)

32 MACCORMACK 因此，从t = 0到t = At的解的精确放大是“At exp[iAt(ki Jp + k2 JG)\} = exp I-\textasciicircum{}(JF£ + JGV) 两者之间的差异是£和r/的三阶。 因此，由Eqs定义的方法。 （2.2，2.3）是二阶精度[10]。 G也可以写成G = I + M-M*- 2M*M，其中2 Az u e«) JF+(1- eir>) JG)，M*是M的复共轭。 如果能够证明G的所有特征值Aj满足冯诺伊曼条件|Aj|<1 + 0（At）[11]，则保证了稳定性。 在相似性变换下，G的特征值是不变的。 使用这种变换S，M可以写成2 Ax，其中S\textasciitilde{}l = s = 其中A = B = •u c 0 0\textbackslash\{\} c u 0 0 0 0 0 u 0 0 0 0 u. /v 0 c 0' v 0 0 c 0 v 0 Vo o o v. c 0 0 0 /3/c 0 0 0 0 0 10,0001 c2 = dp dp pp p2 de'

超高速冲击火山口33是局部隔热声速的平方，e +P 2 2 a = - ul - V P and P 0€ 然后G' = S\textasciicircum{}GS = I + M'- M'* - 2M'*M'，其中M' = S\textasciicircum{}MS是对称的。 对于任何单位向量w，内积（w，G'w）=（w，w）+（w，M'w）-（w，M'*w）-2（w，M'*M'w）= 1 - 2\textbackslash\{\}\textbackslash\{\}M'w\textbackslash\{\}\textbackslash\{\}2 +（w，M'w）-（w，M'*w）其中M'w的范数用|M'w表示；||=（M'w，M'w）1/2。 数量（w,M'w）-（w，M'*w）是纯虚数，并且由于\textbackslash\{\}(w,M'w)\textbackslash\{\} < \textbackslash\{\}\textbackslash\{\}w\textbackslash\{\}\textbackslash\{\}w\textbackslash\{\}\textbackslash\{\}\textbackslash\{\}M'w\textbackslash\{\}\textbackslash\{\} = \textbackslash\{\}\textbackslash\{\}M'w\textbackslash\{\}\textbackslash\{\}由柯西-施瓦茨不等式|K G'w)\textbackslash\{\}2 < (1 - 2\textbackslash\{\}\textbackslash\{\}M'w\textbackslash\{\}\textbackslash\{\}2)2 + (2|||M'H|)2 < 1 + M'wf Gerschgorin定理[12]可用于显示\textbackslash\{\}\textbackslash\{\}M'w\textbackslash\{\}\textbackslash\{\}<\textasciicircum{}(\textbackslash\{\}u\textbackslash\{\} + \textbackslash\{\}v\textbackslash\{\} + 2c) Ax因此，如果v \textasciitilde{} v 4|£(| + \textbackslash\{\}v\textbackslash\{\} + 2c)4 ss 1，那么\textbackslash\{\}\{w,G'w)\textbackslash\{\} < 1 + uO(At)。 特别是，如果w是G'的任何特征向量，G'w = Xw。 然后|A|2 < 1 + vAt或|A| <l + 0\{At）因此，如果v是单位的阶，G的所有特征值都满足|Aj| < 1 + O(At)。 也就是说，如果v保持不变为At、Ax -> O和Eqs。 （2.2,2.3）用于计算时间T的解，n个时间步数为T/At，任何傅里叶分量的放大都小于sup | (w, Gw) |" < (1 + vAt)n/2 < e\{n? 2>At = euT/2 All w, \{w,w)\textasciitilde{}l 因此，在计算给定时间时，通过适当选择v，解决方案的每个组成部分的放大都可以任意小。

34 MACCORMACK p\textbackslash\{\} I m 在一维（即U = | m，F = m2/p + p ]和G = 0）e / \textbackslash\{\}(e+p)m/p/中，很容易证明，如果\textasciicircum{}(|u| + c) < 1，则该方法的放大矩阵的特征值小于单位。 这种情况是众所周知的Courant-Friedrichs-Lewy（C.F.L.）条件，经常出现在流体动力学中。 这是在数值方法中可以实现的最佳边界。 矩阵A和B的非交换性目前阻碍了二维特征值的计算。 从派生边界获得的条件2\textasciicircum{}(M + \textbackslash\{\}v\textbackslash\{\} + 2c) « At1/4比二维C.F.L.条件更具限制性。 虽然可以证明衍生的特征值边界不是最小上限，但也可以证明特征值确实存在，因此稳定性需要比C.F.L.条件更具限制性的条件。 然而，由Eqs定义的方法。 （2.2,2.3）只是本质上相同形式的四种二阶精度方法之一。 例如，如果不是首先使用两个正向空间差异，然后是两个向后差异，可以遵循反向程序，或者一个正向和一个后向差分，然后是相应的向后和正向差异。 对于解的同一傅里叶分量，每个放大矩阵都有不同的特征值和特征向量。 因此，如果方程定义的方法的索引。 （2.2,2.3）被排列，因此四种方法相互循环地相继，预计放大会更小；即，尽管每个Gi的最大特征值大小|Amax（Grj）|在所有£、r]、U和v的集合上最大化，因此|M| + \textbackslash\{\}v\textbackslash\{\} <常数和c是相同的，但£、r/、u和v的单一选择通常不会最大化所有Gi（即|Amax（Gi C\textasciicircum{}G\textasciicircum{}G\textasciicircum{}G\textasciicircum{}）！ < |Amax(Gi)4，i = 1，2，3和4）。 推测，如果At和Ax满足接近C.F.L.条件的条件，则\textbackslash\{\}Xm\textasciicircum{}\{G1G2G3G4)\textbackslash\{\} < 1 + O(At)。 如前所述，如果黏性和热传导项的大小不是太大，则添加黏性和热传导项不会干扰数值稳定性（也就是说，如果选择At和Ax，使\textasciicircum{}\textasciicircum{}-，k-\textasciicircum{}足够小于1）。 如果术语不同，因此它们的截断误差也是二阶的，那么二阶精度也将保持不变。 例如，黏性项-gf5\textasciicircum{}，如果中心差异，（Hi + l+IJ-i Uj + l-Uj \_ Mi+Mi-1 M.-Mj-l \textbackslash\{\} \textasciicircum{}2 Ax 2 Ax J，或者如果p是恒定的Ax Ui+i - 2ui +Ui\_i \textasciicircum{}2 Axz将足以达到二阶精度。

超速度撞击 CRATERING 35 稳定性分析在轴对称圆柱坐标上也基本相同。 例如，对应于Eq的方程组。 （2.4）是8U T ldrU T dU T U -7\textasciicircum{} + JF-\textasciicircum{}- + JGlr = JH-（2.6）在T或oz r，其中r和z是径向和轴向坐标，Jp、Jo和U是这些坐标中定义的对应矩阵和向量，JH是H的雅各布，HT =（0，p，0，0），关于U。 在远离轴的流动区域，r » Ar，可以显示从Eq中删除术语JH\textasciitilde{}的效果。 （2.6）导致仅At阶的相关放大矩阵的特征值发生变化。 因此，为了分析应用于Eq的差分方法的稳定性。 （4）在远离轴线的区域，Eq的右侧。 （4）可以设置为零向量。 这个方程的解的傅里叶分量与之前考虑的相同波数fci和k2是-exp[it(kiJF + k2Ja)\} exp[i(kxr + k2z)]同样，解与差分方程的对应分量，其中£ -j\textasciicircum{}-正向差为1 (i + 1/2)ArUi+itj - (i - 1/2)Arty M (t - l/2)Ar Ar，向后差为1 (i - l/2)Artyd -(i- 3/2)Art/i\_1J (i - l/2)Ar Ar是-W\{t) exp[i\{kir + k2z)]该分量的放大矩阵由Eq. （2.6），在微分后，与之前得到的相同，只是现在x和y被r和z取代。 在轴r=0附近，轴对称性诱导的边界条件（即u\textbackslash\{\}j = -u-ij、vij = i>-i、j、Pxj - P\textasciitilde{}i、j等）预计会影响稳定性，因此上述线性分析是不够的。 尚未分析该区域的数值稳定性。 此外，来自径向动量方程的H的非零分量不允许以发散形式表示方程[13]。 因此，差分法适用于等式。 （2.6）严格只保存质量、轴向动量和能量，而不保留径向动量。 同样，微分方程的二阶项预计不会干扰数值稳定性，如果也与二阶精度不同，

36 MACCORMACK方法本身将具有二阶精度。 例如，将术语\textasciicircum{}yr>与iAr fw+ij+wjj（Ui+1fcUiA - (t - l)Ar fm1i±ti=i\textasciicircum{}\textbackslash\{\} (»M-A»-'\textasciicircum{} (i-l/2)Ar-Ar）区分，将保持准确性。 与颗粒入细胞方法相比，所述方法的优势在于其二阶精度。 在数值计算中，将讨论在计算包括黏度影响在内的超高速影响问题时使用比一阶更准确的方法的必要性。 与其他二阶精度[10、11、13、14]相比，其优势是：（a）扩展到任何欧拉坐标系都是直接的；（b）对于不黏性差分方程，在一点上推进解的计算。 （2.2,2.3），只需要了解七个相邻点，而不是通常的九个；（c）如果网格按行扫荡（x方向），并且解仅由x方向的差值来修改，例如，-\textasciicircum{}（Fj+\textbackslash\{\}tk - Fj,k），那么对于每个j只需要计算Fj+i,k，因为Fj\textasciicircum{}是从“cellj\_ifc”之前的计算中知道的。 同样，在完成此扫描后，网格将逐列扫描，以解释y方向的差项，再次计算和保存每个k只有一个“单元面”的传输、应力和传导项的值，从而大大减少了计算量。 其他Lax-Wendroff方法可以在不同程度上遵循此程序，有些方法要求保存前两个单元格面的值，而另一些则只能再次使用每个面的计算部分。 该方法的缺点是，如上所述，放大矩阵的特征值是未知的。 如果满足冯·诺伊曼条件所需的限制很严重，那么在某些问题上，优势（b）和（c）获得的效率可能会抵消。 对于本文所考虑的问题，At被简单地选择为两个值\textasciicircum{}\textasciitilde{}和\textasciicircum{}\textasciicircum{}-中较小的一个，其中vp是弹丸撞击速度，po是初始密度。 有了这个选择，并且没有方程定义的方法索引的排列。 （2.2,2.3），没有观察到数值不稳定的迹象。 每个问题都有32 x 33个单元格的计算网格，大约需要130个时间步骤才能完成。 在IBM 7094上，机器时间大约是15分钟。

\section{2.4 数值计算}
超速度冲击CRATEPJNG 37 2.4 数值计算由Eqs定义的方法。 （2.2,2.3）用于求解圆柱坐标中可压缩、非导热黏性流体的Navier-Stokes方程。 假设静压等于平均法向应力（即“第二黏度系数”设置为-2/3为“第一系数”。 见参考。 15）。 这些方程的解不像其非黏性对应物那样与特征大小线性缩放。 然而，如果得到一个特征大小d和黏度/i的解，那么所有特征大小和黏度d'和//的解都是已知的，所有其他参数保持相等。 也就是说，时间t、距离和黏度缩放为t -> st d -» sd和jl -» S/J，其中s是任何实数。 因此，JJL的特定选择不如比率\textasciicircum{}的选择重要。 对于所有研究的案例，弹丸是一个长度等于其直径的铝制右圆柱体，以1厘米//zsec的速度撞击铝制目标。 计算中使用的状态方程是蒂洛森[16]为铝制定的。 萨哈罗夫[17]从激波实验中推断出，铝合金（90\%铝）在0.3 Mb（兆巴）时的黏度系数\textbackslash\{\}x约为0.02兆秒（兆巴），并微弱增加，但在高达1兆巴的冲击压力下不超过0.1兆瓦。 在本文中，假设0.01 Mp的恒定值代表了压缩阶段的p值，从初始冲击，冲击压力为1.54 Mb，直到计算停止，冲击衰减到大约比铝材料强度（2或3 kb）大一个数量级。 如前所述，/x的特定选择不如比率\textasciicircum{}重要，所选fi值的结果可以缩放到任何其他选项。 在计算过程中，膨胀区域被视为隐性流动。 更明确地说，当p小于po/1.1时，其中po是初始密度，n设置为零。 一般来说，所选的p值与网格间距Ax的数字大小相同。 黏性应力项的大小与Az成正比，而最后一节中描述的二阶精度方法的截断误差与Ax2成正比。 因此，如果只使用一阶精度的方法，即单元格代码中的粒子，网格间距相同，黏性应力和截断

38 MACCORMACK错误将具有相同的数量级。 网格间距，例如Ax « p2，选择确保黏性应力与截断误差相比占主导地位是不切实际的（在ss--\&p3）。 此外，存在粒子细胞方法的稳定性被这种选择破坏的危险（即，在P.I.C.中截断引入的术语本身具有黏性作用）。 因此，由于铝黏度系数的大小，至少需要一种二阶精度的方法。 每次目标激波或弹射物接近网格边界时，计算网格都会被重新划分为Ax -> 2Ax和Ay ->• 2Ay）。 在每次计算的间隔中，确定了轴向动量Z+的总正分量和总径向动量R。 也就是说，Z+= \textasciicircum{}2 pijui:j（细胞体积）itj Uij>0和R- \textasciicircum{}2 Aj\textasciicircum{}.j（细胞体积）\textasciicircum{}-所有细胞轴向动量Z\_的总负分量，通过动量守恒等于mvp-Z+，其中mvp是弹丸动量。 准确地说，与Z+和Z-不同，R不是向量，因为数量PijVij（细胞体积）\textasciicircum{}•是代数求和的。 由于轴对称性，向量和将消失。 2.4.1 半无限目标 研究了直径为0.16、0.32和0.64厘米的弹丸对厚目标的冲击，p-0.01 Mp = 0，以确定与实验观察到的相当的动量尺度效应是否是由黏度引起的。 由初始弹丸动量归一化的Z+和R的值如图所示。 1与非维时间r相比，其中r = vp |。 黏度的影响在这里清楚地显示了每个撞击案例都有不同的曲线。 还观察到，在后期，每个病例的Z+和R几乎呈线性增加。 图。 2显示了r = 8的Z+和R与d的关系，在接近计算结束时。 从数量上讲，比例效应由曲线的斜率显示在这里，Z+和R的斜率都与零不同。 随着d的减少，两个曲线的坡度都略有增加。 例如，在d = 0.32厘米和0.64厘米处穿过Z+点的直线的斜率是0.113，而对于d = 0.16厘米和0.32厘米，斜率是0.146。 同样，R曲线的相应斜率为0.082和0.1126。 这些值是计算接近尾声的典型值，似乎没有明显变化。 需要强调的是，Z+与实验测量的目标动量不同。 期间的目标

超速度冲击陱头 39 10 r § 5 a s 3 o s 2 - vp= 1.0cm/|Asec x = vpTIME/d (i = 0Mp jl - 0 01 1 总径向动量，d = J / 0.64 cm / /, 0.32 cm / / / '/ / //。 0.16 cm "P ////// / '////// ¥// 总正向V// 轴向动量'/ y\textasciicircum{}C\textasciicircum{} 032 cm '\textasciicircum{}\textasciicircum{}\textasciicircum{} 0.16 cm I 1 1 5.0 7.5 图1 d = 0.16、0.32和0.64厘米的动量与r的计算观察到轴附近包含大量正轴向动量，以及形成陨石坑唇区域的可观的负轴向动量。 净效果将是目标的势头。 尽管如此，在数值计算中观察到的动量简单线性缩放的偏差预计将反映在实验测量中。 直径指数（图中曲线的斜率。 2）比实验中发现的大致相同尺寸范围内的球形弹丸略低。 两条曲线的斜率变化表明，d的指数取决于fi/d，在相同报告的d值下，/i值稍大，预计与实验的相关性会更好。 此外，在线粒体的动量测量中，与实验室尺寸的弹丸相比，预计与简单线性缩放的偏差更大。 2.4.2 有限靶点 还研究了黏度对薄片靶点的影响。 在每种情况下，弹丸直径为0.16厘米，撞击速度为

\section{2.5 结束的讲话}
40 MACCORMACK O P 7 Si 5 UJ S o S 4 0.1 0.2 0.3 0.4 0.5 0.6 0.7 0.8 0.9 1.0 弹丸直径，d. cm 图2 总正轴向动量Z+和总径向动量R与弹丸动量mvv与弹丸直径d在非维时间r -8 1.0 cm//isec的动量比。 厚度等于0.08、0.16和0.24厘米的板材冲击的动量结果如图所示。 3（a）和（b）。 最重要的特征是，与总径向动量相比，黏度引起的总正轴向动量衰减很大。 事实上，对于th/d < 1的情况，R几乎没有减少。 由于黏度，轴向动量减弱较大的预期后果是，由弹丸和目标材料组成的喷雾的动量在撞击板中移动时，强度会更小，发散，因此对任何后续撞击结构的破坏性较小。 对于有限目标，由于自由表面引起的压力的快速衰减，冲击过程预计将由初始流体动力学阶段主导。 在喷雾量测量和喷雾分布的固体角度中实验发现的有限目标尺度效应将增加令人信服的证据，证明在半无限目标中发现的尺度效应是由黏度引起的。 相反，这种效应的缺失将为半无限目标尺度效应是在后期强度依赖阶段引起的理论的可信度。 2.5 结束 讲话 1. 虽然尚未得出结论，在半无限靶标中实验中发现的尺度效应是由黏度引起的，但已经表明，在火山口初始阶段，总正轴向和径向动量表现出由预期黏度水平引起的效应，一致的R/mv„ J I I I I I I

B INS总正轴向动量/弹丸动量总径向动量//P

\section{参考资料}
42 MACCORMACK与实验动量测量。 预计随着弹丸尺寸的缩小，这种效应将变得越来越重要。 2. 薄靶标的黏度有望降低喷雾通过穿孔靶标的动量。 一项实验研究，其中弹丸薄片几何形状随着尺寸的变化而保持不变，可以证实黏度在超高速冲击中的重要性，并为在超高速冲击条件下有效黏度的实验评估提供一种方法。 参考资料1。 德纳多，B。 Pat和Nysmith，C. 罗伯特，与铝球体以每秒24,000英尺的速度撞击到厚铝目标相关的动量转移和火山口现象。 AGARDograph 87，卷。 1. Gordon and Breach，科学出版社，纽约，1966年。 2. 德纳多，B。 帕特，萨默斯，詹姆斯·L。 和Nysmith，C。 罗伯特，弹丸尺寸对铝中超高速撞击坑的影响。 美国宇航局TN D-4067，1967年。 3. 库恩，保罗和菲格，I。 E.，锻铝合金的统一缺口强度分析。 美国宇航局TN D-1259，1962年。 4. Rich，Marvin和Blackman，Samuel S.，欧拉流体动力学的方法。 洛斯阿拉莫斯科学实验室，LAMS-2826，1963年。 5. 沃尔什，J。 M.，Johnson，W. E.，迪内斯，J. K.，蒂洛特森，J. H。 \& 耶茨，D。 R.，关于超高速影响理论的总结报告。 通用动力公司，通用动力公司，GS-5119，1964年。 6. 里尼，T。 D.，使用PIC WICK代码的理论超高速影响计算。 通用电气，R64SD13，1964年。 7. 比约克，R。 L.，Kreyenhagen，K. N。 \& 瓦格纳，M. H.，应用于流星体危害的冲击效应的分析研究。 冲击流体力学公司，1966年。 8. 利普曼，H。 W。 \& Roshko，A.，气体动力学元素。 约翰·威利和儿子，1957年。 9. Walkden，F.，黏性、可压缩气体的运动方程指的是任意移动坐标系。 皇家飞机机构，英国，技术。 代表 66140，1966年。 10. 拉克斯，彼得·D。 \& Wendroff，Burton，具有高精度的双曲方程的差分方案。 通信。 纯和应用。 数学。，卷。 XVII，1964年，pp. 381-398。 11. Richtmyer，Robert D. 和莫顿，K。 W.，初始值问题的差分方法。 第二版。 Interscience出版社，1967年。 12. 艾萨克森，尤金和凯勒，赫伯特，数值方法分析。 John Wiley and Sons，1966年。 13. Burstein，Samuel Z.，含有不连续性的流体动力学流动的有限差分计算。 J。 补偿。 物理，2，1967年。 14. 鲁宾，Ephraim L. \& Burstein，Samuel Z.，可压缩气体的无黏性和黏性方程的差分方法。 J。 补偿。 物理，2，1967年。 15. Pai，Shih-I，黏性流动理论。 D。 Van Nostrand Co.，纽约，1956年。 16. 蒂洛特森，J。 H.，超高速冲击状态的金属方程。 通用原子，通用动力公司，代表。 GS-3216，1962年。 17. 萨哈罗夫，A。 D.，Zaidel，R. M.，Miniev，V. N.，和Oleinik，A. G.，实验性的

超高速冲击力 CRATERING 43 对激波的稳定性和物质在高压和高温下的力学性能的调查。 苏联物理学，Doklady，卷。 9，不。 1965年6月12日，p. 1091.

\clearpage
\chapter{3 MacCormack方法-历史视角}
\section{3.1 介绍}
MacCormack方法-历史视角Ching Mao Hung1，George S. Deiwert1和Mamoru Inouye1 3.1 介绍计算流体动力学（CFD）的重大进步起源于纽约布鲁克林，Bob MacCormack于1940年龙年2月21日出生在那里。 Bob就读于公立学校，并于1961年从布鲁克林学院毕业，获得数学和物理学学士学位。 他响应了肯尼迪总统的呼吁，即“在十年结束前将一个人送上月球”，并决定加入美国宇航局。 对艾姆斯研究中心来说，幸运的是，他接到了早些时候的号召，去西部，而不是在他出生地附近的美国宇航局中心工作。 当鲍勃到达艾姆斯时，他被分配到超高速弹道射程分部，该分部不久后成为高超音速自由飞行分部。 他最初的任务是用轻型气枪研究撞击坑。 部门主管Tom Canning建议Bob使用IBM 7094计算机以数值方式研究这个问题，这一定是他真正才能的第一次认可，因为无论如何，该分支机构都是为中心超级计算机收费的。 Bob在CFD的辉煌职业生涯始于一个周末学习Fortran IV。 与此同时，他获得了M。 1967年从斯坦福大学攻读数学专业。 工程师们知道如何只用边界1 Bob MacCormack的朋友和同事来解决黏性流动问题；美国宇航局艾姆斯研究实验室，加利福尼亚州莫菲特菲尔德，94035 计算流体动力学前沿 - 2002 编辑：David A. Caughey和Mohamed M. 哈菲兹 ©2002 世界科学

\section{3.2 MacCormack方法的演变}
46 HUNG、DEIWERT和INOUYE层理论。 但鲍勃接受过数学家和物理学家的培训，更喜欢将自己视为一名研究科学家，他解决了整个Navier-Stokes方程，该方程已经存在了一个世纪，但只解决了几个简单的问题。 在鼓励这种研究的时代，鲍勃成功地开发了一种方法，以数值求解物体撞击表面的影响方程。 随后是本文将按时间顺序讨论一系列MacCormack方法，然后是对各种问题的应用选择描述。 3.2 MacCormack方法的演变 3.2.1 Mac-0 最初的MacCormack方法，我们将称之为Mac-0，是三十多年前开发的，用于解决撞击坑的不稳定Navier-Stokes 方程。 该方法是明确的，是两步预测器/校正器，在时间和空间上是二阶准确的。 它有一个与Courant-Friedrichs-Levy（CFL）条件相对应的稳定性限制，等于单位。 Mac-0方法比当时现有的方法更简单。 它不需要像Lax-Wendroff方法那样计算雅各布矩阵，它使用一个简单的非交错网格，与Richtmyer的Lax-Wendroff方法的两步版本不同。 Mac-0于1969年4月在俄亥俄州辛辛那提举行的AIAA超高速冲击会议上推出。 AIAA论文69-354[1]题为“黏度对超高速撞击坑的影响”，当时并没有引起航空或流体力学界的注意。 这需要解决空气动力学问题。 3.2.2 Mac-1 当激波对层边界层进行影响时发生的复杂相互作用是Bob选择的问题，向航空和流体力学社区展示他的方法。 在其他人尝试失败的地方，他成功地用数值解决了可压缩的Navier-Stokes 方程。 Mac-1是Mac-0的简单修改。 正如数学家常说的那样，Bob成功的原因是“显而易见的”。 首先，方法很简单。 其次，它被及时分割成两个一维操作，并且仍然保持二阶精度。 第三，它还将计算域分为内黏性区域和外黏性区域，并使用不同的操作序列来提高计算效率（如下文将讨论）。 第四，也是最不被欣赏的，偏置差异化技术被用于扩张区域和

麦考马克方法的历史视角 47 开发了用于振荡的四阶平滑，以将解决方案“粘合”在一起，使方案在接近统一的CFL条件下稳定。 将Navier-Stokes 方程的保守形式写成：Ut + Fx + Gy = 0，并遵循Strang [2]的分时概念，一维差分算子被定义为Lx算子：Ut + Fx = 0 Ly算子：Ut + Gy = 0 每个算子，Lx和Ly，都像Mac-0方法一样，在2步的预测校正子过程中求解。 对于简单的直线运算，如U\{t + dt) = LxLy U\{t），该方案只有一阶准确。 然而，如果将操作序列对称化为U(t + 2dt) = LyLxLxLy U(t)，该方案在时间和空间上都是二阶准确的。 这允许人们有进一步的变化，如U\{t + 2dt) = Ly(dt)Lx(2dt)Ly(dt) U(t)，其不同的稳定性时间步骤取决于每个运算符。 这导致有可能将计算域分割成墙附近的内黏性精细网格和外部黏性粗网格。 与精细网格中Lx运算符相比，i\textasciicircum{}-运算符的稳定时间步数非常小，相反，Lx运算符中的稳定时间步数比Ly运算符中的稳定时间步数小。 因此，人们可以通过有一个方案来提高计算效率，例如内部：U(t + 2dt) = Ly... LyLx(dt)Lx(dt)Ly... Ly U(t) Outer: U(t + 2dt) = Lx(dt)Ly(2dt)Lx(dt) U(t) 计算域中的这种分割避免了各种流动区域中不同稳定时间步骤的差异，并允许该方案处理固体边界附近的精细分辨率。 由于每个算子都是一维的，因此该方案很容易扩展到三维，因为U\{t + 2dt) = Lz(dt)Ly(dt)Lx(2dt)Ly(dt)Lz\{dt) U(t)或该序列的其他三维变体。 这就是Mac-1方法的美妙之处，它真的在CFD世界中起飞了，并取得了巨大的成功。 它最终被全世界用来解决各种问题。

48 HUNG, DEIWERT \& INOUYE Bob对CFD的主要贡献是他使用控制体概念来实现守恒律形式。 他没有直接从控制微分方程中取差异，而是认为流场局部分为受控的有限体积（或细胞），力作用于细胞面，质量、动量和能量在细胞之间传输。 由此产生的控制方程集是守恒律形式。 这个概念最终被CFD社区广泛采用和使用，并被称为“有限体积”公式。 这里值得一提的旁注。 萨瑟兰的经验公式被用来评估他的论文中的分子黏度。 由于MacCormack的成功，从那时起，萨瑟兰的配方被用于空气，甚至一些其有效性值得怀疑的应用。 萨瑟兰本人从未想过他的名字会如此频繁地被引用，并在如此多的报纸上被引用，因为一个人，鲍勃·麦科马克，从NACA报告1135中摘取了它。 Mac-1于1970年9月在加利福尼亚州伯克利举行的第二届流体力学数值方法国际会议（ICNMFD）上发表。 这篇题为《激波与层压边界层相互作用的数值解》的论文[3]，在正确的时间、正确的地点和正确的标题。 它为解决可压缩Navier-Stokes 方程-流体力学运动的控制方程打开了大门。 它引起了航空航天业人员的注意，此时他们正在寻找使用计算机来帮助他们解决复杂的流动问题的方法。 更重要的是，它引起了博士的注意。 美国宇航局艾姆斯研究中心热和气体动力学司司长迪恩·查普曼。 查普曼预见到了数值应用在流体力学中的重要性。 博士 查普曼是一个“激波”的人，他早期的大部分研究工作都与超声速流动、激波和热物理学有关。 看到这种高速问题可以通过数值模拟，他非常高兴。 由于这项工作，Bob在哈佛Lomax的指导下被选为新成立的CFD分部的助理分部主任（1970年）。 查普曼分部。 他被允许继续开发他的方法，培训其他人使用该方法，并将该方法应用于当前感兴趣的问题。 Bob被授予著名的H。 1973年朱利安·艾伦奖最佳艾姆斯科学论文奖。 3.2.3 Mac-2 随着内黏性区域时间稳定性的重要性增加（为了允许处理更复杂的几何形状、更高的雷诺数等），MacCormack进一步推进了分割的想法。 1976年[4]，开发了两个新的运营商来取代Ly运营商

MACCORMACK方法49在精细网格计算中的历史视角，如Ly(dt)-> Lyh(dt)Lyp(dt) 运算符LVh求解G的不黏性（交）项。 它是明确的，并使用特征关系来预测对流和压力场。 运算符LVp求解G中的黏性（半向）项。 它是隐含的，并使用简单的三对角线反转。 它是无条件的稳定。 精细网格中所有（i，j）的运算符序列是m是一个小整数，通常等于2的值，这是可能的，因为稳定性要求大大放宽。 这种方法在这里被称为Mac-2，是时间分割的终极用途。 它不仅将方程分成一维运算，还根据对流和耗散项划分一维算子。 如果Claude Navier或George Stokes在天堂看到这个公式，他们不会认出以他们命名的方程。 通过这种方式，方法大大加快了速度。 然而，编程变得非常复杂。 在世界各地，只有少数与MacCormack密切相关的人使用过这种方法。 有趣的是，这种方法的应用，即使只有少数人，也非常成功。 3.2.4 Mac-3 - 显式-隐式方案 由于显式方法中的时间步骤限制，在20世纪70年代末，正在开发隐式方法来提高计算效率。 通量分裂、迎风差分和总变异减小（TVD）的想法也在开发中。 1981年[5]，MacCormack将一个对角隐式程序纳入显式预测-校正方法中。 在论文中，MacCormack在发展稳态计算的数值方法方面提出了一个非常重要的概念。 他建议，数值方法的理想形式是，以delta形式书写，将数字放在一边，例如左侧，并将物理方程的精确局部近似放在另一边，例如右侧，作为残差。 左侧（数字）的责任是在不违反物理定律的情况下，以稳定的方式在全球范围内传达本地确定的解决方案变化。 为了数值效率，左侧应尽可能简单直截了当，并应尽可能快地将右侧的残差驱动为零。 三角洲形式早些时候在Beam and Warming的隐式方法中引入。 这是

\section{3.3 应用程序}
50 HUNG, DEIWERT \& INOUYE MacCormack，他们让CFDer理解和遵循得如此清晰和容易。 Navier-Stokes 方程仍然像以前一样被划分为本地一维运算符。 正向和向后差在隐式运算符中的通量项的局部线性化矩阵中使用。 预测器步进的正向差异和校正器的向后差异导致了隐式程序的简单对角矩阵反转。 两个（预测器和校正器）标量双边形矩阵的组合有效地产生了一个对角线显性矩阵算子。 随着隐式程序的加入，该方法在理论上在任何时间步骤上都是稳定的。 它不需要标量或块状三对角矩阵反转。 因此，它非常高效，速度加快了大约两个数量级。 这是“预测器-校正器”MacCormack方法的顶峰，此处称为Mac-3。 在这篇论文之后，鲍勃离开艾姆斯前往华盛顿大学，成为年轻学生的导师，并继续传播CFD的种子。 此时，随着数值方法和计算机能力的不断进步，CFD已成为流体力学的一个分支。 各种隐式方案，加上通量分裂，本地时间推进已经出现。 就在那时，用户开始常规地解决复杂的真实几何形状的Navier-Stokes 方程。 MacCormack所做的贡献的重要性仍然在制定后续数值方案的各个方面。 3.3 应用 3.3.1 Inviscid 3.3.1.1 Mac-0，三维，超音速 尽管MacCormack方法最重要的影响是解决Navier-Stokes 方程，但inviscid 钝体和超音速太空分流解决方案的早期即时应用首先引起了人们对他的方法的关注。 在20世纪70年代初，航天飞机是该机构的“太空项目”。 Rizzi和Inouye [6]在三维钝体上将该方法应用于超声速流动。 库特勒[7]用Mac-0方法取代了他之前的非中心方案，在航天飞机配置上计算无黏流动，后来还处理了其他几个超声速流动问题。

MACCORMACK方法的历史视角51 3.3.1.2 Mac-0、跨声速流动、亚音速边界条件 Deiwert（1974）[8]，在开发使用Mac-1方法研究跨声速流动超过翼型配置的代码时，首先在双凸翼型形状上获得了空气、氟利昂和低温氮中的非黏性跨声速流动的结果。 对在高雷诺数下进行风洞测试的担忧越来越大，这促使人们考虑使用空气以外的测试气体，以增加气体密度，同时保持可管理的停滞压力水平。 增加密度的一个方法是使用高分子量的测试气体，如氟利昂12。 另一个是通过显著降低气体温度，例如使用低温氮气。 当时，该机构正在制定建造国家跨声速风洞的计划，该风洞将使用低温氮气作为测试气体。 在前一种情况下（氟利昂），可以将气体视为理想，但其等熵指数（伽马）与空气有很大不同。 在第二种情况下（低温氮），气体并不总是以理想的方式表现，膨胀区域可能非常接近两相区域。 被问到的问题之一是，低温超音速风洞中是否会出现条件，从而导致液化。 使用范德华状态方程来描述热和热量不完美的低温氮，并对流马赫数为0.775和0.95的18\%双凸圆弧的流动进行了模拟。 该研究结果于1974年3月在洛杉矶举行的超音速风洞协会第41届半年会上发表，结果表明，在模拟条件下无法预测液化，事实上，低温氮的使用似乎是可行的。 研究中的其他信息量化了流场结构上不同气体的等熟指数（γ）的一些差异。 当时，研究人员正在使用“有效伽马”的概念来“匹配”他们的黏性解与实验观察到的升力和阻力结果。 这与“有效攻击角度”一起，实际上是解释现实世界实验中发生的黏性位移现象的人为方法。 3.3.2 黏性跨声速流3.3.2.1 Mac-1，广义曲线坐标，对称Deiwert（1974）[9]用时间拆分扩展了基本的MacCormack 显式方法（Mac-1），以处理任意配置的非直角计算网格，以应用于黏性流动一般曲线形状的过去物体。 目标是通过翼型配置“捕捉”冲击，并模拟其与边界层的相互作用。 博士 迪安·查普曼（Dean Chapman）是最早认识到强大潜力的人之一

52 HUNG、DEIWERT和INOUYE将MacCormack方法应用于此类实际应用，确定了这种特别重要的应用。 事实上，是博士。 查普曼问过博士。 Deiwert调查了这个问题领域，甚至在Bob的办公室旁边提供了办公空间，以帮助促进研究。 Deiwert当时又年轻又天真，没有意识到如此复杂的应用程序是闻所未闻的，他欣然同意承担这项任务，因此成为Bob的第一批“学生”之一。 为这项研究选择的配置是18\%厚的双凸圆形弧翼型形状。 开发了用于区分黏性项的坐标变换，并实现了可压缩的湍流运输模型。 此配置的边界条件都是亚音速的。 发现了一个跨声速马赫数，该马赫数将产生足够强的冲击力，足以在冲击的脚下诱导流动分离。 还启动了一个配套的实验计划，通过Ames High 雷诺数通道中的此类配置获取详细数据。 1974年6月，在加利福尼亚州帕洛阿尔托举行的AIAA第七届流体和等离子动力学会议和同月在科罗拉多州博尔德举行的第四届流体动力学数值方法国际会议上，首次研究结果显示了黏性/不黏性相互作用与冲击诱导分离[10]。 开发了一个广义变换，将Lx和Ly运算映射到黏性流动区域的广义非正交网格上。 一个简单的混合长度模型被用来模拟湍流边界层中的黏性运输。 按照当时的标准，计算资源被认为是相当大的。 使用 Mac-1 方法，在 CDC 7600 电脑（当时最先进的）上需要 2 到 10 小时的解决方案。 当结果在第四次国际会议上发表时，观众提出了关于计算机要求的问题。 当答案是CDC 7600需要2到10个小时才能找到融合解决方案时，会议主席博士。 苏联莫斯科科学院计算中心主任Belotserkovskii对演讲者说：“你一定很富有。” 计算时间取决于雷诺数（因此，网格分辨率）。 代码随后被矢量化，使计算机时间不到一半。 该代码也是在实验性Illiac IV计算机（64处理器并行计算机）上编写和运行的，每个解决方案的计算时间从0.6小时到3小时不等。 事实上，这些解决方案是Illiac IV上获得的第一个已发表的解决方案。 Illiac以11.5兆赫的频率运行。 如今，在以数百兆赫运行的台式或笔记本电脑上，可以在更短的时间内完成相同的计算。 跨声速双凸翼型研究继续开发和进一步评估代数湍流运输模型，包括Shang和Hankey以及Baldwin和Rose提出的，适用于分离流。 这些结果与获得的新实验资料进行了比较

McDevitt和Levy在Ames High 雷诺数频道上对MACCORMACK方法53的历史视角，并于1975年6月在康涅狄格州哈特福德的AIAA第8届流体和等离子体动力学报告[11]。 对跨声速流动的实验研究的一个意想不到的结果是，在翼型的后部区域观察到周期性非定常流动，以进行小型选择马赫数范围（见图。 13，参考。 11）。 计算机代码中强加的对称边界条件，以及Mac-1方法的计算时间约束，排除了这些现象的模拟。 3.3.2.2 Mac-2，提升机翼，自适应网格鲍勃对他的方法的改进，显式/隐式概念（Mac-2），旨在提高代码的速度和计算效率。 实现了95\%的速度提高，同时仍然保持非黏性时间尺度的时间精度（第五届流体力学数值方法国际会议，荷兰恩斯赫德，1976年6月[4]。） 这为消除对称约束、考虑提升翼型配置，并开始解决一些不稳定问题开辟了道路，如自助餐和在高雷诺数通道中观察到的不稳定现象。 Deiwert的代码被修改，以处理近唤醒区域的不对称行为。 以希夫开发的方式实施了网状适应程序，以跟随尾流中的剪切流，并跟随激波，后者在非定常流动中时间移动。 这些增强功能大大提高了代码模拟实际跨声速翼型流的能力。 1976年2月，Deiwert在SQUID起重机翼研讨会上报告了首次对起重机翼的处理[12]。 重大进步是对不对称翼型形状和近唤醒区域的处理，以及使用Mac-2方法大大提高了计算效率。 当时一个特别有趣的配置是Richard Whitcomb提出的超临界配置。 当时，超临界机翼引起了很多人的关注。 特别值得注意的是Garabedian和Korn的分析。 在这些无黏性研究中，“有效攻击角度”的概念被用来将计算结果与实验数据进行匹配。 Deiwert [13]的研究结果确凿地表明，对黏性现象（即边界层位移效应）的适当解释对于准确模拟这些翼型形状的性能是必要的和充分的（见图 2和4，参考。 13）。

54 HUNG, DEI WERT \& INOUYE 3.3.2.3 Mac-2, Unsteady Transonic Flows Levy [14]使用修订后的Deiwert代码和最新的MacCormack（Mac-2）方法和不对称唤醒处理来研究在18\%圆弧翼型的实验研究中观察到的非稳定过程。 在计算和实验之间发现了显著的一致性，无论是在非稳定过程的振幅还是频率上，从而使时间精确的MacCormack方法的力量更加可信。 Levy最终（1979年）被授予H。 这项研究获得了朱利安·艾伦奖。 Levy和Bailey[15]以及Deiwert和Bailey[16]随后的研究进一步深入研究了非定常流动模拟与时间准确的MacCormack方法的适用性，并研究了自助餐现象和副杆嗡嗡声现象。 这种方法的应用被发现对具有单一主导频率的流动非常有效，其时间尺度为非黏性时间的顺序。 1976年，Deiwert与教授合作进行研究。 Peter Bradshaw开发了一个单方程剪切模型，该模型在近唤醒研究的代码中实现。 处理尾流和非稳流所需的动态网格自适应问题最终导致了Nakahashi和Deiwert开发的动态自适应网格方案。 Horstman在研究后缘流时使用了Deiwert的代码，并研究了各种湍流运输模型。 Rose还将该代码用于他的几个交互式流程研究。 所有使用此代码的人的评论普遍一致，即“代码很强大，总是给出正确的答案。” 事实上，MacCormack方法一直非常稳健；代码中的关键起搏项目是，并且仍然是湍流运输模型。 3.3.3 黏性超声速流动 3.3.3.1 二维超声速流动 在1970年Mac-1发布后，求解N-S方程的挑战开始升温，活动向前冲。 在Deiwert的黏性跨声速流动努力的同时，1972年，Baldwin使用MacCormack的代码，并添加了涡黏性项来解决湍流冲击反射问题，并于1974年夏天在帕洛阿尔托发表了一篇AIAA论文[17]。 Hung于1973年加入CFD分部，担任NRC博士后助理，并开始修改MacCormack的代码，以在压缩角研究层流。 另外两个小组也在同一条小路上。 一个是在Hankey主管的Wright-Patterson AFB，Shang正在研究湍流压缩坡道。 另一个是Marvin（也在Chapman's Division）领导的Ames的一个实验小组，为开发和验证高速湍流运输模型开发有据可查的实验数据。

MACCORMACK方法的历史视角 55 除了Deiwert、Levy和McDevitt研究的跨声速数据外，Horstman和Kussoy还进行了实验，Coakley在高超音速下进行了轴对称冲击反射问题的计算。 这些导致了1975年在帕萨迪纳举行的AIAA航空航天科学会议上大张旗鼓。 第一届会议的前四篇论文，MacCormack和Baldwin的AIAA论文75-01[18]，Hung和MacCormack的75-02[19]，Shang和Hankey的75-03[20]，以及Horstman、Kussoy、Coakley、Rubesin和Marvin的75-04 [21]，都使用MacCormack的方法（Mac-1）来解决休克诱导的分离问题。 从那时起，激波诱导分离或激波/边界层相互作用的理论研究属于数值模拟，在很长一段时间里，属于MacCormack的方法。 正如大家所意识到的，大多数实际感兴趣的流程是湍流，而数值模拟的起搏项目是湍流模型。 NASA-Ames投入了大量资源来开发这些模型。 Baldwin和Coakley的结果是合理的，但不是很好。 令人惊讶的是，Shang和Hankey使用了一个简单的松弛湍流模型，结果不仅表明表面压力和分离位置的一致性非常好，而且更显著的是，几个上游和下游位置的密度和速度曲线非常一致（见图 参考的9和10。 20）。 根据他们的结果，人们几乎可以声称湍流问题已经解决了。 这在艾姆斯引起了很大的兴趣，并导致小组的形成，研究“放松”模型。 不幸的是，没有人能获得像Shang获得的好结果。 Hung和MacCormack在1976年夏天发表的一篇论文是对放松模型的研究的结果。 最常见的发现是Shang建议的放松长度太长了。 直到一段时间后，Hung才在Shang的代码中报告了两个程序“错误”。 第一个也是最严重的问题是，在搜索边界层外速来计算Cebeci-Smith模型所需的位移厚度时，由于激波的存在，将获得最大速度，而不是预期的边缘速度。 边缘速度的微小差异可能会导致计算的位移厚度的巨大差异，从而导致其相应的外层涡黏性的巨大差异。 这个“错误”使模型对他们的问题非常有效。 另一个错误对涡黏性的计算产生了轻微影响，这里没有描述。 让我们遵循詹姆森的规则，（引自MacCormack的论文[22]）“在任何由至少一盒IBM卡组成的程序（现代翻译：2000个语句长）中，总有一个错误。”

56 HUNG, DEI WERT \& INOUYE 3.3.3.2 三维超音速流 Mac-2的开发使3D Navier-Stokes 方程的解决方案变得可行。 Shang首先使用Mac-1方案进行高超音速层流的3D压缩角模拟，然后切换到湍流的Mac-2方案。 Hung将Mac-2方案应用于超音速层流和湍流流的3-D压缩角，应用于在攻角处具有耀斑的轴对称体和各种3D问题，例如斜激波在气缸上的碰撞。 1981年Mac-3开发后，Kordulla将其应用于二维超音速翼型流，Hung和Kordulla将其扩展到一般的3D几何，并将其应用于平板上的钝翅片。 他们模拟了钝鳍前存在马蹄涡（见Hung和Kordulla [23]），并与普林斯顿大学Bogdonoff小组获得的实验数据非常一致。 1984年，Hung和Buning[24]在AIAA航空航天科学会议上放映了一部模拟电影。 从那时起，在一段时间内，流场模拟电影变得非常流行。 当时，艾姆斯每年制作十到二十部CFD电影，几乎像“好莱坞-北方”，CFD开玩笑地代表“彩色电影展示”。 MacCormack方法最成功的用户之一是Shang。 他利用了Mac-0和Mac-l的简单性的计算机架构，模拟了许多流程问题，2D和3D。 他甚至将该方法应用于圆柱体周围的二维流动振荡和尖峰体上的三维非定常流动。 他是第一个进行完整飞机模拟的人，X24C-10D计算。 MacCormack方法已被用于解决许多超音速和高超音速问题，并征服那些复杂的三维激波和边界层相互作用，这些相互作用是理论家从未梦想过的，实验家只能进行非常有限的表面测量。 与3D流程的实验数据的一致性通常比2D案例好得多。 有三个重要原因。 首先，除了计算时间和数据内存之外，三维问题比二维问题更容易。 与二维计算相比，在三维计算中更容易与实验数据达成“良好的一致性”。 原因很明显；大多数3-D问题都有一个额外的扰动缓解维度，因此由无黏性机制主导，与黏性机制相比，这种机制更容易解决。 Knight和Horstman [25]在3D扫荡冲击计算中表明，即使涡黏性值在许多地方的差异高达14倍，两种不同的计算仍然非常一致。 这并不意味着湍流建模不重要。 这只是意味着3D问题有更多的错误空间。 接下来，虽然激波和分离给理论家带来了很多问题，但相反，超音速和高超音速问题更容易用数值求解

\section{3.4 结束语}
麦考马克方法57的历史视角比不可压缩或跨声速问题。 干扰只向一个方向（下游）传播，速度很快。 边界条件易于设置，解决方案在相对较短的计算时间内收敛。 最后一个原因是，尽管压缩冲击会造成更多的工程问题，因为它们的不利压力梯度会导致边界层分离和空气动力学加热增加，但压缩流比膨胀流更容易解决。 扩展有机会耗尽计算单元，导致负密度或压力，并使方案不稳定。 同样，这些幸运的原因促成了MacCormack方法的成功应用。 3.4 闭幕词 回顾历史，正是MacCormack方法是解决前沿问题的火炬手，并为CFD的发展铺平了道路。 该方法的成功不仅可以归因于其力量和稳健性，还可以归功于Bob愿意帮助其他人实施该方法来解决问题。 随着数值方法和计算能力的不断进步，CFD已成为流体力学的一个独立分支。 使用计算机作为工具，CFD现在不仅能够模拟真实的工业工程问题，而且能够补充实验和理论研究。 它可用于对其自身优点进行研究，计算流体动力学的发展和进一步推进可以被视为物理科学的一个独立领域。 MacCormack的方法在CFD的诞生、生长和发展中发挥了重要的指导作用。 1981年，美国宇航局授予杰出科学成就奖章，表彰鲍勃的贡献。 Bob离开Ames前往学术界，并继续传播CFD的种子，先是在华盛顿大学，现在在斯坦福大学。 作者试图总结MacCormack方法在早期的发展，并强调这些方法的应用所取得的一些重大成就。 对方法的全面描述及其应用的完整文档超出了本文的范围。 虽然已经尽一切努力确保历史和技术的准确性，但我们认识到可能存在一些不准确之处；对此，我们深表歉意。

\section{参考资料}
58 HUNG, DEI WERT \& INOUYE 参考资料1。 麦科马克，R. W.：黏度对超高速撞击坑洞的影响，AIAA论文69-354，AIAA超高速撞击会议，俄亥俄州辛辛那提，4月。 30 - 1969年5月2日，（编者注：本文转载为本卷第2章。） 2. 斯特朗，G.：关于差异方案的构建和比较，SIAM J. 数字。 肛门。，卷。 5，1968年，第506-517页。 3. 麦科马克，R. W.：激波与层边界层相互作用的数值解，物理学讲义，卷。 8，Springer-Verlag，纽约，1971年，p. 151. 4. 麦科马克，R. W.：双曲系统或方程的快速求解器，物理学讲义，卷。 59，A。 I. van de Vooren \& P. J。 Zandbergen，编辑，Springer- Verlag，纽约，pp. 307-317，1976年。 5. 麦科马克，R. W.：求解可压缩黏性流动方程的数值方法，AIAA论文81-110，AIAA第19届航空航天科学会议，St. 路易斯，密苏里州，1月。 1981年12-15日。 6. 里兹·A。 W。 \& Inouye，M.：3D钝体流动的分时有限体积法，AIAA论文73-133，AIAA第11届航空航天科学会议，华盛顿特区。 7. 库特勒，P.，Lomax，H. \& Warming，R.F.：使用非中心有限差分方案计算航天飞机流场，AIAA论文72-193，AIAA第十届航空航天科学会议，加利福尼亚州圣地亚哥，1月。 17-19，1972年。 8. Deiwert，G。 S.：空气、氟利昂和低温氮气中的跨声速流动模拟。 第41届半年度超音速风洞协会，洛克威尔国际，加利福尼亚州洛杉矶，3月。 28-29，1974 9. Deiwert，G。 S.：高雷诺数超声速流动的数值模拟，AIAA J.，卷。 13，pp. 1354-1359，1975年。 （另外，AIAA论文74-603，在1974年6月17日至19日在帕洛阿尔托举行的AIAA第七届流体和等离子动力学会议上发表。） 10. Deiwert，G。 S.：高雷诺数 跨声速流动模拟，Proc。 第4国际机场 关于Num的Conf. 冰毒在流体中，科罗拉多州博尔德，1974年7月11日。 麦克德维特，J. B.，Levy L. L.，和Deiwert，G. S.：跨声速流动关于厚圆弧翼型，AIAA J.，卷。 14，pp. 606-613，1976年。 12. Deiwert，G。 S.：关于超声速流动中黏性现象的预测，在涡轮机械中的跨声速流动问题中，Adamson，T. C. \& Platzer，M。 F. eds.，Hemisphere Publishing Corp.，pp. 371 391，1977年。 13. Deiwert，G。 S.：跨声速流动中黏性效应的最新计算，在物理学讲义中，卷。 59，A。 I. van de Vooren \& F. J。 Zandbergen，编辑，Springer-Verlag，pp. 158-164，1976年。 14. 利维，L。 L.：关于厚翼型的实验和计算稳定和非稳定跨声速流动，AIAA J.，卷。 16，pp. 564-570，1978年。 15. 利维，L。 L。 \& 贝利，H. E.：翼型 Buffet边界的计算，AIAA J.，卷。 19，pp. 1488-90，1981年。 16. Deiwert，G。 S。 \& 贝利，H. E.：非稳态交互流的时间依赖有限差分模拟，在空气动力学流动的数值和物理方面II，T。 Cebeci，编辑，Springer-Verlag，1983年。 17. 鲍德温，B.S. \& MacCormack，R.W.：强激波与高超音速湍流边界层相互作用的数值解，AIAA论文74-558，AIAA第七届流体和等离子动力学会议，加利福尼亚州帕洛阿尔托，1974年6月17日至19日。 18. 麦科马克，R. W。 和鲍德温，B。 S.：求解的数值方法

MACCORMACK方法59 Navier-Stokes 方程的历史视角，应用于冲击-边界层相互作用，AIAA论文75-01，1月。 1975年。 19. 洪，CM。 \& MacCormack，R.W.：超音速和高音速层流在二维压缩角上的数值解，AIAA论文75-02，AIAA第13届航空航天科学会议，加利福尼亚州帕萨迪纳，1月。 20-22，1975年。 20. 尚，J.S. \& Hankey，W.L.Jr.：压缩坡道上超音速湍流的Navier-Stokes 方程的数值解，AIAA论文75-03，AIAA第13届航空航天科学会议，加利福尼亚州帕萨迪纳，1月。 20-22，1975年。 21. Horstman，C.C.，Kussoy，M.I.，Coakley，T.J.，Rubesin，M.W.和Marvin，J.G.：激波诱导的高超音速湍流边界层分离，AIAA论文75-41 AIAA第13届航空航天科学会议，加利福尼亚州帕萨迪纳，1975年1月20日至22日。 22. MacCormack，R.W.：CFD研究四分之一中心的观点，AIAA论文93-3291，AIAA第11届计算流体动力学会议，奥兰多，1993年7月6日至9日。 23. 洪，CM。 \& Kordulla，W..：用于3-D流场模拟的时间分割有限体积算法，AIAA论文83-1957，AIAA第22届航空航天科学会议，内华达州里诺，1月。 9-12，1984年。 24. 洪，CM。 \& Buning，P.G.：钝鳍诱导激波湍流边界层相互作用的模拟，AIAA论文84-457，AIAA第六届计算流体动力学会议，马萨诸塞州丹佛斯，1983年7月13-15日。 25. 骑士，D。 D.，Horstman，C. C，Shapey，B.和Bogdonoff，S.：超音速湍流通过尖锐鳍的结构，AIAA论文86-343，AIAA第24届航空航天科学会议，内华达州里诺，1月。 6-9，1986年。

\clearpage
\chapter{4 实现教科书多网格效率的一般框架：一维欧拉示例}
\section{摘要}
\section{简单介绍}
实现教科书多网格效率的一般框架：一维欧拉示例James L. Thomas1、Boris Diskin2、Achi Brandt3和Jerry C. South, Jr.4 摘要讨论了一个通用的多网格框架，以守恒律形式获得可压缩欧拉和Navier-Stokes 方程解的教科书效率。 一般方法依赖于分布式松弛程序来减少规则（平滑变化）流动区域的误差；对每个因素（椭圆和/或双曲）进行单独和不同的处理，以实现最佳的误差减少。 接近边界和不连续性（震感），保守方程的额外局部放松是必要的。 计算示例是针对准一维Euler 方程进行的；计算说明了一般程序。 导言计算流体动力学（CFD）已成为飞机设计周期的组成部分，因为更快的计算机可用1美国宇航局兰利研究中心，弗吉尼亚州汉普顿23681 2 ICASE，美国宇航局兰利研究中心，汉普顿，弗吉尼亚州23681 3魏茨曼科学研究所，雷霍沃特76100，以色列4威廉斯堡，弗吉尼亚州23185前沿计算流体动力学-2002编辑：David A. Caughey和Mohamed M. 哈菲兹 ©2002 世界科学

62 THOMAS、DISKIN、BRANDT和SOUTH 更多的内存和改进的数值算法和物理模型。 如果能够为设计外性能设计出可靠的方法，则可能产生更大的影响，通常与非稳定、分离、涡流与强激波有关。 此类计算需要比目前可用的更多的计算资源。 当前具有多网格算法的雷诺平均Navier-Stokes（RANS）求解器需要1500个剩余评估才能将升力和阻力收敛到最终值的百分之一，即使在跨声速巡航条件下相对简单的机翼体几何形状也是如此。 众所周知，椭圆问题可以通过使用全多网格（FMG）过程在更少（2到4次）的剩余评估中实现最佳解决方案。 Brandt [1, 2, 3]将多重网格法定义为具有教科书的多网格效率（TME），如果在离散方程系统中运算计数的小（小于10）倍的计算工作中实现了控制方程组的解。 因此，如果RANS方程可以实现TME，操作计数可能会减少几个数量级。 用于大规模可压缩流动应用的最先进的多网格方法使用块矩阵松弛和/或伪时间依赖方法来求解方程；多网格方法已经证明了重大改进，但这些方法不是最佳收敛的。 RANS方程集是耦合非线性方程组，即使对于亚音速马赫数，这些方程组也不是完全椭圆的，但包含双曲分割。 勃兰特的分布式松弛方法[1, 2]将方程组分解为可分离的，通常是标量，可以用最佳方法处理的因子。 几年前，开始进行一项调查，将这种方法扩展到大规模应用；当时，已经完成了几个不可压缩模拟的TME演示，Ta'asan为亚音速Euler 方程[13]展示了有希望的结果。 在将方法扩展到黏性可压缩流动应用[12]和使用紧凑微分方案可压缩Euler 方程方面取得了进展[15]。 已经进行了进一步的不可压缩流动应用，包括复杂的几何形状[14]和二维[17]和三维[18]的高雷诺数黏性流动。 Brandt总结了TME中流体动力学方程的进展和剩余障碍[3]。 本文的目的是讨论大规模可压缩流动应用所需的一般框架。 求解准一维Euler 方程来说明框架。 显示了全亚音速和超音速应用，以及具有捕获冲击的跨声速应用。

\section{总体框架}
教科书多网格效率的一般框架 63 一般框架 质量、动量和能量的时间依赖守恒的黏性可压缩方程可以写成dtQ + R = 0, (4.1)，其中守恒变量是Q = (p,pu,pv,pw,pE)T，表示单位体积的密度、动量向量和总能量，R(Q)是表示对流、黏性和热传递效果的向量函数的空间发散。 一般来说，微分方程的最简单形式对应于用原始变量表示的非保守方程，这里取而代之，由速度、压力和内能组成的集合，q -（u，v，w，p，e）T。 通过将与时间相关的守恒方程转换为与时间相关的原始变量方程，很容易找到这些方程。 同样，可以推导出一组非保守校正方程，其右侧向量由保守残差项的组合组成，给出为L*q=-\textasciicircum{}R。 （4.2）在等式中。 （4.2），-\textasciicircum{}是变换和校正<5q = q"+1 - q™的Jacobian 矩阵，其中n是迭代计数器。 对于稳态方程，时间导数被丢弃。 在离散水平上，校正方程（4.2）的右侧是保守离散化残差的组合，而左侧是非保守算子的主要线性化。 请注意，这不是牛顿线性化；只有R线性化中的主项被保留。 主要项是那些对q的每单位变化的残差做出重大贡献的项。 因此，主要术语通常取决于感兴趣的尺度或网格大小。 对于标量方程，离散的最高导数项在网格大小足够小的网格上是主要的。 在推导高雷诺数模拟的主要线性化时，必须考虑非黏性和黏性尺度--非黏性尺度在大部分流场中占主导地位，黏性尺度在物体附近及其尾部的薄黏性层中很重要。 请注意，对于离散化微分方程组，如R = 0，主项是那些有助于矩阵算子\textasciicircum{}行列式的主项的项。 L中主项的系数是从电流近似值中评估的。 主线性化应用于基于非保守近似的校正方程。 因此，我们期望校正在远离流场中的不连续性（休克、滑线）时是好的。 它在这些

64 THOMAS、DISKIN、BRANDT和SOUTH 规则（平滑变化）流动区域，我们应用分散式松弛。 分布式松弛方法用M<5w替换Sq，使生成的矩阵L M变成对角线或下三角矩阵，因为LM<5w = -§[R。 （4.3）LM的对角元素理想地由矩阵L的行列式的可分离分量组成，并表示方程的椭圆或双曲因子。 Brandt将5w变量称为“幽灵变量”，因为它们不需要在计算中明确出现。 如果用快速方法求解L M中的组成标量对角线算子，分布式松弛方法可以为稳态和非稳态模拟产生快速收敛。 该方法可以应用于非常通用的方程；Brandt推导出了一组矩阵M，为流体动力学的可压缩和不可压缩方程（包括可变状态方程）提供了方便的下三角形式[2]。 对于可压缩的Euler 方程，构成L M主对角线的标量因子是对流和全电位算子。 前者的有效求解器可以基于下游行进，并具有用于再循环流动的额外特殊程序[6, 7, 9]；后者是一个可变类型的运算符，其解决方案需要在亚音速、超音速和超音速区域使用不同的程序。 在深亚音速区域，全电位算子是均匀的椭圆的，因此标准多网格方法产生了最佳效率。 当马赫数接近单位时，运算符变得越来越各向异性，由于一些平滑特性误差分量无法在粗网格上充分近似，经典的多网格方法严重退化。 特征成分是那些在特征方向上比其他方向平滑得多的成分[3，10，11]。 在深超音区域，全电位算子是均匀的双曲的，流方向作为类似时间的方向。 在这个区域，可以通过下游行进程序获得高效的求解器。 然而，这个程序对马赫数下降到统一来说是个问题，因为与下游行进程序相关的Courant 数很大。 因此，需要一个特殊的程序来为跨声速区域提供有效的解决方案。 这个局部程序[4，5，8]基于分段半粗化和一些在粗网格水平上增加耗散的规则。 边界在分布式放松中引入了一些额外的复杂性。 LM的行列式通常比L的行列式更高。 因此，作为一组新变量，5w通常需要额外的边界条件。 在放松中，由于幽灵变量可以添加到域的外部部分，通常可以确定满足原始边界的5w的合适边界条件

\section{准一维方程}
教科书多网格效率的一般框架 原始变量的65个条件。 Thomas、Diskin和Brandt[17]在不可压缩流动中举例说明了进入和防滑边界。 然而，总体而言，构建这种补救措施可能很困难和/或耗时。 此外，强制执行这些边界条件会导致松弛方程在边界附近耦合，而不是像域内部那样解耦。 因此，在边界和不连续性附近，一般方法[2，3]是在原始变量方面直接放松控制方程。 在这个区域可以进行几次强大的（但可能缓慢收敛）的松弛，如牛顿-卡克马克斯松弛。 额外的扫描不会影响整体复杂性，因为与内部点的数量相比，边界和/或不连续点的数量通常可以忽略不计。 下面使用这个一般框架来解决全亚音速、全超音速和超音速（有和不带冲击）流动模式下的准一维Euler 方程。 常规流动区域通过分布式放松而放松。 边界和冲击是通过应用局部松弛来处理的--这里对应于通过保守方程的近似牛顿线性化的直接解来更新。 在所有情况下，都使用FMG算法；在每个级别，方程都使用FV（2,1）全近似方案（FAS）[1，17]多网格循环求解。 尽可能在循环中使用六个级别。 准一维方程 准一维方程将质量、动量和总能量的守恒表示为3t(Q) + R = 0, (4.4) 其中Q = Q<T = \{p,(m,PE)T<r, (4.5) R = 0x(F<r) + S, (4.6) and a = cr(x)是面积项。 通量F和源项S被定义为F=l pu2+p, (4.7) \textbackslash\{\} puE + up J

\section{放松计划}
THOMAS、DISKIN、BRANDT和SOUTH（4.1压力p、内部（热）能e和声速c通过状态方程相关，如（4.9）（4.10），7是比热的比率。 对于一阶或二阶精度，R的离散保守上风偏置微分近似可以定义为R, = -[(Fa)j+i - (Fa) p e (7-l E-u- IP/p, )pe, -, - h\textasciicircum{}u>j+ - (0,Pj,0)T[aj+i -<7.,\_i]。 （4.12）在这里，使用了有限体积离散，其中下标j + \textbackslash\{\}和j - \textbackslash\{\}分别表示以位置j为中心的单元格的右和左界面，x方向的单元格间距是h。 Roe [20]的通量差分裂用于构建界面通量F+i；附录一描述了相关细节。 我们只考虑一阶或二阶精度，不区分Q的平均值和点数值。 面积分布定义为<J(X) - 1 - 0.8a;(l - x)。 对于下面提出的所有结果，我们从位于域0 < x < 1之外的细胞中心的微分问题的精确解中高估了边界值。 放松方案 下面考虑了几个方案来放松稳态方程Rj = 0。 （4.13）这些方案都是以delta形式编写的，因此对解的校正在每次迭代时都得到解决，表示为£Q = Q“+1 - Q”。

教科书多网格效率的一般框架67保守松弛方案离散保守方程的精确牛顿线性化通常会导致一个过于复杂的线性系统需要求解。 在实践中，准牛顿（近似）线性化用于降低构造的代数复杂性和/或生成的线性系统的带宽。 为了构建保守的松弛方案，运算符（4.12）的近似但保守的线性化可以写成[*-A+ + <J+A-](Q)。 （4.14）运算符5\textasciitilde{}和<5+分别表示向后和正向一阶准确差分算子，A+和A-表示正负特征值对相似性矩阵的贡献（见附录一）。 保守松弛更新SQ是从[8- A+ + 5+ A-](6Q) = --Rj计算的。 （4.15）等式右侧的残差。 （4.15）是针对目标离散化（4.13）计算的。 松弛程序要求直接求解三对角方程：（-A+。 JA\textasciicircum{}AT\textasciicircum{}Q) = -AR.. （4.16）Barth [19]用定点分析分析了这种松弛，并表明这种松弛几乎和全牛顿迭代一样好。 放松不足参数，例如伪时间步进，可以在松弛方案中引入，以提高稳健性和稳定性。 在本次调查过程中进行的广泛计算证实，这种放松方法可以在整个领域中使用，包括冲击和边界。 在本文中描述和使用的一般框架中，这种放宽只用于局部地减少边界和冲击附近的残差，并求解最粗的网格方程。 非保守放松方案 非保守放松方案来自守恒格式，Eq。 （4.15），假设雅各布矩阵是局部常数，[A+S- + A-6+]\{6Q) = --R,, (4.17)

\section{分布式放松}
68 THOMAS, DISKIN, BRANDT \& SOUTH 这个公式有望很好地近似等式。 （4.15），除了接近不连续的。 相应的三对角方程是（-A+.lAk.ATWQ）= -AR.. （4.18）减去两个三对角方程，（4.16）和（4.18）得出（-At\_！ + A+,0, AJ+1 - AT)(<JQ) = 0, (4.19) 其中系数- -Aj\_1 + A\textasciicircum{}和A\textasciitilde{}+1 - A\textasciitilde{} -在规则流动区域足够细的网格上小0(h)。 近冲击，方程中使用的非保守线性化。 （4.17）不满足秩序属性，非保守的放松方案将无效。 通过保守和非保守放宽方案的计算，更新了Eqs，这些期望得到了证实。 （4.16）和（4.18）得到了精确的解决。 Barth [19]用定点分析分析了类似的非保守性松弛方案，并表明非保守性方案对带冲击的流动表现不佳。 分布式放松远离边界和冲击，线性更新方案Eq。 （4.17）通过分散的放松而放松。 转换Eq中的校正。 （4.17）到原始变量q = (u,p, e)T给出[A+S- + A-\#](«kO = -(\textasciicircum{}. IR,., (4.20) 带有由保守残差项组合组成的右侧向量。 雅各布矩阵通过与保守的雅各布矩阵的相似性变换相关：或者，根据定义，原始变量校正方程是L8q=-rj，（4.22），其中dq J. fj = S)r\%（4-23）矩阵运算符L的元素在附录二中给出。

教科书多网格效率的一般框架69如附录一和附录二所示，深亚音速或超音速区域中（u > 0）的L的行列式是det（L）= [\{u2-c2）5xx]u S\textasciitilde{}，（4.24），其中6 = \{ S* 5*完全亚声速流动，。 \textasciicircum{}\textasciicircum{} XX \textasciitilde{} I \&x \$x 完全超声速流动。 \textasciicircum{} ' ' 方程中的第一个项。 （4.24）表示全电位算子的近似值，第二个项表示熵对流的上风近似值。 在亚声速流动中，5XX是与全电位算子的椭圆性相关的中心离散化。 在完全超声速流动中，根据全电位算子的双曲性质，6XX是一个单边或上风偏置的近似值。 对于一阶迎风差分，全电位因子是3点算子；对于二阶上风偏置微分，因子是7点算子（见附录II）。 在这个一维的情况下，矩阵运算符L中的前两个方程与第三个方程解耦。 因此，我们只需要考虑前两个方程的分布式松弛；一旦从分布式松弛中找到5u和5p，第三个方程就可以以原始变量形式求解Se。 表示L的上2x2块为L，M的一个明显选择是L的协因子矩阵，它给出了LM的对角矩阵，沿对角线全势因子，即„2 (u2 - c2) Sxx 0 LM=( " o'" („»->)*„ I" \textasciicircum{} 边界处理 正如Sidilkover [15]所指出的，分布松弛和特征变量之间存在联系。 对于亚声速流动，对LM中运算符的Jacobi更新等同于特征方程的Kacmarcz放松，根据修正的特征组合，如（u + c）5\textasciitilde{}（pcSu + 5p）\textbackslash\{\} \_ / peri +f2 \textbackslash\{\} \textasciicircum{} 2?，（u - c）6+（pc5u - Sp）J \textbackslash\{\} pcf1-f2 ）•通过一阶微分，对下游传播特征变量pc5u + Sp的等同和相反的校正在第一方程的Kacmarcz放松中发送到局部单元j和上游单元j - 1；对单元j的校正正好是从相应点Jacobi松弛点发现的一半 特征方程。 同样，一个平等的

70 THOMAS、DISKIN、BRANDT和SOUTH以及对上游传播特征变量pc5u - 5p的相反校正在Kacmarcz中将第二个方程的放松发送到局部单元j和下游单元j +1。 通过在每个点对解决方案进行局部更新，平滑率应等同于标准3点拉普拉斯算子的高斯-塞德尔松弛。 然而，与边界相邻的单元格的松弛与原始问题的边界条件不一致；例如，特征边界条件将分别在上游和下游边界将pcSu + Sp和pcdu-dp规定为零。 这个例子展示了前面提到的一般结果：分布式松弛需要在边界附近进行修改。 在这种情况下，可以找到幽灵变量5w = \{5wx,8w2)T的简单边界规范；在亚音速松弛的流入（流出）时，对5w\textbackslash\{\}（8w2）执行诺伊曼条件，对苏>2（5wi）执行狄利克雷条件，这与原始变量的特征边界条件一致。 相反，我们使用更一般的公式，并在边界点应用另一种局部保守松弛程序，以减少残差。 然后，分布式松弛在内部应用，使用幽灵变量Sw的简单狄利克雷条件。 松弛序列一般解程序是在流入边界、流出边界，然后在冲击区域将局部残差至少减少两个数量级--在用分布式松弛扫荡整个域之前和之后。 局部保守松弛在与边界相邻的前两个单元格上进行。 在冲击区域，局部保守松弛被应用于以冲击上游（超音速）侧为中心的九个细胞。 对应于100的Courant 数的伪时间步进被用作局部保守松弛中的放松不足系数；否则纯粹的稳态方程被放松。 在分布式松弛中，变量（5wi,8w2）T在实现中被明确使用；它们首先在内部细胞上放松，然后向（Su，5p）T进行分布，然后计算对Se的校正。 作为检查，单网格计算验证了具有一阶微分的亚音速方程（4.13）每次松弛的误差的收敛与应用在域一端的诺伊曼边界条件的一维3点拉普拉斯算子的高斯-塞德尔松弛相同。

\section{计算结果}
教科书多网格效率的一般框架71图1 马赫数分布在全亚音速、全超音速和跨声速流动的257点网格上计算，冲击为x = 0.75。 计算结果 本节显示了全亚音速、全超音速和带和无电击的跨声速流动的计算结果。 图中显示了257个点的网格上具有二阶精度离散化的马赫数分布。 1用于全亚音速、全超音速和带电击的跨声速流动；电击位置指定为x = 0.75。 在超音速情况下没有适用限制，因此在冲击时有一些振荡。 在这种一维的情况下，解决方案可以通过使用本质上非振荡（ENO）[21]方法（见附录三）轻松修复，但这项工作的重点是解决方案程序，而不是稳态结果。 完全亚声速流动 原始变量方程的分布式松弛平滑率预计将等于在L行列式的阶乘数的松弛中获得的平滑率中最差的平滑率，在我们的案例中，对流和全电位因子。 使用全电位运算符

72 THOMAS, DISKIN, BRANDT \& SOUTH h 1/32 1/64 1/128 1/256 l|ed||| -P 0.2654xl0-3 0.6341xl0"4 0.1547xl0\textasciitilde{}4 0.3817xl0-5 IMHMI IMI •P <0.01 <0.01 <0.01 <0.01 <0.01 表1 表1 收敛的p中的离散化误差和FMG-1算法后的相对Li-Norm完全亚声速流动.放松（高斯-塞德尔）和从流入边界下游行进行进的对流算子，一阶精确离散化的每个松弛扫描的多网格周期的收敛速率应接近一维3点高斯-塞德尔松弛的平滑系数0.447 拉普拉斯操作员。 缺陷校正通常用于解决高阶隐式离散化，以减少算术运算计数，同时保持目标精度。 缺陷校正的一个看似可能的实现是对L中的对流算子使用一阶精确离散化。 相应的分布矩阵M也由一阶算子组成。 然而，这种实现对于获得良好的平滑率不是一个好主意，因为就高频分量而言，低阶运算符和相应的高阶运算符不一定相互近似。 事实上，这种实现导致了收敛的放缓，而且未能达到与拉普拉斯方程相对应的平滑率。 为了说明这种现象，请仅考虑拉普拉斯算子。 将原始变量方程与此处使用的二阶上风偏置算子区分产生的椭圆算子5xx(u)j)是一个7点算子（见附录II）。 在高斯-塞德尔缺陷校正松弛中计算的当前近似解Wj的修正（Sw）j，以3点拉普拉斯为左侧（驱动）算子，可以从\textasciicircum{}[(ftiVi - 2(Sw)j] = -8xx(wj)中找到。 （4.28）使用局部模式分析计算的这种缺陷校正松弛（4.28）的平滑系数为gs = 0.525，这大大小于收敛速率，通过应用于Eq的多网格FV（2,1）循环观察到0.9的分布松弛。 （4.22）。 方程的分布式松弛的平滑系数。 （4.22）是通过使用局部模式分析作为矩阵G的规范来计算的，定义为

教科书多网格效率的一般框架 73 h 1/32 1/64 1/128 1/256 IMI -P 0.8890xl0"4 0.2151xl0\textasciitilde{}4 0.5278xl0\textasciitilde{}5 0.1306xl0-5 IMHMI IMI •P <0.01 0.02 <0.01 <0.01 表2 收敛时p中的离散化误差和完全超声速流动的FMG-1算法后的相对Li-norm误差。 G = /-(l-fli)(McJ)-1Mt) (4.29) 其中与目标（二阶）和驱动（一阶）运算符的分布矩阵关联的矩阵符号分别表示为M4和M\textasciicircum{}。 计算出的平滑系数为0.752，明显高于0.525的预期值。 平滑系数的增加是由（Md）\_1Mt与恒等矩阵大不相同这一事实来解释的。 通过二阶分布矩阵运算符（M<j = Mt）分布的校正，方程的分布松弛平滑系数。 （4.22）应该与拉普拉斯方程的松弛相同。 实现拉普拉斯方程的缺陷校正松弛（4.28），但用二阶分布分配校正，每次平滑的FV（2,1）周期收敛速率为0.51-接近缺陷校正松弛（4.28）的模式分析预测的0.525。 表1显示了完全亚声速流动的FMG-1算法的结果，表示每个级别有一个FV（2,1）周期的完整多网格算法。 本表和下表中使用的所有规范都是Li规范。 压力的离散化误差按点定义为e\textasciicircum{} = pj - p(xj)exa,ct，其中p(x)exact是Eq的稳态解。 （4.4）和pj是Eq的离散解。 （4.13）。 e<i的Li规范中的二阶空间收敛是显而易见的。 压力的总误差定义为et = p(xj)exaxit - pj，其中pj是FMG-1算法中得到的近似解。 总的Li规范和离散化误差之间的相对差异很小（见表1）表明已经达到了最佳效率。 对于四个网格中的任何一个，FV(2,1)周期的渐近收敛速率为每次放松0.50-0.52。

\section{跨声速流动}
74 THOMAS, DISKIN, BRANDT \& SOUTH Fully 超声速流动 在完全超声速流动中，通常可以构建行进方法来解决全势流方程。 对于一阶微分，单边算子可以通过域轻松求解。 对于这里考虑的二阶算子，一阶上风驱动的缺陷校正放宽可以定义为\textasciicircum{}[(<Hj-2 - 2(<Hj-i + (<Hj] = Sxx(wj)。 （4.30）这个方案可以通过行进来解决，但不幸的是，对于一些错误组件，其放大系数大于1。 一个合适的全电位驱动运算符来取代Eq的左侧运算符。 （4.30）可以在[6- + haS-d-]2的形式中找到，（4.31）其中后向差分运算符在上述表达式中被定义为一阶运算符。 这是一个5点运算符，可以通过行进来解决。 超音速缺陷校正松弛与驱动器（4.31）和a - 0.23的最大放大系数为0.55。 使用二阶分布算子，完全超声速流动的计算对应于流入时的马赫2，以这种速度渐近收敛。 FMG-1算法的结果如表2所示。 离散化压力误差的ii-规范显示了预期的二阶空间收敛。 使用FMG-1算法实现了最佳效率。 在这里，每个级别的单网格迭代的性能与FV（2,1）周期相同，因为全电位因子是通过行进解决的（而不仅仅是平滑）。 对于四个网格中的任何一个，每个松弛的渐近收敛速率为0.53-0.56。 表3显示了平滑跨声速流动的FMG-1算法的跨声速流结果。 离散化压力误差的Li-norm显示了预期的二阶空间收敛。 总误差的L\textbackslash\{\}-规范与离散化误差的Xi规范的相对偏差从全亚音速或全超音速水平增加，但仍然表现出最佳性能。 对于四个网格中的任何一个，每次放松的渐近收敛速率都是0.52-0.56。 表4显示了跨声速流动电击的FMG-1算法的结果。 离散化压力误差的Li-norm显示一个

\section{结束语}
教科书多网格效率的一般框架 75 h 1/32 1/64 1/128 1/256 Ikdll -P 0.1075xl0\textasciitilde{}3 0.2666xlCT4 0.6635xl0-5 0.165xl0-\& IMHMI IMI •P 0.05 0.09 0.03 0.11 表3 收敛中的p中的离散化误差和跨声速流动的FMG-1算法后的相对Li-norm误差。h 1/32 1/64 1/128 1/256 IMI:P 0.4396x10-\textasciicircum{} 0.2218xl0"2 O.llllxlO"2 0.5558xl0"3 Wet \textbackslash\{\}-\textbackslash\{\}\textbackslash\{\}ed\textbackslash\{\}\textbackslash\{\}ed\textbackslash\{\}\textbackslash\{\} Ml •P <0.01 0.015 0.014 <0.01 表4 收敛时的离散化误差和跨声速流动的FMG-1算法后的相对Li-Norm误差。一阶空间收敛， 因为冲击时插值没有限制；即，冲击时发生最大局部误差，并且不会随着网格加密而减少。 FMG-1算法的总误差的Li规范与离散化误差的Li规范的相对偏差再次表现出最佳性能。 在10个周期内，每次放松的平均收敛速率从最粗的网眼的0.35到最精细的网眼的0.51不等。 附录三介绍了本情况下具有ENO差异的计算，并显示出与冲击区域单调溶液行为的类似性能。 结论 讨论了一个通用的多网格框架，用于获得守恒律形式的可压缩欧拉和Navier-Stokes 方程解的教科书效率。 一般方法依赖于分布式松弛程序来减少规则（平滑变化）流动区域的误差；对每个因素（椭圆和/或双曲）进行单独和不同的处理，以实现最佳的误差减少。 在边界和近冲击附近，有必要对保守方程进行额外的局部松弛。 计算示例是针对准一维Euler 方程进行的

\section{参考资料}
76 THOMAS、DISKIN、BRANDT和SOUTH适用于全亚音速和全超声速流动的情况，以及有和没有冲击的跨声速流。 所有计算都表明，FMG-1算法提供了非常准确的近似值。 参考资料1。 Brandt，A.，多电网开发指南。 多网格方法，数学讲义960，Springer-Verlag，1982年。p. 220-312。 2. Brandt，A.，多网格技术：1984年流体力学应用指南，计算流体动力学讲义，Von-Karman流体力学研究所系列讲座，魏茨曼科学研究所，以色列雷霍沃特，1984年，191页，ISBN-3-88457-081-1，GMD Studien Nr. 85. 可从GMD-AIW，Postfach 1316，D-53731，St. 德国奥古斯丁1号。 3. Brandt，A.，在CFD中实现教科书多电网效率（TME）的障碍。 ICASE临时报告编号 32，美国宇航局CR-1998-207647，1998年；更新版本可作为高斯中心报告WI/GC 10，魏茨曼科学研究所，以色列雷霍沃特，12月。 1998年。pp. 1-23。 4. 勃兰特，A。 \& Diskin B.，非对齐声波流的多网格求解器：恒定系数情况，计算机与流体28，1999年，pp. 511-549；还有高斯中心报告WI/GC-8，以色列魏茨曼科学研究所，10月。 1997年。 5. 勃兰特，A。 \& Diskin，B.，非对齐声波流的多网格求解器，SIAM J. 科学。 补偿。 2000年21日，pp. 473-501。 6. 勃兰特，A。 \& Yavneh，I.，关于高雷诺德不可压缩输入流的多网格解决方案。 J。 补偿。 物理。 101，1992年，pp. 151-164。 7. 勃兰特，A。 \& Yavneh，I.，加速多重网格收敛和高雷诺德再循环流，SIAM J. 科学。 补偿。 14，1993年，pp. 607-626。 8. Diskin，B.，线性化跨声速全势流方程的高效多网格求解器，博士论文，魏茨曼科学研究所，1998年。 9. Yavneh, I., Venner, C, \& Brandt, A., 具有闭合特性的对流问题的快速多网解决方案, SIAM J. 科学。 补偿。 1998年，pp. 111-125。 10. 迪斯汀，B。 \& 托马斯，J. L.，解决上风偏置离散化：缺陷校正迭代，ICASE报告99-14，NASA CR-1999-209106，1999年。 11. 迪斯汀，B。 \& 托马斯，J. L.，对流的Fromm 离散化缺陷校正方法的半空间分析，SIAM J. 科学的。 Comp.将于2000年出现。 12. 托马斯，J。 L., Diskin, B., \& Brandt, A., 应用于可压缩Navier-Stokes 方程的分布式松弛和缺陷校正，Proc. AIAA第14届CFD会议，1999年7月，弗吉尼亚州诺福克。 AIAA论文99-3334。 13. Ta'asan，S.，稳态Euler 方程的规范变量多重网格法，ICASE报告94-14，1994年。 14. 罗伯茨，T。 W。 \& 斯旺森，R. C.，将理想收敛的多网格方法扩展到翼型流量，Proc. AIAA第14届CFD会议，1999年7月，弗吉尼亚州诺福克。 AIAA论文99-3337。 15. Sidilkover，D.，为无黏流动方程构建最佳高效多网格求解器的一些方法。 计算机与流体，1999年28日，pp. 551-571。 16. 罗伯茨，T。 W.，Sidilkover，D.，和Thomas，J. L.，可压缩Euler 方程的可分解保守离散化的多网格松弛。 AIAA论文2000-2252，2000年6月。 17. 托马斯·J。 L., Diskin, B, \& Brandt, A., 教科书多电网效率

\section{附录}
\subsection{I - 保守通量}
教科书多网格效率的一般框架77用于不可压缩的Navier-Stokes 方程：高雷诺数唤醒和边界层。 ICASE报告99-51，美国宇航局CR-1999-209831，1999年；还有计算机和流体，将出现在2000年。 18. 蒙特罗，R。 S。 \& Llorente，我。 M.，不可压缩Navier-Stokes 方程的强大多网格算法。 ICASE报告2000-27，美国宇航局CR-1999-210126，2000年。 19. Barth，T。 J.，上风和TVD算法的隐式局部线性化技术分析。 AIAA文件87-0595，1987年。 20. 罗，P。 L.，近似黎曼求解器、参数向量和差分方案。 J。 补偿。 物理。 43（2），1981年，pp. 357-72。 21. Shu C-W.，双曲守恒定律的基本非振荡和加权基本非振荡方案。 ICASE报告97-65，美国宇航局CR-97-206253，1997年11月。 附录I - 保守通量 通量是使用Roe的通量差分分裂构建的[20]，并给出为Fj+i = i[F(Qfl) + F(QL) - |A|(Qfl - QL)] (4.32)，其中A = dF/dQ用左和右状态、Qjr和QR的特定平均值进行评估，以准确满足冲击的跳跃条件[20]，左右状态为QL = Qj，（4.33）QR = Qi+i，（4.34）和Fromm的离散化的二阶精度为QL = Q\textasciicircum{} + +i-Q,--!)，(4.35) QR = Qj+1-7（Qj-Qj+2）。 （4.36）耗散矩阵|A| = T |A| T：其中T，T x是对角化矩阵，特征值矩阵为A = Xl 0 0 0 A2 0 0 0 A3（4.37）

\subsection{II - 分布矩阵}
78 THOMAS, DISKIN, BRANDT \& SOUTH and (Ai,A2,A3)T = \{u + c,u-c,u)T。 （4.38）波浪号上标表示特征值限制在零之外，|Ai \_ / |Ai| - I (A? 如果|A;| >€, (4.39) i2 + (?) 2）/（2？） 否则，主要是为了防止一阶离散化精度的声波点的膨胀冲击，但也是为了使全电位操作员的声波点平稳过渡。 值Z取为/3max\{|Ai|，j Aa|，I A31J-。 这种“熵修复”在任何意义上都不是最佳的，因为更大/？ 较粗网格上的一阶方案需要声波点的值。 J3的值名义上设置为0.1，对于隐性一阶雅各布矩阵中9点或更少的网格，j3的值被任意增加到0.2。 II - 分布矩阵系数矩阵A\textasciicircum{}被定义为对A的正负特征值贡献，其中Xf = (\textbackslash\{\}i±\textbackslash\{\}Xi\textbackslash\{\})/2和A以其特征值A(Ai, A2, A3)的形式写成|(Ai + A2) f(Ai-A2) £(Ai + A2) 0 23E(Ai-A2) £(Ai-A2) \textasciicircum{}(A1 + A2-2A3) A3 L = A+5\textasciitilde{}+ A\textasciitilde{}\#+的上2x2块可以写成pc(t2) h，相应的协因子矩阵是M = h -pc(t2) -£(*») (4.40) (4.41) (4.42) (4.43) 其中h = i[(A+-A+)C + (Ar-A2-)5+]。 （4.44）（4.45）

\subsection{III - 跨声速冲击 - ENO分化}
教科书多网效率的一般框架79 L的行列式是t\textbackslash\{\}-l\textbackslash\{\} = (XtXt)6\textasciitilde{}5- + (ArA2-)<5+<5+ (4.46) 二阶上风偏置微分的算子定义为K CK) = w(-«»j-3 + 2\textasciicircum{}-2 + VIWj-x -2,\&Wj + I7wj+1 + 2wj+2 - wj+3), (4.47) -28u)j\_i - Wj + 6WJ+I + Wj+2）。 (4.48) III - 超音速休克 - ENO Differencing h 1/32 1/64 1/128 1/256 lled||:P 0.3597xl0-3 0.9711xl0-4 0.5133xl0"4 0.2665xl0-5 l|e«l -lledll eM •P <0.01 0.01 0.07 0.13 表5 离散化在收敛处有ENO差异的p误差和跨声速流动的FMG-1算法后的相对Li-norm误差，以及跨声速流动的FMG-1算法后的相对Li-norm误差。 对附录I（Fromm方案）的状态变量的上风偏置插值的均匀应用导致冲击时的振荡。 这些振荡可以通过限制程序消除，以局部减少冲击区域的近似顺序。 在这里，使用ENO方法[21]来防止状态可变插值穿过冲击。 在冲击界面，左（右）状态变量QL（Q\_R）通过单侧二阶外推找到；在冲击的第一个界面上游（下游），状态变量QL（QR）通过二阶中央平均找到。 图。 2展示了冲击区域的马赫数分布，并显示了具有单点表示的非鼻振荡行为

80 THOMAS, DISKIN, BRANDT \& SOUTH 1.3 r o is 0.9 0.8 0.7 0.6 精确离散解 -O FMG-1 算法 0.65 0.7 I l 0.75 0.8 图 2 马赫数 与精确离散解（h - 1/256）相比，用FMG-1算法计算的冲击附近的马赫数分布。冲击。 表5显示了FMG-1算法的总误差的Lj规范的离散化误差和相对偏差；与表4的比较表明，ENO 离散化误差小于无限插值的离散化误差，尽管仍然发现一阶行为。 在这两种情况下，对于这种一维情况，如果Li-norm仅限于冲击的远上游或远下游的区域，则可以达到二阶精度。 虽然没有显示，但我们注意到，在这种情况下，如果在冲击处丢弃熵固定，可以通过上述ENO程序找到统一的二阶精度；然后以相同的方式恢复冲击跳跃，并恢复零点冲击。 这项调查的重点是收敛和图。 2和表5表明FMG-1算法已经达到了最佳效率。

\clearpage
\chapter{5 二维和三维流动的Cauchy-Riemann 方程的数值解}
Cauchy-Riemann 方程的二维和三维流动的数值解Mohamed M. Hafez1和J。 Houseman1摘要对于二维流动，质量守恒和涡量的定义构成了假设涡量给出的速度分量的广义柯西-黎曼系统。 如果流动是可压缩的，密度是速度和熵的函数，后者被认为是已知的。 引入人工时间，可以很容易地构建一个对称的双曲系统。 数值稳定性需要人工黏度，从最小二乘法公式中获得。 增强系统用高度可并行化的标准点松弛算法明确解决。 对于三维流的扩展，连续性方程与两个涡量分量的定义相结合，并为三个速度分量求解。 二阶准确的结果与圆柱体和球体周围的不可压缩、非旋转和旋转流动的精确解决方案进行比较。 还包括可压缩（亚音速）流动的结果。 1 加州大学戴维斯分校机械和航空工程系，加利福尼亚州戴维斯，95616 计算流体动力学前沿 - 2002 编辑：David A. Caughey和Mohamed M. 哈菲兹 ©2002 世界科学

\section{5.1 介绍}
82 HAFEZ \& HOUSEMAN 5.1 介绍 数值流模拟通常基于欧拉和Navier-Stokes 方程的原始或保守变量的解。 由于相关的困难和限制，基于向量电位或速度涡流方程的替代公式并不流行。 最近，有一些努力在运动学和运动动力学解耦合的基础上构建数值方案。 参考文献中讨论了此类配方的动机和优势[8]。 在这项工作中，我们对使用连续性方程和涡量的定义来计算速度场感兴趣。 对于稳定、不可压缩、无黏性、非旋转的流动，可以从标准柯西-黎曼系统中获得两个速度分量。 就势函数或流函数而言，两个一阶方程等同于一个二阶拉普拉斯方程。 此外，最小二乘法程序为两个速度分量产生两个拉普拉斯方程。 如果流动不是非旋转的，并且场中也有源，那么非均匀的Cauchy-Riemann 方程不能还原为单个二阶方程。 速度向量可以表示为电位的梯度加上另一个向量的卷曲。 第一个成分是无卷曲的，第二个成分是无发散的。 根据亥姆霍兹定理，这种分解总是在非常普遍的条件下可能的。 然而，对于一般的三维流，边界条件很复杂。 另一方面，最小二乘法程序为三个速度分量生成三个二阶方程。 然而，在这种公式中，如何轻易地强加熵条件，以排除跨声速流模拟的膨胀冲击，还不清楚。 最小二乘法公式也存在其他问题，即使是不可压缩的流动。 首先，必须通过在边界上强加一阶方程来排除伪解。 此外，质量可能无法在离散水平上保持。 特殊安排，例如，交错网格可能是必要的，以确保保持数值通量。 对于从动量方程的卷曲中获得的涡量，需要进行特殊处理，以保证其发散在离散水平上完全消失。 在下面，我们将解决嵌入由对称双曲系统控制的人工时间过程中的广义Cauchy-Riemann 方程，导致一个良好的初始边界值问题。 由于我们只对稳态解决方案感兴趣，因此瞬态行为的准确性不是问题。 可以使用多网格等标准收敛加速技术来提高计算效率（见附录A）。

\section{5.2 控制方程和边界条件}
Cauchy-Riemann 方程 83的解决方案 为了避免奇偶解耦合并确保数值稳定性，必须引入人工黏度。 为了保证二阶精度（对于不可压缩和亚声速流动），使用了基于最小二乘公式的人工黏度（见Hughes等人[3]和Tezduyar和Hughes[9]）。 该方案也与Lerat最近的工作[5]相似，至少在亚音速流动上是这样。 二阶项在连续级别上完全消失，因为它们与Cauchy-Riemann 方程一致。 在离散水平上，它们产生更高阶的耗散。 这种结构可以被视为一种紧凑的方法，以取代McCormack[6]和Jameson[4]提出的常用的四阶耗散方案。 本文分为四个部分；治理关系、数值算法、数值结果和结论以及一些一般性评论。 5.2 控制方程和边界条件 给出了以圆柱形和球形坐标书写的圆柱体和球体上的流动方程。 在圆柱体连续性方程上的流动ld(rpvr) l\_d(pve) p [r dr r 39 = 0 (5.1) 涡量定义1 d(rve) 1 d(vr) r dr r d0 围绕球体的三维流动 (5.2) 连续性方程 1 d(r2 pvr) 1 d(sin 9 pvg) 1 dpv<i, dr r sin i 86 r sin 9 d<p 0 (5.3) 涡量 定义 1 dvr 1 d(rv,p) r sin 9 d<j> r dr 1 8(rvg) 1 dvr r dr r 39 (5.4) (5.5)

\section{5.3 数值方法}
84 HAFEZ \& HOUSEMAN和边界条件是：vr - 0在r - Ti（5.6）vg和i>0在r = rg，（5-7）给出，其中r和r0是内球和外球的半径。 在非旋转、等熵流动的情况下，涡量同样消失，根据伯努利定律，密度与速度有关：•y-1（5.8），其中M\textasciicircum{}是自由流马赫数。 对于不可压缩的流动M\textasciicircum{} - 0。 5.3 数值方法 圆柱体和球体周围的不可压缩和可压缩（亚音速）流量使用最小二乘公式计算，以及嵌入人工黏度增强的对称双曲系统中的Cauchy-Riemann 方程。 所有空间导数都使用二阶精确有限体积方案离散化。 时间导数项通过一阶时差离散化。 离散化系统使用点高斯-塞德尔松弛方案求解。 源代码以C++实现，并在HP-Visualize C3000工作站上运行。 下述了我们解决的问题。 （带冲击的跨声速流在[7]中使用类似方案计算，但使用线性化边界条件）控制方程，表示为L（w），（5.9）可以嵌入以下系统dtw = L（w）+ eL*L（w）。 （5.10）其中L*是运算元L的伴随。 对于笛卡尔坐标和不可压缩的非旋转流，L*L(w)是w分量的拉普拉斯。 我们强加的边界条件是：在固体表面上：径向速度vr使用Dirchlet条件，从涡量方程中得出的Neumann条件用于角速度vg（和v\textasciicircum{}）。 远场：角速度vg（和v<j）使用Dirchlet条件，从连续性方程得出的Neumann条件用于径向速度vr。-MliW

\section{5.4 数值结果}
Cauchy-Riemann 方程 85流场的解决方案：周期性边界条件用于vr和v\$（和v\textasciicircum{}）的流场。 5.4 数值结果 5.4.1 不可压缩流动 我们用分析解决方案测试了当前的公式，例如，在圆柱体上循环流动，均匀和剪切，以及旋转流体中的圆柱体（见Batchelor [1]）。 我们还使用轴对称和全三维方程计算了球体上的流量。 在图中。 1，流线是为圆柱体和球体上的不可压缩流动绘制的。 在上述所有情况下，结果在准确性和收敛方面都是令人满意的。 例如，球面上的流量计算结果如下。 示例1：在本例中，我们假设围绕球体的轴对称非旋转流，但我们求解每个速度分量（vr、vg、v\textasciicircum{}），就像在全三维流中一样。 数值解是在具有三个空间维度的球面坐标系中计算的。 使用的网格间距是Ar = 0.1、A9 - 7r/18和A<f> = 7r/18。 网格尺寸为36 x 19 x 37，球体的半径非维化为1。 程序运行，直到达到10-5的剩余公差，结果如表1所示。 5.4.2 可压缩流动 示例2：我们用e = 1和e = \textasciicircum{}求解两个空间维度的圆柱坐标系，在最小二乘公式中有和没有密度的圆柱体上求解可压缩流动。 网格间距为Ar = 0.1，Ad - 7r/18。 网格尺寸为36 x 37，圆柱体的半径非维度化为1，圆柱体周围的循环保持在T = 0，涡量 LJZ = 0.0，使用0.1和0.2的马赫数。 程序运行，直到达到10\textasciitilde{}5的剩余公差。 （见表2和图。 2a.） 示例3：我们使用e - I和e = \textasciicircum{}在两个空间维度中的球面坐标系求解可压缩流动。 网格间距为Ar = 0.1和A9 = 7r/18。 网格尺寸为36 x 37，球体的半径非维度化为1，涡量 wz = 0.0，使用0.1和0.2的马赫数。 该程序运行，直到达到10"5的剩余公差。 （见表3和图。 2b.）

\section{5.5 结束的讲话}
\section{5.6 附录：多重网格收敛 结果}
\section{参考资料}
86 HAFEZ \& HOUSEMAN 上述所有结果都是使用恒定的人工黏度获得的，e。 然而，一般来说，e可以因点而异，也可以因迭代而异，并且可以根据局部流量条件进行优化，以加速计算的收敛。 5.5 结论 鉴于涡量和密度，二维和三维流动的速度分量可以从广义的柯西-黎曼方程组中计算出来。 标准数值算法适用于实现预期的精度和效率。 对于完整的流动模拟，必须包括运动的动力学，以提供熵和总焓，从而提供涡量。 完整的模拟仍在进行中，将单独报告。 5.6 附录：多重网格收敛结果点松弛方案的收敛可以通过实施多网格V-Cycle方案来增强。 Briggs、Henson和McCormick [2]的V-Cycle方案被修改，以计算使用e = §的Cauchy-Riemann 方程在循环管上的不可压缩、非旋转的流动。 用于计算的精细网格是N = 322,642,1282,2562。 每次运行使用的最粗网格是N = 82。 点松弛和V-Cycle(fl,\textasciicircum{}2)的比较如表4所示。 参数v\textbackslash\{\} = 1是沿着网格使用的松弛数，v2 = 2是沿着网格使用的松弛数量。 比较包括HrHoo < 10\textasciitilde{}5的收敛之前的迭代次数，计算解决方案和分析解决方案之间的绝对误差||e|loo，以及CPU运行时。 参考资料1。 巴切勒，G.K. 流体力学简介。 剑桥出版社，1967年。 2. 布里格斯，W.L.，亨森，V.E. 和麦考密克，S.F. 多网格教程。 SIAM，2000年。 3. 休斯，T.J.R.，弗朗卡，L.P. \& Hulbert，G.M. 计算流体动力学的新有限元公式：Viii对流扩散方程的加勒金/最小二乘法。 应用力学和工程中的计算机方法，73：173-189，1989年。 4. 詹姆森，A.，施密特，W. \& Turkel，E. 使用Runge-Kutta 时间推进方案的有限体积方法对Euler 方程进行数值解。 AIAA，（-

Cauchy-Riemann 方程 87 * \textasciicircum{} ' '' ' (a) ' \textasciicircum{} ' ' (b) ' 图1 不可压缩流动经过（a）圆柱体和（b）球体的流线，W0 = 0。 示例2。 收敛历史或残点M = 0.1示例3：收敛历史或残点M = 0.1迭代次数迭代次数（a）（b）图2（a）示例2和（b）示例3的系统1和2的||r|oo规范的收敛历史。 1259），1981年6月。 5. 莱拉特，A。 \& Corre，C。 基于残数的紧致方案，用于守恒定律的多维双曲系统。 计算机和流体，出现。 6. 麦科马克，R.W. \& Paullay，A.J. 计算网格对不连续或非线性解的初始值问题的准确性的影响。 计算机和流体，2:339-361，1974年。 7. 罗伊，J.，哈菲兹，M. \& Chattot，J. 求解广义科西·里曼方程和模拟隐式旋转流的显式方法。 计算机和流体，出现。 8. 唐，C。 和M。 哈菲兹，M。 使用区域公式对稳定可压缩流动进行数值模拟。 计算机和流体，出现。 9. Tezduyar，T.E. \& 休斯，T.J.R. 对流主导流的有限元公式化，特别强调可压缩的Euler 方程。 AIAA，（83-0125），1983年1月。

HAFEZ \& HOUSEMAN 表1显示的是计算解决方案的迭代总数和\{vr,ve,V4,)的最终最大规范误差。 其中误差是精确分析解和数值结果之间的差额。 系统1 系统2 系统3 e Ar 迭代次数 绝对误差 0.0084 0.0147 0.0309 表2 所示是e = 1和e = \textasciicircum{}和自由流马赫数0.1和0.2的计算解决方案达到10\textasciitilde{}5的剩余公差的迭代总数。 系统1系统2 e Ar迭代（M = 0.1）迭代（M = 0.2）表3显示的是计算解决方案的剩余公差达到10 5的迭代总数，e - 1和6 =数字0.1 amd 0.2。\textasciitilde{}和自由流马赫系统1系统2 e Ar迭代（M = 0.1）迭代（M = 0.2）N 32\textasciicircum{} 64\textasciicircum{} V-循环迭代表4 收敛结果IMIoo 1.40e-2 3.52e-3 8.80e-4 2.19e-4时间（秒）点松弛迭代11\textasciicircum{}1 |oo 1.40e-2 3.52e-3 8.89e-4 2.27e-4时间（秒）

\clearpage
\chapter{用于多维可压缩流的非均匀网格上的6种高效高阶方案}
\section{6.1 导言}
用于多维可压缩流的非均匀网格上的高效高阶方案A。 勒拉特，C。 科雷和G。 汉斯：6.1 导言 在参考[l]-[2]中为可压缩欧拉和Navier-Stokes 方程提出了一种三阶精度的耗散方案。 这个方案很紧凑，但它不需要线性代数系统的解（对于其显式形式），它涉及一个简单高效的一阶数值耗散。 这些属性是由于该计划的残余建设造成的。 在本文中，我们的目标是研究该方法在非均匀网格上的实现，该方法严格保持三阶精度，而对网格没有任何平滑性假设。 这项研究涉及2D和3D不规则笛卡尔网格，其步骤根据一些拉伸或随机方式而变化。 首先，第6.2节简要介绍了正笛卡尔网格上的欧拉求解器，从两个设计原则中称为基于残差的紧凑性和基于残差的耗散。 然后，欧拉求解器被扩展到第6.3节中完全不规则的笛卡尔网格，使用有限体积公式、精度的适当定义和一些基于残差的校正。 第6.4节考虑了对Navier-Stokes 方程的概括。 最后，数值实验在第6.5节中介绍了三个系列1 SINUMEF实验室，ENSAM，151 Bd de l'Hopital，75013 Paris，Prance。 Frontiers of 计算流体动力学 - 2002 编辑：David A. Caughey和Mohamed M. 哈菲兹 ©2002 世界科学

\section{6.2 常规笛卡尔网格上的欧拉求解器}
90 LERAT、CORRE和HANSS的网格和五个测试案例：二维旋转对流、三维螺旋对流、二维波伊苏伊勒模型问题、二维布罗杰斯问题和二维可压缩射流相互作用。 这些测试用例允许在解决方案平滑时比较形式和计算精度顺序，并展示该方案如何在没有限制器的情况下捕获不连续性。 6.2 正笛卡尔网格上的欧拉求解器 6.2.1 基于残差的紧致度 我们首先考虑二维可压缩Euler 方程：v>t + fx+9y = 0（6.1），其中w是保守变量的向量，/ - f（w），g = g\{w）是x和y方向的通量分量。 等式中的空间导数。 （6.1）在规则笛卡尔网格上近似：Xj - jdx, yk = kSy，与5x PS 6y - 0(h)，使用基本的离散运算符：\{Siv)j+i<k = Vj+1<k - Vjtk (<Mj,fc+± = v3,k+l \textasciitilde{} Vj,k (Ml\textasciicircum{})j+l,fe = UV3 + + Vj,k) (/i2«)j,fc+l = \textbackslash\{\}\{Vj,k+l + vj,k)-空间导数的经典紧近似在3点模版上是四阶准确的。 这个紧凑的模版要付出的代价是每个空间方向的线性代数系统的数值解。 在使用基于残差的紧致度计算稳态解时，可以避免这种额外成本（通过引入对偶时间，对于非稳态问题也是可能的）。 这个想法不是在高阶下离散化每个空间导数，而是考虑稳态r = fx + 9y的整个残差。更准确地说，fx和gy由二阶精度的特殊中心差公式近似，因此全局截断误差只能用r的导数表示，即：1 1 5x2 Sy2 - Sim(I + -5\textbackslash\{\})f = fx + -jrfxxx + -jrfxyyy + 0(fl4) 1 1 Sy2 5x2 -S2H2(I + Q\&\textbackslash\{\})9 = 9y + -£-9yy + -\textasciicircum{}9xxy + 0(hA)。其中/是恒等式运算符。 在Eq中插入这些近似值。 （6.1），我们得到：wt + \textasciicircum{}i»i(I + l\%)f + \textasciicircum{}<W/ + \textbackslash\{\}*i)9 = 0。 （6-2）

非均匀网格上的高阶方案91此方案有一个3 x 3点模版，不需要线性代数。 为了Eq的平滑解。 （6.1），其截断错误为：ST SV e = wt+r+ -rxx + -\textasciicircum{}-ryy + 0(h4)，因此e = 0(h4)对于精确的稳态解 (r = 0)。 因此，紧凑方案（6.2）在稳态下是四阶精度。 基于残基的紧凑性可以很容易地扩展到三个空间维度，也可以扩展到更高的精度顺序（见[1]，[2]）。 6.2.2 基于残差的耗散数值耗散也是基于残差的，即由r的衍生物构建。 因此，它可以是瞬态阶段的一阶，因此可以以简单的方式提供稳健性，同时在稳态下保持高精度。 更准确地说，方案的耗散形式是：m + \textasciicircum{}SiHi(I+Ul)f + \textasciicircum{}-52\textasciicircum{}\{I+\textbackslash\{\}sl)g = i<Ji(*ir!) + \textasciicircum{}2(\textasciicircum{}2) (6.3) ox 0 oy 6 I 2，其中ri和r2分别是j + \textbackslash\{\}, k和j,k+ \textbackslash\{\}的r的二阶中心近似。 系数为1美元和<3？ 2是矩阵，仅取决于8x和Sy以及通量Jacobian矩阵A-df/dw和B-dg/dw。 它们已被确定为确保最佳的多维耗散（见[1]）。 对于标量二维双曲方程，它们可以写成：1。 5x\textbackslash\{\}B\textbackslash\{\} \$1 = sgniA) vnmil, -), \$2 = sgn(B) mm (1, a) 其中a =. a oy\textbackslash\{\}A\textbackslash\{\} 对于守恒定律的二维双曲系统，\$1和\$2在[1]中给出了相当简单的表达式。 耗散方案（6.3）的截断误差是：Sx\textasciicircum{} V -\textasciicircum{}(\$ir) -\textasciicircum{}' r。 ' XX T" - I yy - \textbackslash\{\}\textasciicircum{}\textbackslash\{\}l)x „ e = wt + r + -rxx + -ryy - -\{>f>ir)x - - -\textasciicircum{}(<f>2r)y + 0(h3)。 对于精确的稳态解，我们得到e = 0(h3)，即三阶精度。 此外，方案（6.3）是紧凑的，因为它的数值耗散与其非耗散部分在同一3x3点模版上离散化。 6.2.3 紧凑的三阶方案 这里只对稳态解感兴趣，（6.3）中的时间导数只是近似于一阶。 完全离散方案的显式版本如下：

92 LERAT、CORRE和HANSS a）非耗散数值通量：Fli* = w+1f)i*niH*（6.4）b）数值耗散：5x（J \textbackslash\{\}n \_ 6x r*。 A/, S2HlH29,]n (di)j+i,k \textasciitilde{} T [*i (\textasciicircum{} + -\textasciicircum{}-)Ui,* (d2)jik+i - -f [*2 (-\textasciicircum{}- + •\textasciicircum{}-)],-,fc+I c) 总数值通量：(6.6) d) 新单元值：A-]J' = -Ai(\textasciicircum{} + \textasciicircum{} \textasciicircum{} (6.7)?;in+1 - •)/>" -I- A? /ieXp' j,k \textasciitilde{} Wj,k + \textasciicircum{}Wj,k ' 其中w\textasciicircum{}k表示时间nAt和位置（jSx，kSy）的数值解。 请注意，上述方案不包含限制器。 图1显示了方案（6.4）-（6.7）的线性不稳定域，用于二维标量问题，以每个空间方向的CFL数计算。 为了提高解决稳定问题的计算效率，已经开发了一个隐式版本的方案（见[1]，[2]）。 该方法的隐式阶段是：A«iM +At[\textasciicircum{}(A\textasciicircum{}) " \textasciicircum{}i(*i A\textasciicircum{})\textbackslash\{\}lk +At[v2(Bs-if) - ±82(*2B\textasciicircum{})]lk = Awe3Jl, W W1M1 = wlk + Awi* 已经表明，隐式方案(6.4)-(6.8)总是线性稳定，当与大CFL数一起使用时提供良好的阻尼特性。 在实践中，代数线性系统（6.8）是通过Jacobi或Gauss-Seidel类型的交替线松弛求解的。 只需要几次放松过程的内部迭代（通常一到两次）才能确保无条件的稳定性和非常接近（6.8）的阻尼。

\section{6.3 不规则笛卡尔网格上的欧拉求解器}
非均匀网格上的高阶方案93 -1 -0.5 CFL，=Ait/5x图1显式三阶方案的二维稳定性（隐性问题）备注：在常规网格上，可以给出非耗散性数值通量（6.4）的各种等效表达式，所有这些都导致相同的方案；例如：（6.9）组合再次给出（6.2）。 （6.4）和（6.9）的任何线性组合也是等价的。 以下组合将出现在下一节中：1 Fj+htk = \{(I + -SD\textasciicircum{}f + -SxS, <*2M2fls „ 24 6y Jj+\textasciicircum{}'k GiJc+i = [(I + n*i)M + tfyh-\textasciicircum{})"\textasciicircum{} (6.10) 6.3 不规则笛卡尔网格上的欧拉求解器 6.3.1 精度顺序的有限体积定义 我们现在考虑方程的积分形式。 (6.1): - / w dxdy + f fdy - gdx - 0, dtJn Jon (6.11) 其中J7是一个边界为9f2的有界域。 让£ljtk是结构网格和Tj+i k的四边形单元格，Tjk+i是作为递增网格索引的单元格边缘（见图2）。 单元格区域用|fij,fc|表示，边缘长度用|rj+ifc|表示，网格中最大的边缘长度用h表示。 应用于单元格£lj,k，确切的形式（6.11）为：

94 LERAT, CORRE AND HANSS 图 2 单元格 ft j,k- |fii,fel(\textasciicircum{})i,* +|r,+i,fc|5.+iifc-|ri. 其中条形表示精确平均值：\textbackslash\{\}Fj,k-±\textbackslash\{\}Gj,k- = 0 (6.12) (\textasciicircum{}*)j,fc = ITS-r 3i / w dxdv F j+hk - f /dy - c;dx, Gj)fc+i = p= L gdx - fdy。 对于（6.12）的空间近似，我们考虑了有限体积（FV）方案，形式为：l\%*|(\textasciicircum{});,fc +|ri+iife|Fi+iifc 1 J 2 'Kl J fl* "i.fc-il\textasciicircum{}'.fc-j 0 (6.13)，其中F,, i i.和G,-1.. i是根据位于相邻单元中心的值计算出的数值。 让我们在此回顾一下，可压缩CFD中FV方案的首次开发是由于MacCormack等人在七十年代初（例如，见[3]）。 遵循Rezgui [4]，如果Fj+i,k = Fj+i,k + 0(hP) Gjik+i=GjikH+o(hr)，则方案（6.13）在任何平滑的精确解（6.14）中，则在FV意义上的顺序p上是准确的。 应该注意的是，对于一般网格来说，这种精度要求比基于整个方案截断误差的通常定义更容易满足。 在FV意义上，在p阶精确的方案在通常意义上在p-1阶（至少）是准确的。 然而，通常的定义并不适合FV方法。 例如，考虑将经典的上风方案应用于不规则网格上的一维双曲方程。 众所周知（见Turkel [5]），在使用时，这个方案在通常意义上是不一致的

非均匀网格上的高阶方案 j-l,k jl±L j,k j.k-l j+l.l «yt 5x。 J图3细胞Sljik。一个直截了当的FV配方。 但它在实践中效果相当好，Sanders [6]用适当的规范证明了其收敛。 在本文中，我们将从各种多维数值实验中表明，FV意义上的三阶方案会产生一个误差（数值解与精确解的距离），对于完全不规则的笛卡尔网格来说实际上是三阶。 6.3.2 数值通量的基于残数的校正 现在让我们在一般笛卡尔网格上扩展三阶方案（6.4）-（6.7）。 当前单元格fi\textasciicircum{}是一个以\{XJ, yu)为中心的矩形，边缘长度为Sx\textasciitilde{}和Syk（见图3）和h = max \{max5Xj, max Syk）。j k 直接扩展（通常的）包括使用具有可变空间步进的公式（6.4- 6.7），即在（6.5）和<\&i的定义中，2美元：1（Sx）j+ik = -（6xj + 5xj+1），（sv）j，k = < Syk \{Syk + Syk+i）\{Sx）jtk+L = SXJ，以及（6.7）：\{Sx）jtk = SXJ 这个简单的扩展将被称为方案的非加权版本。 对于不规则的笛卡尔网格，它在FV意义上只有一阶准确，即它满足p = 1的条件（6.14）。 我们现在用p - 3构建满足条件（6.14）的数值通量F和G。 在一般的笛卡尔网格上，精确通量的表达式为：Fi+i，k Syk Jf. fdy，G 3，k+i j+i，* SXJ Jf g dx。 （6.15）i，fc+*

96 LERAT、CORRE和HANSS在\textasciicircum{}方向上执行泰勒展开，我们得到：-5v2 Fj+hk = fj+hk + -\textasciicircum{}(fyy)j+i,k + 0(h4), (6.16)，其中/J+I\textasciicircum{}表示边缘T+i k中间的精确通量/-要完成离散化，我们必须使用六个单元格中心（j, k)、(j+1, k)、(j, k±l)的值近似/和fyy。 在一般的笛卡尔网格上，让我们在边缘中间定义2点平均值和数值导数：(Vlf)j+Xik - 2(fj,k + fj+l,k), (Ml/)j+Iife = (1 - Oj)fj,k + Qjfj+l,k I。 U)j+±,k XJ+I-XJ Sxj+i + Sxj 其中9j = (a; -+i - XJ)/(XJ+I - Xj) - 8XJ/(6XJ + SXJ+I)。 运算符/ii和Vi是一阶准确，/ii是二阶准确。 在y方向上，运算符/i2、J12和V2的定义相似。 让我们也使用实际的单元中心位置在单元格中心引入3点数值导数：在（j，k），Vi是\{fj+P，k）p=-i，o，i在二阶近似fx的线性组合，Vi2是在lst阶近似fx的类似组合，V2，V2是y方向的类似算子。 让我们回到（6.16）中fj+ik的近似值。 我们首先考虑经典平均值n\textbackslash\{\}。 Taylor展开给出：(Mi/)j+i,fc= /j+i.fc + iOfoj+i -Sxj)(fx)j+i,k,fiir> +U5x2i+i + **;)(/**),•+*,* + o(h3). iU/j 在基于残差的方法中，想法不是直接取消误差，而是添加新项，以便在截断误差中恢复残差r（r - fx + gy）的导数。 因此，对于（6.17）中涉及fx的第一个误差项，我们在2阶处添加g\textasciicircum{}项的离散化，对于涉及fxx的（6.17）中的第二个误差项，我们在1阶处添加gxy-teim的离散化。 因此，我们考虑fj+ik的修改近似：2 > f*+±tk = [filf + 4(\textasciicircum{} + 1 " Sxj)fiiV29 + jg\textasciicircum{}+l + fa2)VlV2\textasciicircum{} + l,fe-显然，这个运算符可以扩展为：f;+i>k = /,-+!,* + \textbackslash\{\}(Sxj+1 - 5Xj)rj+r\_tk + \textasciicircum{}(Sx2j+1 + 6x])(rx)j+hk + 0(h3)。 我们现在转向（6.16）的第二学期。 由于一阶精度在这里已经足够了，我们近似（fyy）j+itk wrtn tne运算符V22£ti-最后，

非均匀网格上的高阶方案97数值通量的非耗散部分可以写成：o，G-可以找到类似的表达式。因此，非耗散数值通量读作：F]HM = [(/ + \textasciicircum{}V22)/ii / + \textbackslash\{\}(Sxj+1 - fojO/TiVas ° n £ 2 1 Gj,fc+i = [\{1+ \textasciicircum{}-VI2)A*2 3 + -\textasciicircum{}\{Syk+i -Syk)faVif + \textasciicircum{}(\textasciicircum{}+i + \textasciicircum{})V2V1/]\textasciicircum{}+, (6.18)通过构造，对于光滑的精确稳态解 (r = 0)，我们得到：\textasciicircum{}+iifc = \textasciicircum{}+J,fc + 0(/l3) Glk+h=Glk+,+0(h3) 其中F和G是（6.15)定义的精确值。 备注：a）对于常规笛卡尔网格，通量（6.18）准确地返回形式（6.10）。o b）F的构造也可以从（6.16）中fj+i.\textasciicircum{}的第一次评估开始完成，使用/2i而不是（ii。 由于fii是二阶精确的，这将保存一个修正项，并导致：Fj+i>k = [£i/ + \textasciicircum{}V22Mi/ + -fej+1\textasciicircum{}V1V2ff]j"+,]fe o和G-的类似表达式，我们没有选择这种方法，因为它比上面详述的方法更难推广到Navier-Stokes 方程。 我们现在考虑一般笛卡尔网格上基于残差的耗散的离散化。 由于该项是一阶的，其导数只需要以二阶精度近似。 这是通过再次使用基于残差的校正来实现的。 更准确地说，数值耗散的计算方式如下：Sxj+\textbackslash\{\} + 8XJ 4 3+2,k 4 <• 3 + 2,k /fiiq\textbackslash\{\} Syk+i + Sykf - - \textasciicircum{}-w> JJ,k+± A \textasciicircum{}\textasciicircum{}"1J "'»»j,k+ji

\section{6.4 Navier-Stokes求解器}
98 LERAT、CORRE和HANSS 为了解释它是如何工作的，让我们考虑例如d-\textbackslash\{\}\_。 当然，运算符Vi只有一阶准确：(Vi/)j+i,jt = (fx)j+i,k + 4(\textasciicircum{}+1 - 5xj)Uxx)j+i,k + 0(h2)，但m也是如此，我们得到：(MlV25)i+I,fc = \{9y)j+\textbackslash\{\},k + 4(\textasciicircum{}+1 - Sxj)(9yx)j+±,k + 0(h2)，因此对于光滑的精确稳态解，我们得到：(rfi)”+i k = 0(h3)。 总数值通量F和G仍然由（6.6）定义，新的单元值是根据有限体积方案计算的，即对于显式版本：Aw expl (5XjSyk)-\textasciicircum{}- + SykF? +iik - 6ykF？ \_Hk + SxjG]\textasciicircum{} - SxjG\textasciicircum{}, = 0 也可以写成：Awfkpl \{hF)njk (52G)"k / x J'k = - f J'k - f •?' Fc。 6.20 在OXJ oyk，即使在完全不规则的笛卡尔网格上，该方案在FV意义上的稳态下也是三阶精度。 当网格是规则时，它的3 x 3点模版不会大于所需的。 在下文中，它将被称为加权版本，以便与第6.3.2节开始时描述的非加权版本进行比较。 6.4 Navier-Stokes求解器6.4.1 常规笛卡尔网格上的三阶方案 通过添加残差中的黏性项，包括数值耗散中的残差，可以将基于残差的紧凑性扩展到Navier-Stokes 方程。 为了简洁起见，我们在这里介绍了简化系统的方法：u>t + fx + 9y = v wyy (6.21)

非均匀网格上的高阶方案99 CFL,, = A A t / 8 图4显式三阶方案（黏性问题）的二维稳定性，其中/-f（w）和g = g（w）是欧拉通量，v是正常数系数。 在常规笛卡尔网格上，方案可以写成：a）没有数值耗散的数值通量：Gjlk+i = K1 1 8\%w。 (6.22) b) 数值耗散: <*>",*+* \% [*2( Sx ' 5y 5y 5x (6.23) 总数值通量和新单元格值仍然由(6.6)和(6.7)定义。 这个方案涉及（3 x 3）+ 2 = 11分（完整的Navier-Stokes 方程为13分）。 为了均衡的顺利解决。 (6.21)，方案截断误差为：wt 6x\textasciicircum{} O + g ryy 5JL\&ir)x-6-l.(*2r)y + 0\{h*)，其中r = fx + gy - vwyy。 因此，在稳态下可以获得三阶精度。 稳定性域如图所示。 4用于标量方程wt + Awx = vwyy。 可以将隐式阶段添加到方案中，以获得无条件的线性稳定性（见[2]）。

\section{6.5 数值实验}
100 LERAT、CORRE和HANSS 6.4.2 不规则笛卡尔网格上的三阶方案，如第6.3.2节所述，我们可以将Navier-Stokes求解器扩展到不规则笛卡尔网格。 该方案现在定义为：a）没有数值耗散的数值通量：fj+\textasciicircum{}k = l(I+S-iv22)fnf + -(8xj+1 - SXJ)\textasciicircum{}(V2 9- \textasciicircum{}V» + ±(6x*j+1 + SxpV\textasciicircum{}g- \textasciicircum{}»TjHM O n X 2 1 „ Gj)fc+i = [(/+ \textasciicircum{}-Vi2)(/i25 - \textasciicircum{}V2w) + -(Syk+1 - 5yk)il2Vif + U\&VI+1 + \textasciicircum{})V2V! / - \textasciicircum{}(Syk+1 + 6yk)2V23w\}lk+k (6.24) 其中(y\textbackslash\{\}w)jtk表示从细胞中心（j, k±2)、(j, k±l)和(j, k)的值计算的wyy的二阶近似值，而(V23w)j,fc+i是从(j, k + 2)、(j,k + l)、(j,k)和\{j,k-l推断出的wyyy的一阶近似。b) 数值耗散: (di)]+hk = SXj+\textbackslash\{\}+ 5Xj [*i(V! / + MiV2 克 - »nMw）]？ +iJt (d2)lk+h = 6yk+1 +6yk [\textasciicircum{}(\textasciicircum{}Vi / + v2 g - \textasciicircum{}lw)\}lk+h (6.25) 其中(V2iu),- d2中的fc+i表示从（j, fc+|)的单元中心（j, k - 1）、(j, k)、(j, k + 1)和(j, k + 2)的值计算出wyy的一阶近似值，使其一阶误差与来自非黏性项的误差相结合，产生残差的y导数，即：(V2w)iifc+i = (wyy)jtk+x + - -(6yk+1 - 5yk)\{wyyy)ik+i\_ + 0(h2)。 请注意，当v = 0时，公式（6.24）和（6.25）分别减少到（6.18）和（6.19）。 此外，当网格规则时，公式（6.25）减少到（6.23）。 6.5 数值实验 在本节中，我们将第6.3.2节中定义的基于残差的方案的非加权和加权FV版本应用于Euler 方程和第6.4.2节中为Navier-Stokes 方程的测试案例。 首先考虑具有已知分析解决方案的黏性和黏性问题（[7]中提出了测试案例6.5.2和6.5.3）。 他们将允许演示，

非均匀网格101的高阶方案通过计算精确解和数值解之间的误差的L2规范，加权版本即使在完全不规则的笛卡尔网格上也有三阶误差，而非加权版本在同一网格上严重损失了准确性。 两个版本还将针对涉及不连续性的问题进行比较。 6.5.1 三个网格系列 三种类型的笛卡尔网格用于以下计算：均匀（图。 5（a）），几何拉伸（图。 5（b））拉伸系数为1.11，并随机扰动（图。 5（c））。 在后一种情况下，我们从统一网格开始，在x和y方向上对每个单元格应用随机缩放系数；然后将网格重新缩放到[0, l]2，并且相邻单元格之间不存在一般关系。 对于每种类型的网格，首先生成一个39 x 39的网格，然后通过在每个空间方向上将每个单元格除以2，从粗网格中推导出一个78 x 78和一个156 x 156的网格。 每个3个网格系列都允许计算精度顺序。 请注意，对于随机扰动的网格系列，特征网格大小h的值可能会从一个计算系列更改为另一个计算系列。 6.5.2 2-D旋转对流 此测试案例包括光滑高斯剖面在正方形域上的旋转对流[0, l]2。 更准确地说，我们查看以下问题的稳态解 f ff+yif + a-*)\textasciicircum{}0' (\textasciicircum{},2/)G]0,l[2 i>0 w(x,y,0) = 0, (x,y)e]0,l\{2 < w(x,0,t) = e[-50(x-o.5)2]))是[0,1], w(x,l,t) = 0, iG[0,l], w(0,y,t) = 0, J/e[0,l]，其确切的稳态解由wexact\{x,y) = et-\textasciicircum{}o.a-Va\textasciicircum{}\textasciicircum{})))\textasciicircum{} (ar>y) e [0> 1]2，这意味着w沿y = 0的初始分布在任何中心圆（1,0）上都守恒的。 首先对均匀网格进行一系列计算，我们观察到非加权和加权方案都是三阶准确的，实际上产生相同的误差（见图。 6（a））。 这是意料之中的，因为两个版本都导致了统一网格上的相同方案。 接下来，对非规则网格（拉伸或随机扰动）的计算表明，非加权版本的准确性严重损失，其误差顺序下降到约1.4，而t > 0 t > 0 t > 0

102 LERAT、CORRE和HANSS -I Tmtl（b）（c）图5 39 x 39个单元格的笛卡尔网格：（a）均匀（b）几何拉伸（c）随机扰动加权版本，满足FV精度标准（6.14），p = 3有三阶误差（见图。 6（b）和（c））。 这在解等值上也清晰可见（图。 7）：虽然随机网格上的加权溶液叠加在确切的溶液上，但非加权溶液表现出由网格不规则引起的强烈扰动。 此案例的收敛历史通过域上的Aw/At的L2规范进行监控，并绘制在图6（d）上；两个版本共享相同的非加权隐式阶段（6.8）。 人们可以注意到，加权方案比非加权方案更快地产生收敛到稳态。

非均匀网格上的高阶方案 103 6.5.3 3D螺旋对流 此测试案例是上一个问题的3D扩展；我们在域-1 < a上求解wt + z wx + Q.l-Wy - xwz = 0; < 0, 0 < y < 1，0 < 2 < 1，具有适当的初始和边界条件。 更准确地说，我们在z = 0平面（x和y中的2-D高斯分布）中设置了一个平滑的初始条件，该条件既围绕y轴旋转，又沿着同一轴对流（见图8（a））。 使用一系列32 x 32 x 32、64 x 64 x 64和128 x 128 x 128均匀、几何拉伸和随机扰动的笛卡尔网格进行计算。 然而，从现在开始，为了简洁起见，我们将把计算的精度顺序和数值解的呈现限制在随机扰动的笛卡尔网格上获得的结果；事实上，观察到后一个系列使网格不规则性的影响比几何拉伸系列更明显。 在图8（b）中可以观察到，FV加权方案在随机扰动网格上仍然是真正的三阶精度，而非加权版本的精度顺序下降到1.79。 当与精确解相比，这种精度损失在流出平面x = 0中数值解的等值上可见（见图。 9）：加权溶液没有非加权溶液中存在的所有失衡。 6.5.4 2-D Poiseuille流模型 我们现在考虑由以下方程定义的2-D黏性问题：' ff + «tr = \textasciicircum{} (*,y)e]o,i[2, t>o w(x,y,0) = l, \{x,y)e\}0,l[2 w(x,0,t) = w(x,L,i) = 0, ze]0,l[, i>0. ti>(0,j,0 = sin(7r£), ye]0,l[, i>0 其确切的稳态解由wexact\{x, y) = e[\_7r «x] sin(Try), (Try) £ [0, l]2 这个问题可以被视为对浮石流建模：在流入y - 0处规定的初始轮廓在两个实心之间进行对 在这里，我们再次观察到，基于残差方案的加权版本可以在随机扰动的笛卡尔网格上保留真正的三阶误差，而通常的非加权版本则简化为一阶准确方法（见图 10（a））。 加权版本提供的精度要高得多，如图ll所示，其中绘制了数值解的等值：加权版本的等值是

104 LERAT、CORRE和HANSS叠加在确切的解决方案上，而非加权版本显示由网格不规则引起的强烈失真。 另请注意，隐式加权版本的收敛到稳态再次比隐式非加权版本的收敛更快（图。 10（b））。 6.5.5 2-D Brgers问题 此测试用例处理以下非线性问题：t > 0 t>0 t > 0，其中选择左右状态wi，wr，以产生压缩解析为正态（案例1：wi - 1，wr - -1）或斜（案例2：W[ - 1.5，wr = -0.5）激波。 我们观察到，在这两种情况下，加权方案在随机扰动网格上的压缩区域产生直等线，而非加权版本提供的特征由于网格不规则性而更加扰动（见图。 12和13（a）-（b）-（c））。 在案例1中，当冲击与网格对齐时，加权解非常接近精确的解，几乎没有显示振荡，而非加权解略有振荡（见图 12（d））。 在案例2中，当冲击不再与网格对齐时，两个版本都显示一些振荡，但加权方案产生的振荡在本例中比非加权方案显示的振荡要弱得多（见图 13（d））。 请注意，被随机扰动的网格，这些振荡的振幅在另一次运行中不一定相同。 然而，从对该测试案例进行的许多计算得出结论，加权方案的振荡性系统性低于非加权方案。 6.5.6 2D可压缩射流相互作用 这个问题[8]由两个水平超音速射流的相互作用组成。 由M-4、p-0.5、p=0.25定义的上流和由M=2.4、p=1、p=1定义的下流突然在（x=0，y=1/2）处接触，并研究它们的相互作用在域[0，l]2中。 这种相互作用产生一个膨胀扇和一个激波，分别在高压和低压区域传播，以及前两个波后面的不同密度和速度导致的接触不连续性。\%+dW + i\% = \textasciicircum{} (*,y)e]o,i[2, w(x, y, 0) = (1 - x)wi + xwr (x, y) e [0, l]2 \textbackslash\{\} w(x,0,t) - (1 - x)u>i + xwr, a;€[0,l], w(0,y,t) = wi, ye [0,1], w(l,y,t) = wr, ye [0,1],

\section{6.6 结论}
\section{参考资料}
非均匀网格上的高阶方案 105 从随机扰动网格上非加权和加权解决方案的马赫轮廓来看，后者始终更好地表示了流动中存在的波（见图。 14）。 从沿流出边界的马赫数分布可以看出（见图15），加权方案产生更尖锐的剪切层和冲击（尽管略有振荡），以及更接近确切解决方案的膨胀扇；波之间的高原也更好地表现。 6.6 结论 提出了一种方法来设计一个三阶准确方案，该方案保持其精度顺序，而无需对网格平滑度做出任何假设。 遵循[4]，这个想法是在FV框架中，在三阶处接近控制体表面的精确通量；通过在不规则的笛卡尔网格上使用基于残差的校正，以紧凑的方式实现了这一点。 对于非黏性和黏性多维问题，即使在完全不规则的笛卡尔网格上也得到了三阶实际误差。 这项工作的持续发展包括将这些想法扩展到一般的不规则结构网格。 参考资料1。 莱拉特A。 \& Corre C，基于残基的守恒定律多维双曲系统的紧凑方案，座谈会“CFD的艺术”，法国马赛，1999年9月，将发表在Comp. and Fluids上。 2. 莱拉特A。 \& Corre C，一个紧凑的三阶精确方案，使用可压缩Navier-Stokes 方程的一阶耗散，提交给J。 补偿。 物理.. 3. MacCormack R.W. \& Paullay A.J.，通过有限差分算子的时间分割实现的计算效率，AIAA论文72-154，1972年。 4. Rezgui A.，有限体积方法的准确性和收敛分析，CFD杂志，8（3）：369-377，1999年。 5. Turkel E.，可压缩流体流动的非均匀网格方案的准确性，Appl. 数字。 数学，2：529-550，1985年。 6. Sanders R.，关于具有可变空间微分的单调有限差分方案的收敛，数学。 Comp.，40：91-106，1983年。 7. Deconinck H.，Struijs R.，Bourgois G. \& Roe P.L.，非结构化电网上的紧凑对流方案，VKI LS 1993-04，1993。 8. 格拉兹H.M. \& Wardlaw A.B.，稳定超音速气体动力学的高阶戈德努诺夫方案，J. 补偿。 物理，58（2），1985年。

106 LERAT、CORRE和HANSS -2.8 -3.2 -3.4 -3.6 5-3.8 -4.6 -4.8 -5.2 -5.4 T r r 7 / slope = 3.00 B - -e-! • • - 非加权加权 • ' -2.3 -2.2 -2.1 -2 -1.9 -1.8 -1.7 对数（空格步进） -2.2 -2.4 -2.6 -2.8 --3 3 3-3.2 X-6 H' 3 -3.8 -4.2 -4.6 -4.8 I,. iln -1.8 Kill -1.7 slope =1.38 / / \_ \_ ''' -1.6 -1.5 -1.4 - \textasciicircum{}\textasciicircum{} / J2T / 0 / / slope = 2.98 -e 非加权 - © - - - 加权 • •••••••••••>....i... 1.3 -1.2 -1.1 -1 -0.9 Log（空格步）（b）I -3.5 O -4 -B 非加权 - © - - - 加权 -2 -1.9 1.8 -1.7 -1.6 -1.5 -1.4 -1.3 Log（空格步）（c） 5 -（剩余 O CO 8° -12 -14 -16 -18 - >\textbackslash\{\} \textbackslash\{\}\textbackslash\{\}> " \textasciicircum{} \textbackslash\{\}: \textbackslash\{\} 7 \textbackslash\{\} - \textbackslash\{\} -- 非加权加权 0 50 100 迭代（d）图 6 旋转对流。 计算误差顺序（a）均匀网格（b）拉伸网格（c）随机网格（d）收敛到随机78 x 78网格上的稳态

非均匀网格上的高阶方案107（a）（b）图7随机39 x 39网格上w（x，y）的等值（a）非加权计算（b）加权计算>dslope = 3.02.. p -S非加权-O- -加权-1.7 -1.6 -1.5 -1.4 -1.3 -1.2 Log（空格步）（a）（b）图8（a）精确解（b）计算误差顺序

108 LERAT、CORRE和HANSS（a）（b）图9随机32 x 32 x 32网格上的流出（a）精确解决方案（b）非加权方案（c）加权方案

非均匀网格上的高阶方案 109 坡度 = 1.02 N £ -5.5 4.5 - - -5 -a 非加权 -Q 加权坡度 = 2.94 -2 -1.9 -1.8 Log (步幅) (a) 3. \_ -6 o 3 w -8 0) S> -10 o \_l -12 -14 -1R ry " X \textbackslash\{\}\textbackslash\{\}: - '•,, 加权非加权 \textbackslash\{\} \textbackslash\{\} \textbackslash\{\} \textbackslash\{\} \textbackslash\{\} V \textbackslash\{\} \textbackslash\{\} \textbackslash\{\} \textbackslash\{\} \textbackslash\{\},, '。 \textasciitilde{}Xl\textasciitilde{}，\textasciitilde{}，100次迭代（b）图10 2-D Poiseuille流模型，（a）计算误差顺序（b）收敛到78 x 78网格上的稳态（a）（b）图11随机39 x 39网格上w（x，y）的等值（a）非加权计算（b）加权计算

110 LERAT、CORRE和HANSS 0.9 0.8 0.7 0.6 >- 0.5 0.4 0.3 0.2 0.1 " M ' // ' /// - ////: ////,/ I.I \textbackslash\{\}m \textbackslash\{\}\textbackslash\{\}\textbackslash\{\}v\textbackslash\{\} I \textbackslash\{\} \textbackslash\{\}. \textbackslash\{\}W\textbackslash\{\} 0.9 0.8 0.7 0.6 >- 0.5 0.4 0.3 0.2 0.1 \textasciicircum{} 111, ' //\textbackslash\{\} ' /// \textbackslash\{\} \textbackslash\{\}////f k I, m \textbackslash\{\}\textbackslash\{\}xk \textbackslash\{\}m (a) (b) 精确非加权加权 (c) (d) 图 12 随机 39 x 39 网格上的等值（easel) (a) 精确解，(b) 非加权解 (c) 加权解 (d) y = 0.7 处的解

非均匀网格上的高阶方案111（a）（b）精确非加权加权（c）（d）图13随机39 x 39网格上的等值（案例2）（a）精确解，（b）非加权解（c）加权解（d）y = 0.7的解

112 LERAT、COR.R.E和HANSS（a）（b）图14随机78 x 78网格上的马赫轮廓（a）非加权解（b）加权解0.9 0.8 0.7 0.6 >• 0.5 0.4 0.3 0.2 0.1非加权精确马赫（a）0.9 0.8 0.7 0.6 >。 0.5 0.4 0.3 0.2 0.1 加权精确（b）图15沿流出边界a的马赫轮廓；= 1，在随机78 x 78网格上（a）非加权（b）加权

\clearpage
\chapter{7 计算可压缩流的未来方向：高阶居中与多维上卷}
\section{7.1 介绍}
计算可压缩流的未来方向：高阶居中与多维上卷M。 那不勒斯1，A。 Bonfiglioli2，P。 Cinnella3，P。 De Palma4和G。 Pascazio1 7.1 介绍 在过去的几十年里，计算机性能在速度和内存大小方面都得到了显著提高，因此给定计算的相对成本每十年降低了大约十倍[1]。 与此同时，计算流体动力学（CFD）经历了指数级增长，因此现代飞机、先进涡轮机械和内燃机的设计和开发发生了巨大变化。 事实上，现在有可能在CPU时间内使用数百万个计算单元，在CPU时间内计算一个非常复杂的流场（即在整个飞机周围或涡轮机的一个或多个叶片通道内）。 因此，在通过非常快速和廉价的计算机运行执行所有初步设计后，仍然必要但非常昂贵的实验仅限于最终设计选择。 如今，涡轮机械应用的CFD代码基于数字方法1 DIMeG，Politecnico di Bari，70125 Bari，意大利，napolita@poliba.it。 2 DIFA Universita della Basilicata, 85100 Potenza, Italy。 3 SINUMEF实验室，ENSAM，75013巴黎，法国。 目前在巴里的DIMeG。 4 DIM，罗马大学“Tor Vergata”，00133罗马，意大利。 Frontiers of 计算流体动力学 - 2002 编辑：David A. Caughey和Mohamed M. 哈菲兹 ©2002 世界科学

114 NAPOLITANO等人最初用于空气动力学应用。 他们通过时间行进显式（Runge-Kutta）[2，3，4]或隐式（近似因数分解[5]，线松弛[6]）方案解决稳态雷诺平均可压缩Navier-Stokes 方程，他们的收敛速率通过各种技术加速，如局部时间推进、隐式残差平滑（对于显式方案）、多网格等。 [2，3，4，7]。 就空间离散化而言，有限体积（或元素）方法几乎普遍适用于方程的守恒定律形式，通过保守的离散化，以便正确捕获冲击和接触表面等流动不连续性。 如果方程中的对流项使用中心差分[2]离散化，则避免伪振荡所需的人工耗散要么被添加到方案中，要么“自然”包含在方案本身中，如果“上风”离散化用于此类术语[8、9、10、11、12]。 最近，这些方法已被扩展，通过适当的技术，如双时间推进 [13]来处理时间依赖性问题。 最后，为了解决工程兴趣的流程，上述方法采用复杂度递增的湍流模型（代数模型[14]、微分模型[15、16、17、18]、大涡流模拟[19]）。 事实上，迄今为止，对于实际问题来说，解决非稳定Navier-Stokes 方程的所有时间和长度尺度仍然不可行，即执行直接数值模拟 [20]。 因此，由于过渡和湍流建模的严重限制，以及不能正确解释可压缩流动的多维性质的数值方法，CFD仍然远未成为航空航天和涡轮机械应用中的“独特设计工具”。 在最后一个方面，作者最近为解决可压缩欧拉和Navier-Stokes 方程的两种最先进的方法的开发做出了贡献。 第一种方法是有限体积中心方案，它使用加权平均五点离散化作为非黏性通量，在均匀和轻度非均匀网格上实现三阶精度。 第二种方法是多维上风波动分裂（FS）方案，它“只有”二阶准确，但允许对可压缩流动传播现象进行更真实的建模。 这两种方法都是专门为可压缩的非黏性流动而设计的，并通过使用黏性通量的标准二阶准离散化来扩展到Navier-Stokes 方程的溶液。 本文提供了这两种方法的有价值的数值比较，使用它们来计算从不黏性平滑流到带冲击的黏性流的既定测试案例。 在所有情况下，都使用了三个相同的网格，以便始终实现网格收敛，并评估两种方法的网格灵敏度。 在对两种方法进行简要描述后，将为每个选定的测试案例提供详细的数值结果。 最后，一些结论

\section{7.2 高阶居中数值法}
未来方向将提供115个评论。 7.2 高阶中心数值方法 当前方法[21]是可压缩欧拉和Navier-Stokes 方程的高阶以细胞为中心的有限体积方案（HO）。 该方案的主要特点是引入对Euler 方程的经典二阶方案[2]的空间离散化产生的二阶色散误差进行校正，并在评估非黏性通量时使用加权平均值，以便在轻度不规则曲线网格上保持有效的三阶精度（对于稳定的Euler 方程）。 在这里，将介绍方案的主要特征，适用于以守恒形式编写的一维Euler 方程：Ut + Ft = 0。 （7.1）通过标准的泰勒级数展开，人们可以很容易地表明：1 Sr2 Ut\textbackslash\{\}j + fo8nF\textbackslash\{\}j - -Fxxx = Ut + Fx + 0（5x4），（7.2）其中（<ty）j+i/2 = 4>i+i - <t>j，QWOj+i/2 = \textasciicircum{}（4>j+i + <t>j）-因此，如果在方程（7.2）的左侧（LHS）引入Fxxx的二阶精确中心五点离散化，它提供了方程（7.1）的四阶精确非耗散离散化。 然后，将耗散校正引入为标准的詹姆森人工耗散 [22]，以给出：Ut\textbackslash\{\}j + j\textasciicircum{}S UF - \textbackslash\{\}PliF\textbackslash\{\} - \textasciicircum{}-8 [e2p(A)SU + e4p(A)S3U] = 0, (7.3) with £2|j+i = fc2max\{i/j,i>j+i\}, £4|J+I - max\{0,fc4 -\textasciicircum{}Ij+i\}, v = pj+i-2pj+pj\_1 3 Pj+i + 2Pj+Pj-i 在上述方程中，p(A)表示平均Jacobian 矩阵 A的光谱半径，p是压力，而k2和k\textbackslash\{\}常数是参数（k\textasciicircum{} = 0.032和k2 = 0，分别为超音速和超声速流动）。 在U光滑的区域，e2 = 0（5x2）和£4 = O（l），

\section{7.3 波动分裂方法}
116 NAPOLITANO等人。因此方程（7.3）LHS中的第三个项是0（5x3），方案是三阶准确的。 上述方法通过以细胞为中心的有限体积公式扩展到多维结构网格。 细胞重心处的保守变量被认为是因变量，数值通量使用适当称重的离散化公式进行评估，该公式考虑到了网格的拉伸和倾斜度，详情见[23]。 数值方法最终使用黏性通量的二阶准确中心离散化扩展到Navier-Stokes 方程。 对于目前的稳态计算，使用了四阶段Runge-Kutta时间积分方法[2]，系数：a,\textbackslash\{\} = 1/4, a\textasciicircum{} - 1/3, a\textasciicircum{} = 1/2和a4 = 1。 隐性残留平滑[24，25]、局部时间推进和V循环多重网格法[7]被实施，以加速收敛到稳态。 标准特征边界条件强加于流入点和流出点。 在无黏流动计算的情况下，不透气条件施加在墙壁上，使用格林定理评估的细胞平均压力梯度，从内细胞点外推压力。 在黏性流动的情况下，计算出壁上的压力施加零压力梯度，评估温度执行零热通量。 7.3 波动分裂方法 Euler 方程在由线性有限元（三角形）组成的计算域上离散化。 每个三角形T的离散保守通量平衡，即波动，可以用通过每个三角形边的适当通量来表示（例如，见[26，27]），如下：3，3 \$UIT = R*T，\$T = -Y，iA-niQi = -I2KiQi-（7-4）3 = 1 J=l在方程（7.4）中，A是特征变量Q的雅各布张量，R是从Q到保守变量U的单元平均投影矩阵，lj是与节点j相反的三角形边的长度，rij是垂直于lj的内向单位向量，Ki = 2li iAnU + BnvJ），（7.5）£和rj是自然坐标。 由于系统的双曲性质，Kj可以写成：Kj = \{RKKKlK)j = (RKA + LK)j + \{RKA-KLK)j = K+ + Kj。 （7.6）

未来方向117在方程（7.6）中，RK、J和LK、J是Kj的左右特征向量矩阵，而A\textasciicircum{}-•和A\textasciicircum{}•是相应的正负特征值矩阵。 引入以下向量，（7.7）线性矩阵低扩散A（LDA）方案，是线性保持[28]，得到如下：\$,- = -K+ [Qout - Q-m] • (7.8) 方程的LDA方案（7.8）对于亚音速平滑流非常准确，可以被认为是此类条件下的最佳紧致波动分裂（FS）方案。 对于超音速流，一组不同的特征变量W允许将欧拉系统重新转换为一组等效的四个标量对流方程。 然后，这些由非线性正流不变（PSI）方案[28]离散化，它可以写成\$,- = -Kf+ (Wj - W£), (7.9) -\textasciicircum{}nl，其中K? 使用[29]中定义的非线性雅各布A进行计算。 迄今为止，PSI方程（7.9）是带或不带冲击的超声速流动的最佳FS方案。 最后，为了计算具有强冲击的跨声速流，有必要使用局部单调方案，即[27]的矩阵N方案，给出为：\$,- = -Kl+ (Wj - Wfn)，（7.10），其中Kj和W\{n使用雅各布Aw计算•为了确定需要使用这种低阶方案的“冲击单元”--忽略那些从超声速流动条件到亚音速条件的过渡足够平稳，可以由LDA方案正确处理的冲击单元--有必要对它们进行独特的表征。 通过对正常冲击的流动特性的仔细分析，得出的结论是“冲击细胞”的特点是：i）平均细胞马赫数，M，低于一个；ii）至少一个超音速节点；iii）至少一个局部马赫数，Mj，低于0.9的亚音速节点。 总之，在计算过程的每个步骤中，目前的混合方法标记了计算域的所有单元，并使用超音速单元的方程（7.9）、冲击单元的方程（7.10）和其余单元的方程（7.8）来分配每个残差。 然后通过收集来自计算域的每个顶点j处的Euler 方程残差进行评估

\section{7.4 结果和讨论}
118 NAPOLITANO等人。网格64 x 16 128 x 32 256 x 64 •Mmin HO 0.4159 0.4244 0.4228 FS 0.4189 0.4181 0.4177 M-max HO 0.9163 0.9253 0.9333 FS 0.9403 0.9401 0.9401 Li(s)(xl0-a) HO 53.33 8.166 1.476 FS 197.8 62.55 27.58 Loo(s)(xl0-a) HO 318.0 39.80 5.220 FS 126.0 36.03 11.05 表1 通道流精度研究。周围三角形，如：\textbackslash\{\} / j •> T 数值方法最终使用标准的二阶精确加勒金有限元方案扩展到Navier-Stokes 方程。 对于目前的稳态计算，使用了两种不同的时间积分方法，即[27]的显式Runge-Kutta方案和[30]的隐式Newton-GMRES方案。 标准特征边界条件强加于流入点和流出点。 在显式非黏性代码中，使用一组辅助幽灵细胞强制执行壁边界条件：等熵简单径向平衡与特征校正一起使用，以强制执行无注射条件并评估壁压力[31]。 在隐式代码中，在非黏性流动的情况下，动量通量在每个壁细胞侧强制执行，压力近似为两个节点值的平均值；在黏性流动的情况下，壁节点的两个速度分量都设置为零，零热通量条件通过省略边界热通量贡献来自然执行。 7.4 结果和讨论 考虑了五个有据可查的测试案例，以评估两种数值方法的准确性-性能，使用三种分辨率不同的网格。 第一个测试案例是不黏性亚声速流动通过一个带有协弦形壁的通道，20\%限制和出口马赫数等于0.5。 图1提供了使用FS方法的精细（256 x 64）网格计算的马赫数等高线，相应的HO解实质上是一致的。 表1显示了最小和最大马赫

未来方向119图1通道流马赫数等高线（AM = 0.02）。数字和熵误差的L\textbackslash\{\}和Loo规范（S =（（p/' p1）-\{vlP1）入口）/\{p/P\textasciicircum{}）inlet）分别为HO方法和FS方法。 Zoo规范表明，这两种方法几乎分别是三阶准确和二阶准确，但L\textbackslash\{\}规范对两者的减少速度较慢。 此外，在使用粗网格和中网格时，FS解决方案似乎更准确，对网格大小不太敏感，而HO方法的L\textbackslash\{\}规范随着网格的细化而明显降低。 图2显示了使用三个网格获得的沿底壁的熵误差分布：对于粗网格，FS解决方案提供了较低的误差；在中等网格上获得了非常接近的误差分布；HO方法在精细网格上提供了较低的误差。 此外，与HO方法不同，FS方法给出了单调的解决方案。 值得注意的是，FS方案的出色性能是由于在显式代码中实现的非常准确的墙边界条件。 事实上，隐式代码的熵误差值高出一个数量级，并提供以下Mmin和Mmax值：（0.3955，0.9306），（0.4059，0.9370），（0.4119，0.9392），分别用于粗、中和细网格。 第二个测试案例是经过NACA0012 翼型的非黏性亚声速流动，Moo = 0.63和a = 2°。 采用的三个网格分别沿着轮廓有96、192和384个单元格，自由流边界位于距离身体20个弦的地方。 表2显示了升力和阻力系数（Cx，Co）以及两种方案熵误差的L\textbackslash\{\}和Loo规范。 L\textbackslash\{\}规范表明，这两种方法几乎实现了其设计精度顺序，而Loo规范仍然没有显示渐近行为，可能是因为存在停滞流区域。 两种方案计算的升力系数往往与精炼的网格值相同，而FS方案提供了较低的阻力系数，尽管其熵误差分布较高（见图3）。 最后，图4给出了沿着轮廓的马赫数分布。除了吸侧前部的粗网格HO溶液外，所有曲线都非常接近。 总之，本测试获得的结果表明，这两种方案的准确性相当。

120 NAPOLITANO等人。 0.004 0.003 0.002 W 0.001 图2 沿底壁的通道流熵误差分布。 第三个测试案例是经过NACA0012 翼型的非黏性跨声速流动，MQO-0.85和a = 1°。 使用了与上一个测试案例相同的网格。 对于这样的问题，显式FS代码只能在粗网格上收敛，因此必须使用隐式（使用不太准确的墙边界条件处理）。 图5和图6显示了精细网格上的两个方案获得的总体解决方案。 它们几乎相同，并显示了两种方法非常好的激波捕捉能力。 图7给出了压力系数（Cp - -2(p - POO)/(PCO«TO)）沿剖面的分布。 两个网格136 x 20 272 x 40 544 x 80 HO 0.3176 0.3207 0.3215 FS 0.3348 0.3296 0.3279 CD(xl0-4) HO 14.51 5.918 4.060 FS 12.08 4.274 3.020 In(S)(xl0-9) HO 154.9 27.54 4.266 FS 7313。 1872年。 1099。 Loo(s)(xl0-4) HO 38.61 12.51 6.420 FS 130.5 25.06 9.313 表2 Inviscid 亚声速流动经过NACA0012 翼型：准确性研究。 FS 64x16 FS 128x32 FS 256 x 64 HO 64x16 HO 128x32 HO 256x64

未来方向121图3熵误差图4沿剖面的马赫数分布。沿剖面的分布。网格136 x 20 272 x 40 544 x 80 HO 0.3841 0.3830 0.3728 FS 0.3988 0.3844 0.3702 CDCXIO-*) HO 5.748 5.793 5.742 FS 6.081 5.716 5.646表3 Inviscid 跨声速流动经过NACA0012 翼型：准确性研究。方法提供单调冲击和非常接近的解决方案。 表3显示了使用两种方案获得的升力和阻力系数。 使用精细网格，CL值的差异约为1\%，而C\&值的差异约为2\%。 结果与[32]中提供的数字数据非常一致。 在这种情况下，正如预期的那样，HO方法的准确性明显优越。 为了更好地了解隐式代码中使用的不太准确的边界条件处理的不良影响，图8显示了配置文件上的马赫数分布。 由于沿着身体表面产生过多的熵，所有网格上都低估了冲击后的马赫数值。 相比之下，显式代码的粗网格解更正确。 请注意，此类问题与以下 黏性流动 计算无关。 下一个测试案例是M\textasciicircum{} = 0.5，a = 0和雷诺数的NACA0012 翼型上有据可查的层亚声速流动，基于

122 NAPOLITANO等人。 图5 HO-scheme Mach-mumber 图6 FS-scheme Mach-number等高线（AM - 0.05）。等高线（AM = 0.05）。网格132 x 34 266 x 68 532 x 136 HO 2.097（-2）2.204（-2）2.256（-2）FS 2.251（-2）2.244（-2）2.237（-2）/\textasciitilde{}1VIS HO 3.825（-2）3.384（-2）3.3416（-2）FS 3.502（-2）3.313（-2）3.2777777。 （x/c）HO 0.9178 0.8284 0.8227 FS 0.8642 0.8255 0.8186 表4黏性亚声速流动经过NACA0012 翼型：准确性研究。弦长度和自由流条件，Re\textasciicircum{} - 5000。 使用了三个网格，分别沿着轮廓有112个、224个和448个单元格，自由流边界位于距离身体约20个弦的地方。 这种流动的主要特征是靠近后缘的分离。 表4提供了不黏性\{C™v）和黏性（C]ps）阻力系数，以及使用两种方案计算的分离点。 两组结果彼此非常接近，与文献中报告的数据非常一致[33]。 FS的网格敏感度似乎略低，因此可能更准确。 图9和图10分别显示了压力系数轮廓和皮肤摩擦系数（Cf - 2TW/\textbackslash\{\}poou\textbackslash\{\}o）的分布。 关于压力系数，除了粗网格上的HO解外，所有曲线在绘图精度范围内都是一致的。 此外，图10再次显示了两组结果之间的一致性，但差异很小

未来方向 123 FS 136x20 FS 272x40 FS 544x80 HO 136x20 HO 272x40 HO 544x80 FS 136x20 明确 i i i i i i i „£a**\textasciicircum{}\textasciicircum{}\_\textasciicircum{}\_™BB I - v\textasciicircum{}j\textasciicircum{}\textasciicircum{} \textbackslash\{\}i \textasciitilde{}rJT i B FS 136x20 FS 272x40 „。 \_ FS 544x80 HO 136x20 HO 272x40 o HO 544x80。 FS 136x20 明确的.,.,,., f it 0.5 0.5 图 7 沿剖面的压力系数分布。 图8 马赫数沿剖面分布。网格132 x 34 266 x 68 532 x 136 HO 0.3396 0.3379 0.3391 FS 0.3420 0.3421 0.3397 HO 0.2796 0.2758 0.2754 FS 0.2761 0.2746 0.2733 表5黏性超声速流动经过NACA0012 翼型：精度研究。峰值，HO最大值略低（精细网格上的0.1475与0.1480）。 对于这个问题，HO方案的结果是在半网格上获得的，强制对称性，以实现收敛到机器精度。 否则，在中细网格上，由于周期性涡流脱落现象，残差在下降到约10-3后停滞不前。 最后一个测试案例是NACA0012 翼型上的层超声速流动，Moo = 2，a = 10°，Re\textasciicircum{} \textasciitilde{} 1000。 表5显示了使用两种方案获得的升力和阻力系数。 两组结果是可比的，与[34]中提供的数字数据非常一致。 使用精细（532 x 136）网格使用HO和FS方案计算的马赫数等高线分别显示在图11和图12中。 这两种方法都相当清晰和单调地捕捉冲击。 图13和14分别给出了沿着压力系数和皮肤摩擦系数轮廓的分布。 所有解决方案在绘图中都是一致的

\section{7.5 结论}
124 NAPOLITANO等人。t i iii 0 0.25 0.5 0.75 1 图9 压力系数图10 沿着轮廓的皮肤摩擦系数分布。沿着轮廓的分布。准确性。 7.5 结论 这项工作对解决稳态可压缩欧拉和Navier-Stokes 方程的两种最先进的方法的精度性能进行了非常仔细的一对一比较。 第一种方法是加权平均有限体积的方法，它以三阶精度近似不黏性通量，以二阶精度近似黏性通量；第二种方法是混合多维上风波动分裂方案，以二阶精度近似不黏性通量和黏性通量；两者都在冲击时仅局部精确一阶。 在非黏性和黏性流动计算中，低阶FS方案的表现与HO方案一样好，如果不是更好的话，主要是由于三个原因：i）非黏性通量多维性质的正确建模；ii）非黏性壁边界条件的非常准确的处理；iii）黏性通量的Galerkin近似。 另一方面，HO方法成本更低，并且可以提高黏性通量的准确性。 总之，这两种方法都值得追求，以开发更准确、更坚固、更高效的CFD工具，用于先进的航空航天和涡轮机械应用。i.... i..... i 0 0.25 0.5 0.75

\section{7.6 致谢}
\section{参考资料}
未来方向125图11 HO方案图12 FS方案马赫数马赫数等高线（AM = 0.1）。等高线（AM - 0.1）。 7.6 致谢 本研究得到了MURST/COFIN99的支持。 参考资料1。 Tannehill J.C.、Anderson、D.A.和Pletcher R. H.，计算流体动力学和传热，第二版，泰勒和弗朗西斯，1997年。 2. Jameson，A.，Schmidt，W.，Turkel，E.，使用Runge-Kutta 时间推进方案的有限体积方法对Euler 方程进行数值模拟，AIAA论文81-1259，1981年。 3. 詹姆森 A.，使用Euler 方程的跨声速翼型计算，在航空流体力学中的数值方法，P.L. Roe（编辑），学术出版社，1982年。 4. 詹姆森 A.，计算空气动力学的成功与挑战，AIAA论文87-1184，1987年。 5. Beam，R.M.，变暖，R.F.，可压缩Navier-Stokes 方程的隐式因数方案，AIAA杂志16，1978年，pp. 393-402. 6. Napolitano，M.和Walters，R.W.，Navier-Stokes 方程的增量块线-高斯-塞德尔方法，AIAA杂志24，1986年，pp. 770-776。 7. Brandt，A.，边界值问题的多层次自适应解决方案，数学。 计算机。 31，1977年，pp. 333-390。 8. Steger，J.L.，Warming，R.F.，应用于有限差分法的非黏性气体动力学方程的通量向量分割，J. 计算机。 物理。 40，1981年，pp. 263-293。 9. van Leer，B.，Euler 方程的通量向量分割，Proc。 第8届ICNMFD，1982年，Springer Verlag。

126 纳波利塔诺等人。 图13 压力系数 图14 沿着轮廓的皮肤摩擦系数分布。沿着轮廓的分布。 10. Roe，P.L.，近似黎曼求解器，参数向量和差分方案，J. 计算机。 物理。 43，1981年，pp. 357-372。 11. Harten，A.，双曲守恒定律的高分辨率方案，J. 计算机物理。 49，1983年，pp. 357-393。 12. 哈滕，A.. Osher，S.，均匀高阶精确的非鼻振荡方案，SIAM J. 数字。 肛门。 1987年24日，第279-309页。 13. 詹姆森，A.，使用多网格的时间依赖性计算，应用非稳定流动通过机翼和机翼，AIAA论文91-1596，1991年。 14. Baldwin，B.，Lomax，H.，分离湍流流的薄层近似和代数模型，AIAA论文78-0257，1978年。 15. Launder，B.E.，Spalding，B.，湍流的数学模型，学术出版社，1972年。 16. Wilcox，D.C.，CFD的湍流建模，DCW工业，1993年。 17. Patel，V.C.，Rodi，W.，Scheurer，G.，近壁和低雷诺数流动的湍流模型：综述，AIAA杂志23，1985年，pp. 1308-1319。 18. Yakhot，V.，Orszag，S.A.，湍流的重新归一化群分析，I基本理论，J. 科学。 计算机。 1，1986年，pp. 3-51。 19. Rogallo，R.S.，Moin，P.，湍流流动的数值模拟，流体力学年度评论16，1984年，pp. 99-137。 20. Abraham, J., Magi, V., 混合层中速度比和密度比效应的直接数值模拟，超级计算机研究所研究报告UMSI 95/108，明尼苏达大学，1995年。 21. Huang Y.、Cinnella P.和Lerat A.，空气动力学中湍流 可压缩流动计算的三阶精确中心方案，数字。 流体力学VI，Will Print，1998年，pp. 355-361。 22. Lerat A.和Rezgui A.，可压缩流量的高阶精确紧凑和非紧凑方案，第7届ISCFD会议记录，9月。 1997年，pp. 99-104。 23. Rezgui A.、Cinnella P.和Lerat A.，曲线网格上欧拉计算的三阶有限体积方案”，2000年，将出现。

未来方向127 24。 Jameson A.和Baker T.，复杂配置的Euler 方程解决方案，AIAA第6届计算流体动力学会议，1983年。 25. Lerat A.，Sides J.和Daru V.，用于解决Euler 方程的隐性有限体积法，物理学讲义，170，Springer Verlag，1982年，pp. 343-349。 26. van der Weide E.和Deconinck H.，双曲系统的正矩阵分布方案，应用于Euler 方程，第三届ECCOMAS CFD会议记录，John Wiley \& Sons，9月。 1996年，pp. 747-753。 27. 卡塔拉诺·L。 A.，De Palma P.，Pascazio G.和Napolitano M.，跨声速流动的准确解决方案的矩阵波动分裂方案，物理学讲义，490，Springer Verlag，1997年，pp. 328-333。 28. Struijs R.、Deconinck H.和Roe P. L.，2D Euler 方程的波动分裂方案，VKI LS 1991-01，von Karman研究所，1991年。 29. De Palma P.、Pascazio G.和Napolitano M.，跨声速非黏性流动的混合波动分裂方案，第四届ECCOMAS CFD会议记录，John Wiley \& Sons，9月。 1998年，pp. 579-584。 30. Bonfiglioli A.，非结构化网格上伪可压缩欧拉和Navier-Stokes 方程的多维残差分布方案，物理学讲座笔记，515，Springer Verlag，1998年，pp. 254-259。 31. 卡塔拉诺·L。 A.，De Palma P.，Napolitano M.和Pascazio G.，级联流的细胞顶点自适应欧拉方法，AIAA杂志，33，12月。 1995年，pp. 2299-2304。 32. Dervieux A.、van Leer B.、Periaux J.和Rizzi A. （编辑）”，可压缩欧拉流的数值模拟，数值流体力学笔记，26，Vieweg，1989年。 33. 克鲁普顿P。 I.，Mackenzie J. A.和Morton K. W.，可压缩Navier-Stokes 方程的细胞顶点算法，J。 计算机。 物理，109，1993年，pp. 1-15。 34. 布里斯托·M。 O.，Glowinski R.，Periaux J.和Viviand A. （编辑），可压缩Navier-Stokes流的数值模拟，数值流体力学笔记，18，Vieweg，1987年。

\clearpage
\chapter{8 三维曲线移动网格的高效低耗散高阶方案的扩展}
3D曲线移动网格M的高效低耗散高阶方案的扩展。 Vinokur1和H.C. Yee2摘要 Yee等人[24]提出的效率低耗散高度并行化激波捕捉方案，本质上是四阶或更高，是为有限差分框架中的三维曲线移动网格制定的。 这些方案由高阶紧凑或非紧凑非耗散碱基方案组成，结合自适应非线性特性滤波器，以尽量减少数值耗散远离冲击和剪切区域的使用。 通过将这些方案应用于隐性通量衍生物的熵分裂形式，进一步最小化了数值耗散的数量。 分析了热完美的气体。 将高阶方案扩展到曲线移动网格的主要困难是对坐标变换产生的几何项进行高阶数值评估。 这些术语的数值评估以确保自由流的保护是以协调不变的方式进行的。 在处理轴对称流时，这避免了虚假的数值误差，这些误差是由先前的非不变公式引起的。 1名艾姆斯助理和2名高级职员科学家；美国宇航局艾姆斯研究中心，莫菲特菲尔德，加利福尼亚州94035。 Frontiers of 计算流体动力学 - 2002 编辑：David A. Caughey和Mohamed M. 哈菲兹 ©2002 世界科学

\section{8.1 介绍}
130 VINOKUR \& YEE 8.1 介绍 大多数可用的高阶高分辨率激波捕捉数值方案对于实用的3D复杂模拟来说太CPU密集型。 尽管这些方案具有快速演化流量和短期整合的高分辨率能力，但在航空声学、旋翼、湍流和一般长波传播计算中，这些方案在长时间整合中往往表现出不必要的振幅误差。 这里的高阶是指本质上是距离冲击和剪切区域四阶或更高的空间方案。 数值耗散的微妙平衡稳定性，在长时间集成后不会造成流量物理的过度涂抹，这对这种类型的非定常流动模拟构成了严峻的挑战。 Yee等人[24]及其配套论文Yee等人[25]和Sjogreen \& Yee[14]最近开发的高阶低耗散冲击捕获方案旨在开发克服一些困难的方法。 为了提高效率，Yee等人[24]提出了简单的高度可并行化的空间方案，由基本方案和非线性滤波器组成。 基础方案由窄网格模板高阶紧凑或非紧凑居中非耗散性经典空间差异组成。 滤波器由低阶总变异减小（TVD）的耗散部分、本质上是非振荡（ENO）或加权ENO（WENO）方案和人工压缩方法（ACM）传感器的乘积组成。 ACM 传感器的作用是减少远离冲击和剪切区域的 数值耗散 量。 作为ACM传感器的替代方案，Sjogreen \& Yee [14]利用非正交小波基函数作为多分辨率传感器，以动态确定在每个网格点添加的非线性数值耗散量。 与常用的网格自适应指示器相比，由此产生的小波传感器作为更理想的网格自适应指示器（Gerritsen \& Olsson [4]）随时可用。 与在光谱或类似光谱的非激波捕捉方案和高阶ENO方案之间切换的混合方案相比，高阶非耗散基方案总是被激活的。 如果使用二阶TVD方案作为滤波器，这些方案的最终网格模板在每个空间方向上是五点，如果二阶ENO方案用作四阶基方案的滤波器，则为七点。 研究表明，与标准高阶TVD、正、ENO或WENO方案相比，CPU时间和网格点更少，实现了更高的精度。 这些方案能够准确模拟广泛的流动条件，包括波传播、计算空气声学、燃烧和三维可压缩湍流的直接数值模拟（DNS）的长时间集成。 见Yee等人。 （[24，25，26]，Sandham \& Yee [13]，Sjogreen \& Yee [14，15]，Miiller \& Yee [8]和Polifke等人。 [10]。 表8.1显示了方案的流程图。

曲线移动网格的低耗散方案131高效低耗散高阶方案高阶基方案（所有时间激活）T非耗散（紧凑或非紧凑高阶方案）非黏性和黏性通量非线性特征滤波器（尽量减少Num的使用。 分散。） 传感器“ACM”或“小波的非线性耗散（TVD、ENO或WENO的耗散部分）全强度：冲击和高梯度降低强度：平滑区域（或零）Fux限制器Roe的约。 Riemann 求解器抑制了Spurious Oscil。 （高梯度）提高非线性稳定性满足冲击条件（精确到1-D）静止膨胀-作为膨胀冲击（易于纠正）将方案应用于通量导数的“熵分裂形式”（与非分裂方法比较-j3 = °°）表8.1 F，=标准新III 1 S-F + -J-F W使用相同的基方案在Yee等人[25]中，这些方案应用于非黏性通量导数的熵分裂形式（例如，Fx；见表8.1的底部。） 进行了研究，以确定通量导数的熵分裂形式在多大程度上有助于最小化数值 耗散，或同等地，在改进非线性稳定性中。 他们认为熵分裂是原始守恒定律的一种条件形式。 总体而言，Yee等人的数值结果。 [25]、Sandham \& Yee [13]、Sjogreen \& Yee [14]和Polifke等人[10]表明熵分裂有积极的好处。 分裂可以稳定非耗散性或低耗散性空间离散化产生的虚假噪声，这是非线性不稳定的主要原因。 他们的研究还表明，即使使用不满足Strand [16]条件的数字边界条件，仅熵分割也可以改善非线性稳定性。 熵分割的这种稳定性属性不仅对有关方案类别有价值，而且也可以应用于实践中常用的其他方案。 这强调了一个事实，即人们应该始终尝试将数值方案应用于控制方程的更有条件的形式。 在Sandham \& Yee [13]中，Carpenter等人[1]开发的Strand边界差算子变体的熵拆分与3D无冲击可压缩湍流通道流的DNS的高阶非耗散中央基础方案结合使用。

132 VINOKUR \& YEE 使用粗到中等网格尺寸，不使用过滤器，获得非常准确的全面开发的湍流统计数据。 他们的结果与不可压缩的Navier-Stokes模拟的最佳光谱方法进行了很好的比较。 现代高分辨率数值耗散是改善短期或中度时间积分（非稳态）非线性不稳定性的主要因素。 大多数情况下，新增数值耗散对于更长的时间积分是必要的，而不必诉诸于更精细的网格和极小的时间步进，而牺牲了流物理的过度涂抹。 使用通量导数的熵分裂形式已被证明能够最大限度地减少数值耗散的使用。 在Yee等人[25]中，研究了熵分裂概念与其他状态方程和进化方程集的物理方程的扩展性。 他们的研究表明，熵分裂可以正式扩充套件到热完美气体，内能是温度的任意函式。 虽然将熵分裂扩展到非平衡流动在实践中是不可行的，对于求解不可压缩Navier-Stokes方程的平衡实气体和人工可压缩性方法以及磁流体力学（MHD）流动来说是不可能的，但有关的高阶数值方案适用于这些类型的流动。 请注意，由于麦克斯韦方程是一个容易对称化的双曲方程的线性系统，斯特兰德的数值边界算子仍然有效，有关数值方案也适用。 麦克斯韦方程不需要熵分割。 对于非平衡流动，如果以松散耦合的方式分别求解物种和流动方程，那么流动方程有效地满足了局部热完美气体定律，并且适用“局部”形式的熵分裂。 为了在实践中应用这个方案，它必须以任意的3D 曲线坐标来制定。 有限体积公式通常优于有限差分公式，特别是对于三阶或更低的方案。 对于四阶或更高方案，有限体积公式非常复杂。 因此，为了提高效率，我们只考虑有限差分公式。 因此，一个重要的问题是坐标变换产生的几何项的处理。 这些是变换通量和变换保守变量定义中的体积（或雅各布）中出现的表面积向量（或度量）的分量。 从分析上讲，这些几何量满足某些守恒定律。 它们是表面积守恒律（或公制恒等式），对于时变网格，额外的体积守恒律（有时称为几何守恒律）。 最好在数量上满足这些规律。 这将导致均匀流动的区域完全保留（在四舍五入误差内）。 在非均匀流动的区域，希望通过消除网格造成的误差来提高准确性，特别是如果它高度失真。

曲线移动网格的低消散方案133 在Vinokur [20]中，展示了如何利用有限体积概念在数值上满足这两个定律。 对于表面积向量分量，发现离散化与Pulliam \& Steger [11]在二阶有限差分框架中提出的平均程序相同。 不幸的是，这种有限体积概念的使用仅适用于二阶精度中心差分。 它们不能扩展到高阶紧凑或非紧凑差异。 Thomas \& Lombard [18]提出了一种将表面积向量成分离散化的替代方法，首先以等效的“保守”形式重写它们。 对于二阶精确的中心差分，然后对表面积守恒律进行数值满足。 Gaitonde \& Visbal [3]通过对两个不同的曲线网格的数值实验发现，应用于Thomas和Lombard形式的高阶紧和非紧微分也满足了守恒律的数值。 所用表达式的一个不良特征是它不是坐标不变的。 作者于2000年6月和Thomas \& Neier[19]（P.D.）独立提出了当前坐标不变形式。 托马斯，私人通讯，2000年11月）。 8.1.1 目标 在本文中，我们制定了Yee等人[24]方案，在3-D 曲线坐标中具有熵分裂（Yee等人[25]），用于有限差分方法的热完美气体。 通过将托马斯和伦巴第公式的坐标不变版本离散化，表面积守恒律在数值上得到满足。 这样做的重要性在于，它使我们能够将该方法扩展到轴对称流动，并适当处理生成的源项。 这消除了由于原始的Thomas和Lombard表示式而导致的虚假数字错误。 我们还提供了一个严格的分析证明，即数值混合偏导数对任何网格进行交换，使用任何高阶紧致或非紧致微分，与任意数值边界条件。 这对于确保表面积守恒律在数值上得到满足是必要的。 8.1.2 大纲 在第8.2节中，我们回顾了适用于热完美气体的Euler 方程的熵分裂方法。 这些方程是为任意的3-D 曲线坐标制定的，并专门用于轴对称流动。 本节还包括对Navier-Stokes 方程的简要讨论。 在第8.3节中，我们描述了处理空间项的数值方法。 特别注意几何项的处理。 关于Roe近似Riemann 求解器的部分，这是数值方法的一部分

\section{8.2 方程的制定}
134 VINOKUR和YEE过滤器，包括讨论定义广义三维坐标的通量Jacobian 矩阵的特征向量所涉及的一些问题，以及热完美气体的罗伊平均状态的表达式。 附录B列出了非平衡流动的相应形式。 附录A提供了数值混合偏导数与广义坐标中高阶基数交换的证明。 8.2 方程的制定 8.2.1 守恒定律的规范分割 标量守恒定律系统可以写成Qt + V-F = 0，（2.1.1），其中Q和F（Q）是代数向量，但F的分量是物理向量。 字母下标表示部分分化。 为了获得非线性求解双曲守恒定律系统（2.1.1）的初始边界值问题（IBVP）的非线性稳定方法，我们变换控制方程，使生成的PDE非线性稳定，包括物理边界条件的影响（Olsson \& Oliger [9]）。 我们引入了新的向量W(Q)，使Fw对称，Qw对称且正定。 在这里，矩阵Fw也有物理向量的分量。 此外，W的选择是F和Q都是W的次数/3的齐次函式，即有一个常数0，对于所有a Q(aW) = (TpQ(W), (2.1.2a) F(aW) = <rfiF(W)。 （2.1.2b）通过取方程的偏导数。 （2.1.2）关于a，然后设置a = 1，我们得到欧拉定理，即QWW = 0Q，（2.1.3a）FWW = \{3F。 （2.1.3b）为了获得IBVP稳定性的能量估计，我们引入了规范拆分，并将V • F写成V • F = -£-V • F + -\textasciicircum{}-FW • VW, /3\#-l, (2.1.4) p + i P + i，其中VW是一个向量，其分量是W分量的梯度。 （2.1.4）中第二项中的“•”表示，当采取

曲线移动网格的低消散方案 135 一行Fy/和向量WW的内积，乘分量时应用向量点积。 请注意，V•F已分为保守部分和非保守部分。 如果我们取W和V • F的内积，并使用F的同质性属性和Fw的对称性，我们得到WTV-F = -\textasciicircum{}--[WT(3\textbackslash\{\}7-F + WTFwVW] (2.1.5a) = -\textasciicircum{}-:[WTV • \{FWW) + VWT • FWW] P + 1 (2.1.5b) = \textasciicircum{}V • (WTFWW)。 （2.1.5c）将规范拆分应用于Qt，我们以类似的方式获得WTQt = -J-l\{WTQwW)t。 （2.1.6））取W的内积与（2.1.1），替换（2.1.5c）和（2.1.6），在具有边界表面元素ndS的固定域上积分，其中n是外法线，并应用发散定理，导致j f WTQwWdV = -i WTn • FwWdS。 （2.1.7）其中n • Fw是矩阵，其分量是F\textasciicircum{}的外向法向量分量。 Prom（2.1.7）可以严格地根据与域边界的传入特征变量相对应的绝对特征值来严格建立能量规范增长率的界限（Olsson \& Oliger [9]，Gerritsen \& Olsson [4]）。 在这样做时，人们必须利用Euler 方程的Qw- 8.2.2熵分裂的正定属性。在本节中，我们考虑了三维可压缩Euler 方程的守恒定律形式。 为了最大程度地概括，分析是针对任意的三维、时间变化的网格提出的。 第2.1节中讨论的规范分割基于满足熵条件的凸熵函数。 因此，分裂被称为熵分裂。 虽然这种分裂最初是针对完美气体（Gerritsen \& Olsson [4]）得出的，但它可以扩展到仅在热上完美的气体的流动，内能是温度的任意函数（Yee等人。 [25]）。 这个定律适用于由单一化学物种组成的稀释气体，也是氧气开始解离（约2000°K）时温度以下的空气的一个很好的近似值。

136 VINOKUR \& YEE 让i、j和k是沿着笛卡尔坐标系轴的单位向量。 然后，位置向量r和流速u表示为r = xi + yj + zk (2.2.1)和u = ui + vj + wk。 （2.2.2）时变网格中的网格点速度由v表示。 让我们把相对于运动网格的速度。 然后Euler 方程的向量Q和F的形式是Q = p ' pu pv pw e., F = pu' puu' + pi pvu' + p] pwu' + pk.(e + p)u' +p\textbackslash\{\} (2.2.3) 其中p是密度，p是密度，e - p(e + q2/2)是单位体积的总能量，e是比内能，q2 = u2 + v2 + w2是速度的平方。 请注意，F现在是相对于运动网格定义的。 热完美气体的状态方程为P = pT，（2.2.4），其中T = RT是归一化温度，R是气体常数。 T与e具有相同的维度。 从热力学的第一和第二定律中，我们只能证明e = e(T)。 所有实气都满足条件e > 0和e > 0，其中e = de/dT。 温度T(Q)是通过求解方程\textasciitilde{} e \_ 1 [(pu)2 + (pv)2 + (pw)2\} 6[1) p 2 p2 (2.2.5) 方程(2.2.5)有一个唯一的解，因为e > 0。 根据热力学定律，我们可以通过Pf将无维熵S = S/R与p和f联系起来，其中/(f) = exp(-|4df)。 (2.2.6) (2.2.7) 积分中的任意常数可以被吸收到S的定义中。 遵循Harten [6]，我们从W = 8Q'（2.2.8）中得到向量W（Q）

曲线移动网格137的低散散方案，其中凸函数T(Q)由r = pV(5)给出。 (2.2.9) W的分量有时被称为“熵变量”，而r被称为“熵函数”。 我们在Yee等人[25]中表明，为了满足同质性和正定性条件，ip(S)由\$ = pe-VP给出，（2.2.10），其中（3是一个常数。 然后，这给出了--S//3 \_ <Pf) 1//3 (2.2.11)，我们使用(2.2.6)来获取第二个表达式。 我们从（2.1.4）中注意到，W的定义允许任意乘法常数。 相应地，在if的定义中有一个任意的乘法常数）。 我们选择了这个常量来获得（2.2.11）的最简单形式。 转换变量W现在可以写成W = -[e-2?-p(l+/3) P -pu -pv -pw p\textbackslash\{\} (2.2.12) where? = pe是单位体积的内能。 在（2.2.4）的帮助下，我们可以结合W的分量，以表明T和e（T）是W的0度的齐次函数。 然后从（2.2.3）、（2.2.4）、（2.2.11）和（2.2.12）可以得出，向量Q和F都是W的等级/3的齐次函数。 虽然我们无法获得Q(W)的显式公式，但我们可以推导出对称矩阵Qw的显式表示式，作为Q的函数。 上三角形部分可以写成Q w V> ap apu apv apw apu2 - p apuv apuw apv2 - p apvw apw2 - p where and ae + bp u[ae + (b - l)p] v[ae + (b - l)p] w[ae + (6 - l)p] (2.2.13) a\{T,P) 1 6+1+/3 - 1 b\{T,p) = \{l + 0)a + p=-e + l + /3 (2.2.14) (2.2.15)

138 VINOKUR \& YEE 我们在Yee等人[25]中证明，Qw上的正定条件要求tp < 0，一个条件已经满足（2.2.11），如果/3 > 0或/3 < -(1 + e)，则\textbackslash\{\} > ITI-（2-2J6）条件（2.2.16）满足。 由于I > 0，k.的nicixiniuni值出现a,t Trnax-因此，对于/3 < 0，f3 < -[1 + e(fmax)]。 （2.2.17）通过用e（oo）替换e（Tmax）获得与流动问题无关的充分条件。 完美气体的专业化：对于比热比为7的完美气体，状态的热量方程为f 1 e = -. -和e = -。 (2.2.18) (7-1) 7-1 Pf = (pp\textasciitilde{}'f)\textasciicircum{}, (2.2.19) \textasciicircum{} = -(pp-i)V\textasciicircum{}m, (2.2.20) 它遵循它和0(7.« = 7T\textasciicircum{}TF P-\textasciicircum{}2) /3的正定条件然后变成/3 > 0或/？ < -\textasciicircum{}。 请参阅Yee等人[25]和Sandham和Yee[13]，了解各种流量的有益范围/3的研究。 8.2.3 广义曲线坐标中的公式化 到目前为止提出的方程原则上可以通过任何类型的网格的任何数值方法实现。 由于我们的兴趣在于高效、高阶准确的解决方案，我们将把自己限制在结构化网格上的有限差分公式。 因此，在本节中，我们考虑广义曲线坐标中方程的公式化。 我们首先提出了三维流动的方程。 然后，它们将专门针对轴对称流动的重要情况。 本节最后简要讨论了Navier-Stokes 方程。

曲线移动网格的低消散方案139 8.2.3.1三维流从曲线坐标到物理空间的任意、与时间相关的变换写成r = r(£,r?,C,T) t = T。 （2.3.1a）（2.3.1b）对于计算单元d£、dn和d£，£、n、C方向的归一化表面积向量由S€=r、xrc、S'？给出。 = rcxrc，Sc = re x r，。 （2.3.2）归一化的单元体积由V = rrr,xrc（2.3.3）给出，网格点速度由v = rT给出。 （2.3.4）将变换（2.3.1）应用于运动网格版本（2.1.1），我们得到QT + \% + Fn + Gc = 0，（2.3.5a），其中Q = VQ，E = S*-F，F = S"-F，G = Sc • F。 （2.3.5b）在下文中，我们将使用数字下标来表示笛卡尔分量。 因此，S« = S\textasciicircum{}i + S|j + S|k，（2.3.6）与Sv、S1*和v的定义相似。 让tf = S« • v = S\{v！ + S|«2 + Sf w3, (2-3.7) v\textasciicircum{}和v\textasciicircum{}的定义相似。 对于Euler 方程，变换通量E由puU + Sfp E= pvU + S\textbackslash\{\}p, (2.3.8) pwU + Sf p给出。 (E + p)U + v\textasciicircum{}p. 其中U = S€ • u' = Sfu + S\%v + Slw - ««. （2.3.9）

140 VINOKUR \& YEE 变换的通量导数E\textasciicircum{}现在被分割为E, P '\textasciicircum{}JTiEt + JT-iE\textasciicircum{} (2.3.10)，其中Elfc = S\textasciicircum{} • Fw由E\textbackslash\{\} w和apU apuU给出 apuU - S\textbackslash\{\}p apvU - S\textasciicircum{}P apwU - S\textbackslash\{\}p ei '3f <315 (apu2 - p)U - 2uSfp apuvU - p(5|« + Sf v) e24 e25 (apv2 -p)U - 2vS\%p e34 e35 e45 essss (2.3.11 ei5 = [oe + (b - l)p]U - v\textasciicircum{}24 = apuwU - p(S\textasciicircum{}u + Sf u>), e25 = \{[ae + (b L\textasciicircum{}+p\{2(b-l)--q2\} + \textasciicircum{}\{b(l + (3)-2\}\}U-2vS\textasciicircum{}(e + p)。 （2.3.12h）P P P P F和Fv的类似表达式从（2.3.8）到（2.3.12）通过将U替换为V，将£替换为rj。 同样，G和G\textasciicircum{}的表达式是通过将U替换为W，将£替换为C来获得的。 这里是V和t？ 是（2.3.9），S\textasciicircum{}分别被Sn和S\textasciicircum{}取代，W与（2.2.12）中的熵分割向量W没有关系。 通常，我们需要计算Qw f°r QT = TJTIQT + -gr\textasciicircum{}QwWT的分割形式。 然而，我们只考虑应用时间离散化的半离散方法。 除了使用非黏性通量导数E\textasciicircum{}、Fn和GQ的拆分形式外，我们不必使用QT的拆分形式来实现。 因此，半离散熵分割方法的最终形式仍然可以用保守和原始来表达

曲线移动网格141个变量的低消散方案，使现有计算机代码实现简单高效。 从定义（2.3.2）中，我们可以推导出表面积守恒律（S«）€ +（S”）„ +（S«）c = 0，（2.3.13），它对每个笛卡尔分量都有效。 对于时间变化的网格，通过将（2.3.13）与（2.3.5）相结合，并假设一个均匀的流量，我们推导出体积守恒律 VT = (\textasciicircum{})c + (a"), + (\textasciicircum{})c- (2.3.14) 请注意，我们没有在其分量形式中写出关系（2.3.2）。 我们将在第3.3.1节中表明，为了满足（2.3.13）数值，必须修改这些关系。 最后，我们将我们的符号与Steger[17]介绍的更熟悉的符号联系起来。 这些是V = J\textasciitilde{}\textbackslash\{\} U = J\textasciitilde{}lU，Sf = J-\textasciicircum{}x，\& = -J-1£u（2.3.15），其他量具有类似关系，其中Steger将J定义为变换的雅各布，U是他在£方向的逆变速度。 8.2.3.2 轴对称流动 为了获得轴对称流动的方程，我们首先引入圆柱坐标x，r，（，其中a>轴是极轴，£是极角。 在x-r平面上引入曲线坐标 £，r\textbackslash\{\}，我们有变换方程x = x(Z,r),r)和r = T-(£,T/,T), (2.3.16)和y - r cosC和z = rsm(。 （2.3.17）然后可以从（2.3.2）中获得表面积向量。 特别是，S\textasciicircum{}的成分变成S< = 0，s£ = -S<sinC，S£ = SccosC，（2.3.18），其中5C = x\textasciicircum{} - r\textasciicircum{}Xr，。 （2.3.19）细胞体积由（2.3.3）给出，如V = rSc。 （2.3.20）

142 VINOKUR \& YEE 对于轴对称流，C分量0f u和v，以及物理量的£导数，都等于0。 因此，v\textasciicircum{} = W = 0。 然后，变换的通量G变为<S = SS[0 0 -sinC cosC 0]T，而其£导数为Gc = -Scp[0 0 cosC sinC 0]r。 （2.3.21）（2.3.22）轴对称方程是通过设置（=0）获得的。 然后，u的£-分量变成w，设置为0。 从（2.3.17）可以得出，z = 0和r可以在（2.3.16）、（2.3.19）和（2.3.20）中被y取代。 Gc in (2.3.4)然后成为源项<gc = -S<p[0 0 1 0 Of (2.3.23) 最后，所有向量的第四个分量，以及所有矩阵的第四行和列都可以消除。 请注意，表面积成分的所有关系必须在数值实施前进行修改，如第3.3.2节所示。 我们还应该提到，某些关系的某些极限形式必须用于对称轴上的点。 8.2.3.3 2.3.3. Navier-Stokes 方程 到目前为止，我们只处理了Euler 方程，因为规范分割仅适用于非黏性术语。 为了完整起见，我们对Navier-Stokes 方程提出了一些评论，其中额外的黏性术语以标准方式处理，这些术语有据可查。 守恒定律系统现在写成Qt + V • (F - F„) = 0。 （2.3.24）相对于运动网格，F由（2.2.3）给出，黏性通量向量Fv由F = 1 • T k-T -q + T -u 牛顿流体的假设产生关系T = A(V-U)I + /J[VU+(VU)T]和q = -kVT。 （2.3.25）（2.3.26）（2.3.27）

\section{8.3 数值方法}
曲线移动网格的低消散方案143在这里，fi和A是第一和第二黏度系数，k是导热系数。 I是身份张量。 为了在广义坐标中表达（2.3.25），我们使用关系v = i[s« v-| + < + sC|]- <2328）曲线坐标中方程的最终形式和相应的高阶基方案将不在此给出。 它们可以在标准参考资料中找到。 在这里，我们只会提到轴对称流动的源项的修改。 必须修改（2.3.23）中的系数，以解释作用在方位方向上的法向黏性应力。 因此，压力p in（2.3.23）必须被p-AV-u--取代。 （2.3.29）8.3 数值方法 Yee等人[24]提出的空间离散化由两部分组成，即基本方案和滤波器。 当不使用过滤器时，该方案仅由基本方案组成，第8.3.1节对此进行了讨论。 基础方案涉及紧凑或非紧凑的高阶中心差分近似，用于评估通量导数E(\_,FV,G\textasciicircum{}、它们所包含的表面积分量以及导数W\textasciicircum{}、Wv和W\$。 第3.2节讨论了适用于不黏性术语的过滤器的形式。 第3.3节和第3.4节描述了在数值上满足基本方案和滤波器表面积守恒律所需的几何项的特殊处理。 请注意，如果滤波器的精度空间顺序与基本方案相同，则其表面积分量离散化与基本方案相同。 如果滤波器的精度低于基本方案，为了一致性，表面积组件的相应离散化顺序只需要与滤波器的顺序相匹配。 8.3.1 空间基数方案让<Pi,j,k表示函数<p在（iA£，jAr]，kAQ的离散近似值，其中A£，A77和AC是£，rj，(\textasciicircum{}-方向和（i，j，k）是相应的空间指数中的网格间距。 ip可以是Q和/或（£，ry，<\textasciicircum{}，T）的函数。 （P\$ at（iA£，jArj，kAQ（ipv和ip\textasciicircum{}的类似表达式）的可能高阶中心差近似近似值可以有以下两种类型：

144 VINOKUR \& YEE 中心差分：（四阶和六阶）<Pi Y\textasciicircum{}T ( V«+2,i,fc - \&iPi+l,j,k + 8<Pi\_l,j,fc - <Pi-2,j,k J, (3.1.1) (3.1.2) 紧凑中心差分：（四阶和六阶）其中四阶近似n\textasciitilde{}lTc\textbackslash\{\}AiBw), (3.1.3a) (\textasciicircum{}y,j,fc = g ( V»+i,j,fe + 4\textasciicircum{}i,j,fc + <Pi-ij,k ), (3.1.3b) iB\&)\%,]* = j ( V»+i>j',fe \textasciitilde{} V*-i,i,fc ) > (3.1.3c) 和近似 (A;</>)ij,fc = g f 对E\textasciicircum{}等项应用非耗散性高阶基方案。 对于Navier-Stokes 方程，近似值首先用于评估黏性应力张量T和热通量向量q，然后再次用于获得黏性通量F„的导数。 所有这些项中出现的表面积向量分量也使用这些差分近似进行评估，但等式。 2.3.2必须首先修改，以便完全满足（2.3.13）。 这将在第3.3节中讨论。 对于可压缩湍流计算的直接数值模拟（DNS），Sandham和Yee使用拉普拉斯式的黏性和热传导项，以确保与中央方案相关的任何奇偶解耦合倾向都可以被流体黏度所抵消。 他们还根据类似的论点使用连续性方程的特殊公式。 处理Navier-Stokes 方程中黏性应力张量和热通量的另一种方法是在半网格点进行评估。 使用适当的高阶中心差分公式，人们将获得模板

曲线移动网格145的低消散方案，用于不比欧拉通量宽的黏性通量。 这将需要高阶插值来评估半网格点的表面积向量分量和运输系数。 移位的表面积向量分量不会满足（2.3.13），但这不会影响均匀流动的保持，因为黏性项的贡献为零。 目前正在研究这种替代方法。 8.3.2 过滤器 在本节中，我们首先回顾了将特征过滤器应用于多阶段Runge-Kutta类型和线性多步法（LMM）类型的时间离散化（Yee等人。 [24]）。 显式LMM的例子是前尤拉和亚当斯-巴什福斯。 隐式LMM的例子有后退欧拉、梯形规则和三点后退微分。 Dahlquist的LMM的单腿公式[2]也适用。 然后我们讨论特征过滤器的形式。 8.3.2.1 应用过滤步骤的程序如果需要Runge-Kutta方法等多阶段时间离散化，则在Runge-Kutta方法的每个阶段都应用上一节中讨论的空间微分基础方案。 如果存在黏性术语，我们对黏性术语使用与不黏性术语相同的顺序和类型的基础方案。 应用特征滤波器有两种方法。 方法1是在Runge-Kutta步骤的每个阶段应用过滤器。 方法2是在整个Runge-Kutta步骤结束时应用过滤器。 对于黏性和强冲击相互作用，方法1可能更稳定。 如果想要属于LMM类的时间离散化，那么滤波器可以作为数值耗散向量与基本方案一起应用。 在这种情况下，过滤器在Qn处评估显式LMM。 对于隐式LMM，涉及在n+1时间级别评估的额外类似过滤器。 或者，方法2也可以应用于LMM。 在这种情况下，我们在完成隐式时间步骤后应用过滤器。 通过使用Dahlquist的LMM的单腿公式，可以最大限度地减少通量评估。 唯一的非耗散（时间）二阶、双时间级单腿方法是中点隐式方法。 请注意，中点隐式公式的非迭代线性化形式简化为正则非迭代线性化梯形公式。 对于使用隐式LMM的时间走向稳定状态，某些流动物理学只需要一个明确的耗散项。 此外，隐式运算符可能与显式运算符不同。 请参阅Yee [22]和Yee等人[23]，了解一些高效的保守线性化隐式形式。

146 VINOKUR \& YEE 8.3.2.2 非线性特性滤波器 滤波器运算符与高阶基本方案相结合，有许多可能的候选者。 在这里，我们更喜欢使用网格模板的宽度与基本方案相似的过滤器操作器，以提高数值边界处理的效率和便利性。 当然，高于三阶滤波器运算符是适用的，但它们的CPU更密集，需要在边界点附近进行特殊处理，以保证稳定性和准确性。 滤波器操作器通常由传感器和非线性耗散的乘积组成。 考虑了两种可能的传感器，ACM传感器（Harten [5]）和小波传感器（Sjogreen \& Yee [14]）。 路线图见表8.1。 在这里，我们简要回顾了ACM传感器，有兴趣的读者可以参考Sjogreen \& Yee [14]来了解小波传感器。 Yee等人[24]将非线性耗散项与应用于每个特征波的Harten ACM传感器结合使用，作为滤波器向量。 从本质上讲，非线性耗散项充当二阶或三阶类似ACM的算子，与Harten的一阶ACM（Harten [5]）不同。 该传感器用于向高阶非耗散方案添加的非线性耗散量发出信号，一次一个波。 因此，目前的方法在精神上也与原始的Harten [6]二阶TVD方案不同，后者使用ACM来锐化接触不连续性。 让\textasciicircum{}方向的滤波向量为E*+hjk = \textbackslash\{\}Ri+i\&+,(SS)i+k<jtk的形式。 （3.2.1）E* i.，是曲线坐标中标准数值通量非线性耗散部分的修改形式。 (S\textasciicircum{})i+i -fe是（i + 1/2，j，k）平均状态下的表面积向量S\textasciicircum{}的大小。 数量Ri+i是第3.4.1节中给出的右特征向量矩阵，使用Roe的近似求解器（Roe的平均状态），并抑制了（k，I）指数。 我们以同样的方式定义F*., i.和G* •,,i。 \$* j的元素，用(<f>l,i)*表示，并抑制了(k,l)指数，是Oi+i)* = K0\{+i*i+i。 （3.2-2）（fil. i in（3.2.2））是使用TVD、带坡度限制器的MUSCL、ENO或WENO方案产生的高分辨率*+ 2方案耗散部分的元素。 下文，我们将（3.2.2）称为ACM过滤器。 公式（fi'i涉及第3.4.1和3.4.2节中定义的特征值、特征向量和Roe的平均状态1+2，可以在文献中找到（Yee [21，22]和Yee等人。 [23]）。 有关详细信息和其他可能过滤器的讨论，请参阅Yee等人[24，25]。

曲线移动网格的低消散方案147功能K91，！ 是一种控制高分辨率冲击捕捉方案耗散部分固有的多余数值耗散的机制。 换句话说，\$*，i的元素与TVD、ENO或WENO方案的非线性耗散部分相同，乘以K91。i。 参数n取决于问题。 对于平滑的流动，K可以是1+2非常小。 它用于最大限度地减少产生非线性不稳定的伪高频振荡。 由于流动特性的变化，不同的物理问题可能需要不同的K值。 K值可能从一个特征波到另一个特征波，从流场的一个区域到具有不同流结构的另一个区域。 Yee等人[24，25]中数值实验的K范围为0.0 < K < 0.7。 函数6\textbackslash\{\}k是Harten ACM传感器。 对于一般的2m + 1点基方案，Harten推荐0< =max\$\_m+1,..,\textasciicircum{}+m)> (3-2-3) al la'., 11 - \textbackslash\{\}a\textbackslash\{\} 1 «+*' i i- \textbackslash\{\}a1., 11 + \textbackslash\{\}al. 1 1+51 ' »- (3.2.4) a', i是R的元素。 \textbackslash\{\} \{Qi+i j k - Qijk）在第3.4.1节中定义。 它们用于（3.2.4），也用于（3.2.2）中的（j）1. k。 (3.2.4)中的参数p是指数> 1，不是“压力p”。 使用MUSCL公式的相应数值耗散的\textasciicircum{}i+ii+ii和\textasciicircum{}t+i现在是Q的左右状态的函数。 人们可以改变p，而不是为特定的流动问题改变K。 对于较大的p，添加的数值耗散更少。 请注意，通过改变p > 1 in（3.2.4），人们基本上可以提高耗散项的准确性顺序。 耗散的顺序取决于p的值。 人们可以从激波位置附近的p = 1切换到光滑区域的p > 1。 对于Yee等人的所有数值示例[24，25]，他们使用p = 1和0j+i=max\$,£！ +1）。 （3.2.5）冲击-湍流相互作用问题似乎有利于这种形式的6l。 1. 1+2 为了避免调整任意参数K的需要，Sjogreen \& Yee [14]建议用多分辨率非正交小波传感器l-Y 2 u)1 i替换n9l k。 他们的研究表明，小波传感器确实以与Yee等人所示的相同数值示例的准确性相当地消除了参数依赖性K和p。 [24，25]。 与ACM传感器不同，小波传感器可以检测包括湍流在内的大多数不同的流动特征，从而自动选择数值耗散的适当分布和良好的网格自适应指示器。

148 VINOKUR \& YEE 为了避免实际计算机代码中的一些条件语句，并促进矢量化，修改了过滤器内的几个具有除以零潜力的函数。 特别是，p - 1的传感器（3.2.4）被1 I i+l/2l *i-l/2\textbackslash\{\} \textbackslash\{\}a l i+l/2 + «:• -1/2 + e (3.2.6) 在所有计算中，它们都取e - 10\textasciitilde{}7。 （实际上，e应该与a'+1/2-具有相同的尺寸）这里强调的是，ACM和小波传感器都无法提高冲击和剪切位置的准确性，而不是非线性耗散的固有激波捕捉能力。 冲击和剪切的准确性取决于所选的非线性耗散。 传感器的作用是允许在冲击和剪切区域中充分使用数值耗散，并限制紧靠在远离冲击和剪切位置以及流场其余部分的区域的数值耗散量。 因此，有了合适的传感器，人们不必使用CPU密集型高阶高分辨率激波捕捉 数值耗散，因为这种类型的耗散通常只能提供离不连续性稍有更准确的解决方案，同时表现出与高分辨率二阶或三阶数值耗散相同的冲击和剪切分辨率。 最后，我们注意到，ACM滤波器（3.2.2）或小波滤波器可能不足以（a）向稳态进行时间，以及（b）由于电网分辨率不足和非线性不稳定而导致的虚假高频振荡。 请参阅Yee等人[25]进行讨论。 完整离散化的形式：在这里，我们使用Rung-Kutta方法与\textasciicircum{}在完成完整Runge-Kutta时间步骤时应用的过滤器来说明完整离散化的形式。 让Qn+1成为使用非耗散基方案完成一个Runge-Kutta时间步骤后的解决方案。 如果采用熵拆分，则基本方案适用于黏性通量导数的拆分形式。 下一个时间级Qn+1的解决方案是At Qn+i = Qn+i + - rp* 771* i+\textbackslash\{\},j,k i-i,j,fc I At '"AJJ "AC F*。 \_Ll.-FT。 1。 Gh,k+\textbackslash\{\} Gh,k-\textbackslash\{\} (3.2.7) 这里，E*。 1.,, F*., r, 和 G*.,, 1在Qn+1进行评估。

曲线移动网格的低消散方案149 8.3.3 离散化基方案的几何项对于基方案，所有与变换变量相关的第一偏导数，如（2.3.10）中的偏导数，由第3.1节的高阶紧和非紧中心差公式近似。 但变换通量中的表面积向量分量和变换保守向量中的细胞体积也是以笛卡尔坐标的偏导数来定义的。 表面积向量的数值近似值最好完全满足表面积守恒律（2.3.13）（在四舍五入误差内）。 对于时变网格，我们还希望以数值满足体积守恒律（2.3.14）。 这将消除均匀流动区域的数值误差，并减少电网诱导的误差，即使是非均匀流动。 此外，在推导拆分公式（2.3.10）时，（2.3.13）用于获得第二项的表达式。 Vinokur [20]表明，通过诉诸有限体积概念，可以在数值上满足两个几何守恒定律到二阶精度。 不幸的是，这种方法不能扩展到更高的准确度。 通过应用第3.1节到（2.3.2）的高阶差分公式来计算表面积向量，在评估（2.3.13）时会导致截断误差。 Gaitonde \& Visbal [3]通过进行一些数值实验表明，（2.3.13）可以通过以Thomas \& Lombard[18]首次提出的等效“保守”形式重写（2.3.2）的组件版本来数值上满足。 该形式不是坐标不变的。 在第3.3.1节中，我们提出了三维流动的托马斯和伦巴第关系的坐标不变形式。 然后，我们在第3.3.2节中推导出轴对称流的相应正确关系。 这无法使用原始的托马斯和伦巴第公式来实现。 我们还将在附录A中从理论上证明，将任何任意的紧致或非紧致差分公式应用于新形式的表面积向量定义（2.3.2）将完全满足表面积守恒律（2.3.13）。 8.3.3.1 三维流动 通过检查托马斯和伦巴第公式，人们可以很容易地获得S\textasciicircum{}的坐标不变形式，如S4 - \textasciicircum{}[(r„ x r)c - (rc x r)„]，（3.3.1）通过分别用C和f替换77和£获得的Sn的相应表达式，以及通过分别用£和77替换r\textbackslash\{\}和£获得的S"*。 从分析上讲，（3.3.1）给出了与（2.3.2）提供的相同答案，指出由于混合偏导数交换，r,\textasciicircum{} xr - r\textasciicircum{},, xr = 0。

150 VINOKUR \& YEE 通勤属性还保证了S\textasciicircum{}的答案与r的原点选择无关，并且（2.3.13）是相同的。 为了证明（2.3.13）在数值上完全满足，我们必须证明混合偏导数在数值上交换。 我们在附录A中做了这个。 在分量形式中，（3.3.1）写成si = 2\textasciicircum{}yr>z \textasciitilde{} Zriy\textasciicircum{}i \textasciitilde{}\textasciitilde{} (y<z \textasciitilde{} zCV)r,], (3.3.2a) 52 = 2\textasciicircum{}z,nx \textasciitilde{} x-nz)i \textasciitilde{} (ZCX " \textasciicircum{}C*)»?] > (3.3.2b) 53 = 2\textasciicircum{}xiy \textasciitilde{} Vrix)i \textasciitilde{} (xSy \textasciitilde{} yCx)vl- (3.3.2c) S71和S\textasciicircum{}分量的表达式以类似的方式获得。 托马斯和伦巴第的形式在每个括号中使用了两次第一个术语。 虽然这令人满意（2.3.13），但它不会给出轴对称流动的正确公式，正如第3.3.2节将讨论的那样。 从（2.3.7）中，我们看到变换的网格速度涉及表面积分量和网格速度分量的乘积，不能用“保守”形式表示。 因此，当表面积分量由（3.3.2）给出时，我们无法用数值表达V，从而完全满足（2.3.14）。 因此，我们从不变表达式（2.3.3）中将V计算为V = xv yn zv。 （3.3.3）H XJJ zz ZJI zc 8.3.3.2 轴对称流 轴对称表面积分量是通过将（2.3.16）和（2.3.17）代入（3.3.2），然后设置< = 0获得的。 从（2.3.17）来看，z = 0和r可以被y取代。 结果表达式是si = 2(v2)v S2 = -j\textbackslash\{\}.yvx \textasciitilde{} xvV - ixv)ri\} S3 = °> (3.3.4a) 5? = -\textbackslash\{\}\{y\textbackslash\{\} S？ = -\textbackslash\{\}[yiX - xiV - (xy)t\textbackslash\{\} S\% = 0，其中Sf = 0 S\$ = 0 5| = Sc, s<° = 2\textasciicircum{}xiy \textasciitilde{} ytx)v \textasciitilde{} (xiy \textasciitilde{} yvx)t\textbackslash\{\}- 我们也有关系（3.3 (3.3 (3- (3...4b).4c) 3.5) 3.6)

曲线移动网格的低消散方案151将（3.3.4）、（3.3.5）和（3.3.6）替换为（2.3.13），我们看到表面积守恒律是满足的。 细胞体积由V-ySc给出。 （3.3.7）如果我们使用（3.3.2）的托马斯和伦巴第形式，结果方程（3.3.4）到（3.3.6）将有不同的表达式。 虽然它们仍然会满足（2.3.13），但它们会导致违反轴向对称性要求的数值误差。 这是其形式中缺乏坐标不变性的直接结果。 方程（3.3.4c）仍然成立，（3.3.5）被每个括号中第一个项的两倍取代；即SC = (a*j/)„ - (x\textasciicircum{}. （3.3.8）他们的£衍生品的表达式在物理上是不正确的。 前两个现在将是（Si）c = (Mt)r> - (MV)S (52C)c = (m)v - (*Vvh- (3-3-9) 分析，(S£)c - 0和 (S\textasciicircum{}c = -£"*，对应于常数£曲面是平行于z轴的平面。 但从数值上讲，这两个条件只满足该方法的截断误差顺序。 出于这两个原因，这将导致在连续性、x动量和能量方程中对源项Gc的贡献很小。 通过使用坐标不变表达式（3.3.2），我们完全满足（3.3.6），并避免了这些虚假的数值结果。 8.3.4 Riemann 求解器 8.3.4-1 通量雅各布和特征向量基于特征的滤波器需要在平均状态下使用通量雅各布矩阵的特征值和特征向量。 最好使用原始的保守向量Q而不是Q来处理它们。 设S\textasciicircum{} = S\textasciicircum{}ir*，其中n\textasciicircum{}是单位法线，S\textasciicircum{}是表面积向量S\textasciicircum{}的大小。 类似的定义也适用于Sn和S\textasciicircum{}。 从现在开始，我们省略曲线坐标上标。 在任何平均状态下，只有适当的正态n-n\textbackslash\{\}i + nj + 713k的方向发挥作用。 让un = u • n，vn = v • n，以及u' = un - vn。 变换的通量向量可以写成S-F = SFn，其中Fn = n-F。 然后将通量Jacobian 矩阵定义为A = (Fn)Q = dFn/dQ。 我们将首先发展一般平衡气体的关系。 压力p由P = P(P,«)-（3.4.1）形式的一般状态方程给出

152 VINOKUR \& YEE 导数将用- |和K = dp（3.4.2）表示。如果h = (e + p)/p是比焓，那么音速c可以表示为c2 = X+Kh。 （3.4.3）然后矩阵A可以写成A = (3.4.4)，其中K = \textbackslash\{\}nq2 + x，H - h+ \textbackslash\{\}q2是单位质量的总焓。 A的三个不同的特征值是-vn Kn\textbackslash\{\} - unu Kn2 - unv Kns - unw (K - H) un (1 - K)UTI\textbackslash\{\} + u' VTl\textbackslash\{\} - KT12U Wn\textbackslash\{\} - KT13V, Hn\textbackslash\{\} - Kunu uri2 - Kn\textbackslash\{\}v (1 - K)VTI2 + u' wri2 - nnzv Hn.2 - Kunv unz - KTIIW VT13 - KU2W (1 - K)WU3 + U1 Hnz - Kunw K«l K.U2 KT13 u' + KUH A1 X' X' and XJ c, (3.4.5)，其中A1是重复的特征值。 为了获得相应的线性独立特征向量集，我们引入了满足a* • aJ = 5U的任意正交基a'，其中5i:>是Kronecker delta。 然后可以定义am = n • a*和bl = nxa'。 右特征向量矩阵R可以写成R = anl anlu + cb\textbackslash\{\} anlv + cb\textbackslash\{\} anlw + cb\textbackslash\{\} anlK + c(b1 -u)，其中K=\textbackslash\{\}q2-X/K，a"2 an2u + cb\textbackslash\{\} an2v + cb\textbackslash\{\} an2w + cb2an2K + c(b2 • bi = 6ii + 6*aj an3 an3u + cbf an3v + cb3，an3w + cb\textbackslash\{\} an3K + c(b3 • u) + 63k，以及b' • u = 1 1 u + cn\textbackslash\{\} u - v + CT12 V - CU2 w + cn.3 w - cn3 H + cun H - cun (3.4.6) b[u + biv + biw。 1 7he数量a = R 1AQ最简单地表示为a anl(Ap - Ap/c2) + p\{b\textbackslash\{\}Au • (Ap - Ap/c2) + p(b\textbackslash\{\}Ai,n2 b\textbackslash\{\}Av-\textbackslash\{\}-b\textbackslash\{\}Aw)/c b\textbackslash\{\}Av + blAw)/c an3(Ap - Ap/c2) + pibfAu + b\textbackslash\{\}Av + blAw)/c \textbackslash\{\}\{Ap/c2 + p[rnAu + n2Av + n3Aw]/c) I (Ap/c2 - p[niAu + n2Av + nzAw\textbackslash\{\}/c) (3.4.7) 请注意，在计算过滤器时，我们需要a，我们需要i？ \_1 明示。 为了效率，人们应该使用3.4.7，而不是用向量AQ做矩阵R\textasciitilde{}x的实际乘积，特别是对于非平衡流动（Yee \& Shinn 1989，Yee [22]）。

曲线移动网格的低消散方案153和有三种方法来计算am和b\textbackslash\{\}它们都需要每个平均状态下的n知识，通过计算S = 5n作为两个相邻网格点的S的矢量平均值获得。 这与过滤程序的准确性一致。 每种方法都有一定的缺点。 如果存储不是问题，那么首选第三种方法。 在第一种方法中，笛卡尔单位向量被用作基数a*。 因此，我们定义了ax = i，a2=j，a3 = k。 （3.4.8）然后a"1 = m，a"2 = n2，a"3 = n3，(3.4.9) b\textbackslash\{\} = 0 b\textbackslash\{\} = n3 b\textbackslash\{\} = -n2 b\textbackslash\{\} = -m b\textbackslash\{\} = Q bl = nx (3.4.10) b\textbackslash\{\} =n2 b\textbackslash\{\} = -m \&3 = 0。 这种方法的缺点是，R和R\textasciitilde{}1AQ都涉及许多需要评估的项。 如果正态n是基数a1的一个成员，则其中一些术语将不存在。 这种选择是另外两种方法的底力。 然后，我们可以定义a1 = n，a2 = t，a3 = s，（3.4.11），其中n，t，s形成正交三和，t尚未定义，但满足n • t = 0，s = n x t。 有了这个选择，我们有flnl = Xj fln2 = Q) an3 = \textasciicircum{} (3.4.12)和b\textbackslash\{\} = 0 b\textbackslash\{\} = Q b\textbackslash\{\} = 0 b\textbackslash\{\} = sx b\textbackslash\{\} = s2 bl = s3 (3.4.13) bl = -h b\textbackslash\{\} = -t2 63 = -t3。 定义t的最佳方法是c x n |c x nl（3.4.14），其中c是一些任意的向量。 C的选择区分了另外两种方法。 在第二种方法中，c被当作一个常数。 一个常见的选择是让c = i+j + k。 （3.4.15）

154 VINOKUR \& YEE 如果存储是个问题，这涉及少量额外的操作，并且必须在每个时间步骤中重新计算t和s。 这种方法的缺点是，当c和n平行时，它会崩溃。 为了避免这种发生，在第三种方法中，c被选为相邻网格点的向量S之一。 这种方法比第二种方法涉及更多的操作，但更稳健。 例如，对于（i + 1/2，j，k）处的£曲面，我们可以让c在（i，j，k）处为S\textasciicircum{}。 8.3.4.2 Roe平均值 在各种近似的黎曼求解器中，最常见的求解器使用Roe平均值，因为它的简单性和能够完全满足跨不连续的跳跃条件。 在那些基于局部线性化的求解器中，分隔两个状态QLand QR的点的通量是基于一些平均值A的特征值和特征向量。 A的最佳选择是满足AFn = AAQ，（3.4.16），其中A（-）-（-）R-（-）L。 A的这种选择准确地捕捉到了不连续性。 获得A的一种方法是寻求平均状态Q，这是QL和QR的函数，这样\_\_ A = A(Q)。 （3.4.17）这种状态被称为罗埃平均状态[12]。 人们可以很容易地证明u = vuL + (1 - u)uR (3.4.18)和H = vHL + (1 - u)HR，其中\_ 1 "\textasciitilde{} 1 + y/pR/pL ' 从H的定义中得到h = H u • u。 在（3.4.7）中，我们还需要p，由p= \textbackslash\{\}/pRph给出。 （3.4.22）我们仍然需要评估x和K，它们用于从（3.4.3）中确定c，从其定义中确定K。 对于一般平衡气体，提出了各种近似方法来计算x和K。 由于熵分裂仅适用于热完美气体，我们只考虑这种情况，其中（3.4.19）（3.4.20）（3.4.21）

\section{8.4 结束说}
曲线移动网格的低消散方案155和方程和（3. 4.23a,b)是Ae eRfL \_ eLfR X\textasciitilde{} Ae被K = l/e X = T - «e取代，有一个精确的解。 详情见Yee等人。 [25]。 当Ae -> 0时，确切的公式是\_ AT（3.4.23a）（3.4.23b）（3.4.24a）（3.4.24b）。 对于完美的气体（3.4.23a，b）减少到X = 0和7c = 7-1。 （3.4.25）为了获得过滤器（3.2.1）的Ri+i和\$i+i+i，右状态和左状态（（3.4.16）-（3.4.23）中的上标R和L）是网格索引（i+ l,j,k）和（i,j,k）。 8.3.4-3 非平衡流动 在Yee等人[25]中，我们表明，将熵分裂扩展到完全耦合的非平衡流动是不可行的。 但有关方案是可用的，此外，对于非平衡流动，人们可以获得Roe的Riemann 求解器的确切扩展。 这在附录8.4中介绍了结论。在本文中，我们制定了Yee等人[24]方案，在3D曲线移动网格中具有熵分裂（Yee等人[25]）的热完美气体。 为了提高效率，我们选择了有限差分公式。 通过将托马斯和伦巴第公式的坐标不变版本离散化，表面积守恒律在数值上得到满足。 作者于2000年6月独立提出这种形式，Thomas \& Neier [19]（P.D. 托马斯，私人通讯，2000年11月）。 虽然熵分裂的正式扩展仅限于热完美气体，但数值方案本身没有这种限制。 因此，这里讨论的方案适用于平衡实气、非平衡

\section{认可}
\section{附录A：一类数值混合偏导数的交换性}
156 VINOKUR \& YEE和人工可压缩性处理不可压缩流、MHD和Maxwell方程的方法。 此外，Sjogreen和Yee提出的双重用途小波传感器（动态数值耗散控件和网格自适应指示器）可以作为这里讨论的以外的各种方案的独立选择。 在Miiller \& Yee [8]和Polifke等人中可以找到一般坐标变换中由与流动变量相同的高阶差分算子离散化的度量项的数值实验。 [10]。 说明新三维度量公式性能的数字示例将在未来的论文中报告。 致谢 特别感谢Tom Coakley和Dennis Jespersen对手稿的批判性审查。 附录A：一类数值混合偏导数的交换性 在本附录中，我们证明了数值混合偏导数的交换，因此表面积守恒律完全满足。 我们要感谢美国宇航局艾姆斯研究中心的Dennis Jespersen提供了证明的基本要素。 我们发现使用与正文不同的符号很方便。 大写字母表示矩阵，小写字母表示代数向量，拉丁语下标表示其组成部分。 我们首先介绍张量积（或克罗内克积）的概念。 给定两个任意矩阵A和B，张量积A <g）B被定义为块矩阵，其每个分量都是矩阵B乘以矩阵A的相应分量。 在组件形式中，它写成\{A®B)ij = AijB。 （A.l）从定义（A.l）中，我们可以得出一个重要的混合产品规则。 如果A和C是符合矩阵，B和D也是符合矩阵，那么（A ®B）（C®D）= AC® BD。 （A.2）

曲线移动网格的低消散方案157 证明紧接着写（A. 2）以组件形式。 因此[(A ® B)(C ® D)]ik = J\textasciicircum{}(A ® B)y(C ® D)ifc = £)AyflC,-fc£> = (AC)ikBD = \{AC®BD)ik。 （A.3）让£、77、C的计算空间分别在£、77、£方向上用/、m、n点离散化。 对于固定的77和£，£导数最普遍的有限差分近似是A\textasciicircum{}ui: = B\textasciicircum{}u，（A.4），其中u和u\textasciicircum{}是/维向量，A\textasciicircum{}和B\textasciicircum{}是/by/矩阵。 一些例子是中央非紧凑和紧凑的空间方案（3.1.1）-（3.1.3）。 对A\textasciicircum{}和B\textasciicircum{}的性质没有限制，它们在计算区域的£边界上包含任意的边界条件。 假设整个区域的离散未知数是按以下顺序排列的：£值先变化，77值变化：77值变化，（值）最后变化。 如果相同的有限差分近似（A.4）适用于每个77和£（这意味着沿每个£边界的相同边界条件），那么所有未知数的£导数的近似值是A\textasciicircum{}ut：= B\textasciicircum{}u，（A.5），其中u和u\textasciicircum{}是Ixmxn-维向量，A和B arelxmxn由]\textasciicircum{} = /"®(Jm®\textasciicircum{})，BC = /" ® (7m ® £«)给出的I x xm x n矩阵。 （A.6）7"和7m分别是n乘n和m乘m的恒等式矩阵。 请注意，（A.6）中的括号可以消除，因为从其定义中可以看出张量乘法是关联的。 对于固定的£和£，77导数的有限差分近似可以写成A"vv = B\textasciicircum{}v，（A.7），其中v和v\textasciicircum{}是m维向量，A71和B71是m矩阵的m。 请注意，An和Bn又是任意的，77边界上的边界条件与允许的£边界不同。

158 VINOKUR \& YEE 如果相同的有限差分近似（A.7）适用于每个£和（，那么所有未知数的I x m x n-维向量v和vn由A\textbackslash\{\} = Wv，（A.8）相关联，其中Ixmxn矩阵A和B的Ixmxn由A\textasciicircum{} = In®\{Ar>®Il），\textasciitilde{}W = In®\{B\textasciicircum{}®Il）给出。 （A.9）I1是/by/恒等矩阵。 同样，对于固定的£和77，C导数的有限差分近似可以写成Acwc = B\textasciicircum{}w，（A.10），其中w和w\textasciicircum{}是n维向量，A\textasciicircum{}和B\textasciicircum{}是n乘n矩阵。 A''和B1*又是任意的，£边界上的边界条件与允许的其他边界不同。 如果相同的有限差分近似值（A. 10）适用于每个£和77，然后所有未知数的Ixmxn-维向量w和w\textasciicircum{}由ZCWC = BCw相关，（A.ll），其中Ixmxn矩阵A和B的Ixmxn由AC = A<:®\{Im®Il）给出，BC = Bt®（Im®Il）。 （A.12）让向量u\textasciicircum{}是所有混合£77偏导数的有限差近似值，通过首先使用（A.5）获得u\textasciicircum{}，然后使用（A.8）获得u\textasciicircum{}来计算。 如果网格点在计算空间中间距相等，那么矩阵A\textasciicircum{}、B\textasciicircum{}、An、B71、A1*和B\textasciicircum{}都是常数。 然后最终结果变成A"\textasciicircum{}\textasciicircum{}),, = TA\textasciicircum{}r, = B\textasciicircum{}B\textasciicircum{}u。 （A.13）如果我们按相反的顺序执行操作，我们得到AiAnuvi = BiBr'u。 （A. 14）使用（A.2）、（A.6）和（A.9），它遵循TA？ = [In ® (A" ® /')][/" ® (Im ® A*)] = J"®[(i4,»®j')(JTn®A€)] = 7" ® \{A71 ® \textasciicircum{}). （A.15）

曲线移动网格的低消散方案159以同样的方式，我们表明A* A* = [In ® (Im ® At)][r <8> (A" ® /')] = In®[(Im®AS)(Ar>®I1)\} = In<8)\{AJ'®At)。 （A. 16）因此从（A. 15）和（A. 16）我们已经确定了T\textasciitilde{}A< = A<T。 （A.17）以类似的方式，我们表明B“B€ = BCB”。 （A. 18）来自（A. 13）和（A. 14）它遵循Hv = *\%•（A. 19）以类似的方式，使用（A.5）和（A.11），所有混合££-部分导数的有限差分近似可以证明满足AC(A\%)C = A\textasciicircum{}u\textasciicircum{} = B\textasciicircum{}u。 （A.20）如果我们按相反的顺序进行运算，我们得到A\textasciicircum{}u\textasciicircum{} = B\textasciicircum{}u。 （A.21）使用（A.2）、（A.6）和（A.12），可以得出A\textasciicircum{} = Ai®（Im®AZ）= A<A<，（A.22）具有涉及B和B的类似关系。 从（A.20）和（A.21）中，u\textasciicircum{} = Htf。 （A. 23）使用（A.8）和（A.11）的相同程序可用于表明*V = Hv（A.24）这完成了所有数值混合偏导数交换的证明。

\section{附录B：非平衡流动的Riemann 求解器}
160 VINOKUR \& YEE 附录B：非平衡流动的Riemann 求解器 在本附录中，我们介绍了化学非平衡流动的Riemann 求解器。 分析可以很容易地扩展到热化学非平衡。 我们考虑了ns化学物种的混合物，每个物种都遵循热完美气体定律。 让ps、Rs和es(T)分别是物种s的密度、气体常数和内能。 T是混合物的温度，假定处于热平衡状态。 请注意，我们没有引入规范化的T，因为每个物种的R都不同。 所有物种都满足条件es > 0，其中es = des/dT。 我们现在可以将混合物的密度p定义为物种质量分数为（B.2）压力p，其中P = pRT，混合物的单位体积内能为（B.3）（B.4）（B.5）守恒律的形式为Qt + V • F = Sc（B.6），其中<SC是由每个物种的源项组成的向量。 向量Q和F的形式为Q（ps）l pu pv pw e。，F = r（PS»'）puu' + pi pvu' + pj pwu' + pk。 (E +p)u' +pv (B.7) 请注意，F再次相对于运动网格定义。 Q和F是（ns + 4）-维向量。 我们使用紧致符号“（ps）”和“（psu'）”，其中（B.7）中每个向量写的第一个项是一般成分

曲线移动网格的低消散方案161 ns维子向量对应于ns物种连续性方程。 温度T(Q)是通过隐式求解方程5>V(T) = e - \textasciitilde{}[\{puf + \{pvf + \{pwf\textbackslash\{\}\}得到的。 （B.8）s“方程（B.8）有一个独特的解，因为ps > 0和es > 0。 为了推导通量Jacobian 矩阵，压力p必须以p = p（ps，Z）的形式表示。 （B.9）衍生品将用H£L md *-（£）表示。 •（B-M）在（B.l）到（B.5）的帮助下，我们可以表明V asRs «=\textasciicircum{}\textasciicircum{} (B.ll)和XS = RST-K€S。 （B.12）然后矩阵A可以写成' (6sru'- a"un) (asnx) \{asn2) \{asn3) 0 \{KTn\textbackslash\{\} - unu) (1 - K)un\textbackslash\{\} + u' un2 - KTI\textbackslash\{\}V unz - nn\textbackslash\{\}w KUI \{Krn2 - unv) vni - KU2U (1 - K)VU2 + v! Vni - KU2W reri2 \{Krnz - unw) wn\textbackslash\{\} - Kn\textasciicircum{}u wn2 - KT13V (1 - K)WU3 + u' nnz ((Kr - H) un) Hm - nunu Hn2 - KUnv Hn3 - Kunw u' + KM" (B.13) 其中 KT = \textbackslash\{\}\&q2 + xr and <5sr是Kronecker delta。 所有其他数量在第3.4节中定义。 再次，A的第一行中写入的术语再次使用紧凑符号。 用A写的第一个项是ns子矩阵的ns的sr分量，第二个到第四个项是ns维度的列子向量。 此外，写在A第一列中的最后四个项是维度ns的行子向量。 A的三个不同的特征值又是A1 = u'，A2 = v！ + c，和A3 = u \textasciitilde{} c，（B.14），其中c是c2 = Y\textasciicircum{}aSXS + Kh-（B-15）给出的冻结音速

162 VINOKUR \& YEE A1是多重性ns+2的重复特征值。 选择笛卡尔单位向量（3.4.8）作为基数a*以获得线性独立的特征向量集，这将导致R和R\textasciitilde{}lAQ中对于非平衡流动的额外项太多。 因此，我们把自己限制在三和n、t、s作为基础（3.4.11）。 然后，右特征向量矩阵R可以写成R = (5sr) u w Kr CS\textbackslash\{\} cs2 cs3 c(s • u) -ct\textbackslash\{\} -ct2 -ct3 -c(t • u) («s) u + cni V + CJl2 W + C7l3 H + cun K) l u - cn\textbackslash\{\} V - C712 W - C«3 H -cun。 （B.16）其中Kr - \textbackslash\{\}q2 - xr/K- 数量R AQ最简单地表示为R\textasciitilde{}XAQ = (Aps - asAp/c2) p(siAu + s2Av + s3Aw)/c -p(hAu + t2Av + t3Aw)/c \textbackslash\{\}\{Ap/c2 + p[nxAu + n2Av + n3Aw]/c)。 \textbackslash\{\}\{Ap/c2 - p[n\textbackslash\{\}Au + n2Av + n3Aw]/c) (B.17) 定义Roe平均变量的方程（3.4.16）到（3.4.22）仍然有效。 此外，我们将T = vTL + (1 - v)TR，以及a* = uasL + (1 - p)定义为i* = uesL + (1 - v)e"R。 然后，如果我们定义x* = RST-K€S，我们可以表明（3.4.16）是满足的。 当AT -)• 0时，方程（B.21）被-=Zs*sRs Es aS£s（B.18）（B.19）（B.20）（B.21）（B.22）（B.23）替换。

\section{参考资料}
曲线移动网格的低消散方案163参考1。 M.H. 木匠，J。 诺德斯特龙和D。 Gottlieb（1998），任意空间准确性的稳定和保守界面处理，ICASE报告98-12。 2. G。 达尔奎斯特（1979年），常微分方程的线性多步和单腿方法的一些属性，会议记录，1979年关于数值ODE的SIGNUM会议，香槟，111。 3. D.V. Gaitonde和M.R. Visbal（1999），进一步开发基于高阶公式的Navier-Stokes解程序，AIAA论文99-0557，1999年1月。 4. M。 Gerritsen和P。 Olsson（1996），为流体流动设计高效解决方案策略：I. Euler 方程的稳定高阶有限差分方案和锐利的冲击分辨率，J。 计算机。 物理。 129，245。 5. A. Harten（1978），计算冲击和接触不连续的人工压缩方法：III。 自我调整混合方案，数学。 补偿。 32，363。 6. A. Harten（1983），计算双曲守恒定律弱解的高分辨率方案，J. 计算机。 物理。 49，35。 7. A. Harten（1983年），关于熵守恒定律系统的对称形式，ICASE报告81-34，美国宇航局兰利研究中心，1981年，J。 计算机。 物理。 49，151。 8. B. Miiller和H.C. Yee（2001），Kirchhoff Vortex生成的声音的高阶数值模拟，RIACS技术报告01.02，1月。 2001年，美国宇航局艾姆斯研究中心，提交给J。 科学中的计算和可视化。 9. P。 Olsson和J. Oliger（1994），非线性守恒定律的能量和最大规范估计，RIACS技术报告94-01，美国宇航局艾姆斯研究中心。 10. W。 波利夫克，B。 Miiller和H.C. Yee（2001），发电机外壳转子诱导变形的声音发射，AIAA/CEAS航空声学会议记录，2001年5月28日至30日，荷兰马斯特里赫特；还有，RIACS技术报告，2001年3月，美国宇航局艾姆斯研究中心。 11. T.H. Pulliam和J.L. Steger（1978年），关于三维流动的隐式有限差分模拟，AIAA论文78-10，1978年1月。 12. P.L. Roe（1981），近似黎曼求解器、参数向量和差分方案，J。 计算机。 物理。 43，357。 13. N。 D。 桑德姆和H.C. Yee（2000），可压缩湍流高阶数值模拟的熵分割，RlACS技术报告00.10，2000年6月，美国宇航局艾姆斯研究中心；还有，2000年7月10日至14日，日本京都，第一届CFD国际会议记录。 14. B. Sjogreen和H.C. Yee（2000），基于小波的自适应数值耗散控制冲击-湍流计算，RIACS技术报告01.01，2000年10月，美国宇航局艾姆斯研究中心。 15. B. Sjogreen和H.C. Yee（2001），小规模黏性可压缩流动的激波捕捉方案的网格收敛研究，AIAA CFD会议记录，6月11日至14日，加利福尼亚州阿纳海姆。 16. B. 斯特兰德（1994），d/dx的有限差近似的部分和，J。 计算机。 物理。 110，47。 17. J.L. Steger（1977），关于任意几何的流动的隐式有限差分模拟，应用于机翼，AIAA论文77-665。 18. P.D. 托马斯和C.K. 伦巴第（1979年），几何守恒律及其

164 VINOKUR和YEE在移动网格上的流动计算中的应用，AIAA J. 17，1030。 19. P.D. 托马斯和K.L. Neier（1990），Navier-Stokes与消融三维高超音速平衡流模拟，J. 航天器和火箭27，143。 20. M。 Vinokur（1989），守恒定律的有限差和有限体积公式分析，J. 计算机。 物理。 81，1。 21. H.C. Yee（1985），显式和隐式对称TVD方案的构建及其应用，J. 计算机。 物理。 68,151（1987）；也是美国宇航局TM-86775，1985年7月。 22. H.C. Yee（1989），一类高分辨率显式和隐式冲击捕捉方法，VKI系列讲座1989-04，1989年3月6日至10日，还有NASA TM-101088，2月。 1989年。 23. H.C. 是的，G.H. Klopfer和J.-L. Montagne（1990），黏性和黏性高超音速流的高解析度休克捕捉方案，J. 计算机。 物理。 88，31。 24. H.C. 耶，N.D. Sandham和M.J. Djomehri（1999a），使用基于特征的滤波器的低耗散高阶激波捕捉方法，RIACS技术报告98.11，1998年5月；还有，J. 计算机。 物理学，150，199。 25. H.C. 是的，M。 Vinokur和M.J. Djomehri（1999），熵分裂和数值耗散，美国宇航局技术备忘录208793，1999年8月，美国宇航局艾姆斯研究中心；也J。 计算机。 物理，162，33，2000。 26. H.C. 是的，B。 Sjogreen，N.D. Sandham和A。 Hadjadj（2000）开发一类高效低耗散高阶激波捕捉方法的进展，RIACS技术报告00.11，2000年6月，美国宇航局艾姆斯研究中心，21世纪CFD会议记录，2000年7月15日至17日，日本京都。

\clearpage
\chapter{交错网格上的斯托克斯和Navier-Stokes 方程的9个四阶方法}
\section{9.1 介绍}
交错网格上的斯托克斯和Navier-Stokes 方程的四阶方法 Bertil Gustafsson1和Jonas Nilsson x 献给Robert MacCormack 60岁生日 9.1 介绍 在本文中，我们考虑了2D Stokes方程Ut+Px = v\{uxx+uyy），Vt+Py = V（VXX + Vyy），UX + Vy = 0，以及不可压缩的2D Navier-Stokes 方程 Ut + UUX + VUy +PX - V（UXX + Uyy），Vt + UVX + Wy +Py = U（VXX + Vyy），UX + Vy = 0。 在这里，u和v分别是x和\textasciicircum{}方向的速度分量，p是压力。 我们为两组方程开发了帕德类型的四阶紧差分方法。 我们将近似应用于原始形式，并明确使用发散条件，作为乌普萨拉大学信息技术1科学计算系的一部分，信箱120，SE-751 04乌普萨拉，瑞典。 Frontiers of 计算流体动力学 - 2002 编辑：David A. Caughey和Mohamed M. 哈菲兹 ©2002 世界科学

166 GUSTAFSSON \& NILSSON系统。 边界条件基于稳态斯托克斯方程的[2]结果。 在大多数应用中，速度场v = (u, v)T = v0在整个边界<9Q给出，几乎所有的分析都专门用于这种情况。 然而，有时需要规定压力，至少在部分边界上，我们也将考虑此案。 发散条件引入了对边界数据的限制/vo-ndS = 0，（9.1）JdQ，其中n是dtt的外法线。 这种限制可以被看作是边界值问题是奇异的这一事实的影响。 其结果是，任何一致近似求解的代数系统也必然是奇数的。 通过以特定方式为斯托克斯方程组制定边界条件，使边界值问题变得非奇异，我们能够以直截了当的方式构建非奇异近似。 近似是基于交错网格上的二阶差分近似。 这种类型的近似已被其他人使用，例如见[3]、[6]、[5]和[7]，其优点是不存在寄生解。 在我们的案例中，它还允许进行与连续案例分析非常接近的分析，我们能够得出相同类型的估计值。 其中一些结果在第9.2节中进行了总结。 在第9.3节中，我们描述了斯托克斯方程的交错网格上的四阶紧致差分方法。 通过在边界附近使用非对称近似值，差分运算符在所有网格点上都定义得很好，与手头的特定初始边界值问题无关。 代数系统是通过Krylov类型的迭代方法求解的。 在每次迭代中，通过应用物理边界条件来修改中间溶液。 事实证明，稳定斯托克斯方程的结果很容易延续到与时间相关的方程上，我们应用了完全相同的技术。 此外，这些结果对Navier-Stokes 方程来说也很重要。 为了方便起见，我们以vt + Stokes(v,p) - 0, div v = 0的形式写出斯托克斯方程，其中斯托克斯(v,p)表示动量方程的空间部分。 时间-离散化是通过使用向后微分公式|v«+i + AtStokes(vn+l,pn+l) = 2v" - iv""1, divvn+1 = 0, l\textasciicircum{}j获得的

\section{9.2 稳定的斯托克斯方程和交错网格}
Navier-Stokes 方程 167的四阶方法，其中上标表示时间级别。 接下来考虑Navier-Stokes 方程，我们以vt + Navier(v) + Stokes(v,p) = 0, div v = 0的形式写它们，其中Navier\{v)表示对流项。 对于时间-离散化，我们使用Per Lotstedt提议的半隐式方法。 它基于上述隐式方法方程的斯托克斯部分：-vn+1 + AtStokes(vn+\textbackslash\{\}pn+1) = 2v"-iv"-1 -At(2 Navier\{vn) - Navier(vn-1), (9.3) divvn+1 = 0。 这允许大时间步骤，而不必在每个时间级别解决完整的非线性系统，这是完全隐式方法所必需的。 假设我们为斯托克斯方程有一个高效且条件良好的求解器，我们也为Navier-Stokes 方程提供了一个高效且条件良好的求解器。 众所周知，对于湍流的直接模拟，需要高阶方法。 在第9.4节中，我们使用与斯托克斯方程相同类型的紧凑高阶方法来近似Navier-Stokes方程，并为一个简单的模型问题提出了一些初步结果。 9.2 稳定斯托克斯方程和交错网格 在本节中，我们将总结[2]中的一些结果，并讨论消除系统中奇点的特定类型的边界条件。 我们考虑稳态方程Px = v\{uXx + uyy) i Py = V\{VXX + Vyy), (9.4) ux + vy = 0。 理论结果是在y方向的周期解的情况下获得的：XJ(x,y) = XJ(x,y + 2ir)，其中U = (u,v,p)T，方程定义在矩形域fi = \{0 < a: < 1, 0 < y < 27r\}中。 如果速度u，v在两边x - 0,1上规定，则数据上的限制（9.1）是/*27r /»27r / u(0,y)dy= / u(l,y)dy, Jo Jo

168 GUSTAFSSON \& NILSSON 为了消除作为此限制基础的奇点，我们将边界条件制定为1 f2* "(0,2/)-\textasciicircum{}/ u(0,y)dy = w0(y), u(l, y) = ux\{y), v\{0,y) = w0(y), w(l,? /) = ui(y), (9.5) 2*- i o 2TT，其中假定J0 wo\{y)dy = 0。 然而，应该注意的是，即使这个积分中存在小扰动，问题仍然可以通过其结果解决的小扰动来解决。 使用符号VB = \{WQ,U\textbackslash\{\},VQ,V\textbackslash\{\})T，以及由2TT /-l 2定义的规范，2 U + V. l|v(-,-)ll = \{/ \textbackslash\{\}v(x,y)\textbackslash\{\}2dxdyy'2, \textbackslash\{\}v\textbackslash\{\}2 = Jo Jo |VB(0II = \{ /\textasciicircum{} |v*(l/) W2, |vB|2 = K|2 + \textbackslash\{\}Ul\textbackslash\{\}2 Jo \textbackslash\{\}\textbackslash\{\}p(;-)\textbackslash\{\}\textbackslash\{\} = \{r flpfrytfdxdy\}1'2, Jo Jo我们有定理9.1 假设边界数据u>o, UI, vo, v\textbackslash\{\}在y和J0 wo\{y)dy = 0中是2rr-周期性的，但其他方面是任意的。 然后，边界条件（9.5）的系统（9.4）有一个独特的解决方案，有一个估计||v(-,-)ll2 < co\textasciicircum{}(M2 + ||vB(-)||2), lb(-,-)ll2 < con5i(k0|2 + ||vB(.)||2+|\textasciicircum{}(-)l|2)。 （9'6）我们可以用压力代替其中一个成分，而不是规定两个速度成分。 条件是1 rv u(0,y)-- / u(0,y)dy = w0\{y), u\{l,y) = ux\{y), 1 r2* P(0,y)-\textasciicircum{}/ p\{0,y)dy = qo(y), p(l,y) = Pi(y), (9.7) /»27T p2lT / v(0,y)dy= z0, / v(l,y)dy = z1。 Jo Jo 在这种情况下，在估计解决方案时不需要边界数据的导数。 我们有||U||2 = ||v||2 + ||p||2

Navier-Stokes 方程 169定理9.2的四阶方法假设边界数据WQ、U\textbackslash\{\}、qo、p\textbackslash\{\}在y和fQnWo\{y)dy = |0'qo(y)dy = 0中是2n-周期的，但其他方面是任意的。 然后，带有边界条件（9.7）的系统（9.4）有一个独特的解决方案，有一个估计值|U(-,-)||2 < const (\textbackslash\{\}z0\textbackslash\{\}2 + \textbackslash\{\}zi\textbackslash\{\}2 + + / (\textbackslash\{\}w0(y)\textbackslash\{\}2 + M2/)|2 + \textbackslash\{\}q0(y)\textbackslash\{\}2 + \textbackslash\{\}Pi(y)\textbackslash\{\}2)dy)。 (9.8) Jo 边界条件的第三组是v\{0,y) = vo(y), v\{l,y) = vi(y), 1 f2v p(0,? /)--/ p(0,y)dy = qo(y), P(l,y) = Pl(y), (9.9) r2iv u(0,y)dy = wQ. Jo >0 和上述方式相同，我们可以证明定理9.3，假设边界数据VQ、V\textbackslash\{\}、qo、p\textbackslash\{\}在y和JQnqo(y)dy = 0中是2it-周期性的，但其他方面是任意的。 然后，边界条件（9.9）的系统（9-4）有一个独特的解决方案，有一个估计||U(-,-)||2 < const(\textbackslash\{\}w0\textbackslash\{\}2 + + f\textbackslash\{\}\textbackslash\{\}My)\textbackslash\{\}2 + \textbackslash\{\}vi\{y)\textbackslash\{\}2 + \textbackslash\{\}q0(y)\textbackslash\{\}2 + \textbackslash\{\}Pi(y)\textbackslash\{\}2)dy).( 9.lO）Jo对于空间中的离散化，我们根据图1使用交错的N x M网格。 标准二阶方案适用于内点D-tXpitj = v(D+tXD\_tX + D+\textasciicircum{}D-\textasciicircum{}u\textasciicircum{}ij, D-zyPij = v\{D+,xD-tX + D+,yD-ty)vitj\_\%, (9.11) D+! xUi\_ij + D+tyvitj\_i =0，其中D+tXUij = (v,i+ij - Uitj)/Ax和D-tXUij - (uij - Ui-ij)/A.x，D±\textasciicircum{}y也是如此。 离散边界条件是连续的直接近似值，除了一个速度分量，我们取平均值。

170 GUSTAFSSON \& NILSSON * O * G * G * * * G * O * G * *-e *-G * • • • 图1 A 交错网格。* u • v O P 第一组条件（9.5）近似为2 iJ 2 'J 1 \textasciicircum{}U\_l,•+«!, A x 2 'J 1'J 2TT \textasciicircum{} j=o Ay = tvoivj) lN- + UN+\textbackslash\{\},j ']! ' V0J\_l =V0(j/j\_j), "JV,j- = Wl(\%), M-l (9-12) 其中X\textasciicircum{}=(/ wo(llj)Ay - 0. 对于这个近似值，人们可以得出一个与连续估计值（9.6）完全相似的估计值：•vAf-l |v||2 < const (M2 + ||vB||£), ||p||fc < const (\textbackslash\{\}qo\textbackslash\{\}2 + \textbackslash\{\}\textbackslash\{\}vB\textbackslash\{\}\textbackslash\{\}l + \textbackslash\{\}\textbackslash\{\}D+,yVB\textbackslash\{\}\textbackslash\{\}l), (9.13) 其中常数独立于Ax, Ay。 这里的离散规范|| • ||h是上述定义的连续规范|| • ||的直接离散化。 周期解的假设使使用傅里叶技术推导出上述估计值成为可能。 在大多数应用中，解决方案是非定期的，我们也必须在y方向上规定边界条件：«i-I,0 =U°0ci-±)> wi-i,M =U1(xi\_i), VL\_l + l>, 1 \textbackslash\{\}xt) Vi,M\textasciitilde{}h "*" Vi,M+k - v1\{xi)。 （9.14）

\section{9.3 斯托克斯方程的四阶方法}
Navier-Stokes 方程 171的四阶方法，在这里，u°、v°和u1，vl分别在y = 0和y = 1处给出数据。 [2]中的数值实验表明，在周期性情况下得出的估计值也适用于非周期性情况。 对于时间依赖性问题（9.2），（9.3），这里介绍的边界条件的交错网格和特殊形式也是相同的。 此外，它是下一节中要讨论的高阶近似的基础。 9.3 斯托克斯方程的四阶方法 四阶紧致方法可以得出正则网格和交错网格，见[4]和[1]。 对于斯托克斯方程，所有一阶导数都以两个网格点之间的中间为中心，而所有二阶导数都以一个网格点为中心。 对于函数/和网格间距h，我们有公式24 12 i/iJ*- 24 "+1 \_ l'Jl+1/2 -A\_ f" 4. f" \_|\_ J\_ f" = -tnJi - l\textasciitilde{}Ji \textasciitilde{} -iriJi+1 5\textasciicircum{}2 /i\_i/2), (9.15) (/i+i - 2fi + /4\_i)。 （9.16）为了解决未知数/'和/"，我们需要Pf = Qf，Rf" = Sf形式的闭合系统。 衍生品没有可用的边界条件，我们使用单边公式。 对于右侧的功能本身，我们可以使用上面给出的物理边界条件，因为右侧完全对应于上述二阶方案。 然而，独立于特定问题定义P、Q、R、S很方便，因此我们也为/使用单边公式。 第一导数有两种选择Pi、Qiand P2、Qi，这取决于这些算子是用于压力还是速度。 完整的公式是h -> Ax Pi/' = Qif: Pif = Qif = Bb(/l*+23/i), 552 \textasciicircum{}(/;\_:+22\textasciicircum{}+l), B1J\textasciicircum{}(-25/o + 26/1-/2), 2£i(/i+i-/i),; = 0,...,iV\_i; I\textasciicircum{}(/iv-2-26/\textasciicircum{}\_1+25/w)。 0，...，JV-1，（9.17）

172 GUSTAFSSON \& NILSSON Pif = Q*f: Pif = Q2f = 24 V /o + /l) > 24UJV-1 \textasciitilde{}\textasciitilde{} /jv) > 24Ax V JN /§). i = l,...,7V,JV-1, (9.18) iv+i)- i? /" = Sf: Rf" = \{ XOO(10/JV-I + /AT) ',N-1。 (9.19) l \textasciicircum{}\textasciicircum{}(145/0-304/x • 174/2 - I6/3 + h), -i - Vi + fi+i), i = l iV-1, I2OOA\textasciicircum{}(/\textasciicircum{}-4 - 16/JV-3 + 174/w\_2 - 304/w\_1 + 145/JV) 这些公式是一维的，适用于\textasciicircum{}方向的所有网格线。 它们也适用于y方向，但现在Ax和N分别被Ay和M取代。 我们在运算符上引入额外的指数x和y来表示坐标方向（分别为空间步数Ax和Ay），并将方案写成Prx1QixP-v\{R-1Sx + R\textasciicircum{}1Sy)u = 0，P\textasciicircum{}QlyP-viR\textasciicircum{}Sx+R\textasciicircum{}Sy\textasciicircum{} = 0，P\textasciicircum{}xQ\textasciicircum{}U + P\textasciicircum{}yQlyV = 0。 （9.20）请注意，这个用ATJ = 0表示的系统并不是定义我们边界值问题的正确解决方案的系统。 它只是一个在Krylov型迭代方法的每次迭代中计算的中间解。 在每次迭代中，我们需要计算A\textbackslash\{\}J\textasciicircum{}k\textbackslash\{\}，这反过来又需要求解一维系统r\textbackslash\{\}xl>x QixP\{k）等。 为了完成迭代，应用了真正的边界条件，例如（9.12），并得到了Tj（fc+1）。 请注意，二阶平均被四阶平均取代。 此外，求和公式被（9.5）中积分的四阶近似所取代。 为了说明方案的准确性，我们构建了一个测试问题。 我们在第二个方程中添加一个强迫函数，以获得一个简单的

Navier-Stokes 方程 173 NxM 10x10 20x20 40x40 eio/e20 e2o/e40 eio/e40 I \textbackslash\{\}\textbackslash\{\}u-u*\textbackslash\{\}\textbackslash\{\}h 2.3e-3 6.7e-5 2.3e-6 34.7 28.7 2-norm, ejv, \textbackslash\{\}\textbackslash\{\}v-v*\textbackslash\{\}\textbackslash\{\}h 2.3e-3 6.7e-5 2.3e-6 34.7 28.7 N=[10,20,40] I|P-P*IU 3.5e-3 2.2e-4 1.2e-5 16.0 18.9 表1 表1 v = 1.分析解的稳定斯托克斯方程的数值结果。 问题是Vx - Py xac V\textbackslash\{\}U>xx i \textasciicircum{}yy) - V(VXX + Vyy) UX + Vy t解是u* = v* = - p* = 2v = o, = - 4i\textasciicircum{}cos(a;) sin(y) = 0. sin(x) cos(y), cos(:r) sin(y), cos(x) cos(y)。 (9.21) (9.22) 计算域是fl = \{0 < x, y < 6\}，v = 1。 用迭代的GMRES求解器求解了代数方程组。 在表1中，我们显示了离散\textasciicircum{}-规范的计算误差，即\textbackslash\{\}w\textbackslash\{\}\textbackslash\{\}h =,£i 2AxAy。 （9.23）对于解决方案的每个组成部分。 我们可以看到，在速度分量和压力方面，我们实现了四阶精度（或更好）。 请注意，在非常粗糙的10 x 10网格上，精度已经是\textasciitilde{}10\textasciitilde{}3。 接下来，我们考虑与时间相关的斯托克斯方程和时间离散化（9.2）。 对应于（9.20）的系统是§u"+1 + At (P1-1Qiapn+1 - v\{R\textasciitilde{}xSx + R-1Sy)un+l) = 2un + At(P\{ylQiyPn+1-v\{R-1Sx + R-1Sy)vn+1) = 2vn P2-1Q2x"n+1 + P2yQ2yVn+1 = 0。 3„n+l 2U it,""1, (9.24)

174 GUSTAFSSON \& NILSSON 对于每个时间步骤，该系统都由如上所述的Krylov型方法求解。 请注意，真正的边界条件在Krylov方法的每次迭代时都会应用，而不仅仅是在每个时间步骤的完成时。 第一个测试问题在域£1 = \{0 < x,y < 6\}中得到了解决。 至于稳态方程，我们在方程中添加强迫函数以获得简单的分析解。 问题是（9.25）Ut+Px - V(lLXX +Uyy) Vt+Py- V(vxx + Vyy) UX +Vy，其确切解u* = v* = - sin (a;) cos(y) cos(t), = - COS(:E) sin(y) cos(t) -Au cos(i) sin(y) sm\{t) = o, sin(a;) cos(y) sin(t), - cos(x) sin(y) sin(i), (9.26) p* - 2;/cos(:r) cos(y) sin(f)。 时间依赖性斯托克斯方程的第二个测试问题是直通道SI中的流动-\{0 < x < 1，-1 < y < 1\}：Ut+Px \textasciitilde{} V(UXX + Uyy) - 0，Vt + Py \textasciitilde{} V(VXX + Vyy) = 0，(9.27) Ux+Vy - 0。 对于这个问题，我们通过求解超越方程，推导出一个分析解，包括ansatz u* = U(y) eax-\textasciicircum{}, v* = V(y) eax-ut, (9.28) p* = P\{y) eax-<Jt \textasciicircum{}和边界条件 U(-l) = f/(l) = V(-1) = V(l) = 0。 我们得到的数值v U(y) V(y) P(y) lues v = cx sin(ay) + \textasciicircum{}c2 sm(ny), = c\textbackslash\{\} cos(o:y) + 2c2 cos(«;y), = \textasciicircum{}c2 sin(ay), \textbackslash\{\}fv2 + a2vw K = = 1, a - 1 and to = 11.6347883720355431, d = 0.9229302839450678, c2 = 0.272128679072。

\section{9.4 Navier-Stokes 方程的四阶方法}
Navier-Stokes 方程 175 euw - \textbackslash\{\}\textbackslash\{\}uw -u*\textbackslash\{\}\textbackslash\{\} evw = 11\textasciicircum{}10-\textasciicircum{}*11 epio = HPIO -P*\textbackslash\{\}\textbackslash\{\} eu20 = \textbackslash\{\}\textbackslash\{\}u20 \textasciitilde{}u*\textbackslash\{\}\textbackslash\{\} e\textasciicircum{}20 = ||«20 \textasciitilde{}V*\textbackslash\{\}\textbackslash\{\} eP20 = ||P20-P*| eU40 = ||U40 - U*||| ev40 = \textbackslash\{\}\textbackslash\{\}v40 \textasciitilde{}v40 = ||P40-P*|| euw/eu2o ev2o epio/ep2o/ep2o/e«40 e«40 ev-2o/ev4o/ep2o eu40 evio/evi0 epio/epio 10 000 2.3e-4 2.3e-4 2.0e-3 6.4e-6 6.4e-6 7.7e-5 2.2e-7.2e-7 3.2e-6 35.4 28.5 28.5 22.1 40 000 8.9e-4 8.9e-4 2.7e-3 2.5e-5 2.5e-5 1.3e-4 8.9e-7 8.9e-7 6.0e-6 35.0 35.0 30.5 20.5 28.5 28.5 21.6 50 000 l.le-3 l.le-3 2.9e-3 3.1e-5 3.1e-5 1.5e-4 l.le-6 l.le-6 7.0e-6 35.0 35.0 19.6 28.5 28.5 20.7 表2 求解v - 1、At = l.Oe - 5和N - M = [10,20, 40]的时间依赖性斯托克斯方程时的测试问题（9.25）的数值结果。 测试问题（9.25）和（9.27）都根据四阶方案（9.24）离散化，用迭代GMRES求解器解决了代数方程组。 为了证明空间中精度的正确顺序，选择了一个小的时间步骤，At = l.Oe - 5。 表2和表3显示了错误。 我们可以看到，在两个测试问题中，速度分量和压力都比四阶精度更好。 9.4 Navier-Stokes 方程的四阶方法 在本节中，我们介绍了不可压缩时间依赖性Navier-Stokes 方程的一些初步结果。 边界条件的分析

176 GUSTAFSSON \& NILSSON ewio = \textbackslash\{\}\textbackslash\{\}u10 - u*\textbackslash\{\}\textbackslash\{\} ev10 = \textbackslash\{\}\textbackslash\{\}vw -v*\textbackslash\{\}\textbackslash\{\} epio = ||Pio-P*|| e«20 = 11\textasciicircum{}20 - u*ll eV2Q - 11\textasciicircum{}20 \textasciitilde{}V*\textbackslash\{\}\textbackslash\{\} eP20 = \textbackslash\{\}P*\textbackslash\{\}\textbackslash\{\} eu4o = || "40 - "*| et>40 = |w40 - w*||| eP40-P*||| euio/evw/eV2o epio/ep20 eu2o ev2o/ev2o/evA0 ep2o/eP40 euw/euw/evw/ep4o 10 000 8.5e-4 3.7e-3 2.5e-2 2.3e-5 1.8e-4 1.4e-6 5.9e-5 37.4 29.8 1 L.le-4 7.7e-4 7.0e-7 5.6e-6 2.6e-5 4.3e-8 3.0e-7 1.8e-6 37.4 20.2 29.8 16.3 18.5 14.3 50 000 8.1e-6 3.5e-5 2.4e-4 2.2e-7 1.7e-6 8.1e-6 1.3e-8 9.4e-8 5.6e-7 37.4 20.2 29.8 16.3 18.5 14.5 表3 测试问题的数值结果（9.27），求解v = 1，At - l.Oe - 5和N = M = [10，20，40]。是第9.2节中斯托克斯方程分析的直接推广，边界条件的形式完全相同。 四阶近似值也是一样的，但我们还需要两个对向条件的成分。 由于u和v不存储在同一点，因此需要四阶平均公式Eu和Ev。 此外，用于上述一阶导数的紧致方案必须修改，使其以网格点为中心。 我们使用

Navier-Stokes 方程 177 Pf = Qf的第四阶方法：Pf Qf = \&(/o + 2/\{),!( /;\_！ +4/;+ //+1)，i = i，...，N-i。 2i\textasciicircum{}(-5/o + 4/1 + /2), ihifi+i - fi-i), t = l,...,JV-l, 242\textasciicircum{}(-/JV-2 - 4/JV-I + 5/JV) • (9.29) 使用Px,QX,Py, Qy来指示坐标方向，我们得到系统（对应于（9.24）l\textasciicircum{}1 + At \{P\textasciicircum{}QixP\textasciicircum{}1 - v\{RzxSx + R-1Sy)un+1) = 2un - iu"""1 - 2At(unP-1Qxun + \{Evn)P-1Qxun-1 + \{Ev\textasciicircum{}Qyu"-1), «- - - + Atpn (P1-1Qiy+1 - i/\textasciicircum{}S* + R-Sy)vn+1) = 3„n+l 2 2 请注意，未知U“+1的协效矩阵与上述斯托克斯方程完全相同。 系统在每次迭代中都由真正的边界条件补充，完全如上所述。 最后一个测试问题展示了该方案能够为时间依赖的Navier-Stokes 方程产生四阶精确解决方案。 我们考虑方程Ut + uux + vuy + px - v\{uxx + uyy) = sin(:r) cos(a:) sin (t) + sin(a;) cos(y) cos(i), vt + uvx + Wy +py- v\{yx + vyy) = sin(y) cos(y) sin2(i) - 4k-cos(a:) sin(y) sin(i) - cos(a;) sin(y) cos(i), ux + vy = 0, (9.31)在域f2 = \{0 < x,y < 6\}中。 确切的解与时间依赖性斯托克斯方程（9.26）的第一个测试问题相同。 数值实验的结果如表4所示。 对于这个计算，雷诺数是Re - \textbackslash\{\}jv = 2000，时间步骤At = l.Oe - 5。 我们可以看到，我们取得了比四阶精度更好的成绩

\section{参考资料}
178 GUSTAFSSON \& NILSON euw = ||uio-u*|| evw = \textbackslash\{\}\textbackslash\{\}vw -v*\textbackslash\{\}\textbackslash\{\} epw = ||PIO-P*|| eu2o = | M20 \textasciitilde{}u*\textbackslash\{\}\textbackslash\{\} ev20 = \textbackslash\{\}\textbackslash\{\}v2o -v*\textbackslash\{\}\textbackslash\{\} eP20 = \textasciitilde{}P \textbackslash\{\}\textbackslash\{\} eit40 = ||l/40 - It* 11 ev40 = \textbackslash\{\}\textbackslash\{\}v4o -v*\textbackslash\{\}\textbackslash\{\} eP40 = ||P40 -P*\textbackslash\{\}\textbackslash\{\} -P*\textbackslash\{\}\textbackslash\{\} euw/ev20 epio/ep2o/eM40 ev2o/evio ep4o/eui0/eui0/ew4o/ew4o enXo ep4o ep4o 10 000 2.3e-4 1.8e-3 6.3e-6 6.3e-5 8.8e-4 1.6e-3 2.5e-5 2.5e-5 5.1e-5 8.6e-7 8.6e-7 2.3e-6 35.7 35.8 30.6 28.6 28.7 22.2 50 000 l.le-3 l.le-3 1.4e-3 3.1e-5 3.0e-5 4.7e-5 l.le-6 l.le-6 2.2e-6 35.6 35.8 30.3 28.5 28.6 21.6 表4 用雷诺数 Re - 1/v - 2000、At = l.Oe - 5和N = M = [10,20,40]求解时间依赖性Navier-斯托克斯方程时，问题（9.31）的数值结果。对于速度分量和压力。 至于之前的情况，我们注意到在非常粗糙的10 x 10网格上已经存在小错误。 致谢：Navier-Stokes结果是作为在曲线网格上直接模拟湍流的更大项目的一部分获得的。 该项目的其他参与者包括Arnim Briiger、Dan Henningsson、Arne Johansson、Wendy Kress和Per Lotstedt。 此外，Carl Adamsson、Stefan Engblom和Anders Goran也完成了一些编程工作。 参考资料1。 Fornberg，B.，k Ghrist，M.，波型方程的空间有限差分近似，SIAM J. 数字。 肛门。 37，1999年，pp. 105-130。

Navier-Stokes 方程 179 2的第四阶方法。 古斯塔夫森，B。 \& Nilsson，J.，边界条件和交错网格上稳定斯托克斯方程的估计，技术报告1999-015，乌普萨拉大学信息技术系，11月。 1999年。 提交发表在《计算机与流体》上。 3. 哈洛，F。 H。 \& Welch，J. E.，自由表面流体的时间依赖黏性不可压缩流动的数值计算，物理。 流体8，1965年，pp. 2182-2189。 4. 勒勒，S。 K.，具有光谱样分辨率的紧致有限差分方案，J. 补偿。 物理。 103，1992年，pp. 16-42。 5. 森西，Y.，伦德，T。 S.，瓦西里耶夫，O. V。 \& Moin，P.，不可压缩流动的完全保守高阶有限差分方案，J. 补偿。 物理。 143，1998年，pp. 90-124。 6. 陶，E。 Y.，稳斯托克斯方程的数值解，J. 计算机。 物理。 90，1992年，pp. 190-195。 7. Wesseling，P.，Segal，A. 和卡塞尔，C。 G。 M.，一般三维非平滑交错网格上的计算流，J. 补偿。 物理。 149，1999年，pp. 333-362。

\clearpage
\chapter{10个可扩展的非稳定不可压缩流的并行隐式多网格解决方案}
\section{10.1 摘要}
不稳定不可压缩流的可扩展并行隐式多网格解决方案R。 潘卡雅克山，L。 K。 泰勒，C。 盛，W。 R。 布里利，D。 L。 Whitfield1 10.1 摘要 提出了一种可扩展的并行迭代隐式多网格算法，用于包含旋转和移动组件的复杂非稳不可压缩黏性流，使用动态相对运动多块结构网格。 该算法结合了离散状态可变通量线性化、每个时间步骤的非线性多网格迭代，以及每个多网格级别的块-Jacobi LU/SGS方案引入的可扩展并发性。 为现有和假设计算平台上的并行CPU、内存和成本效率开发了半经验性能估计。 使用这些估计分析了可扩展性，结果以内存受限尺寸放大和常量问题尺寸放大模式的性能景观的形式给出。 包括MPI软件带宽和架构特定软件调整等参数的影响。 这些结果表明，该方法在实际意义上可以扩展到大规模问题。 还讨论了分级收敛速率和多算法变体，并给出了说明由十度舵偏转引起的潜艇机动的大规模模拟的综合结果。 密西西比州立大学ERC计算模拟和设计中心。MS 39762-9627。 Frontiers of 计算流体动力学 - 2002 编辑：David A. Caughey和Mohamed M. 哈菲兹 ©2002 世界科学

\section{10.2 介绍}
\section{10.3 基本非定常流动求解器}
182 PANKAJAKSHAN ET AL 10.2 介绍并行计算既可以缩短运行时间，又可以访问计算流体动力学中大型问题所需的大型但分布式全局内存。 可扩展的并行解算法已成为大规模问题和长期瞬态演变问题的实际解决的关键要素。 本研究开发并分析了一种可扩展的隐式多网格算法。 Brandt [1]广泛讨论了多网格方法，最近关于结构网格并行多网格流求解器的一些工作可以在[2-5]中找到。 三维Euler 方程的第一批平行多块多网格方案之一是Yadlin和Caughey [2]的方案，他们引入了水平和垂直多网格的概念。 在水平模式下，所有块接口边界都会在每个多网格级别后更新，在垂直模式下，在更新接口边界之前，在每个块内完成多网格周期。 在[2]中提出了一个异步垂直实现，其中接口边界使用相邻块的最新数据进行更新，并且该方案在最多八个处理器上成功演示。 这里使用的多网格方案是Sheng、Taylor和Whitfield的全近似方案（FAS）多网格方案的横向改编[6-7]，该方案以垂直模式实现，以减少单处理器代码的内存要求。 Pankajakshan和Briley [8]开发了一种没有多网格的可扩展并行算法，作为Taylor和Whitfield [9-10]提出的离散牛顿/放松（DNR）方案的并行适应。 在[8]中引入了可扩展的concurrency，使用Block Jacobi Lower/Upper Symmetric Gauss-Seidel（BJ-LU/SGS）松弛作为最内部迭代，以解决牛顿迭代。 本研究在每个时间步骤使用非线性多网格迭代周期，在每个多网格级别使用BJ-LU/SGS，并通过进一步研究并行性能和可扩展性来扩展以前的工作。 国防部挑战项目[11]关于移动控制面诱导的潜艇机动的最新结果说明了该方法对复杂非定常流动应用的能力。 10.3 基本非定常流动求解器 目前开发高效并行算法的策略是从有效的顺序算法开始，然后引入可扩展的并发修改，这些修改不会显著降低收敛速率或固有的串行算法的效率。 基本流动求解器是泰勒和惠特菲尔德的流动求解器[9-10，12]，由迭代隐式有限体积方案、Roe/MUSCL通量、数值计算的状态向量通量接近化和使用LU/SGS放松求解的近似牛顿迭代组成。

可扩展并行解决方案183 10.3.1 上风有限体积方案 三维非稳定不可压缩雷诺平均 Navier-Stokes 方程通过引入人工可压缩性 [13]来解决，以促进每个物理时间步骤的迭代解。 以单元格为中心的有限体积方案近似于时间依赖性曲线坐标系的人工可压缩性公式，可以写成dq/dr = -\textbackslash\{\}d.(f-fy) + dj(g-gv) + dk(h-hv)] = -R(q) (1)在这里，R(q)是稳定残差向量，q = J(p,u,v,w)是解向量，其中p是压力，u, v和w是笛卡尔速度分量，J是逆坐标变换的雅各布。 黏性通量向量是/、g、h，包括建模湍流应力的黏性通量分量是/v、gVy hv，x是时间。 中心差分运算符被定为d,•(•) = (•),•+1/2 \textasciitilde{}\textasciitilde{} ('),--1/2 f°r eacn '>./> \textasciicircum{}方向，对应于各自的曲线§, r]和£坐标方向。 人工可压缩性参数是/3 = 5。 详细的定义见[9]。 10.3.2 数值通量 每个细胞面的非黏性通量向量是使用Roe的[14]近似Riemann 求解器和van Leer的MUSCL对左右状态向量的推断，q R和q L获得的，如Anderson、Thomas和van Leer[15]的三阶无限形式实现的那样。 I方向的通量近似值可以写成/.+1 = /(+i(«t+.) + *"(<i. <。） '（<. - <.） (2) ' + 2 ' + 2 < + 2, + 2 ' + 2 ' + 2 + 2 在这里，A是对A \textasciitilde{}的罗氏平均评估，它由不黏性通量Jacobian A = df/dq的特征系统形成。 矩阵A”由A\textasciitilde{} = SA\textasciitilde{} S\textasciitilde{}l定义，其中S是一个矩阵，其列是A的右特征向量，A\textasciitilde{}是一个对角矩阵，包含A的负特征值作为非零元素。 他们和k方向通量的类似定义是通过将[i,f, A]替换为[/', g, B = dg/dq]或[k, h, C = dh/dq]，并使用相应的特征系统来定义B\textasciitilde{}和C\textasciitilde{}。 [9]中给出了仅使用一个方向的度量信息的细节和非奇异特征系统。 10.3.3 迭代隐式非定性算法 求解（1）的迭代隐式非线性方案由Aqn-'+l/Ar = - R(qn + Ul+l) (3)给出，其中A(-)" = (-)" + ' \textasciitilde{} (•)">和s = 0,1,...是迭代索引。 对于非线性空间残差，关于上一个迭代状态qn+Us的离散线性化以R(qn+hs+l) = R(q" + Us) + L"+h\textbackslash\{\}A sqn + hs) + O (As qn+Us f (4)的形式书写

184 PANKAJAKSHAN等人，其中As(-) = (-)I+1 \_ ('/>和£n + ls(-)是一个由通量导数组成的线性空间差分子，随后定义。 这导致了以下迭代线性化隐式方案：[zlr-/ + JL" + 1-( • )](J,g"+u) = Ru\{q"*1") (5) 其中物理非稳残差Ry被定义为Ru(q"+1) = [Ar-1Ip(q" + l-q") + R \{q+x )] (6)，通过在非稳定残差中用Ip = [0,1,1, l]r替换恒等式矩阵/。 然后，收敛解满足物理非稳定不可压缩近似Rv = 0，连续性方程中没有人工可压缩性时间导数。 10.3.4 数值状态-向量通量线性化 通量线性化矩阵是按数值计算的，正如Whitfield和Taylor提出的那样[9-10]。 作者发现，与更近似的通量雅各布相比，准确的线性化矩阵提供了更好的稳定性和迭代收敛速率。 Whitfield和Taylor [12]还提出了一种新的数值通量线性化，其中数值通量（即fj+v）与左右解变量状态向量qR和qL区分，而不是节点值qh和qi+1。 这种技术更经济，避免了是否在高阶通量中省略与qi+2相关的导数的问题，并且似乎在实际计算中表现良好。 这些数值状态向量通量线性化由*+ dfi+i/2 *" dfi-i/2 A定义[12]。 =, A. = (7a,b) d<7,L+1/2 dqf\_l/2，每个矩阵的第k列被评估为A i = [fi+./2 (9,+1/2 - <l\textbackslash\{\}+ x/2 + hek) - /,.+,/2 \{qf+1/2, q\textasciicircum{} 1/2) ]/h (8)，其中ek是\&th单位向量，h = \textasciicircum{}机器零。 类似的定义适用于B和C.. 线性化通量导数运算符<£"(•)现在可以定义为r(-) = -A\textasciicircum{}(-),.\_, + (i,+ -A\textasciitilde{},K-), +i;+1(-),。 +1 - BU • ),\_, + (Bj - B- )( • )j + «T+1( • ),+ 1 (9) -clxC)k\textasciicircum{} + \{tl -C\textasciitilde{}k)\{-k+l(-)t+l 10.3.5 使用LU/SGS松弛的线性化方案的解 Asqn+l的（5）解是通过下上对称高斯-塞德尔（LU/SGS）获得的。 引入第二个次级指数m

\section{10.4 可扩展的并行隐式多网格算法}
可扩展并行解185以及矩阵D和运算符L\{和L2的适当定义，LU/SGS方案可以写成[Dn + U' + ri + 1-'(-)](Aiq'+l-')' + ri + hs\{Asq" + x-s)m=Ru\{q" + x-s) [Dn + hs + ln2+h\textbackslash\{\}-) ](Asqn + x-sT+ x + r + hs(Asq" + hs )* = Ru(qn + U') LU/SGS方案的D、JL、JL和X2的定义是D(-) = [Ar-Il+(AJ-A\textasciitilde{} ) + (Bj-B: ) ) +(C\textasciitilde{} ),,,,.,., <£.,<£ 在LU/SGS迭代的指定次数后，更新为q" s+ «- qn+ s + (Asqn+'s)。 当\textasciicircum{}-迭代收敛时，量Asq"+,s接近零，给出了R v(q" + 1) = 0的解。 关于该算法的进一步讨论可以在[16]中找到。 10.4 可扩展并行隐式多网格算法空间区域分解是一种引入并发性的自然方法，该方法利用子域之间的粗/中粒度并行性，作为签名到多个处理器。 子域等同于多块网格的网格块，该网格块要么生成，要么重新分区，为指定数量的处理器提供负载均衡。 MPI消息传递方法因其在保持控制、通用性（特别是任意子域连接）和跨分布式和共享内存多处理器平台的可移植性方面的优势而使用。 基本流动求解器中顺序数据依赖性的主要来源是LU/SGS迭代算法，但使用多网格加速时会出现其他并行化问题。 建立了一个用于实现嵌套迭代的并发子迭代层次结构[17]，以探索并行性能问题，如高效的CPU利用率、减少通信开销、algorithmic 收敛速率以及稳定和时间准确解决方案背景下的可扩展性。 子迭代层次结构由多个FAS V-Cycle迭代、每个多网格级别的牛顿迭代和LU/SGS迭代组成，如[6-7]的基本多网格方案。 通过将子域接口上的块-Jacobi（BJ）更新与每个子域主内的LU/SGS迭代相结合，在最内部的亚比迭代级别引入了并发性。 虽然增加子域的数量需要更多的BJ-LU/SGS迭代，但成本的增加很小，而且收敛速率经常

186 PANKAJAKSHAN等人超过了通过按顺序更新接口或在更高的迭代级别获得的界面。 算法MGh（qh，Qh）[如果（iih是最好的网格）（计算/？ *，和jh\textasciicircum{}A±，B*，C使用LU/SGS初始化解决qh / = 0 While（Oh不是最粗的网格）\{ r" =/-«*,(«*) r2h = \%2hrh alh \_ alh h J™ = <\&?\}" Fh = R*(qf) + r2h 解决R2£\{q2h) = fh forq2h使用LU/SGS调用MG2h(q2h,Q2h)而(Qh不是最好的网格) (Aqh = q" - qho Smooth Aqh h Aqi h qi = = <\$\textbackslash\{\}Aqh h h qi + Aqi 对于所有子域Q\textbackslash\{\}呼叫MG\textbackslash\{\}qhk,Qhk) k k 1 1,2. 图1伪代码：并行非线性多网格算法图2并行亚迭代序列子doraain 1序列图3一个BJ-LU/SGS迭代的操作序列10.4.1非线性多网格周期非线性多网格方案由图中作为伪代码给出的算法表示。 1. 在这里，3lh是解向量的限制算子，而9kf是残差向量和通量线性化矩阵的限制算子。 同样，9\textasciicircum{}/2是将校正从粗网格插值到下一个更细网格的延长运算符。 每次给rou-的电话

图中的可扩展平行解决方案187齿。 1.代表一个V周期。 多网格限制、延长和隐式校正平滑运算符定义如[6-7]。 10.4.2 并行多网格子迭代序列 在实现FAS多网格方案时，牛顿迭代可以在每个多网格级别使用，[6-7，17]中讨论了此选项。 通过省略牛顿迭代并使用限制运算符将离散线性化矩阵传输到较低网格水平，可以构建一个更经济的算法。 图中显示了该分级序列和主要计算内核的示意图。 2. 通常使用三个多网格级别，非稳态流的典型算法参数是每个时间步骤有三个多网格周期，每个多网格级别有八个或更多BJ-LU/SGS迭代。 对于稳定流量，使用稳定残差R代替RUt，本地时间推进用于单个多网格循环和每个多网格级别五个或更多BJ-LU/SGS迭代。 10.4.3 BJ-LU/SGS子迭代 求解线性化方程的BJ-LU/SGS子迭代被实现为通过每个子域的下域和上（LU）扫描序列，由相邻子域接口上的解数据的消息传递交换相配。 对LU扫描和界面数据交换进行排序有三种可能性。 前两个变体涉及界面数据的双向交换，要么在每次LU扫描之后，要么在两个LU扫描都完成后。 第三种方法在每次下扫和上扫后执行单向数据传输，如图所示。 3. 这第三个变体结合了Aq的快速全球传播和最小的信息传递，通常是首选方法。 10.4.4 算法可扩展性：BJ-LU/SGS 收敛速率 总计算工作受到分级收敛速率的影响，因此不是网格大小的线性函数。 可扩展性直接受到BJ-LU/SGS方案的收敛速率的影响，子域数量增加，时间步骤不同。 这在[17]中进行了实证检查，使用（10）的线性解的结果由具有100,000个网格点的稳态解的测试案例来说明，使用本地时间推进和接近最佳的CFL = 45。 单个域（LU/SGS）需要两个顺序的剩余减少5次迭代，使用BJ-LU/SGS对200个子域逐渐增加到8次迭代，每个子域有500分。 为了比较，红-黑高斯-塞德尔方案，其（较慢的）收敛独立于子域的数量，需要14次迭代。 由于不需要携带最内侧的迭代来完成收敛，Rv = 0的非线性解也取决于添加子域时的BJ-LU/SGS迭代次数。 结果显示了将剩余值降低到六个BJ-LU获得的基线值所需的迭代次数/

\section{10.5 并行性能估计和可扩展性}
188 PANKAJAKSHAN ET AL SGS迭代和八个子域如图所示。 CFL=10和CFL=45为4。 测试案例是一个附加的船体，入射率为18.11度，有80万个点，使用三个多网格水平和一个多网格周期。 子域总数 子域总数 10 20 30 40 I 方向的子域 10 20 30 40 I 方向的子域 图 4 分解对 BJ-LU/SGS 迭代的影响 指定残差减少所需的 10.5 并行性能估计和可扩展性 可扩展性一词有许多定义和解释。 一个简单的定义（在最小化运行时很有用）是，可扩展性在有效Mflops和程序数量之间存在线性依赖关系。 实际CFD模拟的可扩展性目前目标是为大规模实际问题实现可接受的多处理器运行时，同时保持CPU和内存等计算资源的成本效益使用。 10.5.1 并行CPU、内存和成本效益 并行效率和可扩展性在这里使用基于特定流量案例和并行计算平台的输入参数对并行运行时的先验半经验估计进行分析。 每个时间步骤的总运行时间表示为TRuntime = TCPU + TComm，其中TCPU是花在浮点操作上的时间，TComm是花在处理器间通信上的总时间，包括消息传递和加载缓冲区数组。 还引入了闲置或未使用内存的概念，以表示分配给处理器但在执行期间没有实际使用的内存（即保留且其他用户不可用）。 未使用的内存可用于分配给共享内存系统的其他用户，但不能分配给分布式内存系统。 并行CPU效率t]CPU、通信开销yComra和内存效率t]Mem由\_TCPU\_TComm\_MbUsed V CPU J1 ' I Comm -J* » V Mem \textasciitilde{}Mh定义

可扩展的并行解决方案189最后，使用CPU \{CCPU）和内存（CMem）成本的粗略估计值作为并行计算机总成本的一小部分来定义成本效益t]成本。 在这里，CCPU-50\%和CMem-30\%。假定是当前计算机的典型，其余20\%归因于支持硬件。 以下总体成本效率Vcost = 1 - C-C/»(/(l - VcPUJ \textasciitilde{} \textasciicircum{}Mem ( 1 \textasciitilde{}\textasciitilde{} VMem)提供了计算资源利用率成本的总体衡量标准。 10.5.2 成本效率和运行时间 通常观察到，通过选择具有大点体积和小表面积的分解，导致通信/计算比小，以及通信与计算的异步重叠，可以减少消息传递开销。 根据目前的性能分析，分布式内存计算机的最佳成本效益是通过选择提供必要的全局内存所需的最小处理器数量（Pmin = MbProblem/MbProcessJ）获得的，因为这提供了100\%的内存效率，并最大限度地减少了消息传递开销。 然而，由于此运行时在用户的工作环境中可能不被接受，因此可以通过增加使用的处理器数量来减少运行时，尽管这本质上会增加消息传递，并降低分布式内存效率。 目前的性能分析可以帮助指导更快的运行时间和效率较低的计算资源利用率之间的权衡。 对于均匀处理器上的粗粒负载均衡，域应被划分为大小相等的块。 如果由于复杂的网格拓扑和/或非均匀处理器，负载均衡的分区无法适当调整大小，则可以通过用一些较小的子域超载选定的处理器来提高负载均衡效率。 性能分析还表明，内存小的处理器往往需要出色的通信速度才能有效地进行这些模拟。 10.5.3 CPU时间的估计 为了估计CPU时间，首先使用硬件性能监视器计算了四个主要算法组件中每个组件的浮点运算（OP）的数量。 这些在每个网格点的操作中给出如下：OPMaric = 481（网格度量评估）、OPRes = 1230（残差计算）、OPUn = 5284（数字通量线性化）和OPLU = 330（一个LU亚迭代）。 请注意，与剩余和Jacobian评估相比，BJ-LU/SGS迭代的成本很小。 最佳多网格级别的总CPU时间（以每个时间步骤为单位）由TCPU = WMG[OPMetHc + NFAS\{OPRes + Nw X OPw + OPLin/NLin)）给出-\textasciicircum{}-

190 PANKAJAKSHAN等，其中WMC是表示每个多网格迭代周期的多网格工作单元的因子[例而，（1 + 1/8 + 1/64 = 1.14）对于三个多网格级别]，通量线性化（可选）在每个NUn多网格迭代周期中更新。 表1定义了其他参数。 由于RCPU在实践中实现的价值高度依赖于计算机架构、编译器效率和个人编程，因此本分析使用目标架构的经验基准定时，并抑制消息传递。 有效的处理器速度可以通过缓存感知编程（即最大数据重用）来提高。 由于当前代码跨多处理器架构是MPI可移植的，因此它通常不会针对任何特定的机器进行调整，尽管它针对基于缓存的架构进行了通用优化。 目前的算法每百万个网格点需要1.26 Gb的内存，使用3个多网格级别和64位算术。 此内存主要由以下数组组成：通量（24\%）、解变量（18\%）、网格度量（22\%）、32位通量导数矩阵（34\%）和缓冲数组（2\%）。 该算法在FORTRAN-90和MPI中实现，具有单程序多数据（SPMD）结构。 表1估计性能参数p RCPU RBuf PMPI °MPI计算机参数处理器数量有效CPU速率（Mflops）有效缓冲率（Mb/s）MPI软件带宽（Mb/s）MPI软件延迟（s）1 + log2 P Npts NMG NLU Nm Nq算法参数网格点（百万）多网格迭代周期LU/SGS每个NMG迭代平均消息长度（Mb）每个子域的接口数量qn+1更新10.5.4通信时间估计处理器间通信估计是根据表1中定义的参数开发的，并考虑1）点对点消息交换数据 子域接口，2）消息传递缓冲阵列的加载和卸载，以及3）代数湍流量模型所需的全局还原操作。 由于解决方案流程在LU迭代级别同步，因此所有消息大约同时发送。 通过将二分割带宽估计为fp，即二维互连网状架构的确切值，包括消息带宽的自我争执的可能性。 每次迭代都涉及一个接口曲面中Aq的两次更新消息，每个时间步骤都需要在两个曲面和一组全局操作中更新q。 每一步的总通信时间估计为

可扩展并行解决方案191 = WMr,NMr.Nm 2NU, + Na ' MGly MGlym Gu + LmjP PMPI + R Buf + 4 Co* 8 PMPI 参数fiMPI和aMP，从现有文献中获得，以及缓冲率/? Bu/是从校准运行中的直接测量中获得的。 10.5.5 性能分析与直接测量的比较 估计运行时间已在[8,17]中与IBM-SP2上少数案例的实际运行时间进行了比较。 SP2有66Mhz POWER2宽节点，128或256Kb数据缓存，内存带宽从1.07Gb/s到2.13Gbytes/s。 在[8,17]中测量的有效处理器速度RCPU为49 Mflops，测量的大型消息的MPI软件延迟和带宽为oMP，=62fis和（5MP，=34Mb/s。 一些基于缓存的通用调整随后将SP2的RCPU提高到72 Mflops，相同的代码在Cray-T3E上以65-68 Mflops运行。 在图中。 5、从[17]中预测和实际运行时间对330万个点和500个时间步骤的样本案例进行了比较；良好的一致性在很大程度上反映了分析的半经验内容。 缓冲阵列的加载约占通信成本的40\%。 与计算重叠的异步消息通信将消息时间减少了25\%，但没有节省大量运行时间，因为线性求解器速度较慢，可能是因为非本地数据的访问增加。 TRunlime、TCPU和TComm的直接测量也进行常规监控，并给出了与本分析一致的T]CPU和yComm。 恒定问题大小缩放？ 10 o I Runti V 3.3 M 点，500 步 X. i 17 处理器 @ 256Mb \textbackslash\{\}。 •需要\textbackslash\{\}'全局内存\textasciicircum{}N实际运行时间\textasciicircum{}Rfc（32个处理器）IBM SP2 NX估计运行时间\textasciicircum{}O-»x线性加速\textasciicircum{}v J >• o S •r. £ ory E b 10 100处理器图5测试问题的估计和实际运行时间（330万点，500个时间步）图6恒定问题大小缩放的内存效率10.5.6可扩展性：MPI/SPMD软件实现可扩展性可以通过使用当前性能分析来研究在实际和 假设的计算平台。 可调的关键计算机参数-

192 PANKAJAKSHAN等人是RCPU、f\}MP„和RBuf；MPI延迟可以忽略不计，因为只有少量的大消息。 当前的可扩展性结果以在具有以下参数的通用计算机上运行的算法的性能景观的形式给出：RCPU=100Mflops，f\}MPI=130Mb/s，RBuf=30Mb/s，aMPi=15fis，每个处理器内存为512Mb。 这些参数的范围与Cray T3E和SGI Origin 2000类机器一致。 算法参数是用于时间加速度解决方案的典型参数：三个多网格级别，NMG = 3和NLU = 7，500个处理器增加到9。 性能景观包括a）内存受限尺寸的效率，其中问题大小增加，以在添加处理器时保持P = Pmin和rJMem \textasciitilde{} 100\%，以及b）具有不同网格点数量的常量问题大小缩放。 对于恒定的问题大小，通过添加超出最低要求的处理器来减少运行时，导致闲置或未使用的分布式内存。 图中显示的内存效率迅速下降说明了这一点。 6. CPU和成本效率如图所示。 7，这些也会随着问题规模的持续扩大而减少。 然而，除了极端情况外，CPU效率仍然高于线性加速的40\%。 恒定问题大小缩放1.0: o.8 r 0.2 恒定问题大小缩放Vv - 网格点（百万）vv\textbackslash\{\}r\textasciicircum{}\textasciitilde{}\textasciitilde{}-\textasciicircum{}r\textasciitilde{}\textasciitilde{}\textasciitilde{} Hcpu=100，P=130 512Mb处理器-\_\_\_™o ---\_\textasciicircum{}0 ----1° •"""""---3- \textasciitilde{}--\textasciicircum{}1网格点（百万）100 200 300处理器400 500 0.2 512Mb处理器100 200 300处理器400 500图7处理器速度和MPI带宽指定值的恒定问题大小缩放的CPU和成本效率如前所述，内存受限的大小保持最佳效率，图。 8表明，CPU和成本效率分别保持在80\%和90\%以上，最高可达1亿点（实线）。 对CPU速率和MPI带宽的灵敏度也如图所示。 8用于选定的值。 虽然效率对RCPU不太敏感，但任何特定架构的软件调整来增加RCPU都会直接减少运行时间，这满足TRumime <x（PRCPUVCPU）'1-请注意，提高RCPU会降低效率，因为TComm不会受到影响，除了可能的改善

\section{10.6 演示：舵诱导机动模拟}
可扩展的平行溶液193缓冲率RBuf。 总体而言，目前的分析表明，该方法在实际意义上是可扩展的，适用于大规模问题。 内存受限 Sizeup 内存受限 Sizeup 1.0 I ' 1 1 1 1 1 1.00 1。 1。 1 1 1。 1个网格点（百万）网格点（百万）图8处理器速度和MPI带宽的指定值的内存受限尺寸增加的CPU和成本效率10.5.7算法变体已经测试了几种算法变体，以节省内存或运行时。 通量导数矩阵占内存利用率的34\%，并提供了通过重新计算减少内存的选项。 如果对角矩阵D被计算并保存，如果每次LU/SGS迭代重新计算JL和JL2，那么内存将减少约30\%，随后的每次LU/SGS迭代将运行时间增加约16\%（SGI PCA/R10000）。 减少记忆的第二种方法是Frechet变体，其中矩阵-向量积直接使用（df/dq）Aq = \textbackslash\{\}f\{q + eAq）-f（q）]/e计算，而不形成线性化矩阵。 三次LU/SGS迭代的运行时间是均匀的，但十五次迭代增加了37\%。 另一个实现选项是使用符号数学程序（MAPLE 5.0）获得的通量导数矩阵的分析公式。 此选项在SGI PCA/10000上减少了22\%的运行时间，但在IBM SP2上增加了2\%的运行时间。 最后一个选项是预乘等式。 （10）由D\textasciitilde{}l，保存D\textasciitilde{}lLv D\textasciitilde{}lL2和D\textasciitilde{}lRv，以便在后续迭代中重复使用。 通过避免（4x4）系统的解决方案，第一次LU/SGS迭代的更高成本在后续迭代中摊销。 虽然收支平衡点取决于处理器，但在一个例子中，它是五次迭代，十五次迭代节省了10\%。 10.6 演示：使用k-e 湍流模型和1.2 X 107的雷诺数计算了舵诱导的机动模拟机动潜艇模拟[11]，用于完全配置

194 PANKAJAK8HAN ET AL SUBOFF/推进器配置，在两个帆平面和四个船尾平面表面具有逼真的控制表面间隙。 通过将流动求解器与车辆运动的6-DOF求解器耦合来确定车辆轨迹。 整合的黏性应力和压力分布为6-DOF方程提供了流体力和力矩，其解得出了车辆速度、旋转率和轨迹的时间历史。 多生物动态结构网格拓扑可以具有任意的非结构化模式，包括可能与多个相邻网格块连接的块表面。 可移动控制面由网格处理，网格在移动控制面附近的块中局部变形。 网格有450万个点，其子层分辨率使墙壁上的网格间距对应于表面上的y+ < 1。 网格运动技术的详细信息见[18]。 首先计算了恒定速度下的直线运动的启动解，给出了一个由转子引起的周期性运动的无瞬态解。 然后计算出由舵偏转10度诱导的自推进机动的解决方案，在大约四分之一的船体行程中作为时间上的线性运动施加，之后舵保持固定。 说明此解决方案的计算结果如图所示。 9. 启动解决方案显示了船尾附近的轴向速度轮廓，-操纵解决方案在四个时间点上显示，对应于3、9、20和30度的横向（水平）偏转。 这种假想的潜艇没有惯性矩等参数的现实设计值，因此，结果展示了预测机动的能力，但并不代表特定潜艇的实际机动。 •该解决方案需要每螺旋桨旋转350步，每船体长度需要7000步。 使用50个T3E/256Mb处理器的解决方案，每个船体长度需要160小时（8000个处理器小时）。 启动解决方案在大约2个船体长度后变得周期性。 图9示例；由舵运动引起的启动解决方案和潜艇机动

\section{10.7 致谢}
\section{参考资料}
可扩展的平行解决方案 195 10.7 致谢 这项工作由博士赞助。 L。 Patrick Purtell和博士 埃德温·P。 海军研究办公室的Rood，部分由美国宇航局艾姆斯研究中心，由博士监督。 罗杰·斯特劳恩。 它还得到了北极地区超级计算中心在国防部HPC挑战项目下授予的HPC时间的部分支持。 参考资料1。 Brandt，A.，多网格技术：1984年流体力学应用指南，CFD系列讲义：von-Karman流体力学研究所，比利时Rhode-Saint-Genese，1984年3月26日至30日。 2. 亚德林，Y.和D. A. Caughey，并行计算 Euler 方程的块状多网格隐式解决方案策略，AIAA杂志30（8），1992年，pp. 2032-2038。 3. Shreck，E.和Peric，M.，用平行多网格求解器计算流体流，Int。 J。 数字。 流体方法16，1993年，pp. 303-327。 4. 德加尼，A。 T.和Fox，G. C.，并行多网格方法在非定常流动中的应用：性能评估，并行计算流体动力学-使用并行计算机的实现和结果，Ed。 S。 泰勒，A。 埃瑟，J。 Periaux和N。 Satofuca，爱思唯尔科学，B。 V.，阿姆斯特丹，1996年。 5. 卢，J。 Z.和Ferraro，R.，带有并行多网格椭圆内核的并行不可压缩流动求解器包，J. 补偿。 物理学125（1），1996年，pp. 225-243。 6. 盛，C，泰勒，L。 K.和Whitfield，D. L.，三维不可压缩高雷诺数 湍流流的多网格算法，AIAA杂志33（11），1995年，pp. 2073-2079。 7. 盛，C，泰勒，L。 K.和Whitfield，D. L.，非稳定不可逆压欧拉和Navier-Stokes流计算的多网格算法，第六届计算流体动力学国际研讨会，1995年9月4日至8日，内华达州太浩湖。 8. 潘卡贾克山，R.和W. R。 Briley，使用MPI的多块结构网格上黏性不可压缩流动的并行解，并行计算流体动力学-使用并行计算机的实施和结果，Ed。 S。 泰勒，A。 埃瑟，J。 Periaux和N。 Satofuca，爱思唯尔科学，B。 V，阿姆斯特丹，1996年，pp. 601-608。 9. 泰勒，L。 K.和D. L。 Whitfield，用于静止和动态网格的非稳态三维不可压缩欧拉和纳维-斯托克斯求解器，AIAA论文编号 91-1650，1991年。 10. 惠特菲尔德，D。 L.和泰勒，L. K.，高解通量差分割方案的离散化牛顿松弛解，AIAA论文编号 91-1539，1991年。 11. Pankajakshan R.，泰勒，L. K.，liang，M.，Remotigue，M. G.，布里利，W. R.，惠特菲尔德，D. L.，控制表面诱导的海底机动的并行模拟，AIAA论文2000-0962，内华达州里诺，2000年。 12. 惠特菲尔德，D。 L.和泰勒，L. K.，二维时间不可压缩Euler 方程的数值解，MSSU-EIRS-ERC-93-14，1994年4月。 13. 乔林，A。 J.，解决不可压缩黏性流动问题的数值方法，J. 补偿。 物理2，1967年，pp. 12-26。 14. 罗，P。 L.，近似黎曼求解器、参数向量和差分方案，J。 补偿。 物理。 43，1981年，pp. 357-372。 15. 安德森，W。 K.，托马斯，J. L.和van Leer，B，Euler 方程的有限体积通量向量分割的比较，AIAA杂志24（9），1986年，pp. 1453-1460。 16. 布里利，W。 R.和McDonald，H.，隐性Navier-Stokes算法和近似因数分解的概述和推广，将出现在《计算机与流体》中，2000年。 17. Pankajakshan R.，使用多块结构网格的非稳态不可压缩黏性流动的平行解，博士论文，密西西比州立大学，1997年。 18. liang, M., Pankajakshan, R., Remotigue, M. G.，泰勒，L. K.，用于模拟潜艇机动的动态网格生成，第7 Int. 2000年9月，加拿大不列颠哥伦比亚省惠斯勒，计算场模拟中的网格生成会议。

\clearpage
\chapter{11 涡量约束在复杂体流动预测中的应用}
涡量约束在预测复杂体上流动中的应用J。 斯坦霍夫\textbackslash\{\}Y。 Wenren 2，C。 布劳恩3，L。 王4和M。 范5摘要描述了一种技术“涡量约束”，它代表了一种非常有效的、统一的方法，用于处理具有薄对流涡流的复杂、高雷诺数分离流，以及具有薄附着边界层的复杂固体体。 首先，描述了传统欧拉计算方法的缺点，以及涡量约束（也是欧拉计算方法）如何改善这些缺点。 然后回顾了涡量约束中的基本假设。 描述了该方法的一些细节。 在描述之后，提出了一系列结果：首先，提出了对流涡流和圆柱体上的柯西-黎曼流的二维结果。 这些描述了对流涡流和嵌入均匀笛卡尔网格中的固体表面流动方法的突出特征。 然后，展示了包括旋翼飞机在内的复杂物体上的流动的3D结果。 1 田纳西大学空间研究所教授，田纳西州图拉霍马，2研究科学家，流动分析公司，田纳西州图拉霍马，3研究助理，技术大学，亚琛，德国4研究助理，田纳西州图拉霍马，田纳西大学空间研究所5研究科学家，田纳西州图拉霍马，田纳西州前沿计算流体动力学 - 2002编辑：David A. Caughey和Mohamed M. 哈菲兹 ©2002 世界科学

\section{11.1 介绍}
198 STEINHOFF 11.1 导言强调我们的观点，大多数高雷诺数不可压缩流的特点是涡流结构，这些结构要么是固定的，如体符合边界层，要么是分离和对流的，如涡流片和长丝。 这些结构通常可以是湍流，通常由偏微分方程（pde）进行近似建模。 这种建模可以包括一个详细的结构，如涡黏性方法，零厚度不连续性，如固体表面的非黏性离散欧拉pde所暗示的，或由非黏性方程暗示的薄湍流涡流丝核心的粗略模型，由数值离散化误差引起的发散。 不幸的是，这些结构通常非常薄，由于解析问题，这些模型pde非常难以解决。 这些困难导致成本高昂的解决方案策略，涉及身体贴合网格，即使是非黏性处理，以及在身体表面附近进行广泛细化，以及在棚涡流片和长丝内进行广泛细化的自适应网格，如果要假定准确解决这些区域内模型pde的结构'1]。 一旦我们意识到这些pde只是这些涡流区域的近似模型，我们就会想到使用（非线性）差分方程直接在网格上对它们进行建模，而不是使用近似解析模型pde的有限差分方程。 这种方法允许我们将这些结构视为几乎奇异物体，在基本统一的笛卡尔计算网格上仅分布在几个网格单元上。 当然，这个想法是用于冲击捕获算法的，作为计算内部冲击结构的详细Navier Stokes解决方案的替代方案。 然而，与漩涡结构不同，冲击涉及指向内向的特征，使其更容易捕捉。 涡量约束方法，实现了这种涡流结构的方法！ 2'3'，已被证明是治疗身体符合边界层和脱落涡流片和长丝的一种特别有效的技术。 在当前的发展阶段，为这些涡流区域建模了非常简单的结构。 这些足以解决迄今为止所处理的问题，这些问题只涉及细涡流区域和与尖锐角的分离。 此外，目前正在开发初始应用，以处理从光滑表面分离的湍流边界层。 使用该方法，在粗（通常均匀的笛卡尔）网格上计算解。 尽管该方法完全是欧拉式，没有拉格朗日标记阵列，但即使它们只有几个直径的网格细胞，也可以无限期地对流，没有数字扩散。 因此，涡量约束可以控制拉格朗日“涡流跟踪”方案所拥有的涡流结构。 另一方面，就像在传统的尤拉技术中一样，该方法允许涡流片

\section{11.2 传统的欧拉方法}
涡量约束 199的应用是脱落或重新连接和涡流长丝合并或重新连接，无需逻辑操作或重新分配标记阵列，这在拉格朗日方案中通常需要。 本文首先概述了方法。 然后，使用涡量约束为具有上述特征的流动提供了一些代表性的结果，这些特征涉及对流涡流和2D的气缸流，以及3D流在椭圆体和复杂的旋翼上。 11.2 传统的欧拉方法 通常，偏微分方程（pde）首先被制定，用于表达流体的运动方程。 这些方程可以包括湍流区域的显式模型。 虽然以下几点可能看起来很微不足道，但重要的是要强调：如果我们考虑高雷诺数的薄涡流区域，无论是连接（边界层）还是分离（对流涡流），对于现实配置，目标始终是对这些区域进行建模。 当然，这是因为它们大部分可能是湍流，直接的Navier-Stokes计算，包括小的湍流刻度，是不可能的。 这种建模可以是显式的，例如，具有涡黏性模型的Navier-Stokes-like pde和可能还有其他相关的传输pde，或者它可以是隐式的--如果涡流区域被认为非常薄，以至于内部结构的细节对确定整体的“外部”非旋转流并不重要，“隐式建模”则意味着pde仍然是离散的（例如，inviscid Euler 方程），但计算的结果仅被视为涡流区域的粗略模型。 事实上，通常，目标是为这些区域使用<10个网格点。 由于数值黏性，结果是顺利的解决方案。 我们的主要观点是，无论是使用具有模型黏度的Euler 方程（一阶）还是具有模型黏度的Navier-Stokes样方程（二阶），结果仍然是高雷诺数的涡流区域的模型，（只要只用少量专门用于涡流截面的网格单元来计算现实的“大规模”问题）。 当然，拉格朗日“涡流晶格”方法\textasciicircum{}4'5'6]也构成了例如卷起涡流片的内部结构的模型。 此外，拉格朗日“涡流斑点”方法对内部结构进行建模\textasciicircum{}。 此外，使用分析方法的早期方法！ 8'9！，涉及将涡流片视为精确的接触不连续性。 这些方法也可以被认为是实际内部涡流结构的模型。 我们想提出的传统欧拉方法的主要特点是，尽管它们涉及高效的离散化和流场的非旋转或“外”部分的解，但它们也涉及薄涡流区域的低效模型：当使用pde（或一组pde）时

\section{11.3 涡量约束号}
200 STEINHOFF作为一个近似模型，必须在稀疏区域中离散化并准确求解，然后可能会出现巨大的计算困难。 表现包括必须使用非常密集的网格进行对流涡流，或需要大量计算时间才能“生长”和跟随涡流区域的精细自适应网格，以便可以解析离散化的pde！ 10]。 此外，对固体边界的传统处理要求表面的高精度：否则，数值误差，通常以数值生成的涡量的形式产生，并从表面扩散或对流，污染溶液。 对计算要求的影响是巨大的：要么必须使用表面符合性网格，对于复杂的几何形状来说很难生成，要么非符合性笛卡尔网格可以在边界处进行广泛的（非均匀的）细化。 另一个问题是，很难看到这些模型如何离散化并准确解决（即在网格独立的极限中），将流动与光滑表面分开，因为即使是最简单的情况：薄边界层近似中的层2-D Navier-Stokes 方程，已被证明非常难以解决'1 \textasciicircum{}此外，由于大型涡流结构在边界层中很重要，很难看出基于pde模型的基础涡黏性方法如何适当。 因此，由于计算和物理原因，当前近似雷诺平均湍流的pde模型可能不是一个好的方法，也许应该考虑不同类型的建模。 最后一点涉及涉及附近分离涡量的流动精度的一致性，其中表面附近的流动是分离的、对流的涡量以及绑定的、附着的涡量的函数：在传统方法中，结合和分离的涡量的处理方式不同（一个有符合的网格，一个没有）。 这种差异可以抵消在表面附近使用高精度符合格栅的任何好处，因为溶液的最终精度不能比附近的对流涡量更好，而涡量不涉及符合格栅。 11.3 涡量约束 11.3.1 - 要求 涡量约束的主要目标是仅使用横截面中的几个网格点对薄涡流区域进行建模，而不需要它们与网格对齐。 这似乎与这样一个事实一致，即我们没有精确的方程来解决这些涡流区域的结构，因此，只能对它们进行近似建模，此外，如果它们足够薄，可能不必准确建模。 此外，如果有很多漩涡区域，分配不止几个

即使使用自适应网格，将涡量约束 201网格单元应用于其横截面在计算上也可能不可行。 因此，我们的方法必须允许涡流在没有数值扩散的情况下进行长距离对流，但必须允许以受控的方式建模形状的变化。 它还应该允许复杂的固体表面，如薄涡流板，很容易地嵌入到均匀的笛卡尔网格中，而不需要生成符合体的网格。 涡量约束的最后一个要求对配方有很强的影响，它允许合并对流涡流和拓扑结构中的其他变化，例如在物体上对流和重新连接的能力，或者对于涡流片，分离和重新连接。 11.3.2 - 基本概念 涡量 Confinment背后的基本思想是在固定网格（通常是均匀笛卡尔）上开发一组满足上述要求的差分方程。 这意味着方程必须是欧拉pde在外部非旋转区域的准确离散化，但在几个网格单元中，流量变化为0（1）的涡流区域的一组差分方程。 这种涡流区域属性意味着方程不会是简单的pde的准确离散化。 相反，它们将是（非线性）差分方程，在网格节点上产生所需的模型值。 此外，为了允许分离、重新连接、合并等，无法指定涡流结构。 相反，结构必须放松到所需的轮廓。 上述要求意味着该方法中应该有两个基本参数：长度尺度和时间尺度。 这些与计算的网格单元格大小和时间步骤直接相关，即生成的涡流轮廓应为几个网格单元格宽，并且松弛应在少量时间步骤中进行。 当然，如果相关，可以实施更复杂的模型（例如，包括边界层动力学或长期黏性扩散），这将涉及更多的参数。 这些目前正在为与光滑表面分离的案例制定。 11.3.3 - 公式化 对于T035，涡量约束最简单的公式化涉及在原始变量公式中将两个项添加到离散动量守恒方程中，这些项类似于参考中讨论的对流短脉冲的扩散和非线性反扩散项。 [12]。 这些术语本质上是多维和伽利略不变的，只依赖于局部变量，并在涡流区域之外消失。 砷

202 STEINHOFF表示，这些项的方程的解没有指定，而是放松成一个结构，例如，涡流可以合并。 对于一般的非稳性不可压缩流，涡量约束项的控制方程是以下方程的离散化：V -q = 0 dtq = -(q- V)q + V(p/p) + [\textasciicircum{}2q- es\textbackslash\{\}，其中q是速度向量，p压力和p，密度，括号中的两个项是禁闭项。 两个数值系数e和H控制对流涡流区域或涡流边界层的大小及其向准稳定形状的松弛时间。 第二个禁闭期限有很多可能的形式。 最简单的似乎是s = n x w where。 \_ Vr/ n \textasciitilde{} m，涡量向量由Q = V x q给出。标量场，r]以两种方式定义：\textbackslash\{\}uj\textbackslash\{\}：场禁闭\textbackslash\{\}F\textbackslash\{\}：曲面禁闭涡量约束最简单的实现，用于对流涡流，称为“场禁闭”。 下一节将介绍边界层的简单修改“表面限制”。 对于场限制，单位向量n指向涡流区域的局部重心，限制项用于在扩散时将涡量对流回重心。 这种对流增加了扩散项，当两个项平衡时（对于fi和e的任何合理值），就会产生稳态分布。 关于配方的更多讨论可以在参考中找到。 [13,3,14]。 涡量约束方法的一个重要特点是，额外的项仅限于涡流区域：扩散项和限制项在这些区域之外都消失了。 另一个重要特征涉及质量校正引起的总变化，涡量和

涡量约束 203动量的应用，集成在对流涡流的横截面上。 Refs. [13,3,14]显示，质量和涡量是明确守恒的，动量几乎完全守恒。 参考[15]中描述的方法的小扩展，允许它明确地保存动量。 这对本文中描述的结果没有可观察到的影响，除非涉及弱背景速度场中强二维涡流的长期对流。 在这些中，使用动量守恒延伸来确保准确的轨迹。 一般来说，对于一系列值，计算的流并不敏感地依赖于参数e和/x。 因此，设置它们所涉及的问题与其他标准计算流体动力学方案中设置数值参数所涉及的问题相似，例如许多捕获冲击的传统可压缩求解器中的人工耗散。 这种缺乏灵敏度的原因是，例如，如果涡核接近轴对称，核心外的速度对涡量分布不敏感，只要半径保持小，并且由于数值效应而防止变大。 类似的考虑也适用于薄边界层。 （这与人工休克厚度效应相似，它取决于耗散参数）。 除了2D和3D自由对流涡流（固定形状的对流）的涡量分布的孤立波状特征外，两项研究，Refs. [3,14]证明了对流3D涡流丝的能力，最初以环的形式，可以融合和重新形成。 将这些结果\textasciicircum{}与在参考[14]中发表的实验测量结果进行比较，显示出非常接近的一致性。 这表明，通过扩散和非线性项的作用放松到准稳态涡流状态的基本计算概念自动允许逼真的涡流丝重新连接，同时防止因数值效应而扩散。 数值表明，尽管涡流区域只有几个细胞厚，但离散化方程的涡解在定性上与连续方程的预测解接近。 粗略地说，禁闭术语似乎将离散化误差对流到涡流中心。 这一点应该通过分析离散方程本身来解决，我们目前正在进行分析：必须强调的是，我们不是在边界层的涡核中准确求解连续pde方程，而是求解一组非线性离散方程，这些离散方程放松到仅限于几个网格单元的核心的所需结构（即使这些方程成为涡流区域外连续欧拉pde的精确近似）。 最后，应该提到，这些解决方案应该被视为“零阶”解决方案，这些解决方案非常经济，但没有考虑到涡流核心的动力学，如湍流效应。 这个想法

204 STEINHOFF是使用禁闭方法的扩展，以扰动方式包括此类效应（如果这些效应是显著的）。 通过这种方式，涡量约束可以被视为涡量动态计算的新型框架。 11.3.4 使用均匀笛卡尔网格的实体表面建模 涡量约束 最近在Refs. [16,12]和[17,18,19]中介绍了涡量约束在非符合正笛卡尔网格中表示固体表面的固定涡流板。 这代表了一种非常简单、经济的处理复杂物体的方法，因为它不需要身体符合或自适应网格生成，并且可以使用快速的笛卡尔网格设置和流动求解器。 此方法的步骤概述如下：•物体或自由曲面的几何形状以常规方式指定--例如通过曲面上一组点的坐标。 •从这组点中，在正笛卡尔计算网格的每个点上计算一个平滑函数。 该函数的值F(x)是网格点到定义表面的（有符号）距离。 因此，x值的“水平集”使F(x) - 0隐式定义了要求解流的表面。 •F-0表面的流量在一系列时间步骤中按时间准确计算。 11.3.5 计算细节 对于每个时间步骤（n），执行以下计算：步骤a：身体中的速度阻尼 速度，q"乘以F，\textbackslash\{\}(F)的函数，因此在F < 0时，它被减小。 这个因子在表面附近增加到1，在更远的距离上没有减少新速度：q' = X(F)qn 对于呈现的椭圆体，这代表了反射条件，而对于呈现的圆柱体和旋翼以及许多早期的结果，这个因子被简单地设置为零。 步骤b：对流 进行类似对流的计算来处理动量方程的一部分，就像传统的不可压缩“分裂速度”方法一样。 （见参考。 [20]）。 这是q" = q' -Atq'-Vq'的空间离散版本，步骤c：禁闭

涡量约束 205 涡量约束的应用用于计算速度增量，从而获得准稳定的薄涡流结构：q"> = q" + At(efhf xw-eui xw + /iV2g”) (11.1) 这里u是涡量和V|F| \textasciitilde{} \_ V|tU| 这是一个关键步骤：它将涡量向F = 0表面对角，并对流区域-没有它，涡量将继续扩散，有效地导致高黏度低雷诺数溶液，因为只使用大型常规网格单元，而不是像传统中非常薄的身体贴合或自适应精炼的细胞 计算。 （在早期的研究中，上述扩散没有明确添加，而是来自步骤b的离散化）。 在此步骤中，€f(F) - e，一个常数，接近F - 0，对于距离身体超过2个单元格的点来说，它变得小。 还有，eu-t-tf。 对于二维对流涡流结果，如上所述，使用了addes项的简单保守扩展，即动量被明确守恒。 步骤d：压力计算 计算压力时，时间步骤n + 1的速度是无发散的：V-qn+1=0。 这涉及求解泊松方程，就像在传统的“分割速度”程序中一样。 步骤f：速度更新 计算下一个时间步骤的速度，qn+1=q"" + Vcj>。 这与动量方程（At中的一阶）一致，其中<f>与压力有关：<t> = -AtP/p。

\section{11.4 当前结果}
206 STEINHOPF 11.3.6 收敛溶液的性质 在收敛时，V • <f和V x <fare的离散近似驱动到物体内部的0。 此外，q在体内被迫归零。 由于V • q = 0无处不在，因此没有流经身体表面，涡量仅限于身体表面的薄区域和与流动分离和对流的薄区域。 这些区域是2-3个网格单元厚，与整个问题（在精细网格限制中）中的网格单元大小或网格点数无关。 此外，由于计算速度的发散为零，如果分离位置准确，计算的速度场在远离对流涡流和物体表面的网格单元大小上至少将是一阶准确。 一阶误差是由于计算的涡流区域的厚度：通过将误差视为位移厚度效应，可以在方法的框架内制定简单的扰动校正，使解达到二阶精度。 在许多情况下，这个过程应该比诉诸于生成符合网格或自适应高阶细化的身体更简单。 然而，对于带有分离的钝体流，一阶精度应该与分离涡量的计算精度更一致。 11.4 当前结果 11.4.1 对流涡流 作为对涡量约束准确有效地对流集中涡流的能力的基本测试，我们在2-D中测试了两种情况的方法：在均匀的外部场中对流的单个涡流和一对在相互速度场中对流强度相反的涡流。 这两种情况都是强有力的测试，因为每个涡流的转换速度都比围绕涡流（在核心边缘）的最大流动速度低得多。 因此，这种“自诱导”速度的数值效应必须非常准确地抵消，才能获得准确的平移速度。 在这两种情况下，都使用了128 x 128单元格。 此外，外部边界条件是根据涡流的计算重心设置的。 II.4. I.I 单涡旋 在第一种情况下（未显示），单个涡旋在均匀速度场中移动：Uoo = 0.03，v\textasciicircum{} = 0.04（由涡核边缘的最大速度归一化）。 基于相同速度的CFL数字为0.4。 经过5000个时间步数后，涡核的大小与最初相同--1/4最大涡量的轮廓（保持不变）保持了\textasciitilde{}3网格单元的直径。 在计算过程中，漩涡应该传播60

在x方向应用涡量约束 207个单元格，在y方向应用80个单元格（总行程为100个单元格）。 在计算中，等高线图中最大涡量的点距离预测位置小于1个单元格。 这是意料之中的，因为可以很容易地证明，如果该方法保持动量（它确实如此），重心将以正确的速度移动。 11.4.1.2 涡对 在第二种情况下，两个涡流处于45°，中心最初相距10个网格单元（见图。 1，其中绘制了涡量轮廓）。 CFL数字为0.2，再次使用了5000个时间步骤。 根据它们的分离（没有外部流动），它们在x和y方向上预测的总平移量为70.7个网格单元格。 实际的转换，同样基于最大涡量的点，在预测的网格单元格内。 此外，如图所示，它们的分离保持不变。 此外，如图所示，外部（1/4最大涡量）轮廓的直径也为\textasciitilde{}3个网格单元。 重要的是要认识到，等高线的不对称外观是等高线绘图软件的虚构，因为中心不在细胞节点或细胞中心，涡量在1\textasciitilde{}2个单元中变化0（1）。 在中心靠近细胞中心或节点的时间步骤中，轮廓围绕涡流连接线和移动线对称，接近轴对称。 应该提到，如果实际的涡流完全服从非黏性Euler 方程，它们会因为相互感应而变形。 正如已经强调的那样，涡量约束方法旨在表示小型子网格尺度涡流的影响，这些涡流应该像点涡流（在2-D中）一样移动。 显示的涡量轮廓可以被认为是在网格上绘制的这些小核心的传播图像。 这类似于大型涡流模拟的理念。 11.4.2 机身流量的计算 下一个目标是量化用于在笛卡尔网格中处理身体表面的模型：特别是准确预测表面压力分布的能力。 首先，将预测的流量与（简化的）二维圆形圆柱体的确切解决方案进行比较。 然后，使用复杂网格和湍流模型以及风洞数据将椭圆体各种流站的预测表面压力与其他计算结果进行比较。

208 STEINHOFF ft 10 • • -10 0 0 -30 -50 -60 (a) 全网格视图 (b) 特写视图 图 1 涡量 翻译涡流的轮廓 11.4.2.1 二维圆柱体 首先，对流动方程的基本研究结果：动量和质量守恒，没有对流项，呈现在二维圆柱体上的流动。 这本质上是一个具有附着流的柯西-黎曼解，其确切的解是已知的。 这项研究很重要，因为它验证了将约束涡量嵌入正笛卡尔网格中的基本方法。 在这些涉及均匀笛卡尔网格的情况下，气缸是冲动地从休息开始的。 在收敛时，速度场q的发散处处为零，涡量除表面附近的窄带外为零。 此外，圆柱体内的q为零。 计算解决方案的简化如图所示。 2用于气缸直径的三种不同网格分辨率：分别为16、32和64个单元格。 可以清楚地看到涡流层的（数字）位移厚度效应。 计算的表面速度（通过伯努利关系与压力有关），以及精确的解，如图所示。 3作为围绕圆柱体中心角度的函数。 对于这些结果，使用了下文讨论的从涡流层外到表面的外推。 可以看出，计算结果非常准确，即使气缸直径有16个网格单元格。 由于涡量集中在圆柱体表面（F - 0），因此速度在那里迅速变化。 由于我们使用伯努利关系来计算压力，我们必须有一个方法，从流动平稳变化的涡流区域外推断速度（或计算压力）到表面。 此外推使用涡量约束，并且是-,,, / //.'/'//,,,,, lir/\% Kit is,, /...,, \textbackslash\{\}ii"i bd

涡量约束 209（a）16个单元格（b）32个单元格（c）64个单元格的应用图2参考中描述的各种网格分辨率的圆形圆柱体的计算流线。 [16,21]。 然后，速度可以插值到一组位于实际身体表面的点和计算出的压力。 可以实施高阶插值和外推来提高表面压力计算的准确性。 尽管如此，可以看出，二维圆形圆柱体的预测结果非常准确，即使圆柱体直径有16个网格单元。 11.4-2.2 椭圆流的计算 该方法适用于高攻角下6:1椭圆的三维流的计算。 这个问题以与光滑表面和大涡流区域的分离为主，因此它是该方法的一个很好的测试案例。 通过将该方案的计算结果与现有测量数据'22'和传统CFD方案的结果进行比较，进行定量分析，该方案基于雷诺平均Navier-Stokes 方程的有限体积离散和k-e 湍流模型'23'。 对五个攻击角度进行了计算：10°、15°、20°、25°和30°。 •计算网格图。 4显示了本研究中使用的计算网格。 只使用一个规则、均匀的笛卡尔网格，在表面附近没有细化。 不同方向的网格单元格数量为：imax = 188，jmax = 70，kmax = 100。 由于计算资源有限，必须严格限制细胞的最大数量。 椭圆体的长度相当于120个网格单元格，直径等于20个网格单元格。 这意味着，网格只向前延伸四分之一的身体长度和四分之一的身体长度

210 STEINHOFF i.J - 2.0 - 1.5 - 1.0 • 0.5 - 0.0 < -0.5 - -1.0 - 1.5 - -2.0 - 2.5 - • Exact \textasciicircum{}F \textasciicircum{}L -16 cells \textasciicircum{}r \textasciicircum{}c 32 cells f \textbackslash\{\} -64 cells f \textbackslash\{\} i i i i i 60 120 180 240 circenterentiallocation, 6 300 360 图3 Surface速度预测与身体下游二维圆形圆柱体的精确解的比较。 到左右边界的距离是一个直径，到上下边界的距离是2.5个直径。 使用涡量约束发现的结果和比较数据之间的一些差异可能是由于配置和外部边界之间的距离较小。 这个网格可以与传统的k-e 湍流模型计算中使用的网格进行比较！ 23]，如图所示。 5. 涡量约束参数 调查问题的雷诺数决定了禁闭参数e和数值扩散参数fin的选择。 到目前为止，考虑到使用符合网格，对e和鳍和雷诺数之间的关系进行了很少的调查。 在这项研究中，所有计算都只使用了一组特定的参数。 参数是通过尝试确定的，直到-25°的压力分布显示出最佳一致性。 只需要进行三次计算就能找到合适的集合。 使用了以下一组参数：e = 0.065，\textbackslash\{\}xn = 0.0025。 时间步骤

应用涡量约束 211 300个时间步骤，时间增量为At = 0.4。 随着这个时间的增加，每个时间步骤都遵守了CFL条件。 这个时间步数足以大致达到一个融合解，已透过再进行150个时间步进行检查。 与300步后的解决方案相比，450个时间步数后的解决方案没有显著差异。 •表面处理 表面边界算法被证明对计算结果的质量有影响。 在这里，我们将外部值外推到正文内的节点，并反转符号，以强制执行反射边界条件。 这导致对压力外推和分离的敏感性降低。 •结果的一般属性 在本节中，介绍了使用涡量约束计算的一些结果，以显示流动问题的物理属性在多大程度上得到了定量解决。 在图中。 6，配置显示为a = 25°的lee侧交叉分离。 自由涡流由恒定纵向涡量的等面表示。 它由一对纵向涡流组成，这些涡流与物体后端的涡流片相连。 需要注意的是，只能看到涡量的一个值的等面。 这就是为什么自由涡流似乎从远离身体表面的特定点开始的原因。 图。 7显示了椭圆体的三维视图，包括x/L = 0.25、x/L = 0.33、x/L = 0.58和x/L - 0.8的次级流线型图案，这是使用涡量约束计算的。 可以看出，随着纵向位置的增加，流线的分离更早发生。 流线型卷起到主 涡核 中。 这符合流场的物理行为。 •流程模式图。 8 - 图。 10显示了三个不同的粒子痕迹视图，分别使用涡量约束和FV-方法计算，表明在a = 30°时交叉流分离。 在迎风一侧，形成了一个附着的三维边界层。 在背侧，由于不利的周长压力梯度，水流与身体分离，并卷成一对纵向对称涡流。 用本方法发现的流线型模式与蔡英文和惠特尼发现的模式相当好。 在配置的后端附近可以看到差异。 这种差异是由于这样一个事实，即使用FV方法的解决方案是在网格上发现的，其边界呈球形，图。 5. 由于可能只有内网格的网格点被用于数据可视化，流线被内网格的外部边界切断。 分离和重新连接区域也如图所示。 11，它显示了在x/L = 0.77计算的交叉流流模式

212 STEINHOFF 30°攻角。 图中的FV-方法解决方案。 11显示了初级分离和二级分离。 初级涡旋的涡核位于局部坐标中大约（0.04,0.125）的位置。 涡量约束 解决方案无法解决二次分离问题。 其中一个重要原因可能是网格单元格大小。 二次分离的横截面大约对应于笛卡尔网格的3个网格单元格。 在涡流区域的位置和形状方面可以观察到进一步的差异。 与FV-方法结果相比，涡核的位置是（0.04,0.16），其形状是拉伸的。 同样值得注意的是，在分离点流线上表现出喷发样的行为。 所述属性可能是由于数值边界层比物理边界层厚得多这一事实造成的。 这可能会特别影响涡核的位置。 涡核的位置也可以受到以下事实的影响：施加的自由流速度的流入边界非常接近配置。 虽然并非所有细节都可以解决，但涡量约束解决方案清楚地代表了流程最重要的特征。 •表面压力 在笛卡尔网格上使用涡量约束方法，表面点的压力只能根据周围网格点的值来计算。 由于压力场是为计算域中的所有网格点计算的，无论它们是在配置内部还是外部，表面上都不会发生不连续性，分布是平滑的。 因此，无论是从固体外的点推断压力，还是使用内部和外部点内推断压力，都无法观察到差异。 在这项研究中，首先计算了表面被调查点的笛卡尔坐标，然后根据这些信息，确定了周围的网格点。 然后，应用线性插值来计算固体上点的压力。 图。 12 - 图。 16 展示了用所提出的方案获得的计算表面压力的比较，与FV方法的实验数据和结果，轴向位置的攻角为10°到30°，77\%x/L =椭圆长度的0.77。<j> = 0°和<j> = 180°分别表示迎风侧和背侧的点。 分离主导区域从<j>-90°到<j>-180°。 FV-方法结果显示，对更高的攻角有很好的一致性。 例外情况是160°方位角=20°的攻角和30°攻角的近0°方位角。 对于20°和25°的攻角，可以看到涡量约束结果和实验数据之间的最佳一致性。 对于任何攻角，大约155°方位角的局部最小压力都无法解析。 在较低的入射率下，特别是对于a = 15°，压力分布略微波纹。 这种现象的一部分可能归因于

\section{11.5 结论}
涡量约束 213的应用，将压力线性插值从周围的网格点到表面的调查点。 对于涡量约束结果和数据之间的a = 30°差异增加，也发生在边界层附着的迎风区域（0 = 0° - cf> - 90°）。 同样，流动的主导特性，即导致分离的不利周向压力梯度，已经正确预测，并且与实验数据得到了相当一致的一致性。 11.4.3 向前飞行中完整的阿帕奇直升机的计算 参考[21]中描述的计算靠近叶片的流量的过度程序用于计算前向飞行中阿帕奇直升机的流量，其中主流场计算（包括机身）使用统一的笛卡尔网格。 这些计算是针对带有四叶片转子的Apache-A直升机机身进行的。 身体表面的处理和之前描述的一样。 提出了两个计算案例，上升和水平向前飞行。 参考[16]也提出了水准前方飞行解决方案，以说明基本方法的涡流合并属性。 内叶片安装的C网格有120 x 26 x 24个单元格，背景笛卡尔网格有162 x 108 x 54个单元格。 在CRAY YMP上计算需要2.5小时/次。 尖端涡流通过图中涡量大小的等表面来可视化。 17\textasciitilde{}18。 在这些图片中，涡量等级的等高线水平设置为最大涡量等级的30\%。 可以在www.flowanalysis.com上观看一部展示尖端涡流发展的短片。 11.5 结论 本文提出的方法基于涡量约束，已被证明可以解决两个重要的不可压缩流动问题：计算长时间的集中涡流，没有数值扩散，以及计算在没有物体符合网格或长时间设置的复杂物体上的流动。 两种计算都是完全欧拉式的，只需要一个粗略的、规则的计算网格。 新方法涉及在计算晶格上生成孤立的波状配置，以表示作为非线性差分方程的直接解的涡流区域，而不是对这些区域的偏微分方程模型的传统离散有限差分近似的解。 因此，新方法为高效计算大量流量提供了一个新的、高效的框架。

\section{参考资料}
214 STEINHOFF 当然，与传统方法一样，在验证和校准湍流区域的雷诺平均表示模型方面仍有许多工作要做。 我们目前正在开发此类模型，这些模型与传统的pde模型截然不同。 致谢 我们要感谢莫菲特机场陆军航空动力学局的弗兰克·卡拉多纳就旋翼飞机空气动力学问题提供的广泛建议。 方法的开发和这里介绍的工作得到了陆军研究办公室、陆军航空飞行动力学局、美国宇航局和陆军SBIR赠款以及田纳西大学空间研究所的部分支持。 参考资料1。 J.U. 艾哈迈德和R.C. 斯特劳恩，“悬停转子和Overset-Grid Navier-Stokes求解器的计算”，会议记录，美国直升机协会第55届年度论坛，1999年5月。 2. J。 斯坦霍夫，H。 Senge和Y。 Wenren，“计算涡旋捕获”，UTSI预印本，1990年。 3. J。 Steinhoff和D。 Underhill，“涡量约束应用的Euler 方程修改，用于计算相互作用涡旋环”，流体物理学，6，1994年。 4. 罗森黑德，Proc。 罗伊。 Soc，A，134，1931年，第170页。 5. 韦斯特沃特，报告和备忘录，1692年，1936年。 6. R。 克拉斯尼，“涡流片计算：卷起、唤醒、分离”，应用数学讲座，1991年28日。 7. A. 伦纳德，“流动模拟的涡流方法”，J. 计算机。 物理，37，289（1980）。 8. H. von Helmholtz，《关于流体的不连续运动》，Phil。 Mag.，Ser. 4，36，11月 1868年。 9. G。 Kirchhoff，“论自由流体喷射的理论”，Jour，fur die Reine und Angew。 数学，70，1869年。 10. M。 丁达，A.Z. Lemnios，M.S. 谢法德，J.E. Flaherty和K。 Jansen，“旋流空气动力学的自适应解程序”，AIAA论文98-2417。 11. L。 L。 Van Dommelen和S. F。 沈，“分离层边界层中奇点的自发生成”，J. 计算机。 物理，38，1980年。 12. J。 斯坦霍夫，E。 普斯卡斯，S。 巴布，Y。 Wenren和D。 Underhill，“使用孤波在长距离上计算薄特征”，会议记录，AIAA第13届计算流体动力学会议，1997年7月。 13. J。 Steinhoff，Clin Wang，D。 安德希尔，T。 Mersch和Y。 Wenren，“计算涡量约束：涡流主导流的非扩散欧拉方法”，UTSI预印本，1992年。

涡量约束 215 14的应用。 J。 Steinhoff，“涡量约束：计算涡流主导流动的新技术”，计算流体动力学的前沿（D. Caughey和M。 Hafez，编辑），John Wiley \& Sons，1994年。 15. S.V. 佩夫钦，J。 斯坦霍夫和B。 Grossman，“使用不连续性禁闭程序捕获接触不连续性和激波”，AIAA论文97-0874，1997年。 16. J。 斯坦霍夫，Y。 Wenren和L。 王，“使用涡量约束分离高雷诺数不可压缩流的高效计算”，会议记录，第14届AIAA CFD会议，弗吉尼亚州诺福克，1999年6月。 17. J。 Steinhoff等人，“涡量约束：近期结果调查”，会议记录，第一届ARO年度涡量约束和相关方法研讨会，UTSI，1996年5月。 18. J。 Steinhoff等人，会议记录，第二届ARO年度涡量约束和相关方法研讨会，UTSI，1997年5月。 19. J。 Steinhoff等人，会议记录，第三届ARO年度涡量约束和相关方法研讨会，UTSI，11月。 1998年。 20. J。 Kim和P。 Moin，“分数步法对不可压缩Navier-Stokes 方程的应用”，计算物理学杂志，59，1995年。 21. J。 斯坦霍夫F，Y。 文仁，L。 王，M。 范，M。 肖和C。 Braun，“涡量约束在预测直升机转子和复杂体的尾波中的应用”，将发表在计算流体动力学杂志上，2000年。 22. C.J. 切斯纳卡和R.L. 辛普森，“三维交叉分离的详细调查”，AIAA杂志，卷。 35，1997年6月6日。 23. C.Y. 蔡英文和A.K. 惠特尼，“6:1椭圆体的三维流动分离的数值研究”，AIAA论文99-0172，1999年。

216 STEINHOFF ••Jfc m\textasciicircum{}ammsm m w HgE m 55 图4 嵌入式椭球体的计算网格的侧视图（顶部）和前视图（底部）

涡量约束 217的应用图5蔡英文和惠特尼使用的多块有限体积法的计算网格

218 STEINHOFF 图6 使用涡量约束计算的a = 25°配置的lee侧分离涡流。 箭头表示传入的流量。 图7椭圆体的三维视图，在x/L = 0.25，x/L = 0.33，x/L = 0.58和x/L - 0.8时为a = 30°（涡量约束）。 箭头表示传入的流量。

应用OP 涡量约束 219侧视图图8在30°攻角下流线型流线型流动模式的侧视图。 顶部：蔡和惠特尼，1999年。 底部：涡量约束

STEINHOFF图§FYont ¥iew在30°攻角下低6:1椭球体的流线型模式。 左：蔡英文和惠特尼，1999年。 右：涡量约束。 \textbackslash\{\} 图10 30°攻角的6:1椭圆体流线型流动模式的后视图。 顶部：蔡和惠特尼，1999年。 底部：涡量约束

涡量约束 221 0.15 0.05 图11 x/L - 0.77 a = 30°的交叉流线型的应用：左：蔡英文和惠特尼，1999年。 SI和S2表示初级和次级分离的位置，Rl和R2表示初级和二级重新连接的位置。右：涡量约束 0.6 -i 0.4 - 0.2 - o -0.4 -0.6 A实验数据（Chesnakas和Simpson，19 FV-方法（Tsai和惠特尼，1999年）涡量约束 a = 10° 30 60 90 -1-120 150 0 [度] 图12 压力系数数据在x/L = 0.77和a - 10°

222 STEINHOFF 0.6 0.4 0.2 - o -0.4 -0.6 A 实验数据（Chesnakas和Simpson，1998年）FV-方法（Tsaiand Whitney，1999年）涡量约束 /AAA a=15° 30 60 90 120 0 [cleg]图13 压力系数数据在x/L = 0.77和a - 15° 0.6 -i A实验数据（Chesnakas和Simpson，1998年）FV-方法（Tsai和Whitney，1999年）涡量约束 O \&AA£g\textasciicircum{}T\textasciitilde{}\textasciicircum{}\textasciicircum{} AAA 2\textasciicircum{}A <t> [cleg]图14 压力系数数据在x/L = 0.77和a = 20°

涡量约束 223 0.6 -i 0.4 - 0.2 O -0.2 --0.4 -0.6 A实验数据（Chesnakas和Simpson，19 FV-Method（Tsai和Whitney，1999年）涡量约束 a = 25° -r- 30 60 90 -1- -1- <t> [cleg]图15 压力系数数据在x/L - 0.77和a - 25° 0.6 0.4 - eT o - --0.6 A实验数据（Chesnakaspson，1998年）FV-方法（Tsai和Whitney，1999年）涡量约束 a = 30° -1-60 90 120 150 <t> [cleg]图16 压力系数数据x/L = 0.77和a - 30°

224 STEINHOFF J@\$\$"。 |£|?:*I|:S\textasciicircum{}\textasciicircum{} Jilw\textasciicircum{}M' §11?' ':':':':::\textasciicircum{}:\&¥x "\textasciicircum{}\textasciicircum{}s\textasciicircum{}。 J •'l;4.v. 图17 阿帕奇直升机的计算涡量等表面，包括主旋翼叶片（高前进比）（顶部：顶视图；底部；侧视图）

VORT1C1TY限制225 Jf JF' W Jf 0 J|f \#'的应用。 \# \#' M' \textasciitilde{} J§ 如果.. 如果。 如果II。' \&C\textasciicircum{} "i n \textasciicircum{}J\textasciicircum{} pi ii..ii it "1 11:\#P||, IJ\textasciicircum{}JWW\textasciicircum{}JF' '1 1. j\#pr»,\#\textasciicircum{}iiii\textasciicircum{}\textasciicircum{}i\textasciicircum{}\textasciicircum{}i£*\textasciicircum{} f C if.1 H|Mlt! N\# \% l\textbackslash\{\} WW** *！ “*\%* '\%* '\%,:i\%.'\%- - '\textasciicircum{}ik *\%*, Ii* v •\%.” Ik f\textasciicircum{}h\$\textasciicircum{}V*09\textasciicircum{} BR9V Ji fW JW JiF。 Jf J 8*.. \textasciicircum{}j\textasciicircum{}J\textasciicircum{}\textasciicircum{}\textasciicircum{}S "*"**""""*"•*' "Wfth'i\textasciicircum{} """df"""*" •»' \textasciicircum{} jp*f\textasciicircum{} \% 1 j '11;:l 1 111 f\textbackslash\{\} \textbackslash\{\} 11II '.;;.! Iff 1:: 1 图 18 阿帕奇直升机的计算涡量等曲面，包括主旋翼叶片（低前进比）（顶部：顶部视图；底部：侧面视图）

\clearpage
\chapter{12 不可压缩流动的晶格玻尔兹曼模拟}
\section{12.1 简介}
不可压缩流动的晶格玻尔兹曼模拟N。 里福，M。 Ishikura1 12.1 导言 CFD有两种方法。 传统的是根据连续体假设求解Navier-Stokes 方程。 在不可压缩计算的情况下，必须将动量方程与泊松方程一起求解，以满足无发散条件。 第二种方法从使用晶格气体方法的玻尔兹曼方程开始。 玻尔兹曼方程可以使用查普曼-恩斯科格展开来恢复Navier-Stokes 方程。 该方案非常适合大规模并行计算，因为节点的更新只涉及其最近的邻居。 边界条件易于实施。 因此，代码很简单，可以以适合并行处理的形式轻松编写。 在晶格气体方法中，存在晶格气体自动机（LGA）[1]。 LGA在模拟服从Navier-Stokes 方程的现实流体流动方面存在根本性的困难。 除了其内在的噪声特性使计算精度难以实现外，即使在流体极限中，它也包含某些属性。 由于两个基本问题，晶格气体流体动量方程不能简化为Navier-Stokes 方程。 第一个是由于对流系数的密度依赖性而导致的非伽利略不变性。 这仅限制了LGA方法的有效性，只有严格的不可压缩区域。 其次，压力具有明显的非物理速度依赖性。 为了避免其中一些问题，日本京都左京区五所海道町松崎京都工业大学机械与系统工程系的几个606-8585。 Frontiers of 计算流体动力学 - 2002 编辑：David A. Caughey和Mohamed M. 哈菲兹 ©2002 世界科学

\section{12.2 格子 Boltzmann 方法用于二维}
已经提出了228个SATOFUKA和ISHIKURA晶格玻尔兹曼（LB）模型。 [2] -[4] LB方法的主要特点是将粒子占用变量n\textasciicircum{}（布尔变量）替换为单粒子分布函数（实变量）ft =（n\$），其中（）表示进化方程中的局部集合平均值，即晶格玻尔兹曼方程。 Chen等人[5]和Qian等人[6]提出的LB模型应用了Bhatnager、Gross和Krook在1954年首次引入的单松弛时间近似[7]，大大简化了碰撞算子。 这个模型被称为晶格BGK（LBGK）模型。 在本研究中，LBGK方法用于模拟大规模并联计算机（日立SR2201和惠普Exempler V级）上衰变的二维和三维均匀各向同性湍流。 还计算了有/没有突然膨胀的二维通道流和三维管道流。 详细比较了LBGK方法和传统有限差分法（FDM）之间的准确性、物理保真度和效率。 在三维计算中，一维区域分解的加速和CPU时间呈现。 本文的组织结构如下。 第12.2节介绍了本文中使用的方晶格模型的晶格玻尔兹曼方法。 第12.3节介绍了二维均匀各向同性湍流的晶格玻尔兹曼模拟。 还讨论了格子 Boltzmann 方法与传统高阶FDM方法相比的准确性和效率。 第12.4节介绍了有/没有突然膨胀的二维通道流的应用。 三维计算的扩展和应用从第12.5节到12.7节。 第12.8节介绍了并行化的一些结果，最后一节包含总结性评论。 12.2 格子 Boltzmann 方法用于二维12.2.1九速平方晶格模型在本节中概述了碰撞算子的LB方法和BGK模型。 使用具有单位间距的方形晶格，每个节点有八个最近的邻居，由八个链接连接，如图所示。 1. 粒子只能驻留在节点上，并在单位时间内沿着这些链接移动到最近的邻居。 因此，有两种类型的移动粒子。 1型粒子以速度|ei*| = 1沿轴移动，2型粒子以速度|e2j| = y/2沿对角方向移动。 每个节点也允许速度为零的静止粒子。 三种粒子的占用由单粒子分布函数/\textasciicircum{}（x，t）表示，其中下标a和i表示

晶格BOLTZMANN模拟229 <'-U>\textasciicircum{}*P± *SS-<i+lJ）图1方形晶格模型。分别是粒子类型和速度方向。 当o = 0时，只有/oi。 分布函式fai是在速度eai的节点x和时间t处找到粒子的机率。 粒子分布函数满足晶格玻尔兹曼方程/CTi(x + eai, t + 1) - /CTi(x, t) = nai (12.1)，其中f2CT»是表示因碰撞而导致的粒子分布变化率的碰撞算子。 根据Bhatnagar、Gross和Krook（BGK）的说法，碰撞算子使用单一时间松弛近似法简化。 因此，晶格玻尔兹曼BGK（LBGK）方程（以晶格单位为单位）是fai(-x. + eai,t + l) - fai(x.,t) = - /cri(x,i)-/V;(x,i) r L (12.2)，其中f\textasciicircum{} (x,t)是x, t的平衡分布，r是控制接近平衡率的单松弛时间。 每个节点的密度，p，宏观速度，u，由P = \textasciicircum{}2\textasciicircum{}2Ui，pu = Y\textasciicircum{}\textasciicircum{}f<rie°i（12.3）的粒子分布函数定义，可以为每种类型的粒子选择合适的平衡分布，f(0) J01 pa - -pu" /i(-\} = PP + \textbackslash\{\}p(eu • u) + ip(eii • u)2 - ipu2 (12.4)

230 SATOFUKA \& ISHIKURA (0) (1 - 4/3 - a) 1 1 2 1 2 Hi = P\textasciitilde{} 1 + i2P(e2i • u) + -p(e2i • u)2 - - pu2 松弛时间与黏度有关，黏度为6\textasciicircum{}+1 (12.5)，其中v是以晶格单位测量的运动黏度。 Hou等人[8]使用了a = 4/9和/3 = 1/9的值。 平衡种群是通过假设它们可以以速度和密度的形式表示为功率级数来确定的；fl? = Aai(P) + Bai(p)(eai • u) + Cai(p)(eai • u)2 + Dai(p)u2 (12.6) 然后应用Chapman-Enskog程序来确定该模型的宏观行为。 选择A„i、Bai、Cai和Dai的值，以便宏观行为与Navier-Stokes 方程匹配，尽可能高阶。 结果的连续性和动量方程如下；\textasciicircum{}+«S\textasciicircum{}+i("(£ + \textasciicircum{})))+0(e!) +0(M")(128» 特征无维参数是马赫数，Ma = \textbackslash\{\}/Wjc，其中U是特征的宏观流速，Knudsen数与e = CT/L成正比，其中L是宏观流长，以及雷诺数，Re - pUL/p。 选择合适的晶格大小和特征速度后，v可以为给定的Re数计算，然后可以使用Eq来确定松弛时间。 12.5. 从初始密度和速度场开始，可以使用Eq获得平衡分布函数。 12.4，并且/CTi（x，t）可以初始化为/\textasciicircum{}（x，t）。 对于每个时间步骤，粒子分布的更新可以分为两个子步骤：碰撞和对流。 长期以来，哪个是第一个无关紧要。 位置x的碰撞过程根据等式给出的晶格玻尔兹曼方程的右侧发生。 12.2. 结果在x处的粒子分布是原始分布和碰撞项的总和，然后被对流到x + eai的最近邻居，对于粒子速度。 然后p，u可以使用Eq从更新中计算。 12.3. 当满足某个标准时，更新程序可以终止稳态问题。 该方法也可用于瞬态问题。

\section{12.3 二维均匀各向同性湍流}
晶格玻尔兹曼模拟231 12.3 二维均匀各向同性湍流 12.3.1 初始和边界条件 涡量 u的初始条件通过满足关系随机确定，E(k) = \textbackslash\{\} E Mkuk2)\textbackslash\{\}2/k2 = \textasciicircum{}kexp(\textasciitilde{}k) (12.9) |fc-fc|<l/2 其中ui表示傅里叶空间中的涡量，k2 = k\textbackslash\{\} + k2，k\textbackslash\{\}和fc2是波数。 周期性边界条件施加在x和y方向上。 计算域是一个平方，(0,0)<(a;,y)<(27r,27r)。 12.3.2 高雷诺数模拟 作为高雷诺数均质各向同性湍流的大规模直接数值模拟，对v = 0.0001的案例进行了模拟。 这对应于初始积分尺度雷诺数 RL = 25500，表示为RL = Q/vr)1/3。 Q和rj表示总能量和entrophphy耗散率，定义为J\textasciicircum{}(t) = / E(k)dk (12.10) Jo /•OO T]\{t) = 2v / k4E(k)dk (12.11) Jo 晶格节点的数量是1025 x 1025。 图。 2（a）显示了t=3.0的涡量等高线图在方晶格BGK方法和lOth-order FDM之间的比较，而图。 2（b）显示了稍后时间t=6.0的比较。 虽然涡量轮廓在t=6.0时有明显的轻微差异，但与lO阶FDM的解相比，LBGK模拟可以发现非常相似的特征。 在t=3.0时比较波数光谱k3E(k)表明，两种方法在流动的统计行为方面得出了非常相似的答案。 在当前1025 x 1025节点的晶格下，可以解析二维湍流的惯性范围。 我们还发现，有一个波数k < 50的范围，其中k3E(k)大致是恒定的，因此E\{k)与k\textasciitilde{}3成正比。 表1比较了SGI POWER ONYX 10000上这种高雷诺数模拟的LBGK方法的计算成本

232 SATOFUKA \& ISHIKURA (a) t = 3.0 (b) t = 6.0 图2 LBGK和FDM之间的涡量 cotours的比较（lOth-order）atv = 0.0001。 表1计算成本比较LBGK FDM（10阶）\textasciitilde{}v/（u2）51）4 LO时间步骤4047 1000（At = 0.001）TNS 1.0 1.0 CPU时间63006136639（比率）（1）（2.17）CPU时间/时间步骤15.47 136.64（比率）（1）（8.83）计算机：硅图形ONYX 10000晶格节点数：1025x1025动黏度：0.0001 雷诺数：25500

\section{12.4 突然膨胀的二维通道}
LOth-order FDM的LATTICE BOLTZMANN模拟233。 就效率而言，LBGK方法每特征时间的CPU时间不到FDM的半部分。 这些结果表明，LBGK方法与使用相同晶格大小的lOth-order FDM一样准确，可以作为求解Navier-Stokes方程的替代方案。 12.4 突然膨胀的二维通道 12.4.1 问题描述 突然膨胀的对称通道中的流动是一个流体力学系统，具有各种有趣的现象。 该系统已经通过实验和数值进行了研究。 这里研究了一个二维通道，膨胀比为1:2，膨胀的每边长宽比为1:4。在雷诺数的中等值下，如Re = 10，通道中的流动是二维的。 不可压缩流动的9速平方晶格BGK模型用于模拟通道中的二维流动。 系统的Rynolds数定义为Re = HU0（12.12），其中H是入口部分的高度，Uo是最大入口速度，v是运动黏度。 在这个模拟中，Rynolds数被选择为10。 图中显示了突然膨胀的二维通道的几何配置。 3. A •« AH • •* AH A 图 3 突然膨胀的二维通道的几何配置

234 SATOFUKA \& ISHIKURA 12.4.2 边界和初始条件 在入口（上游），强制执行速度水平分量u的抛物线轮廓，最大UQ-0.1，速度的垂直分量v设置为零。 在出口（下游），施加恒定压力边界条件，p-1.0。 在墙壁上应用了防滑边界条件。 在通道内部，平均密度p0设置为1.0，速度场的初始值设置为零。 12.4.3 非均匀网格 模拟中使用了两种类型的计算网格：图的下半部分显示的均匀网格。 4，以及图的上半部分所示的非均匀矩形网格。 4. 在许多实际应用中，非均匀的网格肯定是可取的。 抑制在对数LB方法中使用非均匀网格的是物理空间和动量空间的离散化耦合。 密度分布必须在单个时间步程中从一个晶格位置移动到另一个晶格位置，以保证随后的碰撞过程。 在目前的方法中，计算网格与动量空间的离散化脱钩，它可以有一个任意的形状。 碰撞仍然发生在计算网格的网格点上。 碰撞后，晶格位置的密度分布现在可能无法准确确定，它们始终可以使用插值来计算。 插值碰撞和对流步骤重复后。 模拟中使用了128 x 33 + 129 x 65的均匀矩形网格和27 x 33 + 28 x 65的非均匀矩形网格。 12.4.4 模拟结果 图。 5显示等线（图。 5（a））和流函数的等高线（图。 5（b））用于突然膨胀的对称通道中的流动。 在这些图中，上半部分对应于使用非均匀网格模拟的，而下半部分表示均匀网格。 很明显，该协议相当出色。 使用非均匀网格模拟的计算成本小于使用均匀网格的计算成本的1/5。

\section{12.5 格子 Boltzmann 方法用于三维}
晶格玻尔兹曼模拟235 si\textasciicircum{}fflfflllllllllllllffffifflM 图4 Nnuniform mesh（上）和均匀网格（下）«5 TTC r k i \textasciicircum{}）i x/h（a）等压线f x/h（b）流函数图5 三维12.5.1十五速度立方晶格模型的模拟结果12.5 格子 Boltzmann 方法的简单方法，将正方形晶格扩展到正方形的边角，向量到正方形的边角，是使用向量到立方体的边和角。 这定义了具有eu e (±1,0,0)、(0,±1,0)、(0,0,±1)和e2i e (±1,±1,±1)的体心立方晶格。 这些移动种群和非移动种群的平衡分布函数由f(o) \_ J01 - pa -pu* Ji (0) pP + 3\textasciicircum{}(eH • u) + 2\textasciicircum{}(eii u)5给出，用于沿晶格轴的eij，以及(0) (1-6/3-a) 1 1 2 Hi = P § + \textasciicircum{}P\textasciicircum{}i • u) + YQp(G2i ' U\textasciicircum{};PU 48 pilfer e\textasciicircum{}i沿立方体角的链接给出。 使用a ft = 1/9的值。 （12.13）（12.14）2/9和

\section{12.6 三维均质各向同性湍流}
236 SATOFUKA \& ISHIKURA AZ (i-i,j-i,k+i) T (i-lj+U+l) (i+ij-ijc+n..-,-»(i,j,k+l) ";V(i+lj+U+l) '(i-lj-l,k-l),f. X dj£'.'. fiu,K-;i:" (i+i,j-i,k-i)...-•(i-lj+U-l) (i+ij+u-i) 图6 立方晶格模型。 放松时间定义为6v+l（12.15）12.6三维均匀各向同性湍流 12.6.1初始和边界条件三维衰变湍流用具有能量谱的随机初始条件模拟；E(k) = 16V2/Tr\textasciicircum{}k\textasciicircum{}axk4 exp[-2(k/kmax) (12.16)，其中我们设置v\$ - 1.0和峰值波数kmax = 2.37841。 周期性边界条件在x、y和z方向上施加。 计算域是一个立方体，0 < (x, y, z) < 2TT。 12.6.2 低雷诺数模拟 运动黏度选择为v = 0.0025。 初始积分尺度和微尺度雷诺数为340，对应于中等雷诺

\section{12.7 三维管道流量}
晶格玻尔兹曼模拟237（a）LBGK（b）FDM图7獟的轮廓表面\{y = 0.0025，t = 0.7,65 x 65 x 65，fipiot = 100）。数字模拟。 晶格节点的数量是65 x 65 x 65。 立方晶格BGK方法和FDM在t=0.7时的熍的轮廓表面如图所示。 分别是7（a）和（b）。 图值为f2 = 100。 与涡量一样，两种方法之间的熵等高线表面相似。fl表示熵，定义为to(x,y,z) = ±(u>Z + wl + w2z) (12.17) 立方晶格BGK方法和FDM之间的比较表明，立方晶格BGK方法可以作为求解Navier-Stokes方程的替代方案。 12.7 三维管道流量 12.7.1 问题描述和边界条件 层流发展的数值模拟在方形管道中进行，方形管道突然膨胀，均匀的台阶高度等于入口管道宽度的0.5倍，如图所示。 8. 防滑墙边界条件，u-v=w=0，用于二维模拟。 在入口，我们假设均匀轴向速度u = 1.0，其他速度分量设置为零。 流出边界的压力保持在p = 1.0的恒定值。

SATOFUKA \& ISHIKURA 图 8 突然膨胀的方管 0.00 0.05 0.10 0.15 0.20 图9 速度的轴向分量的发展 Re - 100 12.7.2 模拟结果 模拟是通过使用均匀的立方晶格对入口和出口管道进行的。 入口管道使用160 x 33 x 33的晶格节点，而出口管道使用161 x 65 x 65。 速度的轴向分量的发展如图所示。 9. 在这种情况下，由于管道长度不足，速度曲线尚未完全开发。

\section{12.8 并行化}
格子玻尔兹曼模拟区域3 /\textbackslash\{\} regwn2区域•i！ I<>n(l / 图 10 区域分解 12.8 并行化 12.8.1 区域分解 用于三维二维模拟显示，水平方向的子域大小越长，CPU时间越短。 换句话说，数据的外环尺寸越长，CPU时间越短[9]。 基于这种二维经验，我们选择了图中所示的区域分解类型。 10. 虽然立方晶格模型的每个节点上都放置了十五个分布函数，但与方形晶格模型的情况类似，只需要通过每个域边界发送五个变量。 12.8.2 加速和CPU时间 加速和CPU时间如图所示。 11和表2。 尽管晶格节点的数量很少，但与二维模拟相比，我们可以获得更高的速度。 表2 CPU时间和加速 CPU数量 12 4 8 CPU时间 18336 8921 4585 2660 加速 1.0 2.0 4.0 6.9 11.8

\section{12.9 结论}
\section{参考资料}
240 SATOFUKA \& ISHIKURA 4.0 \textasciicircum{}3.0 -§2.0 <*> 1.0 0.0 N = log2C S = log2T! / y '"7¥ //; // '• 0 12 3 4 CPU数量（N）图11 加速12.9 结论 使用LBGK方法对衰变均匀各向同性湍流和二维通道流动进行突然膨胀的二维模拟表明，该方法与使用相同晶格尺寸的传统方法一样准确。 LBGK方法能够再现不可压缩流动的动态，可以作为解决Navier-Stokes 方程的替代方案。 三维流动的模拟也显示了LBGK方法的准确性和效率。 在格子 Boltzmann 方法的并行计算中，发现水平方向子域的大小越长，CPU时间越短。 在三维模拟中，与二维模拟相比，观察到的速度相对较高。 需要进一步研究立方晶格BGK模型的准确性和效率。 参考资料。 弗里施，U。 \& Hasslacher，B。 \& Pomeau，Y.，Navier-Stokes方程的晶格-Gas自动机，物理。 牧师。 Lett。 56，1986年，pp. 1505-1508.. Heiguere，F。 \& Jimenez，J.，用晶格玻尔兹曼方程模拟圆柱体周围的流动，Europhys，Lett。 1989年9月，pp. 663-668.. Heiguere，F。 \& Succi，S.，用晶格玻尔兹曼方程模拟圆柱体周围的流动，Europhys，Lett。 8，1989年，pp. 517-521。

格子玻尔兹曼模拟241 4。 麦克纳马拉，G。 \& Zanetti，G.，使用玻尔兹曼方程模拟晶格-气体自动机，物理。 牧师。 Lett。 61，1988年，pp. 2332-2335。 5. 陈，H。 \& 马修斯，W. H.，使用晶格-气体玻尔兹曼方法恢复Navier-Stokes方程，物理。 牧师。 A. 45，1992年，pp. 5539-5542。 6. 钱，Y。 H。 \& D'Humieres，D. \& Lallemand，P.，Navier-Stokes方程的晶格BGK模型，Europhys。 Lett。 17，1992年，pp. 479-484。 7. 巴特纳加尔，P. L。 \& Gross，E。 P。 \& Krook，M.，气体碰撞过程的模型。 1. 充电和中性单组分系统中的小振幅过程，物理。 牧师。 94，1954年，pp. 511-525。 8. 侯，S。 \& 邹，O。 \& Chen，S. \& Doolen，G。 和Cogley，A。 C，格子 Boltzmann 方法对腔流的模拟，J。 计算机。 物理。 118，1995年，pp. 329-347。 9. 里福，N。 \& 西冈，T。 \& Obata，M.，不可压缩流的晶格玻尔兹曼方程的并行计算，Proc。 平行CFD '97会议，1997年5月19日至21日，pp. 601-608。

\clearpage
\chapter{13 入口弱电离气体对高超声速流动的MHD效应的数值模拟}
\section{13.1 摘要}
Ramesh K入口弱电离气体的MHD对高超声速流动的影响的数值模拟。 Agarwal and Prasanta Deb 13.1 摘要 近年来，通过在微电离等离子体上施加磁场来减少超音速阻力的可能性受到了极大的关注。 俄罗斯文献中报告的一些结果表明，在相同温度下，微电离气体（等离子体）的冲击结构明显弱于非电离气体的冲击结构。 美国目前正在调查超声速流动中激波修改/消散的这个概念和其他概念。 AJAX计划下的空军。 本文通过数值模拟来评估这个概念。 使用具有双温度模型的可压缩黏性MHD方程，研究了磁场对弱电离流动的影响。 广义坐标中的二维MHD方程使用二阶精确空间离散化的修改后的Runge-Kutta 时间积分方案求解。 对称的Davis-Yee总变异减小（TVD）通量限制器用于抑制冲击区域的振荡。 对二维scramjet入口进行了数值模拟，并与Park、Bogdanoff和Mehta的一维理论分析进行了比较；得到了良好的一致性。 数值结果还表明，通过在弱电离气体中以超音速移动的物体表面施加强磁场来减少超音速阻力的可行性。 国家航空研究所，威奇托州立大学，堪萨斯州威奇托，67260-0093 计算流体动力学的前沿 - 2002 编辑：David A. Caughey和Mohamed M. 哈菲兹 ©2002 世界科学

\section{13.2 命名法}
244 AGARWAL 13.2 命名法Q场向量E通量向量分量x方向的F通量向量分量x笛卡尔座标y笛卡尔座标p密度u速度分量x方向v速度分量\textasciicircum{}方向w速度分量z方向V速度向量Bx磁场分量x磁场分量y方向的磁场分量z方向p静压Pi离子压力pa中性原子压力pe电子压力Po Pa+Pi j-离子气体的比热比热比热比 电子气体y„比热比（离子+中性原子）气体Se电子熵e，每单位质量的总能量c.5声速Vox Alfven波速度在x方向vay阿尔夫温波速度在j方向X 特征值“t，广义坐标T）广义坐标J雅各布变换v\textasciicircum{}快波速度在\textbackslash\{\}-方向v\textasciicircum{}快波速度在T|-方向v\textasciicircum{}慢波速度在£-方向v\textasciicircum{}慢波速度在T\textbackslash\{\}-方向L\textasciicircum{}逆右特征向量与Jacobian 矩阵 A R\textasciicircum{}右特征向量与Jacobian 矩阵 A D\textasciicircum{}对角线相关 特征值矩阵与Jacobian 矩阵 A L\textasciicircum{}逆右特征向量与Jacobian 矩阵 B相关的R\textasciicircum{}右特征向量与Jacobian 矩阵 B相关

MHD效应的数值模拟 8, 5 Vf a Re« IvCmoc Pr ux H Tg Te ne K D12 me e Qa a J IBI E Pe Q Sc q k a 对角线特征值矩阵与Jacobian 矩阵相关的矩阵 B通量限制器函数向量与B与B通量限制器函数相关的熵校正变量通量限制器流体电导率的磁渗透性流体电导率马赫数 雷诺数磁性雷诺数磁性Prandtl数自由流速度分子黏度气体温度电子温度电子数密度离子密度中性原子密度 rij+na玻尔兹曼常数，1.38xl0"23焦耳/°K二元扩散系数电子质量电子电荷碰撞横截面电离程度电流密度向量磁通量向量电场向量电子密度霍尔参数介电常数自由空间源项黏性应力张量热通量矢量导热电率电离程度

\section{13.3 导言}
246 AGARWAL 13.3 介绍 近年来，通过在弱电离等离子体上施加电磁场来降低超音速阻力的可行性备受关注。 美国空军与美国宇航局合作，目前正在对这一概念进行数值和实验研究。 对这一概念的兴趣源于俄罗斯文献中报告的结果，即弱电离等离子体的冲击结构被证明比相同温度下的非电离气体的冲击结构弱[1-3]。 已经提出了几种假设来解释气体放电等离子体的观测特性，这些特性可以分为两组。 第一个假设是基于一个假设，即弱扰动以大于热声速的速度在等离子体中传播；第二个假设（“能量假说”）是基于这样的假设，即能量在冲击前线后面的冲击层中释放（要么能量释放集中在分子自由度中，要么是电流在激波中流动的结果）。 美国空军目前正在AJAX计划下研究超声速流动中激波修改/消散的这些概念和其他概念。 此外，正在评估MHD旁路概念，以提高效率并扩大传统喷气发动机的续航里程。 MHD旁路scramjet概念（图1）包括一个外部MHD发电机，用于从进入发动机的空气中提取电力，一个电离器，用于控制进入发动机的流动电导率，一个与内部MHD发电机串联的MHD加速器，用于绕过发动机燃烧器周围的能量，并将能量返回发动机的排气。 内部MHD发生器用于增加压力，防止流动分离，并将部分气流焓转换为电力。 在20世纪50年代末和60年代，进行了几项实验和分析研究[4-7]，以研究MHD对高超声速流动场的影响。 其中，特别值得注意的是Kantrowitz的论文[8]，其中他建议使用MHD来修改弱电离等离子体中强激波的结构。 最近，Mnatsakanian等人。 [9]，Bityurin等人[10]和Cole等人[11]研究了MHD对空气动力学和推进系统流场的影响。 MHD方程与麦克斯韦方程一起表征了磁场和电场存在的导电流体的流动。 Park等人[12]最近进行了理论计算，以预测MHD旁路scramjet发动机的最大理论性能。 Bityurin等人[13]通过应用基本的热力学原理，报告了对MHD旁路高超音速空气呼吸发动机的性能潜力和科学可行性的评估。 Chase等人[14]讨论了各种技术问题，如与AJAX系统相关的平衡电离和非平衡电离。 Kuranov等人[15]探索了通过正确选择离子发生器和MHD系统其他组件的参数来提高scramjet效率的可能性。 Laux等人[16]开发了一个空气等离子体的双温度动力学模型，用于对MHD系统的数值模拟，该模型根据实验数据进行了验证。

\section{Electro-磁流体动力学的13.4 控制方程}
MHD效应的数值模拟247 本研究的目的是研究电磁场对以高超音速速度移动的scramjet入口中弱电离等离子体流动的影响，使用具有双温度模型和可变电导率模型的可压缩二维电泳道方程。 鲍威尔[17]提出的二维可压缩黏性MHD方程与布拉西尔和加利斯[18]提出的双温度模型，使用二阶精确空间离散化的修改四级Runge-Kutta（R-K）时间积分方案求解。 使用具有总变异减小（TVD）通量限制器的守恒定律（非MUSCL）的非单调上风方案来稳定修改后的R-K方案。 控制方程在广义坐标中得到解决。 为实现数值方案开发了一组特征向量。 Refs. [19, 20]中报告的2D非稳态可压缩黏性MHD代码WSUMHD2D被修改为包括electroMHD模型以及双温度模型和可变电导率模型。 Park等人[12]在理论分析中使用的流动条件，2D电MHD代码用于计算弱电离气体的高超声速流动，该Scramjet入口与MHD发生器集成在一起。 13.4 Electro-磁流体动力学的控制方程 Electro-磁流体动力学的控制方程包括流体流动的质量、动量和能量方程以及电磁学的麦克斯韦方程[参考。 21、22、23]。 这些方程是：连续性方程动量方程\textasciicircum{}+V.(pV)=0 (1) \textasciitilde{}\textasciicircum{}p + V«[pVV -BB+(p+B«B/2)I]=V«T (2) 能量方程 -*jjp +V.(pVet) =-V«(V/7)+V.( V.T)+V.q+E«J (3) 麦克斯韦方程 VxB = n„J (4), VxE = -\textasciicircum{} (5), V.B = 0 (6), V«E = \textasciicircum{} (7)

248 AGARWAL 压力-能量关系 热通量-温度关系 q= -kVT（9）电荷守恒方程-\textasciicircum{} + V.J=0（10）电流密度欧姆定律J =<T（E + VxB）-§j（JxB + Vpe）（11）电子和离子数密度关系pe = neKTe（12）Pe \textasciitilde{}（rn+rie+na）p°\textasciitilde{}（2ne+na）Po（l3）电离程度ne ne ne... pe a=K\textasciicircum{}:r\textasciicircum{}\{lfne<<na）\textasciitilde{}j；（14）可变电导率模型等离子体的电导率由混合物中带电粒子的数量决定。 对于弱电离气体，a\textasciitilde{}a？ *“0'5，其中弱电离气体的a < le-4。 对a对J的依赖性进行建模，对于准确捕捉MHD相互作用的强度至关重要。 以下模型用于此目的。 查普曼和考林给出了微电离气体（a<le-4）的电导率表达式（参考。 22）作为e2neD12 \textasciicircum{}\_ <TL = \textasciitilde{}KfT' (15)

\section{弱守恒律形式的13.5 控制方程}
MHD效应的数值模拟249 339\{KT)0S，其中Dn= g \textbackslash\{\}05（16）在方程（15）和（16）中替换各种常数的值，得到以下方程。 OL = OJ，（17）QaT，其中Qa是从参考中给出的碰撞曲线中获得的。 23和常数ca是根据a和a的初始值进行评估的。 高度电离气体（a »le-4）的电导率表达式由Spitzer-Harm（参考。 23）关系：1.56 10-4?;'5 qtf =. /,<. n<\textbackslash\{\} (18) " ln\{l2300Te,5/ne05) 使用坎特罗维茨的假设（参考。 8），任何电离程度的电导率都可以通过表达式I=-U-L获得。 （1„双温度模型双温度模型包含在MHD方程中，通过包括电子熵的对流方程来模拟离子和电子的电离效应。 在这个模型中，假设导电流体由离子和电子组成，其特征是其自身的平移温度。 电子熵方程由M给出。 + 3P\textasciicircum{} + \textasciicircum{} = o, (20) dt dx dy 其中Se=\%- (21) 13.5 控制方程 在弱守恒律形式中，在笛卡尔坐标系中具有双温度模型的二维黏性电MHD方程可以用以下守恒律形式表示：

250 AGARWAL M + M + \textasciicircum{}\textasciicircum{}H, Bt Bx By Q = ]p pu pv (BBX BBy Bx By Bx By E = pw Bx By Bz pe, pSe]T pw +p+- puvpu -Bt+tf+B: \textasciitilde{}x \textasciitilde{}y \textasciitilde{}z puw- SreHy 4n\textbackslash\{\}if BB X Z uBy-vBx uB-wB, pe,+p+-\%+\% 871(1, 7 J u -\textbackslash\{\}uBr +vBv +wB7 4n\textbackslash\{\}ifK " ' M puSe J (22) (23) (24) pv pvw- BA 471\textasciicircum{} Bf+Bi+B? Pv2+p+-\textasciitilde{}y '\textasciitilde{}x'\textasciitilde{}z 8jt\textbackslash\{\}lf pvw- BA 47I|X V vBx-uBy vB-wB pe,+p+ Bl+fy+Bl V 8;t|i B 7; 4m M«B,+VBV+WBJ P\textasciicircum{} J (25) HM=-\textbackslash\{\} By Bz uBx+vBy+wBz T -J- -- u v w: Q 4n\textbackslash\{\}if |j 4n\textbackslash\{\}if 4n\textbackslash\{\}yf 4n\textbackslash\{\}if (26)

MHD效应的数值模拟E =（J. /CT dx 1 dBy \textbackslash\{\}ifa dx 1 dBz [ifa dx o (27), F = J 1 3BV 1 3BZ MT„+VT,T+WTre-\textasciicircum{},. 0 J Sc = [0 0 0 0 S5 S6 57 58 Of c Q 3 L 3\textasciicircum{} „ \&»J „ Q 9 [ „ 3ff,, 3fl jl 其中S5 = 7\textasciicircum{} T"? \#vT\textasciicircum{}+\#fTT。 S<5=\textasciicircum{}- T\textasciitilde{} i ByT\textasciitilde{}+Bx-\textasciicircum{}r',r, \textbackslash\{\}B\textbackslash\{\}\textbackslash\{\}ltGdy\textbackslash\{\}( ydy 3x|J ' ° \textbackslash\{\}B\textbackslash\{\}\textbackslash\{\}lfOdx\textbackslash\{\}[ ydy 3x|J ' Q d (a aBv TiaBx xaBA Q d (aaB? dBv aB* b7\textasciitilde{} IBI\textasciicircum{}dx[Bxdx'Bxdy'Btdyy + IBIjifO dy [Bxdx + Bydx 'Bydy s s =J-(Mi*Mx+Mi*<)IhJiBx MxYMi MJ 8 |IfC7 vdx dx dy dy \textbackslash\{\}dy )\textbackslash\{\} dx dx p p(«2+v2+w2) B2x+B2y+B2z Y-l 2 8jt|x/ In Equations (27) and (28): 2 du dv * 3x dy > \textasciicircum{}rv - Tw \textasciitilde{}\textasciitilde{} M-: 3« 3V 3>> dx dw \textasciicircum{} 3 = 3»i dv du 3y 3x dw «x=-\&-r-

\section{13.6 控制方程在广义坐标}
252 AGARWAL 注意P =P*+Pi+Pe = Po+Pe，J = y0 (l+b)/(l+a)，a = pe (y. -l)/p0(y0-l)，b = ayjy0。 这些方程被非维化并转换为体拟坐标系（£，T|）。 非维化完成方式如下：- tu„, x * y, p, u, v w p t =--, x =-, y = -. p =- u = -, v = -, w = -, p = -, L L L L l px u„ «„ p„ p„u„ r=f, H'=±,B; =»>,*>,*«,se=*L T- v\textasciicircum{}\textasciitilde{} v\textasciicircum{}/p\textasciitilde{}"- v\textasciicircum{}/p-"\textasciitilde{} v\textasciicircum{}/p- M\textasciitilde{} M\textasciitilde{} 得出的非维参数是：雷诺数: Re = p°°°°° Prandtl数: Pr = -- k Magnetic 雷诺数: Remx =a\textbackslash\{\}lfuxL Freestream 马赫数: M\textasciicircum{}-u\textasciicircum{}l •yjyop0x Dt, d\textbackslash\{\}] a\textasciicircum{} a\textasciicircum{} (3i) x, 5n Q JxE + ZyF T\textbackslash\{\}zE + l\textbackslash\{\}yF J' J J \$xBx+\$yBy \_r\textbackslash\{\}xBx+r\textbackslash\{\}yBy \_ \%xEv+\%yFv \_ r\textbackslash\{\}xEv+T\textbackslash\{\}yFv \#,:,By -, Ev =, Fv =- J

MHD效应的数值模拟E:,= --- \{aVi,+<h\textbackslash\{\} +c\textasciicircum{} +c3vj Re. \textasciicircum{} 回复。 7 l4v\textasciicircum{}+c5wJ Rem / [a4(Bj,+c5(Bjl sb\textasciicircum{}+\textasciicircum{}W Rem 7 k(Bz\textbackslash\{\}+4Bz\textbackslash\{\}] \textasciicircum{}7 N«4 +1 <h\{v2\textbackslash\{\} +\textbackslash\{\}aj\textasciicircum{}\textbackslash\{\} +4uv\textbackslash\{\} +-\textasciicircum{}T a,T\% +lci("2lci\{v2\textbackslash\{\} +lc\&\textbackslash\{\} + ciu\textbackslash\{\} +c4v"n +-S\textasciicircum{}IC5:r'' 0 J F„ = - (k3un + V„ + c3"\textasciicircum{} + c2v4) \textasciicircum{}-j\{b\textasciicircum{}+ciW\%) l--k(BX+c5(Bx\textbackslash\{\}] mj • i • « 7z„y s Re. / Re\_ Re„7 Rem Rem, Rem, 2*,(A 2 -\&,(«')„ + -t>M>l +-t>,\{w2i +\textasciitilde{}bJuv), +-;-\textasciicircum{}, „,.„ Pr(y-l)M\textasciicircum{} W„ + \textasciicircum{}(H2)4 +-c2\{v2\textbackslash\{\} +-c5\{w2\textbackslash\{\} +c3vu\% +c4uvk +\textasciicircum{}-\textasciicircum{}T\textasciicircum{}i 0 J

\section{13.7 数值方法}
254 AGARWAL在方程（32）和（33）中：«。 =\#,+\$,. a,=\%\textbackslash\{\}\textasciicircum{}\textbackslash\{\}\textbackslash\{\}, 0,=\textasciicircum{},. aA=\textbackslash\{\}\textbackslash\{\}+l,\textbackslash\{\} * =f M» + £/n,. c2 =\textasciicircum{}T1, +K,Tly, c3 =\textasciicircum{}11, -f \textasciicircum{}ri,., C4=§xTly-K,Tlx • 和C5=\$xT\textbackslash\{\}x+\$yT\textbackslash\{\}y 定义雅各布矩阵A和B如下图所示，MHD方程可以写成f + jf + 5f=|,+§, (34, dt dt, 3i] 3\textasciicircum{} 3ti 其中A=\textasciicircum{}L-HM\textasciicircum{}- = kxA+lB), B ===-HM\textasciicircum{} = I\$\textbackslash\{\}ZA+T\textbackslash\{\}XB), dQ dQ dQ dQ dQ dQ 本节提供了特征值和特征向量。 Jacobian 矩阵 A的特征值是V =\textbackslash\{\}x" + \%yv。\textasciicircum{}± =\$*("±vJ+Mv±vJ \textasciicircum{}=\textasciicircum{}x" + 5,v，和Xhk=Z,xu + £,yv。 Whefe v\textasciicircum{} = ik(v.2+c»)+zt]' v\textasciicircum{}=i[a4(v>cs2)-ZJ'

MHD效应的数值模拟 255 h=4\textasciicircum{}\textbackslash\{\}\textbackslash\{\}aM+\textasciicircum{})2fcxBx+!-yBy)2yp 7Cp2 JP\_ P vi = Bl 47Cp 4np vi = 4rcp v2=v2 +y2 +y2 "a K ax a\textbackslash\{\} az 对角线特征值矩阵，D\textasciicircum{} = Z\textasciicircum{}A/\textasciicircum{}然后变成K o 1 n where and Vlro r r-v, H -va5 +vft \textasciitilde{}vf£, I,.. 1„. +*\% -v55 r\&, hr\textbackslash\{\} -H hva\% l-,a!. l+vf!. l-vA '+vJ? '-V4 ld\% '/hi J是左右特征向量。 同样，Jacobian 矩阵 B的特征值是\textasciicircum{}on=Tlx"+1,v, lar]±=r\textbackslash\{\}x\{u±vJ + T\textbackslash\{\}y\{v±vay), kdr]=T]xu + J]yv, 和Xkr]=T]xu + T]yv (35) (36) (37) 其中vWMv.2+*;)+*„]• <=\textbackslash\{\}[bM+c])-z,]

256 AGARWAL和zn = A -A \textbackslash\{\}va +c,) -2 V "P 对角线特征值矩阵，Dn =LnBRn然后变成A K \textbackslash\{\}m+ •STJ+ K ol (38)，其中*n = 'ro \textasciicircum{} -"A (39)和4o A,=l'o '+,., \textasciicircum{} \textasciicircum{} \textasciicircum{} '• '*, h (40) 13.7.2 修改后的Runge-Kutta (R-K) 时间积分方案用于时间积分的数值方案是修改后的R-K方案。 广义坐标系中二维黏性MHD方程的R-K公式如下所示：Qy=Q-j-At(LQ«\% Q\textasciicircum{}=Q't\textasciitilde{}\textasciicircum{}m)i.j (41a) (41b) (41c) (4 Id) 其中+1 ri.j-l (LQ)... = EMJ E'-1J +\textasciicircum{}I K *'••> 2A£ 2An

MHD效应的数值模拟 257 -H„ 2A\% 2An j 3\textasciicircum{} df\textbackslash\{\} 13.7.3 总变异减小（TVD）模型 上述修改后的R-K方案对于解决黎曼型问题变得不稳定，因为它使用了对流项的中心差分。 因此，有必要增加一个人工耗散条款来稳定该计划。 本文中使用的人工耗散术语基于总变异减小（TVD）模型。 TVD通量限制器被纳入后处理器方法的数值方案中。 通过在原始方案中添加一个额外的步骤，后处理器方法可以轻松应用于许多单步和多步显式方案。 这种方法变得方便，因为不需要修改现有的数值方案。 在广义坐标系中，将TVD模型添加到方程（4 Id）中可以表示为：（4）'K），J4Kl+4-K），。 香港），。 H 2A£ At (42) 2AnL，其中R\textasciicircum{}和R\textasciicircum{}分别由方程（36）和（39）给出。 通量限制函数向量<\&和O是：KL，=feU，teL，teU（43）K），，.=fc），，+i fcU fc），J+i（44）（43）和（44）的向量元素称为通量限制函数，如下节给出。 下标i+\textbackslash\{\},j和i-\textbackslash\{\},j表示所有其他相关变量是根据变量（p、u、v、w、Bx、By、Bx、p、\%x、\%y、r\textbackslash\{\}x和f\textbackslash\{\}y分别在网格点（i + l、j）、（i、j）和（i、j）、（i、j）、（i-l、j）平均）计算的。

258 AGARWAL 13.7.4 Davis-Yee二阶对称TVD模型作为TVD通量限制器的示例，我们在这里描述了Davis-Yee二阶对称TVD模型。 文献中提出了许多模型。 WSUMHD2D有几种TVD模型，例如Harten-Yee模型、Roe-Sweby模型、Davis-Yee模型等。 目前的计算是用Davis-Yee模型进行的。 对于Davis-Yee对称TVD模型，/*对称通量限制函数由以下方式给出：feL,=-\textasciicircum{}U fcU--\textasciicircum{}U\{(°'W'-feUi (45) (<U=-\textasciicircum{}[feU fcU-\{(K\textbackslash\{\}J[WU-fc)„HI (46) 其中(X,\textasciicircum{}) ±. 和\textbackslash\{\}k'\textasciicircum{})..，是对角矩阵D\textasciicircum{}和D\textasciicircum{}的/*特征值。 函数（p（z）是对z的熵校正，M |（z2 + 82）/28 |z|<8，其中0 < S < 0.125，并给出：5！ =8 U，+ 8，=8 1。 U..，V 2,7 2（47）（48）（49）与U =\%xu + \%yv，V = r\textbackslash\{\}xu + r\textbackslash\{\}yv和0.5< 8 < 0.7。 Davis-Yee对称限制器是：fel\textasciicircum{} =™nmo\textasciicircum{}2(a0,4,4,;,2(a'),4,2(a'),.+|,,\textasciicircum{}[(a'),4,, +(a'),+|,]\} (50) fc),,+i = rninmod\textasciicircum{}2(a'),;4,2(a'),H,2(a'),y+|,|[(a'),j4 + (a'),W| (51)

\section{13.8 重要参数}
MHD效应的数值模拟259，其中（a'），+i。和（a'），；+i是<*，，j+i =2\{l\textasciicircum{}\textbackslash\{\}j+i\{QiJ+l-QiJ）/\{jiJ+l + JiJ），（53）和Lk = R?，L„ = J?-\textbackslash\{\} minmod（a，b，c，...n）= S-max[0，min（a，S b，S-c，...，S• n）]和S = sgn（a）。 13.7.5本地时间步进（LTS）对于稳态解决方案，使用本地时间推进来增强收敛是有利的。 LTS背后的主要思想是在不违反稳定性条件的情况下，在每个单独的网格点以不同的时间步骤进行解决方案，即CFL < 1.0。 每次迭代都会更新每个时间步骤。 计算域中的最小时间步进由以下方式给出：. [cFLAfc CFL.Ar|j Af = min\textasciicircum{}-j-|--,-j-:-- \}。 Ii \textbackslash\{\}K\textbackslash\{\} \textbackslash\{\}K\textbackslash\{\} li •\textasciicircum{} I \textasciicircum{}Imax I Umax jj 13.8 显著参数 电磁场对激波的影响可以通过以下参数来最好地描述：焦耳加热焦耳加热（Jh）有助于冲击的熵增加。 如果等离子体具有完美的导电性，则可以忽略它。 它与流量的总焓有关。 焦耳加热的表达式由Jh= Jj2/adx给出，其中J是电流密度，a是可变电导率霍尔参数霍尔参数与霍尔电流相关联，霍尔电流在导电流体中带来不稳定和扰动。 霍尔参数（Q）的表达式是Q = JB！ 2

\section{13.9 入口超声速流动的数值模拟}
260 AGARWAL MHD相互作用参数MHD相互作用参数决定了实现有意义的电导率所需的磁场强度。 它被定义为入口中超声速流动的13.9数值模拟。 Park等人[12]对采用磁流体动力学（MHD）能量旁路概念的scramjet推进系统进行了理论计算。 燃烧室上游的MHD发生器减缓了流量，使燃烧室入口处的马赫数保持在指定值以下。 理论计算表明，在火箭发动机的飞行速度下，MHD旁路方案可以提高比传统scramjet的比冲程。 目前的计算模拟了Park等人[12]的理论计算，用于集成MHD发电机的scramjet入口。 下表1给出了当前模拟中使用的物理模型与Park等人[12]使用的理论模型之间的差异。 表1：物理模型采用Park等人的理论计算。 [12]一维分析无黏流动没有激波和黏性耗散均匀轴向电流密度热平衡当前计算2D数值模拟黏性流动存在激波和黏性耗散不均匀轴向电流密度热非平衡本研究中使用的流动模拟参数如下表2所示：

MHD效应的数值模拟261表2：流动模拟参数参数流体力学压力电子压力气体温度电子温度速度马赫数 雷诺数磁性雷诺数Prandtl数气体比热比电子比热比电流密度跨电极电压电压横向电压梯度轴向电压梯度磁场霍尔参数磁相互作用参数电导率电率电点值5.372 x 103帕斯卡41.9帕斯卡3592 °K 3592 °K 3240 m/s 3.072 19.8xl0b 0.306 0.72 1.4 1.67 9.16xl04A/m2 1596 V 24,650 V/m 5,000 V/m 8.01 T 4.056 2.865 75.43 mhos/m 7.8 x 10"5 图1显示了scramjet入口的各种流动变量（总压力、气体温度、电子温度、电导率、电子数密度、霍尔参数、电流密度、轴向电压梯度和横向电压梯度）的轮廓。 所有流量变量的轮廓中都有一些共同的特征。 从图1可以看出，冲击是由入口处入口的两个角产生的，它们在入口的颈部后立即融合。 此外，与没有电离的流动相比，这些电离的冲击在流动中略弱。 图2显示了scramjet入口不同垂直部分的各种流动变量（总压力、气体温度、电子温度、电导率、电子数密度、霍尔参数、电流密度、轴向电压梯度和横向电压梯度）的变化。 图2（a）和图2（b）分别显示，与非电离气体相比，微电离气体的总压力和气体温度略有下降，这是一个预期结果。 由于从壁向流体传热，气体温度在壁附近区域最高。 入口核心区域的气体温度更接近壁温度。 另一方面，电子温度在入口颈部下游的核心区域增加[图2（c）]。 如图2（d）所示，导电率在入口附近的墙壁区域最高，下游大幅下降，这反映了流场中电子数密度变化，如图2（e）所示。 其他数量，如霍尔参数（图2（f））、电流密度

262 AGARWAL（图2（g））和轴向电压（图2（h））表现出相似的行为，因为它们直接依赖于电导率。 图3显示了沿scramjet入口中心线的各种流动变量（总压力、气体温度、电子温度、电导率、电子数密度、霍尔参数、电流密度、轴向电压梯度和横向电压梯度）的变化。 除电子数密度外，所有变量的中心线图中都有一些共同的特征。 除电子数密度外，所有变量都达到最大值，而电子数密度在冲击合并的位置达到最小值。 在冲击的背后，变量的值急剧下降，最终变得几乎恒定。 从气体温度（7\textasciicircum{}）和电子温度（Te）的图中可以观察到，自Te»Tg以来，热不平衡在冲击之前就存在，它随着Tt>Tg在冲击后向热平衡（未完全实现）而放松。 从图3（a）可以看出，由于电离，总压力峰值降低。 图4显示了当前数值模拟与Park等人[12]沿轴向的理论计算之间的各种流动变量（总压力、气体温度、流速、马赫数、轴向电压梯度、霍尔参数和电流密度）的比较。 数值模拟中所有流动变量的轴向变化是通过沿着每个部分的垂直方向积分值来获得的。 在数值模拟和理论计算中，霍尔参数以外的所有流动变量都随着轴向距离的增加而减少，而霍尔参数随着轴向距离的增加而增加。 从流量变量中可以观察到，一维理论分析的值和二维数值预测之间存在一些差异。 图4中显示的所有流动变量的轴向变化已沿轴向方向进行积分，积分值如表3所示。 从表3可以看出，所有流量变量的理论分析和数值预测之间所有参数的百分比差异都在5\%以内。 这可以被认为是数值模拟和Park等人的理论分析之间的合理一致性。 [12]。 表3：一维理论分析和二维数值模拟参数表面压力（帕斯卡）气温的比较。 （°K）流速（m/s）马赫数轴向电压梯度（V/m）霍尔参数电流密度（A/m2）Park等人。 计算。 [12] 4.65x10" 2.60 4.45 7.65xl04 当前模拟 4.48 xlO6 2.59 4.23 7.303x10" \% 差异 +3.67 -1.65 -4.94 +0.176 +4.89 +4.95 +4.54

\section{13.10 结论}
\section{13.11 致谢}
\section{参考资料}
MHD效应的数值模拟 263 13.10 结论 采用双温模型和可变电导率模型的二维电磁流体动力学方程来计算磁场（由物体表面的磁铁产生的）对scramjet入口内弱电离气体的高超声速流动的影响。 计算与1D理论分析进行比较；得到了合理的一致性。 计算还表明，通过在弱电离等离子体上施加电磁场来减少scramjet入口超音速阻力的潜力。 13.11 致谢 这项工作得到了空军科学研究办公室（AFOSR）的支持。 作者感谢博士的支持和鼓励。 伦·萨克尔。 参考资料1。 米辛，G。 I.，Serov，Yu。 L。 \& Yavor，我。 P.，在气体放电等离子体中超音速移动的球体周围的流动，Sov。 科技。 物理。 Lett.，卷。 17，不。 1991年6月6日。 2. 米辛，G。 I.，弱电离等离子体中的Sonic \& 激波传播，气体动力学，Yu。 我。 Koptev编辑，新星科学出版社，1992年。 3. 戈尔迪耶夫，V。 P.，克拉斯林科夫，A. V.，拉古丁，V. 我。 \& Otmennikov，V. N.，利用等离子技术降低超音速阻力的可能性的实验研究，流体力学，卷。 31，不。 2，1990年，pp. 313-317。 4. 约瑟夫，W。 C，对钝化AMS的压力分布和流场的实验和理论调查，美国宇航局技术报告，第 D-2969，1965年。 5. 齐默，R.，W. 和布什，W。 B.，磁场对弓形冲击距离的影响，物理评论信，卷。 1，1958年，p。 58. 6. 雷斯勒，E。 L。 和西尔斯，W。 R.，磁空气动力学的前景，航空航天科学杂志，卷。 25，1958年，p. 235。 7. 利维，R。 H.，Gierasch，P. J。 \& 亨德森，D. B.，带或不带钝体的高超音速磁流体动力学，A1AA杂志，卷。 2，1964年，p. 2091。 8. 坎特罗维茨，A。 R.，对极速飞行中发生的物理现象的调查，Proc。 会议。 关于高速飞行员，纽约，1955年，p. 335. 9. Mnatsakanian，A。 克。 \& Naidis，G. V.，关于MHD对地球大气中航天器运动的影响，研究报告92/11，IVTAN，1991年，p. 26.

264 阿加尔瓦尔 10. Bityurin，V。 A. \& 伊万诺夫，V. A.，MHD发电机的替代能源利用，第33届MHD英语方面研讨会，田纳西州图拉霍马，1995年，第X-2.1页。 11. 科尔，J.，坎贝尔，J. \& Robertson，A.，火箭诱导的磁流体动力学高级推进概念，AIAA论文编号 95-4079，1995年。 12. 朴，C，博格丹诺夫，D。 和Mehta，U。 B.，无摩擦MHD旁路Scramjets的理论性能，美国宇航局报告编号 NAS2-14031，1999年。 13. Bityurin，V。 A.，Lineberry，J. T.，利奇福德，R. J。 \& 科尔，J. W.，AJAX推进概念的热力学分析，AIAA论文编号 2000-0445，2000。 14. Chase，R。 L.，梅塔，U. B.，博格达诺夫，D. W.，Park，C，Lawrence，S. L.，Aftosmis，M. J.，Macheret，S. \& Shneider，M.，对MHD能源旁路引擎驱动的Spaceliner的评论，AIAA论文编号 99-4975，1999年。 15. 库拉诺夫，A。 L。 \& 谢金，E. G.，MHD控制提高Scramjet性能的潜力，AIAA论文编号 99-3535，1999年。 16. 劳克斯，C。 O.，Yu，L.，Packan，D. M.，Gessman，R. J.，皮耶罗，L. \& 克鲁格，C. H.，双温空气等离子体的电离机制，AIAA论文编号 99-3476，1999年。 17. 鲍威尔，K。 G.，磁流体动力学的近似Riemann 求解器，ICASE报告编号 94-24，1994年。 18. Brassier，S。 \& Gallice，G.，磁流体动力学双温模型的蝾蝾方案，第十六国的Proc. 关于Num的Conf. 流体力学方法，法国阿卡雄，布鲁诺（编辑），1998年。 19. 德布，P。 \& Agarwal，R. K.，超音速钝体流的双温模型对可压缩黏性MHD方程的数值研究，AIAA论文编号 2000-0449，2000。 20. 阿加瓦尔，R。 K。 \& Augustinus，J.，可压缩黏性MHD流动的数值模拟，以减少钝体的超音速阻力，AIAA论文编号。 99-0601，1999年。 21. Pai，Shih-L，磁气动力学和等离子动力学，Springer-Verlag，维也纳，1962年。 22. 考林，T。 G.，磁流体动力学，Interscience出版社有限公司，伦敦，1957年。 23. 坎贝尔，A。 B.，等离子物理学和磁流体力学，McGraw-Hill Co. 公司，纽约，1963年。

MHD效应的数值模拟265图。 1（a）无电离的总压力轮廓”（N/m）图。 1（b）带电离的总压力轮廓（N/m2）图。 1（c）没有电离的气体温度轮廓（°K）图。 1（d）带电离的气体温度轮廓（°K）

阿加尔瓦图。 1（e）电子温度轮廓（K）图。 1（f）电导率轮廓（mhos/m）图。 1（g）电子数密度轮廓图。 1（h）霍尔参数轮廓

MHD效应的数值模拟267图l（i）电流密度轮廓（A/nr）图l（j）轴向电压梯度轮廓（V/m）\textasciicircum{}。 \textbackslash\{\}\textasciicircum{}，图l（k）无电离的横向电压梯度轮廓（V/m）图。 1（1）带电离的横向电压梯度轮廓（V/m）

268 AGARWAL 1e+7 2e+7 3e+7 3e+7 总压力（N/m1）图。 2（a）总压力的变化0.25 - 0.20 - 0.15 - X«0.2 n 400 600 800 1000 1200 40 50 60 70 80 电导率（mhos/m）图。 2（d）电导率的变化 3400 3600 3800 4000 4200 4400 4600 4800 5000 气体温度（°K） 800 1000 1200 3400 3600 3800 4000 4200 4400 4600 4800 气体温度（°K）3400 3600 3800 4000 4200 4400 4600 4800 5000 5200 气体温度（°K）图。 2（b）气体温度变化50 60 70电导率（mhos/m）图。 2（d）电导率的变化

MHD效应的数值模拟 269 0.8446 0.8448 0.8450 0.8452 0.8454 0e+0 4e+5 8e+5 1e+6 2e+6 2e+6 0.8447 0.8448 0.8449 0,8450 0.8451 电子数密度x (10"20) 图。 2（e）电子数密度的变化 2e+4 3e+4 4e+4 5e+4 6e+4 7e+4 8e+4 9e+4 1e+5电流密度（A/m2）图。 2（g）电流密度霍尔参数的变化图。 2（f）霍尔参数的变化 1000 2000 3000 4000 5000 6000 轴向电压梯度（V/m）图。 2（h）轴向电压梯度的变化

270 AGARWAL 6000 13000 18000 23000 28000 33000 横向电压梯度（V/ml 图。 2（0横向电压梯度的变化Fi9' 3<b>电导率的变化，电3 w a电流密度和沿轴向距离的霍尔参数图。 3（a）总压力、气体温度和电子温度沿轴向距离的变化i 5e+6 -, f 4e+6 - I 3e+6 - c 2e+6 - a 1e+6 - O r 45000 *i| 40000 - " 35000 = 30000 - S 20000没有ionteat'on 图3(c) 电流密度、横向电压梯度和轴向电压梯度沿轴向距离的变化

MHD效应的数值模拟 271 jj 5.0e*5 - £ 4.8o*5 - I 4.6e+5 - S °- 4.4e+5 - •§ 4.2e+5 - 0 1 «S 3720 | 3640 | 3560: 3480 0 1 jj 3000 - - 2800 - •2 2600 - | 2400 - 0 1 -。 \textasciicircum{}J - \textasciicircum{} Park'aCalc。 [12] Praaant Ctle。 \textasciicircum{}\textasciitilde{}-\textasciicircum{}-\textasciitilde{}轴向Dlatanca（m）Park'aCalc。 [12] 普拉桑特计算。 2 3 轴向 Diatanca (m) Park's Calc. [12] PiawntCale。 /'/ 2 3 轴向对坦卡（m 轴向距离（m）图。 4（c）当前模拟和Park计算之间的电流密度比较[12]图。 4（a）当前模拟和Park计算之间的表面压力、气体温度和流速的比较[12]轴向距离（m）5.1-；4.9-|4.7-=4.5i 4.3 i 4.1-3.9--rjllllll\textasciicircum{}-Parti的计算。 [12]现在计算。 轴向距离（m）图。 4（b）当前模拟和Park计算f12l之间的马赫数、轴向电压梯度和霍尔参数的比较。

\clearpage
\chapter{14 计算磁流体空气动力学的进展}
\section{14.1 导言}
计算磁流体空气动力学 Joseph S的进展。 Shang1 Patrick W. Canupp2 Datta V。 Gaitonde3 14.1 引言 在跨学科建模和模拟技术的发展中，电磁现象作为额外的流场控制机制。 事实上，四十年前，Resler和Sears认识到电磁效应在空气动力学中的潜在应用。1他们的观察得到了Bush、Ziemer和Meyer的开创性研究的大力支持。2-4磁空气动力学确实是一项跨学科的努力；相互作用的物理现象需要空气动力学、电磁学、化学物理学和量子物理学的相互作用来描述电磁场存在下的电离气体流动。 这种跨学科的努力不仅提出了极其复杂的科学问题，而且还需要大量的知识库。 然而，技术突破的前景太大了，不容忽视。 众所周知，电磁力可以显著改变导电介质的流场。5-7对于大多数高超音速飞行，由激波和飞行器所界定的空气混合物由高度1名高的1名高阶科学家组成，卓越计算科学中心，空军研究实验室，赖特-帕特森空军基地，俄亥俄州45433-7913。 2 访问科学家，计算科学卓越中心，空军研究实验室，赖特-帕特森空军基地，俄亥俄州45433-7913。 3 高级研究航空航天工程师，计算科学卓越中心，空军研究实验室，赖特-帕特森空军基地，俄亥俄州45433-7913。 Frontiers of 计算流体动力学 - 2002 编辑：David A. Caughey和Mohamed M. 哈菲兹 ©2002 世界科学

274 SHANG、CANUPP和GAITONDE激发了内能模式。 当空气温度超过5000°K时，一小部分解离分子会失去电子。8然后，电离空气混合物的特点是电导率的有限值，根据飞行速度和高度，电导率可能超过100 mho/m的值。 移动的带电粒子和电磁场的相互作用将产生洛伦兹力，作为放松传统空气动力学限制的额外机制。 Resler和Sears研究了一个无量值参数，该参数决定了洛伦兹和动力学力的相对大小。 洛伦兹力似乎在高海拔和高速下占主导地位。1一般来说，这种流动条件在高超音速飞行中自然发生。 洛伦兹力的可能用途之一是控制气体在亚音速或超音速下连续加速或减速而不窒息，即使在恒定的横截面通道中也是如此。 通过焦耳加热导电介质的速度和温度的耦合使这种流场操纵成为可能。1'5'6带电介质的流动方向也可以通过霍尔电流和磁场中电子轨迹的曲率的内在关系来操纵。 霍尔效应已被广泛用作驱动力，以转移带电流体颗粒的方向以增强燃烧或减少入口压缩和失真的长度。 洛伦兹力总是有一个分量，a（U x B）x B，它将使流动减速。1'4在剪切层的非均匀速度分布中，减速力随着距离固体表面的距离而增加。 一般来说，最终结果是剪切层的速度和温度梯度降低。 速度和温度梯度的扁平化减少了与接触表面的传热和皮肤摩擦。1'4'5在某些情况下，湍流流体颗粒穿过磁感应线引起的滗流与磁场相互作用，并倾向于抑制湍流运动。 即使在哈特曼数的中等值下，湍流的壁剪切应力也会降低。6然而，激波在等离子体场中的传播和分叉揭示了导电气体介质中最明显的现象。 最引人注目的观察是弓形冲击的对峙距离在钝体.2上急剧增加。这种行为与所有没有遇到电磁场的非平衡高超声速流动的实验数据和计算形成鲜明对比。8在高超音速钝体流动中，激波距离距离随着马赫数和非平衡效应的增加而减少。 1959年，齐默通过将磁场应用于非平衡高超音速钝体流动来增加激波距离。2他的实验观察得到了布什研究的理论验证。3

\section{14.2 控制方程}
计算磁流体空气动力学 275的进展 最近在修改导电流体高速流动方面的创新涉及空气动力学和电磁场相互作用，如AJAX概念所示。9,10关于激波在20世纪80年代在俄罗斯恢复的等离子体中传播的研究，11-13，而Ganguly等人14最近研究了这种现象。 他们的集体发现表明，非平衡、弱电离气体中的过程大大改变了移动的激波。 具体来说，当激波在等离子体中传播时，振幅会减少，波锋也会分散。 从本质上讲，这些发现具有复合波运动的所有特征。 如果可以在激波船头的上游引入导电介质并量化分叉，这些新发现的物理现象将具有高速飞行的革命性潜力。 尽管在确定控制可压缩流动等离子体修饰的主要机制（例如，热力或磁力效应是负责）方面存在一些争论，但磁空气动力学方程的框架描述了电磁对流动的一阶效应。 研究这种效应是当前努力的主要重点。 14.2 控制方程 电流感兴趣的等离子动力学的一个独特特征是介电常数相对较小。 即使物理现象在微波频率范围内，与传导电流相比，位移电流仍然可以忽略不计。5'6这种简化将广义安培电路定律的制定简化为泊松类型。 对于目前的调查，不包括过于复杂的非平衡化学动力学，控制方程组简化为：| + V.（ \textasciicircum{}) = 0 (14.1) as at d(PU) dt + V • \{UB - BU) = -V x 1 (\textasciicircum{} B - V x - a \textbackslash\{\} n + V TTTT BB ( B-B, r = V-r (14.2) (14.3) 9(PZ), \textasciicircum{} at \textasciitilde{}t v pZU + U-(p+! ff)l-2(B-U) V-[q + U-r] Vx *(**!) （VxB）2（14.4）

276 SHANG, CANUPP \& GAITONDE，其中p表示磁透射性（磁通密度B与磁场强度H有关），pZ = pe+ \textasciicircum{}-（14.5）经典理想磁流体动力学（MHD）控制方程可以通过其他假设从磁空气动力学系统中推断出来。 首先，无限电导率的概念意味着运动诱导的磁场的强度压倒了施加的磁场的强度。5'6其次，在许多流动中，惯性效应大大超过了介质中的黏性耗散和传热。 第三，介质被认为是各向同性的。 然后，只需设置Eqs的右侧，即可轻松获得生成的控制方程。 (2)-(4)到零。ft+V-\{pU) = 0 \textasciicircum{}• + V-(UB-BU) = 0 mi + v.[puuu-\textasciicircum{} + (P+\textasciicircum{})i] = o (14-6) gt -TV iHv V,, T\textasciicircum{}T 2(J d(PZ) at +V< pZU + U-(p+\textasciicircum{})l-B BU = 0 上述理想的MHD方程构成了双曲偏微分系统。 通过对每个空间时间空间中的控制方程的分析，发现了特征向量和特征值。15因为洛伦兹力垂直于磁场强度和电流，因此与磁场强度的法向分量相关的特征值为空值。 经典MHD方程的其余七个特征值也可以局部退化，相互重合，这取决于磁场的相对大小和方向。 七个特征值是u、u±Ca、u±Cs和u±Cf。当扰动在等离子体场中传播时，这些特征值反映了四种不同的波速：声学、Alfven、慢速和快等离子体波速度，5'7 '2 \_ s <2 \_ C2 = dpjdp Cl = B2Jpp = \textasciicircum{}2+s+\textbackslash\{\}/(c2+s)2- -4C*C2 (14.7) (14.8) (14.9) (14.10)

\section{14.3 数字程序}
进展在计算磁流体空气动力学 277中，其中Bn表示磁通密度相对于波前的横向分量。 因此，并非所有在均匀等离子体中传播的波速都是各向同性的；有些波速在很大程度上取决于磁场的极性。 图中的图表。 1展示了波前和施加/诱导磁场之间从0°到180°的相位关系上的四种等离子体波速度。 该计算假设等离子体场均匀，并将所有波速与声速C相对归一化。 参数B2/（ppC2）可以用磁性相互作用参数S = LaB2/' pu和磁性雷诺数，Rm = apuL表示如下：5 \textasciicircum{} = ftT（14J1），其中M是马赫数。 对于等离子室中发光放电产生的电磁场，6参数B2/（ppC2）的值接近1。 为了说明目的，理论结果是使用0.5的值生成的。 在这种情况下，快速等离子体波可以在横向磁场中以声速的1.225倍传播。 在整个相位角范围内，快速等离子波的速度大于或等于音速。 另一方面，缓慢的等离子体波在整个相位角范围内以低于音速的速度传播。 当波运动与磁场平行时，慢等离子波的最高速度是声速的0.7071。 当横向磁分量消失时，缓慢的等离子体波运动停止。 快和慢等离子波速度代表阿尔夫文波和声波传播速率的上限和下限。 由于阿尔夫文波不携带能量，因此没有熵产生与其他波运动模式的相互作用。7图2描绘了马赫数为10时发光放电的波速。 现在，快速等离子波的速度比音速快100倍，而慢等离子波的速度仍然受到声速的限制。 鉴于波速的巨大差异，波干涉仅限于慢等离子波和声波之间的干扰。 然而，必须从详细和系统的实验和计算研究努力中得出对这种现象的全面理解。 14.3 数字过程笛卡尔框架中的控制方程以通量向量形式显示为：dV\_ OF* dFy dFz \_ dFx,r dFy,r 3FZ)r dt dx dy dz dx dy dz

278 SHANG, CANUPP \& GAITONDE - 快等离子波慢等离子波阿尔夫文波声波 -V- 90.0相位角图1波传播速度，SM2/Rm = 0.5，其中V是保守依赖变量的向量。 与理想MHD方程Fx、Fy和Fz相关的通量向量的分量，以及与电阻率、黏性耗散和传热、Fx\textasciicircum{}r、Fy、r和Fz>r相关的通量向量的分量可以从Eqs中轻松获得。 （L）-（4）。 这里只给出x坐标方向的分量：16 V = (p, pu, pv, pw, pZ, Bx, By, Bz) (14.13) Fx = pu pv? +p + B-B/2p-B2x/p, puv - BxBy/fi puw - BxBz/p, \{pe+p + B- B/n) «+([/"• B/n) Bx uBy - vBx uBz - wBx (14.14)

进展 计算磁流体空气动力学 279 10 3| 10 2| 10 *| Fast Plamia Wav« \textbackslash\{\} / 10-1; Ml/ 10'2 1 10 "3 "I 90.0 相角图 2 波传播速度，SM2/Rm = 104 • xy \{U-T + KdT/dx +By [d\{Byln)/dx - d(Bx/fi)/dy] liio \textasciicircum{}i0j +BZ [d\{Bz/ii)ldx - d(Bx/fi)/dz] lixo\} [d(Bv/ri/dx-d(Bx/ii)/dy]/v [d(Bz/v)/dx - d\{Bx/ii)ldz\textbackslash\{\} /a 理想的MHD的方程不是严格的双曲系统17,18，因为不同的特征值在场中退化可能与其他特征值 因此，弱解的分析结构在退化发生的地方是未知的。 然而，解决该系统的大多数数值方法是基于特征的、近似的黎曼，17-21或TVD方案。22在这些方法中，控制方程的守恒形式的双曲分量（理想的MHD）表现为：16'19,20 8t + V-F = 0（14.16）为了提高数值稳定性和准确性，理想的MHD通量向量根据其各自特征值的符号进一步分割。 特征值和特征向量来自雅各布系数

280 SHANG、CANUPP和GAITONDE矩阵、原始变量中的dFx/dW、dFy/dW和dFz/dW。dV dFxdW dFvdW dFz dW 1 1 1 =0（14 17）在dW dx dW dy dW dz \{ '为了克服退化特征值的数值难度，几位研究人员修改了雅各布系数矩阵。 Powell et a/.19成功地将高斯定律V • B = 0作为此目的的修饰词。 最近，MacCormack20表明，通量向量可以被扩展，以保持一度的均质性。 他的方法采用以下方法。 首先，一维MHD方程的离散形式是\textasciicircum{}Vi+£F'e.5 = 0（14.18）边，其中5是细胞界面表面积，V»是第i细胞的体积。 引入原始变量和守恒变量S和S\textasciitilde{}l之间的转换雅各布，该方程也可以写成\textasciicircum{}Vi+£W\textasciicircum{}W5 = 0（14.19）边。接下来，MacCormack将重新定义的理想MHD通量向量拆分为F'e = F'e' - P-\textasciicircum{}（14.20），其中a是一个加的守恒变量，取为单位，因此Eq. （14.19）的形式是\textasciicircum{}Vi+\textasciicircum{}5-\textasciicircum{}5U.5-X;P\textasciicircum{}(f)5U.5 = 0 (14.21) sides sides Eq中的最后一个和。 （14.21）是相同的零，因为数量Bx/a相对于守恒变量是零度的齐次。 矩阵A'e作为一个子矩阵包含Powell使用的Jacobian 矩阵。 因此，这种新方法使用与Powell等人19所示的相同特征系统来评估这部分通量。 由于A'e的最后一列，产生了额外的通量成分。 数值方法使用该通量分量的中心评估。 因此，离散方程显示为\textasciicircum{}Vi+ Y\textasciicircum{} S-1C-1ACSU-S+ \textasciicircum{}b-5 = 0 (14.22) 边

计算磁流体空气动力学 281的进展，其中b向量是rp（，>B2。 „。 B2u uBT vBx WBT\textbackslash\{\} b= (0,(\textasciicircum{}-1)\textasciicircum{},0,0\textasciicircum{}-1)\textasciicircum{}-,-Y.-\textasciicircum{}.-V) (14-23) 等式中第一个求和中的通量分量。 （14.22）根据位于矩阵A的对角线上的特征值的符号进行分割。 因此，拆分通量为F± = S-1C\textasciitilde{}1A±CSXJ（14.24）。上述描述假设控制方程的一维形式，但它很容易扩展到多个维度。 MacCormack的原始论文实际上在三维中发展了该方法，从而为求解理想MHD方程的数值技术做出了重大贡献。17-20人们还可以使用有限差分fv + «w + ".-) + «w + jy) + w; n)的形式编写通量分割方法。 0（14.25）St ox oy oz 另一个用于求解磁空气动力学控制方程的创新数值程序分别采用紧凑差分23和经典Runge-Kutta方法的组合，用于空间和时间导数。 这种方法成功应用于求解时间依赖的三维Maxwell24'25和Navier-Stokes 方程.26'27 Gaitonde成功实现了这种求解磁空气动力学方程的高阶方法。16对于紧微分方案，Lele给出的以下公式近似了标量变量的£导数，</> "(\textasciicircum{})i-l + \{4>i)i + Ot\{<t>i)i+\textbackslash\{\} - 3 (\&+2 - <Ai-2) (14.26) + 2\textasciicircum{}t2l(\textasciicircum{}+1-\textasciicircum{}\_1) 四阶，Pade型紧微分是通过将系数分配为1/4的值来恢复的，是当前数值方法的选择。 当前方法的关键组成部分是低通滤波器。24-28 盖通德在1996年引入了滤波器，作为求解偏微分方程的紧凑微分方案的实际应用的关键元素。28 他证明了空间滤波器可以通过抑制通过网格非均匀性、边界条件实现和非线性现象捕获产生的高频、虚假的傅里叶分量来保持复杂几何图形的计算稳定性。 过滤器的一般配方，基于Lele的

\section{14.4 兰金-胡戈尼奥特跳跃条件}
282 SHANG、CANUPP和GAITONDE框架，是：0\#\_！ + 4>i + M'i+l = hn\{P, 4>i-n, •••, 4>i+n) (14.27) 其中参数/? 必须在-0.5和0.5之间完全有界，上标'表示过滤变量。 右侧多项式的系数可以通过参数（3通过泰勒级数和傅里叶级数分析得出。16滤波器在2n+l节点的模板上运行，从而产生2n阶公式。23'25在应用中，滤波器充当计算解的后处理器。 对于多维计算，过滤器的应用是顺序的和交替的坐标方向。 近似黎曼或紧差分方案的公式通过引入坐标变换，£ = £(aj,y,z)、rj - r)(x, y, z)和£ = ((x,y, z)，轻松扩展到曲线框架。 结果方程是：3V' OF' OF' dF' 8F' 8F'r 8F'r 8F'r -ar + \textasciicircum{} + W + \textasciicircum{} = -\# + \textasciicircum{}f + \textasciicircum{}f (14-28)，其中变换的因变量定义为V - V/Jc，Jc是坐标变换雅各布。 与理想MHD方程相关的通量向量分量以及电阻率、黏性耗散和传热项以类似于曲线框架Ft = \{ZxFx + tvFv + ZzFz)/Jc F'r, = (r)xFx + r)vFy + VzFz)/Jc (14.29) F's = (((XFX + CvFy + CzFz)/Jc 14.4 Rankine-Hugoniot Jump条件电磁力对激波结构的影响可以通过在冲击中修改Rankine-Hugoniot 间断关系来检查。 等离子体中沿着x坐标的正常冲击跳跃条件是从磁空气动力学方程中推导出来的。 跳跃条件变成：29

计算磁流体空气动力学 283的进展 [Jy/a - [Jzl<J - [pu] [p + Pu2] [puv] [puw] [puh] [Bx] - wBx + uBz] - uBy + vBx] J(JyBz - JzBy)dx J(JZBX - JxBz)dx f(JxBy - JyBx)dx fj2/adx (14.30) 其中h表示流动的总焓，括号表示激波的最终状态的变化。 如果等离子体具有完美的电导率，而霍尔电流可以忽略不计，上述跳跃条件将减少到Sutton和Sherman5得出的。 在这些条件下，J = k-\textasciicircum{} \textbackslash\{\}\textasciitilde{}！ L) \textasciitilde{} J\textasciitilde{}3x\textasciitilde{} (\textasciicircum{})' an\textasciicircum{} tne 结果的跳跃条件是通过与x的直接积分获得的。[pu] [p + pu2\} [puv] [puw] [puh] [Bx] [uBz - wBx] [vBx - uBy] -[(B2y + B2)/2p] [BxBy] [BxBz\textbackslash\{\} [-u(B2y+B2)/2p, Bx(vBy + wBz)/p] (14-31) 上述冲击跳跃条件，无论基本假设如何，表明电磁场的存在改变了气体动力学的Rankine-Hugoniot关系。 修改的要素在于两个额外的熵变化机制。29 其中之一是焦耳加热，它是积极的，将有助于整个冲击的熵增加。 另一个机制是电磁力可以在移动的气体粒子上执行的工作。 根据感应或施加的电磁场的极性，开放系统的总熵可以降低。 作为熵减少的直接结果，激波的波阻将减少。 为了深入了解磁场对导电流体流动的影响，本研究在马赫数为1.5的横向电磁场中检查了正常激波的结构。 为了说明基本概念，假定等离子体介质是一种热量和热完美的氦，具有恒定的电导率。 施加的横向磁场的自由流值为By = 0.03 T，Bx = Bz = 0，电导率为105 mho/m。 10 Torr的压力和300 K的温度完全指定了自由-

284 SHANG、CANUPP和GAITONDE气体流状态。 在这些条件下，在一维流动的假设下，控制方程减少到以下形式iUH pu2 +p + 2/u. puh + "rc\textasciicircum{}re By -UT - q A JLiB\textasciicircum{}\_uB aji dx "•\textasciicircum{}y (14.32) Canupp应用Runge-Kutta-Fehlberg方法求解这些具有可变步进大小的耦合常微分方程。30为了确保数值稳定性，在亚音速、后冲击区域启动计算，并在逆流方向向方向上集成。 通过用激波 雷诺数缩放坐标捕获精细冲击结构，x* = Poouooox/fJ-oo-通过与Muntz和Harnett的数据达成出色一致，对气态计算（无磁场）求解方案的验证。31为了比较电磁场和声场中的冲击结构，图。 3显示了两种情况的数值计算结果。 流动从左到右，图中显示了马赫数和熵通过冲击的变化。 对于每种情况，计算出的马赫数和熵变化都正确到达x* = -7500附近的相同上游渐近线，尽管该图的水平轴仅延伸到x* = -3000以突出显示激波内部结构。 气体动力冲击的压缩过程是突然的，并集中在薄薄的冲击结构中。 另一方面，等离子波的信号在气体动力学冲击的上游传播，仿佛压缩分两个不同的阶段发生。 对于相同的迎面而来的马赫数，等离子体中的冲击相对于气态冲击向前移动。 如图所示。 3表明，气体动力学冲击中的气体熵在声点迅速上升到最大值，然后迅速放松到最终的下游值。 薄薄的冲击结构限制了此案例的所有熵变化。 等离子体中休克的熵表现出类似的整体行为，但变化更渐进，不那么强烈。 熵的增加也会在电磁场中比其气动力学对应物更上游地启动电击。 对于这两种冲击，熵在冲击内的声波点上升到最大值，然后渐近地放松到不同的冲击后值。 对于这个计算，等离子体中冲击的熵增量比气动力学冲击低2.6倍。 因为等离子体中冲击的熵变化在很大程度上取决于横向磁场的极化，这

\section{14.5 理想的MHD 激波管模拟}
进展 计算磁流体空气动力学 285 -10.06 - 0.05 - 0.04 - 0.03 • I - 0.01 - 0 - -0.01 ";3ooo -2000 -1000 0 图3 声学和等离子冲击结构简化的一维冲击计算并不能说明所有可能的磁空气动力学现象。 然而，它的效用在于它的简单性和展示气态力学中基本磁场效应的能力。 14.5 理想的MHD 激波管模拟虽然理想化的MHD 激波管不一定可以通过实验再现，但它已成为开发求解理想MHD17'21,22'34和磁空气动力学方程的数值程序的基准。16处理退化特征值的困难或MHD中中间（化合物）冲击的可接受性仍未完全解决。17'18基于特征值和特征向量分析的近似Riemann 求解器产生的数值结果，在未来不可避免地将受到额外的审查。 与基于特征和紧差分方案的计算直接比较是有价值的，这些计算不需要对特征值和特征向量有具体知识。 气态动力学和MHD 激波管模拟的并排比较对于评估激波在等离子体中的突出特征更有用。34为了实现这一目标，通过紧凑差分和通量向量分割方案进行了计算，复制了Brio和Wu17的开创性努力的网格和时间步长大小。16'34总共使用了800个网格点来定义计算域，恒定的空间和时间增量分别为Ax = 1和At = 0.2。 对于1.5 1.4 1.3 i！ 1.2 E z 1.1 8 1 s 0.9 o.s 0.7: „' r ""-\textasciicircum{} M, B. =0 M, B =0.03 (s-sj/c„ B.=0 (s-s\_)/cp, B„=0.03 N \textbackslash\{\} \textbackslash\{\} \textbackslash\{\} •... • 1。 >..

286 SHANG, CANUPP \& GAITONDE 1.1 o i- i i i i i i i i i i i i i i i i i 0 200 400 600 800 图4 通量-向量分割（FVS）和紧凑差分差（CDS）方法的比较 紧凑差解程序，需要局部滤波器切换程序将高阶滤波器更改为本地二阶滤波器，以进行冲击捕获。 基本冲击检测公式采用于Harten和Yee32'33的工作，详细的实施情况见参考。 16. 对于基于特征的计算，数值方法是Steger-Warming通量-向量分割的修改形式，使用通量Jacobian的细胞界面值。34在强压力梯度的区域，计算恢复到原来的Steger-Warming通量-向量分割方案。35 激波管内密度分布的比较如图所示。 4. 简而言之，两种完全不同的解决方案之间的一致性非常好。 结果也与Brio和Wu的结果一致。17正如人们所期望的那样，在强冲击跳跃位置的两个数值结果之间存在微小到可以忽略不计的差异。 同样，在气体动力学和MHD 激波之间的差异尺度上，差异几乎无法检测到。 图5描绘了磁场强度的横向分量，By。 横向磁场分量的剧烈变化在计算域上被清楚地表明。 瞬时和瞬态现象受到缓慢冲击向低压部分的右行、快速稀疏前体和向激波管的高压部分的左行、快速稀疏波的约束。 By场分量在接触不连续性和快速稀疏之间发生极化反转。 与此同时，流向

计算磁流体空气动力学 287 1.0 0.5 0.0 -0.5 -1的进展。 带B场的Q-没有B场O u '6 200 400 600 "\$6o图5 磁场强度的MHD 激波管分量中的横向磁场强度，Bx坚定地保持在0.75的恒定值，以确保高斯定律的共平面磁场，V • B - 0被严格执行。36理想化激波管内的密度变化，图。 6个呈现，最能描述复杂波系统。 快速稀疏波以高于声速的速度向管的两端传播。 缓慢向右运行的冲击的前体，被确定为快速等离子波，进入管子的低压部分。 这个波代表了等离子波运动的独特现象之一。 在横向磁场反转其极性的位置，会出现中间或复合电击。 复合波的存在表明，在等离子体场中，缓慢的稀疏波可以暂时附着在缓慢的冲击上。 当横向磁场在极性反转期间消失，不同的等离子体波退化时，就会发生这种独特的事件。 复合波的存在和物理意义仍然不确定。17'18 至少，这些方法证明了数值中间冲击的存在，作为非凸双曲系统的可能弱解。 这个结论与用于解决空间通量的程序无关。16'17'34图7显示了理想化MHD 激波管内的压力分布。 压力变化的行为为声场和等离子体场的不同波速提供了进一步的确认。 左行的快速稀疏波向高压区域移动，然后是缓慢的冲击，紧接着是稀疏。 在......之间

288 SHANG, CANUPP \& GAITONDE 快速稀疏 慢复合接触不连续性 200 400 600 800 图6 MHD 激波管复合波和右行慢冲击的密度分布，压力分布在接触表面具有恒定值。 在激波管的压缩侧，当右行的快速稀疏之前，右行的慢速冲击产生比气体动力学冲击更大的压力跳跃。 图8显示了理想化的MHD 激波管中的流向速度分量u。 再次，左行的稀疏波迅速加速了流动。 复合波首先使通过冲击的流量减速，并在波的尾部立即加速。 通过压缩膨胀过程实现的最大流向速度低于气态动力学激波管计算。 正如理想的MHD配方所假设的那样，在大速度梯度的高度封闭区域，分子耗散过程不会被忽视。 因此，Myong和Roe认为可接受的冲击必须具有黏性轮廓18，可能会从磁空气动力学方程的解中得出强有力的证实。 与气动学激波管一样，u速度分量在接触表面保持恒定值。 然而，慢速休克减速比其气体动力学减速更大，以恢复右行、快速稀疏前体增加的膨胀。 图9显示了理想化的MHD 激波管中的法向速度分量v。 等离子体冲击系统的另一个独特特征是电磁场中额外的涡量生成机制。 在由右行、稀疏前体和左-所界定的区域

\section{14.6 高超音速MHD 钝体模拟}
计算磁流体空气动力学 289 0.8 0.6 0.4 0.2 w0 200 400 600 800 图7 MHD 激波管中运行快速膨胀波的压力分布，等离子波系统诱导大负法向速度分量v。 这一特征与相应的声波系统明显不同，后者使零法向速度分量不受干扰，v - 0。 对于声场来说，众所周知，直的激波不会产生任何涡量。 然而，这一结果表明，两个缓慢的等离子体激波在垂直于平面冲击的方向上产生涡量 kdv/dx的强反向旋转涡流。 直等离子体冲击和波结构的涡量生成不仅提供了电磁场可以改变熵产生机制的明确证据，而且还指向了流场操纵的新机制。 下一节将讨论这种机制在高超音速钝体流动场中的作用。 14.6 高超音速MHD 钝体模拟目前对应用电磁场来修改空中动态性能的兴趣集中在高超声速流动上。 声学和等离子体场中的钝体计算阐明了这种磁空气动力学流场的基本特征。 本分析计算了在5.85的自由流马赫数下，在声学和理想化等离子体场中经过二维圆柱形鼻体（Rn = 0.1 m）的高超声速流动。 有限体积数值方法使用60 x 50椭圆平滑网格系统和MacCormack通量向量分割方案。20'34数值方法是第一-有B域，没有B域

290 SHANG, CANUPP \& GAITONDE 0.8 0.6 0.4 0.2 -0.2 -0.4 0 200 400 600 800 图8 MHD 激波管顺序中的流向速度分量在空间和时间上都准确。 在自由流边界，流量属性在每次计算期间保持不变。 通过无黏性边界条件，身体表面是不可穿透的边界。 用于MHD计算的额外边界条件包括身体表面恒定正态磁通量密度，该密度在介质界面建立了磁场的连续正态分量，n • 6B = 0，以及远场强加的常数值。 为了说明磁场对流动的影响，图。 10显示了两个钝体流动场的计算密度轮廓的并排比较。 对于左边的案例，自由流中不存在磁场。 对于右侧的案例，自由流边界上存在一个均匀的Boo = 5j域。 相同的模拟代码计算了这两种情况的结果。 对于每个流量，溶液残差下降五阶的收敛标准决定了何时停止计算。 溶液向量变化的I/2-范数作为每个时间步骤后溶液残差的保守度量。 图中的密度轮廓。 10表明在这两种情况下都存在分离的弓激波。 结果表明，磁场有效地将弓激波推到钝体的上游。 这一结果与奥古斯丁等人/.22的发现一致，因为目前的分析只是为了说明理想化的电磁场可以改变分离的冲击包络的距离，所以它不寻求对各种强加的磁场对流动的影响的进一步详细描述。 尽管弓的对峙距离有差异激波 o有B场没有B场

计算磁流体空气动力学 291的进展 图9 MHD 激波管的正常速度分量 在这两种情况下，两个波前在整个流场上显示出一些全球亲和力。 每个流在中心线上都包含一个基本正常的激波，然后是身体肩部周围的快速扩张。 气动力学计算的对峙距离与声场的相关实验结果非常吻合，A/Rn分别为0.471和0.442。37从研究的案例来看，等离子体场的冲击包络对气动冲击表现出1.233倍的向外位移。 反过来，弓形冲击的曲率半径比声场的曲率增加了1.103倍。 图11描绘了每种情况下沿对称线的密度分布。 声场（Boo - 0）的数值结果显示，密度根据Rankine-Hugoniot 间断关系跳过激波，随后随着流量进一步压缩到圆柱形鼻的驻点，密度单调上升。 相比之下，沿等离子体场对称线（-Boo = 5j）的冲击后密度分布遵循非单调变化。 在这种情况下，冲击后密度的较低值为激波对峙距离的增加提供了明确的原因。 连续性考虑要求更大的流管区域通过相同的质量流通过冲击层。 无花果的比较。 10和11表明，震后等离子体场密度的降低在身体表面附近迅速发生。 在这个结构中，压力和速度也会下降，而温度逐渐上升。 这种结构的来源不确定，因为它与V • B / 0的区域重合。 违反高斯磁场定律是由于

292 SHANG, CANUPP \& GAITONDE 0.4 0.2 -0.2 •0.4 0.0 0.5 1.0 x(m) 图10 5.85马赫时圆柱形钝体的弓形冲击包络与该区域通量中的数值误差的比较。 为了减轻这些错误，计算采用了参考中描述的松弛技术。 20. 虽然该算法减少了V•B中的误差，但它并没有在激波和停滞区域中完全消除它们。 未来的工作将继续解决这种数值不确定性。 图12显示了等离子体和声场在二维物体表面的静压分布。 变量s测量沿身体表面与驻点的距离。 两个数值结果都以从驻点向下游后体快速膨胀的形式表现出相似的行为。 目前的计算还表明，在整个计算域中，声场的表面压力均匀高于等离子体的表面压力。 然而，数值结果之间的巨大差异并不完全代表等离子体场中可能的波阻减少。 对于包含有限电磁场强度的流场，洛伦兹力必须进入动量平衡考虑，以进行阻力评估。 在理想的MHD公式中，洛伦兹力分为两个术语，通常称为麦克斯韦应力，BB和磁压，B • B/2fi。 更重要的是，该配方还忽略了剪应力张量。 因此，准确且具有物理意义的阻力计算需要采用完整的磁空气动力学方程。 M=5.85

\section{14.7 结束词}
计算磁流体空气动力学 293 6 - 5 - 3 - 2 - 1 - •3.0 -2.5 -2.0 -1.5 -1.0 x/Rn 图11 圆柱形钝体在5.85马赫14.7时沿停滞流线的密度变化结论性评论计算磁流体空气动力学被公认为跨学科技术发展的新前沿。 这一技术要求的一个关键要素是将时域中的计算电磁学与计算流体动力学和计算化学动力学相结合。 这种跨学科努力对高速飞行的影响可能是革命性的。 与非导电流体流动相比，磁场产生的多种波速允许导电流体系统的行为更加复杂。 理想化的一维正态激波计算表明，在存在磁场的情况下，熵在冲击中的熵跳跃会减少，从而证明了钝体波阻还原的理论潜力。 目前的结果还指出，横向电磁场可以立即在下游产生直的正态冲击涡量。 此外，本研究确定了磁空气动力学方程预测的钝体流动的弓形冲击距离变化的机制。 具体来说，冲击背后密度的非单调变化导致更大的流管面积要求。 最后，这项研究发现，通量向量分割方法在V•B中引入了很大的误差。 未来需要解决的问题包括了解V•B中的错误对高超音速钝体的影响

\section{14.8 致谢}
\section{14.9 参考资料}
294 SHANG, CANUPP \& GAITONDE 0 12 3 s/Rn 图 12 圆柱形钝体在5.85马赫计算下的表面压力分布的比较，以及将数值方法扩展到三个空间维度和高阶精度。 14.8 致谢 博士的赞助 A. 纳赫曼，R。 坎菲尔德和S。 感谢AFOSR的Walker。 这项工作部分得到了WPAFB国防部HPC共享资源中心授予的HPC时间的支持。 14.9 参考资料1Resler，E.L.和Sears W.R.，磁空气动力学的前景，J. 航空。 科学，1958年第25卷，pp. 235-245和258。 2Ziemer R.W.，磁空气动力学实验研究，美国火箭协会。 J.，第29卷，1959年，pp. 642-647。 3布什·W。 B.，磁流体动力学-高超声速流动 过去一个钝体，J。 航空。 科学第25卷，1958年11月，pp. 685-690和728。 4Meyer，R.C.，关于通过磁流体力学技术降低空气动力学热转移率，J. 航空。 科学，1958年3月，pp. 561-

计算磁流体空气动力学 295 566和572的进展。 5Sutton，G.W.和Sherman，A.，工程磁流体动力学，McGraw-Hill，纽约，1965年。 6Mitchner，M.和Kruger，C.H.，部分离子气体，John Wiley and Sons，纽约。 1973年。 7Cabanese，H.，理论磁流动力学，学术出版社，纽约，1970年。 8Shang，J.S.，高超声速流动的数值模拟，高超音速空气动力学的计算方法，Comp. 机械。 出版物，英国南汉普顿，1992年。 9Gurijanov，E.P.和Harsha，P.T.，AJAX：高超音速技术的新方向，AIAA Preprint 96-4609，第七届国际太空平面和高超音速会议，Norfolk VA，1996年11月18日至22日。 10Bityurin，V.A.，Velikodny，V.YU.，Klimov，A.I.，Leonov，S.B.和Potebnya，V.G.，激波与脉冲放电的相互作用，AIAA 99-3533，第30届AIAA等离子动力学和激光会议。 6月28日至7月1日，弗吉尼亚州诺福克，1999年。uMishin，A.P.，Bedin，N.I.，Yushchenkova，G.E.和Ryazin，A.P.，气体中激波的异常松弛和不稳定，Sov。 物理。 技术，第26卷，1981年，pp. 1363-1368。 12Basargin，I.V.和Mishin，G.I.，激波 Argon中横向发光放电的等离子体中的传播，Sov。 科技。 物理。 Lett。 1985年11月，pp. 85-87。 13Voinovich，P.A.，Ershov，A.P.，Ponomareva，S.E.和Shibkov，V.M.，空气中纵向流放电等离子体中弱激波的传播，高温。 第29卷，1990年，pp. 468-475。 14Ganguly，B.N.，Bletzinger，P.和Garscadden，A.，激波非平衡低压氩气等离子体中的阻尼和分散，物理。 Lett。 第230卷，1997年，pp. 218-222。 15Jeffery，A.和Taniuti，T.，非线性波传播，学术出版社，纽约，1964年。 16Gaitonde，G.V.，3D非理想磁气动力学求解器的开发，AIAA预印本99-3610，第30届等离子动力学和激光会议。 1999年6月28日至7月1日，弗吉尼亚州诺福克。 17Brio，M.和Wu。 C.C.，理想方程迎风差分方案磁流体动力学，J。 补偿。 物理学，第75卷，1988页。 400-422。 18Myong，R.S.和Roe，P.，关于Magnetohy drodynamics的Godnunov型方案，J. 补偿。 物理学，第147卷，1998年，第545-567页。 19Powell，K.G.，Roe P.L.，Myong，R.S.，Gombosi，T.和Zeeuw，D.D. 磁流体动力学的逆风方案，AIAA-95-1704-CP，1995年，pp. 661-671。 20MacCormack，R.W.，理想磁流体动力学方程的上风守恒形式方法，AIAA预印本99-3609，第30版

296 SHANG, CANUPP \& GAITONDE 等离子动力学和激光会议。 弗吉尼亚州诺福克，1999年6月28日至7月1日。 21Zachary，A.L.和Colella，P.，理想方程的高阶戈迪诺夫方法磁流体动力学，J. 补偿。 物理学，第95卷，1992年，pp. 341-347。 22Augustinus，J.，Hoffmann，K.A.和Harada，S.，磁场对高速流动结构的影响，J. 航天器和火箭，卷。 35，不。 5，9月至10月。 1998年，第639-646页。 23Lele，S.K.，具有光谱式分辨率的紧凑有限差分方案，J. Comp Physics，第103卷，1992年，第16-42页。 24Gaitonde，D.V.和Shang，J.S.，线性波现象的基于紧凑差分的优化有限体积方案，J. 补偿。 物理学第138卷，1997年，pp. 617-643。 25Shang，J.S.，时间依赖性麦克斯韦方程的高阶紧致差分方案，J. 补偿。 物理学，第153卷，1999年，pp. 312-333。 26Gaitonde，D.V.和Visbal，M.R.，进一步开发基于高阶公式的Navier-Stokes解决方案程序，AIAA预印本99-0557，第37届航空航天科学会议，内华达州里诺，1999年1月11日至14日。 27Visbal，M.R.和Gaitonde，D.V.，使用紧凑差分和过滤方案计算一般几何的空气声场，AIAA预印本99-3706，弗吉尼亚州诺福克，1999年6月。 28Gaitonde，D.V.和Shang，J.S.，波传播现象中的高阶有限体积方案，AIAA Preprint 96-2335，第27届AIAA等离子动力学和激光会议，新奥尔良，1996年6月17日至20日。 29Shang，J.S.，CEM多学科应用展望，AIAA预印本99-0336，第37届航空航天科学会议，内华达州里诺，1999年1月11日至14日。 30Canupp，P.W.，磁场对激波和高超声速流动的影响，AIAA预印本2000-2260，第31届AIAA等离子动力学和激光会议，科罗拉多州丹佛，2000年6月19日至22日。 31Muntz，E.P.和Harnett，L.N.，正常激波中的分子速度分布测量，物理。 流体第12卷，1969年，pp. 2027-2035年。 32Harten，A.，计算冲击和接触不连续性的人工压缩方法：III自调整混合方案，计算数学，第32（142）卷，1978年，第363-389页。 33Yee，H.C.，使用基于特征的滤波器的低耗散高阶激波捕捉方法，J. 补偿。 物理学，第150卷，1999年，pp. 199-210。 34Canupp，P.W.，使用通量向量分割方法解决磁气动力学现象，AIAA预印本2000-2477，流体2000会议，科罗拉多州丹佛，2000年6月19日至22日。 35Steger，J.L.和Warming，R.F.，Invscid气态方程的通量向量分割与有限差分法的应用，J. 补偿。 物理学，第40卷，1981年，pp. 263-293。

计算磁流体空气动力学 297 36Brackbill，J.U.和Barnes，D.C.，非零V • B = 0对磁流体力学方程数值解的影响，J。 补偿。 物理学，第35卷，1980年，pp. 426-430。 37Ambrosio，A.和Wortman，A.，驻点球体和气缸周围流动的冲击分离距离，美国火箭Soc。 J.，第32卷，1962年，p. 281。

\clearpage
\chapter{15 复杂几何的3D DRAGON网格方法的开发}
\section{15.1 介绍}
复杂几何的3D龙网格方法的开发 Meng-Sing Liou1 Yao Zheng2 15.1 导言 一个有效的CFD系统，用于通常涉及复杂几何的工程问题，必须具备某些特征，例如：（1）快速周转和（2）准确可靠的解决方案。 第一点包括人机的努力，需要短的计算设置时间、最小的内存要求以及高效和稳健的解决方案算法。 第二点需要明智地选择离散化程序。 网格方法的选择将极大地影响上述两个标准是否令人满意地满足。 事实上，多年来，主要涉及网格生成的大部分人力工作量没有显著减少，相反，由于微芯片技术的惊人进步，机器工作量急剧减少。 推进系统是复杂几何的一个例子，涉及许多具有奇异形状和急转弯的独立几何实体。 这种拓扑结构给网格生成带来了挑战，特别是对于黏性流动计算。 对于典型的三维流量计算，生成网格所花费的时间约为总模拟工作的三分之二，代表了位于俄亥俄州克利夫兰刘易斯菲尔德的1 NASA格伦研究中心，俄亥俄州44135。 2 Taitech, Inc.，美国宇航局格伦研究中心，位于俄亥俄州克利夫兰刘易斯菲尔德，44135。 Frontiers of 计算流体动力学 - 2002 编辑：David A. Caughey和Mohamed M. 哈菲兹 ©2002 世界科学

整个分析周期中的300 LIOU \& ZHENG瓶颈[28]。 因此，网格生成仍然是实用的CFD分析的起搏技术，也是可以实现重大回报的领域。 此外，用于覆盖黏性区域的高质量网格对于产生准确高效的解决方案至关重要。 为了应对复杂几何在生成电网方面施加巨大限制和困难的情况，目前复合结构化网格方案和非结构化网格方案是解决各种CFD问题的两种主流方法。 嵌套网格方案[26]和类似方案[1，6]使用超集网格来解析复杂的几何形状或流动特征，通常被归类为复合结构网格类别。 过量网格允许结构化网格以良好的质量使用，例如正交性和平滑性，并易于控制网格间距。 虽然经常用于复杂配置的流[4]，但它也被用于分析相对运动物体的流[9]。 此外，超集网格可以用作解决方案适应程序[18，2，22]。 然而，在没有严格满足控制方程的情况下，更新重叠区域变量的非保守插值可能会产生虚假解决方案，特别是通过陡峭梯度的区域。 非结构化网格方法也非常灵活，可以围绕复杂的几何形状生成网格。 特别是，解决方案的适应性也许是它最大的优势。 然而，非结构化网格方法已被证明是内存和计算密集型的[13]。 此外，高效流求解器的选择有限，因此进一步影响了该方法的计算效率。 在实践中，在空间中实现比二阶精度更高的方案比结构化网格更不那么容易。 此外，根据我们的经验，生成3D黏性网格，例如围绕尖锐凹角的领域，并不像看起来那么容易，稳健性是问题所在。 显然，每种方法都有自己的优势和劣势。 因此，我们试图开发一种方法，将结构化和非结构化网格正确结合在一起，并最大限度地发挥其自身优势。 事实上，一些混合方案已经出现[23，16]。 今天，大多数混合网格方法都来自非结构化网格社区，人们认识到结构化网格（棱镜网格）应嵌入非结构化网格下方，以便更好地解决黏性区域[16]。 这样做只是为了解决准确性问题，但事实上，它很难在凹角附近生成这种类型的网格，因此缺乏稳健性。 在这里，大部分地区都充满了非结构化网格，网格数据仍然是非结构化内存问题仍然存在。 相反，我们的方法将产生主要的结构化网格

\section{15.2 DRAGON Grid}
3D DRAGON GRID METHOD 301 域的部分，只有填充了非结构化网格的小区域。 因此，DRAGON网格[21，17]是通过非结构化网格直接取代重叠的任意网格来创建的。 虽然DRAGON网格自适应了嵌套网格的思维，但它有三个重要优势：（1）保留嵌套网格的优势，（2）消除嵌合计划中遇到的困难，以及（3）就控制方程而言，以完全保守和一致的方式实现电网通信。 此外，DRAGON方案旨在实现以下目标：（1）效率，（2）高质量网格，以及（3）复杂多组分几何的稳健性（通用性）。 15.2 龙网格 我们得出结论，（1）保持网格灵活性的属性和嵌合方法的质量肯定要保留，（2）专注于改进（选择）插值方案可能只会导致更复杂，这似乎不是一个富有成效的方法。 一种完全避免插值并严格执行通量守恒的替代方法是，在与域的其余部分相同的基础上解决有关区域。 由于重叠区域的形状必然是不规则的，非结构化网格方法最适合对该区域进行网格化。 此外，这个区域通常远离身体，黏性效应不如非黏性效应重要，就溶液精度而言，粗网格就足够了。 这种情况将使用非结构化网格，从而最大限度地减少其与内存要求相关的惩罚。 两种类型的网格的组合产生了混合网格。 我们的方法与其他混合方法的主要区别是：（1）我们保留了嵌套网格方法的吸引力；（2）我们仅在有限的区域使用非结构化网格。 换句话说，大部分区域是DRAGON网格中的结构化六面体网格，非结构化网格覆盖了黏性效应通常不太重要的区域，从而最大限度地减少了非结构化网格的缺点。 为了保持解决方案方案的效率，我们选择整合两个独立的求解器--一个是结构化的，另一个是非结构化的。 虽然在求解器的开发阶段，一开始需要额外的努力，但努力是一劳永逸的。 此外，如今获得两个这样的代码不再是障碍，因为许多代码都可用，并且它们已经单独验证并在生产模式下使用。 因此，代码的适当整合本质上是要完成的任务。 在下文中，我们将分别描述在DRAGON网格方法中生成结构化和非结构化网格所采取的步骤。 给

302 LIOU \& ZHENG 图1 用非结构化网格直接替换任意网格重叠：（a）嵌套网格，（b）DRAGON网格。DRAGON网格技术的要点，我们将只考虑二维拓扑，但不会失去一般性。 15.2.1 结构化网格区域1。 与嵌套网格一样，整个计算域被分成子域。 我们通常将包含完整计算域的主要（或背景）网格指定为次要网格。 2. 创建孔区域，参考图1（a），如果重叠，否则空隙区域就是孔。 例如，小网格的外部边界可以用作孔的创建边界。 3. 孔边界点现在正在形成非结构化网格区域的新边界。 由于网格不需要在DRAGON网格框架下重叠，因此其程序比嵌套网格方法更灵活。 15.2.2 非结构网格区域 间隙区域不可避免地具有不规则的形状，三角形单元格特别适合适应。 回想一下，DRAGON网格中的一个重要功能是消除任何导致通量不守的繁琐插值。 仅靠非结构化网格不足以完成任务。 施加了额外的约束，要求结构化网格的边界节点与边界三角形单元格的顶点重合。 换句话说，可能有多个三角形连接到一个结构化网格单元，但反之亦然。 幸运的是，这种约束非常适合非结构化网格生成。 Delaunay三角测量方案[3、29、-30]

\section{15.3 三维龙 网格生成}
在本工作中，3D龙网格方法303用于生成非结构化网格。 图1（b）描绘了DRAGON网格，非结构网格填充了孔区域。 在下文中，如果使用嵌套网格方案的框架，我们将给出采用非结构化单元格的步骤。 1. PEGSUS程序[27]提供的边界节点根据其几何座标重新排序。 2. 然后执行Delaunay三角测量法来连接这些边界节点。 3. 除了包含基于单元格和基于边缘的信息的标准连接矩阵外，我们还需要引入额外的矩阵来描述结构化和非结构化网格之间的连接。 15.3 三维DRAGON 网格生成 将上述DRAGON网格的概念扩展到三维空间的概念是直接的。 然而，预计在实际实施方面会有更多的困难。 三维案例的独特挑战存在于各个方面，例如非结构化网格生成算法和代码实现。 这项工作强调了三维DRAGON网格方法的开发。...3 5 /L-A A Structured Grid Faces of Nonstructured Grid Figure 2 三维DRAGON网格中结构化和非结构化网格之间的接口。 参考图1，虽然结构网格的边界边缘与二维边界三角形单元格的边缘重合，但总体而言，结构网格的面不一定与图2所示的三维情况下同一位置的非结构化网格的面相匹配。 这种限制可能源于网格生成

304 LIOU \& ZHENG阶段或适应物理现象的要求。 在图2中，结构化网格上的1、2、3、4、5和6点与非结构化网格上的相应点重合。 因此，可以很容易地制定一个方案，以满足普通面的通量守恒，而无需诉诸通量插值。 例如，图3描绘了压缩机鼓腔问题的DRAGON网格。 此网格由基础结构网格创建，其中有重叠的域和空隙。 结构化网格是车身贴合的，质量很高，适合黏性流动模拟。 非结构化网格可以轻松生成，以填补重叠的域和空隙。 图3提供了DRAGON网格的侧视图，其中外缘几乎没有露出第三维度。 图3压缩机鼓腔问题的DRAGON网格（外缘几乎没有显示出三维）。 15.3.1 编程方面 除了PEGSUS代码[27]外，在三维DRAGON 网格生成中创建了两个名为DRAGONFACE和MGEN3D的主要模块，图4中描述的流程图也显示了相关模块之间的关系。 DRAGONFACE，一个FORTRAN代码，读取由PEGSUS代码分配的IBLANK值的复合结构网格，并写出面的拓扑信息作为非结构网格的曲面描述。 DRAGONFACE程序涉及的主要步骤包括：识别新面孔，创建表面的孤立边缘，以及填充

3D龙网格方法305复合结构网格域连接（PEGSUS）非结构网格面（DRAGONFACE）非结构化网格生成（MGEN3D）表面和体积网格DRAGON网格图4涉及三维DRAGON 网格生成的程序模块。表面的开口，以确保表面的封闭性。 MGEN3D是一个C代码。 它从DRAGONFACE程序中输入曲面描述，生成体积非结构网格，并提供相应曲面网格的拓扑信息，在结构化和非结构化网格之间的接口上提供数据通信。 MGEN3D程序涉及的主要步骤包括：间距源创建、三维表面三角测量、面重新定位、域分类和体积网格的生成。 目前，四面体网格生成方案有两种基本型别：前进前进[25]和Delaunay三角测量[19]。 此外，还提出了一种耦合方法[10]，以利用两种方案的优势。 在对这些计划进行调查后，选择了Delaunay三角测量，部分考虑到了开发所需的时间框架。 另外两种方法，有其自身优势，将来将予以考虑。 MGEN3D程序的算法是根据之前工作[29、30、19]的框架演变而来的，是Delaunay型方法。 Steiner点[24]创建算法已被纳入此网格生成器中。 MGEN3D程序中涉及的算法总结如下。 1. 通过DRAGONFACE程序生成的曲面描述被读取。 2. 表面三角测量是在考虑基线网格间距、点和线源分布情况下进行的。 3. 进行域分类是为了提供域信息

306 LIOU \& ZHENG几何由上述表面定义指定。 4. 为域生成体积网格。 最终创建了单个体积网格的集合数据，包括节点坐标和连接性。 为了展示DRAGON电网代码的能力，我们考虑了典型的薄膜冷却涡轮叶片，其示意图如图5所示。 这种几何形状包括叶片、冷却液和叶片内的33个孔。 图6描绘了由此产生的DRAGON网格，其中33个孔和流域之间的连接区域被非结构化网格填充。 它通过切割网格来深入了解这个DRAGON网格，其中注意到了前缘区域。入口“°w\textasciicircum{}领先\textasciicircum{}\textasciicircum{}”边缘\textasciitilde{}um入口流-垂直到墙后缘弹射堵塞图5叶片横截面和冷却液孔的示意图。 在DRAGON 网格生成期间，为孔创建了33个单独的结构网格，而没有试图在拓扑上将它们与代表外部流域的网格连接，这是多块网格生成 [15]中需要的。 从这个角度来看，DRAGON网格方案可以被认为是一种更灵活、更容易的方法。 然而，由于物理和数值特性，用户可能更喜欢在特定区域有结构化网格。 例如，在本例中，用户可能希望使用结构化网格来填充孔和流动域之间的连接区域。

3D龙网格方法307图6薄膜冷却涡轮叶片前缘区域龙网格的切口视图。 15.3.2 通过网格接口的数据通信 对于DRAGON网格来说，守恒定律在结构化和非结构化网格上都是在相同的基础上解决的。 单元格中心方案用于存储变量和通量评估，其中分别使用四边形和三角形单元格。 图7显示了连接结构化和非结构化网格的接口。 如前所述，在细胞界面评估的数值通量是基于相邻细胞（分别表示为L和R细胞）的条件。 对于非结构化网格，接口通量Fn[，将使用结构化单元格值作为右（R）状态，非结构化单元格值作为左（L）状态进行评估。 因此，已对非结构化网格进行评估的接口通量现在可以直接应用于计算结构化网格的单元体积残差。 因此，通过DRAGON网格中的网格接口进行数据通信保证了满足守恒定律，并在结构化和非结构化网格区域的相同基础上获得解决方案。

\section{15.4 流动求解器}
308 LIOU \& ZHENG 非结构化网格 *"•«.:; / \textbackslash\{\}l y o R Fnl C\textbackslash\{\} ••-L f 结构化网格 O i,i \textbackslash\{\} \textbackslash\{\} \textbackslash\{\} \textbackslash\{\} \textbackslash\{\} \textbackslash\{\} 结构化/非结构化网格接口 图 7 连接结构化和非结构化网格的单元格面上的通量。 15.4 流动求解器 方程（15.1）给出的与时间相关的可压缩Navier-Stokes 方程，在任意控制体 fl上以整数形式求解：F • dS = 0，（15.1）其中保守变量向量U =（p，pu，pv，pw，pet）T，比总能量为e* - e+ \textbackslash\{\}V\textbackslash\{\}2/2 = ht -p/p。 通量向量F包括黏性通量和黏性通量，其中也包括湍流变量。 基于以细胞为中心的有限体积法，半离散化形式，通过所有包围面的通量平衡来描述CI中U的变化时间率，Si，I - 1，---，LX，无论它们是在结构化还是非结构化网格区域，都可以作为OTT LX 1=1（15.2），其中Fn；= F；•n；和Hi是Si的单位法线向量。 流程代码称为DRAGONFLOW，由两个经过验证的美国宇航局代码组成，OVERFLOW [5]和USM3D [11，12]，它们分别是结构化和非结构化网格代码。 在添加接口和修改两个代码方面已经做出了重大努力，更是如此在OVERFLOW代码上。

\section{15.5 测试案例}
3D DRAGON GRID METHOD 309 在OVERFLOW代码中应用了空间因子隐式算法，而在USM3D代码中应用了点隐式算法。 由于我们只对寻求稳态解决方案感兴趣，因此可以容忍由于结构化和非结构化网格区域之间不兼容而导致的收敛速率的差异。 然而，如果非稳态解决方案感兴趣，可以通过双重时间积分方法克服这一困难，其中每个物理时间步骤中的解决方案将迭代以实现收敛。 这个主题留给未来的研究。 至于空间离散化，准确性顺序不一致，因为在结构化基于网格和非结构化基于网格的代码中分别使用了三阶和二阶准确方案。 但应用了类似的限制函数。 我们使用新的通量方案AUSM+，如[20]中详细描述，在两个代码中表达细胞面的非黏性通量。 此外，黏性通量像往常一样，由中心方案近似。 15.5 测试案例 15.5.1 案例 1：激波管 问题 首先，考虑了 2D 激波管 问题。 为了只关注网格问题，在时间和空间上都通过基本的一阶精确离散化解决了流量。 这个案例用于显示嵌套网格中插值的影响，以及当平面冲击穿过嵌入式网格区域移动时，DRAGON网格方法对瞬态问题的有效性。 激波正在以设计的冲击速度Ms = 4的恒定面积通道中进入静止区域。 如图8所示，使用三个网格系统获得解决方案，即（1）单网格、（2）嵌套网格和（3）DRAGON网格。 单网格解决方案用于基准比较。 如图9所示，沿通道中心线的压力分布表明，嵌合方案预测了管子中移动速度更快的冲击，而当前的DRAGON网格和单个网格结果是一致的，表明在穿过嵌入式DRAGON网格区域时，冲击被准确捕获，并很好地保留了守恒特性。 15.5.2 案例2：收敛管道中的超声速流动 接下来，我们考虑透过收敛通道的3马赫的超声速流动。 图10勾勒了几何形状，其中顶部和前侧墙都弯曲了10度的楔形角度，从而形成了两个楔形

初始激波位置图8移动冲击问题的网格（Ms = 4）：（a）单网格，（b）嵌套网格和（c）DRAGON网格。强度相等的冲击以及它们之间的后续相互作用。 一条由阴影区域表示的非结构网格区域的条带被放置在通道的中间部分。 此测试用例旨在通过冲击对DRAGON网格的守恒属性的三维代码进行检查。 图11显示了顶部和正面的密度轮廓以及出口平面的透视图。 两个楔形冲击在角边缘AD附近相交并产生一个角冲击区域（见图10），其表现为从三点发出的滑线。 如图12所示，随着楔形冲击向对面的墙壁扫荡，这个区域逐渐变大。 最终，这两个楔形减震器反射并与之前由角减震器产生的流量相互作用，使流场更加复杂。 图13提供了中平面EFGH上马赫数轮廓的内视图（如图10）。 这揭示了管道内冲击-冲击相互作用的丰富特征，而管道壁上的冲击配置（图11）相对简单。 从冲击三点发出的滑动表面在图13中很明显，它们随后与

3D DRAGON GRID METHOD 311 AO \textbackslash\{\} a SlngtoOrid DChlmMaOrU • DRAOONQrid I I., 1.,,,. 1 180 190 200 210 220 230 x-距离图9三个网格的中心线压力的比较。下游冲击。 很明显，冲击轮廓无缝地穿过DRAGON网格接口，而不会产生假波，保证了保护特性。 在上述每个图中，我们还包括基于单个结构化网格进行验证的结果。 我们看到，两组结果本质上是相同的，除了核心区域的一些微小变化，这是非结构化网格的残余结果。 图14清楚地表明了这一点，在非结构化网格区域的冲击变薄，因为网格大小大致减少了\textbackslash\{\}。 图15给出了单网格和DRAGON网格解决方案在退出平面上的差异的定量测量，表明两个解决方案的一致性。 图10 收敛管道的草图，非结构化网格位于DRAGON网格的阴影区域。 15.5.3 案例3：黏性流动通过环形级联 最后，我们展示了美国宇航局刘易斯[14]开发和测试的涡轮机定子叶片环形级联的流动模拟结果。 图16显示了DRAGON网格，其中放置了一个背景H型网格来覆盖叶片之间的通道，O型黏性网格用于解决

\section{15.6 结束的讲话}
312 LIOU \& ZHENG 图11 收敛管道表面流动的密度分布：（a）单网格；（b）DRAGON网格。 图12 马赫数 各个站点的流量分布：（a）单网格；（b）DRAGON网格。叶片周围的区域，以及它们之间的非结构化网格区域。 之前的计算是使用单个结构化网格获得的[8，7]。 图17绘制了当前DRAGON网格解决方案的静压分布，它们与三个跨度位置的测量数据[14]非常一致，分别为13.3\%、50\%和86.6\%。 15.6 结论 在目前的工作中，我们专注于将DRAGON网格方法扩展到三维空间。 已经调查了扩展的各个方面，以及三维案例的新挑战。 这种方法试图保留结构化和非结构化网格的有利特征，同时消除或尽量减少其各自的缺点。 因此，该方法非常适合为各种单个组件快速创建高质量的黏性网格

\section{致谢}
3D龙网格方法313图13 马赫数典型切割平面EFGH的流量分布：（a）单网格；（b）龙网格。 图14对称平面ABCD上流动的密度分布：（a）单网格；（b）DRAGON网格。在工程系统中发现的复杂形状。 在计算上，生成的网格大大减少了非结构化网格所需的内存，并使更高效的求解器可访问，因为DRAGON网格的主要部分是结构化网格的形式。 此外，高质量的黏性网格在过程中继承，与非结构化网格生成不同，需要添加结构化网格状网格（如棱镜层），但仍然基于非结构化网格数据结构。 此外，流动解决方案通过结构化-非结构化区域的界面证实了守恒属性的满意度，结果与可用的测量数据非常一致，从而证明了该方法的可靠性。 未来计划包括进一步完善非结构化网格生成，以提高平滑度和稳健性，并进一步验证和应用复杂几何形状的工程问题。 致知 这项工作由美国宇航局格伦研究中心罗伯特·科里根管理的涡轮机械和燃烧技术支持。 我们感谢J。 D。 海德曼和R。 V。 Chima分别提供薄膜冷却涡轮机和环形级联的叶片几何形状。

\section{参考资料}
314 LIOU \& ZHENG 图15（a）密度，（b）马赫数和（c）HG线上流动的压力分布。 框和圆圈分别表示基于单网格和DRAGON网格的结果。 参考资料1。 E。 H。 阿塔。 组件自适应网格接口。 AIAA论文81-0382，AIAA第19届航空航天科学会议，St. 路易斯，密苏里州，1981年。 2. M。 J。 Berger和J。 奥利格。 双曲偏微分方程的自适应网格细化。 计算物理学杂志，53:484-512，1984年。 3. A. 鲍耶。 计算狄利克雷镶嵌。 计算杂志，24（2）：162-166，1981年。 4. P。 G。 布宁，我。 T。 奇乌，F。 W。 小马丁，R. L。 米金，S。 小林，Y。 M。 里兹克，J。 L。 Steger和M。 蓍草 航天飞机上升中的流场模拟。 第四届超级计算国际会议，第20-28页，加利福尼亚州圣克拉拉，1989年4月。 5. P。 G。 Buning等人。 OVERFLOW用户手册，版本1.8f。 未发表的美国宇航局报告，美国宇航局，1998年。 6. G。 切什郡和W。 D。 亨肖。 用于求解偏微分方程的复合重叠网格。 计算物理学杂志，90：1-64，1990年。 7. R。 V。 奇玛，P。 W。 Giel和R。 J。 博伊尔。 三维黏性流动的代数湍流模型。 美国宇航局TM 105931，美国宇航局，1993年。 8. R。 V。 Chima和J。 W。 横田。 三维黏性内流的数值分析。 美国宇航局TM 100878，美国宇航局，1988年。 9. F。 C. 多赫蒂，J。 A. Benek和J。 L。 斯特格。 关于应用嵌套网格计划来存储分离。 美国宇航局TM 88193，美国宇航局，1985年10月。 10. P。 J。 弗雷，H。 Borouchaki和P.-L. 乔治。 3D Delaunay网格生成与前进的前沿方法相结合。 应用力学和工程中的计算机方法，157：115-131，1998年。 11. N。 T。 弗里克。 在非结构化四面体网格上求解Euler 方程的上风方案。 AIAA杂志，30（l）：70-77，1992年。 12. N。 T。 弗里克。 湍流流的四面体非结构化Navier-Stokes方法。 AIAA杂志，36（11）：1975-1982，1998。 13. F。 加法里。 关于非结构化网格尤拉求解器的涡流预测能力。 AIAA论文94-0163，AIAA第32届航空航天科学会议和展览，内华达州里诺，1994年1月。 14. L。 J。 戈德曼和R。 G。 Seasholtz。 核心涡轮叶片环形级联中的激光风速计测量以及与理论的比较。 美国宇航局技术论文2018，美国宇航局，1982年。

3D龙网格方法315图16环形涡轮机级联的龙网格。 15. J。 D。 海德曼，D。 Rigby和A。 A. 美国。 薄膜冷却涡轮叶片的三维耦合内部/外部模拟。 翻译。 ASME，涡轮机械杂志，122：348-359，2000年。 16. D。 G。 福尔摩斯和S。 D。 康奈尔。 非结构化自适应网格上的2D Navier-Stokes 方程解决方案。 AIAA论文89-1932，AIAA第9届CFD会议，纽约州水牛城，1989年6月。 17. K.-H. Kao和M.-S. 利欧。 推进重叠网格方案：从嵌合到DRAGON网格。 AIAA杂志，33（10）：1809-1815，1995年。 18. K.-H. 高，M.-S. Liou和C。 Y。 周。 网格自适应使用嵌合体复合重叠网格。 AIAA杂志，32（5）：942-949，1994年。 19. R。 W。 刘易斯，Y。 郑和D。 T。 Gethin。 三维非结构网格生成：第3部分。 体积网格。 应用力学和工程中的计算机方法，134（3/4）：285-310，1996年。 20. M.-S. 利欧。 继续寻找近乎完美的数值通量方案，第一部分：AUSM+。 美国宇航局TM 106524，刘易斯研究中心，俄亥俄州克利夫兰，1994年。 还有《计算物理学杂志》，129：364-382，1996年。 21. M.-S. Liou和K.-H. 考。 网格生成的进展：从嵌合到龙网格。 美国宇航局TM 106709，刘易斯研究中心，俄亥俄州克利夫兰，1994年8月。 还有第21章，在计算流体动力学的前沿1994，ed。 D。 A. Caughey和M。 M。 哈菲兹，约翰·威利和儿子，1994年11月。 22. R。 L。 米金。 关于自适应细化和过度结构化网格。 AIAA论文97-1858，AIAA第13届CFD会议，科罗拉多州斯诺马斯，1997年6月。 23. K。 中桥和S。 小林。 黏性流动的FDM-FEM区域方法

316 LIOU \& ZHENG对多体的计算。 AIAA论文87-0604，AIAA第25届航空航天科学会议，内华达州里诺，1987年1月。 24. A. 冈部，B。 靴子和K。 杉原。 空间镶嵌：沃罗诺伊图的概念和应用。 John Wiley \& Sons，奇切斯特，1992年。 25. J。 佩雷尔，J. Peiro和K。 摩根。 用于三维可压缩流动计算的自适应重新正片。 计算物理学杂志，103：269-285，1992年。 26. J。 L。 Steger和J。 A. Benek。 关于复合网格方案在计算空气动力学中的使用。 应用力学和工程中的计算机方法，64（l/3）：301-320，1987年。 27. N。 E。 Suhs和R。 W。 特拉梅尔。 PEGSUS 4.0 用户手册，AEDC-TR-91-8。 卡尔斯潘公司/AEDC运营，田纳西州阿诺德空军基地，1991年11月。 28. R。 塔加维。 在Cray Research超级计算机上从CAD自动生成、并行和容错网格。 技术报告，Cray用户小组会议，法国图尔，1994年。 29. Y。 郑，R。 W。 Lewis和D。 T。 Gethin。 三维非结构网格生成：第1部分。 三角测量和点创建的基本方面。 应用力学和工程中的计算机方法，134（3/4）：249-268，1996年。 30. Y。 郑，R。 W。 Lewis和D。 T。 Gethin。 三维非结构网格生成：第2部分。 表面网格。 应用力学和工程中的计算机方法，134（3/4）：269-284，1996年。

3D龙网格方法317轴向弦的分数轴向弦的分数-0.7 0-65 0.6 - 0.55 0.5（c）轴向弦的分数图17（a）13.3\%，（b）50\%和（c）86.6\%的轴心跨度的静态压力分布

\clearpage
\chapter{16 将多块、补丁网格拓扑结构应用于Navier-Stokes对陆军炮弹空气动力学的预测}
\section{16.1 介绍}
将多块、补丁网格拓扑结构应用于海军炮弹空气动力学的Navier-Stokes预测沃尔特·B。 Sturek，Sr.1和David J. Haroldsen2 16.1 简介 陆军研究实验室有兴趣应用最先进的高性能计算工具来预测陆军炮弹在攻击角度的空气动力学，适用于从高亚音速到高超音速的广泛范围马赫数。 在本文中，WIND 流动求解器被用来研究两种导弹配置的空气动力学。 GridPro 网格生成软件已被用于生成计算网格。 讨论了与WIND一起使用的电网生成以及获得解决方案的过程相关的几个方面。 显示了几个湍流型号的比较。 美国陆军研究实验室计算流体动力学的研究人员有兴趣调查一系列复杂的流体流动问题。 这些问题包括围绕复杂物体的流动，中度和高马赫数的流动，以及中度到高攻击角度的流动。 最近的一项研究考察了几种不同的Navier-Stokes流求解器的预测能力，这些求解器适用于跨声速和超声速流动速度下14度攻角下的Ogive气缸配置的情况[l]。 这项研究通过检查WIND 流动求解器预测流量1美国的预测能力，扩展了之前的工作。 陆军研究实验室，阿伯丁试验场，MD 21005-5067 2 美国军事学院，西点，纽约 10996 Frontiers in 计算流体动力学 - 2002 编辑：David A. Caughey和Mohamed M. 哈菲兹 ©2002 世界科学

\section{16.2 导弹配置}
320 STUREK和HAROLDSEN攻击角鳍片和中度马赫数导弹的问题。 WIND软件包具有许多功能，使其对计算研究人员具有潜在吸引力。 这些功能包括众多的湍流型号、易用性、便携性、并行处理能力，以及将网格与广义拓扑相结合的能力。 这一特殊功能使WIND在与GridPro 网格生成软件包一起使用时具有吸引力。 GridPro生成具有非重叠块接口的结构化多块网格。 这项工作的重点是调查：1）WIND在复杂流量问题的应用；2）多块、修补的网格拓扑的使用。 本研究考虑将WIND 1.0应用于研究两种不同的导弹配置，攻角为14度和40度，马赫数接近2.5。 讨论了使用GridPro包生成多块、修补的网格，并介绍了不同湍流模型的结果。 16.2 导弹配置 本研究检查了两种导弹配置。 两枚导弹都由3口径的机头和10口径的圆柱体组成。 每枚导弹都有四个鳍片，与俯仰平面对称。 图1和图2显示了特定的鳍形状和位置。 Y 图1 导弹1配置。

MULTIBLOCK 321图2导弹2配置的应用。 导弹1[2]在0和45度、2.5马赫和14度的攻角下进行了研究，雷诺数为1.12 x 10。 导弹2[3]在45度的滚动角度、1.6和2.7马赫、40度的攻角下进行了研究，雷诺数为25万。 16.3 网格生成 本次调查的网格生成是使用GridPro 网格生成软件完成的。 GridPro是位于纽约白原市的Program Development Corporation的产品。 这个软件包很有趣，因为它结合了基于拓扑的网格生成方法。 这种方法强调几何形状和任何流动特征的基本拓扑结构，而不是关注问题的几何形状。 该软件包由用于拓扑设计的GUI、网格生成软件和用于操作网格的实用程序组成。 GridPro生成多块结构化网格，能够以各种格式输出资料。 使用GridPro时的一个重要考虑因素是相邻区域相邻但不重叠。 用户还可以自定义GridPro，以输出与特定流动求解器相关的初始边界数据。 用户通过在感兴趣的区域构建一个粗略的、无结构的六面体网格来设计拓扑结构。 六面体元素成为最终多块网格的单个块。 用户只控制网格的拓扑结构；网格生成器自动计算网格线的精确位置。 由于网格化过程在很大程度上是自动的，用户在设计网格的基本拓扑结构方面具有相当大的灵活性。 这种灵活性包括能够在感兴趣的区域本地定义网格，同时在流量变化不显著的区域留下更粗的网格。 拓扑设计完成后，用户调用网格生成软件。 复杂形状的网格可能会产生数百个块。 软件包中包含的公用设施允许对最终电网进行各种操作。 两个特别感兴趣的是块合并实用程序和聚类实用程序。

\section{16.3 边界/初始条件}
322 STUREK和HAROLDSEN块合并实用程序将（通常）大量块合并为更易于管理的数字。 用户可以控制最终配置的每个块中允许的网格节点数量。 因此，合并实用程序与初始拓扑设计一起可以用作先验的“区域分解”工具。 另一个感兴趣的实用程序是将网格线聚集在特定表面上。 在实践中，网格生成软件包总是用于生成欧拉网格，聚类实用程序用于获取黏性网格。 这大大减少了获得黏性网格所需的时间。 对于正在考虑的导弹，拓扑设计需要几天时间来构建一个合理的拓扑。 网格是使用带有R12000 CPU的SGI ONYX生成的。 网格生成通常需要6到8小时的CPU时间。 每种情况下的初始网格都有375个块，流量的一半（假设对称）。 通过合并到更易于处理的数量，减少了区块的数量。 最终的网格配置列在表导弹1（0度滚动）1（45度滚动）块网格点（以百万计）4.1 4.3 4.2表1网格详细信息中。 为导弹设计的拓扑没有旋转对称性；因此，必须以不同的滚动角度为1号导弹生成不同的拓扑和网格。 图3、4和图5显示了拓扑结构的示例。 图3说明了鼻尖拓扑的特写视图。 图4显示了鳍面拓扑的特写。 图5显示了整个导弹拓扑结构。 生成的网格样本如图6和图7所示。 在这些图中，较深的线表示块边界，并暗示了用于创建网格的拓扑设计。 出于显示目的，图中显示了黏性边界层添加到网格之前的最终网格的粗略版本。 16.4 边界/初始条件 在所有情况下，自由流流入条件用于流入边界，反射边界条件用于对称平面，自由流流出条件用于流出平面，黏性墙条件用于导弹体的黏性表面。 对于导弹1，总自由流压力为20.628磅/平方英寸，总自由流温度为554.4度。 对于导弹2，静态自由流压力为0.5637磅/平方英寸，静态自由流温度为248.4度。 WIND默认初始化用于初始化流场变量。

\section{16.4 性能/收敛标准}
\section{16.5 结果}
多块323的应用使用以下湍流模型进行运行：Baldwin-Lomax（BL）、Baldwin-Barth（BB）、Spalart-Allmaras（SA）和剪应力传输（SST）。 Baldwin-Lomax在有和不有选择最大网格点数以搜索Fmax的选项的情况下运行。 对于前一种情况，研究了最大网格点10和30（分别为BL10和BL30）。 为了避免瞬态不稳定，使用了FIXER关键字。 对于导弹1，在低攻角下计算了初始解决方案，该解决方案被用作在高攻角下计算解决方案的初始解决方案。 对于导弹2，TVD系数减少到1，CFL交叉系数设置为1，CFL数减少到0.4。 16.5 性能/收敛标准运行是在具有多个处理器的Silicon Graphics Origin 2000或Onyx平台上进行的。 运行通常使用8个处理器进行，融合解决方案可以在8-12小时内获得。 在每种情况下，在数千个周期内，残差减少不超过3个数量级。 为了测试收敛，对解决方案进行了监控，直到判断它们是收敛的。 就导弹2而言，使用WIND中的LOADS关键字计算了身上的载荷，当载荷收敛并保持了几百次周期的稳定时，该解决方案被认为是收敛的。 获得的并行性能因使用的网格而异。 用于导弹1的网格被减少到最终的方块数，同时试图平衡每个方块的节点数量。 使用8个处理器获得了高达7.5的加速系数，16个处理器获得了高达14的加速系数。 导弹2的性能要差得多--8个处理器的加速系数约为5--因为块合并过程是为了尽量减少块的数量，而不是优化并行性能。 最佳并行性能的示例如图8所示。 16.6 研究感兴趣的结果数量是身体和鳍片上不同站点的压力系数，以及几个轴向站外流场的皮托压力曲线。 所提供的数据显示了使用WIND获得的导弹1的结果示例。 选择图中显示的站点和数据是为了最终将计算数据与实验数据进行比较。 图9所示的卷角为0和轴向站X/D=11.5的导弹1的皮托压力预测对SA、BB和SST模型显示了类似的结果。 BL10 湍流模型预测了一个更小、更强烈的初级涡流。 此外，该模型预测了导弹机身附近的更结构化的解决方案。 BL 湍流模型预测的解决方案与预测更相似

\section{16.6 结束的讲话}
\section{16.7 致谢}
\section{参考资料}
324 STUREK和HAROLDSEN的一和二方程模型。 不同轴向站导弹体表面压力预测的比较显示，除BLIO外，不同湍流型号之间的差异很小。 样品比较如图10和图11所示，卷=0度。 在所需的攻角下，导弹2无法获得合理的结果。 试图通过提高网格质量和密度，减少CFL数量，以及使用WIND提供的平滑选项来提高问题的稳定性。 在每种情况下，由于流动场的奇点，运行都中止了。 困难很可能是由于极端的40度攻击角度。 使用相同的网格、相同的流动条件执行，但以20度的攻角成功收敛。 图12和图13显示了预测皮托压力的三维视图。 还显示了导弹1在两个滚动角度上涡流芯位置的线。 这些可视化是使用博士开发的PV3软件包获得的。 麻省理工学院的罗伯特·海姆斯。 16.7 结论性评论 WIND 流动求解器已被证明是提高和扩展计算流体动力学研究人员预测能力的有效工具。 事实证明，WIND对具有复杂几何形状的流动问题特别有用，尽管极端的流动条件造成了一些困难。 WIND能够容纳多块、修补的网格，使GridPro 网格生成软件的使用成为可能。 该软件包被发现可以以适度的努力生成高质量的结构化网格，并且非常适合用于陆军导弹配置。 与本研究中使用的其他模型相比，Spalart-Allmaras 湍流模型的表现相当好。 特别是，应用高阶模型没有带来任何好处。 致谢 这项工作得到了陆军研究实验室主要共享资源中心授予的计算机时间的支持。 陆军研究实验室的科学可视化实验室也提供了帮助。 这项研究由陆军研究实验室高性能计算部门资助。 参考资料1。 Sturek, W.B., Birch, T., Lauzon, M., Housh, C, Manter, J., Josyula, E., Soni, B., “CFD在导弹体流预测中的应用”，第35届航空航天科学会议和展览，内华达州里诺，1997年1月6日至10日。

应用OP MULTIBLOCK 325 2。 伯奇，T.，私人通信，1996年。 3. Stallings, R., Lamb, M.5 Watson, C, “雷诺数对超音速十字形翼体稳定性特性的影响”，美国宇航局TP 1683, My 1980. x\&iirfi\textasciicircum{}J *i wl»rf,::--. > •••？ 图3 鼻尖拓扑结构。

326 STUREK和HAROLDSEN \textasciicircum{} 4i。 *\textasciicircum{}'\textasciicircum{}？ \#*\&•'•\% Sg\textasciicircum{}i"-* j\textasciicircum{}jf\textasciicircum{}sl\textasciicircum{}V' p-\textasciicircum{}y ••！ ><>•„••C：。 \$图4鳍拓扑，；S\textasciicircum{}F"** \textasciicircum{} ild图5导弹拓扑。

多块327的应用图6导弹1的粗欧拉网格示例。 图7 鳍网格。

328 STUREK和HAROLDSEN lb r z A B 计算 / / /..... //* C //' ya As i....1. 10 20 30 处理器 图 8 使用导弹1的网格测量加速，滚动角度为0。 ' 0 0.2 0.4 0.6 0.8 1 1.2 1.4 1.6 图 9 导弹1的X/D = 11.5时的皮托特压力预测的比较。

多块的应用 329 X/D = 5.5 X/D = 7.93 X/D = 8.79 0.3 r- 0.25 '- 0.2 - 0.15 - 0.1 - a O 0.05 F- -0.05 -0.1 -0.15 -0.2。 X/D = 9.93 50 P\textasciicircum{},00 BB SST 2.5 > 2 5。 5 7.93 i/ 'i '1|J1 9.93 图10 导弹机身表面压力预测的比较 1。

330 STUREK和HAROLDSEN X/D = 5.5 X/D = 7.93 50 100 PHI X/D = 8.79 3.5 >-2.5 1.5 0.5 或PHI X/D = 9.93 7.93 8.79 PHI 9.93 BL10 BL30 10 11 12 13 图11 导弹机身表面压力预测的比较。

多块331图12皮托特压力和涡核预测在滚动角度0的导弹1上的应用，SA 湍流模型*图13皮托特压力和涡核预测在滚动角度45的导弹I，SA 湍流模型。

\clearpage
\chapter{17 雷诺方程方程的空气动力学预测-平均Navier-Stokes 方程}
\section{17.1 介绍}
关于雷诺平均 Navier-Stokes 方程 M解的空气动力学预测。 G。 Hall1 17.1 导言 介绍了通过翼型翼的二维流动的数值实验。 目标是不仅要识别和减少湍流模型的局限性，还要确定雷诺平均近似本身的实际空气动力学局限性。 这里考虑了雷诺平均限制在多大程度上，即使是雷诺平均Navier-Stokes（RANS）方程的理想解也能实现的。 当然，对于整个实用的、通常更复杂的配置来说，这是一个重要的问题，但机翼或机翼部分可能是可以学到有用经验的最简单的。 机翼性能可预测性的任何严重限制也必须适用于有翼飞机。 这里，事先是数值实验背景的概要。 普遍的看法是，在具有实际利益的高雷诺数下，湍流刻度的范围非常大，以至于在某些年内，直接和大涡模拟都不太可能可行。 对于实际的空气动力学预测，人们普遍认为，RANS方程的数值解，以及改进的湍流建模，目前提供了最好的替代方案。 Speziale最近对该主题进行了深入的审查[9]。 非常需要改进湍流建模，但它1 Hall C。 F。 D。 有限公司，8 Dene Lane，Farnham，Surrey GU10 3PW，英国。 Frontiers of 计算流体动力学 - 2002 编辑：David A. Caughey和Mohamed M. 哈菲兹 ©2002 世界科学

334 HALL提出了一个巨大的挑战。 即使经过相当大的努力，几年来，只有稍微偏离平衡状态湍流的流量才能获得令人满意的结果。 然而，例如，典型飞机机翼上的湍流在正常运行条件下--在过渡、分离或重新连接、与冲击的相互作用以及后缘--都明显偏离了平衡。 现在的努力似乎主要针对开发双方程和应力传输模型，这些模型是应力张量全运输方程的更好模型，以改善非平衡流的模拟，重点是尽量减少经验假设的数量。 可以理解的是，在RANS平均过程中，很少关注物理流动的一些特征的损失。 冲击与边界层相互作用的RANS模拟为RANS固有的局限性提供了一个说明性的例子。 模拟冲击的厚度将主要与网格的单元大小有关，因此，如果网格足够细，边界层外的冲击压力上升将发生在与边界层厚度相比较小的距离。 然而，在现实中，互动流程非常不同[8，5]。 边界层中较大的涡层在单个涡层的纵向维度的距离上诱导了冲击的不规则振荡。 因此，平均冲击压力上升将分布在距离上，该距离也与大堤尺寸相同，这通常与边界层厚度相同。 无论采用湍流模型，模拟流量在质量上都会与平均实际流量不同，通过完善网格来提高RANS解决方案的准确性只会增加这种差异。 例如，如果边界层处于不利的压力梯度中，并且容易分离，可能会产生严重的实际后果。 通过平均简化Navier-Stokes 方程的一个惩罚是失去模拟大涡对相邻无黏流动的任何特定影响的能力。 对于通过起重机翼后缘的流动来说，准确的模拟尤为理想，因为后缘流动的变化会影响整个机翼的性能。 在数值实验中特别关注这一点。 大涡的影响不像冲击/边界层相互作用那样容易识别。 在本实验中，选择了具有大量后部弧度的机翼，即RAE5225；它代表了现代客机的机翼部分。 这种机翼已经参与了不同RANS代码的比较测试，风洞测试结果可供比较[1]。 然而，尚未对湍流进行测量。 后弧度增加了通过后缘的流动的重要性和复杂性。 根据其他航空翼型的经验，至少可以列出该流程的一些主要特征。 首先，有一个

空气动力学预测335 边界层在上表面向后缘方向的非常明显增厚。 可以推断，正常应力在这个区域会很大。 接下来，流动是弯曲的，因为在通过后缘时，它必须从沿着翼形表面转向自由流动方向。 对于后弧，这种曲率是明显的；此外，它在后缘的上游之前是相反符号的曲率。 这些曲率预计会影响湍流的生产或阻尼。 同样，在湍流通过后缘时，由于固体表面的存在，对湍流波动的约束突然放松；大涡的结构一定有突然的、明显的变化。 这意味着剪切层的位移效应发生了明显的变化，这将改变压力场和升力。 最后，大涡流和无黏流动之间的后缘存在不稳定的相互作用，如果后缘上游的平均流量接近分离，则可以明显。 为了准确的模拟，所有这些特征都应得到充分的处理。 后两个的处理，如冲击/边界层相互作用的处理，似乎超出了严格的RANS框架的范围。 目前的数值实验从一组跨声速流动经过RAE5225的T17常规模拟开始，以及升力适中且没有冲击的入射率。 对于这些模拟，Hall [2]提出的顶点-重心RANS方案的改进版本与Menter [6]提出的流行的k-u SST 湍流模型耦合。 这个湍流模型被广泛接受为空气动力学流动性能最好的模型之一。 结果与相应的风洞结果进行比较。 注意到了显著的差异。 努力确定这些原因。 这表明，通过修改涡黏性的Menter公式，差异可能会减少，该公式仅限于后缘区域。 最后，为了测试这种改进的可能性，对Menter公式进行了简单的修改，并在给定的马赫数和发生率下通过试验RANS解决方案调整修改函数中的常数，以与隧道结果保持良好的一致性。 然后，在保持修改功能固定的情况下，尝试对更高的入射率进行模拟，其中升力高出60\%，超声速流动的范围要大得多，有明显的冲击，并且风洞结果再次可供比较。 本案获得的同意程度出乎意料地高。 实验以对结果的讨论结束。 用于生成此贡献结果的计算机代码是由作者为法恩伯勒的DERA（原RAE）编写的。 尽管作者不再与DERA有任何正式联系，但非正式交流仍在继续，DERA的代码和计算设施的使用得到了正式承认。 然而，这里表达的观点是作者自己的观点。

\section{17.2 RANS方案和Menter 湍流模型}
336 大厅 17.2 RANS方案和Menter 湍流模型 以下是RANS方案和湍流模型的主要特征摘要。 这里采用的RANS计划基本上没有什么新东西。 这是一个具有顶点-重心离散化的有限体积方案，由Hall [2]于1991年推出，目的是将提高的准确性与提高的稳健性相结合，相对于当时最好的细胞中心和细胞-顶点方案。 自变量在网格单元的顶点上指定，但由单元重心定义的控制体用于整合无黏性通量和黏性通量。 人工耗散的处理是基于第二和第四差异的标准组合，除了使用分割差异外，该组合很好地适用于Euler 方程的解决方案。 该处理现在包括一种减少人工耗散的方法，其中物理耗散本身有助于抑制数值不稳定性；它是Hall [3]之前提出的设备的改进版本。 为了加速收敛，使用标准多网格和残余平滑方案，在精细网格上进行Lax-Wendroff计时，用于上绕元素，在粗网格上进行Runge-Kutta计时。 RANS方案的数值准确性通过比较Hall [4]获得的结果来证明，RAE5225机翼在自由流马赫数 M\textasciicircum{} = 0.735和入射角a = 1.57度，Benton（BAe）[12]、Radespiel（DLR）[12]和Swanson（NASA）[10]在同一C网格上获得的结果。 对于这个RANS验证，采用了简单的Baldwin-Lomax 湍流模型。 下表1显示了每个贡献者的计算升力和阻力系数，用于三个级别的网格细度。 对于最好的网格，不同的贡献者获得的结果几乎相同，微小的差异可能归因于其代码的差异。 检查结果的变化与网格的细化表明了数值方法的准确性。 在这个标准上，可以看到霍尔在256 x 64网格上的结果的准确性与其他网格获得的相应精度相比是有利的；事实上，其他人只能通过在512 x 128网格上计算来匹配霍尔，更精细一层。 应该注意的是，虽然表格结果显示，在RANS计算中通常可以实现非常令人满意的数值精度，但它们也表明，至少在Baldwin-Lomax 湍流模型中，RANS计算对真实流量的模拟非常差。 计算出的压力分布显示中弦周围有明显的冲击，在风洞测量中没有迹象。 M\textasciicircum{} = 0.735和a = 1.59（在未发表的补充[1]中）的测量升力和阻力系数分别为0.4057和0.01068，而不是计算平均值[12]、0.5505和0.00978。 计算的升力高了大约36\%！

关于空气动力学预测337贡献者本顿·拉德斯皮尔·斯旺森大厅网格128 x 32 256 x 64 512 x 128 1024 x 256 128 x 32 256 512 x 128 x 128 1024 x 256 128 x 64 512 x 128 x 128 1024 x 128 x 32 256 x 32 256 512 x 128 1024 x 256 0.5538 0.5546 0.5436 0.5499 0.5499 0.5537 0.5589 0.5550 0.5493 0.553 0.5513 0.5504 0.010291 0.009911 0.009825 0.010460 0.009887 0.009799 0.012892 0.010069 0.009730 0.009745 0.009771 0.009803 表1 Baldwin-Lomax 湍流型号的RAE5225 气翼的升力和阻力系数。 Menter的湍流模型提供了一些改进。 这个线性模型从两个偏微分方程中产生了一个等向同性涡黏性，用于RANS方案，用于湍流能量的对流传输及其比耗散率，分别为k和w。 该模型类似于威尔科克斯[11]的标准k-w 湍流模型，但有两个重要的区别特征。 Menter指出，虽然Wilcox模型比边界层内部广泛使用的k-e模型具有更好的数值稳定性，但其结果对u>小自由流值的变化过分敏感，这是k-e模型自由的一个缺点。 然后，他使用恒等式e = ku>在k-e模型中更改为k,us变量，并将后者版本用于边界层的外部部分，同时保留Wilcox的k-w模型用于内部部分。 通过引入混合函数，公式被简化为一对方程，对于k和u，就像在Wilcox的模型中一样，但menter的w方程包含一个涉及混合函数和交叉导数的额外项。 第二个特征是对Wilcox对涡黏性公式的修改。 Wilcox的公式pk先生=-

\section{Menter 湍流模型的17.3 RANS结果}
338 HALL意味着主剪应力与剪应力成正比。 它不提供剪切应力传输（SST），这是在模拟有分离的边界层或不利压力梯度中的严重缺陷，其中观察到剪切应力与湍流动能大致成正比。 Menter建议，作为补救措施，a\textbackslash\{\}pk »T= r OFV (17.1) max \textbackslash\{\}aiw\textbackslash\{\}ilr \textbackslash\{\}，其中a\textbackslash\{\} - 0.31，O是涡量的大小，F是一个从边界层内部1.0平稳变化到外部零的函数。 对于这里介绍的计算，Menter的k和w方程，以及他用于涡黏性 /x\textasciicircum{}的公式（17.1），用于为Navier-Stokes解的每个多网格周期更新一次变量。2在更新中，当前的Navier-Stokes变量保持固定，而湍流变量则使用顶心离散化和两阶段时间步进方案。 由于在Menter模型的数值实现中，单个程序员之间存在差异，因此与Baldwin-Lomax模型相比，使用Menter模型的RANS解决方案的单个结果之间的差异会更大。 表1所示的标准化验证结果尚未提供。 在目前的实施中，考虑到这样一个事实，即在固体表面附近，耗散w与与表面距离的平方成反比，因为正如将要证明的那样，这种测量提高了数值精度。 17.3 Menter 湍流模型的RANS结果 在本节中，对M\textasciicircum{} = 0.735、a = 1.59和自由流雷诺数 i的RAE5225机翼提供了计算结果与风洞测量的比较？ Eoo - 6 x 106。 表2显示了三个网格级别的计算升力和阻力系数，首先是湍流模型的简单实现，不考虑w的反平方行为，然后实现明确允许行为。 测量值包含在表格的底部。 可以看出，随着网格的完善，简单的实现产生了升力的显著增加和阻力的显著减少。 这些变化比表1中Baldwin-Lomax 湍流模型的变化要明显得多。 另一方面，在明确允许u>的反平方行为时获得的结果显示了网格级别2的变化。这里采用的方程是Menter和Rumsey[7]给出的方程，这些方程在前导、对流项上与Menter [6]最初提出的方程不同。

空气动力学预测339 案例简单配额测量网格128 x 32 256 x 64 512 x 128 128 x 32 256 x 64 512 x 128 0.38222 0.4126 0.4567 0.4671 0.4675 0.4688 0.4057 0.01481 0.01217 0.01099 0.01180 0.01070 0.01049 0.01068 表2计算和测量的升力和阻力系数。 Menter 湍流型号。 RAE5225翼油，Moo = 0.735，a = 1.59，Re\textasciicircum{} = 6 x 106。接近表1中的那些。 事实上，随着网格的完善，“简单”结果倾向于“允许”结果。 对于256 x 64网格，“简单”的升力确实与测量值令人满意地一致，但与网格加密的变化和阻力的差异暴露了这是虚假和误导性的；这里不再考虑这种实现。 因此，对于上述测试条件，虽然Menter 湍流模型的阻力结果非常令人满意，但升力结果仍然高出约15\%。 图1显示了用于比较的机翼表面的计算和测量压力分布。 可以看出，在256 x 64网格上获得的结果与在512 x 128网格上获得的结果基本上无法区分。 然而，上表面的计算和测量压力分布有很大不同，计算的吸力压力明显高于测量值，出现微弱但明显的计算冲击。 只有在后缘才有公平的协议。 事实上，计算压力分布可能是更高的发生率的预期。 详细检查后缘附近上表面的压力分布提供了可能的解释。 请注意，计算的压力保持到后缘的陡峭梯度，而x/c > 0.95时，测量的压力梯度明显降低。 尾缘计算和测量压力分布的这种微小差异是Menter 湍流模型可能存在缺点的线索。 上表面测量的压力梯度的降低通常与边界层的快速增厚有关。 这种增厚将减少翼型流向下的偏转，这可以显著降低升力，就像后缘襟翼向上运动减少升力一方式。 那么，这样的机制是什么

340 HALL -1.4 -1.2 -1.0 -0.8 Cp -0.6 -0.4 -0.2 0.0 0.2 0.4 ' 0.0 0.2 0.4 0.6 0.8 1.0 x/c 图1 计算和测量的表面压力分布。 Menter 湍流型号。 RAE5225翼油，Mx - 0.735，a = 1.59，Re\textasciicircum{} =6x 106. 尚未包含在Menter模型中的快速增厚？ 回忆一下Ashill、Wood和Weeks是如何解决与相同风洞测量有关的类似问题的，这很有帮助[1]。 这些作者报告了通过为该方法的黏性部分引入高阶边界层模型来预测机翼流的交互式黏性-非黏性方法的改进。 他们最初的边界层模型与Menter的湍流模型一样，结合了应力运输的津贴，以应对不利的压力梯度，未能重现后缘观察到的压力梯度降低，并产生了0.477的升力系数，高出约17.6\%。 然后，他们对黏性模型进行了六次不同的半经验修改，试图重现观察到的结果。 其中两个旨在减少边界层近似的基本限制，使其方法具有接近RANS方法的能力。 升力减少了约3\%，降至0.465，正如预期的那样，这接近表1中使用Menter 湍流模型的RANS解决方案的结果。 接下来的三个修改，对于低雷诺数的影响，正常512x128网格256x64网格o实验

\section{17.4 对Menter 湍流型号的修改}
空气动力学预测341应力和流曲率是研究人员纳入涡黏性新非线性模型的类型。 他们一起将升力进一步降低了11\%，达到0.421。 最后，对分离附近速度曲线的特殊形状有特殊津贴，这将升力降低4\%至0.406，即测量水平。 现在，作者将是第一个同意每个修改都有相当大的不确定性，因此其总和的累积不确定性将是相当大的。 此外，可能还有完全无法解释的影响，例如后缘大涡结构的突然变化，以及大涡和无黏流动之间的不稳定相互作用，这些在导言中提到。 异常速度曲线可能是后者不稳定相互作用的症状。 一些效果的模拟将超出严格的RANS公式的范围，尽管可以想象，此类效果的一些考虑可能被纳入湍流模型中。 作者自己提供了证据，证明任何用他们的下一个测试案例RAE5230创建通用模型的尝试都是徒劳的，这是RAE5225的修改，后弧度增加足以在风洞测试条件下产生明确定义的后缘分离。 对于这种情况和类似的情况，他们建议关闭流曲率修改，该修改为RAE5225提供了所需的升力减少的三分之一，以获得最佳效果。 上述情况表明，任何试图通过在方程中添加非线性项，或诉诸应力张量的全传输方程，以通常的方式缩小数值模拟与物理现实之间的差距，都不太可能毫无条件地成功。 如果有足够的努力和计算能力，差距无疑可以缩小，但物理非线性和偏离平衡是明显的，在雷诺平均中失去的物理效应仍然存在。 在这种情况下，对湍流模型的简单修改，在RANS框架中可能是非物理的，但专注于后缘，并考虑到真实、不稳定、大波结构、流程而设计，可能会更有效。 它将把Menter未涵盖的所有重要后缘效应捆绑在一个简单的软件包中，其中一些可以通过使用非线性湍流模型或对应力张量的全传输方程进行建模来处理，其中一些在雷诺平均中丢失了。 这些影响都不会单独考虑。 这里探讨了这种可能性。 17.4 对Menter 湍流模型的修改 为了以简单的方式将Menter模型的拟议修改限制在后缘，控制对流方程保持不变。 相反，

342 HALL Case Menter修改测量网格128 x 32 256 x 64 512 x 128 128 x 32 256 x 64 512 x 128 0.4671 0.4675 0.4688 0.4131 0.4070 0.4069 0.4057 0.01180 0.01070 0.01049 0.01175 0.01051 0.01034 0.01068 表3 计算和测量的升力和阻力系数。 导师和改装的湍流模型。 RAE5225 aerofoil, Moo = 0.735, a - 1.59, Re0 6 x 10b. 将一个修改因子FTE添加到Menter公式（17.1）中，以产生FTBaipk VT = 7 FTFr (17-2> max \{aiw; ill*\} 修改因子FTE可以采取多种形式。 做出选择的唯一标准是，它应该代表真实、复杂、高度互动的流程的平均值。 所选形式没有特别的优点。 它是FTE = 1-A(I)\{1 + G(I,J)\}。 （17.3）这里的I和J分别是网格C线及其横线的间隔计数器。 函数A(I)是上表面0.975 < x/c < 1.0的常数0 < A < 1，并在两个方向上与I线性减少，因此A - 0在大部分翼翼上。 它的作用是减少后缘附近上表面的涡黏性。 添加函数G（I，J）以给出横向方向的适当变化。 涡黏性的变化大小和覆盖区域的空间范围可以通过A和G中的可调常数来改变。 一些试验，以测量的升力水平为目标，然后产生表3所示的升力和阻力结果。 从表3可以看出，计算的升力与测量值进行了令人满意的匹配。 Menter的湍流模型预测的阻力没有严重的错误，也没有试图通过修改来改进它。 256 x 64网格的相应压力分布如图2所示。 不出所料，计算出的压力分布现在与其测量的对应物保持了令人满意的一致。 仔细检查表明，经过修改的湍流模型，一个

航空动力预测343修改-1.2-导师°实验0.0 0.2 0.4 0.6 0.8 1.0 x/c图2计算和测量的表面压力分布。 导师和改装的湍流模型。 RAE5225翼油，Moo = 0.735，a = 1.59，fieoo = 6 x 106。上表面向后缘的压力梯度降低与隧道测试中测量的降低相似。 这表明，修改后的模型在后缘上产生了所需的边界层的明显增厚。 为了测试修改后的湍流模型，它现在用于不同的流程，对涡黏性公式（17.2）中FTE因子的表达式（17.3）中的函数A和G没有任何更改。 计算了在较高入射率下经过RAE5225翼的流量，a = 2.763。 在这个事件发生时，升力高出60\%，整体压力分布非常不同，激波在中弦左右。 在表4中，将由此产生的升力和阻力系数与相应的隧道测试值进行比较，并与使用涡黏性的Menter原始公式（17.1）获得的计算值进行比较。 256x64网格的相应压力分布如图3所示。 计算结果和测量结果之间的一致性质量，

344 霍尔案例 门特修改测量网格 128 x 32 256 x 64 512 x 128 128 x 32 256 x 64 512 x 128 0.6885 0.6946 0.7015 0.6440 0.6572 0.6597 0.6616 0.01387 0.01396 0.01394 0.01280 0.01285 0.01277 0.01293 表4 计算和测量的升力和阻力系数。 导师和改装的湍流模型。 RAE5225翼型，Mx = 0.737，a = 2.763，Reoo = 6 x 106。如表4和图3所示，总体良好。 计算出的升力和阻力误差分别只有0.03\%和1.2\%。 计算出的压力分布与测量的分布总体上非常匹配，但显示后缘的梯度过度减少。 有人可能会认为，这种良好的协议是预期的，因为Menter的湍流模型提供了一个提升，在这种情况下，只高了6\%（与a = 1.59的15\%相比），因此只需要相对较小的修正。 另一方面，在这种情况下，门特模型给出的阻力过高了8\%，与a = 1.59时1.8\%的差异相比，这很大，因此需要对a = 2.763的阻力进行相对较大的校正。 鉴于在涡黏性的修改公式（17.2）的校准中没有努力匹配阻力，当Menter的公式给出糟糕的结果时，阻力预测得如此之好，这似乎令人惊讶。 这当然可能是偶然的，但有一个可能的解释。 在较低的发病率下，没有激波，或者只有非常微弱的电击。 在这样的流动下，对于给定的升力变化，阻力变化相对较小。 在较高的入射点，上表面有一个激波，它为总阻力贡献了自己的波阻。 对于给定的升力变化，减震器将改变其位置和强度；波阻将发生重大变化，因此，总阻力也会发生重大变化。 因此，Menter的模型在较高的入射率下对阻力的估计似乎很差，因为它在阻力对升力敏感的情况下对升力的估计不准确；一旦在后缘使用修改后的涡黏性准确估计了升力，RANS与Menter的模型一起确保了冲击设置在正确的位置，并且很好地预测了阻力。

\section{17.5 结束的讲话}
空气动力学预测345 -1.4 -1.2 -1.0 -0.8 Cp -0.6 -0.4 -0.2 0.0 0.2 0.4 0.0 0.2 0.4 0.6 0.8 1.0 x/c 图3 计算和测量的表面压力分布。 导师和改装的湍流模型。 RAE5225机翼，Mx = 0.737，a = 2.763，Reoo = 6 x 106。 17.5 结论 对Menter的SST 湍流模型进行简单修改的第一次测试与风洞测量有很好的一致性。 然而，这只是湍流模型的实际使用认真开发的第一步。 到目前为止，拟议的修改只针对一种新条件进行了测试，即新的入射角，自由流马赫数和翼型的形状保持在完全指定的修改状态。 需要对一系列条件和一系列航空翼进行测试。 这些可能暗示了进一步的修改。 他们肯定会暴露出这种方法的局限性。 显然，当前类型的修改可以修改到任何涡道黏度类型的湍流模型；它们不限于Menter的涡黏性模型。 也许应该优先考虑开发一个修改了比耗散率u的双方程湍流模型

\section{参考资料}
346 HALL被一个更适合数值计算的变量所取代。 通过求解一个量ui的非线性偏微分方程，推导出涡黏性似乎是反常的，该方程与与表面距离的平方反平方相等。 最后，目前的结果表明，在标准湍流模型中包括一个简单的津贴，用于模型未涵盖的后缘物理效应，可以在空气动力学预测方面产生有价值的改进。 津贴将单个效应捆绑在一起，其中一些可以通过使用非线性湍流模型或对应力张量的全传输方程进行建模来处理，其中一些在雷诺平均中丢失了。 这些效应都没有单独考虑；它们的总和被视为单一的后缘效应。 这种方法的要求是，首先，识别它们的综合物理效应，然后，设计一个适当的数值表示。 参考资料1。 阿希尔，P。 R.，伍德，R. F。 和周，D。 J.，黏性、加拉贝迪安和科恩方法（VGK）的改进半逆版本，RAE TR87002，1987年1月。 2. 霍尔，M。 G.，Navier-Stokes 方程改进有限体积解的顶点-重心方案，AIAA论文91-1540，1991年6月。 3. 霍尔，M。 G.，关于黏性流动解决方案中人工耗散的减少，计算流体动力学-1994的前沿，编辑D。 A. Caughey和M。 M。 哈菲兹，威利，1994年，pp. 303-317。 4. 霍尔，M。 G.，计算的升力和阻力系数，RAE5225，以及RAE2822的计算时间，未发表的DERA承包商说明，1997年12月。 5. 马歇尔，T。 \& Dolling，D。 S.，湍流的计算，分离，未扫过压缩坡道相互作用，AIAA杂志30日，8月。 1992年，pp. 2056-2065年。 6. 门特，F。 R.，空气动力学流动的区域二方程k-u> 湍流模型，AIAA论文93-2906，1993年7月。 7. 门特，F。 R。 \& Rumsey，C。 L.，跨声速流动的双方程湍流模型的评估，AIAA论文94-2343，1994年6月。 8. Muck，K.-C，Andreopoulos，J. \& Dussauge，J.-P.，激波/湍流边界层相互作用的不稳定性质，AIAA杂志26日，2月。 1988年，第179-187页。 9. 斯佩齐亚尔，C。 G.，湍流 时间依赖性RANS和VLES的建模：回顾，AIAA杂志36，2月。 1998年，pp. 173-184。 10. Swanson，R. C，RAE5225 翼型的结果，矩阵消散，私人通信，1995年。 11. 威尔考克斯，D。 C，湍流 CFD建模，DCW工业公司，5354 Palm Drive，La Canada，加利福尼亚州，1993年。 12. 威廉姆斯，B。 R.，2D Navier-Stokes 方程的计算，GARTEUR/TP-067，1月。 1995年。

\clearpage
\chapter{计算空气动力流的算法的18个进步}
\section{18.1 导言}
计算空气动力学流动算法的进步 David W. Zingg,1 Stan De Rango1 \& Alberto Pueyo1 18.1 简介 过去三十年在计算流体动力学（CFD）领域取得的成功往往掩盖了21世纪开始时CFD社区面临的巨大挑战。 如果我们专注于CFD在飞机设计中的应用，或者更具体地说，在这种情况下，雷诺平均Navier-Stokes（RANS）方程的解，可以在以下三个领域确定挑战：计算效率。 必须减少实现适当解决的解决方案所需的计算时间。 由于一体化产品和流程开发环境的趋势[34]以及在空气动力学和多学科设计优化的背景下，这一需求尤为紧迫。 为了完全集成到设计过程中，在三维配置上解决RANS方程所需的时间必须是几分钟。 这大约比当前能力快两个数量级。 虽然提高计算机速度，特别是并行架构，无疑会有所帮助，但也需要改进算法。 在设计优化背景下，算法可靠性也越来越重要。 如果流动求解器没有在设计空间的相关区域收敛，现代设计优化算法，如伴随方法，就不能有效。 人类的效率。 需要减少人力努力和1个航空航天研究所，多伦多大学，多伦多，多伦多，加拿大M3H 5T6。电子邮件：dwz@oddjob.utias.utoronto.ca Frontiers of 计算流体动力学 - 2002 编辑：David A. Caughey和Mohamed M. 哈菲兹 ©2002 世界科学

348 ZINGG、DE RANGO和PUEYO在复杂配置上计算流程所需的专业知识可能更紧迫，因为人类不受摩尔定律的支配。 虽然期望没有CFD和空气动力学知识的用户进行有用的计算是不明智的，但必须尽量减少选择求解器和网格参数所需的专业知识。 目前，从以与计算机辅助设计包相关的格式存储的复杂三维几何中生成计算网格所需的时间和知识是过多的。 尽管非结构化解决方案自适应网格技术在这方面具有巨大的潜力，但在实现其承诺之前，有许多障碍需要克服。 估计全球数值误差是另一个基本问题，它将大大减少用户专业知识要求，但仍有待充分解决。 物理建模的准确性。 虽然可以通过适当的网格分辨率和其他方式仔细控制数值误差，但物理模型（包括湍流模型和层-湍流过渡的预测）产生的误差更难估计和控制。 目前流行用于计算空气动力学流量的涡流黏度湍流模型通常无法准确预测微妙现象，如雷诺数或高升配置的襟翼间隙效应。 看来，雷诺应力模型（或第二时关闭）可能是最有前途的方法。 应该加快将此类模型纳入空气动力学流求解器的努力，至少是为了评估和指导此类模型的开发，如果不是为了现阶段的生产使用。 湍流模型的进一步开发取决于对各种流的全面评估和测试，这反过来又需要减少计算三维流的良好解决方案所需的时间。 此外，需要更多高质量的实验数据集，其中包括计算所需的所有边界条件数据。 在本章中，我们描述和讨论了最近旨在提高RANS求解器计算效率的两项最新进展。 第一个是不精确的Newton-Krylov算法，它减少了实现稳态解所需的计算时间[31]。 第二个是高阶空间离散化，它降低了给定数值精度水平的网格分辨率要求，从而也降低了计算成本[12，51,13]。 这两个主题在格式相似的独立部分中进行涵盖。 每个部分都包含背景材料和算法概述，然后是结果和对所提出问题的讨论。 本章的一个基本主题是需要客观地衡量算法性能。 为了取得进步，人们必须能够衡量它。 尽管总体目标是最大限度地减少计算时间，

\section{18.2 牛顿-克里洛夫算法}
空气动力流的算法349可能很简单，2算法评估是一件复杂的事情，因为依赖硬件和写程序。 无论如何，关键指标是计算时间，而不是经常报告的迭代次数。 在下一节中，我们使用计算时间的规范化度量，虽然它没有考虑到所有因素，但有助于比较在不同计算机上运行的不同算法。 测量空间离散化的性能更加困难，因为它需要计算数值误差的能力。 在描述高阶空间离散化的部分，我们为此广泛使用了网格收敛研究。 18.2 Newton-Krylov算法18.2.1 背景 在稳态RANS方程中的空间导数离散化后，无论是通过有限体积、有限差或有限元法，都可以得到一个非线性代数方程的耦合系统。 Q表示网格每个节点上包含保守变量的向量，我们将此系统写成R\{Q）=0。 （18.1）对于非线性代数方程，考虑牛顿方法是很自然的，它需要在每次迭代时求解一个线性系统，并在特定条件下收敛二次。 这种方法对相对较小的问题很有效[3]，但随着问题规模的增加，与线性问题的直接解决方案相关的缩放和内存使用变得禁止。 这促使人们使用不精确的牛顿方法，在这种方法中，使用迭代方法求解每次牛顿迭代时产生的线性系统。 不精确的牛顿迭代可以写成\textbackslash\{\}\textbackslash\{\}R(Qn) +,4(Q„)AQn|| < Vn\textbackslash\{\}\textbackslash\{\}R(Qn)\textbackslash\{\}\textbackslash\{\}, (18.2)，其中A是R的雅各布，AQn = Qn+i - Qn，Qn是当前解，Qn+i是更新的解。 参数r]n决定了线性系统的迭代解的收敛度，它控制了不精确的牛顿方法的收敛。 方便将不精确的牛顿迭代定义为外部迭代，并将线性系统解所需的迭代定义为内部迭代。 如果r\textbackslash\{\}n对所有n都等于零，即线性系统被精确求解，那么牛顿的方法就恢复了。 可以选择r]n的非零值，即使这是过度简化，因为内存使用也是一个重要的考虑因素。

如[11]所示，保留了350 ZINGG、DE RANGO和PUEYO 收敛。 [14]讨论了选择rjn值序列的进一步考虑因素。 如果使用Krylov子空间方法求解不精确的牛顿框架中的线性系统，则该组合称为Newton-Krylov算法。 尽管非对称线性系统存在许多迭代方法，但广义最小残差算法（GMRES）[39]可能是在这种情况下最受欢迎的。 为了求解线性系统Ax = b，GMRES利用vi, Av\textbackslash\{\},A2vi,...给出的子空间，其中向量vi由初始猜测XQ形成，为b - Axn "' = F\textasciicircum{}y- (18'3) 在每次GMRES迭代时都会发现子空间中最小化残差的向量的线性组合，如果没有达到足够的收敛，则将额外的向量添加到子空间中。 在实践中，一旦子空间达到指定大小，GMRES通常会重新启动，以避免过度使用内存。 Wigton等人[48]和Shakib[40]是第一批在CFD中使用GMRES的人。 [23、45、9、38、30、4、2、10、5、31、27]报告了牛顿-克里洛夫算法的进一步值得注意的发展。 最近的贡献涉及与并行化[24]和多网格整合[25，19]相关的问题。 牛顿-克里洛夫主题出现了几个不同的变体。 我们将在这里描述其中的三个。 第一个与[48]相关，将牛顿-克里洛夫方法应用于应用迭代求解器产生的非线性代数方程组的求解。 例如，如果迭代解技术通过Qn+l = M（Qn）给出的更新公式收敛到R（Q）=0的解，其中运算符M是迭代技术的函数，那么Newton-Krylov方法适用于Q-M（Q）=0。 （18.5）实际上，原始求解器被用作GMRES的预处理器。 这种方法很容易添加到现有的求解器中。 此外，GMRES只需要矩阵向量积，这些积可以使用近似的Frechet衍生物形成，而无需明确定义Jacobian 矩阵。 由于[48]中使用的求解器不需要Jacobian 矩阵的形成和存储，因此生成的算法确实是无矩阵的，并且具有适度的内存开销。 对额外内存的主要需求与GMRES所需的搜索方向相关联，该搜索方向可以使用重新启动的GMRES进行控制。

空气动力学流动算法351在第二种方法中，被称为近似牛顿，Jacobian 矩阵被简化，通常使用与一阶空间离散化相关的雅各布。 Rogers [38]举了这一策略的一个例子。 由于一阶空间离散化只涉及最近邻居，简化的雅各布比完整的雅各布所需的存储要少得多。 此外，它通常有更好的条件，并且可能对角线占主导地位，因此内部迭代的收敛是快速的。 然而，二次收敛的可能性消失了，所需的外部迭代次数大大增加。 在计算时间要求方面，近似牛顿方法不如我们下面描述的无雅各布不精确牛顿方法[31, 5]。 在第三种方法中，GMRES所需的矩阵向量积再次使用近似Frechet导数形成，因此不需要形成完整的Jacobian。 然而，与近似牛顿框架中使用的Jacobian 矩阵近似值的近似值被用作GMRES预调节器的基础[31]。 因此，与所描述的第一个方法相比，这种方法是无雅各布的，但不是无矩阵的。 通常，近似雅各布的某种形式的不完全下上（ILU）因数分解被用作预调节器。 虽然这种方法存在与记忆惩罚相关，但它可以非常快地进行，我们将在下面看到。 18.2.2 算法 在本小节中，我们介绍了牛顿-克里洛夫算法的各种组件，该算法属于上述雅各布自由类。 虽然基本算法相对简单，但几个细节对其成功至关重要。 为了提供易于解释的结果，我们通过对残差的单次评估的成本将计算时间归一化。 诚然，这是一个不完美的尝试，允许硬件差异，在较小程度上，允许不同的空间离散化。 空间离散化空间离散化基于通过广义曲线坐标变换应用的有限差异。 二阶中心差分与二阶和四阶差分阶序量人工耗散的混合结合使用。 二次差分耗散仅在近冲击时是显著的。 Baldwin和Lomax的代数湍流模型用于确定涡黏性。 获得的稳态溶液与成熟的流动求解器 ARC2D的溶液相同[32]。 对于所有示例，网格具有C拓扑结构。

352 ZINGG, DE RANGO \& PUEYO 初始阶段 牛顿方法的一个缺陷是它不是全局收敛的。 有几种策略可以解决这个问题。 一种方法是基于这样的观察，即牛顿的方法是从隐含的欧拉时间行进方法中获得的，当时间步骤走向无穷大时，局部时间线性化应用于常微分方程dft = R(Q) (18.6)的极限。 因此，隐式欧拉方法最初可以与有限时间步数一起使用，随着残差的减少，时间步进会增加。 通常，使用近似牛顿方法更有效，直到时间步骤有效无限，此时可以启动无雅各布策略[5]。 对于目前的结果，我们使用近似因数分解算法[32]的对角线形式和网格测序来处理初始迭代。 残差首先在粗网格上减少两个数量级（或最多150次迭代），然后在精细网格上减少5次迭代。 这为所研究的所有案例产生了一个稳定的算法，初始阶段的计算时间非常短。 不精确的牛顿策略在上述初始阶段之后，我们使用以下j]n值：•前10个外部迭代的Vn-0.5，•剩余迭代的Vn = 0.1。 虽然这种方法产生了线性而不是二次收敛，但它在计算时间方面非常有效。 图。 1显示了j]n值对计算时间的影响，i]n保持不变，对于一个具有代表性的测试案例。 很明显，线性迭代的严格容忍度对计算时间没有好处，尽管它减少了外部迭代的次数。 超过一定点，线性系统的残差的进一步减少对非线性系统的残差没有影响。 为r]n选择的策略避免了这种情况，这被称为过度解决线性系统[14]。 Jacobian-Free GMRES为了形成Krylov子空间，GMRES需要Jacobian 矩阵 A（Qn）和任意向量v的乘积。 这可以在不实际计算A(Qn)v=R\textasciicircum{} + ev)e-R\textasciicircum{}\textbackslash\{\}（18.7）的雅各布时获得

空气动力学流动算法353，其中e从[30] e|M|2 = v\textasciicircum{}> (18-8)计算，em是机器零的值。 预处理 使用近似 Frechet 导数提供了 Jacobian 矩阵 和任意向量的乘积的出色近似值。 与所用空间离散化产生的残余函数R（Q）相关的Jacobian 矩阵既条件差，又对角线外显性。 这些特性导致GMRES停滞不前，几乎没有残留物减少。 收敛可以通过预处理这个基本系统得到极大的改进。 事实上，预处理器的选择可能是决定整体算法性能的最重要的单一因素。 使用正确的预处理，系统Ax = b变为AM\textasciitilde{}lMx = b，（18.9），其中M通常是对A的近似值，比A更容易反转。 这个想法是，AM-1的特征值比A的特征值更集群，从而提高了GMRES的性能。 请注意，预处理器可以用作迭代求解器，反之亦然。 例如，在Wigton [48]的方法中，迭代求解器，如近似因子隐式方法或多级多重网格法，用于系统预处理。 这完全避免了Jacobian 矩阵的形成和储存。 相比之下，另一种基于不完全上下（ILU）分解的流行预处理器需要存储Jacobian 矩阵或某种合理的近似值。 流行的是，ILU因式分解基于与一阶空间离散化相关的近似Jacobian 矩阵，而不是与高阶离散化相关的真正的雅各布。 最初的动机是减少存储和操作计数，但Pueyo和Zingg [31]表明，使用近似的Jacobian也会导致更有效的预调节器。 使用真正的雅各布的困难在于，不完整的L和U因子可能条件很差[16, 6]。 因此，尽管M可以很好地近似于A，但数值计算的乘积AM-1可能与恒等矩阵相去甚远[31]。 由一阶空间离散化形成的近似雅各布，它非常耗散，倾向于更对角线显性，并产生更好的条件不完整的L和U因子。 当使用上风空间离散化时，定义一阶离散化和相应的雅各布是直接的

354 ZINGG、DE RANGO和PUEYO用于形成预调节器。 在当前的空间方案中，由于第四差分人工耗散项，对最近邻居的依赖性产生了。 在形成近似的雅各布时，仅使用二差分耗散，系数el2，由el2 = er2 + aer4（18.10）确定，其中eJj和e\textbackslash\{\}分别是实际的第二差和第四差耗散系数，a是用户指定的系数。 增加a使近似的雅各布在对角线上更加占主导地位，但可能比真正的雅各布的近似值差。 图。 2显示了将外部残差减少12个数量级所需的GMRES迭代总数，作为下面讨论的四个测试案例的a的函数。 基于这些和其他类似的结果，a - 5被选为良好的通用值。 ILU因数分解的形成方式与LU因数分解相同，但某些非零条目被丢弃。 最简单的这种因式分解是ILU(O)，其中L和U因子中唯一保留的条目是那些被分解的矩阵具有非零条目的位置。 因此，L和U因子所需的存储与矩阵的存储相同。 根据填充水平或阈值策略，允许一定数量的额外填充，可以形成更准确的ILU因数分解。 在前者中，仅使用矩阵的图形来确定要保留哪些条目。 在阈值策略中，也考虑了条目的大小。 经过大量实验，其中一些记录在[31]中，我们选择了具有两个填充级别的填充级别策略，即ILU（2）。 请注意，近似Jacobian中的条目是4乘4块。 我们不使用块ILU，而是使用标量版本，但如果一个块至少包含一个非零条目，我们将其所有条目标记为非零。 我们把结果分解块填充称为ILU（2）或BFILU（2）。 图。 3显示了剩余收敛历史，根据典型测试案例的剩余函数评估中测量的计算时间绘制。 由近似的雅各布Al形成的BFILU（2）产生的收敛比由Al或真正的雅各布A2形成的BFILU（O）快得多。 影响ILU因式分解性能的另一个因素是网格节点的排序。 根据[31]中报告的研究，反向Cuthill-McKee（RCM）排序[8]几乎无处不在，并在这里使用。 然而，我们获得了一些初步结果，这些结果表明，身体动机的订购可能更有效。 最后，我们只计算一次预调节器，在初始阶段之后使用近似分解算法。 这在[31]中显示，在不损害GMRES的收敛的情况下显著减少计算时间。

空气动力学流算法 355 18.2.3 结果和讨论 所描述的Newton-Krylov算法的性能针对四个湍流流，三个关于NACA 0012 翼型，第四个关于RAE 2822 翼型。 流动条件是：案例1-攻角：0度，自由流马赫数：0.3，雷诺数：288万 案例2-攻角：6度，自由流马赫数：0.3，雷诺数：288万 案例3-攻角：1.49度，自由流马赫数：0.70，雷诺数：900万 案例4-攻角：2.31度，自由流马赫数：0.729，雷诺数：650万 关于NACA 0012 翼型的网格有一个C拓扑，在溪流方向有331个节点，在正向方向有51个。 从翼型表面到第一个节点的距离是1 x 10-5弦，导致最大细胞长宽比超过1 x 105。 RAE 2822 翼型的网格在溪流方向有321个节点，在正向有49个节点，具有相同的墙外间距。 图4至图7显示了绘制的剩余历史，作为计算时间的函数，由单次剩余计算所需的时间归一化。 请注意，这只是一种归一化，以便提供与其他求解器进行比较的手段；它并不反映剩余评估的实际数量，后者要少得多。 为了提供参考，还提出了近似因子多网格算法（表示为AF-MG）的结果，该算法通常比仅近似因子化算法快四到六倍[7]。 多网格算法使用三网格锯齿循环。 Newton-Krylov算法在案例1、3和4中收敛了十二个数量级。 对于案例2，在减少大约八个订单后，剩余部分会挂起。 这是由Baldwin-Lomax 湍流型号引起的。 冻结湍流 涡黏性会导致收敛完成。 在所有情况下，Newton-Krylov算法的收敛速度都比近似因子多网格算法快。 预处理器BFILU（2）也可以用作求解器。 图6显示，它可以表现非常好，实际上收敛速度比近似因子多网格算法快。 然而，它并不可靠，在某些情况下有所不同。 图中还绘制了没有多网格的近似因数分解算法的剩余历史。 6. 它被认为比BFILU（2）效率低得多。 因此，很明显，BFILU（2）是Newton-Krylov算法的关键组成部分。 此外，这表明

\section{18.3 高阶空间离散化}
356 ZINGG、DE RANGO和PUEYO使用近似因数分解算法作为GMRES预处理的完全无矩阵策略不太可能具有竞争力。 接下来，我们使用265 x 49的网格重新审视案例2，墙外间距为1.2 xlO-6。 减少的墙外间距可以改善边界层的分辨率，并减少阻力中的数值误差。 此案例的剩余历史如图所示。 8. Newton-Krylov算法在相当于1000次剩余函数评估的计算时间内将残差减少了12个数量级，实际上比图所示的上一个网格上的收敛略快。 5. 相比之下，墙外间距的减少（以及相关的单元长宽比的增加）对近似因子多网格算法的收敛产生了相当大的不利影响。 对于所研究的案例，由于使用了松散的内部公差（r）n），十二阶残差还原的外部迭代次数从17到22不等。 内部迭代总数从227到370个不等，每个外部迭代的平均内部迭代数在12到20之间。 奔腾Pro 180处理器上的总计算时间从不到6分钟到近9分钟不等。 总体而言，Newton-Krylov算法比近似因子多网格算法收敛得更快、更可靠。 18.3 高阶空间离散化 18.3.1 背景 降低计算成本的补充手段是提高空间离散化的准确性，从而能够使用更粗的网格，同时将数值误差保持在适当的阈值以下。 高阶有限差分方法有着悠久的历史，[26]和[43]提供了大部分初始动机。 Gustafsson等人[21]提供了一个有趣的讨论，包括一些早期的参考资料；进一步的分析和比较可以在[17]和[55]中找到。 大部分显示高阶方法效率的分析是基于具有均匀网格和周期性边界条件的简单标量方程。 在具有非均匀网格的RANS方程的解中，更通用的边界条件，并添加了数值耗散，情况并不那么清楚。 因此，在实际使用中，绝大多数RANS流求解器都基于二阶离散化（有时对非黏性项使用三阶方案）。 在解决需要高精度才能使计算可行的问题时，使用高阶方案更受欢迎，例如过渡模拟和湍流 [33，50]、电磁学[54，41]和航空声学[47，15]。 将高阶方法应用于实际空气动力学流量是

空气动力学流量算法357最初受到许多早期RANS求解器使用与[22]相关的标量人工耗散方案来提供稳定性所需的数值耗散这一事实的限制。 例如，如[1]所示，该方案过于耗散，是高雷诺数下边界层流动计算的主要误差来源。 在保留标量人工耗散方案的同时，将离散化提升到更高阶是徒劳的。 因此，向上风方案[37]和先进的人工耗散方案[42]的开发对于成功实施高阶方法至关重要。 这说明了提高空间离散化准确性的一个关键方面：必须解决最大的错误来源。 不幸的是，这不一定容易确定。 自引入复杂的数值耗散方案以来，几位研究人员将高阶方法应用于实际的空气动力学问题。 [35、44、46、12、13]给出了使用结构化网格和曲线坐标变换的示例。 在每种情况下，在效率方面都有显著的好处。 为了实现这些好处，必须解决离散化的以下方面：•不黏性通量，包括人工耗散或过滤，•曲线坐标变换的度量，•黏性通量，•湍流模型中的对流和扩散通量，•近边界算子，•边界外推，•区域界面插值，•力和力矩计算的积分。 确定与高阶空间离散化相关的效率提高需要一些方法来计算在给定网格上计算的解决方案的数值误差。 目前，实现这一目标的最可靠的方法是使用在更精细的网格上计算的解决方案作为参考解决方案。 假设较细网格上的误差比较粗网格上的误差小得多，较粗网格上的误差只是两个解决方案之间的差异，从而提供局部和全局误差估计。 这目前只在二维度上是实用的，不幸的是，它并不像听起来那么简单。 使用极其精细的网格可能会导致收敛的困难，这既是因为数值刚度增加，也因为鳞片可能会被解决，这些刻度往往不稳定，例如在翼型的后缘附近。 人们更喜欢系统的网格加密研究，允许使用Richardson外推[52]，但实现合适的渐近行为可能非常困难（或需要过于精细的网格）[36]。 例如，[51]中提出的一些网格加密研究显示了明确定义的渐近行为，而另一些则没有。 此外，Zingg等人[51]表明，不同的空间离散化可以

358 ZINGG、DE RANGO和PUEYO即使在极细的网格上也能产生令人惊讶的大解决方案差异。 流动求解器的渐近收敛行为很难确定，原因有很多。 其中一个原因是有几个不同顺序的错误来源。 例如，远离冲击和不连续性的数值耗散通常随着网格间距缩放到三阶，而黏性和不黏性通量导数近似值的误差通常是二阶。 此外，可能存在一阶的局部误差源，例如在边界附近添加的数值耗散和在边界附近使用的一些近似值。 在足够精细的网格上，一阶误差源必须占主导地位。 然而，它们可能比网格上的高阶误差更小，这是实用的。 在粗网格上，三阶数值耗散可能占主导地位。 进一步加重难度的是网格和流动奇点的存在，例如翼型的后缘。 当存在这些时，网格加密研究中获得的错误行为可能与离散化的局部截断误差不对应，特别是高阶离散化。 网格的精炼方式和几何基础表示的连续程度也会影响实现的准确性顺序。 奇点和不连续的存在似乎会使使用高阶方法无效，因为高阶行为无法实现。 然而，[35、44、46、12、13]中提出的结果表明，高阶方法仍然有益。 在奇点处引入的误差通常是局部的，对周围的解决方案影响不大。 因此，尽管在存在奇点的情况下，高阶算法可能会产生二阶收敛，但它仍然可能产生比二阶算法更小的误差。 因此，我们的重点是给定网格上误差的大小，而不是准确性的有效顺序。 18.3.2 算法 目前的高阶空间离散化使用对角线形式的近似因数分解算法[32]实现，以迭代到稳态。 湍流使用代数Baldwin-Lomax模型或单方程Spalart-Allmaras模型进行建模。 这里给出了空间离散化的概述；详情见[12]和[13]。 非黏性通量、数值耗散和网格度量 一阶导数的任何有限差分近似都可以写成斜对称算子和对称算子的总和。 当应用于线性对流方程时，误差可以根据方案的傅里叶分析进行分类。 如果谐波函数的振幅衰减，误差

空气动力学流动算法359被归类为耗散性。 如果相位速度与实际相位速度不同，从而成为波数的函数，则误差被描述为色散误差。 耗散误差与运算符的对称部分相关。 如果有限差分方案是纯粹的斜对称，即算子的对称部分为零，那么半分数方案是非耗散的。 色散误差与倾斜对称部分相关，该部分不能为零。 因此，所有有限差分方案都会产生色散误差，除非在理想化条件下，例如特定的Courant数。 人们普遍认为，Euler 方程的空间离散化必须包含一些数值耗散，即对于收敛和稳定性，运算符的对称部分必须非零。3然而，不知道需要多少耗散。 高分辨率方案基于最小化近不连续性增加的一阶耗散，同时满足总变异、正性或单调性方面的一些标准。 在流动的平滑区域中，将数值耗散最小化的可比理论方法是非常可取的。 基于熵的方法显示出一些潜力[29]。 在解决RANS方程时通常也需要数值耗散。 尽管这些方程包含物理耗散机制，但根据网格分辨率可能需要进一步的数值耗散。[18]探索了使用单元格雷诺数作为数值耗散的缩放参数，而[49]出于同样的目的使用Baldwin-Lomax 湍流模型的涡量函数。 尽管Zingg等人[51]表明，当与二阶中心微分相结合时，三阶矩阵耗散并不是主要的误差来源，但当使用高阶微分方案时，可能需要重新审视这种缩放。 领先耗散误差项总是奇数顺序，而领先分散误差项总是偶数。 对于给定的对称模版，色散误差总是可以比耗散误差高一个阶。 然而，对称算子可以乘以小系数，因此耗散误差可能小于有限波数的色散误差。 这是[54]中采用的策略，使用七点模版。 五阶对称算子产生的耗散误差小于六阶斜对称算子的色散误差，但波数非常小（无论如何误差可以忽略不计）。 耗散误差和色散误差作为波数函数的类似行为导致了理想的属性，即只有具有大色散误差的模式才会严重耗散。 在本求解器中，我们使用五点模版，导致三阶耗散误差和四阶色散误差。 斜对称[28]就这一点提供了一些有趣的讨论。

360 ZINGG, DE RANGO \& PUEYO 运算符是 \{5au 基本对称 \{5su)j = •)•- 1 ( 'h" 12Axl ric 运算符 e4, -Uj + 2 是 4uj+i + 8uj+i + 6UJ - - 8Uj。 4uj\_i -i + + Uj。 Uj (18.11) (18.12) 其中e4是一个系数，通常是0.02。 请注意，e4 = 1/12产生流行的三阶上风偏置算子。 所用矩阵耗散方案的详细信息，包括用于在冲击附近添加一阶耗散的压力开关，可以在[42]和[12]中找到。 曲线坐标变换的度量是使用上述倾斜对称运算符计算的。 [51]表明，网格指标运算符的精度应与非黏性通量运算符的倾斜对称部分的准确性相对应。 如果指标使用二阶近似值，准确性可能会大大降低。 黏性通量 黏性通量离散化有许多可能的策略，形式如下：dx\{adxp）。 （18.13）一些作者[46]只是将用于非黏性通量的高阶一导算子两次，没有明显的并发症。 然而，通常应避免这种方法，因为它会产生高波数阻尼不良（对多网格尤其重要）和比必要的更大的模板。 四阶近似值可以用五点模版实现。 然而，为了方便编程，我们使用以下方案，这导致了一个七点模版。 术语dx/3j首先使用以下四阶表达式在j + 1/2近似：(8xP)j+i/2 = 24\textasciicircum{}(\textasciicircum{}-1-27\textasciicircum{} + 27\textasciicircum{}+1-\textasciicircum{}+2)。 （18-14）系数a使用四阶插值公式在j + 1/2处计算：然后完成算子变成：（ij+i/2 = T\textasciicircum{}（-aj-i + 9a,- + 9aj+1 - aj+2）。 (18.15) 8x(ajSxPj) = 24\textasciicircum{}taJ-3/2(\textasciicircum{}\textasciicircum{})j-3/2 - 27<X,\_i/2(£x/3);,-l/2 +27aj+1/2(6xP)j+i/2 \textasciitilde{} aj+3/2(\$x/3)j+3/2\}- (18.16)

空气动力学流量算法361 湍流模型在Baldwin-Lomax模型中，涡量必须以j + 1/2计算。 这是使用Eq完成的。 18.14。 在Spalart-Allmaras模型中，扩散通量像Eq一样离散化。 18.16. 对流项使用一阶上绕进行离散化，以保持涡黏性的积极性。 对流项的三阶近似实验显示准确性没有提高。 尽管如此，随着其他来源的继续得到解决，这种一阶处理最终可能会成为错误的主要来源。 近边界运算符 与使用高阶方法相关的问题之一是需要大模版产生的近边界算子。 很难找到与内部方案一起稳定的适当精度的近边界算子，因为与入射波相关的边界算子变得顺风偏置[53]。 然而，降低边界方案的准确性可能会破坏高阶内部方案的好处。 在目前的演算法中，我们能够在保持稳定性的同时使用三阶边界方案。 以下三阶偏置算子用于近似内部节点的一导数，该节点的边界节点作为其j -1的邻居：（Sxu）j = -- (-2uj\_i - 3UJ + 6uj+i - uj+2）。 （18.17）相同的公式用于计算此类节点的网格度量。 对于边界节点的网格度量，使用以下方程：(M)j = \textasciicircum{}\textasciicircum{}(-UuJ + 18ui+i " 9ui+* + 2\%+s)。 (18.18) 近边界节点的耗散算子是- (-«,--i + 2,Uj - 2,UJ+1 + uj+2)。 （18.19）对于黏性通量，对应于方程的近边界算子。 18.14、18.15和18.16分别是，（<y\%+i/2 = 24\textasciicircum{}(-23\& + 21\&+1 + 3ft+2-\&+3)，（18.20）etj+i/2 = e（3aj + 6aj+i - aj+2），（18.21）

362 ZINGG、DE RANGO \& PUEYO和Sx(aj6x/3j) = \textasciicircum{}-\textasciicircum{}[-23aj\_1/2(5x/3)j\_1/2 + 2laj+1/2(6x/3)j+1/2 +3aj+3/2(5x/3)j+3/2 - aj+5/2(5x/3)j+5/2\textbackslash\{\}。 （18.22）边界外推 在流动为亚音速的远场边界上，解决方案是根据规定的自由流动条件和从内部推断的状态的组合确定的，该状态取决于流动的方向[51]。 外推状态使用以下公式找到：Uj - Siij+i - 2>Uj+2 + \%+3- (18.23) 在翼型表面，压力是从表达式Pl = \textasciicircum{}-（18p2-9p3 + 2p4)，（18.24）计算出来的，该表达式来自dp/dn - 0的三阶近似值。 对于隔热表面，密度是使用相同的公式找到的。 我们还使用了Eq。 18.23将压力推断到翼型表面，结果没有变化。 区域接口插值在[35]中，高阶方法适用于重叠网格，因此必须相当注意区域之间的插值。 在本研究中，使用单块C网格，在尾迹切割上也出现了类似的问题。 虽然唤醒切口可以被视为内部点，但它目前被视为接口。 以下插值公式用于根据上面和下面的解计算尾迹切（用下标wc表示）上的解：Qkwc = g（-9fcu，c+2 + Mkwc+i + 4gfctl）C\_i - qkwc-2）（18.25）力和力矩计算积分 给定使用高阶空间离散化计算的解，应使用高阶积分技术来计算升力和阻力。 这是通过将立方样条插值拟合到翼型几何形状、表面压力和表面的皮肤摩擦力来实现的。 四阶单边运算符用于计算皮肤所需的导数

空气动力学流量363摩擦的算法。 基于样条插值，使用具有全局误差控制策略的自适应两点高斯-勒让德正交来执行积分。 线性化 近似因数分解算法的对角线形式导致标量带线性系统的求解。 黏性通量的处理基于五点模版，导致五对角线系统。 我们的黏性算子基于七点模版，因此我们使用标准二阶黏性算子的线性化，因为黏性项对对角线形式左侧的贡献在任何情况下都是非常近似的。 近似线性化也用于位于五对角结构之外的近边界算子。 De Rango和Zingg [12，13]包括一些剩余的收敛历史，将高阶算法与也使用矩阵耗散的二阶算法进行比较。 在某些情况下，收敛几乎相同，而在其他情况下，高阶算法收敛速度较慢。 虽然这可能与黏性项线性化的近似值有关，但它更可能是由与高阶算法分辨率的提高相关的刚度增加引起的。 18.3.3 结果和讨论 在本节中，我们介绍了一些网格收敛研究的结果，以证明高阶空间离散化的相对性能。 高阶离散化（在图中表示为高阶）与离散化进行了比较，使用矩阵耗散的二阶中心差，以及低阶边界方案和二阶积分程序（在图中表示为二阶）。 在某些情况下，还显示了三阶上风偏置通量差分割方案的结果，该方案具有二阶黏性项、网格度量和积分（在图中表示为三阶上风）。 对于第一种情况，我们考虑了NACA 0012 翼型的亚声速流动，攻击角为12度，基于288万弦的雷诺数，以及0.16的自由流马赫数。 层-湍流过渡分别固定在上表面和下表面的0.05和0.8弦。 流动的特点是靠近后缘的上表面有一个小区域的分离流动。 在这些条件下，实验测量的升力系数为1.240，阻力系数为0.0180 [20]。 表1描述了所使用的网格。 网格A是使用椭圆网格发生器生成的，是一个非常密集的网格，提供了一个参考

364 ZINGG, DE RANGO \& PUEYO 网格 A B 尺寸 1057 x 193 529 x 97 265 x 49 墙外间距（xlO-6）0.23 0.53 1.2 前缘聚类（xlO-3）0.1 0.2 0.4 尾缘聚类（xlO-3）0.5 1.0 2.0 表1亚音速情况的网格。数值误差小的解决方案。 网格B是通过从网格A中删除每个坐标方向的每个第二个节点来生成的，网格C同样是从网格B生成的。 因此，这三个网格形成了一个适合网格收敛研究的序列。 我们主要对网格C计算中的数值误差感兴趣，这是实践中使用的网格的典型特征。 图。 9显示了使用上述Baldwin-Lomax 湍流模型的三种离散化在网格A、B和C上计算的升力、压力阻力和摩擦阻力系数。 结果与1/iV相比绘制，其中N是网格节点的总数。 在网格A上计算的解决方案彼此非常接近，证实了这些解决方案的数值误差很小。 因此，可以通过与网格A解决方案进行比较来估计在网格C上计算的解决方案中的误差。 所有三个离散化在网格C的升力系数中都产生了不到百分之一的误差。 阻力组件中的误差要大得多。 在网格C上，二阶方案产生的压力和摩擦阻力误差分别超过40\%和12\%。 这些错误是相反的符号，但总体阻力误差超过20\%。 三阶上风方案产生的误差比二阶方案小一些，但误差比高阶格式大得多，这表明，除非网格指标和其他近似值的准确性得到提高，否则提高通量外推的顺序并不特别有效。 在网格C上使用高阶格式计算的阻力分量中的误差小于2\%，并且小于使用网格B上其他方案计算的误差，该方案的节点数是网格B的4倍。 仔细检查解决方案后发现，低阶方案导致对上表面边界层厚度的过度预测，这与对压力阻力的过度预测和对摩擦的低估的预测一致。4我们对使用网格独立一词犹豫不决，因为这可能意味着网格独立于机器精度，但情况当然不是。 然而，在网格A上计算的解决方案可以被认为是在合理的公差范围内独立于网格的。

空气动力流的算法365拖动。 这如图所示。 10，它以85\%的弦在上表面显示计算的边界层速度曲线。 使用三个离散化计算的网格A解是无法区分的；显示了二阶解。 网格A的每四点绘制一个解，对应于网格C上的点。 在网格C上计算的高阶解非常接近网格A解，而其他解决方案的误差很大，与阻力分量中的误差一致。 接下来，我们使用Spalart-Allmaras模型重新审视此案例，以检查湍流模型的离散化的影响。 结果如图所示。 11. 这些趋势与使用Baldwin-Lomax模型观察到的趋势相似，表明与Spalart-Allmaras模型中对流项的一阶离散化相关的误差并不大。 再次，在网格C上计算的高阶结果比网格B上的二阶结果更准确。 顺便说一句，两个湍流模型都会导致与实验值相比对升力的高估和对阻力的低估，Spalart-Allmaras结果更受欢迎。 有趣的是，由于数值和物理模型误差之间的取消，使用二阶离散化计算的网格C结果实际上最接近实验值。 这强化了这样的想法，即准确的数值解决方案对于成功的湍流模型评估至关重要。 无花果。 12和13显示了RAE 2822 翼型上两个跨声速流动的类似网格收敛研究。 两种情况的流动条件如下：跨声速情况I-攻角：2.31度，自由流马赫数：0.729，雷诺数：650万跨声速情况II-攻角：2.57度，自由流马赫数：0.754，雷诺数：650万 对于这两种情况，两个情况，过渡固定在0.03弦。 病例II的特点是上表面有更强的激波和更大的冲击诱导边界层分离区域。 再次使用三个网格的序列，其特征如表2所示。 Spalart-Allmaras 湍流型号用于这两种情况。 无花果。 12和13表明，与二阶方案相比，高阶离散化导致误差显著减少，通常在网格C上产生精确到2\%以内的解。 例外是案例2的压力阻力，其中高阶解的误差接近4\%，而使用二阶方案计算的解的误差刚刚超过5\%。 即使对于高阶格式来说，在冲击时分离气泡附近也有可能没有足够的网格分辨率。 无花果。 14显示了计算的一部分

\section{18.4 结束的讲话}
366 ZINGG、DE RANGO和PUEYO网格A B尺寸1025 x 225 513 x 113 257 x 57壁外间距（xlO"6）0.23 0.53 1.2前缘聚类（xlO"3）0.1 0.2 0.4尾缘聚类（xlO-3）0.25 0.5 1.0表2跨声速情况的网格。 案例II的上表面的压力系数分布。 在网格C上使用高阶离散化计算的解决方案比在网格C上使用二阶方案计算的解决方案更接近网格A解决方案。 这表明，尽管压力阻力误差略有减少，但高阶格式导致局部解误差大幅减少。 图证实了这一点。 15，它显示了95\%和弦上表面计算的边界层速度曲线。 使用二阶方案在网格C上计算的速度曲线的误差相当大，而高阶结果的误差虽然可见，但很小。 在图中。 15，我们还绘制了使用高阶离散化计算的速度曲线，使用二阶方案离散的黏性项。 这表明，将黏性术语提高到更高阶约占总误差减少的10\%。 结果表明，与类似的二阶离散化相比，呈现的高阶空间离散化导致给定网格上获得的数值误差大幅减少。 这是通过每次迭代增加约6\%的计算成本和收敛速率的小幅减少来实现的。 在许多情况下，在给定网格上计算的高阶解比在节点数四倍的网格上计算的二阶解更准确。 这转化为计算时间至少节省了四倍。 18.4 结论 我们提出了两个互补的算法进步，这些进步降低了空气动力学配置流动的RANS方程解的计算成本。 基于二维流的结果，这两种算法都显示出足够的潜力来证明扩展到三维和比更复杂的物理学是合理的

\section{致谢}
\section{参考资料}
空气动力学流量算法 迄今为止研究了367个简单的单元素翼型案例。 需要进一步开发和测试，以确保这些算法在实际使用中足够可靠。 致谢 作者很高兴为Robert W贡献这篇论文。 MacCormack，计算流体动力学的真正先驱。 第一作者很荣幸地被算作教授。 MacCormack作为朋友，希望感谢他从他那里得到的帮助。 所描述的工作是在庞巴迪航空航天和加拿大自然科学与工程研究委员会的财政支持下进行的。 参考资料1。 Allmaras，S.R.，Navier-Stokes解决方案中人工耗散对层边界层的污染，流体力学的数值方法，M.J. 贝恩斯和K.W. Morton，编辑，Calerndon，牛津，英国，1993年，pp. 443-449。 2. 安德森，W.K.，劳什，R.D. \& Bonhaus，D.L.，非结构化网格上不可压缩湍流流的隐式/多网格算法，J. 补偿。 物理。 128，1996年，pp. 391-408。 3. 贝利，H.E. \& Beam，R.M.，牛顿方法应用于稳态可压缩Navier-Stokes 方程的有限差分近似，J. 补偿。 物理。 93，1991年3月，pp. 108-127。 4. 巴特，T.J. \& Linton，S.W.，可压缩流体流动的非结构网格牛顿求解器及其并行实现，AIAA论文95-0221，1月。 1995年。 5. 布兰科，M。 \& Zingg，D.W.，非结构化电网的快速牛顿-克里洛夫求解器，AIAA J.，36，1998年4月，pp. 607-612。 6. 布鲁阿塞特，上午，特维托，A。 \& Winther，R.，关于松弛不完全LU因数分解的稳定性，计算数学54，1990年，pp. 701-719。 7. 奇肖姆，T。 \& Zingg，D.W.，空气动力学流动近似因子算法的多网格加速，加拿大航空和太空J. 45，1998年，pp. 239-246。 8. Cuthill，E.H. \& McKee，J.M.，减少稀疏对称矩阵的带宽，计算机协会第24届全国会议论文集，Brondon Press，1969年，pp. 157-172。 9. 德格雷斯，G。 \& Issman，E.，Krylov子空间方法对可压缩流动求解器的加速，讲座笔记，von Karman Inst，流体动力学系列讲座1994-05，1994年。 10. Delanaye，M.，Geuzaine，P.，Essers，J.A. \& Rogiest，P.，非结构化网格上的欧拉和Navier-Stokes 方程求解二阶有限体积方案，AGARD第77届流体动力学小组研讨会，计算流体动力学方法和算法的进展和挑战，西班牙塞维利亚，10月。 1995年。

368 ZINGG, DE RANGO \& PUEYO 11号。 登博，R.S.，艾森斯塔特，S.C. \& Steihaug, T., Inexact Newton Methods, SIAM J. on Numerical Analysis 19, 1982, pp. 400-408。 12. 德兰戈，S。 \& Zingg，D.W.，空气动力学计算高阶算法的进一步研究，AIAA论文2000-0823，1月。 2000年。 13. 德兰戈，S。 \& Zingg，D.W.，湍流空气动力学计算的高阶空间离散化，AIAA J.，2000年3月提交。 14. 艾森斯塔特，S.C. \& Walker，H.F.，在不精确的牛顿方法中选择强制项，SIAM J.关于科学和统计计算17，1996年，pp. 16-32。 15. Ekaterinis，J.A.，气体动力学和空气声学的隐含、高分辨率、紧凑方案，J. 补偿。 物理。 156，1999年，pp. 272-299。 16. Elman，H.C.，不完整的LU分解的A 稳定性分析，计算数学47，1986年，pp. 191-217。 17. Fornberg，B.，关于双曲方程积分的傅里叶方法，SIAM J.关于数值分析12，1975年，pp. 509-528。 18. 弗鲁，K.，Zingg，D.W. \& DeRango，S.，人工耗散黏性翼型计算方案，AIAA J. 36，9月 1998年，pp. 1732-1734。 19. Geuzaine，P.，Lepot，I.，Meers，F. \& Essers，J.-A.，用于计算非结构化网格上可压缩流的多级牛顿-克里洛夫算法，AIAA论文99-3341，1999年6月。 20. 格雷戈里，N。 \& O'Reilly，C.L.，NACA 0012 翼型部分的低速空气动力学特性，包括模拟Hoar Frost的上表面粗糙度的影响，航空研究委员会，报告和备忘录3276，英国，1月。 1970年。 21. 古斯塔夫森，B.，克雷斯，H.-O. \& Oliger，J.，时间依赖问题和差异方法，美国威利，1995年。 22. 詹姆森，A.，施密特，W. \& Turkel，E.，使用Runge-Kutta 时间推进的有限体积方法对Euler 方程的数值解，AIAA论文81-1259，1981年6月。 23. Johan，Z.，Hughes，T.J.R. \& Shakib，F.，流体有限元分析中产生的隐式时间行进方案的全球收敛无矩阵算法，应用机械工程中的计算方法87，1991年，pp. 281-304。 24. 考希克，D.K.，凯斯，D.E. \& Smith，B.F.，非结构化电网上可压缩和不可压缩流动的NKS方法，第11届Int. 关于区域分解方法的会议，C.-H。 Lai等人，编辑，1998年，pp. 286-297。 25. 诺尔，D.A. \& Rider，W.J.，多网格预制牛顿-克里洛夫方法，SIAM J.关于科学计算21，1999年，pp. 691-710。 26. 克里斯，H.-O. \& Oliger，J.，双曲方程积分的准确方法的比较，Tellus 24，1972年，pp. 199-215。 27. 罗，H.，鲍姆，J.D. \& Lohner，R.，用于非结构化网格上可压缩流的快速、无矩阵隐式方法，J. 计算机。 物理。 146，11月 1998年，pp. 664-690。 28. 洛马克斯，H。 \& Liu，Y.，Euler 方程的非耗散稳态解决方案，AIAA论文85-1498，1985年7月。 29. Merriam，M.L.，基于熵的非线性稳定性方法，美国宇航局技术备忘录101086，1989年3月。 30. 尼尔森，E.J.，安德森，W.K.，沃尔特斯，R.W. \& Keyes，D.E.，牛顿-克里洛夫方法论在三维非结构化欧拉代码中的应用，AIAA论文95-1733，1995年6月。 31. 普埃约，A。 \& Zingg，D.W.，高效的牛顿-克里洛夫空气动力学求解器

空气动力学流动算法369计算，AIAA J。 36，11月 1998年，pp. 1991-1997年。 32. Pulliam，T.H.，Navier-Stokes 方程的高效解法，讲义，von Karman Inst，流体力学系列讲座1月。 1986年。 33. Rai，M.M. \& Moin，P.，过渡的直接数值模拟和空间演变的湍流 边界层，J. 补偿。 物理。 109，1993年，pp. 169-192。 34. Raj, P., CFD at a Crossroads: An Industry Perspective, in Frontiers of 计算流体动力学 1998, D.A. Caughey和M.M. 哈菲兹，编辑，世界科学，新加坡，1998年，pp. 105-110。 35. 兰瓦拉，A.A. \& Rai，M.M.，Navier-Stokes 方程的多区高阶有限差分法，AIAA论文95-1706，1995年6月。 36. Roache，P.J.，计算科学和工程的验证和验证，新墨西哥州赫莫萨，1998年。 37. Roe，P.L.，近似黎曼求解器、参数向量和差分方案，J. 补偿。 物理。 43，1981年，pp. 357-372。 38. Rogers，S.E.，不可压缩导航-斯托克斯方程的隐式方案的比较，AIAA论文95-0567，1995年6月。 39. 萨德，Y。 \& Schultz，M.H.，GMRES：求解非对称线性系统的广义最小残差算法，SIAM J.关于科学和统计计算7，1986年，pp. 856-869。 40. Shakib，F.，可压缩欧拉和Navier-Stokes 方程的有限元分析，博士。 论文，斯坦福大学，11月。 1988年。 41. Shang，J.S.，时间依赖的Maxwell方程的高阶紧差分方案，J. 补偿。 物理。 153，8月 1999年，pp. 312-333。 42. 斯旺森，R.C \& Turkel，E.，关于中心差和上风方案，J. 补偿。 物理。 101，1992年，pp. 292-306。 43. 斯沃茨，B。 \& Wendroff，B.，有限差分和有限元方法的相对效率，SIAM J.关于数值分析11，1974年，pp. 979-993。 44. Treidler，E.B.，叶卡捷琳娜，J. \& Childs，R.E.，高精度中心差分CFD方案的高效解算法，AIAA论文99-0302，1月。 1999年。 45. Venkatakrishnan，V.和Mavriplis，D.J.，非结构化网格的隐式求解器，J. 补偿。 物理。 105，1993年，pp. 83-91。 46. 维斯巴尔，M.R. \& Gaitonde，D.V.，复杂非稳态流动的高阶精确方法，AIAA J. 37，10月 1999年，pp. 1231-1239。 47. 威尔斯，V.L. \& Renaut，R.A.，计算空气动力学产生的噪音，Annu。 牧师。 流体机械。 29，1997年，pp. 161-199。 48. 威格顿，L.B.，尤，新泽西州 \& Young，D.P.，计算流体动力学代码的GMRES加速，AIAA论文85-1494，1985年7月。 49. Shalman, E., Yakhot, A., Shalman, S., Igra, O. \& Yadlin，Y.，通过衰减人工耗散准确计算Navier-Stokes 湍流边界层，AIAA论文97-0544，1月。 1997年。 50. Zhong，X.，高超音速边界层转换数值模拟的高阶有限差分方案，J. 补偿。 物理。 144，1998年，pp. 662-709。 51. Zingg，D.W.，De Rango，S.，Nemec，M. \& Pulliam，T.H.，Navier-Stokes 方程的几个空间离散化的比较，J. 补偿。 物理。 160，2000年5月，pp. 683-704。 52. Zingg，D.W.，翼型流场的薄层Navier-Stokes计算的网格研究，AIAA J. 10月30日 1992年，pp. 2561-2564。 53. Zingg，D.W. \& Lomax，H.，关于与双曲方程的数值边界方案相关的特征系统，流体力学的数值方法，M.J. 贝恩斯和K.W. 莫顿，编辑，克拉伦登出版社，牛津，英国，1993年，pp. 471-481。

370 ZINGG, DE RANGO \& PUEYO V ' - VV *\&>'- \textbackslash\{\}£>-., \textbackslash\{\}\textbackslash\{\}B'."' X \textasciitilde{}'\textasciicircum{}A \textbackslash\{\}v。 ''•• \textbackslash\{\} X \textbackslash\{\}\textbackslash\{\} Q N \textbackslash\{\}\textbackslash\{\} \textbackslash\{\} \textbackslash\{\} "\textbackslash\{\} \textbackslash\{\} \textbackslash\{\} \textbackslash\{\} \textbackslash\{\} " \textbackslash\{\}\textbackslash\{\} ''•• '*• \textbackslash\{\}\textbackslash\{\} \textbackslash\{\} \textbackslash\{\} \textbackslash\{\} b \textasciicircum{} * b 1 1 l.e-1 l.e-2 l.e-3 l.e-4 l.e-5 y a -+-. -•\&-•- -. *\_.- 功能评估中的200 400 600 800 1000 1200 CPU时间图1 r\textbackslash\{\}n对收敛的影响。 Rjn的值在图例中给出。 每个符号代表一个外部迭代。s 200 • 案例1 -*- 案例2 --•- 案例3 <D"- 案例4 * 0 2 4 6 8 10 12 14 16 18 20 图2 a对收敛的影响。 54. Zingg，D.W.，Lomax，H. \& Jurgens，H.M.，线性波传播的高精度有限差分方案，SI AM J.关于科学计算17，1996年3月，pp. 328-346。 55. Zingg，D.W.，线性波传播的高精度有限差分方法的比较，SIAM J.关于科学计算，正在印刷。

空气动力学流量算法 371 O -10 -12 <? \textbackslash\{\} \textasciicircum{}\textbackslash\{\} k \textbackslash\{\} 9 \textbackslash\{\} k \textbackslash\{\} 'a \textbackslash\{\} k \textbackslash\{\} a \textbackslash\{\} a in 1 1 BFILU(O) from A2 -»- BFILU(O) from Al -+--- BFILU(2) from Al -\&••- " Xs, 500 1000 1500 2000 CPU 函数评估中的时间 图 3 各种预处理器的性能。-10 - -12 - 14 Newton-Krylov AF-MG f**m*tsJS>j*£ 500 1000 1500 2000 2500 3000 CPU 函数评估时间 图 4 案例1的剩余历史。

ZINGG, DE RANGO \& PUEYO -10 -12 -14 Newton-Krylov AF-MG 0 500 1000 1500 2000 2500 CPU 功能评估时间图5 案例2的剩余历史。-10 -12 -14 Newton-Krylov AF-MG BFILU（2）AF 500 1000 1500 2000 CPU功能评估时间图6案例3的剩余历史。

空气动力学流量算法 373 -14 Newton-Krylov AF-MG 0 200 400 600 800 1000 1200 1400 CPU 函数评估时间图7 案例4的剩余历史。-10 -12 -14 Newton-Krylov AF-MG 500 1000 1500 2000 2500 3000 CPU 功能评估时间图8 案例2的剩余历史，减少墙外网格间距。

374 ZINGG, DE RANGO \& PUEYO (a) £ 1.315 1.31 二阶 -e- 高阶 - *--- 三阶上风 •••••«•• 1e-05 2e-05 3e-05 4e-05 5e-05 6e-05 7e-05 8e-05 1/N fficient 压力阻力 Coe ent 0.012 0.0115 0.011 0.0105 0.009 0.0095 0.008 0.0085 0.008 0.0049 0.0049 0.0048 二阶 -e----*- 三阶上风 -o-05 2e-05 3e-05 4e-05 6e-05 7e-05 8e-05 1/N i= 0.0044 0.0043 0.0042 二阶 -e-e-高阶-x -a- 1e-05 2e-05 3e-05 3e-05 4e-05 6e-05 6e-05 8e-05 （a）升力系数，（b）压力-阻力系数，（c）摩擦-阻力系数。

空气动力学流动算法 375 0.05 0 0.1 0.2 0.3 0.4 0.5 0.6 0.7 0.8 0.9 1 1.1 非维速度图10亚音速情况的边界层轮廓，85\%弦的上表面，Baldwin-Lomax模型。

376 ZINGG, DE RANGO \& PUEYO (a) 1.28 1.27 g 1.265 1.255 1.25 1e-05 2e-05 3e-05 4e-05 5e-05 6e-05 7e-05 8e-05 1/N (b) 0 1e-05 2e-05 3e-05 4e-05 6e-05 7e-05 8e-05 1/N ( C ) 0.0053 0.0052 0.0051 0.005 0.005 0.0049 0.0048 0.0047 0.0046 0.0045 0.0044 二阶 -e-高阶 --K-• 1e-05 2e-05 3e-05 4e-05 6e-05 7e-05 8e 1/N 图11 网格收敛亚音速案例研究，Spalart-Allmaras模型。 （a）升力系数，（b）压力-阻力系数，（c）摩擦-阻力系数。

空气动力学流动算法 377 (b) 0.756 0.754 0.752 I 0-75 | 0.748 O § 0.746 0.744 0.742 0.74 0.008 0.00798 0.00796 0.00794 0.00792 0.00792 0.00788 0.00786 0.00784 0.00782 0.00778 0.00778 二阶 -e- 高阶 ---* - 1e-05 2e-05 3e-05 4e-05 6e-05 7e-05 1/N 二阶 -高阶 -x- 0 1e-05 2e-05 3e-05 4e-05 6e-05 7e-05 1/N 1e-05 3e-05 3e-05 5e-05 5e-05 1/N 0 1e-05 3e-05 5e-05 5e-05 1/N 05 7e-05 图12 图12跨声速箱I的网格收敛研究，Spalart-Allmaras模型。 （a）升力系数，（b）压力-阻力系数，（c）摩擦-阻力系数。

378 ZINGG, DE RANGO \& PUEYO (a) (b) 0.0208 (C) 0.005 0 1e-05 2e-05 3e-05 4e-05 5e-05 6e-05 7e-05 1/N 二阶 -e- 高阶 -x - 0 1e-05 2e-05 3e-05 4e-05 5e-05 6e-05 7e-05 1/N 1e-05 2e-05 3e-05 4e-05 5e-05 6e-05 7e-05 1/N 图13跨声速案例II的网格收敛研究，Spalart-Allmaras模型。 （a）升力系数，（b）压力-阻力系数，（c）摩擦-阻力系数。

空气动力学流量算法 379 o -1.3 -1.2 -1.1 -0.9 0.8 ft -0.7 -0.6 -0.5 -0.4 -0.3 二阶网格 A 二阶网格 C 高阶网格 C 0 0.1 0.2 0.3 0.4 0.5 0.6 0.7 x/c 图 14 跨声速情况II的表面压力分布。 O i foil E o it nee CO • - Q 0.06 0.05 0.04 0.03 0.02 0.01 n o 二阶网格 A 二阶网格 C 高阶网格 C 二阶黏性项 网格 C jS 4sO -- \&*\textasciicircum{} \textasciicircum{} -\textasciicircum{}S y „\textasciicircum{}S /. J\% / J7 s.yo Sjg \textasciicircum{} - - jo I I 0.2 0.4 0.6 0.8 非维速度图15跨声速案例II的边界层轮廓，95\%弦的上表面。

\clearpage
\chapter{19 高超音速边界层稳定性和受体的数值模拟}
\section{19.1 引言}
高超音速边界层稳定性和接受性的数值模拟 Zhong Xiaolin, Chong W. Whang和Yanbao Ma1 19.1 导言 高超音速边界层层湍流过渡的预测是高超音速飞行器空气动力学设计和控制的关键部分[2]。 过渡过程是层边界层对强迫扰动的非线性响应的结果，这可能源于许多差异源，包括自由流扰动、表面粗糙度和振动。 在初始干扰较弱的环境中，过渡路径由三个阶段组成：1）接受性，2）线性特征模式生长或瞬态生长，以及3）非线性分解到湍流。 第一阶段是接受性过程[7]，它将环境扰动转化为边界层的初始不稳定波。 第二阶段是边界层不稳定性波的后续线性发展和增长。 超音速边界层的相关不稳定性波包括第一模式和更高模式不稳定性[5]、凹面上的戈特勒不稳定性[10]、前缘的附着线不稳定性以及三维边界层的交叉流不稳定性。 第三阶段是线性不稳定波的分解，以及线性不稳定波达到一定幅度后过渡到湍流[4]。 1 加州大学洛杉矶分校机械与航空航天工程系，加利福尼亚州90095。 由AFOSR赠款F49620-00-1-0101支持的研究，由博士监督。 伦·萨克尔。 Frontiers of 计算流体动力学 - 2002 编辑：David A. Caughey和Mohamed M. 哈菲兹 ©2002 世界科学

\section{19.2 控制方程和数值方法}
382 ZHONG, WHANG \& MA [5, 8, 9]审查了超音速和高超音速边界层的稳定性和过渡。 我们对高超音速边界层稳定性的大部分知识都是通过线性稳定性理论（LST）获得的[5]。 Mack[5]发现，除了超音速和高音速边界层中的第一模式不稳定波外，还有更高的声学不稳定模式。 其中，第二种模式成为马赫数大于约4的高超音速边界层的主要不稳定性。 实验研究已经验证了第二模式的存在和主导地位[ll]。 实用的高超音速车辆有钝的机头，以减少热负荷。 人们普遍认为，钝鼻前的弓形冲击对其后面的边界层的稳定性和过渡有很强的影响[9]。 由于涉及不稳定性和受体过程的瞬态高超声速流动场的复杂性，研究高超音速边界层稳定性和受体的有效方法是完整的Navier-Stokes 方程的直接数值模拟。 在[13]和[15]中，我们提出并验证了一种新的高阶（五阶和六阶）上风有限差分激波拟合方法，用于具有强弓震和刚源项的高超音速流直接数值模拟。 使用高阶减震拟合方案可以获得高精度的平均流量和不稳定的解，这些解在船头减震器后面没有虚假的数值振荡。 该方法随后被验证并应用于钝体上二维和三维高超声速流动的接受性和稳定性的数值研究。 本文介绍了将新方法应用于四个测试案例的结果，涉及高超音速边界层流动在二维和三维钝体上的接受性和稳定性。 四个测试案例是：1）抛物线上2D Mach 15流动的受力，2）椭圆截面钝锥体上3D Mach 15流的受力，3）凹面上的高超音速Gortler涡流，以及4）平板上4.5马赫流动的波传播。 19.2 控制方程和数值方法 控制方程和数值方法在此简要总结。 详细信息可以在之前的2-D流的论文中找到[13，14]。 控制方程是不稳定的3D Navier-Stokes 方程，如下所示：其中上标「*」表示维度变数，U* = \{p*, p*u\textbackslash\{\}, p*U2, p*u\%, e*\}。 假定这种气体在热和热量上是完美的。 黏度和导热系数是用以下方式计算的

\section{19.3 结果和讨论}
超音质边界层稳定性383萨瑟兰定律与恒定普兰特尔数一起，Pr。 在由弓冲击和身体表面界定的计算域中，方程被转换为身体拟合的曲线计算坐标。 弓形冲击的位置和运动是未知的，需要用流量变量来解决。 我们对自由流速度£/£的非维度化，身体表面方程给出的参考长度d*的长度尺度，p\textasciicircum{}的密度，p\textasciicircum{}的压力，TjJ的温度，d* /U\textasciicircum{}的时间，U/d*的涡量，C*的熵，1/d*的熵，1/d*的波数等。 无维流量变量用相同的维度符号表示，但没有上标“*”。 3-D Navier-Stokes 方程的空间离散化的数值方法是流向和壁法向方向的五阶冲击拟合方案，以及3D楔形几何的周期性跨度流方向的傅里叶搭配方法，或锥形几何的方位方向。 空间离散方程使用Williamson[12]的低存储Runge-Kutta方案在时间上推进，最高可达三阶。 计算结果的数值精度由网格加密研究以及与现有实验或理论结果的比较来评估。 代码验证和错误评估的详细结果可以在[13，15]中找到。 19.3 结果和讨论案例1。 抛物线上15马赫流量的接收性 第一种案例是2-D高超音速边界层流量在抛物线前缘上弱自由流声学扰动的接收性的数值研究。 自由流扰动是弱单色平面声波，波锋垂直于身体中心线，形式如下：Qoo(x,t)' = |\textasciicircum{}|e\textasciicircum{}(-\textasciitilde{}*) (19.2)，其中q表示任何流动变量的扰动，\{q'\textasciicircum{}l是波幅常数，k\textasciicircum{}是波数，CQQ是达到冲击前的波速。 对于线性声波，非维流动变量的扰动振幅满足以下关系：\{s'\textasciicircum{}l - [v'\textasciicircum{}l = 0 and IP'CQI = \textbackslash\{\}p'oo\textbackslash\{\}/lf = lu'ool-\textasciicircum{}oo - eM00，其中e是一个小数字，eMoo表示自由流声波的相对振幅。 在本文中，只考虑自由流中的快速声波。 自由流声波的强迫频率由F-106 77I5-定义的无量值频率F表示。

384 ZHONG, WHANG \& MA 图1 稳定流动解决方案：（a）压力轮廓；（b）沿身体表面局部定义的皮肤摩擦和传热系数，c/和st。 测试案例的流动条件是：M\textasciicircum{} = 15，7 - 1.4，Pr - 0.72，Reoo = 6026.6，T£，= 193K，T\textasciicircum{} = 1000X，r* = 0.0125 m和d* = 0.1 m。 图1（a）显示了稳定的压力轮廓。 图中显示，抛物线上的稳定流动沿着身体表面形成有利的压力梯度。 由于鼻子钝的影响，在靠近前缘的区域，边界层的压力略有变化。 随着水流进一步向下流发展，边界层的压力分布接近常数。 Navier-Stokes 方程的当前数值解与近似的边界层方程结果进行比较，以便部分验证并更好地了解数值模拟结果。 图1（b）显示了由局部边界层变量定义的局部剪切应力系数Cf和斯坦顿数st。 图中显示，Navier-Stokes解决方案与皮肤摩擦和传热率的边界层结果非常一致。 在获得稳定解后，研究了抛物线前缘二维高超声速流动的边界层对自由流声学扰动的接受性。 非定常流动解是通过对自由流中的稳流解施加声学干扰获得的。 使用完整的Navier-Stokes 方程计算了随后扰动与冲击的相互作用以及边界层在抛物线上的接受性。 考虑了七个非稳态情况，非维频率F从531到2655不等，e从5 x 10\textasciitilde{}4到10-1不等。 图2（a）显示了在流场达到时间周期状态后，瞬时扰动v'的轮廓，在e = 5 x 10-4的F = 2655的情况下。 瞬时轮廓显示了自由流扰动与弓形冲击的相互作用以及发展

高超音速边界层稳定性385（a）（b）图2（a）：在e - 5 x 10\textasciitilde{}4时F = 2655的速度分量v'瞬时扰动的轮廓；（b）表面边界层中七例F的不稳定波的水平速度扰动振幅曲线。 边界层外区域的波浪模式与边界层内部的波形不同。 边界层外的波模式是自由流波通过冲击并在流场中传播的结果。 另一方面，边界层内部的波形显示出非常不同的结构。 这些结构主要是由强迫波引起的边界波。 如图所示，墙上边界层区域的波浪。 2（a），包含两个不同波形的独立区域。 第一个区域位于x < 0.2的区域，第二个区域位于x > 0.2的区域。 在第一个区域，边界层的振荡幅度只有一个峰值。 另一方面，第二个区域在远离墙壁的振荡中形成两个幅度的峰值。 因此，这些结果表明，第一波区形成的波对应于第一模式波，而第二波区由第二模式波主导。 在模拟中达到时间周期状态后，对非定常流动变量扰动的数值解进行时间傅里叶分析。 通过将线性稳定性特征函数与接受性问题的纳维埃-斯托克斯解的傅里叶振幅进行比较，进一步研究了自由流扰动引起的波模的识别。 为了比较，特征函数由它们各自的第一个峰值从墙中归一化。 图3显示了F = 2655的情况下，在几个i网格站上沿着垂直于抛物线表面的网格线的温度振幅的这种比较。 图中显示，第一波区早期站点的数值解更接近边界内的线性稳定性第一模式

386 ZHONG, WHANG \& MA 图3 三个网格站沿抛物线垂直于抛物线表面的网格线的温度振幅变化：（a）i = 50 LST：第1模式，（b）i = 100 LST：第2模式，（c）i = 140 LST：第3模式（F = 2665在e = 5 x 10\textasciitilde{}4）。靠近墙壁的层。 在后来的站点，数值结果逐渐变为更接近i = 100的线性稳定性第二模式的轮廓，i = 140的线性稳定性第三模式。 该图还表明，由于弓形冲击与流动扰动相互作用的影响，数值解在边界层外的非黏性区域具有非常复杂的流动分布。 边界层以外的变化是由强迫扰动及其与弓形冲击的相互作用引起的。 数值结果和线性稳定性结果主要在墙上的边界层内一致。 因此，结果表明，在第一个区域激发的不稳定波以第一模式为主，然后逐渐过渡到下游的第二和第三模式。 在高超音速边界层接受性研究中，鼻子的钝度可以通过由S-\%JT\textasciitilde{}定义的非维Strouhal数来表征，其中r*是鼻子半径。 固定鼻半径下强迫频率的增加等同于相对鼻钝度的增加。 本文通过考虑七个不同强迫频率F的测试案例，同时保持所有其他流动参数固定，研究了鼻子钝度和强迫频率对高超音速边界层的接收性的影响。 七个案例的非维频率在F-531到2655的范围内，对应于S = 0.4到2范围内的Strouhal数。 图2（b）显示了不同频率的七个测试用例，沿着身体表面附近的网格线水平速度扰动的基本频率的傅里叶振幅分布，其中图中的情况a到g对应于从F = 531到2655的增加频率。 结果显示，第一个模式的所有频率的接收波模式相似

高超音速边界层稳定性387\textasciitilde{}（b）图4自由流平面声学扰动引起的瞬时熵扰动的稳态马赫数轮廓和非稳态轮廓。生长和衰减，其次是第二和第三模式的生长和衰减。 随着频率的降低，最大第一模式振幅变大，它们的位置向下游移动，第一模式不稳定区域扩展到更长的范围。 对于本例，最大第一模式振幅在F = 500至600左右的非维频率下达到峰值。 第二模式不稳定区域出现在具有较高局部雷诺数的第一模式区域的下游。 与第一种模式类似，第二种模式随着频率的降低而增加强度，尽管在当前流动条件下，它们比第一种模式波弱得多。 由于频率的增加对应于Strouhal数和相对鼻半径的增加，因此鼻钝度的增加会导致接受系数的降低。 案例2。 M = 15时钝椭圆锥的接受性。第二种情况是马赫15自由流中2:1钝椭圆锥在零攻角下的接受性。 身体表面由以下方式指定：“*2 \_*2 X* = V- + - - d* 2r；2r*，其中r*和r*是曲率的主要和次要鼻子半径。 流动条件是：M\textasciicircum{} = 15，e = 5 x 10\textasciitilde{}3，T\textasciicircum{}/T\textasciicircum{} - 0.167，参考长度d* = 0.01m，Re\textasciicircum{} = 0.60266 x 105/m，鼻半径r* =.0125m，以及r*: r* = 4: 1。 本文提出的结果是使用两套网格获得的：92 x 31 x 64和92 x 61 x 64网格。 图4（a）显示了马赫数轮廓的三维稳定解。 主轴和小轴的弓减震强度不均匀，造成了（19.3）

ZHONG、WHANG和MA图5沿主轴（k=1）和次轴（k=17）轴的横向速度扰动振幅分布周角压力梯度和交叉流速。 交叉质量流动使边界层变厚，并扩大了小轴的支架距离。 驻点加热率的数值结果与Fay和Riddell [3]获得的边界层理论结果进行了比较。 Navier-Stokes模拟预测停滞点加热系数为0.0168，而Fay和Riddell理论给出0.0163。 计算结果和理论结果之间的一致性很好。 在获得稳流解后，用无量值频率F = 4247.9的自由流扰动波模拟了自由流声学扰动对边界层波的产生。 图4（b）显示了非稳态计算达到周期状态后瞬时熵扰动s'的三维轮廓。 瞬时轮廓显示了边界层中沿着表面的波浪模式的发展。 图5显示了|u'|沿主轴和次轴在墙面附近的平行网格线上的|u'|振幅分布。 两个无花果。 4（b）和5表明，边界层中的速度扰动在前缘区域附近形成了振幅峰值。 峰值是边界层中第一个模式的最大诱导强度。 自由流强迫波进入边界层，并始终在前缘附近生成第一模式，快速衰减转换为第二模式。 图还显示，第一模式的峰值在小轴（k = 17）比主轴（A；= 1）更早达到。 这是因为小轴有一个更厚的边界层，它更快地在|u'|扰动中形成峰值。 通过与线性稳定性分析获得的第一模式和第二模式结构的比较，可以进一步分析前缘附近高超音速受体的图片。 稳定性分析使用局部平行假设在每个局部平均流量段进行。 图6将模拟获得的扰动振幅与

高超音速边界层稳定性389图6 k - 1（左图，线：LST第一模式，圆：N-S）和k = 17（右图，线：LST第二模式，圆：N-S）网格线的沿墙法线速度扰动振幅的分布。主轴上第一模式和小轴上第二模式的特征函数。 数值解与边界层中的LST结果非常吻合。 比较确实表明，边界层中的感应波以主轴的第一模式为主，而解决方案已经在小轴上发展了第二模式结构。 比较还表明，数值解在边界层区域准确地捕获了模式结构。 案例3。 凹入口表面的高超音速戈特勒涡流 第三个案例是高超声速流动在钝凹体表面，具有3D戈特勒涡流的非线性发展。 除了高超音速边界层中的其他高超音速不稳定波外，高超音速入口前的凹压缩面还容易受到戈特勒反旋转不稳定涡流的影响[10]。 Gortler涡流的非线性分解是凹高超音速边界层层湍流过渡过程中的关键步骤。 实验表明，戈特勒涡旋的非线性生长相互作用和其他形式的扰动在过渡过程中起着重要作用。 在非线性生长过程中，戈特勒涡流发展成为蘑菇状涡流，因为逆向旋转的涡流以低流速度将流体从壁上泵离。 这种蘑菇形的戈特勒涡流的特点是存在两个区域：低速度的峰值区域和高速速度的山谷区域。 考虑的比流量是M\textasciicircum{} = 15在凹面的钝楔上的流量。 基于鼻半径的自由流雷诺数是。 Reoo - 15075.3。 流场的总长度约为40个鼻半径。 足够的

390 ZHONG, WHANG \& MA Zone 9 jshroom-shaped Vortices Zone 11 Zone 10 图 7 沿溪流方向的流平均速度的等形轮廓分布。 每个区域的网格大小为161 x 121 x 64。 由于非线性效应，随着水流向下游移动，蘑菇形的涡流就会发展。

高超音速边界层稳定性391图8四个不同溪流位置的剖面溪流速度分布。稳定平均流量解是二维的，由使用相同的激波拟合方案单独模拟计算。 获得稳态解后，从19.4鼻半径的x坐标开始，模拟了Gortler涡流的线性和非线性生长。 戈特勒涡流是由局部计算域入口稳态解施加的小扰动引起的。 这些入口扰动是基于模拟平均流量从线性稳定性分析获得的第一模式Gortler涡流的特征函数。 表面对戈特勒涡流的影响由G = Re\textasciicircum{}（19.4）定义的非维戈特勒数来测量，其中S - A/\textasciicircum{}T-»和R是物体表面曲率的局部半径。 在本测试案例中，扰动模拟入口处的戈特勒数为6.7，其中线性稳定性分析预测，随着水流向下游移动，戈特勒涡流的强度在空间上增长。 为了研究Gortler不稳定性的非线性效应，我们在计算域的入口引入了相对较大的振幅扰动。 入口Gortler模式中速度扰动的最大振幅为0.251700。 图7显示了在入口扰动的情况下，当水流向下游移动时，顺流速度。 图中显示，随着水流向下游移动，由于非线性效应，最初的线性戈特勒涡流发展成为蘑菇状涡流。 弓形冲击对流场没有太大影响，因为戈特勒涡流在远离冲击的黏性层中发展。 图8显示了四个不同溪流位置的溪流速度的跨度等高等高。 高峰

392 ZHONG, WHANG \& MA '°g,(E.) 图9 图中清楚地显示了各种跨度傅里叶模式和山谷区域的扰动能量的非线性演变。 虽然中间区域（峰值）趋于上升，但其他区域则变窄。 戈特勒涡流的非线性增长可以用傅里叶光谱能模式来表示。 图9显示了每个傅里叶模式的扰动能的流动发展。 最初，Gortler涡流（模式1）线性增长，而其他模式可以忽略不计。 随着模式之间的非线性相互作用在下游变得更加显著，平均流量校正模式（模式0）大幅增长，随着流量向下游移动，它主导了模式1。 在这种非线性模式下，主要相互作用是基本流校正模式和平均流校正模式（模式0）。 然而，随着水流向下游移动，更高模式（模式2、3等）的影响也会增加。 在下游，这些更高的模式将进一步增加并饱和。 案例4。 4.5马赫边界层在平板上 第四种情况是使用五阶激波拟合方案来模拟4.5马赫在平板上的线性和非线性稳定性流动。 这项工作的目的是对完美和非平衡真实气体流动的高超音速边界层中不稳定波的非线性击穿进行数值研究。 本文只介绍了完美的气体解决方案。 为了使用激波拟合方案，使用TVD（总变异减小）冲击捕获方案计算了前缘附近的稳定流场。 靠近前缘的TVD结果被用作下流场中激波拟合方案的流入条件。 对于非定常流动模拟，不稳定波是由通过前缘附近墙壁上的狭窄槽的弱吹气和吸力引起的。 测试案例的具体流动条件是M\textasciicircum{} = 4.5流动

高超音能边界层稳定性393 1.0E-3 F 8.0E-4 6.0E-4 4.0E-4 2.0E-4 0.0EO 0 200 400 600 800 1000 1200 1400 1600 1800 2000 ReL 图10 稳定性的中性曲线（F vs. ReL，）用于M = 4.5的二维一模和二模扰动。平板。 计算域从x = 0.006m开始，到0.25m结束，对应于从Rex = 0.432 x 105到Rex = 1.8 x 106的局部雷诺数。 在边界稳定性研究中，经常使用基于边界层厚度的长度尺度的雷诺数，ReL = U\textasciicircum{}L/v\textasciicircum{} = \textbackslash\{\}/Rex，其中L = y\textasciicircum{}ffr--计算域的ReL值从207.9到1341.6。 计算域沿着溪流方向分为7个区域，溪流方向共有1221个网格点，墙法线方向共有121个点。 根据局部线性稳定性理论，选择计算域入口附近的打击和吸力的非维频率F，以便第一或第二模式不稳定。 使用Malik的[6]局部LST方法获得的4.5马赫平板边界层的第一模式和第二模式不稳定区域如图所示。 10. 根据此图，为吹气和吸力扰动选择了两种不同的频率：分别为F-2.2 x 10-4（第二模式）和F = 0.6 x 10\textasciitilde{}4（第一模式）。 出于代码验证的目的，模拟结果与LST预测进行比较。 在本文中，只提出了强迫波频率为F = 2.2 x 10\textasciitilde{}4的非稳态情况的结果。 图10显示，当波向下游发展时，第二模式波首先经过一个稳定的衰变区域，然后是狭窄的不稳定生长区域，然后又经过一个稳定的区域。 LST预测的第二模式不稳定区域从ReL-836（分支I中性稳定点）到1028（分支II中性稳定点）。 根据LST结果，第二模式是该区域最不稳定的模式，由初始吹气和吸力引起的波场预计将由那里的第二模式波主导。 图12（a）显示了瞬时压力扰动的轮廓

394 ZHONO。 WHANG \& MA 2.4E-3 1.7E-3 1.0E-3 3.4E-4 -3.4E-4 1.0E-3 -I.7E-3 -2.4E-3 2.0E5 4.0E5 6.0E5 8.0E5 1.0E6 1.2E6 1.4E6 1.6E6 1.8E6 (a) Re. （b）1.0E6 1.1E6 Re。 图11非稳定解：（a）瞬时压力扰动轮廓；（b）M = 4.5，F = 2.2 x 10\textasciitilde{}4的局部波模式。 1200 1500 Re，图12在流场达到周期状态后，M = 4.5，F = 2.2 x 10-4的壁压波动的不稳定解。 很明显，吹气和吸力引入了边界层干扰以及边界层外的马赫波辐射。 墙上边界层内部的波浪是吹气和吸力区域附近的几种波浪模式的组合。 当波浪向下游发展到达第二模式的分支I中性稳定点时，第二模式的振幅因不稳定而增加。 在这个区域，边界层的扰动以第二模式波为主。 如等高线所示，第二模式波在通过分支II稳定点后衰减。 图12（b）显示了第二个生长区域的局部波浪模式。 图中显示了典型的第二模式波形。 沿着表面的压力扰动分布被绘制在

\section{19.4 结束评论}
高超音质边界层稳定性395图。 12. 压力扰动的包络（破折号线）显示了随着波浪向下游发展而变化的振幅。 图中显示，扰动在经过初始调制区域后进入第二模式不稳定区域。 在模拟结果中，扰动振幅在大约Rei, = \textbackslash\{\}/Rex - 1052的分支II中性点达到最大值。 通过分支II中性点后，扰动的振幅在第二模式稳定区域迅速衰减。 如图所示，通过数值模拟获得的第二模式不稳定区域。 12大约在Rei，范围为840到1052。 这些值与LST预测的范围从836到1028非常吻合。 从非稳定解的数值结果中，可以计算出扰动的流波数和增长率。 数值模拟的波数和增长率与基于相同平均流量的局部线性稳定性分析获得的特征值进行比较。 图13分别显示了数值模拟和LST结果之间的流波数和增长率的比较。 在图中，a的实数部分和虚数部分分别是波数和增长率。 这些数字表明，在ar和on的数字和LST结果之间存在良好的一致性，在位于840 < Rei < 1052区域的第二模式不稳定区域。 由于吹气和吸力区域的强迫波的影响，模拟结果与入口附近区域的LST结果不一致。 根据Balakumar和Malik的研究[1]，环境强迫将产生几种离散模式以及连续光谱。 不稳定波可能是离散模式和连续光谱以及强迫波的组合。 在第二模式不稳定区域，第二模式波占主导地位。 因此，在这个区域，LST和数值模拟的结果之间存在良好的一致性是合理的。 我们还可以比较模拟结果和LST之间的特征函数。 图14显示了Rei站模拟和LST计算的速度扰动结构，-1007.57。 尽管外部边界层存在差异，但协议非常好。 在这个位置，雷诺数很大，因此，非平行效应很小。 总体而言，数值模拟可以很好地捕捉扰动波的结构。 19.4 结论 使用五阶激波拟合方案研究高超音速边界层的稳定性和受体，以准确模拟激波与流场中物理波之间的物理相互作用。 本文讨论了四个测试案例。 结果表明

\section{参考资料}
396 ZHONG, WHANG \& MA 400 500 0.050 \textasciicircum{} 0.040 0.030 0.020 0.010 0.000 -0.010 •0.020 -0.030 -0.040 -0.050 (a) Re. （b）关于图13在M = 4.5，F - 2.2 x 1CT4时通过模拟和LST计算获得的波数的流实数和虚部分的比较。五阶激波拟合方案适用于对超音速和高超音速边界层的接收性、稳定性和过渡进行详细的数值研究。 对于15马赫流在抛物线和钝椭圆锥上的接受性的情况下，自由流声学干扰在墙上的边界层中产生第一、第二和第三模式波。 结果表明，当相对鼻半径减少时，受体系数会增加。 对于凹陷入口表面的高音速戈特勒涡流的情况，戈特勒涡流中高速和低速区域的非线性发展会产生蘑菇状涡流。 数值方法也在平板上4.5马赫边界层中进行了测试。 我们的数值模拟结果和LST结果之间有很好的一致性。 参考资料。 P。 Balakumar和M.R. 马利克。 超音速边界层中的离散模式和连续光谱。 流体力学杂志，239：631-656，1992年。。 国防科学委员会。 第二国防科学委员会国家航空航天飞机（NASP）工作组的最终报告。 AD-A274530，94-00052，1992年11月。。 J。 A. Fay和F。 R。 里德尔。 驻点解离空气中的传热理论。 航空科学杂志，25：73-85，1958年。 Y。 S。 卡沙诺夫。 层边界层过渡的物理机制。 流体力学年度回顾，26:411-82，1994年。。 L。 M。 麦克。 边界层 线性稳定性理论。 在AGARD报告中，编号 709，第3-1至3-81页，1984年。 M。 R。 马利克。 高超音速边界层稳定性的数值方法。 公司杂志。 物理，86：376-413，1990年。

超音域边界层稳定性397 DNS LST DNS LST图14在M = 4.5、ReL=1007.57和F=2.2 x 10"4处的特征函数结构比较。 7. M。 莫尔科维恩。 关于过渡的许多方面。 在C.S. Wells，编辑，黏性减减阻力，第1-31页。 全体会议，1969年。 8. M。 V。 莫尔科维恩。 以高超音速过渡。 ICASE临时报告1，美国宇航局CR 178315，1987年5月。 9. E。 雷肖特科。 高超音速稳定性和过渡，在高超音速流的再入问题中，编辑。 J.-A. 德赛德里，R。 Glowinski和J。 Periaux，Springer-Verlaq，1:18-34，1991年。 10. W。 S。 讽刺。 戈特勒·沃提斯。 流体力学年度评论，26:379-409，1994年。 11. K。 F。 Stetson和R。 L。 金梅尔。 关于高超音速边界层稳定性AIAA论文92-0737，1992年。 12. J。 H。 威廉森。 低存储Runge-Kutta方案。 计算物理学杂志，35:48-56，1980年。 13. X。 钟。 高超音速边界层过渡在钝前缘的直接数值模拟，第一部分：新的数值方法和验证（邀请）。 AIAA论文97-0755，第35届AIAA航空航天科学会议和展览，1月6日至9日，内华达州里诺，1997年。 14. X。 钟。 高超音速边界层过渡的直接数值模拟在钝前缘，第二部分：对声音的接受性（邀请）AIAA论文97-0756，1997年1月。 15. X。 钟。 高超音速边界层过渡数值模拟的高阶有限差分方案。 计算物理学杂志144：662-709，1998年8月。

\clearpage
\chapter{20 使用重叠网格方法在涡轮泵中对不可压缩流动进行时间依赖模拟}
\section{20.1 介绍}
使用重叠网格方法Cetin Kiris1和Dochan Kwak1在涡轮泵中对不可压缩流动进行时间依赖模拟 摘要本文报告了使用重叠网格系统实现完整非稳定涡轮泵模拟能力的进展。 涡轮泵叶轮的计算模型被用作MPI、混合MPI/Open-MP和INS3D代码的MLP版本的性能评估的测试案例。 通过采用重叠网格技术获得转子-定子相互作用的网格系统的相对运动。 在美国宇航局艾姆斯研究中心的Origin 2000系统上对涡轮泵进行不稳定的计算，涡轮泵包含114个带有3430万个网格点的区域。 介绍了这些模拟所采取的方法，以及代码的并行版本的性能。 20.1 引言 这一努力的动机基于两个主要要素。 首先，整个涡轮泵模拟旨在为液体火箭发动机燃料供应系统的设计和分析提供一个计算框架。 这项工作是高性能计算和通信（HPCC）高级1 M.S.的一部分。 T27B，NAS应用部门，美国宇航局高级超级计算（NAS）部门。 美国宇航局-艾姆斯研究中心，莫菲特菲尔德，加利福尼亚州94035 计算流体动力学的前沿 - 2002 编辑：David A. Caughey和Mohamed M. 哈菲兹 ©2002 世界科学

\section{20.2 数值方法}
400个KIRIS和KWAK技术应用项目。 这项研究的第二个目标是支持用于太空运输的液体火箭系统的设计。 由于在不久的将来，太空发射系统可能会依赖液体火箭发动机，因此提高发动机部件的效率和可靠性是一项重要任务。 液体火箭发动机的主要问题之一是了解燃料和氧化剂从油箱流向羽流的流体动力学。 通过数值模拟了解涡轮泵的流量，将对于找到更简单、更高效、更稳健、制造成本更低的更好设计具有重要价值。 直到最近，泵的设计过程与几十年前没有太大区别。 目前的半经验涡轮机械设计过程没有考虑到泵流量中的三维（3-D）黏性现象。 其中一些三维黏性现象包括尾迹；枢纽、裹尸布和叶片中的边界层；连接流；以及尖端间隙流。 尽管文献中广泛报道了计算流体动力学（CFD）在涡轮机中的应用，但在全泵模拟中的应用却相当有限。 本文的目的是介绍和评估一个通过整个涡轮泵配置模拟不可压缩流动的并行计算程序。 这些3D 非定常流动模拟需要大幅减少计算时间，以减少泵的设计周期时间。 这一速度的一部分是由于计算机硬件平台的增强。 算法的进步和高效的并行实现将为加速的剩余部分做出贡献。 以下部分概述了为实现这一加速而采取的初步努力和步骤。 20.2 数值方法 目前的计算是利用INS3D计算机代码进行的，该代码解决了稳态和非稳态流动的不可压缩Navier-Stokes 方程。 不可压缩Navier-Stokes 方程的数值解需要特别注意，以满足速度场上的无发散约束，因为不可压缩公式没有从状态方程或连续性方程中明确产生压力场。 避免由问题的椭圆性质引起的数值困难的一个方法是使用Chorin'开发的人工可压缩性方法。 人工可压缩性算法将压力项的时间导数引入连续性方程；椭圆-方差分型偏微分方程被转换为双曲-方差型。 基于该算法开发了一个2"3的流量求解器系列。 由于结果方程的对流项是双曲的，迎风差分可以应用于这些项。 当前版本被指定为INS3D代码，使用Roe的通量差异拆分4。 这里使用的三阶和五阶迎风差分是一类高精度通量差分方案的实现

不可压缩流动在TURBOPUMP 401 可压缩流动方程中。 为了获得时间准确解，每个物理时间步骤的方程以伪时间迭代到收敛，直到速度场的发散减少到指定的公差值以下。 所需的子迭代总数因问题、时间步骤大小和使用的人工可压缩性参数而异，通常从10到30个子迭代不等。 矩阵方程通过使用非因数高斯-塞德尔型线松弛方案5来迭代求解，该方案保持稳定性并允许采取较大的伪时间步进。 数值方法的详细信息可以在参考中找到。 2-3。 GMRES方案也被用于求解生成的矩阵方程6。 带有线松弛的流动求解器 INS3D的计算机内存要求是单词网格点数量的35倍，GMRES-ILU(O)方案是单词网格点数量的220倍。 当快速收敛方案，如GMRES-ILU（O）求解器，被实施到人工可压缩性方法中时，以前的计算表明，计算结果和实验数据之间得到了合理的一致性。 人工可压缩性方法中的线松弛方案对于时间精确计算来说可能非常昂贵，如果每个物理时间步骤中没有强制执行不可压缩性，可能会导致错误的解决方案。 TUmOPUMP-INDUCER质量流量：9093 GPM Re：7.99©+7几何表面压力图1。 泵感应器的几何形状和表面压力。

\section{20.3 方法和计算模型}
402 KIMS \& KWAK 20.3 方法和计算模型 液氧泵的几何形状有各种旋转和静止组件，如电感器、定子、踢球器和液压涡轮机，其中流量非常不稳定。 图1显示了稳态分量分析中感应器的几何形状和计算表面压力。 当包括旋转和固定部件时，由于不稳定的交互，需要进行与时间相关的模拟。 为了处理几何复杂性，使用了重叠网格方法。 流模拟的超量结构化网格方法已被用于解决航空航天、海洋、生物医学和气象应用中的各种问题。 流量制度可以从商用飞机的简单稳定流量，到像涡轮泵配置那样具有相对运动的物体的非稳态三维流量。 如图2所示，一个几何复杂的物体被分解成许多简单的网格组件。 允许相邻的.grid任意重叠的自由意味着这些网格可以相互独立创建，每个网格通常都是高质量的，几乎正交的。 相邻网格之间的连通性是通过网格外边界的插值建立的。 通过在不干扰现有电网的情况下建立新的连接，在系统中添加新组件并模拟多个物体之间的任意相对运动来实现。 已经分解的网格系统自然实现了并行计算平台上的可扩展性。 对于某些问题，将网格收集到大小大致相等的组中进行并行处理会更有效率。 ShroiMl图2。 用于叶轮长叶片部分的重叠网格系统，带尖端间隙

不可压缩流动在TURBOPUMP 403 FurVir-o图3中。 泵模型和模拟过程中采取的步骤。 为了计算具有移动边界的网格上的流量，OVERFLOW-D7代码中的重叠网格方案与INS3D求解器相结合，在每个时间步骤中都可以获得新的连接数据。 重叠网格方案允许子域相互相对移动，并在边界移动产生大位移时提供一般灵活性。 图3显示了增压泵的模型和模拟程序中采取的步骤。 图3中的数字表示每个部分的刀片数量。 计算网格是使用OVERGRID8软件生成的。 OVERGWD是一个统一的图形界面，用于执行超集网格生成。 该软件包含一般的网格操作功能，以及专门针对高效创建重叠网格的模块。 一般电网实用程序包括电网转换、再分配、平滑、串联、提取、外推、投影等功能。 对超集网格特别有用的功能包括特征曲线提取、双曲和代数曲面网格生成、双曲体积网格生成和笛卡尔框网格生成。 使用OpenGL实现可视化，而小部件则使用Tcl/Tk构建。 该软件可以在从UNIX工作站到个人电脑的各种平台之间进行移植。 为了展示当前的不稳态解能力，选择了SSME穿梭升级泵配置。 图4显示了在NASA-MSFC设施中测试的该泵的测试平台的几何形状。 在这种特定的配置中，SSME叶轮是解除的。

404 KIRIS和KWAK图4。 SSME-rigl穿梭升级泵叶轮的几何形状 SSME-rigl配置的入口导向叶片、叶轮和扩散器部分的计算网格分别如图596和79所示。 图5。 SSME-rigl入口导向叶片的计算网格，有17个区域和550万个网格点。

不可压缩流动在涡轮泵405图6中。 SSME叶轮的计算网格有60个区域和1920万个网格点。 图7。 SSME-rigl扩散器的计算网格，有24个区域和650万个网格点。

\section{20.4 计算结果}
406 KIMS \& KWAK 压力 111 \textasciitilde{}«.194 图 8。 SSME-HPFT叶轮的计算表面压力。 20.4 计算结果 计算结果是280万和1920万网格点SSME叶轮模型的。 图8显示了罩状SSME叶轮的计算表面压力。 本节报告了两种方法在获得INS3D代码的多级并行性方面的性能。 第一种方法是混合MPI/OpenMP，第二种方法是美国宇航局-艾姆斯研究中心开发的多级并行（MLP）。 第一种方法是使用按摩传递接口（MPI）进行区域间并行，并使用OpenMP指令进行区域内并行。 INS3D-MPI9基于跨MPI组的显式按摩传递接口，专为粗晶粒并行而设计。 主要策略是将区域分布在一组处理器上。 在迭代期间，所有处理器将在处理器之间交换边界条件数据，处理器的区域与其他处理器上的区域共享接口。 选择了一个简单的主-工人架构，因为它的实现相对简单，并且它是并行CFD应用程序的常见架构。 所有的I/O都是

不可压缩流动 IN A TURBOPUMP 407由主MPI流程执行，数据已分发给工人。 初始化阶段完成后，程序开始其主迭代循环。 具有24个区域和280万个总网格点的SSME叶轮模型被用作粗粒INS3D-MPI代码的测试案例。 对于此版本的代码，计算模型中的区域数量限制了CPU计数的最大数量。 图9显示了SGI Origin 2000平台上此计算的每秒浮点计数。 与线性速度相比，24个处理器的平均速度约为65\%。 应该注意的是，本文中报告的MPI组的数量总是包括主MPI流程。 对于前四个处理器，性能非常好。 当CPU计数进一步增加时，由于负载均衡问题，代码的性能会降低。 为了获得细粒并行性，使用了OpenMP指令°。 图10显示了每时间积分步所需的时间（以秒为单位）与混合平行代码的处理器数量。 应该注意的是，本文中报告的“每次迭代的时间”包括获得Ax=b线性方程组的时间以及为整个网格系统求解这个特定方程组的时间。 换句话说，迭代项用于物理时间步骤，而不是线性方程求解器的迭代。 还应该注意的是，隐式线松弛扫描的数量在每个时间步骤中都保持不变。 图10绘制了4、12和24 MPI组的病例。 对于每个MPI组，使用了不同数量的线程，如1、2、4、8和16。 CPU计数的数量等于线程数量乘以MPI组的数量。 当线程数量增加时，代码的性能会变慢，因为对于线程数量较多来说，每个区域的网格大小相对较小。 这如图12所示。 INS3D-MP1 1500-1:::: 3-1 1000o El S 500- 0H 1 i 1 i 0 5 10 15 20 25 CPU 数 图 9。 SGI 02K，SSME-HPFT叶轮上的INS3D-MPI性能。

408 KIRIS \& KWAK 当问题大小从2增加时。 800万网格点到1920万网格点，SSME叶轮有60个区域，其中最小的区域有74K网格点，最大的区域有996K网格点。 图12显示了使用一个OpenMP线程时，MPI组对代码性能的影响。 一个好的负载均衡可以获得多达20个MPI组。 当使用超过20个MPI组时，代码的性能不再得到改善。 对于相同的MPI组，OpenMP线程的数量增加到2、4、6、10和15。 这些案例在图13和图15中绘制。 图14显示了每次迭代的时间与总CPU计数，图14显示了各种OpenMP线程的情况。 20个MPI组获得了最佳表现。 在图13中，可以看到负载均衡对30个MPI组的影响。 OpenMP指令显示了20和30 MPI组的加速非常相似（见图14）。c 100 o 2 so JS 50 I 40 9 '5" 1 20 t- 10 TT II I I |!! 我我•.. 280万点-a-4 MPI组-m 12 MPI组-A-24 MPI组i i 111 J i i i 111 10 20 30 40 100 CPU数量图10。 各种MPI组（INS3D MPI/OpenMP）的每个迭代步骤的时间。

不可压缩流动 IN A TURBOPUMP 409 100 -C o 2 60 g 50 Z 40 CL。 30 "5" * E bs::CL i I I!!! I 2.8M点-D- 4个MPI组-•- 12个MPI组-±- 24个MPI组每个MPI组OpenMP线程图11。 每个迭代步骤的时间与每个MPI组的OpenMP线程。 S0Q0 i 1 1 1-I-I-i i i li 1 1 1 5000 | 2000 n - 1000 8。 77 600 S 500 HL 400 E \_• -X rv -D- 1 OpenMP Thread - - n \textbackslash\{\},-Q D-Q i i n i 2 3 4 6 10 数字ol CPUs 20 30 图12。 当使用一个OpenMP线程（SSME叶轮1920万网格）时，各种MPI组的性能。

410 KIRIS \& KWAK eoo C 400 •\{= 300 •S 200 a a. "5" 100 o a, 60 E 50 i- 40 A"!' “I D\textasciicircum{} II \textbackslash\{\} T I I I III I i i I I 19.2M 点 -a- 6 MPI 组 -•- 12 MPI 组 -Ac- 20 MPI 组 -A- 30 MPI 组 \_ • A • •-l\textasciicircum{}lM - \textasciicircum{}\textasciicircum{}A.., VIPi i i i i 6 10 20 30 40 100 200 300 CPU数量图13。 INS3D-MPI/OpenMP性能与Origin 2000的CPU计数。o ••p CO fc » Q. * \& a X！ •！！！ 我我！ II M-i-\textasciicircum{} - - -X A "m N 19.2M 点 -a- 6 MPI 组 \textasciitilde{} -«- 12 MPI 组 -AT- 20 MPI 组 -A- 30 MPI 组 - •\textasciicircum{}L. m-\textasciicircum{}-n - ' >\textasciicircum{}M - i i i i i i 20 1 2 3 4 6 10 每个MPI组的OpenMP线程图14。 INS3D-MPI/OpenMP性能与各种MPI组的OpenMP线程。

不可压缩流动在涡轮泵flfiQU I中。 80 *\textasciitilde{} 20 -CP n 1 Mill！ 1S.200万点/ Open-MP -a- Gaus-Seidel Lin© Rel. -A- GMRES Solver •\textasciicircum{}s \textasciicircum{}v I ' i \_L I.1-l-LJ-L- 1 2 3 4 8 10 20 30 Number oi CPUs 图15。 两种不同线性求解器的OpenMP性能。 MLP 流程 1 常见 /Seeai/ aa,fct> MLP 流程 2 常见 /local/ aa,M> 区域 1,4:i:tiCj iScmoiy MLP Osfe slmtis * Lm\&Steres *i --'- 区域 2,3,5 C.-u I I! CW/>。 Sftarec Memory Amne frioMs 8C ttsia. etc I I Common /global/ x,y,z 图16。 NAS-MLP n的共享内存MLP组织。

412 KIRJS和KWAK 20 30 40 100 200 300 CPU数量图17。 INS3D-MLP性能与Origin 2000的CPU计数。 图15显示了OpenMP线程对隐式高斯-塞德尔线松弛方案和GMRES-ILU(O)方案的影响。 线松方案进行了四次高斯-塞德尔扫描，对GMRES求解器进行了20次定向搜索。 由于斜平面中的嵌套循环效应，OpenMP指令显示GMRES求解器的速度更好。 多级并行（MLP）的第二种方法是使用Taft为OVERFLOW代码开发的NAS-MLP11例程获得的。 美国宇航局艾姆斯研究中心开发的共享内存MLP技术已被纳入INS3D代码中。 这种方法与MPI/OpenMP方法在根本上有所不同，因为它根本不使用消息传递。 所有最粗略和最精细的数据通信都是通过直接内存引用指令完成的。 这种方法还为转换传统代码提供了更简单的机制，例如INS3D，然后是MPI。 对于共享内存MLP，通过标准UNIX分叉生成独立进程来提供最粗的级别的并行性。 UNIX分叉超过MPI过程的优势在于，用户不必更改大型生产代码的初始化部分。 共享内存MLP以与Taft将MLP插入OVERFLOW代码的方式非常相似的方式插入到INS3D中。 例程库用于启动分叉，以

\section{20.5 总结}
不可压缩流动在TURBOPUMP 413中建立共享内存领域，并提供同步原语。 INS3D的共享内存组织如图16所示。 重叠网格系统的边界数据由每个进程存档在共享内存领域中。 其他进程根据需要从竞技场访问数据。 图17显示了1920万网格点SSME叶轮模型的INS3D-MLP性能与CPU计数的对比。 GMRES-ILU(O)线性求解器用于这些计算。 MLP版本的代码显示了73\%的线性加速性能。 当MLP性能与MPI/OpenMP性能（图13，20个MPI组）进行比较时，使用MLP版本的代码可以观察到19\%的速度。 这种比较可以在图18中看到。 应该注意的是，此比较是初步的，因为目前正在进一步改进MLP代码的精细并行化。c o CO o o SL E - ••• •-1- i i i i i i i i 19.2M Points -•- MPI-OpenMP Hybrid - -A- NAS-MLP \textbackslash\{\}!!! 11 i i; \textbackslash\{\} U, 1::! -I j- \_\_: V°w r - MM iK\textasciicircum{}I 10 20 30 40 100 CPU数量图18。 INS3D-MPI/OpenMP和INS3D-MLP性能的比较。 20.5 摘要 稳定和时间准确公式的不可压缩流动求解器已被用于并行涡轮泵计算。 概述了非稳涡轮泵模拟的网格系统和数值程序。 为INS3D-MPI/OpenMP和INS3D-MLP版本的性能评估提供了280万和1920万个网格点SSME叶轮模型的结果

\section{20.6 致谢}
\section{参考资料}
414 KIRIS \& KWAK代码。 具有60个区域的SSME叶轮模型显示，多达20个MPI组混合代码显示出良好的可扩展性。 OpenMP指令对GMRES（ILU）求解器比线路放松方案更有效。 共享内存MLP版本的代码是使用NAS-MLP例程开发的。 SSME叶轮计算显示MLP版本的可扩展性非常好。 20.6 致谢 作者感谢Tom Faulkner和Jennifer Dacles Mariani提供代码的粗粒INS3D-MPI版本，感谢Henry Jin在OpenMP指令中的支持和有价值的想法和讨论，以及James R. Taft用于提供NAS-MLP例程。 参考资料1。 Chorin，A.，J.，“解决不可压缩黏性流动问题的数值方法”计算物理学杂志，卷。 2，pp. 1967年12月26日。 2. 郭，D.，张，J. L C，Shanks，S。 P.和Chakravarthy，S.，“使用原始变量的三维不可压缩的Navier-Stokes 流动求解器”，AIAA杂志，卷。 24，不。 3，pp. 390-396，1977年。 3. 罗杰斯，S。 E.，Kwak，D.和Kiris，C，“稳定和时间依赖问题的不可压缩Navier-Stokes 方程的数值解”，AIAA杂志，卷。 29，不。 4，pp. 603-610，1991年。 4. Roe，P.L.，“近似黎曼求解器、参数向量和差分方案”，J. ofComp。 支付。，卷。 43，pp. 357-372 1981。 5. MacCormack，R.，W.，“Navier-Stokes方程的数值解的现状”，AIAA论文编号 85-0032，1985年。 6. 罗杰斯，S。 E.，“不可压缩Navier-Stokes 方程和人工可压缩性的隐式方案的比较”，AIAA杂志，卷。 33，不。 10月10日 1995年。 7. Meakin，Robert L.，“复合过度结构化网格”，网格生成手册，CRC出版社，编辑。 汤普森，索尼，韦瑟里尔，1998年。 8. 陈，W。 M.，OVERGRID - 统一的Overset 网格生成图形界面，将出现在2000年网格生成杂志上。 9. Faulkner，T.和Mariani，J.，“不可压缩的Navier-Stokes求解器（INS3D）的MPI并行化”。http://www.nas.nasa.gov/\textasciitilde{}faulkner/home.html 10. H。 金，M。 Frumkin和J。 Yan，“OpenMP指令的自动生成及其在计算流体动力学代码中的应用”，第三届高性能计算国际研讨会论文集，日本东京，10月。 16-18，2000年。 11. 塔夫脱，J。 R.，OVERFLOW-MLP和LAURA-MLP CFD代码以及NASA Ames 512 CPU Origin系统的性能，”HPCC/CAS 2000研讨会，NASA Ames研究中心，2000年。

\clearpage
\chapter{模拟三角翼上涡流的21个方面}
\section{21.1 导言}
三角翼上涡流模拟的各个方面 Arthur Rizzi、Stefan Gortz和Yann LeMoigne1 21.1 简介 当前和未来的战斗机被设计为以非常高的攻角短暂飞行。 例如，像“眼镜蛇机动”这样的斗狗机动现在需要这种超机动性。 这些新型机动涉及高俯仰率和在超出静态失速攻角的入射处飞行。 机动战斗机周围的空气动力学非常复杂，但需要了解，以便优化飞机设计。 由于使用了三角翼和增强的动态失速，战斗机在高攻角机动中通常具有卓越的空气动力学性能。 这些机翼在入射时的空气动力学特点是在机翼的背向一侧形成一对强大的前缘涡流，与传统的矩形相比，产生了额外的升力。 然而，在非常高的发生率下，涡流破裂，出现带有湍流尾迹的再循环区，升力下降。 在机动过程中，会发生同样的现象，但流动必须适应移动的平面，因此在动态响应中观察到时间滞后。 例如，对形成空气动力历史的迟滞回路的预测和理解，对于设计者在机动过程中扩展飞行包络的极限非常有用。 在飞机机头的高攻角下形成的前体涡流也是1航空系，KTH皇家理工学院，瑞典斯德哥尔摩10044。 Frontiers of 计算流体动力学 - 2002 编辑：David A. Caughey和Mohamed M. 哈菲兹 ©2002 世界科学

416 RIZZI、GORTZ和LEMOIGNE在机动特征中起着主导作用。 为了增强飞机在机动中的行为，需要预测和可能控制涡流的不对称发展。 尽管在过去的五十年里，为了阐明高攻角系统中遇到的复杂流体物理学，对三角洲机翼上的涡流进行了实验和计算研究，但稳流的某些细节仍未得到解决，瞬态飞行的许多基本问题仍未得到解答，这共同证明了继续研究兴趣和进一步调查。 例如，SAAB Gripen最近进行了重要的高阿尔法飞行测试，并更好地了解了作用在飞机上的力和力矩的滞后效应背后的不稳定空气动力学。 这可以改善对飞行测试结果的理解，并更准确地调整飞行模拟器中使用的空气动力学数据库。 21.1.1 静止翼 Gortz [10]对70°锐边三角翼的数值研究表明，从定性角度可以准确地模拟涡流的基本特征以及涡流分解现象。 发现欧拉和Navier-Stokes 方程都能够预测涡流分解。 突破位置被认为与物理建模无关，但它似乎取决于使用的不同数值方案--Roe的上风方案预测的上游爆炸位置比中央方案更多。 在Navier-Stokes计算中可以看到黏性效应，如二次分离，但在欧拉计算中没有。 涡旋轨迹在跨度方向上与实验数据显示出良好的一致性，但涡旋轴位于机翼上方太高。 此外，机翼表面的吸力峰值被认为太低。 在Gortz等人的后续研究中。[ll]根据四个具体标准判断了建模的定量准确性：1）综合空气动力学系数与迎角的变化，2）涡旋在机翼上的轴向位置，3）三角翼上的流动拓扑结构，包括表面压力，以及4）沿涡核轴的速度轴的轴向分量。 空气动力学系数可以证明与实验数据的良好一致性。 尽管由于实验数据中的散射，与风洞测量的比较被证明很困难，但计算的故障位置达到了合理的一致性。 然而，流动的细节，如涡核中拖曳流体速度的轴向分量或表面压力分布，对网格拓扑结构非常敏感，包括奇异线和网格细度

DELTA-WING VORTEX在顶点和前缘流动417，以及整个网格单元的大小。 这一结果提出了一个问题，即“最佳网格”应该是什么样子，以及是否可以实现网格收敛。 由于缺乏有意义的实验数据，当时无法确定代表实验流动条件的速度曲线，因此是另一个未决的问题。 给出的例子表明，在高阿尔法下对delta翼的模拟并不像人们所相信的那样直接，因为许多已发表的低到中度攻角的结果，其中计算的流场通常与实验数据非常一致，网格灵敏度很低。 本文的第一部分旨在评估静止-三角洲-翼计算的现状，并尽可能彻底地分析影响预测流场的参数。 网格生成过程正在受到特别关注，这似乎对结果的质量和准确性有重大影响。 此外，我们将尝试回答关于速度轴向分量沿涡核轴的分布情况的问题。 因此，在具有不同网格拓扑结构和局部网格加密的网格上计算的轴向速度曲线将与实验数据相关。 21.1.2 音高振荡三角洲 本研究第二部分的目标是围绕俯仰锋利的70° 三角翼（见图1）模拟非定常流动。 这是战斗机高阿尔法空气动力学研究的第一个也是主要组成部分。 第二个部件可能是接近战斗机前体的细长体。 这种相对简单形状的模拟，无论是单独还是组合在一起，都有助于更好地了解空气动力学，并预测作用在机动战斗机上的力，目的是改进工程设计。 在这项研究中，三角翼在正弦波振荡俯仰运动中被动画化，这是飞行机动的一个相当简单但基本组成部分。 自八十年代末以来，三角翼和细长体都对这种简单运动进行了实验研究。 Brandon [2]记录了涡流爆炸位置的演变和70°三角翼在各种平均攻击角度的俯仰振荡的法向力历史。 Soltani等人[20]的结果表明，雷诺数、侧滑角和降低振荡频率对从0到55°俯仰的三角翼的有力的影响。 Thompson等人在[21]中报告了雷诺数和降低频率对30°-40°振荡的涡流击穿位置的影响，而在[22]中，Thompson等人研究了倾斜三角翼的表面压力分布。 足够的

LeMay等人[17]研究了418 RIZZI、GORTZ和LEMOIGNE涡流的击穿位置和传播速度，用于29°-39°和0°-45°的振荡。 所有这些实验都是在亚音速风洞中对70° 三角翼进行的。 一般来说，这种几何学的数据不多，报告的数据也不完整：通常报告了力的历史，但通常缺少非稳压力和涡流位置的时间历史。 相反，介绍后者的论文不包括部队历史。 俯仰对象的数值模拟，虽然是CFD模拟的最新话题，但也有很好的记录。 关于三角翼，弗里茨[8]围绕9°和15°的平均攻角模拟了一个65°的俯仰振荡的裁剪翼，并报告了Q历史和压力分布。 在[1]中，Arthur等人比较了几个CFD代码的结果（力和压力数据），用于围绕9°和21°的65°裁剪三角翼进行计算。 Kandil \& Chuang [14]分析了平均攻角约为20.5°的三角翼的压力分布。 Ekaterinaris和Schiff[7]还在22.4°左右的音高振荡中模拟了双三角翼。 这些Navier-Stokes计算的结果以Cp图和粒子追踪图片的形式给出。 对于跨声速条件（Mach=0.85），Kandil和Kandil [15]在65°裁剪的三角翼上进行了Navier-Stokes模拟，这导致了压力等高线图和Cj历史。 所有这些论文都为振荡振幅小且限制在±8°的情况下提出了重要和开创性的结果（见[7]）。 结果与实验数据相比相当好，这令人鼓舞地继续通过数值模拟进一步研究大振幅情况。 这里的计算研究涉及一个70°的单斜角三角翼，在低亚音速马赫数下进行正弦波振荡。 机翼的运动定律由攻角的历史给出：a(t) = am + a® sin(wt)。 考虑将低振荡振幅与上述模拟进行比较，但也研究了较大的振幅（在0°-40°左右）。 事实上，由于非常明显的滞后回路，代码模拟俯仰运动的能力在大振幅下得到了更好的评估，它也具有更实用性。 使用围绕俯仰轴与机翼旋转的运动网格模拟俯仰运动和随之而来的非稳定流场。 本文首先对数值方法进行了简要描述，然后是测试用例和生成的网格，然后是对结果和结论的讨论。

\section{21.2 计算方法}
DELTA-WING VORTEX FLOWS 419 21.2 计算方法 所有测试用例都是使用Navier-Stokes Multi-Block (NSMB) 流动求解器 [25, 9, 26, 5]解决的，这是法国、瑞典和瑞士大学和工业界联合欧洲研究项目的产物。 该代码使用单元中心有限体积法在结构化多块网格上解决可压缩的Navier-Stokes 方程，以离散化方程。 在NSMB中实现了几个湍流模型，包括Baldwin-Lomax和Granville的代数模型、Spalart和Allmaras的单方程模型、Chien和Hoffman的双方程re-e模型以及双方程re-r模型。 使用运动网格（所谓的ALE，任意拉格朗日-欧拉方法）解决了俯仰-delta问题。 整个网格与倾斜的物体一起作为刚体旋转，非稳的边界条件用于远场和表面边界。 因此，在控制方程中必须考虑网格速度：dt JJjQ(t) u(x, t) <m + II J J a T-ndS = 0 dQ(t) (21.1)，其中控制体 Q,(t)与时间相关，流体状态向量U和非黏性成分的通量张量T = Tmv + Tms是U P px pE <r-inv P(v - x) /9v(v - x) + pi pE\{\textbackslash\{\} - x) + pv，p，p、v、E和p分别是密度、速度向量、总能量和压力，而n表示具有速度x的移动边界9f2(i)的外法线。 细胞面上的隐形数值通量向量使用詹姆森中心方案或罗氏上风方案的二阶和三阶总变异减小（TVD）版本近似，应用守恒定律的单调上风方案（MUSCL）外推。 黏性通量是在细胞表面中心使用移位控制体上的梯度定理计算的。 由此产生的半离散方程组，dt \{Un) + R(Un) = Q (21.2)，其中现在fl是单元体积，U是单元状态向量，R是显式通量差分算子（残差），可以使用例如显式Runge-Kutta方案在时间上进行积分。 或者，在NSMB中，隐式方程组dt \{Ufl) + R(Un+1) =0 (21.3)

420 RIZZI、GORTZ和LEMOIGNE可以使用下上对称高斯-塞德尔（LU-SGS）方案进行集成。 这里介绍的所有固定翼计算都是使用詹姆森的中央通量方案和LU-SGS时间步进程序的标量版本进行的，以加速收敛到稳态。 NSMB代码的串行版本用于在富士通VX/3并行矢量计算机上进行计算，其三个处理元素的每个元件的峰值性能为2.2 Gflop/s。 该计算机由皇家理工学院（KTH）的并行计算机中心（PDC）运营。 21.2.1 双时间步进 如果要解决的问题因时间而异，那么时间准确性很重要，不能使用本地时间步进和多网格等收敛加速技术。 显式时间步进变得不切实际，甚至隐式方案也不适合，因为由于求解程序中需要的线性化，它们在时间上只有一阶准确。 双时间步进是解决这种困境的一种方式，因为它提供了更高阶的时间精度，同时允许使用高效的收敛程序，如多网格，但以额外的迭代循环为代价。 该技术引入了一个使用完全隐式方案进行实时时间准确时间步骤的外部时间步骤循环，以及一个具有虚构时间步骤的内部循环，以在每个实时时间步骤达到“稳定状态”。 然后，可以在内部循环中使用局部时间步进、多网格或隐式方案等加速技术，将解决方案迭代到“稳态”。 这种实施遵循了詹姆森提出的技术[13]。 将二阶精度的隐式向后差公式（BDF）方案应用于Eq.（ 21.3）产量：\textasciitilde{} [Un+1 ft"+1] - •\%- [Un Qn] + -i- [Un\textasciitilde{}l fi"-1] + R (Un+1) = 0 (21.4) 然后将此离散化的修改后残差R* (U)定义为R*\{u) = db tc/fi"+1] \textasciitilde{} h[un nn] + wt t\textasciicircum{}1 fin-1]+ R (t7) (2L5)，以便当迭代U满足系统R* (U) - 0时，达到二阶时间准确解Un+1。 因此，可以在虚构的时间步进程序中使用这个修改后的残差来制定一个新的常微分方程组，类似于Eq（21.2），将U向前迭代到时间准确水平Un+1 \textasciitilde{}(Un) + R*(U) = 0 (21.6)

\section{21.3 测试案例和网格}
DELTA-WING VORTEX FLOWS 421方程（21.6）在虚构时间t*中进行到稳态，使用流动求解器中实现的任何收敛加速方法，在Eq（21.5）中略微修改的残差。 21.2.1.1 启动程序 由于C/n\_1在第一个时间步骤中没有定义，因此需要特殊程序来启动二阶BDF 时间推进。 这可以通过多种方式完成，例如，先设置Un\textasciitilde{}1 = Un，然后像往常一样使用BDF，或者仅在第一步使用另一个只涉及Un+1和Un的方案，然后继续使用BDF。 后者更可取，因为前者牺牲了全球二阶准确性。 NSMB在第一步使用一阶后退欧拉方案，因为它易于实现，并且比二阶梯形方案具有更好的阻尼特性，此外，它仍然产生了一种全局二阶精确的方法。 21.3 测试案例和网格 最初只调查了姿态静止的独立三角翼平面，即没有对机身进行建模。 然而，已经测试了机翼的几种物理模型，尺寸和边缘对角不同。 21.3.1 模型一 这里处理的第一个（风洞）模型是一个简单的平板三角翼，具有70°前缘扫描，一个20.61英寸（523.5毫米）的根弦和一个15英寸（381.0毫米）的跨度在后缘，图。 1. 机翼厚度为0.5英寸（12.7毫米），前缘和后缘在顶部和底部表面使用25°的斜面进行了磨削。 21.3.1.1 H-H网格：模型1为该模型生成了一个结构化的H-H型网格，由1,168,512卷组成。 机翼法线方向的最小细胞尺寸为0.55 x 10\textasciitilde{}5根弦。 认为有必要使用18块的拓扑结构来正确地代表机翼的整个几何形状，包括所有斜面。 沿着翅膀沿弦方向放置了100多个点，跨度方向上放置了多达46个点，图。 2. 在顶端放置了大量网格点，其形状对初级涡流的发展和涡流分解的开始有重大影响。 由于斜面，通常的顶部网格退化在后部移动到上平面。 远场边界向上游延伸了6.5根弦和10

422 RIZZI, GORTZ \& LEMOIGNE 图1 模型1（根据Soltani \& Bragg，1988年）机翼下游的根弦，以及与机翼表面垂直的每个方向的6.5根弦。 图2 上表面网格，双斜面模型（模型1）21.3.2 模型二 第二个模型厚度为0.25英寸（6.4毫米），边缘仅使用底部表面14°斜面锐化，图。 3. 对于这个模型，生成了两个网格。

DELTA-WING VORTEX FLOWS 423 /\textasciicircum{}.4- 图3 模型2（在Soltani \& Bragg之后，1993年）21.3.2.1 E-C网格：模型2 第一个是Kumar[16]建议的7块嵌入式锥形（E-C）型网格，用于三角翼计算。 机翼由两个块（一个在上面，一个在下面）包围，它们共同形成一个圆锥体，而圆锥体又嵌入到标准H-O网格中，图。 4. 与机翼相邻的单元的高度为1.97 x 10-3根弦（此网格仅用于欧拉计算）。 机翼上方和下方的块分别在溪流、跨度和垂直方向有65点、65点和33点，总共有616,448个单元格。 网格在尖端有一个奇点，即锥形鼻的位置。 远场边界向各个方向延伸5根弦，机翼后部的块向上倾斜22°，以更好地解决尾流。 21.3.2.2 C-0网格：模型2 第二个网格是一个由714,000个单元格组成的6块结构C-0网格，见图。 5. 网格近似于实几何形状，使其在斜面处更简单，即斜面边缘在与菲亚特下表面的交叉处没有完全表示，而是在包括斜角和下表面的表面上的一个块面的投影。 所选的拓扑结构在顶点产生了奇点。 沿着机翼沿弦方向放置了130个点，在跨度方向放置了40个点，试图解决流场

424 RIZZI、GORTZ和LEMOIGNE图4域轮廓，嵌入式锥形（E-C）型网格，7块（模型2）图5 C-0型网格，在天头具有极性奇点（模型2）

DELTA-WING VORTEX 425在机翼上方正常流动。 网格在靠近顶点和前缘周围进行了精炼，图。 6. 远场边界在机翼体的上游延伸1根弦和2根弦下游，在垂直于机身上表面的方向延伸1根弦，在垂直于下表面的方向延伸1根弦。 图6 C-0型网格的顶点细节：网格加密 21.3.3 模型三 第三个模型是65° 三角翼，包括在这里，因为创建了一个C-0网格，在顶点上没有极性奇异线。 21.3.3.1 非奇异C-0网格：模型3 最上的奇异线被一个由4个“虚构”角组成的网格块所取代，一个由所谓的电视屏幕面形成的块，图。 7. 然而，这样做需要稍微圆一点，然后圆化前缘。 如果尖尖是尖锐的，这种方法是不可能的。 图中显示了包含超过70万个单元格的8块C-0网格。 8. 网眼向机翼的上表面和下表面以及前缘和后缘被精炼。 所有网格都是使用ICEM CFD Hexa生成的。 这种商业网状发生器的主要缺点是它不能有效地控制网状线的曲率，即网状线几乎是直的。 这一缺点使得在三角翼的尖锐前缘周围生成高质量的网格是一项困难的任务。 所需的网格正交性靠近

\section{21.4 静止翼计算和结果}
426 RIZZI, GORTZ \& LEMOIGNE Hach=0 1E15 fUphe=35 Re=le+30 Tine=0 图7特写：具有圆形顶点和圆角前缘机翼表面的Model 3的“电视屏幕”解决方案并不总是在顶点或前缘等关键区域实现，比较图。 7. 21.4 固定翼计算和结果 21.4.1 风洞墙的影响 使用前面描述的E-C型网格de Try[6]表明，当风洞墙建模时，压力吸峰值可以与实验数据更好地一致。 图。 9将自由空气的上表面压力分布与风洞壁建模的情况进行了比较。 在三角翼的顶端，自由空气溶液的吸力峰值明显降低。 这一结果部分解释了参考[10]和[11]中发现的计算和测量的表面压力之间的差异。 然而，风洞壁似乎对涡核轴在垂直于机翼表面方向的位置没有任何重大影响。 21.4.2 人工耗散的影响 研究了人工黏度对收敛历史的影响。 图10和图11显示了升力系数与模型3的迭代次数，三角翼的前缘略微圆润，其中只有四阶人工耗散系数k发生变化。 它

DELTA-WING VORTEX FLOWS r1och=0,1615 Fllpho=35 Re=1e+3D Tine=D 图8 带有“电视屏幕”块的C-0型网格（模型3）xfc»O.20EC2sb xfc=0.50 X/c=0.70 x/«0.20 ECWT25 xfc=0.50 x/c-0.70 V/\textbackslash\{\}. v. 风洞 //> '/ 自由空气 »v 图 9 风洞壁对不同弦站表面压力的影响，E-C网格，模型2，欧拉结果，a = 25°

428 RIZZI、GORTZ和LEMOIGNE可以看出，人工黏度越高，涡流的阻尼越多，导致流量的梯度更小，偏差更小，因此升力系数振荡的振幅更小。 0.36 0.35 0.34 O | 0.33 o fc g 0.32 0.31 0.30 0.29 0 5000 10000 15000 20000 迭代 图 10 Lift coeff. 与迭代次数，模型 3，人工耗散 coeff. k = 1/12 21-4-2.1 准稳定 流 虽然我们正试图计算 稳态解，但实际上该流场被更好地描述为准稳定。 当螺旋型涡流破裂发生时，流场会出现振荡。 在击穿点的下游，瞬时涡核轴围绕平均轴旋转，逆着主涡旋的旋转感。 此外，瞬时涡旋轴在主涡流的意义上相对于时间而旋转，频率与力和力矩振荡相同。 如果人工耗散没有抑制，这种准稳定性可以在模拟中捕捉到。 图。 12显示了模型2的正弦波振荡。 流场分析表明，人工耗散也对涡流分解位置有影响。 我们认为，这种准不稳定性不应该被抑制，人工耗散的数量应该只是稳定计算所需的最低限度。 否则，计算的解会随着耗散水平的函数而变化。 或者，人们可以说网格

DELTA-WING VORTEX FLOWS 429 图 11 提升系数与迭代次数，模型 3，人工耗散 系数 k = 1/33 10000 20000 迭代 图 12 升力系数 对比次数，如果迭代，6 块 CO 型网格，模型 2，欧拉解，人工耗散 系数 k = 1/33，a = 35°

430 RIZZI、GORTZ和LEMOIGNE必须进行改进，直到人工耗散的效果变得微不足道。 这不仅适用于尤拉计算，也适用于使用双方程模型的Navier-Stokes模拟。 21.4.3 物理建模的影响 物理建模对沿核心轴向速度的影响可以在图中看到35°的攻角。 13. 分层和湍流情况的击穿前的峰值轴向速度几乎相同，但尤拉情况较低。 但轴向速度沿着涡核的轴线以相同的方式发展，在所有三种情况下（在x/c = 40\%，在欧拉情况下，在42\%，在湍流情况下，在层压情况下，流体在几乎相同的弦位置被带到几乎相同的弦位置）。 尽管如此，可以说，涡流分解在隐性情况下稍早发生，这也可以在其他攻击角度上表现出来。 Darmofal[4]的理论发现证实了这一点。 使用准圆柱形近似，他证明了在低雷诺数下，黏度是方位涡量的耗散的原因。 发现这种黏性效应减缓了负方位涡量的非黏性产生，根据Brown和Lopez[3]的说法，涡核中流体的轴向减速与涡流分解有关。 根据该理论，黏度延迟了低雷诺数时的涡流分解。 然而，Darmofal发现，黏性效应在实际应用中可以忽略不计，即在实际雷诺数上。 对管道故障的数值调查证实了这一点[4]，目前的调查普遍支持这一点。 读者可以参考参考。[10]和[11]，更详细地调查物理建模对三角翼空气动力学的影响。 21.4.4 网格拓扑和机翼几何的影响 我们研究了网格拓扑、块布局和网格奇点对收敛速率、表面压力分布和轴向速度曲线的影响。 当通过模拟所谓的“电视屏幕”块来避免在高点计算网格中的极性奇异线时，收敛速率得到了改进。 实验和计算表面压力分布在嵌入式圆锥形（E-C）网格上实现了最佳一致性，缺点是轴向速度曲线不准确。 前缘几何（模型1与 模型2）对涡核轴上的压力分布、空气动力学系数、击穿位置和最大轴向速度有很小的影响，而事实证明它对轴向速度曲线的类型没有影响。 当使用H-H型网格时，沿速度的轴向分量

DELTA-WING VORTEX FLOWS 431 35 deg. 攻角尤拉 \textasciitilde{}N laminar \textbackslash\{\}\textbackslash\{\} 湍流 -41 »-\$»:• 'o.O 0.1 0.2 0.3 0.4 0.5 0.6 0.7 0.8 0.9 1.0 x/c 图13 物理建模对轴向速度和故障位置的影响（H-H型网格，a = 35°，Re = 1.97 x 106，Mx = 0.1615）涡核轴以自由流速度开始，然后随著涡流的发展在顶端和核心加速，图。 14. 轴向速度在分解前继续增加到一段距离。 最终到达一个位置，速度不会进一步增加，而是达到某种最大值，在某个距离内保持该幅度，直到突然穿过分解区域（“弧型”轴向速度曲线，也比较图。 参考中的10和12。 [11]）。 Hummel[12]、Verhaagen和Ransbeeck[23]在这两种情况下都使用5孔探针实验发现了这种速度曲线（这可能影响了探针在顶点区域的上游流场）。 Visbal[24]在他的一项关于俯仰三角翼的计算研究中发现了一个“弧型”速度曲线。 与此相反，当使用Kumar [16]建议的嵌入式锥形网格时，沿核心轴线的速度轴向分量在顶点急剧增加，随后逐渐减少，直到在突破点达到零轴向速度（“峰值型”轴向速度曲线，比较图。 15，还有图。 参考中的20。 2）。 这一结果似乎与Hummel、Verhaagen和Ransbeeck的实验数据以及其他关于三角翼涡流的数值研究相冲突。 模型2的C-0型网格产生了一个介于前两个之间的轮廓，图。 16. 最后，第三个模型，即具有圆形顶点和前缘的65° 三角翼，产生了第四种速度曲线，图。 17. 我们相信“弧型”配置文件，图。 14，是正确的，因为0.40 0.30 0.20 0.10 0.00 -0.10

432 RIZZI、GORTZ和LEMOIGNE图14“弧型”轴向速度曲线（模型1，双斜角模型，H-H型网格拓扑）0.20 0.10 o.oo 0.0 0.2图15“峰值型”轴向速度曲线（模型2，单斜角模型，E-C型网格拓扑）

DELTA-WING VORTEX FLOWS 433 0.4 0.3 ca - 0.2 1 0.1 0.0 -0.1 0.0 0.2 0.4 0.6 0.8 1.0 x/c图16“扁平型”轴向速度曲线（模型2，单斜面模型，C-0型网格拓扑）遵循物理推理：在突破点的上游，从机翼前缘散发出的自由剪切层沿着整个前缘连续地将涡量馈入涡流，导致轴向速度增加。 当前缘的涡量进给率超过下游可以对流的进给率时，涡流丝收紧，导致涡核的分解和横截面膨胀。 此外，弧型剖面是唯一具有“经典”涡流分解现象的轮廓，即核心轴向和周向速度分量的幅度突然下降。 峰值型剖面，图。 15，可以用圆锥奇点来解释，即由天头的圆锥网块来解释。 假定正确的“弧型”速度曲线仅在七个不同网格中的一个上计算，即模型1的18块H-H型网格。 在这方面，这个网格优于所有其他网格，因为顶点的网格点数量很多，前缘很精致，而且通常的顶点网格退化被移动到上部平面。 21.4.5 网格加密的影响 模型2的C-0型网在前缘进一步完善。 在精炼的网格上，在分解区域计算了更陡峭的核心速度梯度，图。 18. 更陡峭的梯度表示更好的分辨率

\section{21.5 Pitching Delta的初步结果}
434 RIZZI, GORTZ \& LEMOIGNE 0.4 0.3 r i - 0.0 ' 0.0 0.2 0.4 0.6 0.8 1.0 x/c 图17“圆形”轴向速度曲线（模型3，圆形前缘，顶端有“电视屏幕”的C-O型网格）涡流分解现象。 此外，它导致一个更上游的故障位置，即速度的轴向分量为零的点位于更靠近尖点的地方。 在精炼网格上，最大轴向速度也更高一些，这转化为机翼表面的涡流更强和更高的吸力。 最近，NSMB可以处理修补的网格，即具有本地精炼网格的块，其中考虑了跨块边界的单元格不匹配。 使用补丁网格，可以在感兴趣的区域本地提高网格分辨率。 这种可能性将用于细化涡流区域的机翼和机翼体网格，以接近网格分辨率。 对于以涡流分解为特征的情况，网格将在分解点附近进行细化，以解决流动的小规模细节。 预计吸力面积的分辨率会大幅提高，这将导致计算表面压力和实验表面压力之间的更好一致性。 21.5 Pitching Delta的初步结果 对25度和35度的两个静态情况进行了网格和计算的不同参数的第一次测试。 这些是欧拉计算，自由流马赫数为0.1615，隐式时间积分为

DELTA-WING VORTEX FLOWS 435 0.4 0.3 r i - 0.0 -0.1 0.0 0.2 0.4 0.6 0.8 1.0 x/c 图18 网格加密对沿涡核轴的速度轴分量的影响，CO型网格，模型2，欧拉解，k = 1/33，a - 25°（CO =>•非精炼网格；C02 =>在前缘精炼网格）标量LU-SGS方法和人工耗散的中心方案。 图19所示的结果与[2]的静态实验值非常一致，并鼓励我们继续动态振荡。 作为我们对俯仰计算的初始模拟，我们求解Euler 方程，并与[2]中的实验测试案例进行比较。 三角翼的音高围绕其40\%的和弦位置，降低的音高频率（或音高率，通常表示为k，定义为k = u>c/2Uoo）等于0.0376。 计算的自由流马赫数为0.2。 计算了两个不同的条件：一个用于低振幅振荡，另一个用于大振幅振荡。 小振幅振荡 平均攻角为12°，机翼以±3°的俯仰振荡。 大振幅振荡 机翼在22°左右的俯仰振荡，大振幅振荡为±18°。 C02-mesh CO-mesh 所有计算都是使用隐式双时间步进法和不同的Roe上风方案进行的：一阶、二阶和三阶。 足够的

436 RIZZI、GORTZ和LEMOIGNE 静态结果 1.5 1.0 C\_N 0.5 0.0" • ' • ' ' ' 0.0 20.0 40.0 60.0 AoA [deg] 图 19 与静态测量相比，在两个静止a值下计算的法向力Cn法向力结果与实验结果一起显示在图20中（见[2]）。 通过为外部时间步骤使用完全隐式方案，可以进行大型时间步骤。 双时间步进方法的问题在于找到外部时间步进的最佳值，因为外部时间步进的值过大需要大量的内部时间步进才能收敛到稳定状态，而如果外部时间步进的值太小，与Runge Kutta方案相比，双时间推进方法的增益很小。 大振幅情况下的一个振荡周期是用具有60个真实时间步进的双时间步法计算的，其中运行100次重写以获得“稳态”（在一阶和二阶情况下）。 在图21中，绘制了二阶计算的收敛历史。 该图清楚地显示了双时间步法的原理，内部迭代的收敛在每次外部迭代中重复。 使用三阶Roe上风方案时遇到了收敛问题。 将亚迭代次数增加到200次可以提高收敛，但只有一部分周期振荡可以成功计算。 在图20中，一阶和二阶方案的滞后回路与实验值相比非常好。 然而，一阶方案预测法向力低于实验，环路也更薄。 另一方面，二阶结果超过了实验*计算

DELTA-WING VORTEX FLOWS 437 1.50 C\_N 1.00 0.00 俯仰振荡 A,"-''A*ii-\# >\textasciicircum{} A I +W"- -r-f + 4 * •'V -- 实验 + 1阶100内侧 • 2阶100内侧 A 3阶200内侧-低振幅0.0 10.0 20.0 AoA [度] 30.0 40.0 图20与实验数据相比，俯仰振荡的计算法向力Cn 收敛历史（第2阶61外100内侧）upww\textasciicircum{}fPIJl 2000 3000 4000 5000迭代次数图21 收敛大振幅情况的历史（二阶）显示每个外部时间步骤的内时间步 收敛

\section{21.6 结论和展望}
438 RIZZI、GORTZ和LEMOIGNE在高攻击角度下的实验值。 用三阶方案计算的循环的第一部分与实验迟滞循环的一般形状更一致，尽管在这种情况下，值也被过度预测了。 对于所有方案，曲线的斜率都非常相似，比实验数据略高。 即使二阶方案已经产生了有希望的结果，使用三阶方案计算的完整振荡可能会给出更准确的结果。 有趣的是，方案的顺序越高，大攻角的法向力就越高，而在较小的角度下，不同的方案预测的力非常相似。 低振幅振荡仅使用24个实时步进和80次迭代进行计算，因为计算要求较低。 在这种情况下，结果低于实验值，但线的斜率略高。 然而，这里的结果也令人鼓舞，因为这些点以几乎完美的顺序对齐，就像实验所建议的这种攻击角度一样，表明代码NSMB能够再现涡流。 21.6 结论和展望 21.6.1 固定机翼 事实证明，网状生成过程在高攻角下模拟三角机翼的空气动力学中起着重要作用。 菲奥场的细节，如沿涡核轴的速度轴分量或击穿位置，高度依赖于顶点和前缘的网格拓扑结构和网格细度。 风洞壁的建模可以证明对靠近高点的表面压力分布有不可忽视的影响。 当表面压力感兴趣时，嵌入式锥形网格似乎优于其他网格拓扑结构。 然而，沿涡核轴的速度轴向变化被认为对这种网格类型来说是非物理的。 假定正确的轴向速度曲线是在高端没有奇点的精细H-H型网格上计算的。 当Apex奇点被替换为“电视屏幕”块时，收敛速率被认为得到了改善。 人工耗散被证明会影响旋转击穿区域的阻尼。 它必须针对每个单独的测试案例进行调整。 未来的静止-delta模拟工作包括持续努力寻找“最佳网格”。 此外，它将为具有涡流分解的测试案例显示时间准确的解决方案。 流动求解器新可用的补丁网格功能将用于细化感兴趣区域的涡流场。

\section{21.7 致谢}
\section{参考资料}
DELTA-WING VORTEX FLOWS 439 21.6.2 Pitching Delta 在俯仰中振荡的70° 三角翼获得的初步结果非常有希望，因为整体迟滞效应在空气动力的历史中再现。 然而，这些结果是初步的，可以通过进一步开发代码来改进，例如自我适应时间步骤，以及进一步体验这种不稳定流场的物理特性。 例如，需要进一步调查振荡期间实时时间步数的影响。 在获得与周期数无关的结果之前所需的周期数也是如此（第一次振荡不会产生完全相同的结果）。 数值参数，如亚迭代次数、空间离散格式和CFL数也非常敏感。 黏度的影响，已知会影响涡旋在机翼背侧的位置，从而改变升力，也必须通过进行Navier-Stokes计算来检查。 当更好地理解所有这些参数时，下一步将是修改投球条件。 俯仰率、平均攻角、旋转轴位置的影响需要进一步调查。 除了动力外，作者还打算观察机翼周围的非定常流动场和振荡期间的表面压力。 最终目标是计算整个大振幅振荡周期，流可视化突出显示涡流形成和爆裂，以便更好地了解实验中观察到的滞后现象。 21.7 致谢 我们要感谢许多人的帮助：CERFACS的Denis Darracq为我们提供了NSMB-ALE代码，EPFL的Jan Vos在调整NSMB方面的帮助，SAAB的Kari Munukka和Peter Johansson为网格生成提供了非常有帮助的建议，最后是并行计算机中心（PDC）在运行机器方面的帮助。 参考资料1。 Arthur，M.T.，Brandsma，F.，Ceresola，N. \& Kordulla，W. 1999年三角翼在俯仰运动中的自流流的时间精确欧拉计算，AIAA论文99-3110。 2. 布兰登，J.M. 1991年动态失速效应及其对高性能飞机的应用，AGARD-R-776第2-1页2-15 3。 布朗，G。 L.和Lopez，J. M.，轴对称涡流分解，第2部分。 身体的

440 RIZZI、GORTZ和LEMOIGNE机制，流体力学杂志，卷。 221，1990年，pp. 553-576。 4. 达莫法尔，D。 L.，轴对称涡流分解机制研究，博士论文，麻省理工学院航空航天系计算航空航天科学实验室，1993年11月。 5. Darracq，D.，Champagneux，S.和Corjon，A.，用空气弹性AUSM+隐式求解器计算非稳态湍流 翼型流动，AIAA论文98-2411。 6. de Try，F.，风洞中70°扫三角翼的涡流行为计算研究，理学硕士。 论文，KTH，航空系，斯德哥尔摩，1999年。 7. Ekaterinaris，J.A. \& Schiff，L.B. 1991年Navier-Stokes振荡双三角翼解决方案，AIAA论文91-1624 8。 弗里茨，W。 关于高入射率振荡的三角翼的非稳态涡流的数值模拟，数值流体力学笔记，卷。 72页 162-169 9. Gacherieu，C，Weber，C.和Coriolis，G.，围绕完整飞机配置的跨声速湍流的代数和一方程湍流模型的评估，AIAA论文98-2737。 10. Gortz，S.，70°扫描三角洲翼的涡流分解的计算研究，理学硕士。 论文，KTH，航空系，Skrift 98-27，斯德哥尔摩，1998年。 11. Gortz，S.，Rizzi，A.，Munukka，K.，对扫三角翼的涡流破裂的计算研究，AIAA论文99-3118，1999年。 12. Hummel，D.，Untersuchungen iiber das Aufplatzen der Wirbel an schlanken Deltafliigeln，Zeitung fur Flugwissenschaften 13，Heft 5，1965年。 13. 詹姆森，A.，使用多网格的时间依赖性计算，应用于通过翼形和机翼的非稳流动，AIAA论文91-1596，1991年6月。 14. 坎迪尔，O.A. \& Chuang，H.A. 1990年三角翼正在经历俯仰振荡的涡流流流的计算，AIAA杂志，卷。 28 不。 9月9日 1990年第1589-1595页。 15. 坎迪尔，O.A. \& 坎迪尔，H.A. 1994年跨声速涡流分解流中65度三角翼的俯仰振荡，AIAA论文94-1426。 16. Kumar，A.，使用嵌入式锥形网格准确开发前沿漩涡，AIAA期刊，卷。 34，不。 1996年10月10日。 17. LeMay S.P.，巴蒂尔，S.M. \& 纳尔逊，R.C. 1990年俯仰三角翼上的涡流动力学，飞机杂志，卷。 27 不。 1990年2月2日第131-138页18。 Modiano，D.和Murman，E.M. 1993年带涡流分解的三角翼周围流动的自适应计算，AIAA论文93-3400。 19. O'Neil，P.J.，Roos，F.W.，Kegelman，J.T.，Barnett，R.M.和Hawk，J.D.，开发涡流的流动特征调查，Rep No. NADC- 89114-60，1989年5月。 20. Soltani，M。 R.和Bragg，M. B.，亚声速流动中振荡70度三角翼的实验测量，AIAA论文88-2576-CP。 21. 汤普森，S.A.，巴蒂尔，S.M. \& 纳尔逊，R.C. 1989年，纤细的三角翼上的分离流场正在经历瞬态俯仰运动，AIAA论文89-0194。 22. 汤普森，S.A.，巴蒂尔，S.M. \& 纳尔逊，R.C. 1990年三角翼高角度攻击机动的表面压力，AIAA论文90-2813。 23. Verhaagen，N.，Ransbeeck，P.，前缘涡流芯流动的实验和数值调查，AIAA-Paper 90-0384。 24. Visbal，M.，Pitching 三角翼上的涡流分解结构，AIAA论文93-0434。

三角洲翼涡流441 25。 沃斯，J。 B.，Rizzi，A. W.，Corjon，A.，Chaput，E.和Soinne，E.，NSMB（Navier Stokes Multi Block）联盟内空气动力学的最新进展，AIAA论文98-0225。 26. Weber，C.等人，NSMB结构化多块求解器在空气动力学中的最新应用，在计算流体动力学-1998年的前沿中，（编辑）D.A. Caughey和M.M. 哈菲兹，世界科学出版社，新加坡，1998年。

\clearpage
\chapter{DLR的22个选定的CFD功能}
\section{22.1 介绍}
DLR W的选定CFD功能。 Kordulla1 22.1 介绍 计算流体动力学（CFD）自七十年代以来经历了戏剧性的发展，如果将Navier-Stokes的二维解决方案作为起点。 最早的此类解决方案之一是MacCormack提出的基于显式预测校正器Lax-Wendroff公式的变体[1]。 目前的技术状态反映在[2]或下面引用的一些论文中。 原则上，今天几乎任何层流都可以被预测。 瓶颈是，就像一直和可能一样，适用于复杂的3D几何形状的网格生成。 可以使用Reynolds或Favre平均Navier-Stokes（RANS）方程以适当的方式模拟光滑表面上的静止湍流流，而没有大量分离，前提是所使用的代码，包括涉及的湍流模型，已针对所涉问题类别进行了适当的评估。 对于具有大规模分离的湍流流或从层流过渡到湍流的流，仍有基础研究需要进行，尽管可以认识到一些成就[10，11]。 由于在这种流中需要解决的非常不同的长度尺度，尽管算法和网格自适应自大约二十年以来进行了努力，尽管在计算机架构和加速方面取得了进步，但我们仍然离直接解决层压Navier-Stokes 方程还很远。 大涡流模拟等中间步骤受益于较少的网格点，只需要解决大型湍流涡流，而必须以在子尺度层中建模物理学为代价，该层似乎是1 DLR（Deutsches Zentrum fur Luft- und Raumfahrt），37073 哥廷根，德国，电子邮件：Wilhelm.Kordulla@dlr.de 计算流体动力学 - 2002 编辑：David A. Caughey和Mohamed M. 哈菲兹 ©2002 世界科学

\section{22.2 CFD 发展}
444 KORDULLA对湍流来说比对过渡性流动来说更容易[5]。 过渡和湍流流动的建模仍然是将CFD应用于现实高雷诺数的航空流动的主要绊脚石。 在高温或反应性流动的情况下，风洞和自由飞行情况下的反应机制建模增加了这种难度[4]。 在这里，在对工程师来说新的情况下，数值模拟真正具有预测性之前，需要进行大量的工作。 从某种意义上说，这种情况为实验者和地面设施提供了人寿保险。 在DLR，开发数值方法来整合支配各种流动的偏微分方程有着悠久的传统。 考虑的流量包括例如机翼后面的尾迹[7]、涡轮机械中的三维黏性流量[6]、空气弹性现象[9]和燃气轮机燃烧室中的反应流量。 当然，航空航天是一个核心应用领域，因为相应的研究和技术开发领域是DLR的主要工作重点。 22.2 CFD的发展 大约20年前，DLR存在各种CFD方法，在某种意义上相互竞争。 在不伦瑞克中，选择了詹姆森的显式Runge-Kutta时间步进方案[37]，而在Gttingen MacCormack的显式隐式方法[38，39，40]中，使用了Beam-and-Warming的隐式方案[41，42]，以及Runge-Kutta时间步进方法的另一个版本[43，44]。 这些方法基于结构化网格。 稍后，在Gttingen开发了一种使用所谓的非结构网格的有限体积法，名为TAU [45]，缩写为“三角形、自适应和非结构”。 多年来，随着资源的不断减少，幸存下来的方法是在不伦瑞克站点开发的Runge-Kutta时间步进代码的演变，并产生了MEGAFLOW项目，以及混合3D TAU代码。 对于空间流动应用，正在使用MEGAFLOW中早期版本的流动求解器的衍生产品，称为CEVCATS-N。 22.2.1 MEGAFLOW项目 在德国航空航天研究计划的框架内，DLR领导了一个名为MEGAFLOW的国家CFD项目[16，20，21]。 这个首字母缩写词结合了生成MegaCads网格和流体流动FLOWer的控制方程的数值积分系统的名称。 目标是为德国工业提供一个高效的数字工具，用于预测整个飞机在巡航中的性能，以及

在起飞和降落配置中，DLR 445的选定CFD功能。 为了满足工业实施的要求，飞机工业、DLR和几所大学进行了集中的合作努力。 该项目的第一阶段由德国联邦教育、科学、研究和技术部（BMBF）部分资助，于1995年年中开始，1998年底结束。 大约125人工年的总努力花在了五项活动上：一项是开发电网生成器MegaCads，另一项是改进块结构求解器FLOWer，另一项是工业应用软件系统的验证。 由于许多合作伙伴致力于一个产品，模拟系统的软件质量管理是进一步的活动，最后，总努力的一小部分专门用于使用基于TAU代码的非结构化网格的替代方法的调查。 该系统辅以用于计算空气动力学力的后处理工具和空气动力学优化系统。 MegaCads（多块椭圆网格生成和CAD系统）展示了图形用户界面、可视化图形模块以及CAD和网格生成技术模块[48]。 主要致力于改进为工业飞机应用生成表面和体积网格的技术。 这些努力不仅涉及电网的质量[20，21]，还涉及发电过程的加快[16]。 这包括应用高级技术来规范椭圆网格生成的控制项，无论是在网格块内还是跨块边界（例如[47，48]），以及使用多网格和并行化[16]。 由于参数化概念，MegaCads可用于空气动力学优化循环。 例如，结果显示在[20]中。 FLOWer集成了可压缩的3D RANS方程[50]，并基于块结构网格上的显式细胞顶点有限体积公式。 基线方案使用空间中的中央离散化，具有詹姆森型人工黏度和各种显式多阶段时间步进能力。 通过使用局部时间推进、隐式残余平滑和多网格加速收敛到稳态。 根据之前在早期国家项目（POPINDA）的工作，该代码已被扩展为平行代码[16，20，21]。 对于使用消息传递接口MPI的计算机来说，它是完全可移植的。 可实现的加速在很大程度上依赖于适当的处理器负载均衡。 对于黏性（接）不可压缩流动的数值准确模拟，例如起飞和降落情况，可压缩流动方程的预处理在FLOWer中实现。 最初由J发明的嵌合体式网格技术进一步增强了该方案。 Steger和J。 Benek大约在二十年前就解决了商店分离问题。 这允许在描述复杂配置时更加灵活，而牺牲非平凡的任务，为覆盖或修补提供保守的接口边界条件

446个KORDULLA网格，特别是3D网格。 FLOWer也模拟了非稳流动。 网格运动由控制方程中的附加项表示。 在詹姆森之后，通过被称为双时间推进的简单隐式方法规避了与黏性流动模拟的显式方案相关的严重时间步骤限制。 与基本的显式方案相比，与多网格加速相结合，由此产生的方案证明了通过振荡机翼的黏性流量的速度提高了约三个数量级[16]。 结合并行化，这种方法有望处理最复杂的配置流程。 由于代码能够处理网格变形，FLOWer系统可以用于空气弹性流结构模拟。 使用嵌合体技术可以模拟在相对运动中流过物体或组件，例如火车进入隧道或相互通过的情况[46]。 鉴于飞机制造商在保证阻力计数方面的需求，准确模拟黏性流动的一个主要关注点是纳入过渡的物理模型以及湍流。 因此，MEGAFLOW项目与大学合作，采取各种方法来应对这一问题[16、19、20、21]。 由于达到了高水平的质量和效率，以及由于根据工业要求对系统进行了广泛的验证，MEGAFLOW系统正在被戴姆勒克莱斯勒航空航天公司在新飞机的设计过程中密集使用。 然而，在改进物理建模、增强网格生成的自动化、缩短周转时间、进一步验证、为多学科模拟提供接口以及增强空气动力学优化系统的能力方面，仍然需要进一步的开发步骤。 用于确定低超音速和跨声速流动制度中通常需要的再入配置的动态系数，使用FLOWer[29]。 相应的结果是作者已知的第一个完整起重车辆的结果。 22.2.2 CEVCATS-N CEVCATS-N与N代表非平衡相结合的CEVCATS，FLOWer的前身，但为高超声速流动模拟改进和调整，以及用于反应流的黏性流动求解器NSHYP的物理建模[36，42，50]。 代码CEVCATS-N集成了有限体积公式中黏性流动的控制方程。 流体可以是处于热化学平衡或非平衡状态的完美气体或气体混合物。 在基于Tannehill热力学和运输特性近似的矢量化曲线拟合的帮助下，描述了热化学平衡下空气的高温效应。 给

选定的CFD能力在DLR 447非平衡中流动，通常应用五种模型。 适当的墙边界条件允许考虑催化效应、辐射和各种传热条件。 空间离散化基于显式二阶有限体积方案，变量存储在顶点。 热力学和化学源术语以点隐式的方式考虑。 数值方案使用混合上风通量-向量分割：范利尔方案用于强冲击区域，以及根据Liou和Steffen其他地方的AUSM方案。 通过MUSCL外推实现二阶精度，以正确捕捉强冲击和接触不连续性。 局部时间推进、隐性残差平均和全多网格（带半粗化）用于将收敛加速到稳态，类似于FLOWer中使用的程序。 与典型的黏性跨声速流动模拟相比，对于再入流计算，与声波速度相比，数值黏性特征值很大。 在强冲击下，在残余平滑的帮助下获得的大库兰特数将导致溶液发散。 因此，采用了自适应时间步骤，在强冲击下将Courant 数减少到大约一个。 CEVCATS-N在DLR的NEC SX 4/4超级计算机上的计算性能约为每个单处理器1.1 GFLOPS。 平衡计算的每个网格点和多网格周期的CPU时间为30 fxsec，非平衡计算为70 //秒。 并行使用4个处理器将后一个时间减少到20/isec。 CEVCATS-N程序在欧空局和其他组织组织的各种测试案例中得到了广泛的验证，特别是通过与实验获得的数据的比较，例如 [24，25，26，27，28，31，32，34，35]。 22.2.3 TAU TAU的开发始于九十年代初[45]，用于非结构化三角网格，在MEGAFLOW项目以及几个欧盟资助的项目的框架中得到了加强，例如，见 [3]。 参与MEGAFLOW特别能够详细比较结构化和非结构化网格的预测结果。 特别是，非结构化网格的代码必须重现结构化网格的结果，尽管由于更复杂但更灵活的代码结构，计算工作量更大。 一旦复杂配置的表面描述得足够准确（防水），非结构化网格的生成时间很短，包括应用局部网格加密与缺少作用的更大单元。 如果考虑不黏性流动，这当然是真的。 对于高雷诺数下的黏性流动模拟，例如靠近墙壁的边界层的适当分辨率将需要大量的规则形状的细胞或具有

448 KORDULLA 相应边缘之间非常尖锐的角度。 前者导致巨大的计算时间，后者导致热量转移等量的准确性或振荡分布降低。 因此，混合网格被用于远离墙壁的非结构化四面体网格和结构网格，通常是棱镜层，靠近表面[13，15，16，17，18，3]。 原则上，DLR的策略是与市售的3D 网格生成软件生产商合作，因为该领域的内部工作越来越不实惠。 然而，适应性或网格变形等功能是内部构建的，例如 [15]。 今天的TAU代码是一个软件包，用于在上述网格上求解欧拉方程或RANS方程。 它由预处理部分、流动求解器和适应模块组成。 预处理部分目前支持四面体、棱柱形、金字塔形和六面体元素。 复杂几何形状的网格生成，例如机翼/车身/塔架/发动机配置，仅在几天内完成，而关于相同几何形状的块状结构网格需要几周时间。 请注意，求解器完全独立于网格单元格的类型，因为它使用双网格方法采用基于边缘的数据结构。 预处理不仅从初始网格和相应的数据结构中生成双网格。 它还为多网格方法和网格分区建立了聚集的粗网格，即 区域分解，用于并行计算机。 此外，对边缘进行着色，允许在矢量计算机上进行矢量化或在基于缓存的机器上进行缓存优化。 流动求解器尽可能复制FLOWer方法中存在的功能。 因此，求解器对时间梯度的离散化使用显式多阶段时间步进Runge-Kutta方案。 在时间准确模拟的情况下，使用双时间推进方法。 为了将收敛加速到稳态局部-时间推进，采用了残余平滑和多网格技术。 求解器包括用于计算对流通量的各种中央和上风方法，以及用于空气动力学流动的一些湍流模型[14]。 TAU的适应模块使用基于残差和/或梯度的传感器，根据有问题的流动特征提供混合电网的自动局部细化。 结构化子层中的点可以重新分配，以提供改进的边界层分辨率。 由于TAU的一个需要的应用是具有结构动力学代码的多学科链接[12]，因此该系统中增加了一个有效的网格变形工具。 值得注意的是，DLR未来CFD代码的开发旨在将代码系统FLOWer和TAU的良好功能结合到改进的TAU+中。 此外，目前正在转移CEVCATS-N中存在的know-how，以从主要优势中获利

\section{22.3 最近的申请}
DLR 449 TAU的选定CFD功能，即在必须考虑新配置时缩短周转时间，并为空间流动活动TAU+N（单平衡）生成更灵活的代码。 最后，近年来建立了可压缩TAU代码的衍生物，以应对燃气轮机燃烧室中的不可压缩流动[13]。 这项工作是在DLR项目的框架内进行的，其中流动求解器与描述燃烧过程（包括辐射）的模型相结合。 22.3 最近的应用 为以下列表选择的参考资料概述了三维的广泛航空航天应用，从不可压缩到反应性高超声速流动，从稳态到非稳态。 这里描述了几个例子，以给人以DLR存在的能力的印象，第一个例子取自高升空力学，见图。 1和2 [51]。 配置是DLR ALVAST风洞模型，这是一架类似于空中客车A320的双喷气窄体商用飞机的1:10比例模型。 图1机翼-机体交界处附近的表面网格：块状结构（左）和混合非结构（右）。 交互式MegaCads模块已被用于生成具有920万个网格点的结构化网格，用于完整（一半）配置。 经验丰富的用户需要大约3个月的时间来生成网格，包括板条和襟翼间隙。 图。 1显示左侧的表面网格。 在右侧，观察到具有580万网格点的适应性非结构化混合网格的表面网格。 网格是用商业软件CENTAUR [53]生成的。 根据模型的清晰定义的CAD描述，网格生成流程只花了大约一周时间。 图2显示了

450 KORDULLA实验RANS；FLOWer k - ©模型RANS；TAU，S.-A.模型2.0 1.0 v\textasciicircum{}r\textasciicircum{}。 * \%£T® - 实验RANS；FLGWer k - 建模RANS；TAU，S.-A.模式！ 图2 实验观测的升力/阻力曲线与使用结构化和混合网格获得的预测曲线的比较。与DNW的大型低速风洞中获得的实验数据相比，用FLOWer和TAU预测的相应的升力和阻力极性。 无论使用湍流模型，该协议在线性制度中都很好，此后都会如预期的那样恶化。 [22]，图3DLR-F6配置上的适应表面网格和最终压力分布讨论了FLOWer对流过直升机的结果。 图3和图4提供了使用非结构化混合网格的TAU获得的结果的印象[3，16，21]。 巡航条件下的通用飞机配置是DLR-F6机翼-机身-塔架-机舱配置。 从带有棱镜子层和四面体元素的初始网格开始

选定的CFD功能在DLR 451 500 1000 1500 2000多电网循环41evel-W TAU-SA，CI - TAU-SA，10* Cd ® Flower-BL，CI O Flower-kw，CI • Flower-BL，10 * Cd n Flower-kw，10 * Cd 500 1000 1500 2000多网格循环41evel--W......f。 D图4 收敛完整适应周期的历史，与结构化网格获得的积分数据的比较。 在有140万个网格点的外围区域，为一半的配置获得了350万个点的最终网格。 Spalart-Allmaras1 湍流模型已被用于模拟。 图。 4显示适应的表面网格和最终压力分布，而图。 5反映了与基于380万点的结构网格的结果相比，残差以及升力和阻力的收敛的减少[21]。 图5表示不可压缩TAU代码适用于燃气轮机通用燃烧室中惰性流动的能力[13]。 适应特征被整齐地观察到。 图5 通用燃烧室部分表面的适应网格和垂直切割中的轴向速度。

452 KORDULLA 图6 在DNW-RWG风洞条件下（M\textasciicircum{} = 6，Reoo = 3.1 • 1065层）X-38体圈上的热负荷和溪流涡流可视化。 其余示例涉及使用CEVCATS-N模拟升降再入车辆的空气动力学和空气热力学性能。 要计算此类车辆的动态系数，请参阅读者[29]。 超音速运输飞机和航天器优化的结果，包括考虑前体与空气呼吸发动机入口的相互作用，如[30]和[23]所述。 胶囊型车辆设计的优化结果在[31]中可以找到。 图6显示了在哥廷根风洞DNW-RWG条件下，NASA3s船员救援车的X-38演示机体襟翼热负荷的预测。 该实验的主要目标是在ONERA S4 Modane风洞条件下（Moo = 10，a = 40°）X-38背风侧的德国国家技术开发图7 层流模式的框架下进行的。

选定的CFD能力在DLR 453升力阻力俯仰矩jf / \textasciicircum{} /\%/ ¥// \textbackslash\{\} X --:S3ModaneWT / J \textbackslash\{\} -- 图8 X-38在Moo = 6的空气动力学系数，用于不同攻击角度和身体襟翼偏转。TETRA程序旨在证明襟翼上边界层的Gortler型不稳定性。 除了经典的压力和热传递测量外，还使用热电偶液晶来检测襟翼上的涡流痕迹。 在实验中，没有看到明确的证据。 然而，在数值模拟中，襟翼周围的块中有100多万个点，可以清楚地识别出具有流向涡流的流动结构，见图。 6 [35]。 这是第一次在逼真的3D配置的表面上预测流向逆向旋转的涡流。 需要进一步的调查来证实调查结果。 图7和图8显示了在ONERA的风洞中实现的不同风洞条件下的X-38配置的结果。 图8显示了演示器背向一侧的预测和实验观察到的皮肤摩擦线的出色比较，由于那里的密度低和相关的收敛问题，这很难实现。 图。 8展示了该代码在™»»r-图9的X-38表面飞行和HEG条件下的预测压力分布时，再现实验获得的积分量和俯仰矩的能力。

\section{22.4 去哪里}
454 KORDULLA T[K] 图10 M = 25的自由光条件下选择皮肤摩擦线的温度分布。 [28，31，34]。 图9和图10显示了高超音低反应的预测结果，前者是哥廷根高焓设施HEG的风洞条件和相应的自由光条件。 图。 10在M = 25的不对称圈设置下，在自由光下为X-38提供温度分布。 22.4 去哪里 预测CFD的未来，一方面在很大程度上取决于对物理学进行适当建模的能力，并使具有成本效益的模型易于在代码中实现。 另一方面，它依赖于增强的网格生成能力，提高了灵活性和速度，根据“良好的网格是解决方案的一半”的规则。 因此，DLR未来将专注于一个基于TAU的现代数据结构的非结构化混合网格的代码。 除了这个选择之外，游戏的名称是各种学科的多学科网络，这些学科被包裹在一些用户友好的信息技术环境中，包括使用具有超快模拟代码和高效耦合程序的元计算，可能超越了目前松散耦合的第一步。 目前，DLR支持为多学科耦合生产适当外壳的内部项目，例如用于跨声速飞机（AEROSTABIL，AMANDA [52]）或飞行力学（AEROSUM）的空气弹性模拟。 这也要求在优化过程中以更快、更自动的方式使用CFD。 在空间应用中，一个项目通过将热反应流与结构中的传热和相应的响应作为Irst步骤（IMENS）相结合，开始以多学科方向工作。

\section{22.5 致谢}
\section{参考资料}
DLR 455 22.5 的精选CFD能力 致谢 如果没有他在不伦瑞克和哥廷根的同事的支持，作者不可能写这篇论文，因此感谢他们。 也特别感谢S。 Briick对IATpjX的支持。 还要感谢由BMBF、DFG、DLR、ESA/ESTEC和欧洲共同体资助的各种项目，这些项目允许生成数据。 参考资料1。 MacCormack，R.W.，黏度对超高速撞击坑的影响，AIAA论文69-354，1969年。 2. 考伊，D.A. 和哈菲兹，M.M. （编辑），计算流体动力学-2000的前沿，世界科学出版公司PTE。 有限公司，新加坡，2000年。 3. Schwamborn，D.，验证和验证：使用混合非结构化网格的求解器的例子，VKI系列讲座2000-8，比利时布鲁塞尔，2000年6月5日至8日。 4. Sarma，G.S.R.，高超声速流动模拟中的物理化学建模，航空航天科学进步36，281-349，2000年。 5. Jimenez Hartel，C.J.，Analyse und Modellierung der Feinstructur im wandnahen Bereich turbulenter Scherstromungen，DLR-FB 94-22，1994年。 6. Eulitz，F.，Engel，K.，Niirnberger，D.，Schmitt，S. \& Yamamoto，K.，关于用于非稳定涡轮机械流动的并行时间精确Navier-Stokes求解器的最新进展，ECCOMAS 1998，1998年9月7日至11日，希腊雅典，会议记录，John Wiley \& Sons，卷。 1，252-258，1998年。 7. 格兹，T。 \& Holzapfel，F.，翼尖电波，湍流和排放分布，AIAA杂志，卷。 37，1270-1276，1999年。 8. Wegner，W.，使用虚拟网格变形方法进行弹性振荡机翼的空气动力学，RTO-AGARD报告R-822，数值非稳态空气动力学和空气弹性模拟，AGARD结构和材料小组第85次会议，丹麦奥尔堡，1997年。 9. 格里伯，B。 \& Carstens，V.，黏性效应对振动超音速压缩机叶片空气动力学阻尼的影响 一项数值研究，ASME TURBOEXPO 2000会议记录，2000年5月8日至11日，德国慕尼黑，论文2000-GT-0383（将发表在涡轮机械杂志上）。 10. Theofilis，V.，两个空间维度的线性不稳定性分析，ECCOMAS 1998，第四届欧洲CFD会议，1998年9月7日至11日，希腊雅典，会议记录，第1卷，547-552。 11. 海因，S.，哈尼菲，A。 \& Casalis，G.，非线性过渡预测，ECCOMAS 2000，2000年9月11日至14日，西班牙巴塞罗那，会议记录。 12. Beckert，A.，Gerhold，T.，Giinther，G.，Schwamborn，D. \& Weinman，K.，DLR TAU代码在混合电网流体结构相互作用中的应用，ECCOMAS 2000，2000年9月11日至14日，西班牙巴塞罗那，会议记录。 13. 施瓦本，D.，格霍尔德，T。 \& Kessler，R.，DLR TAU代码概述，第一届ONERA/DLR航空航天研讨会，ONERA，巴黎，1999年6月21日至24日。 14. Weinman，K.，湍流建模对湍流空气动力学流动的影响，ECCOMAS 2000，2000年9月11日至14日，西班牙巴塞罗那，会议记录。 15. 格霍尔德，T。 \& Evans，J.，使用自动适应的DLR TAU代码高效计算复杂配置的3D流，注释

456 KORDULLA数值流体力学，卷。 72，Vieweg，178-185，1999（将出现在AST中）。 16. Aumann，P.，Barnewitz，H.，Schwarten，H.，Becker，K.，Heinrich，R.，Roll，B.，Galle，M.，Kroll，N.，Gerhold，T.，Schwamborn，D. \& Franke，M.，MEGAFLOW：平行完整飞机CFD，将出现在并行计算，应用特刊中，“航空航天中的并行计算”。 17. 施瓦本，D.，格霍尔德，T。 \& Hannemann，V.，关于DLR TAU代码的验证，数值流体力学笔记，卷。 72，Vieweg，426-433，1999年。 18. Gerhold, T., Friedrich, O., Evans, J. \& Galle，M.，使用DLR TAU代码的复杂三维配置的计算，AIAA 97-0167，1997年。 19. Monsen，E.，Franke，M.，Rung，T.，Aumann，P. \& Ronzheimer，A.，用于飞机空气动力学性能预测的高级运输方程湍流模型的评估，AIAA 99-3701，1999年。 20. Kroll，N.，Rossow，C.C.，Becker，K. \& Thiele，F.，MEGAFLOW项目，将出现在航空航天科学上。 技术。 5（2000）。 21. Kroll，N.，Rossow，C.C.，Becker，K. \& Thiele，F.，MEGAFLOW数字流动模拟系统，ICAS大会，1998年9月13日至18日，澳大利亚墨尔本。 22. Pahlke，K.，Sides，J. \& Costes，M.，CHANCE 完整的直升机计算流体动力学的法国-德国研究项目，ECCOMAS 2000，2000年9月11日至14日，西班牙巴塞罗那，会议记录。 23. 埃格斯，Th。 \& Novelli，Ph.，开发双模Ramjet飞行测试车的设计研究。 DGLR-Jahrestagung，1999年9月27日至30日，德国柏林。 24. 科杜拉，W。 \& Morice，Ph.，欧洲CFD高超音速外流状态的计算模拟，AGARD-CP-600，卷。 3，C14-1-16，AGARD研讨会，法国帕拉索，1997年4月14日至17日。 25. Kordulla，W.，Kroll，N. \& Schiitz，H.，DLR的精选CFD活动。 第11届NAL飞机公司研讨会。 空气动力学，1993年6月10日，日本东京。 26. 科杜拉，W。 \& Radespiel，R.，DLR的计算高速可压缩流最近发展，DLR-FB 97-34，1997年。 27. 汉尼曼，K。 \& Schnieder，M.，HEG锥形喷嘴的校准结果，ESA SP-426，667-671，1999年。 28. Longo，J.M.A.，Orlowski，M. \& Brack，S.，关于再入车辆设计的CFD建模的考虑，ESA SP-426，35-42，1999年。 29. 吉塞，P.，海因里希，R. \& Radespiel，R.，使用Navier-Stokes求解器对提升体的动态导数的数值预测，数值流体力学笔记，卷。 72，pp. 186-193，Vieweg Verlag，不伦瑞克，1999年。 30. Herrmann，U.，Orlowski，M. \& Radespiel，R.，空气动力学对未来SCT的贡献，ECCOMAS 1998，会议记录卷。 2，649-654，1998年。 31. Radespiel, R., Briick, S., Giese, P., Hannemann, V., Longo, J.M.A., Liideke, H. \& Orlowski，M.，再入车辆周围流动的数值模拟。 第一届ONERA-DLR航空航天研讨会ONERA巴黎，1999年6月21日至24日。 32. 汉内曼，K.，雷曼，B. \& Schnieder，M.，HEG测试部分流量的综合实验和数值表征，第22届。 激波研讨会，英国伦敦帝国学院，7月18日至23日，论文0730，1999年。 33. Galle，M.，Gerhold，T. \& Evans，J.，使用DLR TAU代码对混合电网上复杂几何形状的湍流流的并行计算，第11次平行CFD会议记录，美国威廉斯堡，1999年5月23日至26日。 34. Labbe，S.G.，Perez，L.F. Fitzgerald，S.，Longo，J.M.A.，Molina，R. \& Rapuc，M.，X-38综合空气动力学和空气热力学活动，航空航天科学

选定的CFD功能在DLR 457技术3，485-493，1999年。 35. 利德克，H。 \& Krogmann，P.，层压/湍流边界层过渡的数值和实验调查，2000年ECCOMAS会议记录，巴塞罗那，2000年9月11日至14日。 36. 布拉克，S.，拉德斯皮尔，R. \& Longo，J.M.A.，简化空间-Shuttle配置的非平衡流的比较，AIAA论文97-0275，1997年。 37. Kroll，N。 \& Jain，R.K.，二维Euler 方程的解决方案-有限体积代码的经验，DFVLR-FB 87-41，1987年。 38. KorduUa，W. \& MacCormack，R.W.，使用显式隐式方法的跨声速流计算，物理学笔记，卷。 170，Springer Verlag，420-426，1982年。 39. 洪，CM。 \& KorduUa，W.，三维流场模拟的时间分割有限体积算法，AIAA杂志22，1564-1572，1984年。 40. KorduUa，W.，Vollmers，H. \& Dallmann，U.，通过半球-气缸配置分离的三维跨声速流动的模拟，AGARD-CP-412，31-1至-15，1986年。 41. Miiller，B.，隐式光束和变暖方案的矢量化，在Gentzsch，W. （编辑），应用到计算流体动力学的计算机程序的矢量化，Vieweg Verlag，172-194，1984年。 42. 里德尔包奇，S。 \& Brenner，G.，层压高超声速流动过去钝体的数值模拟，包括高温效应，AIAA论文90-1492，1990年。 43. Schwamborn，D.，使用块状结构显式Navier-Stokes方法模拟DFVLR-F5机翼实验，数值流体力学笔记，卷。 22，1988年。 44. Hilgenstock，A.，Ein Bitrag zur numericen Simulation der transschen Stromung um einen Deltafliigel durch Losung der Navier-Stokesschen Bewegungsgleichungen，DLR-FB 90-13，1990。 45. Sonar，Th.，关于可压缩流动的一般三角测量的上风方案设计，数值算法4，135-149，1993年。 46. Pahlke，K.，标准航空CFD方法FLOWer在ETR 500隧道入口中的应用，Transaero研讨会论文集，法国巴黎，1999年5月4日至5日。 47. Niederdrenk，P.，关于椭圆网格生成的控制，会议记录，第6届国际。 关于Comp中的数字网格生成的会议。 现场模拟，美国密西西比州，257-266，1998年。 48. Brodersen，O.，Hepperle，M.，Ronzheimer，A.，Rossow，C.-C. \& Schoning，B.，参数网格生成系统MegaCads，会议记录，第5届国际。 关于Num的Conf. Comp中的网格生成。 现场模拟，美国密西西比州，353-362，1996年。 49. Kroll，N.，Radespiel，R. \& Rossow，C.-C，结构网格上3D应用的准确高效的流动求解器，AGARD报告R-87，4.1-4.59，1995年。 50. 布拉克，S。 \& Radespiel，R.，将欧拉/Navier-Stokes代码CEVCATS扩展到黏性非平衡流动，DLR-IB 223-95 A64，1996年。 51. Rudnik，R.，Melber，S.，Ronzheimer，A. \& Brodersen，O.，运输机高升配置的3D RANS模拟的各个方面，AIAA论文2000-4326，2000年。 52. Honlinger，H.，KorduUa，W.，Meier，G.E.A. \& Eitelberg，G.，Simulationszentrum fur Stromung und Struktur（SISS），ein innovationsorientiertes Forschungskonzept am DLR-Standort Gottingen，会议记录，DGLR-Jahresbericht，1999年。 53. CENTAUR：http://www.centaursoft.com

\clearpage
\chapter{23 CFD对太空运输系统的应用}
\section{23.1 介绍}
CFD在太空运输系统中的应用 Kozo Fujii1 23.1 介绍 在20世纪70年代初，CFD在航空航天领域首次吸引了人们的注意。 跨声速翼型上的嵌入激波在计算机模拟中自动捕获，商用飞机的设计过程被大大简化。 在80年代中期，航空航天领域的CFD再次吸引了人们的注意。 空间运输系统的创新需要对流动环境进行CFD模拟，这在地面上很难实现。 最近，我们在CFD中没有看到如此具有划时代的话题。 我们想乐观地理解它，并认为它表明CFD已经进入成熟阶段，人们开始将其作为流体动力学研究的简单分析工具之一。 本文展示了CFD最近在空间与宇航科学研究所空间运输系统的开发和基础研究中的一些应用。 选择了三个主题；（1）空气呼吸发动机船尾区域的阻力减少，（2）截断环塞喷嘴的流场分析，以及（3）再入胶囊的动态不稳定性分析。 这些例子展示了CFD技术在空间应用中的能力和有效使用。 '空间和航天科学研究所，相模原，神奈川，229，日本前沿计算流体动力学 - 2002 编辑：David A. Caughey和Mohamed M. 哈菲兹@2002世界科学

\section{23.2 数值方法}
\section{23.3 结果和讨论}
460 FUJII 23.2 数值方法 23.2.1 基本方程 基本方程是用广义坐标系书写的三维可压缩Navier-Stokes 方程。 由于其中一个应用程序需要网格运动，因此有必要在公式中包含时间度量术语。 这些方程是及时解决的，但基本方程被修改为每个网格区之间的信息交换，因为其中两个应用程序使用超集区网格系统。 本研究中使用的接口方法称为“强化解决方案算法（FSA）”，其中将源项添加到基本方程中，作为强制项来强制执行其他区域的解决方案。 该算法的详细信息可以在参考中找到。 1. 23.2.3离散化LU-ADI算法[2]用于时间积分，由Shima和Jonouchi开发的SHUS方案离散化的对流项[3]。 SHUS方案是AUSM型方案之一，修改了溶液的单调性。 三阶精度是通过使用原始变量的MUSCL插值[4]实现的。 黏性术语被中心差分离散化。 Baldwin \& Lomax [5]的代数湍流模型被使用，因为雷诺数大约是106到107，必要时会进行一些修改。 自由流低压区自由剪切层23.3结果和DisCUSSionS膨胀风扇23.3.1空气涡轮RAMjet发动机的船尾阻力减少空气涡轮RAMjet发动机（所谓的ATREX发动机）正在空间和宇航科学研究所（ISAS）开发中，将来用作TSTO（两级轨道）运输系统的助推器[6]。 图中的当前问题之一。 1 ATREX的船尾区域，这个RAMjet发动机是船尾区域诱导的插头喷嘴阻力量。 在超音速或低超音速下，船尾阻力可能高达总推力的20\%以上。 图1显示了原理图

空间系统船尾区域。 插头喷嘴用于在整个飞行路径上实现高推力，它们在自由流的强烈伸缩和从罩唇的羽流边界之间产生低压区域。 这个低压区域导致所谓的“船尾拖曳”。 首先根据模拟讨论基本流动结构，然后尝试次要流动注入以减少阻力。 三维Navier-Stokes方程被用于Mure在at tack角度下计算物体，尽管这里的结果仅限于轴对称流。 2 计算网格（外部区域）图。 3个计算网格（内部喷嘴区域）两个网格系统；一个用于喷嘴的内部区域（10L.17\_81），另一个用于外部流动区域（347\_17\_151），用作覆盖网格。 图2和图3分别显示了这两个网格。 图中的箭头。 3是二次流动注射的指示和位置。 喷嘴和机身的几何形状取自在IS AS为ATREX发动机开发进行的实验。 设计压力比为9.86，喷嘴面积比为1.92。 图。 4 马赫数等高线图\_（4=6.5）与速度矢量图的特写视图

462 富士自由流马赫数是1.5，雷诺数字是1.37.„ 107. 腔室总压力和自由流静压的比率是4.5、5.5和6.5。 二次流注入的总流条件与内部流的条件给出。 计算的马赫数等高线图如图所示。 4对于£=6.5的压力比，没有注射。 沿着船尾区域观察到强烈的膨胀，因此，形成了低压区域，分离激波出现在喷嘴羽流前面，分离激波和羽流之间发生了再循环区域。 图5显示了沿船尾区域和插头的压力分布。 再循环区域的压力低于自由流动压力，它成为阻力的原因之一。 然而，如图所示，分离激波前面的膨胀出现了更宽的低压区域。 5，这个区域的减少是减少船尾阻力的关键。 虽然这里没有显示，但当压力比增加时，分离激波向前移动，船尾阻力会降低。 结果表明，低压区域将变小。并且阻力将因分离激波的向前运动而减少。 为了减少低压区域，尝试了二次流动注入。 图6显示了马赫数等高线图。 与图相比。 4，分离激波向前移动，位于二次喷射喷射的前面。 由于环境压力的变化，喷嘴羽流的膨胀更强。 在图中。 7，绘制了这种情况的压力分布。 低压区域向前移动到身体的平坦区域。 船尾阻力减少所有1.5 O 0.5 -0.5.... j..... j.... I >... j. - S e p -ar ati-cj-n • -sh o -dk w av b- • Stlart point of! \textbackslash\{\} '••; Bfil\textasciicircum{}aiifepdti V;" i i I j | KjJ i:; Co; 插头表面 \textbackslash\{\}] i j \textbackslash\{\}i! F \textbackslash\{\} / " t •! I I: wlli'p;; -1 -0.5 0 0.5 1 1.5 2 2.5 非尺寸tl距离图。 5 船尾区域和插头表面的压力分布图。 6个马赫数轮廓图（5=6.5）-带二次注射

空间系统463通过注射三种不同的压力比，尽管这里没有显示。 船尾阻力的比较表明，二次喷射在任何压力比下都有效地降低了船尾阻力，尽管这里没有显示。 综上所述，澄清了船尾区域的流动特征，并确定了船尾拖曳的原因。 还证明了二次流动注入减少阻力的有效性。 23.3.2 空心截断塞喷嘴的流场分析-Cp（无二次流动）-I'-p i'-iv-.th其次-lov*i压力比6.5 o °5 Blowpointp插头表面起点•船尾区域-1 -0.5 0 0.5 1 1.5 2 2.5非维距离图。 7 沿着船尾区域和塞子表面的压力分布-注入单级轨道（SSTO）车辆是预计下个世纪的两个发射系统之一。 高效的推进系统是实现SSTO的关键因素之一。 插头喷嘴似乎是系统的一个很好的候选者，并且在过去几年里已经进行了研究。 与钟形喷嘴不同，喷管流动不受墙的限制，相反，耗尽的喷射受外部流动的限制。 插头喷嘴被认为具有全球更好的性能，因为喷射边界根据环境压力调整其配置，因此喷射器可以针对整个条件（例如高度）进行最佳地进行曝气。 此外，喷嘴性能不受喷嘴截断（切断后部）的影响。 插头喷嘴的这些有利特征被以前的研究所限制，但由于实验和/或计算研究数量有限，关于该机制的讨论不足。 在本研究中，对插头喷管流动场进行了计算模拟，以调查上述特征。 两个参数；图中的腔室。 8 计算马赫数分配（轮廓插头喷嘴，压力比71.0-最佳）

464 富士的环境压力比和喷嘴的长度被系统地改变。 主要研究轮廓塞子喷嘴的流动结构，并对推力性能进行定量评估，并讨论以下塞子喷嘴的流动机制。 半角25度的环形锥形塞用于验证当前计算[7]。 虽然由于长度限制而没有显示在这里，但计算出的压力分布与Tomita等人的实验非常一致。[8]对于所有压力比，结果表明，目前的方法适合分析。a 200 1 15° I too a \& so v60 '63。 \textasciicircum{} '6。 < Z65.9' V7/ D坡道 • 底座 1J 内喷嘴 20 30 40 SO 100 插头喷嘴长度（\%）图。 9\_每个截断喷嘴的压力比为500的推力部分。接下来，对采用特征方法设计的轮廓插头喷嘴进行性能的详细分析。 截断喷嘴的比较适用于五种类型的插头喷嘴，这些喷嘴保留了20、30、40、50和100\%的总插头长度。 在压力比71.0下实现最佳的放大。 根据插头轴和罩唇之间的长度，雷诺数被设置为5xl05。 其中一个结果如图所示。 8. 对于这个20\%截断（80\%截断）的喷嘴，流量在底部边缘分离并膨胀。 基部区域边缘产生的剪切层在收敛并使流动平行于轴线时诱导拖曳激波。 图9显示了插头喷嘴的每个组件在500的压力比下产生的推力部分。 插头截断导致喷嘴长度变短。 喷嘴的坡道面积减少，坡道产生的压力推力也因截断而减少。 另一方面，基地压力产生的推力因基地面积的增加而增加。 它补偿了坡道压力推力降低造成的总推力损失。 因此，对于任何喷嘴截断，总喷嘴推力几乎相同。o.i.. o.oi o.ooi \textbackslash\{\}|! 设计< • • 20\% m\% 4(1\% 50-* | 1 喷嘴喷嘴喷嘴 NtJiJZh: 10 100 1000 压力比图。 10 平均基础压力.vs.压力比（无外部流量）

SPACE SYSTEMS 465 *ivelop shock trfiliM\textasciicircum{}hock 图10显示了20、30的平均基压？ 与各种压力比相比，绘制了40\%和50\%的插头喷嘴。 破折线显示环境压力，随着压力比变高（高海拔），环境压力会降低。 对于低压比（低海拔），基压随着环境精神压力的降低而线性下降。 在这个区域，喷射不会在尾部收敛，基部区域向外部环境开放。 因此，基地区域受到外部环境的影响，基地的压力在数量上等于环境压力。 施加的推力是E„ 1t D.,. /, ir\textasciicircum{}.u • \_. • J +u u • 钻机。 11 压力等高线图（上半部分）环境和基部在低压比下是相应的亚音速区域。 （压力比：35.0）另一方面，在高压比下，尽管压力比变化，但模具底座上的压力还是恒定的。 在高压比区域，射流在轴的下游扩张和收敛。 再循环区域在底座后面形成，底座压力独立于外部流动条件。 随着高度（压力比）的升高，环境压力降低，基底压力和环境压力之间的差异增加。 因此，基础压力推力增加。 由于基部区域，这种推力增加使插头喷嘴在整个高度上具有高效率。 事实证明，当压力比低时，基础压力几乎等于环境压力，当压力比高时，基础压力独立于环境精神压力。 解释中使用的原因是，唤醒区域对低压比是开放的，对高压比是封闭的，再循环区域是为高压比创造的；这是过去的普遍信念。 然而，一些计算表明，即使尾部关闭，基压也不独立于外部压力。 突然改变基压特征的机制是什么？ 可以找到该机制的两个参考[9,10]。a A • 20\%喷嘴*2（J\%喷嘴<M==压力比图。 12 平均基压.vs.压力比（带外部流量）

466 FUJII 2.4 2.2 vo.o；两篇论文都表明，从罩唇发出的压缩波/包层冲击起着重要作用。 前者根据尾部冲击发电处的尾迹停滞位置定义标准，而后者根据包面冲击冲击的位置定义标准。 为了讨论该现象的详细机制，进行了进一步的计算。 正如我们之前研究的研究表明，除了要加强的包络冲击外，外部流动不会改变基本流动机制，因此讨论了与外部流动的模拟流动机制。 主要讨论了20\%截断喷嘴的结果。 作为计算流场的一个例子，图中绘制了计算的压力轮廓以及显示亚音速区域的马赫数等高线。 压力比为35.0的11，其中显示了由于剪切层碰撞和包层冲击冲击导致的两个高压区域。 0.8压力比“图图。 13 驻点在尾迹区域中心线的位置 首先，对于图中有外部流动的情况，重新计算了基压与压力比的关系。 12. 流量特征在压力比32.5时发生变化。 在压力比30.0时，基压受到外部流动的影响，但值不超过32.5。 在压力比30.0时，尾随的激波在基区的亚音速部分。 在压力比为50.0时，尾随的激波撞上了基区的超音部分。 在压力比35.0下，尾随的激波撞上亚音速区域（见图。 11），但根据图，基压与外部流动无关。 12. 冲击是否冲击基区的亚音速或超音速部分不一定对应于基压的特征变化[ll]。 在研究压力比32.5附近的流场细节之前，根据图中的压力比绘制了尾迹中心线的停滞点的位置。 13. 停滞位置在压力比32.5时改变其位置，图中的值。 14 计算压力等高线图和沿中心线的压力分布（压力比30.0）

空间系统467的强烈变化发生在基本压力上。 这表明驻点的突然移动可能是基本压力突然变化的关键。 沿着j wake中心线的压力分布如图所示。 14对于压力比30.0。 尾部有两个高压区域。 第一个峰值是由于剪切层碰撞造成的，第二个峰值是由于包膜冲击冲击造成的。 \_第二个峰值要高得多。 流量停滞不前，Fi 15 计算压力轮廓和强烈的反向流量。 重排量和压力分布适合更高的压力比（35.0）都是图中绘制的中心线。 15. 前峰现在（压力比35-0）高于第二个峰。 另一个压力峰值出现在前停的正向流和第二停的后流相互撞。 更高的压力比的结果表明，前停压峰值比下游峰值高得多。 靠近基地的流动永远不会感受到下游的压力，因为由于剪切层碰撞产生的高压区域的存在，下游的流动不会进入基地区域。 这里的研究清楚地表明，由剪切层和尾迹中的包膜激波诱导的两个高压区域支配着流动机制。 基压是否受到环境压力的影响取决于由包膜激波和剪切层碰撞引起的停滞区域的反向流动。 分析表明，基压特性不是受压力波（声速）的支配，而是受尾迹区域的对流速度现象的支配[l 1]。 23.3.3 再入间隙的动态不稳定性分析 钝再入胶囊在跨声速下往往具有动态不稳定。 不稳定现象自20世纪60年代以来就为人所知，主要通过实验进行研究。 然而，直到最近才讨论不稳定的机制。 太空与宇航科学研究所有一个样本返回项目，其中装有小行星样本的太空舱经历了高速返回。 减速舱的降落伞将在低速下打开，因此舱在超音速下面临这种不稳定。 最近，ISAS的Hiraki通过实验研究了流场，并表明不稳定是由基压的相位延迟引起的。 以下com-

468 富士！ 当前作者的假定研究也支持这一结果[12]。 本文总结了关于该主题的一系列计算研究[12-14]。 在研究过程中，开发了一种新的后处理方法，这种方法对于很好地理解脚趾流结构来说是不可避免的。 图。 16 ScMieren振荡囊上流场的图像（M=l.3，间距=-20至20度）首先模拟了强制俯仰振荡中囊上方的流场。 最大间距为±20度，振荡频率为20赫兹，这是HiraM等人在风洞实验中观察到的值。 最大直径为100毫米，身体长度为50毫米。 相应的雷诺数是2.5 x 106，用于计算的自由流马赫数是1.3。 研究中使用的网格点总数约为100万。 图16显示了强制俯仰振荡中胶囊上方流场的结构（schlieren图像）。 出现弓激波、剪切层、再压缩激波，这些使流场变得复杂。 正如Hirakfs实验显示的那样，基压存在时间延迟。 因此，绘制了前部和底部代表点压力的时间历史，以确认实验观测。 结果表明，前部没有压力，但基压有大约3毫秒（15度）的时间延迟。 实验中观察到的俯仰振荡的频率为20赫兹，降低的频率为10英寸2。 因此，尽管我们必须考虑基压的时间流逝，但我们可能会认为流场是准稳定的。 图。 17 每个俯仰角度的基本俯仰矩（固定俯仰的结果）

空间系统469 a=10。 Qdeg。 因此，详细研究了固定俯仰角度下胶囊上的流场。 仅由基压分布产生的俯仰矩如图所示。 几个音高角度为17。 在这个情节中观察到了两个重要的特征。 首先，俯仰矩对于负音高角度是正的，对于正的音高天使是负的。 基区产生的投球矩有助于静态稳定性。 其次，在高俯仰角（绝对值）下，俯仰矩几乎恒定，俯仰矩在接近零俯仰角度时突然改变其符号。 结果表明，除了接近零度外，大多数俯仰角的基流场是相似的。 图中绘制的瞬时蒸汽线。 18清楚地证明了这一点。 出现了一对漩涡；底部有一个大漩涡，上面有一个小\textasciitilde{}反向旋转漩涡。 图。 18 固定音高胶囊的瞬时流线型，胶囊现在在几个圆度的音高振荡中的俯仰时刻如图所示。 19. 还绘制了对应于「常量放线模型5」的曲线。 “恒定延迟模型”意味着俯仰振荡中胶囊的流场被解释为准稳态解，基压的时间延迟。 在“恒定延迟模型”中，（基本）俯仰矩会发生一定的时间延迟，迟滞出现在图的周期中。 图19表明，尽管观察到一些数据散射，但“恒定延迟模型”是可以接受的。 在俯仰振荡中，胶囊的基流场中存在一些毫秒延迟，我们可以根据固定俯仰simulac\textasciicircum{}\textasciicircum{}s\textasciicircum{}\textasciicircum{}X\textasciicircum{}-\textasciicircum{}N \textbackslash\{\} t\textbackslash\{\} - 振荡““”固定r.»tch 3T固定间距4 5r St攻击角度图讨论流动结构。 俯仰振荡中胶囊的19个基本俯仰矩-迟滞循环

470 富士！ 考虑到一定的时间延迟。 图中绘制了从下到上沿中心线的时间平均基压分布。 20用于几个固定俯仰角度的模拟结果。 随着俯仰角的增加，压力水平会降低。 对于所有俯仰角，底座的顶部（图中的右侧）都有一个相对高压区域，随着俯仰角的增加，该区域与其他较低区域之间的差异变大。 虽然这里没有显示，但这个高压区域是由回循环区域的强逆流对流产生的（见图。 18）。 还应该注意的是，高压区域只出现在具有动态不稳定倾向的囊中（见参考 13是详细信息）。 由于胶囊的基部压力在俯仰运动中存在时间延迟，在固定俯仰角度下，胶囊底部的较高和下部的压力差比胶囊更小。 这将有助于动态稳定性，因为高压区域不会立即出现在运动中的俯仰角度变化中。 时间平均基压（D45 IVbdel）为了确认这个高压区域存在于俯仰运动中的胶囊的流场中，并且基流场中存在时间延迟，为俯仰运动中的胶囊绘制了时间准确的粒子痕迹。 然而，流场没有显示清晰的结构。 原因是从胶囊边缘流出的涡流的振幅与音高运动导致的流动结构变化的振幅几乎相同，它们不能与动画分开。 定义了数字过滤器的新概念。 对于每个网格点的流动变量的时间序列，运行频率滤波器。 在这种情况下，操作截止频率为150赫兹的低通滤波器（LPF），以消除高频流量结构。 结果清楚地表明，高压区域存在于俯仰运动中胶囊的流场中，并且基流场中存在时间延迟。 在一定的音高角度延迟后，流动结构会发生变化。 图21显示了粒子痕迹的即时镜头。 流程变化的详细信息可以在参考中找到。 14. 这种过滤技术对许多类型的流场都很有用。 该方法在参考中进行了讨论。 15和其他例子。 \_.在D上，...，.... r 图。 基础固定间距垂直中心线上的20个压力分布

\section{23.4 结论}
空间系统471（a）原始原始数据（b）过滤数据JFig。 21 基区的时间准确粒子痕迹' Sawada [16]开发的涡流识别技术用于澄清尾迹区域的涡流结构。 结果如图所示。 22 图中出现的两个反向旋转涡流。 18个不是相互隔离的，但它们是一个组合的涡旋环，如图所示。 下部（当正距角度）位于前方，这个强环在基部附近形成大的再循环区域（见图。 18）。 戒指的另一部分相当薄弱，远离底座。 由于涡流两侧的差异，下游出现了一对两条纵向涡流街道。 这种纵向涡旋与翼翼尖涡相似。 翼尖涡的强度反映了机翼的静态升力系数，如果胶囊也是如此，胶囊的静态升力与一对纵向涡流的形成直接相连。 由于基地上部高压区域的时间延迟，发生了动态不稳定。 这个高压区域是由于胶囊后面的涡流结构而发生的，涡流结构与一对纵向涡流有关。 基于对流动机制的这种理解，我们可能会根据静态升力特性讨论这种类型的胶囊的动态不稳定性，因为纵向涡流的强度是由升力强度决定的（就像翼尖涡一样）。 讨论才刚刚开始，尚未得出结论，但一些结果表明了这一假设的积极原因。 23.4 • 结论CFD最近在空间和宇航科学研究所对空间运输系统的开发和基本图> 22 涡核分布研究的一些应用。 选择了三个主题；（1）减少空气呼吸的船尾区域的阻力-

\section{23.5 确认}
\section{参考资料}
472 富士电子酒，（2）截断环塞喷嘴的流场分析，以及（3）再入胶囊动态不稳定性分析。 澄清了流场和船尾阻力特性，并演示了二次流注入的效果。 截断塞式喷嘴的流量分析显示了塞式喷嘴的特性，并显示了这些特性背后的关键流量特征。 对再入囊动态不稳定性的分析表明，基压的延迟时间是不稳定的关键，并澄清了流动机制。 这些例子展示了CFD技术在空间应用中的最新能力。 23.5 致谢 这里介绍的研究是作者在空间和宇航科学研究所实验室的合作成果。 航空航天系研究生Kazuhiro Imai（现为MHI）对船尾阻力的减少进行了研究。 东京大学。 对RLV的环形截断塞喷嘴的研究主要由Takashi Ito（现任航空航天系研究生）进行。 东京大学），青山学院大学的研究生。 对胶囊动态稳定性的研究主要由实验室的研究助理Susumu Teramoto进行。 参考资料1。 Fujii，K.，基于强化解算法的统一区域方法，计算物理学杂志，卷。 118，pp. 92-108，1995年。 2. Obayashi，S.，Matsushima，K.，Fujii，K.和Kuwahara，K.，使用LU-ADI分解算法提高Navier-Stokes计算的效率和可靠性，AIAA论文86-0338，1986年。 3. Shima E.和Jounouchi，T.，CFD在航空工程中的作用/No. 14 - AUSM类型上风计划 -，Proc。 第14届NAL飞机压缩空气动力学研讨会，pp. 7-12，1997年。 4. 托马斯，J。 L.，van Leer，B.和Walters，B. W.，Euler 方程的隐式通量分割方案，AIAA论文85-1680，1985年。 5. 鲍德温，B。 S.和Lomax，H.，分离湍流流的层近似和代数模型，AIAA论文78-257，1978年。 6. Tanatsugu，N.，ATREX引擎的开发研究，第七届国际太空飞机和高超音速系统与技术会议，弗吉尼亚州诺福克，1996年。 7. Ito T.、Fujii K.和Hayashi A. K.，轴对称塞喷管流动场的计算，AIAA论文99-3211，1999年。 8. Tomita T.，Tamura H.和Takahashi M.，Plug 喷管流动场的实验评估，AIAA论文96-2632，1996年。 9. Hagemann G.，Immich H.和Trehardt M.，先进火箭喷嘴中的流动现象-塞子喷嘴，AIAA Paper 98-3522，1998年。 10. 鲁夫，J。 H.和McConaughey，P. K.，三维的数值分析-

SPACE SYSTEMS 473 sional Aerospike，AIAA Paper 97-3217，1997年。 11. Fujii，K，Ito，T.和Hayashi，K.，环形截断塞喷嘴的流场分析，国际空间技术与科学研讨会，ISTS论文2000-e-10，2000年5月。 12 Teramoto, S., Hiraki, K. and Fujii, K., Transonic速度下再入舱动态稳定性的数值分析，AIAA大气飞行力学会议，AIAA-98-4451，马萨诸塞州波士顿，1998年8月10日至12日。 13 Teramoto，S.和Fujii，K.，Transonic速度下钝胶囊背后的流场的计算研究，第30届AIAA流体动力学会议，弗吉尼亚州诺福克，AIAApaper 99-3414，1999年6月。 14 Teramoto，S.和Fujii，K.，Transonic速度下再入胶囊的不稳定机制研究，流体2000，AIAA论文2000-2597，2000年6月。 15 Teramoto，S.和Fujii，K.，识别复杂非稳流流动机制的可视化技术，Proc。 2000年7月，京都，第一届计算流体动力学国际会议（即将出席）。 16 Sawada K.，识别涡流核心的便捷可视化方法，日本航空航天科学学会交易，卷。 38，不。 120，pp. 102-120，1995年。

\clearpage
\chapter{使用进化算法对超音速机翼进行24个多点最佳设计}
\section{24.1 介绍}
使用进化算法对超音速机翼进行多点最佳设计 Shigeru Obayashi、Yukihiro Takeguchi和Daisuke Sasaki1 24.1 介绍 开发新型超音速运输（SST）的需求预计将增加，因为人们更频繁地出国旅行。 在过去的十年里，美国、欧洲和日本为开发新的SST做出了几次努力。 在日本，国家航空航天实验室（NAL）正在研究和设计SST，以在2002年发射小型超音速实验飞机1。 考虑到新的SST设计，存在许多技术困难需要克服。 L/D必须改进，声爆应该被阻止。 然而，在降低阻力和繁荣之间存在着严重的权衡。 因此，新的SST预计仅在海上以超音速巡航，在陆地上空以超音速巡航。 这意味著重要的设计目标不仅是提高超音速巡航性能，也是提高跨声速巡航性能。 例如，大扫角可以降低波阻，但由于结构问题，它限制了允许的长宽比。 因此，在设计SST时，有许多权衡需要解决。 为了确定全球权衡，这个问题可以被视为多目标（MO）优化。 MO优化旨在优化向量值目标函数的分量。 一般来说，这个问题的解决方案不是与单目标优化不同的单点，而是被称为帕累托最优集的点家族。 帕累托解决方案是帕累托最优集的成员，代表了多个目标之间的权衡。 多目标遗传算法（MOGA）是高效和有效地采样多个帕累托解的独特优化方法。2 GA是模仿自然进化的优化方法。 由于GA使用大量设计候选者并行寻求最佳解决方案，MOGA可以在同一日本仙台东北大学航空航天工程系980-8579确定多个帕累托解决方案。 Frontiers of 计算流体动力学 - 2002 编辑：David A. Caughey和Mohamed M. 哈菲兹 ©2002 世界科学

\section{24.2 最佳化方法}
476 OBAYASHI等人。没有指定目标之间的权重的时间。 本文考虑了超音速和跨声速巡航条件下SST机翼形状的多点空气动力学优化。 在升力约束下，空气动力学阻力将在两种条件下最小化。 根部的弯曲矩也将最小化，以防止所有帕累托解决方案具有不切实际的大长宽比。 在空气动力学优化中，设计变量指定平面形状、弧度、厚度分布和扭曲分布。 在之前的研究中，在无黏流动假设下使用潜在求解器和欧拉求解器执行了相同的MO优化。3为了考虑更现实的流场，如可能的流动分离，在当前优化中考虑了黏性效应。 因此，Navier-Stokes求解器用于评估两种巡航条件下的机翼性能。 最后，对生成的帕累托解决方案进行了分析，并与NAL的设计和之前的结果进行了比较。 24.2 优化方法 将GA应用于MO优化有很多优点。 它们的优势源于算法本身，它模仿了自然进化的机制，生物种群世代相传，通过选择、交叉和突变来适应环境。 在设计优化问题中，健身、个体和基因分别对应于目标函数、设计候选和设计变量。 GAs同时和随机地从设计空间的多点搜索，而不是像基于梯度的方法那样确定性地从单点移动。 此功能可防止设计候选者在本地最佳位置定居。 此外，GA不需要计算目标函数的梯度。 这些特征导致了GA的以下三个优势：1，GAs有能力找到全球最佳解决方案。 2，GA可以并行处理。 3，高保真CFD代码无需任何修改即可轻松适应GA，因为GA仅使用目标函数值。 GA已经扩展，以成功解决MO问题。2 GA使用总体并行寻求最佳解决方案。 此功能可以扩展为并行寻求帕累托解，而无需指定目标函数之间的权重。 由此产生的帕累托解决方案代表了全球权衡。 因此，MOGA是解决MO问题的非常独特和有吸引力的方法。 图1显示了本研究中MOGA的流程图。 以下是这里雇用的遗传操作员的简要描述。 传统上，GA使用二进制数来表示设计参数值。 然而，对于像目前的空气动力学优化这样的实函数优化，使用实数更直接。 因此，这里使用了浮点表示。 选择基于帕累托排名方法和健身共享。2 每个人根据占主导地位的个人数量被分配到其排名。 标准的健身共享功能用于保持人口的多样性。 为了更有效地找到极端的帕累托解，所谓的最佳N选择4也是

\section{24.3 当前优化问题的制定}
多点最佳设计477与。 采用下面描述的混合交叉5（BLX-a）。 此运算符在由两个父端和用户指定的参数a定义的段上生成子。 在这种优化中，新设计变量的加权平均值被使用为Childl = y\_Parentl + (l-y)\_Parent2 Child2 = (l-y)\_Parentl + y\_Parent2 (-11') y=(l +2a)\_ranl -a，其中Childl,2和Parentl,2分别表示孩子（新人口的成员）和父母（老一代的交配夫妇）的编码设计变量。 这里显示的随机数rani在[0,1]中是均匀的。 参数a设置为0.5，六个平面设计变量除外。 由于平面对空气动力学性能有很大影响，因此其设计参数必须保守地给出。 否则，计算会发散，许多孩子无法被评估。 因此，对于这六个设计变数，引数a被设定为0.0。 突变的概率为20\%。 如果发生突变，那么Eqs。 （1）将替换为Childl = y\_Parentl + (l-y)\_Parent2 + m\_(ran2-0.5) \textasciicircum{} \textasciicircum{} Child2 = (l-y)\_Parentl + y\_Parent2 + m\_\{ran2-0.5)，其中ran2也是[0,1]中的均匀数，m设置为每个设计变量给定范围的10\%。 使用CFD求解器对每个人的目标功能进行评估。 为了评估黏性效应，应求解三维Navier-Stokes 方程。 在这项研究中，三维、可压缩、薄层Navier-Stokes求解器用于评估跨声速和超音速巡航条件下的空气动力学性能。 这个Navier-Stokes代码采用总变异减小类型迎风差分和下上因数对称高斯-塞德尔方案。 多重网格法7也用于加速收敛。 本代码中的湍流模型采用Baldwin和Lomax的代数混合长度模型。8 24.3 当前最佳化问题的表述 目前的最佳化问题可以表述如下。 [目标功能] 1. 阻力系数用于跨声速巡航，CD，t 2。 用于超音速巡航的阻力系数，CD，S 3。 超音速巡航条件的机翼根处的弯曲力矩，Mroot [约束] 1。 巡航条件下的升力系数，CL,I = 0.15和CL,S = 0.10 2。 机翼区域，S = 60 3。 最大翼型厚度，t/c > 0.03

\section{24.4 超音速运输翼的优化}
478 OBAYASHI等人。 [流量条件] 1. 跨声速巡航马赫数：0.9 2。 超音速巡航马赫数：2.0 3。 雷诺数基于两种条件下的根弦长度：1.0 x 107 在当前的优化中，所有三个目标函数都要最小化。 使用Navier-Stokes 流动求解器来评估超音速和超音速阻力系数。 弯曲矩是通过在超音速巡航条件下直接积分压力载荷来计算的。 为了保持升力系数恒定，使用从有限差中得到的CLa来预测攻角。 因此，每次评估都执行三次Navier-Stokes计算。 在空气动力学优化过程中，机翼面积冻结在恒定值上。 机翼厚度也受到结构强度的限制。 设计变量分为平面、翼型形状和机翼扭曲。 如图所示，机翼平面由六个设计变量决定。 2和它们的范围写在表1中。 由于固定的翼面积，翼尖的弦长度是相应确定的。 翼型形状由其厚度分布和弧线组成。 如图所示，厚度分布由九个多边形9定义的贝塞尔曲线表示。 3. 表1还显示了他们的设计范围。 厚度分布在翼根、折点和尖端定义，然后沿跨度方向线性插值。 整个厚度分布的多边形总数为27个。 由翼型弧线组成的弧形表面在机翼的内侧和外侧分别定义。 每个曲面由贝齐尔曲面表示，贝齐尔曲面由弦方向的四个多边形和跨度方向的三个多边形定义。 例如，图4显示了弧线及其根部的控制点。 它仅在根部是凹的，在其他跨度位置变得凸，类似于基于线性化理论的经线设计。 定义两个弧形曲面的多边形数量总共为20个多边形。 最后，如图所示，机翼扭曲由带有六个多边形的B-花式曲线表示。 5. 因此，66个设计变量被用来定义机翼形状。 目前的优化是在物理化学研究所的富士通VPP700E超级计算机系统上进行的。 该系统有160个PE，384 GFLOPS和320 GB。 主PE管理MOGA，而从属PE计算Navier-Stokes代码。 人口规模被设定为64，因此该过程根据可用性与8-64 PE并行化。 应该注意的是，平行化几乎是100\%，因为Navier-Stokes的计算主导了CPU时间。 24.4 超音速运输机翼的优化 24.4.1 帕累托解决方案概述 进化计算了30代。 在MOGA的当前优化之后，所有评估的解都再次排序，以尽可能多地找到帕累托解。 结果，最终的帕累托解在三个-

多点最佳设计479维目标函数空间如图所示。 6. 图中显示了具有目标函数的权衡面。 它还显示了四个典型的平面形状；Co,t最小值、CD,S最小值、弯曲力矩最小值和某个帕累托解决方案。 极端的帕累托解决方案，三个最小化各自目标功能的平面形状在物理上看起来很合理。 为了更清楚地呈现目标之间的权衡，将帕累托解投射到二维平面上，如图所示。 7-9。 图7和图8显示了跨声速和超音速阻力系数之间的权衡。 解决方案由长宽比标记，锥度比使用图中的不同符号进行标记。 分别是7和8。 在图中。 7，长宽比较大的机翼在空气动力学理论中实现了较低的阻力系数。 图8显示，锥度比小于0.4的机翼具有良好的空气动力学性能，但锥度比的进一步降低并不直接对应于巡航阻力的降低。 另一方面，如图所示，锥度比大于0.4的机翼具有较低的弯曲矩和较差的空气动力学性能。 9. 24.4.2 与NAL的第二个设计的比较 为了检查当前帕累托解决方案的质量，将两个帕累托解决方案与NAL的第二个设计进行了比较。 NAL SST设计团队已经完成了将于2002年发射的实验性超音速飞机的第四次空气动力学设计。 为了简要总结他们的设计概念，第一个设计确定了99个候选者中的平面形状，然后第二个设计是通过基于线性化理论的经率优化执行的。 第三个设计通过使用Navier-Stokes代码的逆法瞄准了自然层流（NLF）机翼。 最后，第四个设计是针对机翼-机组配置进行的。 由于在目前的Navier-Stokes计算中假设了完全开发的湍流，因此将目前的帕累托解决方案与NAL的NLF机翼设计进行比较是不恰当的。 因此，选择NAL第二设计进行比较。 表2总结了两个帕累托解决方案与NAL第二个设计的比较。 NAL第二种设计的空气动力学计算在这里使用相同的Navier-Stokes求解器进行。 这里介绍的帕累托解决方案A和B在所有三个目标上都优于NAL的第二个设计。 图10显示了这三个的机翼平面，表明当前解决方案和NAL的设计之间的平面形状有很大差异。 目前的平面形式类似于“箭翼”平面形式，NAL的平面形式类似于传统的“三角翼”平面形式。 三个机翼的厚度分布如图所示。 11. 帕累托解A和B的厚度分布趋势非常相似，有钝的前缘和薄的后缘。 NAL设计的厚度分布只是取自现有的NLF 翼型。 相比之下，目前的优化是在厚度受限的完全湍流下进行的。 因此，最大厚度出现在前缘附近。 然后，厚度向后缘减少，以防止边界层的快速增长。

\section{24.5 结论}
480 OBAYASHI等人。 24.4.3 黏性和不黏性计算之间的差异 目前的黏性设计与之前计算的黏性设计进行了比较。3 通过比较两个优化结果，黏性效应导致的机翼形状差异变得清晰。 在这两种情况下，帕累托解决方案在所有三个目标上都优于NAL的设计，被选中进行比较。 平面形状的比较如图所示。 12. 两种平面形状都与“箭头翼”平面相似，但形状略有不同。 与非黏性流动下的优化机翼相比，目前的机翼的扫角和锥度比也更低。 高度扫荡的机翼往往在机翼尖端附近有一个流动分离。 目前的黏性设计看起来比非黏性设计更好，可以防止尖端分离。 图13显示了根部厚度分布的比较。 它显示了截然不同的分布。 在黏性情况下，机翼在前缘附近较厚，后缘附近较薄。 然而，在非黏性的情况下，翅膀非常厚。 图中所示的Cp分布。 14 清楚地解释它们的区别。 在不黏性的情况下，Cp分布在后缘有不连续性，因此即使在后缘也能产生升力。 然而，如此厚的翼型可能会导致流动分离。 另一方面，黏性流动外壳的后缘没有不连续性。 考虑设计厚度分布的黏性效应很重要。 24.5 结论 SST机翼的多点设计优化是使用MOGA执行的。 三个客观函数用于最小化超音速阻力、跨声速阻力和翼根的弯矩。 完整的机翼形状由总共66个设计变量表示。 Navier-Stokes求解器用于评估这些空气动力学阻力。 获得成功的优化结果。 极端帕累托解决方案的方案形式在物理上似乎很合理。 提出了目标之间的全球权衡。 与NAL的第二个设计相比，两个帕累托解决方案在所有三个目标功能上都具有更好的性能。 还将目前的帕累托解决方案与之前在无黏流动下设计的最佳机翼进行比较，以检查黏性效应。 发现黏性效应对厚度分布有很大影响。 发现目前的结果可以更好地防止可能的边界层分离。 对帕累托解决方案的分析表明，理想的平面形状是一种新型的箭头翼，其锥度比相对较大，长宽比相对较小，类似于之前的非黏性结果。 致谢 这项研究部分由日本政府的赠款资助

\section{参考资料}
多点最佳设计481科学研究，编号 10305071。 第二位作者的研究得到了加拿大多伦多庞巴迪航空航天公司的部分支持。 计算时间由物理和化学研究所提供。 作者要感谢国家航空航天实验室的SST设计团队提供了许多有用的数据。 参考资料[1]Iwamiya，T.，NAL SST项目和实验飞机的空气动力学设计，第四届ECCOMAS计算流体动力学会议，2，John Wiley \&Sons，1998年，pp. 580-585。 [2]丰塞卡，C。 M。 和弗莱明，P. J.，多目标优化的遗传算法：制定、讨论和概括，第五届遗传算法国际会议，Morgan Kaufmann出版社，1993年，pp. 416-423。 [3] Obayashi, S., Sasaki, D., Takeguchi Y. \& Hirose，N。 超音速翼外形优化的多目标进化计算，IEEE进化计算交易，印刷版。 [4]小林，S.，高桥S. \& Takeguchi，Y.，MOGA的利基和精英模型，并行解决问题的形式自然-PPSN V，计算机科学讲义，Springer，1998年，pp. 260-269。 [5] 埃谢尔曼 L. J。 \& Schaffer，J. D.，实码遗传算法和间隔模式，遗传算法基础2，Morgan Kaufmann出版社，1993年，pp. 187-202。 [6]大林S. \& Grurswamy，G。 P.，用于高效静态空气弹性计算的Navier-Stokes求解器的收敛加速，AIAA杂志，33，1995年，pp. 1134-1141。 [7]詹姆森·A。 和Caughey，D。 A.，人工扩散方案对多重网格收敛的影响，AIAA论文77-635，1977年。 [8] Baldwin a B. S。 \& Lomax，H.，分离湍流流的薄层近似和代数模型，AIAA论文78-257，1978年。 [9] Grenon，R.，应用于超音速运输飞机的空气动力学设计中的数值优化，国际。 CFD超音速运输设计研讨会，国家航空航天实验室，1998年，pp. 83-104。

482 OBAYASHIETAL。 代际交替i i 初始化 i r 评估选择交叉突变 终止 图1：GAs根的流程图 图2：翼平面定义 平面形状 弦长度 跨度长度 扫角（度）*-TOOt \textasciicircum{}kink 10-20 3-15 2-7 2-7 35-70 35-70 厚度分布 最大厚度（\%） 最大厚度位置（\%） Zp4 Xp4 3-4 15-70 表1：平面和厚度分布设计变量的范围

多点最佳设计 483 0.02 t/c 尖端 0.015 0.01 0.005 P; P4 P3 P4 P5 P6 P7 P -0.2 0 0.2 0.4X/C0.6 0.8 1 1.2 图3：机翼厚度定义 -0.2 0 0.2 0.4 x/(0.6 0.8 1 1 图4：翼弧度定义T.»扭曲10角度（度）j（-0.2 0 0.2 0.42y/\#.6 0.8 1图5：翼扭曲定义

484 奥巴亚希塔尔。 帕累托解弯曲矩最小值B.013.017 CD（超音速）CD（超音速）最小值0.008 0.009 0.016 0.015 0.014 0.0il;' 013CD（超音速）CD（超音速）最小图6：目标函数空间和典型平面形状中的帕累托前

多点最佳设计 485 0.02 0.018 0.016 S 0.014 0.012 0.01 0.008！！！ I i •:.» " ° •;\% o: n „ Hr v h •*\&: t • i <: - = q n „ i W V nNt! “NAL2nd i i;:;: • o "o \textasciicircum{}rTpi ° 0.0-1.0 " 1.0-1.5 ' 1.5-2.0 * 2.0-2.5 * 2.5-3.0 3.0-3.5 • 3.5-4.0 • 4.0-4.5 0.008 0.009 0.01 0.011 0.012 C (跨声速) 0.013 0.014 0.015 图7：帕累托前投影到超音速和跨声速阻力权衡，根据长宽比标记。 NAL的设计在这里进行比较，尽管它不是帕累托最优。 0.02 0.01 0.016 g 0.014 0.012 0.01 0.008 0.008 0.009 0.01 0.011 0.012 0.013 0.014 0.015 C（超音速）图8：根据锥度标记的超音速和超音速阻力权衡的帕累托前投影

486 OBAYASHIETAL. r i? ».... J ft：*;\textasciicircum{}fy-。 B *: * », -i -Va 1 "VU!: H* ° 0.0-0.1 c 0.1-0.2 " 0.2-0.3-0.4 • 0.4-0.5 d.5-0 6 • 0.6-0.7 • 0.7-0.8 "."v 0.008 0.01 0.012 0.014 0.016 C (超音速) 0.018 0.02 图9：帕累托前向弯曲力矩和超音速阻力权衡的投影，根据锥度比标记 AL2nd A B *»r- - -../ *y i/ \&\textasciicircum{}\textasciicircum{} /* /* 4' /*' 图10：选定帕累托解决方案和NAL设计之间的平面形状的比较 A B NAL2nd NAL第2长宽比 2.19 2.20 锥度比 0.12 0.11 0.20 CD（跨声速）（xlO"4）100.40 100.96 100.99 CD（超 18.52 表2：选定的帕累托解决方案和NAL设计的性能比较

多点最佳设计 487 0.2 0.4 x/c 0.6 O.i c）66\%跨度位置0 0.2 0.4 x/c 0.6 b）33\%跨度位置0.8 0.2 0.4 x/c 0.6 a）0\%跨度位置0.8图11：选定的帕累托解决方案和NAL设计之间的厚度分布的比较

488 OBAYASHIETAL。 图12：黏性和非黏性设计的平面形状与NAL设计的比较 0 0.2 0.4 x/c 0.6 0.8 I 图13：黏性和非黏性设计的厚度分布与NAL设计的比较-0.1 0.1, 0.2 0.3 0.4 0.5 • / /•!;: Ji'->\textbackslash\{\} • 黏性黏性0 0.2 0.4 0.6 0.8 1 x/c 图14：黏性和非黏性设计的Cp分布与NAL设计的比较

\clearpage
\chapter{25 信息科学--CFD的新前沿}
\section{25.1 从确定性系统进入复杂系统}
信息科学-CFD Koichi Oshima1和Yuko Oshima2 25.1的新前沿，从确定性系统进入复杂系统，自欧洲文艺复兴以来，在东方，科学的目的是找出“自然规则”并获得“结果”。 “因果关系”，即“因果之间的确定性关系”一直是现代科学的基本假设。 科学家只处理这种“确定性”现象。 作为“主体”，他们只观察“对象”，从不与他们的“对象”互动。 流体力学和CFD也属于“精确科学”的基本方法是：1. 科学家定义了对物体影响最大的参数。 2. 他们开发模型，通过已知的公式和/或对现象的实验观察，将自然描述为这些参数的函数。 3. 他们用数学方法解决这些模型方程，通常是用数值方法，并验证他们使用的解程序。 这是“验证”。 4. 他们通过实验或观察将解决方案与现实世界进行比较。 这是“验证”。 在现代科学的众多领域中，CFD是最成功的领域之一，例如，它已被用于飞机的实际设计。 可能是因为基本模型方程--Navier-Stokes方程--已经建立得很好，管理参数是1东芝先进系统公司2理光公司计算流体动力学前沿-2002编辑：David A. Caughey和Mohamed M. 哈菲兹 ©2002 世界科学

\section{25.2 计算机与人类大脑}
490 OSHIMA和OSHIMA定义明确。 也就是说，飞机设计问题涉及几何简单的形状和流动的物理属性，如密度、黏度等，是明确确定的。 湍流仍然是一个未解决的问题，因为它同时包含宏观和微观现象。 同样，自20世纪中叶以来，我们被迫面对复杂的现象，这些现象包含不同的时间和/或空间尺度。 这样的系统必然是“非决定性的”。 它们涉及我们对待人类社会的案例。 因此，治疗现实世界“复杂系统”的新科学出现了，社会科学与自然科学融合，复杂系统的基本特征是：1. 因果的相互干扰，观察到反馈效应，并且通常是积极的，“主体”和“客体”被认为属于同一类。 2. 微观系统和宏观系统之间的相互干扰。 微观现象中的混乱使「紧急进化」，并成长为宏观现象。 在这样的过程中，“分形字符”经常被识别。 3. 自我组织的性格。 随机运动经常成长为有组织的运动，从某种意义上说，这违背了热力学第二定律。 25.2 计算机与人类大脑 在20世纪50年代，在数字技术出现后，其在实验流体动力学中的应用开始，并随着其进步而不断增长。 如今，它已成为实验流体力学不可或缺的一部分。 数值计算来得有点晚，其增长完全取决于计算机的技术进步。 事实上，流体动力学家没有贡献其硬件，除了20世纪40-50年代的初期。 在整个计算中心、超级计算机和大型平行机器的时代，流体动力学家使用其机构的这些机器，并产生无尽的数据流，求解相同的方程--Navier-Stokes方程。 今天剩下的唯一问题是如何使用这些数据？ 我们根本没有存储空间来保存这些数据！ 补救措施在于流动可视化技术，它可以在三维空间上实时生成计算流动场的虚拟世界。 来自计算机的大量数据流直接实时使用，不需要保存在内存中。 所谓的计算数据的后处理现在是全自动的，必要的设计数据，例如，各种

\section{25.3 信息科学}
信息科学491部队数据直接发送到设计办公室。 在实验室实验和实地观察中，这种情况非常相似。 今天的风洞实验是全自动的，测量数据直接发送到设计办公室，设计软件在那里处理一切。 （几乎具有讽刺意味的是，那些不懂流体力学的人会设计飞机。） 在天气预报局，各种观测工具，从经典气象站到卫星数据，不断产生大量数据流。 它们被转移到中央车站，实时处理，并及时分发。 因此，天气预报的准确性最近得到了极大的提高。 作为实验工具，信息技术的作用远比计算工具更重要。 在这里，为了理解计算机的作用，我们将它们与人类的大脑进行比较。 计算机具有所谓的处理器基础架构，并产生“硬”输出。 数据被输入，冯·诺伊曼处理器进行计算，然后数据输出流出，在使用前必须存储在某个地方。 大脑是基于记忆的架构，并产生“软”答案。 数据输入用于已存储的关联内存的索引搜索。 然后输出是引用内存。 它的功能与所谓的面向对象数据库（OODB）有一些相似之处。 信息技术可能会持续发展到：1. 大规模并行系统，地球模拟器和TERA天文台的项目已经开始。 2. 非冯·诺伊曼计算机，如神经、量子计算机，将很快成为现实。 3. 使用PLD（可编程逻辑设备）的模拟/可视化技术正在进行中。 25.3 信息科学 现在，我们可以将这种基于信息科学和技术的新科学称为“信息”。 它的主要功能是：1. 全球观测站的数据采集，包括卫星和本地实验室，以及实时在线传输到处理器的数据。 2. 使用输入条件对这些数据进行索引，并进行处理，包括CFD计算，并将其存储为OODB。 3. 根据要求，调用这些关联数据库，并以虚拟现实的形式提供结果。 信息科学可以为社会提供1。 风险管理的直接建议，例如，乌苏火山爆发，

492 大岛和大岛地震预测，环境危机。 广泛分布的数据采集系统和实时数据处理将是管理此类灾害的关键因素。 2. 我们可以在不模拟微观现象的情况下解决复杂的流体动态流动。 相反，在真实站点的在线实时数据采集，结合CFD的宏观流量分析，将有助于直接预测流量。 3. 大型系统严重事故的风险控制，如核电站、火箭发射、飞机坠毁等，可以使用这些在线实时数据传输系统，结合计算机分析，如PSA（概率安全评估）、FTA（故障树分析）或FMPA（故障模式和效应分析）来有效管理。 最重要的是，信息技术现在正在部分地取代人类在生活的各个方面的聪明才智。 事实上，信息科学是CFD的新领域。

\clearpage
\chapter{26 将CFD整合到空气动力学教育中}
\section{26.1 介绍}
将CFD融入空气动力学教育Earll M. Murman1和Arthur Rizzi2 26.1 介绍 Robert MacCormack在计算流体动力学的职业生涯始于美国宇航局的艾姆斯研究中心，他在那里制作了1969年关于“MacCormack的方法”的开创性论文。 1981年，他转到学术界，将CFD引入航空航天工程课程和学术研究项目，先是在华盛顿大学，后来在斯坦福大学。 在1981年之前，只有少数学校，特别是爱荷华州立大学，拥有强大的航空CFD项目。 针对这一缺陷，美国宇航局资助了包括斯坦福大学在内的七所大学的“计算流体动力学研究生培训项目”（库特勒，[1]）。 到20世纪80年代中期，研究生CFD课程已经建立，第一波博士出现，许多人成为CFD的学术、行业和政府领导人。 Bob的研究、教学和个人领导风格是新一代CFDers的主要贡献者。 鉴于自1981年以来的许多变化，这次纪念Bob MacCormack的研讨会提供了一个反思航空航天工程CFD课程现状的机会。 CFD已经发展成为空气动力学分析和设计的支柱。 影响航空航天工程和CFD的外部因素发生了巨大变化。 本着研讨会主题的精神，我们挑战我们的同事大胆思考CFD在气功教育中的作用。 我们提出问题，提出想法，并就将CFD集成到1福特工程教授的未来方向提出假设，麻省理工学院航空航天系，马萨诸塞大道77号33-412室。 剑桥，马萨诸塞州 02139 美国。murman@mit.edu 2 皇家理工学院（KTH）航空系教授，SE-10044 瑞典斯德哥尔摩，rizzi@aero.kth.se 计算流体动力学前沿 - 2002 编辑：David A. Caughey和Mohamed M. 哈菲兹 ©2002 世界科学

\section{26.2 从1981年到2000年的变化}
494 MURMAN \& RIZZI空气动力学教育。 我们介绍了针对该主题要素的同事的非正式调查结果，并就我们自己的当前教学法提供了一些简短的评论。 虽然CFD适用于热和质量转移、生产工艺和其他应用可能会提出同样的问题，但我们的范围仅限于我们自己的专业领域。 许多变化也影响了研究方向和机会，但我们把这些主题留给研讨会上的其他论文。 26.2 1981年至2000年的变化 十九年是教育一代。 2000年秋天进入大学的学生出生在鲍勃开始在华盛顿大学学习的那一年。 CFD领域是一代人，信息技术经历了一场革命，航空航天业发生了重大变化。 在考虑教育主题之前，我们简要回顾了这些变化。 表1对比了1981年至2000年的一些CFD属性。 关于1981年的参考点，我们以ALAA第5届CFD会议记录[1]为参考。 对于2000年，我们使用目前的研讨会记录。 在这些中，我们添加我们自己的知识。 1981年，CFD作为空气动力学革命力量的潜力得到了认可（Lomax，Hall [1]）。 到2000年，已经实现了全部影响。 CFD已经从专家的活动扩散到实践空气动力学家的主要工具。 三维稳定或非稳定、非黏性尤拉或雷诺平均纳维尔·斯托克斯求解器使用结构化或非结构化网格进行实际飞行配置是常态。 事实上，2000年6月的《美国航空航天》杂志有来自三家此类能力的供应商的全页彩色广告[2]。 基于线性面板和涡流晶格模型、全耦合非黏性/黏性求解器以及2D和3D非线性非非黏性求解器的方法在桌面上。 扩展基本的CFD方法以包括其他物理学是很常见的。 只有模型以外的湍流流的直接解决方案是遥不可及的。 的确，空气动力学设计和分析发生了一场革命。 表2对比了这两个样本年的一些信息技术属性。 再一次，一场彻底的革命发生了。 1981年，IBM PC和Apple Macintosh还没有推出。 大型计算机通过硬接线“笨字”字母数字或CRT图形终端访问。 ARPA Net用于文件传输或远程用户访问有限时间共享超级计算机。 编程以FORTRAN为导向。 记忆是一种宝贵的资源。 当今信息技术丰富的环境是众所周知的，不需要详细说明。

CFD 495 2D/3D线性方法的集成，应用空气分析的支柱，2D/3D超音速全电位方法用于翼型和机翼体配置分析，基于逆迭代耦合的2D黏性/非黏性相互作用，欧拉和纳维尔·斯托克斯方法对3D配置的有限应用。 用于2D配置的不稳定时间精确的欧拉和Navier-Stokes抛物线Navier-Stokes用于3D高超音速分析“全局”湍流模型。 多网格方法新兴的有限有限元方法网格生成对于复体是一个重大障碍。 单块或块结构网格是规范。 网格自适应仅限于网格点重新分配，单色图形终端，具有“笔式绘图仪”风格输出，仅限于非空气动力学模拟和其他物理学的耦合，主要目标是在合理的时间内执行模拟，用于全3D配置CFD，由内部专家在桌面上执行2D/3D线性方法，桌面上的2D/3D跨声速全电位方法。 由桌面上的直接牛顿求解器用于逼真的设计优化 2D黏性/非黏性交互，Euler和Navier Stokes方法广泛用于逼真的复杂3D配置。 用于复杂的3D配置的不稳定时间精确的欧拉和纳维尔斯托克斯。 用于高超音速分析的全3D欧拉和Navier-Stokes“局部”湍流模型（墙函数）和一些LES大量使用多网格方法，鲁棒性是一个挑战有限元方法广泛使用的鲁棒网格生成方法，用于使用非结构化四边形/三角形或六边形/四面体网格加上各种过集和块状结构网格的复杂物体。 电网质量是个问题。 具有局部特征的单元划分的稳健全自适应网格，视觉丰富、交互式、颜色、以像素为导向的图形输出，经常与结构、控制和其他物理耦合，主要目标是在真实产品CFD的设计周期内以已知的精度、可靠性执行模拟，由新手执行并可从供应商处购买表1-1981年和1990年的代表性CFD属性，航空航天也经历了重大变化，如表3所示。 20世纪90年代初发表的一篇有趣的文章《飞机作为计算机外围设备》，以流行的方式捕捉到了飞机的设计、功能和操作随着

496 MURMAN \& RIZZI 主机- 终端架构 商业超级计算机和专业数字计算器（例如Illiac IV、阵列处理器） 有限 ARPA 网络访问 CRT“笔绘图仪”图形 RAM 一种珍贵的商品 FORTRAN 面向文字处理 客户端-服务器架构 带有并行处理器的商业笔记本电脑和台式机工作站 无限制互联网访问 交互式、3D 彩色图形工作站 RAM 丰富的C、面向对象、MATLAB、集成 笔记本电脑/桌面演示 表 2 - 1981年和1990年信息技术出现的代表性信息技术属性。 冷战后时代，工业和政府实验室经历了重大重组。 如图1所示，新军用飞机的起始量大幅下降。 信息经济及其所有“圆点玉米”已经出现，将投资者的期望设定在贬低航空航天业务的水平。 1999年底，美国十大航空航天公司3的累计市值约为戴尔计算机。 国家和全球联盟正在成为常态，例如与联合打击战斗机计划或星航空联盟的合作伙伴关系。 航空航天车辆的使命已经从移动质量（有效载荷）扩展到成为全球信息基础设施的节点。 里根冷战的建立更高、更快、更远的口号15家美国主要航空航天公司航天飞机刚刚飞行109个收入客运公里，年增长7.8\%，航空航天出口最大的部门“旧经济”占主导地位，很少有联盟或伙伴关系“最后”战斗机（JSF）在竞争中更好、更快、更便宜的口号9家美国主要航空航天公司空间站正在组装3x10收入客运公里5\%年增长370亿美元，特许权使用费和许可出口250亿美元飞机出口“新经济”主导国家、区域和全球联盟和伙伴关系表3-代表性航空航天工业 1981年和1990年的属性3个波音、霍尼韦尔、UTC、通用动力、Textron、洛克希德·马丁、雷神、TRW、诺斯罗普·格鲁曼和利顿工业。

\section{26.3 教育注意事项和问题}
CFD 497 C120 AOl I X266 AV88 F101 X15 „的集成。 F/A18 1950年代 1960年代 1970年代 1980年代 1990年代 2000年代 2010年代 2020年代 图1 - 按典型工程师的十年和职业生涯长度划分的新美国军用飞机计划。 埃尔南德斯[3]。 随着CFD、信息技术和航空航天领域的所有这些变化，人们是否应该期望大学课程能够保留其1981年的结构？ 我们当然不这么认为，其他人也不这么认为（例如 [4]，[5]，[6]）。 现实中，空气动力学在未来的航空航天课程中的作用会较小。 然而，学生需要在学科中达到扎实的工作水平。 目前空气动力学课程的大部分内容是在第二次世界大战和20世纪60年代中期之间建立的。 如图1所示，从那时起，它为开发许多新飞机的一代人提供了很好的服务。 但图1中的前景以及大型商用飞机的类似趋势表明，当今学生的未来会有所不同。 26.3 教育考虑因素和问题 工业趋势远离了工程科学创新所需的专业化，转向了产品开发创新所需的更全面的思想家，例如[5]，[7]。 技术能力是期望和奖励的。 但冷战时期以工程科学为基础的教育范式的狭隘专业化代表，比系统专业知识、小组技能和多功能能力更不具有价值。 无论我们愿否，对空气动力学家的需求有所减少，而对系统和信息工程师的需求有所增长。 尽管这仍然是飞机设计中必不可少的学科，

498 MURMAN \& RIZZI空气动力学不再是航空航天领域的驱动技术。 这些是空气动力学社区需要接受的因素，因为它们超出了我们的控制范围。 此外，随着进入大学的一代学生在信息丰富、以媒体为导向的世界中成长，学习风格也发生了变化。 今天的学生寻求并更好地应对实践体验式学习，而不是绝大多数大学课程的循序渐进的基础到应用方法。 通过课堂练习和课外项目让学生积极参与学习过程，比传统的讲座-听众模式更有效。 在许多方面，我们的教育系统反映了大规模生产模式，即将知识传递给学生，以防他们需要知识。 更适应当今学生的教学法可能是“及时”的方法，即根据学生需要应用知识来寻求和获取知识。 除了这些因素之外，信息技术对教育交付的影响才刚刚开始[6]、[8]、[9]。 桌面计算和基于网络的信息访问为新的材料交付提供了许多机会。 这种能力可以使教科书和黑板讲座等传统教学材料过时或显著改变。 有了随时可用的信息，讲师的角色可能会发生重大变化。 甚至有人预测，我们所知道的大学将过时。 无论我们是否同意这种愿景，我们都应避免被动思考，并有意识地考虑未来的机会。 有了这些框架的想法，我们提出了以下问题：随着航空航天工程师对其他技能的需求增加，对空气动力学专家的需求减少，我们应该如何调整本科和研究生教育的空气动力学课程？ 我们将如何减少一个典型的航空航天工程师学习的空气动力学科目的数量（本科生1-2个，研究生3-4个），同时为他们的职业生涯传授足够的理解水平？ 由于CFD是空气动力学家的主要分析和设计工具，我们是否在经典理论、实验和计算方法之间取得了最佳平衡，为学生未来的职业生涯做好准备？

\section{26.4 非正式调查结果}
CFD 499的整合，我们能否利用CFD方法来解决当今学生的体验式学习风格，在空气动力学方面提供扎实的教育？ 我们将这些问题留给读者考虑。 但我们认为，像20世纪80年代初一样，可能是时候再次暂停并考虑将CFD纳入航空航天课程了。 当时，需要将CFD纳入研究生教育，“通过增加CFD领域的合格专家人数来满足国家航空航天工作的关键人力需求”（Kutler [1]）。 我们推测，目前的CFD课程大部分是之前需求的遗产。 随着上述许多变化，我们推测今天的需求是不同的，并提供以下假设：“今天的空气动力学工程师需要精通现代CFD方法和工具，并且必须知道如何将它们与理论和实验结合使用，进行空气动力学分析和设计。 因此，我们看到越来越需要将CFD完全整合到主流空气动力学教学法中。” 26A 非正式调查的结果 我们认为通过非正式调查来抽样同事的意见会很有趣。 2000年6月，我们通过电子邮件向一些在美国和欧洲教授空气动力学的学术同事发送了一份调查问卷。 我们还向研讨会的与会者分发了调查问卷。 我们从美国、欧洲和亚洲的教师和非教师同事那里收到了40份回复。 我们强调，以下调查结果并不代表科学或统计学上有效的样本，也不一定以避免混淆因素的方式提出问题。 我们的第一个问题涉及目前CFD在空气动力学科目中的使用程度。 所要求的信息包括主题级别、马赫数制度、CFD方法、代码和使用的文本。 表4显示了回复的汇总摘要。 本科生很少使用CFD，代表了空气动力学的第一科。 CFD方法开始进入高年级本科生，在美国的使用比欧洲多。 典型的方法是2D面板、3D涡流晶格和边界层方法。 在研究生阶段，有大量使用CFD方法的报告。 欧洲教育工作者使用的方法比他们的美国同事更先进一些。 不出所料，2D方法主导了3D

500 MURMAN \& RIZZI 在下表中输入有关您在空气动力学科目中使用的CFD方法的信息，以教授基本理解和/或应用空气动力学原理。 学生Yr本科Is'和2年级本科3年级和4年级研究生入门级研究生高级水平美国有限使用亚音速2D面板、涡流晶格和边界层方法。 关于超音速方法的一个回应 大量使用亚音速二维面板和黏性方法。 一些涡流晶格方法的使用。 一些超音速/超音速二维方法的使用。 3D尤拉的大量使用2D尤拉和Navier-Stokes方法，以及一些2D全电位、面板、边界层和3D涡旋晶格方法的大量使用2D欧拉和Navier-Stokes方法，以及3D尤拉和Navier-Stokes方法。 一些使用2D全电位、面板、边界层和3D涡旋晶格方法欧洲没有使用限制使用2D面板方法。 有限使用2D和3D欧拉和雷诺平均纳维尔·斯托克斯方法，各种2D欧拉、Navier-Stokes、面板和边界层方法。 使用了一些3D面板和涡流晶格方法。 大量使用2D和3D Euler和Navier-Stokes方法，包括非稳态方法。 表4-CFD方法在主流空气动力学科目中的使用报告。方法，但后者有合理的表示。 除了表4中报告的信息外，受访者还表示使用了许多不同的CFD软件包，其中大部分是本地开发的，很少有在不止一所学校使用。 XFOIL和MSES是唯一真正的例外。 正在使用各种教科书和笔记。 更常被引用的是Anderson [10]、Anderson、Tannehill和Pletcher [11]以及Hirsch [12]。 第二个问题直接解决了我们的假设。 表5中的回复显示，将CFD纳入研究生水平的空气动力学科目的强烈一致性，而对本科科目的一致性较少-

CFD 501的集成，后者在美国比欧洲更强。 其他观点，包括将CFD与空气动力学分开教学，平均响应低于“3”。 欧洲人比美国反应者更强烈地认为，在学生研究、独立和设计项目中，应该在“需要知道”的基础上接近CFD。 与其他教育空气动力学工程师使用CFD工具和方法的教学方法相比，您对将CFD方法整合到主流空气动力学科目中有什么看法？ 1 = 强烈不同意 3 = 有点同意 5 = 强烈同意 观点 回复数量 CKD 方法应纳入研究生水平的主流空气动力学科目。 CFD方法应该融入本科水平的主流空气动力学科目。 CFD方法应在“需要知道”的基础上教授，作为学生研究或独立项目的一部分。 CFD方法应与主流空气动力学科目分开教授。 作为设计科目的一部分，CFD方法应在“需要知道”的基础上进行教学。 CFD方法最好在毕业后通过“在职”经验来学习。 CMJ方法最好通过在工业或政府的暑期实习来学习。 平均美国 4.3 3.7 2.4 2.9 2.0 1.9 1.9 平均欧尔 4.4 3.1 3.4 2.3 3.0 2.4 1.9 平均全部 4.4 3.5 2.8 2.7 2.3 2.1 2.0 STD 全部 0.8 1.2 1.3 1.3 1.1 1.2 1.1 表 5 - 关于CFD应该被整合到课程中的观点，按总体平均响应排序。 调查回复者就将CFD纳入空气动力学课程提出了一些意见。 这些总结揭示了多种观点，从CFD和空气动力学的强烈整合到CFD主题之前的更传统的基本空气动力学和数值分析主题的连续方法。 我们的下一个问题侧重于在综合方法中强调什么的困难选择。 表6报告了强烈的一致性，即课程应该深入涵盖基本的空气动力学理论以及对输出空气动力学参数的解释。 用户可控制的参数和底层数值方法的覆盖深度较小。 建议很少强调编程

502 MURMAN \& RIZZI实施。 调查响应者增加了评论，重点是数值方法的深度处理，深度和光覆盖之间的视图分布。 假设CFD方法被整合到主流空气动力学科目中，以下主题应该深入到什么程度，以便学生获得有效应用相关CFD工具或方法的必要理解？ 1 = 轻度覆盖 3 = 中等深度覆盖 5 = 深度覆盖主题 响应数量）空气动力学问题的理论表述（例如流体假设，控制方程，边界条件）输出解释（例如力和力矩系数、压力分布、流场特征）用户控制参数（例如，网格分辨率/质量、收敛程度、数值常数）数值解决方案的制定（例如 离散化方法，误差源，离散方程的解）基础数值方法（例如色散和耗散误差，矩阵解法， 稳定性，收敛）编程实现（例如编程语言、程序结构、内存利用率和执行速度、并行化）平均美国4.4 4.2 3.6 3.7 3.4 2.4平均欧尔4.6 4.5 4.2 4.0 3.6 2.5平均全部4.5 4.3 3.8 3.8 3.5 2.5 STD全部0.8 0.8 1.1 1.0 1.3 1.4表6-CFD集成到课程中的空气动力学科目内容的优先级，按整体平均响应排序，最后一个问题解决了将CFD集成到空气动力学课程中的障碍。 表7只显示学者的回复。 有趣的是，调查中提供的选择都没有被判断为“重大”障碍。 所有的选择都在“有点障碍”的范围内。 总的来说，教学问题的排名比技术因素高一些。 然而，鉴于标准差，这些差异的任何统计学意义都是值得怀疑的。 有趣的是，响应者添加的评论侧重于保持软件最新的困难、许可问题、教授和教学助理对软件的熟练程度以及例程或软件包的过时。

\section{26.5 集成示例}
CFD 503的整合 根据您的第一手知识，请评价以下可能的障碍对将CFD方法集成到空气动力学教学中的重要性。 1=不是重要的障碍；3=有点障碍；5=重要的障碍可能的障碍数或响应*！ 学生准备/预处理 技术支持（教学助理或员工） 修改课程，将CFD纳入空气动力学 了解最佳教学方法 CFD方法的教科书覆盖范围 可用的预处理软件 软件文档 示例问题和练习的可用性 CFD求解器的稳健性 可用的CFD应用程序包/程序 支持系主任或同等计算机速度/内存资源 您在空气动力学科目中教授CFD的经验 可用的后处理软件 计算机图形资源 测量学生成绩 平均3.5 3.3 3.4 2.9 2.9 2.8 2.8 2.4 2.5 2.4 2.4 2.1 2.1 2.5 2.1 1.9 Avg Eur 3.6 3.5 2.8 3.2 2.8 3.3 3.1 2.8 2.7 2.9 2.7 2.8 2.2 2.6 2.5 Avg All 3.6 3.3 3.2 3.0 2.9 2.9 2.7 2.7 2.6 2.5 2.4 2.3 2.3 2.1 STD All 1.2 1.2 1.4 1.1 1.1 1.2 1.0 1.4 1.2 1.4 1.0 1.4 1.2 1.0 1.0 表7 - 将CFD整合到空气动力学课程的障碍，按平均响应排序。 26.5 集成示例 作者一直在尝试将CFD整合到本科空气动力学科目中。 我们以简要叙述我们的经历来结束论文。 在过去的学年里，第一作者参与了MTT[13]教授初级-高级空气动力学科目的新方法。 该课程以洛克希德-马丁公司贡献的与F-16机翼体空气动力学建模相关的案例研究为中心。 基于预测相关飞机性能参数（范围、特定超功率、超音速冲刺时间、起飞距离）的能力，采用了综合理论、实验和CFD方法，并采取了“良好”措施。 CFD方法作为工具被引入，背景算法仅涵盖概念层面。 覆盖范围反映了表6中建议的优先事项，但时间压力导致对数值解决方案和方法的强调较少

504 MURMAN \& RIZZI比调查回复者推荐的要多。 引入的CFD方法有2D面板（线性变化涡流）、3D涡流晶格（AVL）、3D欧拉（非结构化网格，Runge-Kutta 时间积分，FELISA）、积分边界层方法和耦合面板/边界层方法（XFOIL）。 总体而言，学生们对这些CFD工具以及如何将它们与其他空气动力学分析方法相结合，取得了入门级的理解。 他们一直希望在CFD和理论方法上有更多的深度，但时间压力将这种短途旅行降到了更高级的科目。 就障碍而言，y似乎与表7中的障碍非常吻合。 我们没有遇到任何重大障碍，但可以肯定的是，表格中的前五名在某种程度上都得到了满足。 总体而言，我们对实验很满意，并计划在即将到来的学年继续完善它。 作为斯德哥尔摩皇家理工学院（KTH）航空学硕士学位要求的一部分，学生必须选修两门空气动力学课程，并可以选择第三门高级课程。 自1993年以来，第二作者一直在教授第二门必修课程（硕士课程的第四年）。 正如具有CFD背景的讲师所期望的那样，课程首次运行时的内容强调了高端方法，即全潜力和欧拉。 在课程结束时，一位以前在航空航天业工作过的学生评论说，给人的印象是，“当面对空气动力学问题时，我们应该立即跑到一台大电脑前”。 他的评论被放在心上，在随后的几年里，讲座和课程内容不断发展，强调对基础知识的理解以及使用一些实用的计算方法。 这与表6中列出的内容优先级相当一致。 无论如何，电脑都无法逃脱。 与风洞一起，它是空气动力学的基本工具，但它不需要很大。 相反，它应该适合放在桌面上，放在学生的指尖上，就像纸和铅笔一样，在解决问题的过程中毫不费力地使用。 计算机在空气动力学中的效用是巨大的，从数字数据通信等数字化半实证方法到欧拉和纳维埃-斯托克斯求解器。 此外，合适的软件越来越容易获得。 例如，公共领域航空软件（见http://www.pdas.com）提供了该范围的一系列程序，其中一些可能适合课程作业。 以目前的形式，KTH的空气动力学课程（请参阅http://www.flyg.kth.se/edu-dir/ugrad-dir/prog-dir/4E1211prog.html获取课程招股说明书）专注于使用2D耦合面板-集成-边界层理论和工业强度工具XFOIL和AVL（雅典娜涡流格子方法）中实现的3D涡旋格子方法对机翼和机翼的分析和设计。 此外，用TSP方法研究了跨声速（隐性）流动中的翼型。 远远超出了黑匣子的方法，在理解上花费了大量的时间和精力

\section{26.6 总结}
CFD 505的集成如何以及这些工具产生什么结果。 事实上，该课程包含表6中的前五个项目，重点是项目1，然后下降到项目5。 然而，几年的课堂经验揭示了实现这一目标的一个重大障碍，即软件文档。 特别是XFOIL是一个非常复杂的软件，人们根本无法教授它如何组合的所有细节。 因此，KTH的MATLAB正在开发教学辅助工具，作为理解XFOIL和AVL等代码过程中的垫脚石。 两个是PABLO（面板/边界层）和TORNADO（涡流晶格），或多或少地逐字逐句地从其教科书描述中编程。 迄今为止的经验表明，这种垫脚石方法是有帮助的。 26.6 摘要 在Bob MacCormack在学术界的这段时间里，新一代学生诞生了，并将于今年秋天进入大学。 在其有生之年，CFD已经从“即将到来的革命”演变为空气动力学从业者的标准分析和设计工具。 除此之外，自1981年以来，信息技术、航空航天业的现状和学生学习风格也发生了许多变化。 作为教育工作者，我们认为这些因素要求我们对空气动力学教学方式和具体课程内容进行重大改变。 因此，我们提出了一个假设，即CFD（在该术语的最广义上）应该完全融入主流空气动力学科目。 我们对美国、欧洲和本次研讨会的同事进行了非正式调查，以提供与这一假设相关的观点的横截面，并可以辨别出一些趋势。 CFD正在开始被整合到更高层次的本科和研究生阶段的传统空气动力学课程中，本科课程采用更简单的方法进入研究生课程的更高级方法。 教授这些方法的重点主要放在基本理解上，而不是实际编程实施上。 虽然将这些方法纳入课程确实存在一些障碍，但没有一个被认为是重要的。 我们还报告了我们自己最近将CFD整合到本科空气动力学科目中的实验，以及我们的个人经验与同事的观点相比如何。 我们看到，随着航空航天课程的扩大，以满足下一代学生的教育需求，我们将面临重新思考学科内容及其授课的挑战。 对我们来说，很明显，传统的空气动力学课程与现在和未来的要求之间存在重大差距。 在我们看来，真正的挑战是“如何”将CFD完全整合到主流空气动力学课程中。 有鉴于此，我们强烈建议教育工作者认真思考

\section{26.7 致谢}
\section{参考资料}
506 MURMAN \& RIZZI对空气动力学课程的结构进行实验，以尝试新的方法。 此外，我们建议通过印刷出版物和互联网广泛分享这些实验的结果。 26.7 致谢 我们想感谢许多花时间回复我们的调查的人，他们太多了，无法说出名字。 参考资料1。 AIAA 计算流体动力学会议，技术论文集，加利福尼亚州帕洛阿尔托，1981年6月22日至23日2。 美国航空航天局，AIAA，2000年6月3。 埃尔南德斯，C。 《知识产权白皮书》，加州工程基金会，1999年12月7日。 此外，Drezner、J、Smith、G、Horgan、L、Rogers、C和Schmidt、R。 “保持未来的军用飞机设计能力”，兰德报告R-4199F，1992年4。 Crawley, E.F, Greitzer, E.M., Widnall, S.E., Hall, S.R., McManus, H.L., Hansman, R.J., Shea, J.F.,Landahl,M., 「麻省理工学院航空航天课程改革」，AIAA Paper 93-0325，1993年1月 5. McMasters，J.H.和Lang，J.D.，“加强工程和制造教育：行业需求，行业角色”，1995年ASEE年会和博览会，加利福尼亚州阿纳海姆，1995年6月25日至28日6。 Mason，W.H.和Davenport，W.J.，“应用空气动力学教育：发展和机遇”，AIAA论文98-2791，1998年6月7。 《回到未来-未来工程师》，国际飞行，2000年1月1日，第131-133 8页。 考伊，D。 A.和Liggett，J. 《流体力学入门计算机教科书》，计算流体动力学前沿1998，世界科学，1998，第465-481页9。 Kroo，Ilan，应用空气动力学-数字教科书，桌面航空学，1997年。http://www.desktopaero.com 10。 安德森，J.D. Jr.，空气动力学基础，第二版，McGraw Hill，11。 Tannehill，J.C.，Anderson，D.A.和Pletcher R.H.，计算流体动力学和传热，第二版，Taylor \& Francis，华盛顿特区，1997年。 12. Hirsch，C，内部和外部流动的数值计算：第1卷和第2卷。 John Wiley \& Sons有限公司 1988年。 13. Darmofal，D.，Murman，E.，Love，M.，“重新工程空气动力学教育”，AIAA论文2001-0870，AIAA 39*1航空航天科学会议，内华达州里诺，2001年1月8日至11日

2002年计算流体动力学的前沿 这一系列关于“计算流体动力学的前沿”的书籍是为了表彰对该领域产生重大影响的贡献者。 第一卷于1994年出版，献给安东尼·詹姆森教授；第二卷于1998年出版，献给厄尔·默尔曼教授。 这卷书献给Robert MacCormack教授。 本卷的26章是由来自学术界、政府实验室和工业界的领先研究人员撰写的。 他们介绍了流体流动问题数值分析技术的最新发展的最新描述，以及这些技术在工业中重要问题中的应用，以及向世界介绍“MacCormack方案”的经典论文。www. worldscientific. com 4887 he ISBN 981-02-4849-0

\backmatter
\chapter*{翻译与版式说明}
全文按原书章节顺序生成，并采用连续的纯中文书籍版式。
\end{document}
